\documentclass[11pt,DIV=11,parskip=false]{scrartcl}
\usepackage{testimony}
\title{Sola fide}
\author{\credit{TestimonyArguments.SolaFide}}
\date{Generated by \texttt{lake exe argtex}}
\begin{document}
\maketitle
\noindent This is one argument of the Testimony library, printed on its own. The library prints every argument together in \texttt{arguments.pdf}. This one is set as its Lean source sets it: the module docstrings as prose, in the order they are written, and under each one the package, line of reason, countermodel or result it introduces. Atoms are numbered propositional variables, numbered once per argument, so a variable means the same claim in every position within it; the legend gives each one's claim, its classification, the tradition that holds it and the source that carries it.

\noindent Nothing here asserts that the premises are true --- only that the conclusion does or does not follow from them.

\tableofcontents
\clearpage
\section*{Sola fide}
\addcontentsline{toc}{section}{Sola fide}

\subsection*{What sola fide claims}
\addcontentsline{toc}{subsection}{What sola fide claims}
\begin{quote}
“By grace you have been saved through faith. And this is not your own doing; it is the gift of God, not a result of works” (Ephesians 2:8–9).

\end{quote}
The Reformers summed up verses like this one in two Latin words, \emph{sola fide}: by faith alone. Rome answered at the Council of Trent, in its decree on justification of 1547, that faith is necessary and is not alone. The two sides have argued over the same verses ever since, and they still do. These chapters set the argument out so that a reader with an English Bible can follow every step of it, and see exactly where the two sides part.

\subsubsection*{What these chapters do, and what they cannot do}
Each position is written down as a list of premises, the things it takes as given, and every premise has a source: the verse, the council, the commentator who holds it. Nothing is assumed silently. Then a machine checks what follows from each list, and what does not. The checking is done in Lean, a program that verifies proofs; the Lean is kept behind each chapter, under “How this is checked”, for anyone who wants it.

What the checking cannot do is tell you which premises are true. Whether Paul meant one thing or another by a phrase is a question for the reader, and for the scholars cited. The checking adds something narrower, and harder to come by in a debate this old: if you grant these premises, this conclusion follows; if you deny that one, here is what you lose and what you keep.

Scripture is quoted from the English Standard Version (ESV). Where a Greek or Hebrew word matters, it is given in brackets after the English. Three words recur throughout:

\begin{itemize}
  \item A \textbf{premise} is something a position takes as given: a verse, or a claim about what a word in it means.
  \item A \textbf{reading} is a premise about what a text means where readers disagree.
  \item A \textbf{step} is an inference from some premises, its \emph{starting points}, to a conclusion. It is rated like a premise.
  \item A \textbf{rating} says how well a premise or a step is supported among the sources cited. It is \texttt{disputed} when a cited source grants the evidence and still denies the claim; above that come \texttt{plausible}, \texttt{well\allowbreak{}Supported} and \texttt{consensus}.
\end{itemize}
Names set like \texttt{this} are the names of premises and results in the source, given in brackets so that a reader can find them.

\subsubsection*{Three claims in one}
“By grace, through faith, not by works” is really three claims: salvation is \emph{by grace}, it is \emph{not by works}, and it is \emph{through faith}. The chapters keep them apart, because the positions disagree about different parts.

\begin{itemize}
  \item \textbf{Trent} denies the second, in one respect. It grants that nothing before justification, “whether faith or works”, merits it (Session VI, ch. 8). But it holds that works done afterwards, in grace, merit an increase of justification (ch. 10; canon 24). So for Trent salvation as a whole is not apart from works.
  \item \textbf{The apocalyptic reading of Paul} (Martyn, Campbell), so called because it reads Paul’s gospel as God breaking into the world to rescue it, does not reach the third. It takes “the righteousness of God” as God’s act of rescue, not as a status given on condition of faith. So it has no ground of its own for “through faith”, though it reaches grace and “not by works” by its own route (\texttt{apocalyptic\_\allowbreak{}establishes\_\allowbreak{}grace\_\allowbreak{}not\_\allowbreak{}works}).
\end{itemize}
Splitting the claim also shows what the hard questions are for. Grace and “not by works” follow from the verses that say so in as many words, once James is answered, with every one of those readings taken away (\texttt{grace\_\allowbreak{}and\_\allowbreak{}works\_\allowbreak{}rest\_\allowbreak{}on\_\allowbreak{}no\_\allowbreak{}lexical\_\allowbreak{}premise}). Those readings decide one thing only: \emph{through faith}.

\subsubsection*{Four witnesses}
The Reformed case reaches “through faith” by four routes, called \emph{strands}:

\begin{itemize}
  \item \textbf{Paul}, answering the teachers in Galatia: “a person is not justified by works of the law but through faith in Jesus Christ” (Galatians 2:16).
  \item \textbf{Jesus, in Luke}, to the woman who anointed his feet: “your faith has saved you; go in peace” (Luke 7:50).
  \item \textbf{Peter, in Acts}, at the Jerusalem council: “we believe that we will be saved through the grace of the Lord Jesus, just as they will” (Acts 15:11).
  \item \textbf{Jesus, in John}: “this is the work of God, that you believe in him whom he has sent” (John 6:29).
\end{itemize}
Each rests on its own texts and on a reading of them that the conclusion needs. For Paul, Peter and John that reading is one a scholar cited here denies; for Luke it is that “saved” at 7:50 means saved, which no scholar found denies. The case does not hang on any one of those readings: take any one away and the others still reach the conclusion. Only when such a reading is taken from every strand at once does sola fide fail to follow (\texttt{lexical\_\allowbreak{}premises\_\allowbreak{}jointly\_\allowbreak{}load\_\allowbreak{}bearing}). The next chapter takes the four in turn, and John has a chapter of his own.

\subsubsection*{What justification is}
Underneath the argument about faith lies a second question, which these chapters keep separate: what justification \emph{is}. Three answers are heard.

\begin{itemize}
  \item \textbf{The Reformed answer} is that justification is a verdict only (the technical word is “forensic”, as in a court): God pardons, and counts righteousness to the believer, and that is not a change worked inside the person (Westminster Confession XI.1).
  \item \textbf{Trent’s answer} is that justification is “not remission of sins merely, but also the sanctification and renewal of the inward man”.
  \item \textbf{The Finnish reading of Luther} (Mannermaa and his school) is that in faith Christ himself is present in the believer, so justification is not a verdict \emph{and nothing more}.
\end{itemize}
The Reformed case with the Finnish denial added, that justification is not only a verdict, still establishes sola fide (\texttt{forensic\_\allowbreak{}justification\_\allowbreak{}not\_\allowbreak{}load\_\allowbreak{}bearing}). Trent denies it and does not. So denying the verdict-only account is not, by itself, what separates Trent from the Reformers. What Trent’s objection needs is its fuller definition: justification includes the inner renewal, that renewal grows through good works, and so works increase justification. Both sides grant the growth. The Reformed call it sanctification, distinct from justification and following it (Westminster XIII.1; Calvin, \emph{Institutes} III.xi.6). Put that distinction in place of Trent’s definition, and Trent’s case no longer denies “not by works” (\texttt{trent\_\allowbreak{}objection\_\allowbreak{}rests\_\allowbreak{}on\_\allowbreak{}its\_\allowbreak{}definition}). Whether works merit anything in justification turns on whether justification includes the inner renewal.

\subsubsection*{James 2:24}
\begin{quote}
“You see that a person is justified by works and not by faith alone” (James 2:24).

\end{quote}
This is the only place in the New Testament where the words \emph{faith alone} occur, and they are denied. The chapters do not assume the verse away. They derive James’s agreement with Paul from two cited premises. The first is that James is attacking a barren faith, mere assent (“even the demons believe”, 2:19). The second is that works are the fruit and evidence of saving faith, and not a ground of justification. Moo and Johnson are cited for the reading of James, and the Westminster Confession for the doctrine:

\begin{quote}
faith “is not alone in the person justified, but is ever accompanied with all other saving graces, and is no dead faith, but worketh by love” (XI.2).

\end{quote}
The case depends on this answer. Take it away and sola fide does not follow (\texttt{james\_\allowbreak{}harmonisation\_\allowbreak{}is\_\allowbreak{}load\_\allowbreak{}bearing}). Grace alone does not need it (\texttt{grace\_\allowbreak{}needs\_\allowbreak{}no\_\allowbreak{}answer\_\allowbreak{}to\_\allowbreak{}james}), because James bears on “not by works”, the part he appears to contradict.

\subsubsection*{The positions weighed against each other}
So far the question has been what follows from each position’s own premises. The chapter “Weighing the positions” asks a different one: when the positions are set against one another, which survive? It uses three words in a sense of its own.

\begin{itemize}
  \item One position \textbf{defeats} another when it contradicts it and rests on premises rated at least as well.
  \item A position is \textbf{forced} when it must be accepted however the open disagreements are settled.
  \item Two positions that defeat each other are a \textbf{tie}. A reader may hold either, and neither is forced.
\end{itemize}
Seventeen positions are heard: the four strands, each argued alone; Trent; the apocalyptic reading; Sanders and his critics; Jervell; Aquinas on John, and Luther’s answer; Chrysostom on John 3:36, and the Reformed answer; Paul’s gospel in Galatians; and three positions on what Paul’s word “justify” means.

One thing prevails outright: Paul’s verb “justify” (δικαιόω) does not mean the renewal of the inward man (\texttt{only\_\allowbreak{}pauls\_\allowbreak{}word\_\allowbreak{}prevails}). It is argued two ways, from the word’s usage and from Romans 4, and those are the only positions whose every step is rated above \texttt{disputed}. Trent’s definition, taken as a claim about what Paul’s word means, cannot be defended against them (\texttt{trent\_\allowbreak{}on\_\allowbreak{}pauls\_\allowbreak{}word\_\allowbreak{}indefensible}).

Everything else ties. Every other position rests on something rated \texttt{disputed}, and every attack between Trent and a Reformed strand runs in both directions. Luke’s case, Paul’s gospel and Trent can each be defended, and none is forced. John’s case ties with Trent, the apocalyptic reading, Aquinas and Chrysostom, as the other strands tie with their rivals.

\subsubsection*{The questions the later chapters take}
Each later chapter takes one question a reader brings to this debate.

\begin{itemize}
  \item \textbf{Where exactly do Trent and Paul part?} At one step: that a verdict on Christ’s finished work excludes the renewal worked in us. Trent can grant every text, and the verdict sense of the verb (“Where Trent parts from Paul”).
  \item \textbf{Is Trent’s definition about Paul’s word, or about what God does?} Taken as a claim about the word, it cannot be held with the evidence for the word’s meaning. Taken as a claim about what God does, nothing about the word reaches it, and it meets Paul at steps rated \texttt{disputed}: Augustine accepts Romans 4 and reads it as God making the ungodly godly (“What Trent means by justification”).
  \item \textbf{What did the Reformation itself decide? And without Trent, would Scripture establish sola fide?} Heard as Rome and the Reformers alone, only the meaning of Paul’s word is forced. Heard without Trent, two things are forced: Paul’s gospel, and the meaning of Paul’s word. \emph{Faith alone} still is not, and neither is Luke’s “faith suffices”. Campbell, Sanders and Dunn, Jervell, and Aquinas each still stand against one strand, and Chrysostom against all four (“Weighing the positions”).
  \item \textbf{What does Jesus himself say?} The tax collector “went down to his house justified” (Luke 18:14). Trent’s case is consistent with every word of the parable, and meets it where it meets Luke 7:50. Matthew 7:21–23 and 19:16–22 are each read both ways, and neither settles its question alone (“What Jesus’ own words say about faith and works”).
  \item \textbf{Is baptism necessary?} Trent makes water baptism the instrument of justification, “or the desire thereof”. Water, the Spirit and “fire” are kept apart, so that no argument slides from one sense of “baptism” to another, and John 3:5 is examined under each of the four readings of its “water” (“Is baptism how we are justified?”).
  \item \textbf{Does justification have degrees?} A verdict acquits or condemns, and has no degrees, so if justification is only a verdict, Trent’s “increase of justification” has nothing to increase. Trent can grant the courtroom texts; the two part only at the definition (“Does justification have degrees?”).
  \item \textbf{What does “justified by works” mean in James?} The same kind of hearing, over James 2:24 and the two other texts Trent cites for the increase of justification (“James 2:24” and “Trent’s two Latin texts”).
  \item \textbf{What does Genesis say about Abraham?} “He believed the LORD, and he counted it to him as righteousness” (Genesis 15:6), and then God alone passes between the pieces of the covenant. Read so, Abraham was counted righteous with no work as a condition; read from his testing in Genesis 22, as 1 Maccabees reads it, he was not. The texts of Genesis alone do not decide (“Abraham, counted righteous”).
  \item \textbf{How firm are the ratings?} Each rating a verdict depends on is recomputed under three standards (“Why each step is rated as it is”).
\end{itemize}
The claims themselves, each with its number, its source and its rating, are listed in the appendix, “The claims and their sources”.

\subsection*{Four witnesses: Paul, Luke, Acts and John}
\addcontentsline{toc}{subsection}{Four witnesses: Paul, Luke, Acts and John}
The first chapter named the four strands. Each begins from a different speaker: Paul, Jesus as Luke reports him, Peter at Jerusalem, and Jesus as John reports him. Each rests on its own texts, and on a reading of them that the conclusion needs. This chapter takes them in turn, then asks what they show together.

\subsubsection*{Paul}
The first is Paul’s answer to the teachers in Galatia who required believers to be circumcised as well as to believe:

\begin{quote}
“A person is not justified by works of the law but through faith in Jesus Christ” (Galatians 2:16).

\end{quote}
Two questions about this verse carry the strand. The first is what “works of the law” (ἔργα νόμου) means. The traditional answer, Calvin’s and among modern scholars Westerholm’s, is human works in general. Dunn and Wright, building on Sanders, answer that it means what the law required of Israel as a people set apart, seen above all in the practices that marked Jews off from gentiles: circumcision, food laws and the sabbath. Their reading is known as the New Perspective on Paul. Only on the first answer does refusing circumcision mean refusing every work as a requirement (\texttt{works\allowbreak{}Of\allowbreak{}Law\allowbreak{}Means\allowbreak{}Works\allowbreak{}Generally}).

The second question is what “faith in Jesus Christ” (πίστις Χριστοῦ) means: our faith in Christ, or Christ’s own faithfulness, as Hays holds. Only on the first does the verse name the believer’s faith as the means (\texttt{pistis\allowbreak{}Christou\allowbreak{}Objective}). This is a premise inside Paul’s strand rather than a strand of its own, because every verse with the phrase sets it against the law, so no route through it avoids the first question. Nor does it divide the same people: Dunn rejects the traditional reading of “works of the law” and keeps “faith in Christ”.

\subsubsection*{Jesus, in Luke}
The second strand is a sentence Jesus spoke to the woman who anointed his feet:

\begin{quote}
“Your faith has saved you; go in peace” (Luke 7:50).

\end{quote}
Luke has the saying four times (7:50, 8:48, 17:19, 18:42). The Greek verb (σῴζω) means both \emph{save} and \emph{heal}, and in the other three places the person has just been healed, which is why most English Bibles print “made you well” there. So the strand rests on 7:50 alone, where the saying follows “your sins are forgiven” and no one is ill. That the word means \emph{saved} there is a reading the strand needs (\texttt{sozo\allowbreak{}Is\allowbreak{}Soteriological}), and it is rated \texttt{well\allowbreak{}Supported}, because no scholar found argues the healing sense at 7:50. As Marshall puts it, the saying is used at 7:50 “in a case where physical healing is not involved”.

The strand’s disputed reading is of the verse before:

\begin{quote}
“Her sins, which are many, are forgiven—for she loved much” (Luke 7:47).

\end{quote}
Read one way, love is the reason she was forgiven. The strand reads it the other way: her love shows her forgiveness and does not earn it (\texttt{luke7\_\allowbreak{}47\allowbreak{}Love\allowbreak{}Is\allowbreak{}Evidence}). The rest of the verse (“he who is forgiven little, loves little”) points that way, and so does the parable of the two debtors just before it (7:41–43). Three readers are cited for it:

\begin{itemize}
  \item Osvaldo Padilla: her actions “stem from the fact that she loves much because she has been forgiven much”.
  \item Melanchthon, in the \emph{Apology of the Augsburg Confession}: “Christ interprets Himself when He adds: Thy faith hath saved thee”.
  \item Marshall, in his commentary on the Greek text, from the grammar: “for” here means “as is evidenced by the fact that”.
\end{itemize}
Marshall also names those who read the love as the cause, most Roman Catholic commentators among them, and that is why this premise is rated \texttt{disputed}.

\subsubsection*{Peter, in Acts}
The third strand comes from the Jerusalem council, where some demanded:

\begin{quote}
“Unless you are circumcised according to the custom of Moses, you cannot be saved” (Acts 15:1).

\end{quote}
Peter answered that God had cleansed the gentiles’ hearts by faith, that the law is “a yoke… that neither our fathers nor we have been able to bear”, and that “we believe that we will be saved through the grace of the Lord Jesus, just as they will” (15:9–11). The strand’s disputed reading is that the yoke is the whole law treated as a condition of salvation (\texttt{acts15\allowbreak{}Yoke\allowbreak{}Is\allowbreak{}Law\allowbreak{}As\allowbreak{}Condition}). Jervell reads it instead as Israel’s law laid on gentiles as a mark of belonging. On his reading, Peter refuses to make gentiles keep Israel’s law, and says nothing about the law as a condition of salvation.

\subsubsection*{Jesus, in John}
The fourth strand is an exchange in John 6:

\begin{quote}
“What must we do, to be doing the works of God?” Jesus answered them, “This is the work of God, that you believe in him whom he has sent” (John 6:28–29).

\end{quote}
Its disputed reading is that this believing is trust, which brings nothing of its own to God (\texttt{johannine\allowbreak{}Believing\allowbreak{}Is\allowbreak{}Trust}). Calvin reads it so. Aquinas reads it as faith formed by charity, that is, faith made alive by love, which is how Trent describes the faith that justifies. Luther answers Aquinas from Galatians 3:11–12, and his answer rests on a step Aquinas denies. The chapter on John sets this out, and takes up John 3:36 as well, “whoever does not obey the Son shall not see life”, with Chrysostom’s reading of it.

\subsubsection*{What the four show together}
The case does not hang on any one of these readings. Take any one away and the other strands still reach the conclusion (\texttt{works\allowbreak{}Of\allowbreak{}Law\_\allowbreak{}not\_\allowbreak{}load\_\allowbreak{}bearing}, \texttt{sozo\_\allowbreak{}not\_\allowbreak{}load\_\allowbreak{}bearing}, \texttt{acts15\allowbreak{}Yoke\_\allowbreak{}not\_\allowbreak{}load\_\allowbreak{}bearing}, \texttt{believing\_\allowbreak{}as\_\allowbreak{}trust\_\allowbreak{}not\_\allowbreak{}load\_\allowbreak{}bearing}). Acts and John each reach it with every other strand’s reading left out (\texttt{apostolic\_\allowbreak{}strand\_\allowbreak{}survives\_\allowbreak{}paul\_\allowbreak{}and\_\allowbreak{}luke}, \texttt{johannine\_\allowbreak{}strand\_\allowbreak{}survives\_\allowbreak{}the\_\allowbreak{}other\_\allowbreak{}three}). Only when such a reading is taken from every strand at once does sola fide fail to follow (\texttt{lexical\_\allowbreak{}premises\_\allowbreak{}jointly\_\allowbreak{}load\_\allowbreak{}bearing}).

So a reader who wins the argument over “works of the law” has not thereby defeated sola fide. The New Perspective rejects the traditional reading of that phrase and still holds justification by faith, and its premises establish the conclusion too (\texttt{new\allowbreak{}Perspective\_\allowbreak{}establishes}). Winning “the faithfulness of Christ” for Hays does not defeat it either: that costs the Reformed case its route through Paul, not its conclusion (\texttt{pistis\allowbreak{}Christou\_\allowbreak{}not\_\allowbreak{}load\_\allowbreak{}bearing}).

None of this shows that the strands are independent in practice. A reader who takes “works of the law” as Israel’s boundary markers will usually read Peter’s yoke the same way. But they are different premises about different texts, and an opponent has to answer each one.

\subsubsection*{Each witness is credited only with what it says}
On the Reformed reading, Paul argues for \emph{alone}. The Galatians already believed, and adding circumcision to their faith is what he says severs them from Christ (5:4). John’s strand, as Calvin reads it, claims the same: asked about works, Jesus names believing and nothing else.

Luke 7:50 and Acts 15 say less. Her faith saved her; the gentiles are saved without the yoke. That is, faith \emph{suffices} (\texttt{faith\allowbreak{}Is\allowbreak{}Sufficient}). Neither says that nothing else could be a condition, and “through faith” needs only that much. The difference matters when positions are weighed against each other, because a case is attacked on what it claims. Credit Luke with \emph{alone}, and the apocalyptic reading would attack Luke for a claim Luke never made.

\subsubsection*{Why the New Perspective reads Paul as it does}
In 1977 E. P. Sanders argued that Judaism in Paul’s day taught that one enters the covenant by grace and stays in by obedience. He called it “covenantal nomism”. If he is right, Paul’s opponents were not trying to earn salvation, and Dunn infers that what Paul refuses is the boundary markers (\texttt{sanders\allowbreak{}Line}). Gathercole, and the \emph{Variegated Nomism} volume edited by Carson, O’Brien and Seifrid, deny Sanders’ starting point: Jewish texts of the period also speak of final vindication according to works, so the demand Paul refuses is a demand for obedience (\texttt{critics\allowbreak{}Line}).

On the critics’ history, Paul’s strand by itself reaches sola fide (\texttt{critics\_\allowbreak{}carry\_\allowbreak{}the\_\allowbreak{}pauline\_\allowbreak{}strand}). On Sanders’, he does not (\texttt{sanders\_\allowbreak{}costs\_\allowbreak{}the\_\allowbreak{}pauline\_\allowbreak{}strand}). Whether Paul’s strand reaches sola fide depends on what Judaism in Paul’s day taught about works.

\subsubsection*{Who wrote the letters}
Ephesians, Titus and 1 and 2 Peter are each disputed as to author, and each dispute is recorded from both sides: Hoehner and Lincoln on Ephesians; Mounce, Dibelius–Conzelmann and Marshall on the Pastorals (1 and 2 Timothy and Titus); Jobes and Achtemeier on 1 Peter; Schreiner and Bauckham on 2 Peter; Ehrman across all four. Granting every critical conclusion costs the case nothing (\texttt{authorship\_\allowbreak{}not\_\allowbreak{}load\_\allowbreak{}bearing}), because it reads these letters as Scripture, not as the testimony of a named apostle. It does cost them their use as evidence of what \emph{Paul} meant by “works”, and that is why that premise is argued from Romans 4, 9 and 11, where the authorship is not in doubt.

\subsubsection*{The positions, written out}
Besides the Reformed case, these positions are written out in this chapter:

\begin{itemize}
  \item \textbf{the New Perspective} (Dunn, Wright), which reads “works of the law” as Israel’s boundary markers;
  \item \textbf{Hays}, who reads “faith of Christ” as Christ’s own faithfulness;
  \item \textbf{Jervell}, who reads Peter’s yoke as Israel’s law laid on gentiles;
  \item \textbf{the Finnish reading of Luther} (Mannermaa), for which justification is not a verdict and nothing more;
  \item \textbf{the critical view of authorship}, which denies that Paul wrote Ephesians and Titus, and that Peter wrote 1 and 2 Peter;
  \item \textbf{the apocalyptic reading} (Martyn, Campbell);
  \item and \textbf{Trent}.
\end{itemize}
Beside them are variants of the Reformed case. Each asks whether the case depends on one premise: with it taken out, either the conclusion still follows, or the checking names a way of taking the texts on which it does not.

\subsubsection*{What follows, and what does not}
Each result takes one position and asks whether its conclusion follows from its premises. Where it does, the result says so. Where it does not, the result names a consistent way of taking the texts on which every premise holds and the conclusion fails. That way of taking them is the rival position, written down, and a reader can check it verse by verse.

Many results ask what the case depends on. A premise is taken out, everything else is kept, and the result says whether the conclusion still follows. If it does, the case never needed that premise.

John’s strand has more to answer than the others, and the next chapter takes it.

\subsubsection*{How this is checked}
\paragraph*{Lines of reason}
\leavevmode\par\nobreak
\par\medskip\noindent\textbf{\texttt{pauline\allowbreak{}To\allowbreak{}Faith\allowbreak{}Alone}}\par\nobreak\noindent
\textbf{Paul’s strand.} From Romans and Galatians to justification by faith alone, with Galatians 2:16 at the centre: “a person is not justified by works of the law but through faith in Jesus Christ.”

Each premise does its own work. Galatians 2:16 excludes works of the law. Read as “faith \emph{in} Christ”, it also names the believer’s faith as the means. The letter’s occasion says what is being refused: circumcision \emph{added to} faith as a requirement. That is where \emph{alone} comes from. The Galatians already believed, and Paul says that adding a work to their faith severs them from Christ (5:4). Only the reading of “works of the law” as works in general lets circumcision stand for every work, and not for one mark of being a Jew.

\begin{derivation}
  \item[] \((\pvar{TestimonyArguments.SolaFide}{2} \land \pvar{TestimonyArguments.SolaFide}{3} \land \pvar{TestimonyArguments.SolaFide}{8} \land \pvar{TestimonyArguments.SolaFide}{7} \land \pvar{TestimonyArguments.SolaFide}{6}) \rightarrow \pvar{TestimonyArguments.SolaFide}{25}\)\gloss{\emph{would give:} Justification is by faith alone}
\end{derivation}
\par\medskip\noindent\textbf{\texttt{dominical\allowbreak{}To\allowbreak{}Sufficiency}}\par\nobreak\noindent
\textbf{Luke’s strand.} From Jesus’ words at Luke 7:50, “your faith has saved you”, to the claim that faith suffices. It needs two readings. The word “saved” (σῴζω) means saved there, not healed. And “for she loved much” (7:47) gives the evidence of her forgiveness, not its cause, so that what Jesus names as saving her is her faith.

The strand says that her faith saved her, and no more. It does not say that nothing else could have. It owes nothing to Paul or to the dispute over “works of the law”.

\begin{derivation}
  \item[] \((\pvar{TestimonyArguments.SolaFide}{17} \land \pvar{TestimonyArguments.SolaFide}{18} \land \pvar{TestimonyArguments.SolaFide}{19}) \rightarrow \pvar{TestimonyArguments.SolaFide}{26}\)\gloss{\emph{would give:} Faith is sufficient: whoever believes is saved}
\end{derivation}
\par\medskip\noindent\textbf{\texttt{apostolic\allowbreak{}To\allowbreak{}Sufficiency}}\par\nobreak\noindent
\textbf{Peter’s strand.} From Peter’s speech at the Jerusalem council to the claim that faith suffices.

The council answers the demand Galatians answers: circumcision and the law of Moses as a condition of salvation (Acts 15:1, 15:5). Peter replies that God cleansed the gentiles’ hearts by faith, that Jew and gentile alike are saved through grace, and that the yoke should not be laid on them. If the yoke is the whole law as a condition of salvation, then the gentiles are saved without it, by faith. That is the reading the strand needs.

\begin{derivation}
  \item[] \((\pvar{TestimonyArguments.SolaFide}{9} \land \pvar{TestimonyArguments.SolaFide}{10}) \rightarrow \pvar{TestimonyArguments.SolaFide}{26}\)\gloss{\emph{would give:} Faith is sufficient: whoever believes is saved}
\end{derivation}
\par\medskip\noindent\textbf{\texttt{johannine\allowbreak{}To\allowbreak{}Faith\allowbreak{}Alone}}\par\nobreak\noindent
\textbf{The Johannine strand.} From John’s Gospel to justification by faith alone. The crowd asks what they must do “to be doing the works of God”, and Jesus answers: “This is the work of God, that you believe in him whom he has sent” (6:28–29). The Gospel says throughout that eternal life comes through believing (3:16–18, 3:36, 5:24, 20:31).

The step needs one more premise, and it is the disputed one: that this believing is trust, which brings nothing of its own to God. Independent of Paul, of Luke 7:50 and of Acts 15.

\begin{derivation}
  \item[] \((\pvar{TestimonyArguments.SolaFide}{36} \land \pvar{TestimonyArguments.SolaFide}{37} \land \pvar{TestimonyArguments.SolaFide}{40}) \rightarrow \pvar{TestimonyArguments.SolaFide}{25}\)\gloss{\emph{would give:} Justification is by faith alone}
\end{derivation}
\par\medskip\noindent\textbf{\texttt{james\allowbreak{}Harmonisation}}\par\nobreak\noindent
The answer to James. James 2:24 says “a person is justified by works and not by faith alone”. If James is attacking a barren faith, and if works are the fruit of saving faith and not its ground, then the verse does not contradict Paul. The agreement is derived from those two premises, not assumed.

\begin{derivation}
  \item[] \((\pvar{TestimonyArguments.SolaFide}{20} \land \pvar{TestimonyArguments.SolaFide}{21}) \rightarrow \pvar{TestimonyArguments.SolaFide}{22}\)\gloss{\emph{would give:} James 2:24 is compatible with Paul, using ‘justify’ in a different sense}
\end{derivation}
\par\medskip\noindent\textbf{\texttt{sola\allowbreak{}Fide}}\par\nobreak\noindent
\textbf{The conclusion}, in three parts: salvation is by grace, not by works, and through faith. They are stated as three claims and not one, so that a position can be shown to hold some and deny another. Trent and the apocalyptic reading each do that, over different parts.

\begin{derivation}
  \item[] \(\pvar{TestimonyArguments.SolaFide}{27} \land \pvar{TestimonyArguments.SolaFide}{28} \land \pvar{TestimonyArguments.SolaFide}{29}\)
\end{derivation}
\par\medskip\noindent\textbf{\texttt{grace\allowbreak{}Not\allowbreak{}Works}}\par\nobreak\noindent
The first two parts: by grace, and not by works. This is what the apocalyptic reading keeps.

\begin{derivation}
  \item[] \(\pvar{TestimonyArguments.SolaFide}{27} \land \pvar{TestimonyArguments.SolaFide}{28}\)
\end{derivation}
\par\medskip\noindent\textbf{\texttt{to\allowbreak{}Grace}}\par\nobreak\noindent
\emph{By grace}, from the verses that say so: “by grace you have been saved… it is the gift of God” (Ephesians 2:8), a gift and not wages (Romans 4:4–5), “according to his own mercy” (Titus 3:5). This part needs neither a strand nor the answer to James.

\begin{derivation}
  \item[] \((\pvar{TestimonyArguments.SolaFide}{1} \land \pvar{TestimonyArguments.SolaFide}{4} \land \pvar{TestimonyArguments.SolaFide}{5}) \rightarrow \pvar{TestimonyArguments.SolaFide}{27}\)\gloss{\emph{would give:} Salvation is by grace, a gift and not wages owed}
\end{derivation}
\par\medskip\noindent\textbf{\texttt{to\allowbreak{}Not\allowbreak{}By\allowbreak{}Works}}\par\nobreak\noindent
\emph{Not by works}, from the same verses, which say “not a result of works” in as many words. It needs two more things: an answer to James 2:24, and that Scripture does not contradict itself. It needs no strand.

\begin{derivation}
  \item[] \((\pvar{TestimonyArguments.SolaFide}{1} \land \pvar{TestimonyArguments.SolaFide}{4} \land \pvar{TestimonyArguments.SolaFide}{5} \land \pvar{TestimonyArguments.SolaFide}{22} \land \pvar{TestimonyArguments.SolaFide}{24}) \rightarrow \pvar{TestimonyArguments.SolaFide}{28}\)\gloss{\emph{would give:} Salvation is not by works: no work is its ground}
\end{derivation}
\par\medskip\noindent\textbf{\texttt{faith\allowbreak{}Alone\allowbreak{}Suffices}}\par\nobreak\noindent
If justification is by faith alone, then faith suffices: if nothing but faith is the condition, faith is enough. This is how Paul’s strand and John’s, which argue the stronger claim, reach the part of the conclusion that Luke’s and Peter’s reach directly.

\begin{derivation}
  \item[] \(\pvar{TestimonyArguments.SolaFide}{25} \rightarrow \pvar{TestimonyArguments.SolaFide}{26}\)\gloss{\emph{would give:} Faith is sufficient: whoever believes is saved}
\end{derivation}
\par\medskip\noindent\textbf{\texttt{to\allowbreak{}Through\allowbreak{}Faith}}\par\nobreak\noindent
\emph{Through faith}, from the claim that faith suffices together with Ephesians 2:8, “through faith” (διὰ πίστεως), which names faith as the means. This is the only part of the conclusion the four strands are needed for, and it needs only what each of them says.

\begin{derivation}
  \item[] \((\pvar{TestimonyArguments.SolaFide}{26} \land \pvar{TestimonyArguments.SolaFide}{1}) \rightarrow \pvar{TestimonyArguments.SolaFide}{29}\)\gloss{\emph{would give:} Salvation is received through faith}
\end{derivation}
\par\medskip\noindent\textbf{\texttt{pauline\allowbreak{}Line}} --- Pauline strand (ἔργα νόμου, πίστις Χριστοῦ)\par\nobreak\noindent
Paul’s line of argument. What is distinctive to it is the reading of Galatians as a letter against requiring circumcision, and two disputed readings: of “works of the law” (ἔργα νόμου) and of “faith in Christ” (πίστις Χριστοῦ). The verses it reads are shared with every other position, which is why denying one of the readings does not cost the argument its texts.

The question of “faith in Christ” is a premise \emph{inside} this strand, not a strand of its own. Every verse with the phrase sets it against the law, so any route through it still has to say what “works of the law” means.

\begin{derivation}
  \item[(1)] \(\pvar{TestimonyArguments.SolaFide}{8}\)\gloss{Galatians opposes requiring circumcision, besides faith, for justification}
  \item[(2)] \(\pvar{TestimonyArguments.SolaFide}{6}\)\gloss{Paul’s ‘works of the law’ means human works in general, not boundary markers}
  \item[(3)] \(\pvar{TestimonyArguments.SolaFide}{7}\)\gloss{In Galatians 2:16 ‘faith of Christ’ (πίστις Χριστοῦ) means our faith in Christ}
  \item[(4)] \((\pvar{TestimonyArguments.SolaFide}{2} \land \pvar{TestimonyArguments.SolaFide}{3} \land \pvar{TestimonyArguments.SolaFide}{8} \land \pvar{TestimonyArguments.SolaFide}{7} \land \pvar{TestimonyArguments.SolaFide}{6}) \rightarrow \pvar{TestimonyArguments.SolaFide}{25}\)
  \item[\(\vdash\)] \(\pvar{TestimonyArguments.SolaFide}{25}\)\gloss{\emph{at issue:} Justification is by faith alone}
\end{derivation}
\par\medskip\noindent\textbf{\texttt{dominical\allowbreak{}Line}} --- Dominical strand (σῴζω at Luke 7:50)\par\nobreak\noindent
Luke’s line of argument, resting on Jesus’ words and not Paul’s. What is distinctive to it is the reading of “saved” at Luke 7:50, and the reading of 7:47 on which the woman’s love is not the cause of her forgiveness.

\begin{derivation}
  \item[(1)] \(\pvar{TestimonyArguments.SolaFide}{18}\)\gloss{‘Saved’ (σῴζω) at Luke 7:50 means saved from sin, not healed}
  \item[(2)] \(\pvar{TestimonyArguments.SolaFide}{19}\)\gloss{Luke 7:47 — her love is the evidence of her forgiveness, not its ground}
  \item[(3)] \((\pvar{TestimonyArguments.SolaFide}{17} \land \pvar{TestimonyArguments.SolaFide}{18} \land \pvar{TestimonyArguments.SolaFide}{19}) \rightarrow \pvar{TestimonyArguments.SolaFide}{26}\)
  \item[\(\vdash\)] \(\pvar{TestimonyArguments.SolaFide}{26}\)\gloss{\emph{at issue:} Faith is sufficient: whoever believes is saved}
\end{derivation}
\par\medskip\noindent\textbf{\texttt{apostolic\allowbreak{}Line}} --- Apostolic strand (Acts 15:7–11)\par\nobreak\noindent
Peter’s line of argument, as Luke reports him. What is distinctive to it is the reading of the yoke.

It shares no premise with Paul’s line. That does not make the two independent in practice. A reader who takes “works of the law” as the marks of Israel will read Peter’s yoke the same way, as Jervell does, so the two readings tend to be won or lost together. The encoding shows only that they are two premises about two texts, and an opponent has to answer both.

\begin{derivation}
  \item[(1)] \(\pvar{TestimonyArguments.SolaFide}{10}\)\gloss{The yoke refused at Acts 15:10 is the whole law as a condition of salvation}
  \item[(2)] \((\pvar{TestimonyArguments.SolaFide}{9} \land \pvar{TestimonyArguments.SolaFide}{10}) \rightarrow \pvar{TestimonyArguments.SolaFide}{26}\)
  \item[\(\vdash\)] \(\pvar{TestimonyArguments.SolaFide}{26}\)\gloss{\emph{at issue:} Faith is sufficient: whoever believes is saved}
\end{derivation}
\par\medskip\noindent\textbf{\texttt{calvin\allowbreak{}On\allowbreak{}John\allowbreak{}Six}}\par\nobreak\noindent
Calvin’s step from John to justification. John never uses Paul’s verb “justify” (δικαιόω), so passing from “eternal life through believing” to “justified by faith alone” is an inference and not the reading of a word. Rated \texttt{plausible}: no cited source grants the strand’s grounds and denies the step, but the step rests on Calvin’s comment on 6:29 and not on a full argument from John’s vocabulary to Paul’s.

{\small\begin{verbatim}
def calvinOnJohnSix : Source :=
  { primary := Reference.work calvinJohn (Bib.Locus.adLoc Scripture.john6_29),
    tradition := Tradition.reformedProtestant,
    confidence := Confidence.plausible }
\end{verbatim}}
\par\medskip\noindent\textbf{\texttt{johannine\allowbreak{}Line}} --- Johannine strand (John 6:29)\par\nobreak\noindent
The Johannine line, resting on Jesus’ words in John’s Gospel. Its texts are its own grounds: no other strand reads them. Its distinctive ground is the reading of believing as trust, which Aquinas denies.

\begin{derivation}
  \item[(1)] \(\pvar{TestimonyArguments.SolaFide}{36}\)\gloss{John 6:28–29 — the work God requires is that you believe in him whom he sent}
  \item[(2)] \(\pvar{TestimonyArguments.SolaFide}{37}\)\gloss{John 3:16–18, 3:36, 5:24, 20:31 — eternal life through believing in the Son}
  \item[(3)] \(\pvar{TestimonyArguments.SolaFide}{40}\)\gloss{The believing of John 6:29 is trust, which brings nothing and receives Christ}
  \item[(4)] \((\pvar{TestimonyArguments.SolaFide}{36} \land \pvar{TestimonyArguments.SolaFide}{37} \land \pvar{TestimonyArguments.SolaFide}{40}) \rightarrow \pvar{TestimonyArguments.SolaFide}{25}\)
  \item[\(\vdash\)] \(\pvar{TestimonyArguments.SolaFide}{25}\)\gloss{\emph{at issue:} Justification is by faith alone}
\end{derivation}
\par\medskip\noindent\textbf{\texttt{james\allowbreak{}Line}} --- Harmonisation of James 2:24\par\nobreak\noindent
The line that answers James. It is not a route to the conclusion. It answers the one verse that appears to stand against it.

\begin{derivation}
  \item[(1)] \(\pvar{TestimonyArguments.SolaFide}{20}\)\gloss{James 2:14–26 targets a barren faith — mere assent — not Paul’s doctrine}
  \item[(2)] \(\pvar{TestimonyArguments.SolaFide}{21}\)\gloss{Good works are the fruit and evidence of saving faith, not its ground}
  \item[(3)] \((\pvar{TestimonyArguments.SolaFide}{20} \land \pvar{TestimonyArguments.SolaFide}{21}) \rightarrow \pvar{TestimonyArguments.SolaFide}{22}\)
  \item[\(\vdash\)] \(\pvar{TestimonyArguments.SolaFide}{22}\)\gloss{\emph{at issue:} James 2:24 is compatible with Paul, using ‘justify’ in a different sense}
\end{derivation}
\par\medskip\noindent\textbf{\texttt{deliverance\allowbreak{}Not\allowbreak{}Faith\allowbreak{}Alone}}\par\nobreak\noindent
\textbf{The apocalyptic reading} (Martyn, Campbell). If “faith of Christ” is Christ’s own faithfulness, and “the righteousness of God” (δικαιοσύνη θεοῦ) is God’s act of rescue, then faith is not the condition on which a person is justified. So justification by faith alone, in the sense the Reformed and the New Perspective share, is denied. Grace and “not by works” are not denied: the reading rejects the condition, not the gift.

\begin{derivation}
  \item[] \((\lnot \pvar{TestimonyArguments.SolaFide}{7} \land \pvar{TestimonyArguments.SolaFide}{16}) \rightarrow \lnot \pvar{TestimonyArguments.SolaFide}{25}\)\gloss{\emph{would deny:} Justification is by faith alone}
\end{derivation}
\par\medskip\noindent\textbf{\texttt{deliverance\allowbreak{}Is\allowbreak{}Grace}}\par\nobreak\noindent
What the apocalyptic reading keeps. If “the righteousness of God” is God’s act of rescue, it is a gift that depends on nothing a person does: by grace, and not by works. This is Campbell’s inference, and his reason for calling the reading more gracious than the one it replaces.

\begin{derivation}
  \item[] \(\pvar{TestimonyArguments.SolaFide}{16} \rightarrow \pvar{TestimonyArguments.SolaFide}{27} \land \pvar{TestimonyArguments.SolaFide}{28}\)\gloss{\emph{would give:} \(\pvar{TestimonyArguments.SolaFide}{27}\) Salvation is by grace, a gift and not wages owed; \(\pvar{TestimonyArguments.SolaFide}{28}\) Salvation is not by works: no work is its ground}
\end{derivation}
\par\medskip\noindent\textbf{\texttt{apocalyptic\allowbreak{}Line}} --- Apocalyptic reading (δικαιοσύνη θεοῦ as deliverance)\par\nobreak\noindent
The apocalyptic line: another route from Paul, arriving at the denial of what Paul’s strand delivers.

\begin{derivation}
  \item[(1)] \(\lnot \pvar{TestimonyArguments.SolaFide}{7}\)\gloss{\emph{denied:} In Galatians 2:16 ‘faith of Christ’ (πίστις Χριστοῦ) means our faith in Christ}
  \item[(2)] \(\pvar{TestimonyArguments.SolaFide}{16}\)\gloss{‘The righteousness of God’ (δικαιοσύνη θεοῦ) names God’s act of rescue in Christ}
  \item[(3)] \((\lnot \pvar{TestimonyArguments.SolaFide}{7} \land \pvar{TestimonyArguments.SolaFide}{16}) \rightarrow \lnot \pvar{TestimonyArguments.SolaFide}{25}\)
  \item[\(\vdash\)] \(\lnot \pvar{TestimonyArguments.SolaFide}{25}\)\gloss{\emph{at issue, the denial of:} Justification is by faith alone}
\end{derivation}
\par\medskip\noindent\textbf{\texttt{sanders\allowbreak{}Line}} --- Covenantal nomism (Sanders, Dunn)\par\nobreak\noindent
\textbf{The New Perspective’s reason} for its reading of “works of the law”. Sanders argued that Judaism in Paul’s day taught entry into the covenant by grace and staying in by obedience (“covenantal nomism”). If so, Paul’s opponents were not trying to earn salvation, and what Paul refuses is not works as such but the practices that kept gentiles out. The starting point is Sanders’. The inference is Dunn’s, and is cited as his.

\begin{derivation}
  \item[(1)] \(\pvar{TestimonyArguments.SolaFide}{31}\)\gloss{Second Temple Judaism: in by grace, staying in by works}
  \item[(2)] \(\pvar{TestimonyArguments.SolaFide}{31} \rightarrow \lnot \pvar{TestimonyArguments.SolaFide}{6}\)
  \item[\(\vdash\)] \(\lnot \pvar{TestimonyArguments.SolaFide}{6}\)\gloss{\emph{at issue, the denial of:} Paul’s ‘works of the law’ means human works in general, not boundary markers}
\end{derivation}
\par\medskip\noindent\textbf{\texttt{critics\allowbreak{}Line}} --- Variegated nomism (Gathercole, Carson et al.)\par\nobreak\noindent
\textbf{The critics’ reply.} If Jewish texts of the period also speak of final vindication according to works, then the demand Galatians refuses, circumcision added to faith as a requirement, is a demand for obedience. In refusing it Paul refuses works as such. So the Reformed reading of “works of the law” is derived here and not assumed: from the occasion of Galatians, and from a claim about first-century Judaism that is argued on the evidence. Cited to Gathercole, and to the volume edited by Carson, O’Brien and Seifrid.

\begin{derivation}
  \item[(1)] \(\lnot \pvar{TestimonyArguments.SolaFide}{31}\)\gloss{\emph{denied:} Second Temple Judaism: in by grace, staying in by works}
  \item[(2)] \(\pvar{TestimonyArguments.SolaFide}{8}\)\gloss{Galatians opposes requiring circumcision, besides faith, for justification}
  \item[(3)] \((\lnot \pvar{TestimonyArguments.SolaFide}{31} \land \pvar{TestimonyArguments.SolaFide}{8}) \rightarrow \pvar{TestimonyArguments.SolaFide}{6}\)
  \item[\(\vdash\)] \(\pvar{TestimonyArguments.SolaFide}{6}\)\gloss{\emph{at issue:} Paul’s ‘works of the law’ means human works in general, not boundary markers}
\end{derivation}
\par\medskip\noindent\textbf{\texttt{jervell\allowbreak{}Line}} --- Law-observant Luke (Jervell)\par\nobreak\noindent
\textbf{Jervell’s reading of the yoke.} Luke shows the Jewish believers keeping the law (Acts 21:20–24), and the council’s decree lays part of it on gentiles (Acts 15:20–21). So, Jervell argues, what Peter refuses at 15:10 is the whole of Israel’s law laid on gentiles, and not the law as a condition of salvation. The step is contested: Bruce accepts both passages and reads the yoke as the law taken as a condition of salvation.

\begin{derivation}
  \item[(1)] \(\pvar{TestimonyArguments.SolaFide}{11}\)\gloss{Acts 15:20–21, 21:20–24: Jewish believers keep the law; gentiles keep part}
  \item[(2)] \(\pvar{TestimonyArguments.SolaFide}{11} \rightarrow \lnot \pvar{TestimonyArguments.SolaFide}{10}\)
  \item[\(\vdash\)] \(\lnot \pvar{TestimonyArguments.SolaFide}{10}\)\gloss{\emph{at issue, the denial of:} The yoke refused at Acts 15:10 is the whole law as a condition of salvation}
\end{derivation}
\par\medskip\noindent\textbf{\texttt{trent\allowbreak{}On\allowbreak{}Increase}}\par\nobreak\noindent
Trent’s step from its definition to merit. If justification \emph{is} the renewal of the inward person, and that renewal grows through good works, then good works increase justification (Session VI, ch. 10 and canon 24). Rated \texttt{consensus}, as \texttt{trent\allowbreak{}On\allowbreak{}Merit} is. The Reformed deny the definition, not the inference, and canon 24 condemns the Reformed alternative by name: that works are “merely the fruits and signs of Justification obtained”.

{\small\begin{verbatim}
def trentOnIncrease : Source :=
  {
    primary :=
      Reference.work tannerDecrees
        (Bib.Locus.sectionRef "Trent, Session VI (1547), canon 24"),
    supporting :=
      [Reference.work tannerDecrees
          (Bib.Locus.sectionRef
            "Trent, Session VI (1547), Decree on Justification, ch. 10")],
    tradition := Tradition.romanCatholic, confidence := Confidence.consensus }
\end{verbatim}}
\par\medskip\noindent\textbf{\texttt{tridentine\allowbreak{}Line}} --- Trent (justification as renewal)\par\nobreak\noindent
\textbf{Trent’s line of argument.} Justification is the renewal of the inward person and not remission of sins only. That renewal grows through good works. So good works increase justification (Session VI, ch. 10, canon 24). The Reformed accept that the renewal grows and deny the definition, so the step carries the objection only as far as the definition does.

\begin{derivation}
  \item[(1)] \(\pvar{TestimonyArguments.SolaFide}{33}\)\gloss{Justification is not remission of sins only, but renewal of the inward man}
  \item[(2)] \(\pvar{TestimonyArguments.SolaFide}{35}\)\gloss{The inward renewal of the justified grows as they do good works in grace}
  \item[(3)] \((\pvar{TestimonyArguments.SolaFide}{33} \land \pvar{TestimonyArguments.SolaFide}{35}) \rightarrow \pvar{TestimonyArguments.SolaFide}{32}\)
  \item[\(\vdash\)] \(\pvar{TestimonyArguments.SolaFide}{32}\)\gloss{\emph{at issue:} Works performed in grace merit an increase of justification}
\end{derivation}
\par\medskip\noindent\textbf{\texttt{trent\allowbreak{}Definition\allowbreak{}Steps}}\par\nobreak\noindent
What Trent’s definition rules out. Justification as a verdict only (canon 11). Justification as distinct from the sanctification that follows it. And faith as sufficient: if justification \emph{is} the renewal of the inward person by charity poured in, then faith “unless hope and charity be added thereto” does not unite a person perfectly with Christ (ch. 7), and whoever says that “nothing else is required to cooperate” is anathema (canon 9). So Trent’s denial that faith suffices is not a separate thesis. It follows from what Trent says justification is.

\begin{derivation}
  \item[(1)] \(\pvar{TestimonyArguments.SolaFide}{33} \rightarrow \lnot \pvar{TestimonyArguments.SolaFide}{30}\)\gloss{\emph{would deny:} Justification is a verdict only: pardon and righteousness counted, not poured in}
  \item[(2)] \(\pvar{TestimonyArguments.SolaFide}{33} \rightarrow \lnot \pvar{TestimonyArguments.SolaFide}{34}\)\gloss{\emph{would deny:} Justification and sanctification are inseparable but distinct}
  \item[(3)] \(\pvar{TestimonyArguments.SolaFide}{33} \rightarrow \lnot \pvar{TestimonyArguments.SolaFide}{26}\)\gloss{\emph{would deny:} Faith is sufficient: whoever believes is saved}
\end{derivation}
\par\medskip\noindent\textbf{\texttt{prooftexts}}\par\nobreak\noindent
The verses the strands read. Every Reformed position shares them, and so does the New Perspective.

\begin{derivation}
  \item[(1)] \(\pvar{TestimonyArguments.SolaFide}{1}\)\gloss{Ephesians 2:8–9 teaches salvation by grace through faith, not of works}
  \item[(2)] \(\pvar{TestimonyArguments.SolaFide}{2}\)\gloss{Romans 3:28 teaches justification by faith apart from works of the law}
  \item[(3)] \(\pvar{TestimonyArguments.SolaFide}{3}\)\gloss{Galatians 2:16 teaches that no one is justified by works of the law}
  \breakrow{(4)} \(\pvar{TestimonyArguments.SolaFide}{4}\)\gloss{Romans 4:4–5 contrasts wages owed with a gift reckoned to the one who believes}
  \breakrow{(5)} \(\pvar{TestimonyArguments.SolaFide}{5}\)\gloss{Titus 3:5 teaches that God saved us not by works done in righteousness}
  \item[(6)] \(\pvar{TestimonyArguments.SolaFide}{17}\)\gloss{Luke 7:50 — Jesus says ‘your faith has saved you’ after declaring sins forgiven}
  \item[(7)] \(\pvar{TestimonyArguments.SolaFide}{9}\)\gloss{Acts 15:9–11 — Peter: hearts cleansed by faith; saved through grace}
\end{derivation}
\par\medskip\noindent\textbf{\texttt{shared\allowbreak{}Grounds}}\par\nobreak\noindent
Everything the strands share: those verses, the two premises that answer James, and that Scripture does not contradict itself.

\begin{derivation}
  \item[(1)] \(\pvar{TestimonyArguments.SolaFide}{1}\)\gloss{Ephesians 2:8–9 teaches salvation by grace through faith, not of works}
  \item[(2)] \(\pvar{TestimonyArguments.SolaFide}{2}\)\gloss{Romans 3:28 teaches justification by faith apart from works of the law}
  \item[(3)] \(\pvar{TestimonyArguments.SolaFide}{3}\)\gloss{Galatians 2:16 teaches that no one is justified by works of the law}
  \breakrow{(4)} \(\pvar{TestimonyArguments.SolaFide}{4}\)\gloss{Romans 4:4–5 contrasts wages owed with a gift reckoned to the one who believes}
  \breakrow{(5)} \(\pvar{TestimonyArguments.SolaFide}{5}\)\gloss{Titus 3:5 teaches that God saved us not by works done in righteousness}
  \breakrow{(6)} \(\pvar{TestimonyArguments.SolaFide}{17}\)\gloss{Luke 7:50 — Jesus says ‘your faith has saved you’ after declaring sins forgiven}
  \breakrow{(7)} \(\pvar{TestimonyArguments.SolaFide}{9}\)\gloss{Acts 15:9–11 — Peter: hearts cleansed by faith; saved through grace}
  \breakrow{(8)} \(\pvar{TestimonyArguments.SolaFide}{20}\)\gloss{James 2:14–26 targets a barren faith — mere assent — not Paul’s doctrine}
  \item[(9)] \(\pvar{TestimonyArguments.SolaFide}{21}\)\gloss{Good works are the fruit and evidence of saving faith, not its ground}
  \item[(10)] \(\pvar{TestimonyArguments.SolaFide}{24}\)\gloss{Scripture does not contradict itself}
\end{derivation}
\par\medskip\noindent\textbf{\texttt{shared\allowbreak{}Grounds\allowbreak{}Without\allowbreak{}James}}\par\nobreak\noindent
The same, with James left unanswered.

\begin{derivation}
  \item[(1)] \(\pvar{TestimonyArguments.SolaFide}{1}\)\gloss{Ephesians 2:8–9 teaches salvation by grace through faith, not of works}
  \item[(2)] \(\pvar{TestimonyArguments.SolaFide}{2}\)\gloss{Romans 3:28 teaches justification by faith apart from works of the law}
  \item[(3)] \(\pvar{TestimonyArguments.SolaFide}{3}\)\gloss{Galatians 2:16 teaches that no one is justified by works of the law}
  \breakrow{(4)} \(\pvar{TestimonyArguments.SolaFide}{4}\)\gloss{Romans 4:4–5 contrasts wages owed with a gift reckoned to the one who believes}
  \breakrow{(5)} \(\pvar{TestimonyArguments.SolaFide}{5}\)\gloss{Titus 3:5 teaches that God saved us not by works done in righteousness}
  \breakrow{(6)} \(\pvar{TestimonyArguments.SolaFide}{17}\)\gloss{Luke 7:50 — Jesus says ‘your faith has saved you’ after declaring sins forgiven}
  \item[(7)] \(\pvar{TestimonyArguments.SolaFide}{9}\)\gloss{Acts 15:9–11 — Peter: hearts cleansed by faith; saved through grace}
  \item[(8)] \(\pvar{TestimonyArguments.SolaFide}{24}\)\gloss{Scripture does not contradict itself}
\end{derivation}
\par\medskip\noindent\textbf{\texttt{conclusion\allowbreak{}Steps}}\par\nobreak\noindent
The steps to the three parts of the conclusion, one for each, and the step from “faith alone” to “faith suffices”.

\begin{derivation}
  \item[(1)] \((\pvar{TestimonyArguments.SolaFide}{1} \land \pvar{TestimonyArguments.SolaFide}{4} \land \pvar{TestimonyArguments.SolaFide}{5}) \rightarrow \pvar{TestimonyArguments.SolaFide}{27}\)\gloss{\emph{would give:} Salvation is by grace, a gift and not wages owed}
  \item[(2)] \((\pvar{TestimonyArguments.SolaFide}{1} \land \pvar{TestimonyArguments.SolaFide}{4} \land \pvar{TestimonyArguments.SolaFide}{5} \land \pvar{TestimonyArguments.SolaFide}{22} \land \pvar{TestimonyArguments.SolaFide}{24}) \rightarrow \pvar{TestimonyArguments.SolaFide}{28}\)\gloss{\emph{would give:} Salvation is not by works: no work is its ground}
  \item[(3)] \(\pvar{TestimonyArguments.SolaFide}{25} \rightarrow \pvar{TestimonyArguments.SolaFide}{26}\)\gloss{\emph{would give:} Faith is sufficient: whoever believes is saved}
  \item[(4)] \((\pvar{TestimonyArguments.SolaFide}{26} \land \pvar{TestimonyArguments.SolaFide}{1}) \rightarrow \pvar{TestimonyArguments.SolaFide}{29}\)\gloss{\emph{would give:} Salvation is received through faith}
\end{derivation}
\par\medskip\noindent\textbf{\texttt{closing\allowbreak{}Steps}}\par\nobreak\noindent
The steps that finish the argument: the answer to James, then the steps to the three parts of the conclusion.

\begin{derivation}
  \item[(1)] \((\pvar{TestimonyArguments.SolaFide}{20} \land \pvar{TestimonyArguments.SolaFide}{21}) \rightarrow \pvar{TestimonyArguments.SolaFide}{22}\)\gloss{\emph{would give:} James 2:24 is compatible with Paul, using ‘justify’ in a different sense}
  \item[(2)] \((\pvar{TestimonyArguments.SolaFide}{1} \land \pvar{TestimonyArguments.SolaFide}{4} \land \pvar{TestimonyArguments.SolaFide}{5}) \rightarrow \pvar{TestimonyArguments.SolaFide}{27}\)\gloss{\emph{would give:} Salvation is by grace, a gift and not wages owed}
  \item[(3)] \((\pvar{TestimonyArguments.SolaFide}{1} \land \pvar{TestimonyArguments.SolaFide}{4} \land \pvar{TestimonyArguments.SolaFide}{5} \land \pvar{TestimonyArguments.SolaFide}{22} \land \pvar{TestimonyArguments.SolaFide}{24}) \rightarrow \pvar{TestimonyArguments.SolaFide}{28}\)\gloss{\emph{would give:} Salvation is not by works: no work is its ground}
  \item[(4)] \(\pvar{TestimonyArguments.SolaFide}{25} \rightarrow \pvar{TestimonyArguments.SolaFide}{26}\)\gloss{\emph{would give:} Faith is sufficient: whoever believes is saved}
  \item[(5)] \((\pvar{TestimonyArguments.SolaFide}{26} \land \pvar{TestimonyArguments.SolaFide}{1}) \rightarrow \pvar{TestimonyArguments.SolaFide}{29}\)\gloss{\emph{would give:} Salvation is received through faith}
\end{derivation}
\paragraph*{Paul’s gospel, against justification as renewal}
\leavevmode\par\nobreak
\par\medskip\noindent\textbf{\texttt{galatian\allowbreak{}Gospel\allowbreak{}Reading\allowbreak{}Source}}\par\nobreak\noindent
The reading of the gospel in Galatians: Christ’s death is the whole ground of justification, and to add a ground is to preach another gospel. Cited to Moo on Galatians and to the Westminster Confession. Rated \texttt{disputed}: Trent accepts every text, and that Christ’s passion merited justification (ch. 7), and denies that nothing is added (canons 24, 32).

{\small\begin{verbatim}
def galatianGospelReadingSource : Source :=
  {
    primary :=
      Reference.work mooGalatians
        (Bib.Locus.adLoc { book := Book.galatians, chapter := 1, verse := 6 }),
    supporting :=
      [Reference.work westminsterConfession (Bib.Locus.sectionRef "XI.1")],
    tradition := Tradition.reformedProtestant,
    confidence := Confidence.disputed }
\end{verbatim}}
\par\medskip\noindent\textbf{\texttt{verdict\allowbreak{}Not\allowbreak{}Renewal\allowbreak{}Source}}\par\nobreak\noindent
The step from “a verdict on Christ’s finished work” to “renewal is no part of justification”. Westminster XI.1: God justifies “not by infusing righteousness into them, but by pardoning their sins, and by accounting and accepting their persons as righteous; not for anything wrought in them”. Rated \texttt{disputed}, and the source that makes it so is Rome’s own. The \emph{Joint Declaration} accepts the starting point, that God “no longer imputes to them their sin” (§22). Its Annex denies the step, defining justification as “forgiveness of sins and being made righteous” (Annex 2 A). Among New Testament scholars, Gorman takes the same side: the verdict is a word that does what it says, and changes the person.

{\small\begin{verbatim}
def verdictNotRenewalSource : Source :=
  {
    primary :=
      Reference.work westminsterConfession (Bib.Locus.sectionRef "XI.1"),
    supporting :=
      [Reference.work Bib.calvinInstitutes (Bib.Locus.sectionRef "III.xi.2"),
        Reference.work morrisApostolicPreaching],
    tradition := Tradition.reformedProtestant,
    confidence := Confidence.disputed }
\end{verbatim}}
\par\medskip\noindent\textbf{\texttt{galatian\allowbreak{}Gospel\allowbreak{}Line}} --- Paul’s gospel (Galatians 1:6–9; 2:21; 5:2–4; 1 Corinthians 15:3)\par\nobreak\noindent
\textbf{Paul’s gospel}, as Galatians states it and 1 Corinthians sums it up.

Galatians is a letter about another gospel. The teachers Paul opposes did not deny Christ. They added circumcision to him, and Paul calls that addition a different gospel and its preachers accursed (1:6–9). Against it he sets what he received: “Christ died for our sins” (1 Corinthians 15:3); “if righteousness were through the law, then Christ died for no purpose” (Galatians 2:21); whoever would be justified by the law is “severed from Christ” (5:2–4).

This line reads those texts as one claim: Christ’s work is the whole ground of justification. With the courtroom sense of Paul’s verb “justify”, it carries that claim to the denial of Trent’s definition. If justification is God’s verdict on the ground of a finished work, then the renewal worked in us is no part of it.

\begin{derivation}
  \item[(1)] \(\pvar{TestimonyArguments.SolaFide}{44}\)\gloss{Galatians 1:6–9: whoever preaches a gospel contrary to the one received is accursed}
  \item[(2)] \(\pvar{TestimonyArguments.SolaFide}{45}\)\gloss{1 Corinthians 15:3: the gospel of first importance — Christ died for our sins}
  \item[(3)] \(\pvar{TestimonyArguments.SolaFide}{46}\)\gloss{Galatians 2:21: if righteousness were through the law, Christ died for no purpose}
  \breakrow{(4)} \(\pvar{TestimonyArguments.SolaFide}{47}\)\gloss{Galatians 5:2–4: you who would be justified by the law are severed from Christ}
  \item[(5)] \((\pvar{TestimonyArguments.SolaFide}{44} \land \pvar{TestimonyArguments.SolaFide}{45} \land \pvar{TestimonyArguments.SolaFide}{46} \land \pvar{TestimonyArguments.SolaFide}{47}) \rightarrow \pvar{TestimonyArguments.SolaFide}{58}\)
  \item[\(\vdash\)] \(\pvar{TestimonyArguments.SolaFide}{58}\)\gloss{\emph{at issue:} Christ’s death for our sins is the whole ground of justification; to add to it is another gospel}
\end{derivation}
\par\medskip\noindent\textbf{\texttt{verdict\allowbreak{}Excludes\allowbreak{}Renewal}}\par\nobreak\noindent
\textbf{The step where Trent parts from Paul.} If justification is God’s verdict, since Paul’s verb “justify” (δικαιόω) is a courtroom word whose opposite is “condemn” (Romans 8:33–34), and if its ground is Christ’s finished work alone, then the renewal of the inward person is no part of it.

\begin{derivation}
  \item[] \((\pvar{TestimonyArguments.SolaFide}{58} \land \pvar{TestimonyArguments.SolaFide}{49} \land \pvar{TestimonyArguments.SolaFide}{48}) \rightarrow \lnot \pvar{TestimonyArguments.SolaFide}{33}\)\gloss{\emph{would deny:} Justification is not remission of sins only, but renewal of the inward man}
\end{derivation}
\paragraph*{What Paul’s word means}
\leavevmode\par\nobreak
\par\medskip\noindent\textbf{\texttt{least\allowbreak{}Meaning\allowbreak{}Source}}\par\nobreak\noindent
The step from three things to “Paul’s ‘justify’ does not mean renewal”: the verb is a courtroom word; Paul has other words for renewal; and a word contributes the least meaning its context requires. The method is Silva’s, applied as Moo applies it to Romans.

Cited at \texttt{well\allowbreak{}Supported}. Trent keeps the first starting point, the verdict, for the justified “are not only reputed, but are truly called, and are, just” (Session VI, ch. 7), and denies that the word means no more; Augustine himself gives the forensic sense beside “made righteous”, “deemed, or reckoned as just” (\emph{On the Spirit and the Letter} 26.45). So a cited reader grants the starting points and denies the conclusion, and the checks count Trent’s reading a credible critique of this step under every standard (\texttt{trent\_\allowbreak{}dissents\_\allowbreak{}from\_\allowbreak{}the\_\allowbreak{}definition\_\allowbreak{}steps}). By the library’s rule that makes the step \texttt{disputed}, and the disputes weigh it at \texttt{well\allowbreak{}Supported}: this is the one place they rest on a rating the rule does not license. Whether the step can rise above \texttt{disputed} turns on Romans 6:7 and 1 Corinthians 6:11, which readers take both ways.

Those who deny the first starting point itself, VanLandingham among them, are weighed in that premise’s rating. Those who hold that God’s verdict renews, Gorman and the \emph{Joint Declaration}, make a claim about what God does, not about the word.

{\small\begin{verbatim}
def leastMeaningSource : Source :=
  { primary := Reference.work silvaBiblicalWords,
    supporting :=
      [Reference.work mooRomans (Bib.Locus.sectionRef "p. 80"),
        Reference.work barrSemantics],
    tradition := Tradition.christianHistoricalGrammatical,
    confidence := Confidence.wellSupported }
\end{verbatim}}
\par\medskip\noindent\textbf{\texttt{lexical\allowbreak{}Line}} --- Paul’s word (δικαιόω, by the rule of least meaning)\par\nobreak\noindent
\textbf{Paul’s word}, read by the rule that a word means no more than its context requires.

Trent’s definition can be read as a claim about Paul’s word: that “justify” (δικαιόω), in Paul, means the renewal of the inward person. This line answers that reading from the word alone.

The verb is a courtroom word: its opposite is “condemn”. Paul has other words for renewal, and sets them beside it. And a word contributes the least meaning its context requires. Nothing in the setting of a verdict requires a second, inward sense, and to supply one from the doctrine is the mistake Barr called “illegitimate totality transfer”: reading into one use of a word everything it can mean elsewhere. So Paul’s “justify” does not mean renewal.

The line says nothing about what God does when he justifies. That is the other way of reading Trent’s definition, and this line does not reach it (the chapter “What Trent means by justification”).

\begin{derivation}
  \item[(1)] \(\pvar{TestimonyArguments.SolaFide}{49}\)\gloss{In Paul, ‘justify’ (δικαιόω) is a verdict: to declare righteous}
  \item[(2)] \(\pvar{TestimonyArguments.SolaFide}{50}\)\gloss{Paul names renewal with other words (ἀνακαίνωσις, ἁγιασμός) and sets them beside ‘justify’}
  \item[(3)] \(\pvar{TestimonyArguments.SolaFide}{51}\)\gloss{A word contributes the least meaning its context requires; one occurrence does not carry the whole concept}
  \item[(4)] \((\pvar{TestimonyArguments.SolaFide}{49} \land \pvar{TestimonyArguments.SolaFide}{50} \land \pvar{TestimonyArguments.SolaFide}{51}) \rightarrow \lnot \pvar{TestimonyArguments.SolaFide}{52}\)
  \item[\(\vdash\)] \(\lnot \pvar{TestimonyArguments.SolaFide}{52}\)\gloss{\emph{at issue, the denial of:} Paul’s word ‘justify’ means the renewal of the inward man, not only a verdict}
\end{derivation}
\paragraph*{Romans 4: the ungodly, justified}
\leavevmode\par\nobreak
\par\medskip\noindent\textbf{\texttt{romans\allowbreak{}Four\allowbreak{}Exegesis\allowbreak{}Source}}\par\nobreak\noindent
The reading of Romans 4:3–8. Paul uses the language of an account (“counted”, λογίζομαι), sets it against wages, and explains what is counted by quoting a psalm. Cited to Moo on Romans. Rated \texttt{well\allowbreak{}Supported}: Augustine and Trent accept the texts and dispute what follows from them, not what Paul says the counting is.

{\small\begin{verbatim}
def romansFourExegesisSource : Source :=
  {
    primary :=
      Reference.work mooRomans
        (Bib.Locus.adLoc { book := Book.romans, chapter := 4, verse := 6 }),
    supporting := [Reference.work bdag (Bib.Locus.sectionRef "s.v. λογίζομαι")],
    tradition := Tradition.christianHistoricalGrammatical,
    confidence := Confidence.wellSupported }
\end{verbatim}}
\par\medskip\noindent\textbf{\texttt{romans\allowbreak{}Four\allowbreak{}Reading\allowbreak{}Source}}\par\nobreak\noindent
The step from Paul’s explanation, with the courtroom sense of the verb, to “renewal is no part of justification”. Cited to Moo on Romans 4:5 and Westminster XI.1. Rated \texttt{disputed}: Augustine accepts the texts and reads “justifies the ungodly” as making the ungodly godly, and Trent’s chapter 7 makes the justice “whereby He maketh us just” the formal cause of justification.

{\small\begin{verbatim}
def romansFourReadingSource : Source :=
  {
    primary :=
      Reference.work mooRomans
        (Bib.Locus.adLoc { book := Book.romans, chapter := 4, verse := 5 }),
    supporting :=
      [Reference.work westminsterConfession (Bib.Locus.sectionRef "XI.1")],
    tradition := Tradition.reformedProtestant,
    confidence := Confidence.disputed }
\end{verbatim}}
\par\medskip\noindent\textbf{\texttt{romans\allowbreak{}Four\allowbreak{}Exegesis\allowbreak{}Line}} --- Romans 4:3–8 (the ungodly, justified)\par\nobreak\noindent
\textbf{Romans 4:3–8, read.} God “justifies the ungodly”, and the righteousness he counts to them is, in Paul’s own explanation, sin not counted.

An argument from one passage, in two steps.

The first step is about what Paul says. God “justifies the ungodly” (4:5). What he counts to them is described in the psalm Paul quotes: “blessed are those whose lawless deeds are forgiven… blessed is the man against whom the Lord will not count his sin” (4:6–8). “Count” (λογίζομαι) is the language of an account, set against wages owed (4:4). So, in Paul’s own argument, the righteousness counted to the ungodly is sin not counted. No source cited here accepts the texts and the sense of the verb and denies that.

The second step puts the first to work, and it goes two ways.

\begin{itemize}
  \item \textbf{Against a claim about Paul’s word.} With the rule of least meaning, the verb Paul explains in that way does not also mean the renewal of the inward person. Rated \texttt{well\allowbreak{}Supported}, as the argument from the word is.
  \item \textbf{Against a claim about what God does.} That what is counted is sin not counted, and that the verdict is a courtroom one, excludes renewal from justification. Augustine accepts every word of Romans 4 and reads the verse the other way: God “justifies the ungodly man, that he may become a godly one” (\emph{On the Spirit and the Letter} 26.45). Trent makes the justice “whereby He maketh us just” the formal cause of justification (ch. 7). So that step is \texttt{disputed}, as the step from Galatians is.
\end{itemize}
\begin{derivation}
  \item[(1)] \(\pvar{TestimonyArguments.SolaFide}{4}\)\gloss{Romans 4:4–5 contrasts wages owed with a gift reckoned to the one who believes}
  \item[(2)] \(\pvar{TestimonyArguments.SolaFide}{55}\)\gloss{Romans 4:6–8: righteousness counted apart from works — sins forgiven, sin not counted}
  \item[(3)] \(\pvar{TestimonyArguments.SolaFide}{56}\)\gloss{‘Count’ (λογίζομαι) in Romans 4:3–11 is the language of an account: credited to one}
  \item[(4)] \((\pvar{TestimonyArguments.SolaFide}{4} \land \pvar{TestimonyArguments.SolaFide}{55} \land \pvar{TestimonyArguments.SolaFide}{56}) \rightarrow \pvar{TestimonyArguments.SolaFide}{57}\)
  \item[\(\vdash\)] \(\pvar{TestimonyArguments.SolaFide}{57}\)\gloss{\emph{at issue:} In Romans 4:6–8 Paul explains righteousness counted apart from works as sin not counted}
\end{derivation}
\par\medskip\noindent\textbf{\texttt{gloss\allowbreak{}By\allowbreak{}Least\allowbreak{}Meaning}}\par\nobreak\noindent
From Paul’s explanation, by the rule of least meaning: the verb he explains in that way does not also mean renewal.

\begin{derivation}
  \item[] \((\pvar{TestimonyArguments.SolaFide}{57} \land \pvar{TestimonyArguments.SolaFide}{51}) \rightarrow \lnot \pvar{TestimonyArguments.SolaFide}{52}\)\gloss{\emph{would deny:} Paul’s word ‘justify’ means the renewal of the inward man, not only a verdict}
\end{derivation}
\par\medskip\noindent\textbf{\texttt{counted\allowbreak{}Excludes\allowbreak{}Renewal}}\par\nobreak\noindent
From Paul’s explanation, with the courtroom sense of the verb: renewal is no part of justification. This is the step Augustine denies.

\begin{derivation}
  \item[] \((\pvar{TestimonyArguments.SolaFide}{57} \land \pvar{TestimonyArguments.SolaFide}{49}) \rightarrow \lnot \pvar{TestimonyArguments.SolaFide}{33}\)\gloss{\emph{would deny:} Justification is not remission of sins only, but renewal of the inward man}
\end{derivation}
\paragraph*{Packages}
\leavevmode\par\nobreak
\par\medskip\noindent\textbf{\texttt{reformed}} --- Reformed (sola fide)\par\nobreak\noindent
The classical Protestant position: all four strands, and every disputed reading they rest on. In Paul, that “works of the law” (ἔργα νόμου) means works in general, and that “faith of Christ” (πίστις Χριστοῦ) means faith in Christ. In Luke, that “saved” (σῴζω) at 7:50 means saved and not healed. In Acts, that the yoke is the law as a condition of salvation. In John, that believing is trust.

\begin{derivation}
  \item[(1)] \(\pvar{TestimonyArguments.SolaFide}{8}\)\gloss{Galatians opposes requiring circumcision, besides faith, for justification}
  \item[(2)] \(\pvar{TestimonyArguments.SolaFide}{6}\)\gloss{Paul’s ‘works of the law’ means human works in general, not boundary markers}
  \item[(3)] \(\pvar{TestimonyArguments.SolaFide}{7}\)\gloss{In Galatians 2:16 ‘faith of Christ’ (πίστις Χριστοῦ) means our faith in Christ}
  \breakrow{(4)} \(\pvar{TestimonyArguments.SolaFide}{18}\)\gloss{‘Saved’ (σῴζω) at Luke 7:50 means saved from sin, not healed}
  \breakrow{(5)} \(\pvar{TestimonyArguments.SolaFide}{19}\)\gloss{Luke 7:47 — her love is the evidence of her forgiveness, not its ground}
  \breakrow{(6)} \(\pvar{TestimonyArguments.SolaFide}{10}\)\gloss{The yoke refused at Acts 15:10 is the whole law as a condition of salvation}
  \breakrow{(7)} \(\pvar{TestimonyArguments.SolaFide}{36}\)\gloss{John 6:28–29 — the work God requires is that you believe in him whom he sent}
  \breakrow{(8)} \(\pvar{TestimonyArguments.SolaFide}{37}\)\gloss{John 3:16–18, 3:36, 5:24, 20:31 — eternal life through believing in the Son}
  \breakrow{(9)} \(\pvar{TestimonyArguments.SolaFide}{40}\)\gloss{The believing of John 6:29 is trust, which brings nothing and receives Christ}
  \breakrow{(10)} \(\pvar{TestimonyArguments.SolaFide}{1}\)\gloss{Ephesians 2:8–9 teaches salvation by grace through faith, not of works}
  \breakrow{(11)} \(\pvar{TestimonyArguments.SolaFide}{2}\)\gloss{Romans 3:28 teaches justification by faith apart from works of the law}
  \breakrow{(12)} \(\pvar{TestimonyArguments.SolaFide}{3}\)\gloss{Galatians 2:16 teaches that no one is justified by works of the law}
  \breakrow{(13)} \(\pvar{TestimonyArguments.SolaFide}{4}\)\gloss{Romans 4:4–5 contrasts wages owed with a gift reckoned to the one who believes}
  \breakrow{(14)} \(\pvar{TestimonyArguments.SolaFide}{5}\)\gloss{Titus 3:5 teaches that God saved us not by works done in righteousness}
  \breakrow{(15)} \(\pvar{TestimonyArguments.SolaFide}{17}\)\gloss{Luke 7:50 — Jesus says ‘your faith has saved you’ after declaring sins forgiven}
  \breakrow{(16)} \(\pvar{TestimonyArguments.SolaFide}{9}\)\gloss{Acts 15:9–11 — Peter: hearts cleansed by faith; saved through grace}
  \breakrow{(17)} \(\pvar{TestimonyArguments.SolaFide}{20}\)\gloss{James 2:14–26 targets a barren faith — mere assent — not Paul’s doctrine}
  \breakrow{(18)} \(\pvar{TestimonyArguments.SolaFide}{21}\)\gloss{Good works are the fruit and evidence of saving faith, not its ground}
  \breakrow{(19)} \(\pvar{TestimonyArguments.SolaFide}{24}\)\gloss{Scripture does not contradict itself}
  \breakrow{(20)} \((\pvar{TestimonyArguments.SolaFide}{2} \land \pvar{TestimonyArguments.SolaFide}{3} \land \pvar{TestimonyArguments.SolaFide}{8} \land \pvar{TestimonyArguments.SolaFide}{7} \land \pvar{TestimonyArguments.SolaFide}{6}) \rightarrow \pvar{TestimonyArguments.SolaFide}{25}\)\gloss{\emph{would give:} Justification is by faith alone}
  \breakrow{(21)} \((\pvar{TestimonyArguments.SolaFide}{17} \land \pvar{TestimonyArguments.SolaFide}{18} \land \pvar{TestimonyArguments.SolaFide}{19}) \rightarrow \pvar{TestimonyArguments.SolaFide}{26}\)\gloss{\emph{would give:} Faith is sufficient: whoever believes is saved}
  \breakrow{(22)} \((\pvar{TestimonyArguments.SolaFide}{9} \land \pvar{TestimonyArguments.SolaFide}{10}) \rightarrow \pvar{TestimonyArguments.SolaFide}{26}\)\gloss{\emph{would give:} Faith is sufficient: whoever believes is saved}
  \breakrow{(23)} \((\pvar{TestimonyArguments.SolaFide}{36} \land \pvar{TestimonyArguments.SolaFide}{37} \land \pvar{TestimonyArguments.SolaFide}{40}) \rightarrow \pvar{TestimonyArguments.SolaFide}{25}\)\gloss{\emph{would give:} Justification is by faith alone}
  \breakrow{(24)} \((\pvar{TestimonyArguments.SolaFide}{20} \land \pvar{TestimonyArguments.SolaFide}{21}) \rightarrow \pvar{TestimonyArguments.SolaFide}{22}\)\gloss{\emph{would give:} James 2:24 is compatible with Paul, using ‘justify’ in a different sense}
  \breakrow{(25)} \((\pvar{TestimonyArguments.SolaFide}{1} \land \pvar{TestimonyArguments.SolaFide}{4} \land \pvar{TestimonyArguments.SolaFide}{5}) \rightarrow \pvar{TestimonyArguments.SolaFide}{27}\)\gloss{\emph{would give:} Salvation is by grace, a gift and not wages owed}
  \breakrow{(26)} \((\pvar{TestimonyArguments.SolaFide}{1} \land \pvar{TestimonyArguments.SolaFide}{4} \land \pvar{TestimonyArguments.SolaFide}{5} \land \pvar{TestimonyArguments.SolaFide}{22} \land \pvar{TestimonyArguments.SolaFide}{24}) \rightarrow \pvar{TestimonyArguments.SolaFide}{28}\)\gloss{\emph{would give:} Salvation is not by works: no work is its ground}
  \breakrow{(27)} \(\pvar{TestimonyArguments.SolaFide}{25} \rightarrow \pvar{TestimonyArguments.SolaFide}{26}\)\gloss{\emph{would give:} Faith is sufficient: whoever believes is saved}
  \item[(28)] \((\pvar{TestimonyArguments.SolaFide}{26} \land \pvar{TestimonyArguments.SolaFide}{1}) \rightarrow \pvar{TestimonyArguments.SolaFide}{29}\)\gloss{\emph{would give:} Salvation is received through faith}
  \item[\(\vdash\)] \(\pvar{TestimonyArguments.SolaFide}{27} \land \pvar{TestimonyArguments.SolaFide}{28} \land \pvar{TestimonyArguments.SolaFide}{29}\)\gloss{\emph{at issue:} \(\pvar{TestimonyArguments.SolaFide}{27}\) Salvation is by grace, a gift and not wages owed; \(\pvar{TestimonyArguments.SolaFide}{28}\) Salvation is not by works: no work is its ground; \(\pvar{TestimonyArguments.SolaFide}{29}\) Salvation is received through faith}
  \keepwithnext
\end{derivation}
\emph{Grounded in scripture alone:} John 3:16–18, 3:36, 5:24, 20:31 — eternal life through believing in the Son.
\par\medskip\noindent\textbf{\texttt{dominical\allowbreak{}Without\allowbreak{}Sozo}} --- Dominical strand (σῴζω at Luke 7:50)\par\nobreak\noindent
Luke’s line without the reading of “saved”. The reading of 7:47 stays: it is not what is being tested.

\begin{derivation}
  \item[(1)] \(\pvar{TestimonyArguments.SolaFide}{19}\)\gloss{Luke 7:47 — her love is the evidence of her forgiveness, not its ground}
  \item[(2)] \((\pvar{TestimonyArguments.SolaFide}{17} \land \pvar{TestimonyArguments.SolaFide}{18} \land \pvar{TestimonyArguments.SolaFide}{19}) \rightarrow \pvar{TestimonyArguments.SolaFide}{26}\)
  \item[\(\vdash\)] \(\pvar{TestimonyArguments.SolaFide}{26}\)\gloss{\emph{at issue:} Faith is sufficient: whoever believes is saved}
\end{derivation}
\par\medskip\noindent\textbf{\texttt{pauline\allowbreak{}Without\allowbreak{}Works\allowbreak{}Of\allowbreak{}Law}} --- Pauline strand (ἔργα νόμου, πίστις Χριστοῦ)\par\nobreak\noindent
Paul’s line without the reading of “works of the law”.

\begin{derivation}
  \item[(1)] \(\pvar{TestimonyArguments.SolaFide}{8}\)\gloss{Galatians opposes requiring circumcision, besides faith, for justification}
  \item[(2)] \(\pvar{TestimonyArguments.SolaFide}{7}\)\gloss{In Galatians 2:16 ‘faith of Christ’ (πίστις Χριστοῦ) means our faith in Christ}
  \item[(3)] \((\pvar{TestimonyArguments.SolaFide}{2} \land \pvar{TestimonyArguments.SolaFide}{3} \land \pvar{TestimonyArguments.SolaFide}{8} \land \pvar{TestimonyArguments.SolaFide}{7} \land \pvar{TestimonyArguments.SolaFide}{6}) \rightarrow \pvar{TestimonyArguments.SolaFide}{25}\)
  \item[\(\vdash\)] \(\pvar{TestimonyArguments.SolaFide}{25}\)\gloss{\emph{at issue:} Justification is by faith alone}
\end{derivation}
\par\medskip\noindent\textbf{\texttt{pauline\allowbreak{}Without\allowbreak{}Pistis\allowbreak{}Christou}} --- Pauline strand (ἔργα νόμου, πίστις Χριστοῦ)\par\nobreak\noindent
Paul’s line without the reading of “faith in Christ”.

\begin{derivation}
  \item[(1)] \(\pvar{TestimonyArguments.SolaFide}{8}\)\gloss{Galatians opposes requiring circumcision, besides faith, for justification}
  \item[(2)] \(\pvar{TestimonyArguments.SolaFide}{6}\)\gloss{Paul’s ‘works of the law’ means human works in general, not boundary markers}
  \item[(3)] \((\pvar{TestimonyArguments.SolaFide}{2} \land \pvar{TestimonyArguments.SolaFide}{3} \land \pvar{TestimonyArguments.SolaFide}{8} \land \pvar{TestimonyArguments.SolaFide}{7} \land \pvar{TestimonyArguments.SolaFide}{6}) \rightarrow \pvar{TestimonyArguments.SolaFide}{25}\)
  \item[\(\vdash\)] \(\pvar{TestimonyArguments.SolaFide}{25}\)\gloss{\emph{at issue:} Justification is by faith alone}
\end{derivation}
\par\medskip\noindent\textbf{\texttt{johannine\allowbreak{}Without\allowbreak{}Trust}} --- Johannine strand (John 6:29)\par\nobreak\noindent
The Johannine line without the reading of believing as trust. Its texts stay: they are not what is being tested.

\begin{derivation}
  \item[(1)] \(\pvar{TestimonyArguments.SolaFide}{36}\)\gloss{John 6:28–29 — the work God requires is that you believe in him whom he sent}
  \item[(2)] \(\pvar{TestimonyArguments.SolaFide}{37}\)\gloss{John 3:16–18, 3:36, 5:24, 20:31 — eternal life through believing in the Son}
  \item[(3)] \((\pvar{TestimonyArguments.SolaFide}{36} \land \pvar{TestimonyArguments.SolaFide}{37} \land \pvar{TestimonyArguments.SolaFide}{40}) \rightarrow \pvar{TestimonyArguments.SolaFide}{25}\)
  \item[\(\vdash\)] \(\pvar{TestimonyArguments.SolaFide}{25}\)\gloss{\emph{at issue:} Justification is by faith alone}
\end{derivation}
\par\medskip\noindent\textbf{\texttt{new\allowbreak{}Perspective}} --- New Perspective on Paul\par\nobreak\noindent
The New Perspective on Paul. It denies the traditional reading of “works of the law” (ἔργα νόμου) and accepts justification by faith, which is what Dunn and Wright hold. It keeps “faith in Christ” as Dunn does, so the one reading of Paul’s that it denies is the first.

The denial is not assumed. \texttt{sanders\allowbreak{}Line} derives it from Sanders’ account of Judaism by Dunn’s inference, which is where its critics engage it.

It takes no position on the yoke of Acts 15:10, or on what believing is in John 6:29. Each of those lines keeps its step and loses its disputed reading, so the position neither grants nor denies the reading on the New Perspective’s behalf.

\begin{derivation}
  \item[(1)] \(\pvar{TestimonyArguments.SolaFide}{8}\)\gloss{Galatians opposes requiring circumcision, besides faith, for justification}
  \item[(2)] \(\pvar{TestimonyArguments.SolaFide}{7}\)\gloss{In Galatians 2:16 ‘faith of Christ’ (πίστις Χριστοῦ) means our faith in Christ}
  \item[(3)] \(\pvar{TestimonyArguments.SolaFide}{31}\)\gloss{Second Temple Judaism: in by grace, staying in by works}
  \breakrow{(4)} \(\pvar{TestimonyArguments.SolaFide}{18}\)\gloss{‘Saved’ (σῴζω) at Luke 7:50 means saved from sin, not healed}
  \breakrow{(5)} \(\pvar{TestimonyArguments.SolaFide}{19}\)\gloss{Luke 7:47 — her love is the evidence of her forgiveness, not its ground}
  \breakrow{(6)} \(\pvar{TestimonyArguments.SolaFide}{36}\)\gloss{John 6:28–29 — the work God requires is that you believe in him whom he sent}
  \breakrow{(7)} \(\pvar{TestimonyArguments.SolaFide}{37}\)\gloss{John 3:16–18, 3:36, 5:24, 20:31 — eternal life through believing in the Son}
  \breakrow{(8)} \(\pvar{TestimonyArguments.SolaFide}{1}\)\gloss{Ephesians 2:8–9 teaches salvation by grace through faith, not of works}
  \breakrow{(9)} \(\pvar{TestimonyArguments.SolaFide}{2}\)\gloss{Romans 3:28 teaches justification by faith apart from works of the law}
  \breakrow{(10)} \(\pvar{TestimonyArguments.SolaFide}{3}\)\gloss{Galatians 2:16 teaches that no one is justified by works of the law}
  \breakrow{(11)} \(\pvar{TestimonyArguments.SolaFide}{4}\)\gloss{Romans 4:4–5 contrasts wages owed with a gift reckoned to the one who believes}
  \breakrow{(12)} \(\pvar{TestimonyArguments.SolaFide}{5}\)\gloss{Titus 3:5 teaches that God saved us not by works done in righteousness}
  \breakrow{(13)} \(\pvar{TestimonyArguments.SolaFide}{17}\)\gloss{Luke 7:50 — Jesus says ‘your faith has saved you’ after declaring sins forgiven}
  \breakrow{(14)} \(\pvar{TestimonyArguments.SolaFide}{9}\)\gloss{Acts 15:9–11 — Peter: hearts cleansed by faith; saved through grace}
  \breakrow{(15)} \(\pvar{TestimonyArguments.SolaFide}{20}\)\gloss{James 2:14–26 targets a barren faith — mere assent — not Paul’s doctrine}
  \breakrow{(16)} \(\pvar{TestimonyArguments.SolaFide}{21}\)\gloss{Good works are the fruit and evidence of saving faith, not its ground}
  \breakrow{(17)} \(\pvar{TestimonyArguments.SolaFide}{24}\)\gloss{Scripture does not contradict itself}
  \breakrow{(18)} \((\pvar{TestimonyArguments.SolaFide}{2} \land \pvar{TestimonyArguments.SolaFide}{3} \land \pvar{TestimonyArguments.SolaFide}{8} \land \pvar{TestimonyArguments.SolaFide}{7} \land \pvar{TestimonyArguments.SolaFide}{6}) \rightarrow \pvar{TestimonyArguments.SolaFide}{25}\)\gloss{\emph{would give:} Justification is by faith alone}
  \breakrow{(19)} \(\pvar{TestimonyArguments.SolaFide}{31} \rightarrow \lnot \pvar{TestimonyArguments.SolaFide}{6}\)\gloss{\emph{would deny:} Paul’s ‘works of the law’ means human works in general, not boundary markers}
  \breakrow{(20)} \((\pvar{TestimonyArguments.SolaFide}{17} \land \pvar{TestimonyArguments.SolaFide}{18} \land \pvar{TestimonyArguments.SolaFide}{19}) \rightarrow \pvar{TestimonyArguments.SolaFide}{26}\)\gloss{\emph{would give:} Faith is sufficient: whoever believes is saved}
  \breakrow{(21)} \((\pvar{TestimonyArguments.SolaFide}{9} \land \pvar{TestimonyArguments.SolaFide}{10}) \rightarrow \pvar{TestimonyArguments.SolaFide}{26}\)\gloss{\emph{would give:} Faith is sufficient: whoever believes is saved}
  \breakrow{(22)} \((\pvar{TestimonyArguments.SolaFide}{36} \land \pvar{TestimonyArguments.SolaFide}{37} \land \pvar{TestimonyArguments.SolaFide}{40}) \rightarrow \pvar{TestimonyArguments.SolaFide}{25}\)\gloss{\emph{would give:} Justification is by faith alone}
  \breakrow{(23)} \((\pvar{TestimonyArguments.SolaFide}{20} \land \pvar{TestimonyArguments.SolaFide}{21}) \rightarrow \pvar{TestimonyArguments.SolaFide}{22}\)\gloss{\emph{would give:} James 2:24 is compatible with Paul, using ‘justify’ in a different sense}
  \breakrow{(24)} \((\pvar{TestimonyArguments.SolaFide}{1} \land \pvar{TestimonyArguments.SolaFide}{4} \land \pvar{TestimonyArguments.SolaFide}{5}) \rightarrow \pvar{TestimonyArguments.SolaFide}{27}\)\gloss{\emph{would give:} Salvation is by grace, a gift and not wages owed}
  \breakrow{(25)} \((\pvar{TestimonyArguments.SolaFide}{1} \land \pvar{TestimonyArguments.SolaFide}{4} \land \pvar{TestimonyArguments.SolaFide}{5} \land \pvar{TestimonyArguments.SolaFide}{22} \land \pvar{TestimonyArguments.SolaFide}{24}) \rightarrow \pvar{TestimonyArguments.SolaFide}{28}\)\gloss{\emph{would give:} Salvation is not by works: no work is its ground}
  \breakrow{(26)} \(\pvar{TestimonyArguments.SolaFide}{25} \rightarrow \pvar{TestimonyArguments.SolaFide}{26}\)\gloss{\emph{would give:} Faith is sufficient: whoever believes is saved}
  \item[(27)] \((\pvar{TestimonyArguments.SolaFide}{26} \land \pvar{TestimonyArguments.SolaFide}{1}) \rightarrow \pvar{TestimonyArguments.SolaFide}{29}\)\gloss{\emph{would give:} Salvation is received through faith}
  \item[\(\vdash\)] \(\pvar{TestimonyArguments.SolaFide}{27} \land \pvar{TestimonyArguments.SolaFide}{28} \land \pvar{TestimonyArguments.SolaFide}{29}\)\gloss{\emph{at issue:} \(\pvar{TestimonyArguments.SolaFide}{27}\) Salvation is by grace, a gift and not wages owed; \(\pvar{TestimonyArguments.SolaFide}{28}\) Salvation is not by works: no work is its ground; \(\pvar{TestimonyArguments.SolaFide}{29}\) Salvation is received through faith}
  \keepwithnext
\end{derivation}
\emph{Grounded in scripture alone:} John 3:16–18, 3:36, 5:24, 20:31 — eternal life through believing in the Son.
\par\medskip\noindent\textbf{\texttt{finnish}} --- Finnish reading of Luther (Mannermaa)\par\nobreak\noindent
The Finnish reading of Luther (Mannermaa): the Reformed case, with the verdict-only account of justification denied. Christ himself is present in faith and is the believer’s righteousness, so justification is not a declaration and nothing more.

On what justification \emph{is}, this stands with Trent against the Reformed account. On sola fide it stands with the Reformed. The encoding shows both, and shows that the two questions are separate.

\begin{derivation}
  \item[(1)] \(\pvar{TestimonyArguments.SolaFide}{8}\)\gloss{Galatians opposes requiring circumcision, besides faith, for justification}
  \item[(2)] \(\pvar{TestimonyArguments.SolaFide}{6}\)\gloss{Paul’s ‘works of the law’ means human works in general, not boundary markers}
  \item[(3)] \(\pvar{TestimonyArguments.SolaFide}{7}\)\gloss{In Galatians 2:16 ‘faith of Christ’ (πίστις Χριστοῦ) means our faith in Christ}
  \breakrow{(4)} \(\pvar{TestimonyArguments.SolaFide}{18}\)\gloss{‘Saved’ (σῴζω) at Luke 7:50 means saved from sin, not healed}
  \breakrow{(5)} \(\pvar{TestimonyArguments.SolaFide}{19}\)\gloss{Luke 7:47 — her love is the evidence of her forgiveness, not its ground}
  \breakrow{(6)} \(\pvar{TestimonyArguments.SolaFide}{10}\)\gloss{The yoke refused at Acts 15:10 is the whole law as a condition of salvation}
  \breakrow{(7)} \(\pvar{TestimonyArguments.SolaFide}{36}\)\gloss{John 6:28–29 — the work God requires is that you believe in him whom he sent}
  \breakrow{(8)} \(\pvar{TestimonyArguments.SolaFide}{37}\)\gloss{John 3:16–18, 3:36, 5:24, 20:31 — eternal life through believing in the Son}
  \breakrow{(9)} \(\pvar{TestimonyArguments.SolaFide}{40}\)\gloss{The believing of John 6:29 is trust, which brings nothing and receives Christ}
  \breakrow{(10)} \(\pvar{TestimonyArguments.SolaFide}{1}\)\gloss{Ephesians 2:8–9 teaches salvation by grace through faith, not of works}
  \breakrow{(11)} \(\pvar{TestimonyArguments.SolaFide}{2}\)\gloss{Romans 3:28 teaches justification by faith apart from works of the law}
  \breakrow{(12)} \(\pvar{TestimonyArguments.SolaFide}{3}\)\gloss{Galatians 2:16 teaches that no one is justified by works of the law}
  \breakrow{(13)} \(\pvar{TestimonyArguments.SolaFide}{4}\)\gloss{Romans 4:4–5 contrasts wages owed with a gift reckoned to the one who believes}
  \breakrow{(14)} \(\pvar{TestimonyArguments.SolaFide}{5}\)\gloss{Titus 3:5 teaches that God saved us not by works done in righteousness}
  \breakrow{(15)} \(\pvar{TestimonyArguments.SolaFide}{17}\)\gloss{Luke 7:50 — Jesus says ‘your faith has saved you’ after declaring sins forgiven}
  \breakrow{(16)} \(\pvar{TestimonyArguments.SolaFide}{9}\)\gloss{Acts 15:9–11 — Peter: hearts cleansed by faith; saved through grace}
  \breakrow{(17)} \(\pvar{TestimonyArguments.SolaFide}{20}\)\gloss{James 2:14–26 targets a barren faith — mere assent — not Paul’s doctrine}
  \breakrow{(18)} \(\pvar{TestimonyArguments.SolaFide}{21}\)\gloss{Good works are the fruit and evidence of saving faith, not its ground}
  \breakrow{(19)} \(\pvar{TestimonyArguments.SolaFide}{24}\)\gloss{Scripture does not contradict itself}
  \breakrow{(20)} \((\pvar{TestimonyArguments.SolaFide}{2} \land \pvar{TestimonyArguments.SolaFide}{3} \land \pvar{TestimonyArguments.SolaFide}{8} \land \pvar{TestimonyArguments.SolaFide}{7} \land \pvar{TestimonyArguments.SolaFide}{6}) \rightarrow \pvar{TestimonyArguments.SolaFide}{25}\)\gloss{\emph{would give:} Justification is by faith alone}
  \breakrow{(21)} \((\pvar{TestimonyArguments.SolaFide}{17} \land \pvar{TestimonyArguments.SolaFide}{18} \land \pvar{TestimonyArguments.SolaFide}{19}) \rightarrow \pvar{TestimonyArguments.SolaFide}{26}\)\gloss{\emph{would give:} Faith is sufficient: whoever believes is saved}
  \breakrow{(22)} \((\pvar{TestimonyArguments.SolaFide}{9} \land \pvar{TestimonyArguments.SolaFide}{10}) \rightarrow \pvar{TestimonyArguments.SolaFide}{26}\)\gloss{\emph{would give:} Faith is sufficient: whoever believes is saved}
  \breakrow{(23)} \((\pvar{TestimonyArguments.SolaFide}{36} \land \pvar{TestimonyArguments.SolaFide}{37} \land \pvar{TestimonyArguments.SolaFide}{40}) \rightarrow \pvar{TestimonyArguments.SolaFide}{25}\)\gloss{\emph{would give:} Justification is by faith alone}
  \breakrow{(24)} \((\pvar{TestimonyArguments.SolaFide}{20} \land \pvar{TestimonyArguments.SolaFide}{21}) \rightarrow \pvar{TestimonyArguments.SolaFide}{22}\)\gloss{\emph{would give:} James 2:24 is compatible with Paul, using ‘justify’ in a different sense}
  \breakrow{(25)} \((\pvar{TestimonyArguments.SolaFide}{1} \land \pvar{TestimonyArguments.SolaFide}{4} \land \pvar{TestimonyArguments.SolaFide}{5}) \rightarrow \pvar{TestimonyArguments.SolaFide}{27}\)\gloss{\emph{would give:} Salvation is by grace, a gift and not wages owed}
  \breakrow{(26)} \((\pvar{TestimonyArguments.SolaFide}{1} \land \pvar{TestimonyArguments.SolaFide}{4} \land \pvar{TestimonyArguments.SolaFide}{5} \land \pvar{TestimonyArguments.SolaFide}{22} \land \pvar{TestimonyArguments.SolaFide}{24}) \rightarrow \pvar{TestimonyArguments.SolaFide}{28}\)\gloss{\emph{would give:} Salvation is not by works: no work is its ground}
  \breakrow{(27)} \(\pvar{TestimonyArguments.SolaFide}{25} \rightarrow \pvar{TestimonyArguments.SolaFide}{26}\)\gloss{\emph{would give:} Faith is sufficient: whoever believes is saved}
  \breakrow{(28)} \((\pvar{TestimonyArguments.SolaFide}{26} \land \pvar{TestimonyArguments.SolaFide}{1}) \rightarrow \pvar{TestimonyArguments.SolaFide}{29}\)\gloss{\emph{would give:} Salvation is received through faith}
  \item[(29)] \(\lnot \pvar{TestimonyArguments.SolaFide}{30}\)\gloss{\emph{denied:} Justification is a verdict only: pardon and righteousness counted, not poured in}
  \item[\(\vdash\)] \(\pvar{TestimonyArguments.SolaFide}{27} \land \pvar{TestimonyArguments.SolaFide}{28} \land \pvar{TestimonyArguments.SolaFide}{29}\)\gloss{\emph{at issue:} \(\pvar{TestimonyArguments.SolaFide}{27}\) Salvation is by grace, a gift and not wages owed; \(\pvar{TestimonyArguments.SolaFide}{28}\) Salvation is not by works: no work is its ground; \(\pvar{TestimonyArguments.SolaFide}{29}\) Salvation is received through faith}
  \keepwithnext
\end{derivation}
\emph{Grounded in scripture alone:} John 3:16–18, 3:36, 5:24, 20:31 — eternal life through believing in the Son.
\par\medskip\noindent\textbf{\texttt{subjective\allowbreak{}Genitive}} --- Subjective genitive (πίστις Χριστοῦ as Christ’s faithfulness)\par\nobreak\noindent
Hays’s reading: “faith of Christ” (πίστις Χριστοῦ) in Galatians 2:16 is Christ’s own faithfulness, so the verse does not name the believer’s faith as the means of justification. Everything else in the Reformed case is granted.

\begin{derivation}
  \item[(1)] \(\pvar{TestimonyArguments.SolaFide}{8}\)\gloss{Galatians opposes requiring circumcision, besides faith, for justification}
  \item[(2)] \(\pvar{TestimonyArguments.SolaFide}{6}\)\gloss{Paul’s ‘works of the law’ means human works in general, not boundary markers}
  \item[(3)] \(\lnot \pvar{TestimonyArguments.SolaFide}{7}\)\gloss{\emph{denied:} In Galatians 2:16 ‘faith of Christ’ (πίστις Χριστοῦ) means our faith in Christ}
  \breakrow{(4)} \(\pvar{TestimonyArguments.SolaFide}{18}\)\gloss{‘Saved’ (σῴζω) at Luke 7:50 means saved from sin, not healed}
  \breakrow{(5)} \(\pvar{TestimonyArguments.SolaFide}{19}\)\gloss{Luke 7:47 — her love is the evidence of her forgiveness, not its ground}
  \breakrow{(6)} \(\pvar{TestimonyArguments.SolaFide}{10}\)\gloss{The yoke refused at Acts 15:10 is the whole law as a condition of salvation}
  \breakrow{(7)} \(\pvar{TestimonyArguments.SolaFide}{36}\)\gloss{John 6:28–29 — the work God requires is that you believe in him whom he sent}
  \breakrow{(8)} \(\pvar{TestimonyArguments.SolaFide}{37}\)\gloss{John 3:16–18, 3:36, 5:24, 20:31 — eternal life through believing in the Son}
  \breakrow{(9)} \(\pvar{TestimonyArguments.SolaFide}{40}\)\gloss{The believing of John 6:29 is trust, which brings nothing and receives Christ}
  \breakrow{(10)} \(\pvar{TestimonyArguments.SolaFide}{1}\)\gloss{Ephesians 2:8–9 teaches salvation by grace through faith, not of works}
  \breakrow{(11)} \(\pvar{TestimonyArguments.SolaFide}{2}\)\gloss{Romans 3:28 teaches justification by faith apart from works of the law}
  \breakrow{(12)} \(\pvar{TestimonyArguments.SolaFide}{3}\)\gloss{Galatians 2:16 teaches that no one is justified by works of the law}
  \breakrow{(13)} \(\pvar{TestimonyArguments.SolaFide}{4}\)\gloss{Romans 4:4–5 contrasts wages owed with a gift reckoned to the one who believes}
  \breakrow{(14)} \(\pvar{TestimonyArguments.SolaFide}{5}\)\gloss{Titus 3:5 teaches that God saved us not by works done in righteousness}
  \breakrow{(15)} \(\pvar{TestimonyArguments.SolaFide}{17}\)\gloss{Luke 7:50 — Jesus says ‘your faith has saved you’ after declaring sins forgiven}
  \breakrow{(16)} \(\pvar{TestimonyArguments.SolaFide}{9}\)\gloss{Acts 15:9–11 — Peter: hearts cleansed by faith; saved through grace}
  \breakrow{(17)} \(\pvar{TestimonyArguments.SolaFide}{20}\)\gloss{James 2:14–26 targets a barren faith — mere assent — not Paul’s doctrine}
  \breakrow{(18)} \(\pvar{TestimonyArguments.SolaFide}{21}\)\gloss{Good works are the fruit and evidence of saving faith, not its ground}
  \breakrow{(19)} \(\pvar{TestimonyArguments.SolaFide}{24}\)\gloss{Scripture does not contradict itself}
  \breakrow{(20)} \((\pvar{TestimonyArguments.SolaFide}{2} \land \pvar{TestimonyArguments.SolaFide}{3} \land \pvar{TestimonyArguments.SolaFide}{8} \land \pvar{TestimonyArguments.SolaFide}{7} \land \pvar{TestimonyArguments.SolaFide}{6}) \rightarrow \pvar{TestimonyArguments.SolaFide}{25}\)\gloss{\emph{would give:} Justification is by faith alone}
  \breakrow{(21)} \((\pvar{TestimonyArguments.SolaFide}{17} \land \pvar{TestimonyArguments.SolaFide}{18} \land \pvar{TestimonyArguments.SolaFide}{19}) \rightarrow \pvar{TestimonyArguments.SolaFide}{26}\)\gloss{\emph{would give:} Faith is sufficient: whoever believes is saved}
  \breakrow{(22)} \((\pvar{TestimonyArguments.SolaFide}{9} \land \pvar{TestimonyArguments.SolaFide}{10}) \rightarrow \pvar{TestimonyArguments.SolaFide}{26}\)\gloss{\emph{would give:} Faith is sufficient: whoever believes is saved}
  \breakrow{(23)} \((\pvar{TestimonyArguments.SolaFide}{36} \land \pvar{TestimonyArguments.SolaFide}{37} \land \pvar{TestimonyArguments.SolaFide}{40}) \rightarrow \pvar{TestimonyArguments.SolaFide}{25}\)\gloss{\emph{would give:} Justification is by faith alone}
  \breakrow{(24)} \((\pvar{TestimonyArguments.SolaFide}{20} \land \pvar{TestimonyArguments.SolaFide}{21}) \rightarrow \pvar{TestimonyArguments.SolaFide}{22}\)\gloss{\emph{would give:} James 2:24 is compatible with Paul, using ‘justify’ in a different sense}
  \breakrow{(25)} \((\pvar{TestimonyArguments.SolaFide}{1} \land \pvar{TestimonyArguments.SolaFide}{4} \land \pvar{TestimonyArguments.SolaFide}{5}) \rightarrow \pvar{TestimonyArguments.SolaFide}{27}\)\gloss{\emph{would give:} Salvation is by grace, a gift and not wages owed}
  \breakrow{(26)} \((\pvar{TestimonyArguments.SolaFide}{1} \land \pvar{TestimonyArguments.SolaFide}{4} \land \pvar{TestimonyArguments.SolaFide}{5} \land \pvar{TestimonyArguments.SolaFide}{22} \land \pvar{TestimonyArguments.SolaFide}{24}) \rightarrow \pvar{TestimonyArguments.SolaFide}{28}\)\gloss{\emph{would give:} Salvation is not by works: no work is its ground}
  \breakrow{(27)} \(\pvar{TestimonyArguments.SolaFide}{25} \rightarrow \pvar{TestimonyArguments.SolaFide}{26}\)\gloss{\emph{would give:} Faith is sufficient: whoever believes is saved}
  \item[(28)] \((\pvar{TestimonyArguments.SolaFide}{26} \land \pvar{TestimonyArguments.SolaFide}{1}) \rightarrow \pvar{TestimonyArguments.SolaFide}{29}\)\gloss{\emph{would give:} Salvation is received through faith}
  \item[\(\vdash\)] \(\pvar{TestimonyArguments.SolaFide}{27} \land \pvar{TestimonyArguments.SolaFide}{28} \land \pvar{TestimonyArguments.SolaFide}{29}\)\gloss{\emph{at issue:} \(\pvar{TestimonyArguments.SolaFide}{27}\) Salvation is by grace, a gift and not wages owed; \(\pvar{TestimonyArguments.SolaFide}{28}\) Salvation is not by works: no work is its ground; \(\pvar{TestimonyArguments.SolaFide}{29}\) Salvation is received through faith}
  \keepwithnext
\end{derivation}
\emph{Grounded in scripture alone:} John 3:16–18, 3:36, 5:24, 20:31 — eternal life through believing in the Son.
\par\medskip\noindent\textbf{\texttt{law\allowbreak{}Observant\allowbreak{}Luke}} --- Law-observant Luke (the yoke as Israel’s law for gentiles)\par\nobreak\noindent
Jervell’s reading of Luke: the yoke Peter refuses at Acts 15:10 is Israel’s law laid on gentiles as a mark of belonging, and not the law as a condition of salvation. Everything else in the Reformed case is granted. This is the rival Peter’s strand was written against.

\begin{derivation}
  \item[(1)] \(\pvar{TestimonyArguments.SolaFide}{8}\)\gloss{Galatians opposes requiring circumcision, besides faith, for justification}
  \item[(2)] \(\pvar{TestimonyArguments.SolaFide}{6}\)\gloss{Paul’s ‘works of the law’ means human works in general, not boundary markers}
  \item[(3)] \(\pvar{TestimonyArguments.SolaFide}{7}\)\gloss{In Galatians 2:16 ‘faith of Christ’ (πίστις Χριστοῦ) means our faith in Christ}
  \breakrow{(4)} \(\pvar{TestimonyArguments.SolaFide}{18}\)\gloss{‘Saved’ (σῴζω) at Luke 7:50 means saved from sin, not healed}
  \breakrow{(5)} \(\pvar{TestimonyArguments.SolaFide}{19}\)\gloss{Luke 7:47 — her love is the evidence of her forgiveness, not its ground}
  \breakrow{(6)} \(\lnot \pvar{TestimonyArguments.SolaFide}{10}\)\gloss{\emph{denied:} The yoke refused at Acts 15:10 is the whole law as a condition of salvation}
  \breakrow{(7)} \(\pvar{TestimonyArguments.SolaFide}{36}\)\gloss{John 6:28–29 — the work God requires is that you believe in him whom he sent}
  \breakrow{(8)} \(\pvar{TestimonyArguments.SolaFide}{37}\)\gloss{John 3:16–18, 3:36, 5:24, 20:31 — eternal life through believing in the Son}
  \breakrow{(9)} \(\pvar{TestimonyArguments.SolaFide}{40}\)\gloss{The believing of John 6:29 is trust, which brings nothing and receives Christ}
  \breakrow{(10)} \(\pvar{TestimonyArguments.SolaFide}{1}\)\gloss{Ephesians 2:8–9 teaches salvation by grace through faith, not of works}
  \breakrow{(11)} \(\pvar{TestimonyArguments.SolaFide}{2}\)\gloss{Romans 3:28 teaches justification by faith apart from works of the law}
  \breakrow{(12)} \(\pvar{TestimonyArguments.SolaFide}{3}\)\gloss{Galatians 2:16 teaches that no one is justified by works of the law}
  \breakrow{(13)} \(\pvar{TestimonyArguments.SolaFide}{4}\)\gloss{Romans 4:4–5 contrasts wages owed with a gift reckoned to the one who believes}
  \breakrow{(14)} \(\pvar{TestimonyArguments.SolaFide}{5}\)\gloss{Titus 3:5 teaches that God saved us not by works done in righteousness}
  \breakrow{(15)} \(\pvar{TestimonyArguments.SolaFide}{17}\)\gloss{Luke 7:50 — Jesus says ‘your faith has saved you’ after declaring sins forgiven}
  \breakrow{(16)} \(\pvar{TestimonyArguments.SolaFide}{9}\)\gloss{Acts 15:9–11 — Peter: hearts cleansed by faith; saved through grace}
  \breakrow{(17)} \(\pvar{TestimonyArguments.SolaFide}{20}\)\gloss{James 2:14–26 targets a barren faith — mere assent — not Paul’s doctrine}
  \breakrow{(18)} \(\pvar{TestimonyArguments.SolaFide}{21}\)\gloss{Good works are the fruit and evidence of saving faith, not its ground}
  \breakrow{(19)} \(\pvar{TestimonyArguments.SolaFide}{24}\)\gloss{Scripture does not contradict itself}
  \breakrow{(20)} \((\pvar{TestimonyArguments.SolaFide}{2} \land \pvar{TestimonyArguments.SolaFide}{3} \land \pvar{TestimonyArguments.SolaFide}{8} \land \pvar{TestimonyArguments.SolaFide}{7} \land \pvar{TestimonyArguments.SolaFide}{6}) \rightarrow \pvar{TestimonyArguments.SolaFide}{25}\)\gloss{\emph{would give:} Justification is by faith alone}
  \breakrow{(21)} \((\pvar{TestimonyArguments.SolaFide}{17} \land \pvar{TestimonyArguments.SolaFide}{18} \land \pvar{TestimonyArguments.SolaFide}{19}) \rightarrow \pvar{TestimonyArguments.SolaFide}{26}\)\gloss{\emph{would give:} Faith is sufficient: whoever believes is saved}
  \breakrow{(22)} \((\pvar{TestimonyArguments.SolaFide}{9} \land \pvar{TestimonyArguments.SolaFide}{10}) \rightarrow \pvar{TestimonyArguments.SolaFide}{26}\)\gloss{\emph{would give:} Faith is sufficient: whoever believes is saved}
  \breakrow{(23)} \((\pvar{TestimonyArguments.SolaFide}{36} \land \pvar{TestimonyArguments.SolaFide}{37} \land \pvar{TestimonyArguments.SolaFide}{40}) \rightarrow \pvar{TestimonyArguments.SolaFide}{25}\)\gloss{\emph{would give:} Justification is by faith alone}
  \breakrow{(24)} \((\pvar{TestimonyArguments.SolaFide}{20} \land \pvar{TestimonyArguments.SolaFide}{21}) \rightarrow \pvar{TestimonyArguments.SolaFide}{22}\)\gloss{\emph{would give:} James 2:24 is compatible with Paul, using ‘justify’ in a different sense}
  \breakrow{(25)} \((\pvar{TestimonyArguments.SolaFide}{1} \land \pvar{TestimonyArguments.SolaFide}{4} \land \pvar{TestimonyArguments.SolaFide}{5}) \rightarrow \pvar{TestimonyArguments.SolaFide}{27}\)\gloss{\emph{would give:} Salvation is by grace, a gift and not wages owed}
  \breakrow{(26)} \((\pvar{TestimonyArguments.SolaFide}{1} \land \pvar{TestimonyArguments.SolaFide}{4} \land \pvar{TestimonyArguments.SolaFide}{5} \land \pvar{TestimonyArguments.SolaFide}{22} \land \pvar{TestimonyArguments.SolaFide}{24}) \rightarrow \pvar{TestimonyArguments.SolaFide}{28}\)\gloss{\emph{would give:} Salvation is not by works: no work is its ground}
  \breakrow{(27)} \(\pvar{TestimonyArguments.SolaFide}{25} \rightarrow \pvar{TestimonyArguments.SolaFide}{26}\)\gloss{\emph{would give:} Faith is sufficient: whoever believes is saved}
  \item[(28)] \((\pvar{TestimonyArguments.SolaFide}{26} \land \pvar{TestimonyArguments.SolaFide}{1}) \rightarrow \pvar{TestimonyArguments.SolaFide}{29}\)\gloss{\emph{would give:} Salvation is received through faith}
  \item[\(\vdash\)] \(\pvar{TestimonyArguments.SolaFide}{27} \land \pvar{TestimonyArguments.SolaFide}{28} \land \pvar{TestimonyArguments.SolaFide}{29}\)\gloss{\emph{at issue:} \(\pvar{TestimonyArguments.SolaFide}{27}\) Salvation is by grace, a gift and not wages owed; \(\pvar{TestimonyArguments.SolaFide}{28}\) Salvation is not by works: no work is its ground; \(\pvar{TestimonyArguments.SolaFide}{29}\) Salvation is received through faith}
  \keepwithnext
\end{derivation}
\emph{Grounded in scripture alone:} John 3:16–18, 3:36, 5:24, 20:31 — eternal life through believing in the Son.
\par\medskip\noindent\textbf{\texttt{critical\allowbreak{}Authorship}} --- Reformed, with the critical view of authorship\par\nobreak\noindent
The Reformed case with every critical conclusion about authorship granted: Ephesians and Titus not by Paul, 1 and 2 Peter not by Peter.

Nothing else changes. The case reads these letters as Scripture, held together by the premise that Scripture does not contradict itself, and not as the testimony of a named apostle. So conceding their authorship costs the argument nothing. It would cost an argument that cited Ephesians 2:9 as evidence of what \emph{Paul} meant by “works”, which is why the reading of “works of the law” leans on Romans.

\begin{derivation}
  \item[(1)] \(\pvar{TestimonyArguments.SolaFide}{8}\)\gloss{Galatians opposes requiring circumcision, besides faith, for justification}
  \item[(2)] \(\pvar{TestimonyArguments.SolaFide}{6}\)\gloss{Paul’s ‘works of the law’ means human works in general, not boundary markers}
  \item[(3)] \(\pvar{TestimonyArguments.SolaFide}{7}\)\gloss{In Galatians 2:16 ‘faith of Christ’ (πίστις Χριστοῦ) means our faith in Christ}
  \breakrow{(4)} \(\pvar{TestimonyArguments.SolaFide}{18}\)\gloss{‘Saved’ (σῴζω) at Luke 7:50 means saved from sin, not healed}
  \breakrow{(5)} \(\pvar{TestimonyArguments.SolaFide}{19}\)\gloss{Luke 7:47 — her love is the evidence of her forgiveness, not its ground}
  \breakrow{(6)} \(\pvar{TestimonyArguments.SolaFide}{10}\)\gloss{The yoke refused at Acts 15:10 is the whole law as a condition of salvation}
  \breakrow{(7)} \(\pvar{TestimonyArguments.SolaFide}{36}\)\gloss{John 6:28–29 — the work God requires is that you believe in him whom he sent}
  \breakrow{(8)} \(\pvar{TestimonyArguments.SolaFide}{37}\)\gloss{John 3:16–18, 3:36, 5:24, 20:31 — eternal life through believing in the Son}
  \breakrow{(9)} \(\pvar{TestimonyArguments.SolaFide}{40}\)\gloss{The believing of John 6:29 is trust, which brings nothing and receives Christ}
  \breakrow{(10)} \(\pvar{TestimonyArguments.SolaFide}{1}\)\gloss{Ephesians 2:8–9 teaches salvation by grace through faith, not of works}
  \breakrow{(11)} \(\pvar{TestimonyArguments.SolaFide}{2}\)\gloss{Romans 3:28 teaches justification by faith apart from works of the law}
  \breakrow{(12)} \(\pvar{TestimonyArguments.SolaFide}{3}\)\gloss{Galatians 2:16 teaches that no one is justified by works of the law}
  \breakrow{(13)} \(\pvar{TestimonyArguments.SolaFide}{4}\)\gloss{Romans 4:4–5 contrasts wages owed with a gift reckoned to the one who believes}
  \breakrow{(14)} \(\pvar{TestimonyArguments.SolaFide}{5}\)\gloss{Titus 3:5 teaches that God saved us not by works done in righteousness}
  \breakrow{(15)} \(\pvar{TestimonyArguments.SolaFide}{17}\)\gloss{Luke 7:50 — Jesus says ‘your faith has saved you’ after declaring sins forgiven}
  \breakrow{(16)} \(\pvar{TestimonyArguments.SolaFide}{9}\)\gloss{Acts 15:9–11 — Peter: hearts cleansed by faith; saved through grace}
  \breakrow{(17)} \(\pvar{TestimonyArguments.SolaFide}{20}\)\gloss{James 2:14–26 targets a barren faith — mere assent — not Paul’s doctrine}
  \breakrow{(18)} \(\pvar{TestimonyArguments.SolaFide}{21}\)\gloss{Good works are the fruit and evidence of saving faith, not its ground}
  \breakrow{(19)} \(\pvar{TestimonyArguments.SolaFide}{24}\)\gloss{Scripture does not contradict itself}
  \breakrow{(20)} \((\pvar{TestimonyArguments.SolaFide}{2} \land \pvar{TestimonyArguments.SolaFide}{3} \land \pvar{TestimonyArguments.SolaFide}{8} \land \pvar{TestimonyArguments.SolaFide}{7} \land \pvar{TestimonyArguments.SolaFide}{6}) \rightarrow \pvar{TestimonyArguments.SolaFide}{25}\)\gloss{\emph{would give:} Justification is by faith alone}
  \breakrow{(21)} \((\pvar{TestimonyArguments.SolaFide}{17} \land \pvar{TestimonyArguments.SolaFide}{18} \land \pvar{TestimonyArguments.SolaFide}{19}) \rightarrow \pvar{TestimonyArguments.SolaFide}{26}\)\gloss{\emph{would give:} Faith is sufficient: whoever believes is saved}
  \breakrow{(22)} \((\pvar{TestimonyArguments.SolaFide}{9} \land \pvar{TestimonyArguments.SolaFide}{10}) \rightarrow \pvar{TestimonyArguments.SolaFide}{26}\)\gloss{\emph{would give:} Faith is sufficient: whoever believes is saved}
  \breakrow{(23)} \((\pvar{TestimonyArguments.SolaFide}{36} \land \pvar{TestimonyArguments.SolaFide}{37} \land \pvar{TestimonyArguments.SolaFide}{40}) \rightarrow \pvar{TestimonyArguments.SolaFide}{25}\)\gloss{\emph{would give:} Justification is by faith alone}
  \breakrow{(24)} \((\pvar{TestimonyArguments.SolaFide}{20} \land \pvar{TestimonyArguments.SolaFide}{21}) \rightarrow \pvar{TestimonyArguments.SolaFide}{22}\)\gloss{\emph{would give:} James 2:24 is compatible with Paul, using ‘justify’ in a different sense}
  \breakrow{(25)} \((\pvar{TestimonyArguments.SolaFide}{1} \land \pvar{TestimonyArguments.SolaFide}{4} \land \pvar{TestimonyArguments.SolaFide}{5}) \rightarrow \pvar{TestimonyArguments.SolaFide}{27}\)\gloss{\emph{would give:} Salvation is by grace, a gift and not wages owed}
  \breakrow{(26)} \((\pvar{TestimonyArguments.SolaFide}{1} \land \pvar{TestimonyArguments.SolaFide}{4} \land \pvar{TestimonyArguments.SolaFide}{5} \land \pvar{TestimonyArguments.SolaFide}{22} \land \pvar{TestimonyArguments.SolaFide}{24}) \rightarrow \pvar{TestimonyArguments.SolaFide}{28}\)\gloss{\emph{would give:} Salvation is not by works: no work is its ground}
  \breakrow{(27)} \(\pvar{TestimonyArguments.SolaFide}{25} \rightarrow \pvar{TestimonyArguments.SolaFide}{26}\)\gloss{\emph{would give:} Faith is sufficient: whoever believes is saved}
  \breakrow{(28)} \((\pvar{TestimonyArguments.SolaFide}{26} \land \pvar{TestimonyArguments.SolaFide}{1}) \rightarrow \pvar{TestimonyArguments.SolaFide}{29}\)\gloss{\emph{would give:} Salvation is received through faith}
  \breakrow{(29)} \(\lnot \pvar{TestimonyArguments.SolaFide}{12}\)\gloss{\emph{denied:} Ephesians was written by Paul}
  \breakrow{(30)} \(\lnot \pvar{TestimonyArguments.SolaFide}{13}\)\gloss{\emph{denied:} Titus was written by Paul}
  \breakrow{(31)} \(\lnot \pvar{TestimonyArguments.SolaFide}{14}\)\gloss{\emph{denied:} 1 Peter was written by the apostle Peter}
  \item[(32)] \(\lnot \pvar{TestimonyArguments.SolaFide}{15}\)\gloss{\emph{denied:} 2 Peter was written by the apostle Peter}
  \item[\(\vdash\)] \(\pvar{TestimonyArguments.SolaFide}{27} \land \pvar{TestimonyArguments.SolaFide}{28} \land \pvar{TestimonyArguments.SolaFide}{29}\)\gloss{\emph{at issue:} \(\pvar{TestimonyArguments.SolaFide}{27}\) Salvation is by grace, a gift and not wages owed; \(\pvar{TestimonyArguments.SolaFide}{28}\) Salvation is not by works: no work is its ground; \(\pvar{TestimonyArguments.SolaFide}{29}\) Salvation is received through faith}
  \keepwithnext
\end{derivation}
\emph{Grounded in scripture alone:} John 3:16–18, 3:36, 5:24, 20:31 — eternal life through believing in the Son.
\par\medskip\noindent\textbf{\texttt{apocalyptic}} --- Apocalyptic reading of Paul (Martyn, Campbell)\par\nobreak\noindent
The apocalyptic reading of Paul (Martyn, Campbell). It rejects the traditional framing of justification and the New Perspective’s as well. For it the question is neither how an individual is accepted nor how Jew and gentile are related. It is God’s rescue of the world in Christ.

It shares Galatians with the Reformed reading: the text of 2:16 and the occasion of the letter. It closes with the same step to salvation. It does not share Luke’s strand, Peter’s or John’s, which neither Martyn nor Campbell argues from. To grant it Luke 7:50, Acts 15 or John 6:29 would be to encode a position nobody holds.

\begin{derivation}
  \item[(1)] \(\lnot \pvar{TestimonyArguments.SolaFide}{7}\)\gloss{\emph{denied:} In Galatians 2:16 ‘faith of Christ’ (πίστις Χριστοῦ) means our faith in Christ}
  \item[(2)] \(\pvar{TestimonyArguments.SolaFide}{16}\)\gloss{‘The righteousness of God’ (δικαιοσύνη θεοῦ) names God’s act of rescue in Christ}
  \item[(3)] \(\pvar{TestimonyArguments.SolaFide}{3}\)\gloss{Galatians 2:16 teaches that no one is justified by works of the law}
  \breakrow{(4)} \(\pvar{TestimonyArguments.SolaFide}{8}\)\gloss{Galatians opposes requiring circumcision, besides faith, for justification}
  \breakrow{(5)} \((\lnot \pvar{TestimonyArguments.SolaFide}{7} \land \pvar{TestimonyArguments.SolaFide}{16}) \rightarrow \lnot \pvar{TestimonyArguments.SolaFide}{25}\)\gloss{\emph{would deny:} Justification is by faith alone}
  \breakrow{(6)} \(\pvar{TestimonyArguments.SolaFide}{16} \rightarrow \pvar{TestimonyArguments.SolaFide}{27} \land \pvar{TestimonyArguments.SolaFide}{28}\)\gloss{\emph{would give:} \(\pvar{TestimonyArguments.SolaFide}{27}\) Salvation is by grace, a gift and not wages owed; \(\pvar{TestimonyArguments.SolaFide}{28}\) Salvation is not by works: no work is its ground}
  \item[(7)] \((\pvar{TestimonyArguments.SolaFide}{26} \land \pvar{TestimonyArguments.SolaFide}{1}) \rightarrow \pvar{TestimonyArguments.SolaFide}{29}\)\gloss{\emph{would give:} Salvation is received through faith}
  \item[\(\vdash\)] \(\pvar{TestimonyArguments.SolaFide}{27} \land \pvar{TestimonyArguments.SolaFide}{28} \land \pvar{TestimonyArguments.SolaFide}{29}\)\gloss{\emph{at issue:} \(\pvar{TestimonyArguments.SolaFide}{27}\) Salvation is by grace, a gift and not wages owed; \(\pvar{TestimonyArguments.SolaFide}{28}\) Salvation is not by works: no work is its ground; \(\pvar{TestimonyArguments.SolaFide}{29}\) Salvation is received through faith}
  \keepwithnext
\end{derivation}
\emph{No premise here rests on scripture alone.}
\par\medskip\noindent\textbf{\texttt{apocalyptic\allowbreak{}On\allowbreak{}Grace\allowbreak{}And\allowbreak{}Works}} --- Apocalyptic reading, on grace and works only\par\nobreak\noindent
The apocalyptic reading, asked only about the first two parts of the conclusion: by grace, and not by works.

\begin{derivation}
  \item[(1)] \(\lnot \pvar{TestimonyArguments.SolaFide}{7}\)\gloss{\emph{denied:} In Galatians 2:16 ‘faith of Christ’ (πίστις Χριστοῦ) means our faith in Christ}
  \item[(2)] \(\pvar{TestimonyArguments.SolaFide}{16}\)\gloss{‘The righteousness of God’ (δικαιοσύνη θεοῦ) names God’s act of rescue in Christ}
  \item[(3)] \(\pvar{TestimonyArguments.SolaFide}{3}\)\gloss{Galatians 2:16 teaches that no one is justified by works of the law}
  \breakrow{(4)} \(\pvar{TestimonyArguments.SolaFide}{8}\)\gloss{Galatians opposes requiring circumcision, besides faith, for justification}
  \breakrow{(5)} \((\lnot \pvar{TestimonyArguments.SolaFide}{7} \land \pvar{TestimonyArguments.SolaFide}{16}) \rightarrow \lnot \pvar{TestimonyArguments.SolaFide}{25}\)\gloss{\emph{would deny:} Justification is by faith alone}
  \breakrow{(6)} \(\pvar{TestimonyArguments.SolaFide}{16} \rightarrow \pvar{TestimonyArguments.SolaFide}{27} \land \pvar{TestimonyArguments.SolaFide}{28}\)
  \item[(7)] \((\pvar{TestimonyArguments.SolaFide}{26} \land \pvar{TestimonyArguments.SolaFide}{1}) \rightarrow \pvar{TestimonyArguments.SolaFide}{29}\)\gloss{\emph{would give:} Salvation is received through faith}
  \item[\(\vdash\)] \(\pvar{TestimonyArguments.SolaFide}{27} \land \pvar{TestimonyArguments.SolaFide}{28}\)\gloss{\emph{at issue:} \(\pvar{TestimonyArguments.SolaFide}{27}\) Salvation is by grace, a gift and not wages owed; \(\pvar{TestimonyArguments.SolaFide}{28}\) Salvation is not by works: no work is its ground}
  \keepwithnext
\end{derivation}
\emph{No premise here rests on scripture alone.}
\par\medskip\noindent\textbf{\texttt{tridentine\allowbreak{}Premises}}\par\nobreak\noindent
Trent’s premises, given the line that carries its definition to merit. Four verses it shares with the Reformed. That Scripture does not contradict itself. What its definition rules out. And the step from merit to the denial of “not by works” (ch. 16).

It is not built from the Reformed lines. It shares four verses with them and nothing else, so building it from them would misdescribe it.

{\small\begin{verbatim}
def tridentinePremises : Line Claim → List (Formula Claim) :=
  fun line =>
    [p Claim.ephesians2_8_9, p Claim.romans3_28,
            p Claim.galatians2_16,
            p Claim.titus3_5,
            p Claim.scriptureSelfConsistent] ++
          line.premises ++
        trentDefinitionSteps ++
      [p Claim.worksMeritIncreaseOfJustification 🡒
          notP Claim.salvationNotByWorks]
\end{verbatim}}
\par\medskip\noindent\textbf{\texttt{tridentine}} --- Tridentine (Council of Trent, Session VI)\par\nobreak\noindent
The position of the Council of Trent (Session VI, 1547). Justification is the renewal of the inward person. That renewal grows through good works. So works done in grace merit an increase of justification, which denies the second part of the conclusion, “not by works”. The same definition denies that justification is only a verdict.

\begin{derivation}
  \item[(1)] \(\pvar{TestimonyArguments.SolaFide}{1}\)\gloss{Ephesians 2:8–9 teaches salvation by grace through faith, not of works}
  \item[(2)] \(\pvar{TestimonyArguments.SolaFide}{2}\)\gloss{Romans 3:28 teaches justification by faith apart from works of the law}
  \item[(3)] \(\pvar{TestimonyArguments.SolaFide}{3}\)\gloss{Galatians 2:16 teaches that no one is justified by works of the law}
  \breakrow{(4)} \(\pvar{TestimonyArguments.SolaFide}{5}\)\gloss{Titus 3:5 teaches that God saved us not by works done in righteousness}
  \breakrow{(5)} \(\pvar{TestimonyArguments.SolaFide}{24}\)\gloss{Scripture does not contradict itself}
  \breakrow{(6)} \(\pvar{TestimonyArguments.SolaFide}{33}\)\gloss{Justification is not remission of sins only, but renewal of the inward man}
  \breakrow{(7)} \(\pvar{TestimonyArguments.SolaFide}{35}\)\gloss{The inward renewal of the justified grows as they do good works in grace}
  \breakrow{(8)} \((\pvar{TestimonyArguments.SolaFide}{33} \land \pvar{TestimonyArguments.SolaFide}{35}) \rightarrow \pvar{TestimonyArguments.SolaFide}{32}\)\gloss{\emph{would give:} Works performed in grace merit an increase of justification}
  \breakrow{(9)} \(\pvar{TestimonyArguments.SolaFide}{33} \rightarrow \lnot \pvar{TestimonyArguments.SolaFide}{30}\)\gloss{\emph{would deny:} Justification is a verdict only: pardon and righteousness counted, not poured in}
  \breakrow{(10)} \(\pvar{TestimonyArguments.SolaFide}{33} \rightarrow \lnot \pvar{TestimonyArguments.SolaFide}{34}\)\gloss{\emph{would deny:} Justification and sanctification are inseparable but distinct}
  \breakrow{(11)} \(\pvar{TestimonyArguments.SolaFide}{33} \rightarrow \lnot \pvar{TestimonyArguments.SolaFide}{26}\)\gloss{\emph{would deny:} Faith is sufficient: whoever believes is saved}
  \item[(12)] \(\pvar{TestimonyArguments.SolaFide}{32} \rightarrow \lnot \pvar{TestimonyArguments.SolaFide}{28}\)\gloss{\emph{would deny:} Salvation is not by works: no work is its ground}
  \item[\(\vdash\)] \(\pvar{TestimonyArguments.SolaFide}{27} \land \pvar{TestimonyArguments.SolaFide}{28} \land \pvar{TestimonyArguments.SolaFide}{29}\)\gloss{\emph{at issue:} \(\pvar{TestimonyArguments.SolaFide}{27}\) Salvation is by grace, a gift and not wages owed; \(\pvar{TestimonyArguments.SolaFide}{28}\) Salvation is not by works: no work is its ground; \(\pvar{TestimonyArguments.SolaFide}{29}\) Salvation is received through faith}
  \keepwithnext
\end{derivation}
\emph{No premise here rests on scripture alone.}
\paragraph*{Variants}
\leavevmode\par\nobreak
Each removes named premises from the Reformed case and keeps the rest. Paul’s strand rests on two disputed readings, so it has two reduced forms.

\par\medskip\noindent\textbf{\texttt{reformed\allowbreak{}Without\allowbreak{}Works\allowbreak{}Of\allowbreak{}Law}} --- Reformed, minus the ἔργα νόμου premise\par\nobreak\noindent
The Reformed case without the reading of “works of the law”.

\begin{derivation}
  \item[(1)] \(\pvar{TestimonyArguments.SolaFide}{8}\)\gloss{Galatians opposes requiring circumcision, besides faith, for justification}
  \item[(2)] \(\pvar{TestimonyArguments.SolaFide}{7}\)\gloss{In Galatians 2:16 ‘faith of Christ’ (πίστις Χριστοῦ) means our faith in Christ}
  \item[(3)] \(\pvar{TestimonyArguments.SolaFide}{18}\)\gloss{‘Saved’ (σῴζω) at Luke 7:50 means saved from sin, not healed}
  \breakrow{(4)} \(\pvar{TestimonyArguments.SolaFide}{19}\)\gloss{Luke 7:47 — her love is the evidence of her forgiveness, not its ground}
  \breakrow{(5)} \(\pvar{TestimonyArguments.SolaFide}{10}\)\gloss{The yoke refused at Acts 15:10 is the whole law as a condition of salvation}
  \breakrow{(6)} \(\pvar{TestimonyArguments.SolaFide}{36}\)\gloss{John 6:28–29 — the work God requires is that you believe in him whom he sent}
  \breakrow{(7)} \(\pvar{TestimonyArguments.SolaFide}{37}\)\gloss{John 3:16–18, 3:36, 5:24, 20:31 — eternal life through believing in the Son}
  \breakrow{(8)} \(\pvar{TestimonyArguments.SolaFide}{40}\)\gloss{The believing of John 6:29 is trust, which brings nothing and receives Christ}
  \breakrow{(9)} \(\pvar{TestimonyArguments.SolaFide}{1}\)\gloss{Ephesians 2:8–9 teaches salvation by grace through faith, not of works}
  \breakrow{(10)} \(\pvar{TestimonyArguments.SolaFide}{2}\)\gloss{Romans 3:28 teaches justification by faith apart from works of the law}
  \breakrow{(11)} \(\pvar{TestimonyArguments.SolaFide}{3}\)\gloss{Galatians 2:16 teaches that no one is justified by works of the law}
  \breakrow{(12)} \(\pvar{TestimonyArguments.SolaFide}{4}\)\gloss{Romans 4:4–5 contrasts wages owed with a gift reckoned to the one who believes}
  \breakrow{(13)} \(\pvar{TestimonyArguments.SolaFide}{5}\)\gloss{Titus 3:5 teaches that God saved us not by works done in righteousness}
  \breakrow{(14)} \(\pvar{TestimonyArguments.SolaFide}{17}\)\gloss{Luke 7:50 — Jesus says ‘your faith has saved you’ after declaring sins forgiven}
  \breakrow{(15)} \(\pvar{TestimonyArguments.SolaFide}{9}\)\gloss{Acts 15:9–11 — Peter: hearts cleansed by faith; saved through grace}
  \breakrow{(16)} \(\pvar{TestimonyArguments.SolaFide}{20}\)\gloss{James 2:14–26 targets a barren faith — mere assent — not Paul’s doctrine}
  \breakrow{(17)} \(\pvar{TestimonyArguments.SolaFide}{21}\)\gloss{Good works are the fruit and evidence of saving faith, not its ground}
  \breakrow{(18)} \(\pvar{TestimonyArguments.SolaFide}{24}\)\gloss{Scripture does not contradict itself}
  \breakrow{(19)} \((\pvar{TestimonyArguments.SolaFide}{2} \land \pvar{TestimonyArguments.SolaFide}{3} \land \pvar{TestimonyArguments.SolaFide}{8} \land \pvar{TestimonyArguments.SolaFide}{7} \land \pvar{TestimonyArguments.SolaFide}{6}) \rightarrow \pvar{TestimonyArguments.SolaFide}{25}\)\gloss{\emph{would give:} Justification is by faith alone}
  \breakrow{(20)} \((\pvar{TestimonyArguments.SolaFide}{17} \land \pvar{TestimonyArguments.SolaFide}{18} \land \pvar{TestimonyArguments.SolaFide}{19}) \rightarrow \pvar{TestimonyArguments.SolaFide}{26}\)\gloss{\emph{would give:} Faith is sufficient: whoever believes is saved}
  \breakrow{(21)} \((\pvar{TestimonyArguments.SolaFide}{9} \land \pvar{TestimonyArguments.SolaFide}{10}) \rightarrow \pvar{TestimonyArguments.SolaFide}{26}\)\gloss{\emph{would give:} Faith is sufficient: whoever believes is saved}
  \breakrow{(22)} \((\pvar{TestimonyArguments.SolaFide}{36} \land \pvar{TestimonyArguments.SolaFide}{37} \land \pvar{TestimonyArguments.SolaFide}{40}) \rightarrow \pvar{TestimonyArguments.SolaFide}{25}\)\gloss{\emph{would give:} Justification is by faith alone}
  \breakrow{(23)} \((\pvar{TestimonyArguments.SolaFide}{20} \land \pvar{TestimonyArguments.SolaFide}{21}) \rightarrow \pvar{TestimonyArguments.SolaFide}{22}\)\gloss{\emph{would give:} James 2:24 is compatible with Paul, using ‘justify’ in a different sense}
  \breakrow{(24)} \((\pvar{TestimonyArguments.SolaFide}{1} \land \pvar{TestimonyArguments.SolaFide}{4} \land \pvar{TestimonyArguments.SolaFide}{5}) \rightarrow \pvar{TestimonyArguments.SolaFide}{27}\)\gloss{\emph{would give:} Salvation is by grace, a gift and not wages owed}
  \breakrow{(25)} \((\pvar{TestimonyArguments.SolaFide}{1} \land \pvar{TestimonyArguments.SolaFide}{4} \land \pvar{TestimonyArguments.SolaFide}{5} \land \pvar{TestimonyArguments.SolaFide}{22} \land \pvar{TestimonyArguments.SolaFide}{24}) \rightarrow \pvar{TestimonyArguments.SolaFide}{28}\)\gloss{\emph{would give:} Salvation is not by works: no work is its ground}
  \breakrow{(26)} \(\pvar{TestimonyArguments.SolaFide}{25} \rightarrow \pvar{TestimonyArguments.SolaFide}{26}\)\gloss{\emph{would give:} Faith is sufficient: whoever believes is saved}
  \item[(27)] \((\pvar{TestimonyArguments.SolaFide}{26} \land \pvar{TestimonyArguments.SolaFide}{1}) \rightarrow \pvar{TestimonyArguments.SolaFide}{29}\)\gloss{\emph{would give:} Salvation is received through faith}
  \item[\(\vdash\)] \(\pvar{TestimonyArguments.SolaFide}{27} \land \pvar{TestimonyArguments.SolaFide}{28} \land \pvar{TestimonyArguments.SolaFide}{29}\)\gloss{\emph{at issue:} \(\pvar{TestimonyArguments.SolaFide}{27}\) Salvation is by grace, a gift and not wages owed; \(\pvar{TestimonyArguments.SolaFide}{28}\) Salvation is not by works: no work is its ground; \(\pvar{TestimonyArguments.SolaFide}{29}\) Salvation is received through faith}
  \keepwithnext
\end{derivation}
\emph{Grounded in scripture alone:} John 3:16–18, 3:36, 5:24, 20:31 — eternal life through believing in the Son.
\par\medskip\noindent\textbf{\texttt{reformed\allowbreak{}Without\allowbreak{}Sozo}} --- Reformed, minus the dominical lexical premise\par\nobreak\noindent
The Reformed case without the reading of “saved” at Luke 7:50.

\begin{derivation}
  \item[(1)] \(\pvar{TestimonyArguments.SolaFide}{8}\)\gloss{Galatians opposes requiring circumcision, besides faith, for justification}
  \item[(2)] \(\pvar{TestimonyArguments.SolaFide}{6}\)\gloss{Paul’s ‘works of the law’ means human works in general, not boundary markers}
  \item[(3)] \(\pvar{TestimonyArguments.SolaFide}{7}\)\gloss{In Galatians 2:16 ‘faith of Christ’ (πίστις Χριστοῦ) means our faith in Christ}
  \breakrow{(4)} \(\pvar{TestimonyArguments.SolaFide}{19}\)\gloss{Luke 7:47 — her love is the evidence of her forgiveness, not its ground}
  \breakrow{(5)} \(\pvar{TestimonyArguments.SolaFide}{10}\)\gloss{The yoke refused at Acts 15:10 is the whole law as a condition of salvation}
  \breakrow{(6)} \(\pvar{TestimonyArguments.SolaFide}{36}\)\gloss{John 6:28–29 — the work God requires is that you believe in him whom he sent}
  \breakrow{(7)} \(\pvar{TestimonyArguments.SolaFide}{37}\)\gloss{John 3:16–18, 3:36, 5:24, 20:31 — eternal life through believing in the Son}
  \breakrow{(8)} \(\pvar{TestimonyArguments.SolaFide}{40}\)\gloss{The believing of John 6:29 is trust, which brings nothing and receives Christ}
  \breakrow{(9)} \(\pvar{TestimonyArguments.SolaFide}{1}\)\gloss{Ephesians 2:8–9 teaches salvation by grace through faith, not of works}
  \breakrow{(10)} \(\pvar{TestimonyArguments.SolaFide}{2}\)\gloss{Romans 3:28 teaches justification by faith apart from works of the law}
  \breakrow{(11)} \(\pvar{TestimonyArguments.SolaFide}{3}\)\gloss{Galatians 2:16 teaches that no one is justified by works of the law}
  \breakrow{(12)} \(\pvar{TestimonyArguments.SolaFide}{4}\)\gloss{Romans 4:4–5 contrasts wages owed with a gift reckoned to the one who believes}
  \breakrow{(13)} \(\pvar{TestimonyArguments.SolaFide}{5}\)\gloss{Titus 3:5 teaches that God saved us not by works done in righteousness}
  \breakrow{(14)} \(\pvar{TestimonyArguments.SolaFide}{17}\)\gloss{Luke 7:50 — Jesus says ‘your faith has saved you’ after declaring sins forgiven}
  \breakrow{(15)} \(\pvar{TestimonyArguments.SolaFide}{9}\)\gloss{Acts 15:9–11 — Peter: hearts cleansed by faith; saved through grace}
  \breakrow{(16)} \(\pvar{TestimonyArguments.SolaFide}{20}\)\gloss{James 2:14–26 targets a barren faith — mere assent — not Paul’s doctrine}
  \breakrow{(17)} \(\pvar{TestimonyArguments.SolaFide}{21}\)\gloss{Good works are the fruit and evidence of saving faith, not its ground}
  \breakrow{(18)} \(\pvar{TestimonyArguments.SolaFide}{24}\)\gloss{Scripture does not contradict itself}
  \breakrow{(19)} \((\pvar{TestimonyArguments.SolaFide}{2} \land \pvar{TestimonyArguments.SolaFide}{3} \land \pvar{TestimonyArguments.SolaFide}{8} \land \pvar{TestimonyArguments.SolaFide}{7} \land \pvar{TestimonyArguments.SolaFide}{6}) \rightarrow \pvar{TestimonyArguments.SolaFide}{25}\)\gloss{\emph{would give:} Justification is by faith alone}
  \breakrow{(20)} \((\pvar{TestimonyArguments.SolaFide}{17} \land \pvar{TestimonyArguments.SolaFide}{18} \land \pvar{TestimonyArguments.SolaFide}{19}) \rightarrow \pvar{TestimonyArguments.SolaFide}{26}\)\gloss{\emph{would give:} Faith is sufficient: whoever believes is saved}
  \breakrow{(21)} \((\pvar{TestimonyArguments.SolaFide}{9} \land \pvar{TestimonyArguments.SolaFide}{10}) \rightarrow \pvar{TestimonyArguments.SolaFide}{26}\)\gloss{\emph{would give:} Faith is sufficient: whoever believes is saved}
  \breakrow{(22)} \((\pvar{TestimonyArguments.SolaFide}{36} \land \pvar{TestimonyArguments.SolaFide}{37} \land \pvar{TestimonyArguments.SolaFide}{40}) \rightarrow \pvar{TestimonyArguments.SolaFide}{25}\)\gloss{\emph{would give:} Justification is by faith alone}
  \breakrow{(23)} \((\pvar{TestimonyArguments.SolaFide}{20} \land \pvar{TestimonyArguments.SolaFide}{21}) \rightarrow \pvar{TestimonyArguments.SolaFide}{22}\)\gloss{\emph{would give:} James 2:24 is compatible with Paul, using ‘justify’ in a different sense}
  \breakrow{(24)} \((\pvar{TestimonyArguments.SolaFide}{1} \land \pvar{TestimonyArguments.SolaFide}{4} \land \pvar{TestimonyArguments.SolaFide}{5}) \rightarrow \pvar{TestimonyArguments.SolaFide}{27}\)\gloss{\emph{would give:} Salvation is by grace, a gift and not wages owed}
  \breakrow{(25)} \((\pvar{TestimonyArguments.SolaFide}{1} \land \pvar{TestimonyArguments.SolaFide}{4} \land \pvar{TestimonyArguments.SolaFide}{5} \land \pvar{TestimonyArguments.SolaFide}{22} \land \pvar{TestimonyArguments.SolaFide}{24}) \rightarrow \pvar{TestimonyArguments.SolaFide}{28}\)\gloss{\emph{would give:} Salvation is not by works: no work is its ground}
  \breakrow{(26)} \(\pvar{TestimonyArguments.SolaFide}{25} \rightarrow \pvar{TestimonyArguments.SolaFide}{26}\)\gloss{\emph{would give:} Faith is sufficient: whoever believes is saved}
  \item[(27)] \((\pvar{TestimonyArguments.SolaFide}{26} \land \pvar{TestimonyArguments.SolaFide}{1}) \rightarrow \pvar{TestimonyArguments.SolaFide}{29}\)\gloss{\emph{would give:} Salvation is received through faith}
  \item[\(\vdash\)] \(\pvar{TestimonyArguments.SolaFide}{27} \land \pvar{TestimonyArguments.SolaFide}{28} \land \pvar{TestimonyArguments.SolaFide}{29}\)\gloss{\emph{at issue:} \(\pvar{TestimonyArguments.SolaFide}{27}\) Salvation is by grace, a gift and not wages owed; \(\pvar{TestimonyArguments.SolaFide}{28}\) Salvation is not by works: no work is its ground; \(\pvar{TestimonyArguments.SolaFide}{29}\) Salvation is received through faith}
  \keepwithnext
\end{derivation}
\emph{Grounded in scripture alone:} John 3:16–18, 3:36, 5:24, 20:31 — eternal life through believing in the Son.
\par\medskip\noindent\textbf{\texttt{reformed\allowbreak{}Resting\allowbreak{}On\allowbreak{}Acts}} --- Reformed, with only the Acts strand left whole\par\nobreak\noindent
The Reformed package resting on Acts alone. The disputed reading is removed from each of the other three strands: ἔργα νόμου from Paul, σῴζω from Luke, believing as trust from John. Peter’s speech is untouched.

\begin{derivation}
  \item[(1)] \(\pvar{TestimonyArguments.SolaFide}{8}\)\gloss{Galatians opposes requiring circumcision, besides faith, for justification}
  \item[(2)] \(\pvar{TestimonyArguments.SolaFide}{7}\)\gloss{In Galatians 2:16 ‘faith of Christ’ (πίστις Χριστοῦ) means our faith in Christ}
  \item[(3)] \(\pvar{TestimonyArguments.SolaFide}{19}\)\gloss{Luke 7:47 — her love is the evidence of her forgiveness, not its ground}
  \breakrow{(4)} \(\pvar{TestimonyArguments.SolaFide}{10}\)\gloss{The yoke refused at Acts 15:10 is the whole law as a condition of salvation}
  \breakrow{(5)} \(\pvar{TestimonyArguments.SolaFide}{36}\)\gloss{John 6:28–29 — the work God requires is that you believe in him whom he sent}
  \breakrow{(6)} \(\pvar{TestimonyArguments.SolaFide}{37}\)\gloss{John 3:16–18, 3:36, 5:24, 20:31 — eternal life through believing in the Son}
  \breakrow{(7)} \(\pvar{TestimonyArguments.SolaFide}{1}\)\gloss{Ephesians 2:8–9 teaches salvation by grace through faith, not of works}
  \breakrow{(8)} \(\pvar{TestimonyArguments.SolaFide}{2}\)\gloss{Romans 3:28 teaches justification by faith apart from works of the law}
  \breakrow{(9)} \(\pvar{TestimonyArguments.SolaFide}{3}\)\gloss{Galatians 2:16 teaches that no one is justified by works of the law}
  \breakrow{(10)} \(\pvar{TestimonyArguments.SolaFide}{4}\)\gloss{Romans 4:4–5 contrasts wages owed with a gift reckoned to the one who believes}
  \breakrow{(11)} \(\pvar{TestimonyArguments.SolaFide}{5}\)\gloss{Titus 3:5 teaches that God saved us not by works done in righteousness}
  \breakrow{(12)} \(\pvar{TestimonyArguments.SolaFide}{17}\)\gloss{Luke 7:50 — Jesus says ‘your faith has saved you’ after declaring sins forgiven}
  \breakrow{(13)} \(\pvar{TestimonyArguments.SolaFide}{9}\)\gloss{Acts 15:9–11 — Peter: hearts cleansed by faith; saved through grace}
  \breakrow{(14)} \(\pvar{TestimonyArguments.SolaFide}{20}\)\gloss{James 2:14–26 targets a barren faith — mere assent — not Paul’s doctrine}
  \breakrow{(15)} \(\pvar{TestimonyArguments.SolaFide}{21}\)\gloss{Good works are the fruit and evidence of saving faith, not its ground}
  \breakrow{(16)} \(\pvar{TestimonyArguments.SolaFide}{24}\)\gloss{Scripture does not contradict itself}
  \breakrow{(17)} \((\pvar{TestimonyArguments.SolaFide}{2} \land \pvar{TestimonyArguments.SolaFide}{3} \land \pvar{TestimonyArguments.SolaFide}{8} \land \pvar{TestimonyArguments.SolaFide}{7} \land \pvar{TestimonyArguments.SolaFide}{6}) \rightarrow \pvar{TestimonyArguments.SolaFide}{25}\)\gloss{\emph{would give:} Justification is by faith alone}
  \breakrow{(18)} \((\pvar{TestimonyArguments.SolaFide}{17} \land \pvar{TestimonyArguments.SolaFide}{18} \land \pvar{TestimonyArguments.SolaFide}{19}) \rightarrow \pvar{TestimonyArguments.SolaFide}{26}\)\gloss{\emph{would give:} Faith is sufficient: whoever believes is saved}
  \breakrow{(19)} \((\pvar{TestimonyArguments.SolaFide}{9} \land \pvar{TestimonyArguments.SolaFide}{10}) \rightarrow \pvar{TestimonyArguments.SolaFide}{26}\)\gloss{\emph{would give:} Faith is sufficient: whoever believes is saved}
  \breakrow{(20)} \((\pvar{TestimonyArguments.SolaFide}{36} \land \pvar{TestimonyArguments.SolaFide}{37} \land \pvar{TestimonyArguments.SolaFide}{40}) \rightarrow \pvar{TestimonyArguments.SolaFide}{25}\)\gloss{\emph{would give:} Justification is by faith alone}
  \breakrow{(21)} \((\pvar{TestimonyArguments.SolaFide}{20} \land \pvar{TestimonyArguments.SolaFide}{21}) \rightarrow \pvar{TestimonyArguments.SolaFide}{22}\)\gloss{\emph{would give:} James 2:24 is compatible with Paul, using ‘justify’ in a different sense}
  \breakrow{(22)} \((\pvar{TestimonyArguments.SolaFide}{1} \land \pvar{TestimonyArguments.SolaFide}{4} \land \pvar{TestimonyArguments.SolaFide}{5}) \rightarrow \pvar{TestimonyArguments.SolaFide}{27}\)\gloss{\emph{would give:} Salvation is by grace, a gift and not wages owed}
  \breakrow{(23)} \((\pvar{TestimonyArguments.SolaFide}{1} \land \pvar{TestimonyArguments.SolaFide}{4} \land \pvar{TestimonyArguments.SolaFide}{5} \land \pvar{TestimonyArguments.SolaFide}{22} \land \pvar{TestimonyArguments.SolaFide}{24}) \rightarrow \pvar{TestimonyArguments.SolaFide}{28}\)\gloss{\emph{would give:} Salvation is not by works: no work is its ground}
  \breakrow{(24)} \(\pvar{TestimonyArguments.SolaFide}{25} \rightarrow \pvar{TestimonyArguments.SolaFide}{26}\)\gloss{\emph{would give:} Faith is sufficient: whoever believes is saved}
  \item[(25)] \((\pvar{TestimonyArguments.SolaFide}{26} \land \pvar{TestimonyArguments.SolaFide}{1}) \rightarrow \pvar{TestimonyArguments.SolaFide}{29}\)\gloss{\emph{would give:} Salvation is received through faith}
  \item[\(\vdash\)] \(\pvar{TestimonyArguments.SolaFide}{27} \land \pvar{TestimonyArguments.SolaFide}{28} \land \pvar{TestimonyArguments.SolaFide}{29}\)\gloss{\emph{at issue:} \(\pvar{TestimonyArguments.SolaFide}{27}\) Salvation is by grace, a gift and not wages owed; \(\pvar{TestimonyArguments.SolaFide}{28}\) Salvation is not by works: no work is its ground; \(\pvar{TestimonyArguments.SolaFide}{29}\) Salvation is received through faith}
  \keepwithnext
\end{derivation}
\emph{Grounded in scripture alone:} John 3:16–18, 3:36, 5:24, 20:31 — eternal life through believing in the Son.
\par\medskip\noindent\textbf{\texttt{reformed\allowbreak{}Resting\allowbreak{}On\allowbreak{}John}} --- Reformed, with only the Johannine strand left whole\par\nobreak\noindent
The Reformed package resting on John alone. The disputed reading is removed from each of the other three strands: ἔργα νόμου from Paul, σῴζω from Luke, the yoke from Acts. John’s strand is untouched.

\begin{derivation}
  \item[(1)] \(\pvar{TestimonyArguments.SolaFide}{8}\)\gloss{Galatians opposes requiring circumcision, besides faith, for justification}
  \item[(2)] \(\pvar{TestimonyArguments.SolaFide}{7}\)\gloss{In Galatians 2:16 ‘faith of Christ’ (πίστις Χριστοῦ) means our faith in Christ}
  \item[(3)] \(\pvar{TestimonyArguments.SolaFide}{19}\)\gloss{Luke 7:47 — her love is the evidence of her forgiveness, not its ground}
  \breakrow{(4)} \(\pvar{TestimonyArguments.SolaFide}{36}\)\gloss{John 6:28–29 — the work God requires is that you believe in him whom he sent}
  \breakrow{(5)} \(\pvar{TestimonyArguments.SolaFide}{37}\)\gloss{John 3:16–18, 3:36, 5:24, 20:31 — eternal life through believing in the Son}
  \breakrow{(6)} \(\pvar{TestimonyArguments.SolaFide}{40}\)\gloss{The believing of John 6:29 is trust, which brings nothing and receives Christ}
  \breakrow{(7)} \(\pvar{TestimonyArguments.SolaFide}{1}\)\gloss{Ephesians 2:8–9 teaches salvation by grace through faith, not of works}
  \breakrow{(8)} \(\pvar{TestimonyArguments.SolaFide}{2}\)\gloss{Romans 3:28 teaches justification by faith apart from works of the law}
  \breakrow{(9)} \(\pvar{TestimonyArguments.SolaFide}{3}\)\gloss{Galatians 2:16 teaches that no one is justified by works of the law}
  \breakrow{(10)} \(\pvar{TestimonyArguments.SolaFide}{4}\)\gloss{Romans 4:4–5 contrasts wages owed with a gift reckoned to the one who believes}
  \breakrow{(11)} \(\pvar{TestimonyArguments.SolaFide}{5}\)\gloss{Titus 3:5 teaches that God saved us not by works done in righteousness}
  \breakrow{(12)} \(\pvar{TestimonyArguments.SolaFide}{17}\)\gloss{Luke 7:50 — Jesus says ‘your faith has saved you’ after declaring sins forgiven}
  \breakrow{(13)} \(\pvar{TestimonyArguments.SolaFide}{9}\)\gloss{Acts 15:9–11 — Peter: hearts cleansed by faith; saved through grace}
  \breakrow{(14)} \(\pvar{TestimonyArguments.SolaFide}{20}\)\gloss{James 2:14–26 targets a barren faith — mere assent — not Paul’s doctrine}
  \breakrow{(15)} \(\pvar{TestimonyArguments.SolaFide}{21}\)\gloss{Good works are the fruit and evidence of saving faith, not its ground}
  \breakrow{(16)} \(\pvar{TestimonyArguments.SolaFide}{24}\)\gloss{Scripture does not contradict itself}
  \breakrow{(17)} \((\pvar{TestimonyArguments.SolaFide}{2} \land \pvar{TestimonyArguments.SolaFide}{3} \land \pvar{TestimonyArguments.SolaFide}{8} \land \pvar{TestimonyArguments.SolaFide}{7} \land \pvar{TestimonyArguments.SolaFide}{6}) \rightarrow \pvar{TestimonyArguments.SolaFide}{25}\)\gloss{\emph{would give:} Justification is by faith alone}
  \breakrow{(18)} \((\pvar{TestimonyArguments.SolaFide}{17} \land \pvar{TestimonyArguments.SolaFide}{18} \land \pvar{TestimonyArguments.SolaFide}{19}) \rightarrow \pvar{TestimonyArguments.SolaFide}{26}\)\gloss{\emph{would give:} Faith is sufficient: whoever believes is saved}
  \breakrow{(19)} \((\pvar{TestimonyArguments.SolaFide}{9} \land \pvar{TestimonyArguments.SolaFide}{10}) \rightarrow \pvar{TestimonyArguments.SolaFide}{26}\)\gloss{\emph{would give:} Faith is sufficient: whoever believes is saved}
  \breakrow{(20)} \((\pvar{TestimonyArguments.SolaFide}{36} \land \pvar{TestimonyArguments.SolaFide}{37} \land \pvar{TestimonyArguments.SolaFide}{40}) \rightarrow \pvar{TestimonyArguments.SolaFide}{25}\)\gloss{\emph{would give:} Justification is by faith alone}
  \breakrow{(21)} \((\pvar{TestimonyArguments.SolaFide}{20} \land \pvar{TestimonyArguments.SolaFide}{21}) \rightarrow \pvar{TestimonyArguments.SolaFide}{22}\)\gloss{\emph{would give:} James 2:24 is compatible with Paul, using ‘justify’ in a different sense}
  \breakrow{(22)} \((\pvar{TestimonyArguments.SolaFide}{1} \land \pvar{TestimonyArguments.SolaFide}{4} \land \pvar{TestimonyArguments.SolaFide}{5}) \rightarrow \pvar{TestimonyArguments.SolaFide}{27}\)\gloss{\emph{would give:} Salvation is by grace, a gift and not wages owed}
  \breakrow{(23)} \((\pvar{TestimonyArguments.SolaFide}{1} \land \pvar{TestimonyArguments.SolaFide}{4} \land \pvar{TestimonyArguments.SolaFide}{5} \land \pvar{TestimonyArguments.SolaFide}{22} \land \pvar{TestimonyArguments.SolaFide}{24}) \rightarrow \pvar{TestimonyArguments.SolaFide}{28}\)\gloss{\emph{would give:} Salvation is not by works: no work is its ground}
  \breakrow{(24)} \(\pvar{TestimonyArguments.SolaFide}{25} \rightarrow \pvar{TestimonyArguments.SolaFide}{26}\)\gloss{\emph{would give:} Faith is sufficient: whoever believes is saved}
  \item[(25)] \((\pvar{TestimonyArguments.SolaFide}{26} \land \pvar{TestimonyArguments.SolaFide}{1}) \rightarrow \pvar{TestimonyArguments.SolaFide}{29}\)\gloss{\emph{would give:} Salvation is received through faith}
  \item[\(\vdash\)] \(\pvar{TestimonyArguments.SolaFide}{27} \land \pvar{TestimonyArguments.SolaFide}{28} \land \pvar{TestimonyArguments.SolaFide}{29}\)\gloss{\emph{at issue:} \(\pvar{TestimonyArguments.SolaFide}{27}\) Salvation is by grace, a gift and not wages owed; \(\pvar{TestimonyArguments.SolaFide}{28}\) Salvation is not by works: no work is its ground; \(\pvar{TestimonyArguments.SolaFide}{29}\) Salvation is received through faith}
  \keepwithnext
\end{derivation}
\emph{Grounded in scripture alone:} John 3:16–18, 3:36, 5:24, 20:31 — eternal life through believing in the Son.
\par\medskip\noindent\textbf{\texttt{reformed\allowbreak{}Without\allowbreak{}Every\allowbreak{}Strands\allowbreak{}Premise}} --- Reformed, minus the ἔργα νόμου, σῴζω, yoke and trust premises\par\nobreak\noindent
The Reformed package with a disputed reading removed from \textbf{every} strand: ἔργα νόμου from Paul, σῴζω from Luke, the yoke from Acts, believing as trust from John. The objective genitive stays.

\begin{derivation}
  \item[(1)] \(\pvar{TestimonyArguments.SolaFide}{8}\)\gloss{Galatians opposes requiring circumcision, besides faith, for justification}
  \item[(2)] \(\pvar{TestimonyArguments.SolaFide}{7}\)\gloss{In Galatians 2:16 ‘faith of Christ’ (πίστις Χριστοῦ) means our faith in Christ}
  \item[(3)] \(\pvar{TestimonyArguments.SolaFide}{19}\)\gloss{Luke 7:47 — her love is the evidence of her forgiveness, not its ground}
  \breakrow{(4)} \(\pvar{TestimonyArguments.SolaFide}{36}\)\gloss{John 6:28–29 — the work God requires is that you believe in him whom he sent}
  \breakrow{(5)} \(\pvar{TestimonyArguments.SolaFide}{37}\)\gloss{John 3:16–18, 3:36, 5:24, 20:31 — eternal life through believing in the Son}
  \breakrow{(6)} \(\pvar{TestimonyArguments.SolaFide}{1}\)\gloss{Ephesians 2:8–9 teaches salvation by grace through faith, not of works}
  \breakrow{(7)} \(\pvar{TestimonyArguments.SolaFide}{2}\)\gloss{Romans 3:28 teaches justification by faith apart from works of the law}
  \breakrow{(8)} \(\pvar{TestimonyArguments.SolaFide}{3}\)\gloss{Galatians 2:16 teaches that no one is justified by works of the law}
  \breakrow{(9)} \(\pvar{TestimonyArguments.SolaFide}{4}\)\gloss{Romans 4:4–5 contrasts wages owed with a gift reckoned to the one who believes}
  \breakrow{(10)} \(\pvar{TestimonyArguments.SolaFide}{5}\)\gloss{Titus 3:5 teaches that God saved us not by works done in righteousness}
  \breakrow{(11)} \(\pvar{TestimonyArguments.SolaFide}{17}\)\gloss{Luke 7:50 — Jesus says ‘your faith has saved you’ after declaring sins forgiven}
  \breakrow{(12)} \(\pvar{TestimonyArguments.SolaFide}{9}\)\gloss{Acts 15:9–11 — Peter: hearts cleansed by faith; saved through grace}
  \breakrow{(13)} \(\pvar{TestimonyArguments.SolaFide}{20}\)\gloss{James 2:14–26 targets a barren faith — mere assent — not Paul’s doctrine}
  \breakrow{(14)} \(\pvar{TestimonyArguments.SolaFide}{21}\)\gloss{Good works are the fruit and evidence of saving faith, not its ground}
  \breakrow{(15)} \(\pvar{TestimonyArguments.SolaFide}{24}\)\gloss{Scripture does not contradict itself}
  \breakrow{(16)} \((\pvar{TestimonyArguments.SolaFide}{2} \land \pvar{TestimonyArguments.SolaFide}{3} \land \pvar{TestimonyArguments.SolaFide}{8} \land \pvar{TestimonyArguments.SolaFide}{7} \land \pvar{TestimonyArguments.SolaFide}{6}) \rightarrow \pvar{TestimonyArguments.SolaFide}{25}\)\gloss{\emph{would give:} Justification is by faith alone}
  \breakrow{(17)} \((\pvar{TestimonyArguments.SolaFide}{17} \land \pvar{TestimonyArguments.SolaFide}{18} \land \pvar{TestimonyArguments.SolaFide}{19}) \rightarrow \pvar{TestimonyArguments.SolaFide}{26}\)\gloss{\emph{would give:} Faith is sufficient: whoever believes is saved}
  \breakrow{(18)} \((\pvar{TestimonyArguments.SolaFide}{9} \land \pvar{TestimonyArguments.SolaFide}{10}) \rightarrow \pvar{TestimonyArguments.SolaFide}{26}\)\gloss{\emph{would give:} Faith is sufficient: whoever believes is saved}
  \breakrow{(19)} \((\pvar{TestimonyArguments.SolaFide}{36} \land \pvar{TestimonyArguments.SolaFide}{37} \land \pvar{TestimonyArguments.SolaFide}{40}) \rightarrow \pvar{TestimonyArguments.SolaFide}{25}\)\gloss{\emph{would give:} Justification is by faith alone}
  \breakrow{(20)} \((\pvar{TestimonyArguments.SolaFide}{20} \land \pvar{TestimonyArguments.SolaFide}{21}) \rightarrow \pvar{TestimonyArguments.SolaFide}{22}\)\gloss{\emph{would give:} James 2:24 is compatible with Paul, using ‘justify’ in a different sense}
  \breakrow{(21)} \((\pvar{TestimonyArguments.SolaFide}{1} \land \pvar{TestimonyArguments.SolaFide}{4} \land \pvar{TestimonyArguments.SolaFide}{5}) \rightarrow \pvar{TestimonyArguments.SolaFide}{27}\)\gloss{\emph{would give:} Salvation is by grace, a gift and not wages owed}
  \breakrow{(22)} \((\pvar{TestimonyArguments.SolaFide}{1} \land \pvar{TestimonyArguments.SolaFide}{4} \land \pvar{TestimonyArguments.SolaFide}{5} \land \pvar{TestimonyArguments.SolaFide}{22} \land \pvar{TestimonyArguments.SolaFide}{24}) \rightarrow \pvar{TestimonyArguments.SolaFide}{28}\)\gloss{\emph{would give:} Salvation is not by works: no work is its ground}
  \breakrow{(23)} \(\pvar{TestimonyArguments.SolaFide}{25} \rightarrow \pvar{TestimonyArguments.SolaFide}{26}\)\gloss{\emph{would give:} Faith is sufficient: whoever believes is saved}
  \item[(24)] \((\pvar{TestimonyArguments.SolaFide}{26} \land \pvar{TestimonyArguments.SolaFide}{1}) \rightarrow \pvar{TestimonyArguments.SolaFide}{29}\)\gloss{\emph{would give:} Salvation is received through faith}
  \item[\(\vdash\)] \(\pvar{TestimonyArguments.SolaFide}{27} \land \pvar{TestimonyArguments.SolaFide}{28} \land \pvar{TestimonyArguments.SolaFide}{29}\)\gloss{\emph{at issue:} \(\pvar{TestimonyArguments.SolaFide}{27}\) Salvation is by grace, a gift and not wages owed; \(\pvar{TestimonyArguments.SolaFide}{28}\) Salvation is not by works: no work is its ground; \(\pvar{TestimonyArguments.SolaFide}{29}\) Salvation is received through faith}
  \keepwithnext
\end{derivation}
\emph{Grounded in scripture alone:} John 3:16–18, 3:36, 5:24, 20:31 — eternal life through believing in the Son.
\par\medskip\noindent\textbf{\texttt{reformed\allowbreak{}Without\allowbreak{}Pistis\allowbreak{}Christou\allowbreak{}Sozo\allowbreak{}Or\allowbreak{}Yoke}} --- Reformed, minus the πίστις Χριστοῦ, σῴζω, yoke and trust premises\par\nobreak\noindent
The same with the objective genitive removed in place of ἔργα νόμου: the Pauline strand loses its other disputed reading, and the other three strands theirs. The name lists Luke’s and Peter’s; John’s reading of believing as trust is removed too.

\begin{derivation}
  \item[(1)] \(\pvar{TestimonyArguments.SolaFide}{8}\)\gloss{Galatians opposes requiring circumcision, besides faith, for justification}
  \item[(2)] \(\pvar{TestimonyArguments.SolaFide}{6}\)\gloss{Paul’s ‘works of the law’ means human works in general, not boundary markers}
  \item[(3)] \(\pvar{TestimonyArguments.SolaFide}{19}\)\gloss{Luke 7:47 — her love is the evidence of her forgiveness, not its ground}
  \breakrow{(4)} \(\pvar{TestimonyArguments.SolaFide}{36}\)\gloss{John 6:28–29 — the work God requires is that you believe in him whom he sent}
  \breakrow{(5)} \(\pvar{TestimonyArguments.SolaFide}{37}\)\gloss{John 3:16–18, 3:36, 5:24, 20:31 — eternal life through believing in the Son}
  \breakrow{(6)} \(\pvar{TestimonyArguments.SolaFide}{1}\)\gloss{Ephesians 2:8–9 teaches salvation by grace through faith, not of works}
  \breakrow{(7)} \(\pvar{TestimonyArguments.SolaFide}{2}\)\gloss{Romans 3:28 teaches justification by faith apart from works of the law}
  \breakrow{(8)} \(\pvar{TestimonyArguments.SolaFide}{3}\)\gloss{Galatians 2:16 teaches that no one is justified by works of the law}
  \breakrow{(9)} \(\pvar{TestimonyArguments.SolaFide}{4}\)\gloss{Romans 4:4–5 contrasts wages owed with a gift reckoned to the one who believes}
  \breakrow{(10)} \(\pvar{TestimonyArguments.SolaFide}{5}\)\gloss{Titus 3:5 teaches that God saved us not by works done in righteousness}
  \breakrow{(11)} \(\pvar{TestimonyArguments.SolaFide}{17}\)\gloss{Luke 7:50 — Jesus says ‘your faith has saved you’ after declaring sins forgiven}
  \breakrow{(12)} \(\pvar{TestimonyArguments.SolaFide}{9}\)\gloss{Acts 15:9–11 — Peter: hearts cleansed by faith; saved through grace}
  \breakrow{(13)} \(\pvar{TestimonyArguments.SolaFide}{20}\)\gloss{James 2:14–26 targets a barren faith — mere assent — not Paul’s doctrine}
  \breakrow{(14)} \(\pvar{TestimonyArguments.SolaFide}{21}\)\gloss{Good works are the fruit and evidence of saving faith, not its ground}
  \breakrow{(15)} \(\pvar{TestimonyArguments.SolaFide}{24}\)\gloss{Scripture does not contradict itself}
  \breakrow{(16)} \((\pvar{TestimonyArguments.SolaFide}{2} \land \pvar{TestimonyArguments.SolaFide}{3} \land \pvar{TestimonyArguments.SolaFide}{8} \land \pvar{TestimonyArguments.SolaFide}{7} \land \pvar{TestimonyArguments.SolaFide}{6}) \rightarrow \pvar{TestimonyArguments.SolaFide}{25}\)\gloss{\emph{would give:} Justification is by faith alone}
  \breakrow{(17)} \((\pvar{TestimonyArguments.SolaFide}{17} \land \pvar{TestimonyArguments.SolaFide}{18} \land \pvar{TestimonyArguments.SolaFide}{19}) \rightarrow \pvar{TestimonyArguments.SolaFide}{26}\)\gloss{\emph{would give:} Faith is sufficient: whoever believes is saved}
  \breakrow{(18)} \((\pvar{TestimonyArguments.SolaFide}{9} \land \pvar{TestimonyArguments.SolaFide}{10}) \rightarrow \pvar{TestimonyArguments.SolaFide}{26}\)\gloss{\emph{would give:} Faith is sufficient: whoever believes is saved}
  \breakrow{(19)} \((\pvar{TestimonyArguments.SolaFide}{36} \land \pvar{TestimonyArguments.SolaFide}{37} \land \pvar{TestimonyArguments.SolaFide}{40}) \rightarrow \pvar{TestimonyArguments.SolaFide}{25}\)\gloss{\emph{would give:} Justification is by faith alone}
  \breakrow{(20)} \((\pvar{TestimonyArguments.SolaFide}{20} \land \pvar{TestimonyArguments.SolaFide}{21}) \rightarrow \pvar{TestimonyArguments.SolaFide}{22}\)\gloss{\emph{would give:} James 2:24 is compatible with Paul, using ‘justify’ in a different sense}
  \breakrow{(21)} \((\pvar{TestimonyArguments.SolaFide}{1} \land \pvar{TestimonyArguments.SolaFide}{4} \land \pvar{TestimonyArguments.SolaFide}{5}) \rightarrow \pvar{TestimonyArguments.SolaFide}{27}\)\gloss{\emph{would give:} Salvation is by grace, a gift and not wages owed}
  \breakrow{(22)} \((\pvar{TestimonyArguments.SolaFide}{1} \land \pvar{TestimonyArguments.SolaFide}{4} \land \pvar{TestimonyArguments.SolaFide}{5} \land \pvar{TestimonyArguments.SolaFide}{22} \land \pvar{TestimonyArguments.SolaFide}{24}) \rightarrow \pvar{TestimonyArguments.SolaFide}{28}\)\gloss{\emph{would give:} Salvation is not by works: no work is its ground}
  \breakrow{(23)} \(\pvar{TestimonyArguments.SolaFide}{25} \rightarrow \pvar{TestimonyArguments.SolaFide}{26}\)\gloss{\emph{would give:} Faith is sufficient: whoever believes is saved}
  \item[(24)] \((\pvar{TestimonyArguments.SolaFide}{26} \land \pvar{TestimonyArguments.SolaFide}{1}) \rightarrow \pvar{TestimonyArguments.SolaFide}{29}\)\gloss{\emph{would give:} Salvation is received through faith}
  \item[\(\vdash\)] \(\pvar{TestimonyArguments.SolaFide}{27} \land \pvar{TestimonyArguments.SolaFide}{28} \land \pvar{TestimonyArguments.SolaFide}{29}\)\gloss{\emph{at issue:} \(\pvar{TestimonyArguments.SolaFide}{27}\) Salvation is by grace, a gift and not wages owed; \(\pvar{TestimonyArguments.SolaFide}{28}\) Salvation is not by works: no work is its ground; \(\pvar{TestimonyArguments.SolaFide}{29}\) Salvation is received through faith}
  \keepwithnext
\end{derivation}
\emph{Grounded in scripture alone:} John 3:16–18, 3:36, 5:24, 20:31 — eternal life through believing in the Son.
\par\medskip\noindent\textbf{\texttt{grace\allowbreak{}And\allowbreak{}Works\allowbreak{}Without\allowbreak{}Any\allowbreak{}Lexical\allowbreak{}Premise}} --- Grace and works, minus every lexical premise\par\nobreak\noindent
The Reformed case with \textbf{every} disputed reading removed: both of Paul’s, Luke’s “saved”, the yoke of Acts 15 and believing as trust in John. It is asked only about grace and works. The reading of Galatians as a letter against requiring circumcision stays, and so does the answer to James.

\begin{derivation}
  \item[(1)] \(\pvar{TestimonyArguments.SolaFide}{8}\)\gloss{Galatians opposes requiring circumcision, besides faith, for justification}
  \item[(2)] \(\pvar{TestimonyArguments.SolaFide}{19}\)\gloss{Luke 7:47 — her love is the evidence of her forgiveness, not its ground}
  \item[(3)] \(\pvar{TestimonyArguments.SolaFide}{36}\)\gloss{John 6:28–29 — the work God requires is that you believe in him whom he sent}
  \breakrow{(4)} \(\pvar{TestimonyArguments.SolaFide}{37}\)\gloss{John 3:16–18, 3:36, 5:24, 20:31 — eternal life through believing in the Son}
  \breakrow{(5)} \(\pvar{TestimonyArguments.SolaFide}{1}\)\gloss{Ephesians 2:8–9 teaches salvation by grace through faith, not of works}
  \breakrow{(6)} \(\pvar{TestimonyArguments.SolaFide}{2}\)\gloss{Romans 3:28 teaches justification by faith apart from works of the law}
  \breakrow{(7)} \(\pvar{TestimonyArguments.SolaFide}{3}\)\gloss{Galatians 2:16 teaches that no one is justified by works of the law}
  \breakrow{(8)} \(\pvar{TestimonyArguments.SolaFide}{4}\)\gloss{Romans 4:4–5 contrasts wages owed with a gift reckoned to the one who believes}
  \breakrow{(9)} \(\pvar{TestimonyArguments.SolaFide}{5}\)\gloss{Titus 3:5 teaches that God saved us not by works done in righteousness}
  \breakrow{(10)} \(\pvar{TestimonyArguments.SolaFide}{17}\)\gloss{Luke 7:50 — Jesus says ‘your faith has saved you’ after declaring sins forgiven}
  \breakrow{(11)} \(\pvar{TestimonyArguments.SolaFide}{9}\)\gloss{Acts 15:9–11 — Peter: hearts cleansed by faith; saved through grace}
  \breakrow{(12)} \(\pvar{TestimonyArguments.SolaFide}{20}\)\gloss{James 2:14–26 targets a barren faith — mere assent — not Paul’s doctrine}
  \breakrow{(13)} \(\pvar{TestimonyArguments.SolaFide}{21}\)\gloss{Good works are the fruit and evidence of saving faith, not its ground}
  \breakrow{(14)} \(\pvar{TestimonyArguments.SolaFide}{24}\)\gloss{Scripture does not contradict itself}
  \breakrow{(15)} \((\pvar{TestimonyArguments.SolaFide}{2} \land \pvar{TestimonyArguments.SolaFide}{3} \land \pvar{TestimonyArguments.SolaFide}{8} \land \pvar{TestimonyArguments.SolaFide}{7} \land \pvar{TestimonyArguments.SolaFide}{6}) \rightarrow \pvar{TestimonyArguments.SolaFide}{25}\)\gloss{\emph{would give:} Justification is by faith alone}
  \breakrow{(16)} \((\pvar{TestimonyArguments.SolaFide}{17} \land \pvar{TestimonyArguments.SolaFide}{18} \land \pvar{TestimonyArguments.SolaFide}{19}) \rightarrow \pvar{TestimonyArguments.SolaFide}{26}\)\gloss{\emph{would give:} Faith is sufficient: whoever believes is saved}
  \breakrow{(17)} \((\pvar{TestimonyArguments.SolaFide}{9} \land \pvar{TestimonyArguments.SolaFide}{10}) \rightarrow \pvar{TestimonyArguments.SolaFide}{26}\)\gloss{\emph{would give:} Faith is sufficient: whoever believes is saved}
  \breakrow{(18)} \((\pvar{TestimonyArguments.SolaFide}{36} \land \pvar{TestimonyArguments.SolaFide}{37} \land \pvar{TestimonyArguments.SolaFide}{40}) \rightarrow \pvar{TestimonyArguments.SolaFide}{25}\)\gloss{\emph{would give:} Justification is by faith alone}
  \breakrow{(19)} \((\pvar{TestimonyArguments.SolaFide}{20} \land \pvar{TestimonyArguments.SolaFide}{21}) \rightarrow \pvar{TestimonyArguments.SolaFide}{22}\)\gloss{\emph{would give:} James 2:24 is compatible with Paul, using ‘justify’ in a different sense}
  \breakrow{(20)} \((\pvar{TestimonyArguments.SolaFide}{1} \land \pvar{TestimonyArguments.SolaFide}{4} \land \pvar{TestimonyArguments.SolaFide}{5}) \rightarrow \pvar{TestimonyArguments.SolaFide}{27}\)\gloss{\emph{would give:} Salvation is by grace, a gift and not wages owed}
  \breakrow{(21)} \((\pvar{TestimonyArguments.SolaFide}{1} \land \pvar{TestimonyArguments.SolaFide}{4} \land \pvar{TestimonyArguments.SolaFide}{5} \land \pvar{TestimonyArguments.SolaFide}{22} \land \pvar{TestimonyArguments.SolaFide}{24}) \rightarrow \pvar{TestimonyArguments.SolaFide}{28}\)\gloss{\emph{would give:} Salvation is not by works: no work is its ground}
  \breakrow{(22)} \(\pvar{TestimonyArguments.SolaFide}{25} \rightarrow \pvar{TestimonyArguments.SolaFide}{26}\)\gloss{\emph{would give:} Faith is sufficient: whoever believes is saved}
  \item[(23)] \((\pvar{TestimonyArguments.SolaFide}{26} \land \pvar{TestimonyArguments.SolaFide}{1}) \rightarrow \pvar{TestimonyArguments.SolaFide}{29}\)\gloss{\emph{would give:} Salvation is received through faith}
  \item[\(\vdash\)] \(\pvar{TestimonyArguments.SolaFide}{27} \land \pvar{TestimonyArguments.SolaFide}{28}\)\gloss{\emph{at issue:} \(\pvar{TestimonyArguments.SolaFide}{27}\) Salvation is by grace, a gift and not wages owed; \(\pvar{TestimonyArguments.SolaFide}{28}\) Salvation is not by works: no work is its ground}
  \keepwithnext
\end{derivation}
\emph{Grounded in scripture alone:} John 3:16–18, 3:36, 5:24, 20:31 — eternal life through believing in the Son.
\par\medskip\noindent\textbf{\texttt{pauline\allowbreak{}Strand\allowbreak{}On\allowbreak{}The\allowbreak{}Critics}} --- Pauline strand alone, on variegated nomism\par\nobreak\noindent
Paul’s strand alone, on the critics’ account of Judaism (Gathercole; Carson, O’Brien and Seifrid). The reading of “works of the law” is derived from their denial of Sanders’ account, not assumed, and the other three strands are left out. The question: on that history, does Paul alone reach sola fide?

\begin{derivation}
  \item[(1)] \(\pvar{TestimonyArguments.SolaFide}{8}\)\gloss{Galatians opposes requiring circumcision, besides faith, for justification}
  \item[(2)] \(\pvar{TestimonyArguments.SolaFide}{7}\)\gloss{In Galatians 2:16 ‘faith of Christ’ (πίστις Χριστοῦ) means our faith in Christ}
  \item[(3)] \(\lnot \pvar{TestimonyArguments.SolaFide}{31}\)\gloss{\emph{denied:} Second Temple Judaism: in by grace, staying in by works}
  \breakrow{(4)} \(\pvar{TestimonyArguments.SolaFide}{8}\)\gloss{Galatians opposes requiring circumcision, besides faith, for justification}
  \breakrow{(5)} \(\pvar{TestimonyArguments.SolaFide}{1}\)\gloss{Ephesians 2:8–9 teaches salvation by grace through faith, not of works}
  \breakrow{(6)} \(\pvar{TestimonyArguments.SolaFide}{2}\)\gloss{Romans 3:28 teaches justification by faith apart from works of the law}
  \breakrow{(7)} \(\pvar{TestimonyArguments.SolaFide}{3}\)\gloss{Galatians 2:16 teaches that no one is justified by works of the law}
  \breakrow{(8)} \(\pvar{TestimonyArguments.SolaFide}{4}\)\gloss{Romans 4:4–5 contrasts wages owed with a gift reckoned to the one who believes}
  \breakrow{(9)} \(\pvar{TestimonyArguments.SolaFide}{5}\)\gloss{Titus 3:5 teaches that God saved us not by works done in righteousness}
  \breakrow{(10)} \(\pvar{TestimonyArguments.SolaFide}{17}\)\gloss{Luke 7:50 — Jesus says ‘your faith has saved you’ after declaring sins forgiven}
  \breakrow{(11)} \(\pvar{TestimonyArguments.SolaFide}{9}\)\gloss{Acts 15:9–11 — Peter: hearts cleansed by faith; saved through grace}
  \breakrow{(12)} \(\pvar{TestimonyArguments.SolaFide}{20}\)\gloss{James 2:14–26 targets a barren faith — mere assent — not Paul’s doctrine}
  \breakrow{(13)} \(\pvar{TestimonyArguments.SolaFide}{21}\)\gloss{Good works are the fruit and evidence of saving faith, not its ground}
  \breakrow{(14)} \(\pvar{TestimonyArguments.SolaFide}{24}\)\gloss{Scripture does not contradict itself}
  \breakrow{(15)} \((\pvar{TestimonyArguments.SolaFide}{2} \land \pvar{TestimonyArguments.SolaFide}{3} \land \pvar{TestimonyArguments.SolaFide}{8} \land \pvar{TestimonyArguments.SolaFide}{7} \land \pvar{TestimonyArguments.SolaFide}{6}) \rightarrow \pvar{TestimonyArguments.SolaFide}{25}\)\gloss{\emph{would give:} Justification is by faith alone}
  \breakrow{(16)} \((\lnot \pvar{TestimonyArguments.SolaFide}{31} \land \pvar{TestimonyArguments.SolaFide}{8}) \rightarrow \pvar{TestimonyArguments.SolaFide}{6}\)\gloss{\emph{would give:} Paul’s ‘works of the law’ means human works in general, not boundary markers}
  \breakrow{(17)} \((\pvar{TestimonyArguments.SolaFide}{20} \land \pvar{TestimonyArguments.SolaFide}{21}) \rightarrow \pvar{TestimonyArguments.SolaFide}{22}\)\gloss{\emph{would give:} James 2:24 is compatible with Paul, using ‘justify’ in a different sense}
  \breakrow{(18)} \((\pvar{TestimonyArguments.SolaFide}{1} \land \pvar{TestimonyArguments.SolaFide}{4} \land \pvar{TestimonyArguments.SolaFide}{5}) \rightarrow \pvar{TestimonyArguments.SolaFide}{27}\)\gloss{\emph{would give:} Salvation is by grace, a gift and not wages owed}
  \breakrow{(19)} \((\pvar{TestimonyArguments.SolaFide}{1} \land \pvar{TestimonyArguments.SolaFide}{4} \land \pvar{TestimonyArguments.SolaFide}{5} \land \pvar{TestimonyArguments.SolaFide}{22} \land \pvar{TestimonyArguments.SolaFide}{24}) \rightarrow \pvar{TestimonyArguments.SolaFide}{28}\)\gloss{\emph{would give:} Salvation is not by works: no work is its ground}
  \breakrow{(20)} \(\pvar{TestimonyArguments.SolaFide}{25} \rightarrow \pvar{TestimonyArguments.SolaFide}{26}\)\gloss{\emph{would give:} Faith is sufficient: whoever believes is saved}
  \item[(21)] \((\pvar{TestimonyArguments.SolaFide}{26} \land \pvar{TestimonyArguments.SolaFide}{1}) \rightarrow \pvar{TestimonyArguments.SolaFide}{29}\)\gloss{\emph{would give:} Salvation is received through faith}
  \item[\(\vdash\)] \(\pvar{TestimonyArguments.SolaFide}{27} \land \pvar{TestimonyArguments.SolaFide}{28} \land \pvar{TestimonyArguments.SolaFide}{29}\)\gloss{\emph{at issue:} \(\pvar{TestimonyArguments.SolaFide}{27}\) Salvation is by grace, a gift and not wages owed; \(\pvar{TestimonyArguments.SolaFide}{28}\) Salvation is not by works: no work is its ground; \(\pvar{TestimonyArguments.SolaFide}{29}\) Salvation is received through faith}
  \keepwithnext
\end{derivation}
\emph{No premise here rests on scripture alone.}
\par\medskip\noindent\textbf{\texttt{pauline\allowbreak{}Strand\allowbreak{}On\allowbreak{}Sanders}} --- Pauline strand alone, on covenantal nomism\par\nobreak\noindent
The same on Sanders’ account of Judaism.

\begin{derivation}
  \item[(1)] \(\pvar{TestimonyArguments.SolaFide}{8}\)\gloss{Galatians opposes requiring circumcision, besides faith, for justification}
  \item[(2)] \(\pvar{TestimonyArguments.SolaFide}{7}\)\gloss{In Galatians 2:16 ‘faith of Christ’ (πίστις Χριστοῦ) means our faith in Christ}
  \item[(3)] \(\pvar{TestimonyArguments.SolaFide}{31}\)\gloss{Second Temple Judaism: in by grace, staying in by works}
  \breakrow{(4)} \(\pvar{TestimonyArguments.SolaFide}{1}\)\gloss{Ephesians 2:8–9 teaches salvation by grace through faith, not of works}
  \breakrow{(5)} \(\pvar{TestimonyArguments.SolaFide}{2}\)\gloss{Romans 3:28 teaches justification by faith apart from works of the law}
  \breakrow{(6)} \(\pvar{TestimonyArguments.SolaFide}{3}\)\gloss{Galatians 2:16 teaches that no one is justified by works of the law}
  \breakrow{(7)} \(\pvar{TestimonyArguments.SolaFide}{4}\)\gloss{Romans 4:4–5 contrasts wages owed with a gift reckoned to the one who believes}
  \breakrow{(8)} \(\pvar{TestimonyArguments.SolaFide}{5}\)\gloss{Titus 3:5 teaches that God saved us not by works done in righteousness}
  \breakrow{(9)} \(\pvar{TestimonyArguments.SolaFide}{17}\)\gloss{Luke 7:50 — Jesus says ‘your faith has saved you’ after declaring sins forgiven}
  \breakrow{(10)} \(\pvar{TestimonyArguments.SolaFide}{9}\)\gloss{Acts 15:9–11 — Peter: hearts cleansed by faith; saved through grace}
  \breakrow{(11)} \(\pvar{TestimonyArguments.SolaFide}{20}\)\gloss{James 2:14–26 targets a barren faith — mere assent — not Paul’s doctrine}
  \breakrow{(12)} \(\pvar{TestimonyArguments.SolaFide}{21}\)\gloss{Good works are the fruit and evidence of saving faith, not its ground}
  \breakrow{(13)} \(\pvar{TestimonyArguments.SolaFide}{24}\)\gloss{Scripture does not contradict itself}
  \breakrow{(14)} \((\pvar{TestimonyArguments.SolaFide}{2} \land \pvar{TestimonyArguments.SolaFide}{3} \land \pvar{TestimonyArguments.SolaFide}{8} \land \pvar{TestimonyArguments.SolaFide}{7} \land \pvar{TestimonyArguments.SolaFide}{6}) \rightarrow \pvar{TestimonyArguments.SolaFide}{25}\)\gloss{\emph{would give:} Justification is by faith alone}
  \breakrow{(15)} \(\pvar{TestimonyArguments.SolaFide}{31} \rightarrow \lnot \pvar{TestimonyArguments.SolaFide}{6}\)\gloss{\emph{would deny:} Paul’s ‘works of the law’ means human works in general, not boundary markers}
  \breakrow{(16)} \((\pvar{TestimonyArguments.SolaFide}{20} \land \pvar{TestimonyArguments.SolaFide}{21}) \rightarrow \pvar{TestimonyArguments.SolaFide}{22}\)\gloss{\emph{would give:} James 2:24 is compatible with Paul, using ‘justify’ in a different sense}
  \breakrow{(17)} \((\pvar{TestimonyArguments.SolaFide}{1} \land \pvar{TestimonyArguments.SolaFide}{4} \land \pvar{TestimonyArguments.SolaFide}{5}) \rightarrow \pvar{TestimonyArguments.SolaFide}{27}\)\gloss{\emph{would give:} Salvation is by grace, a gift and not wages owed}
  \breakrow{(18)} \((\pvar{TestimonyArguments.SolaFide}{1} \land \pvar{TestimonyArguments.SolaFide}{4} \land \pvar{TestimonyArguments.SolaFide}{5} \land \pvar{TestimonyArguments.SolaFide}{22} \land \pvar{TestimonyArguments.SolaFide}{24}) \rightarrow \pvar{TestimonyArguments.SolaFide}{28}\)\gloss{\emph{would give:} Salvation is not by works: no work is its ground}
  \breakrow{(19)} \(\pvar{TestimonyArguments.SolaFide}{25} \rightarrow \pvar{TestimonyArguments.SolaFide}{26}\)\gloss{\emph{would give:} Faith is sufficient: whoever believes is saved}
  \item[(20)] \((\pvar{TestimonyArguments.SolaFide}{26} \land \pvar{TestimonyArguments.SolaFide}{1}) \rightarrow \pvar{TestimonyArguments.SolaFide}{29}\)\gloss{\emph{would give:} Salvation is received through faith}
  \item[\(\vdash\)] \(\pvar{TestimonyArguments.SolaFide}{27} \land \pvar{TestimonyArguments.SolaFide}{28} \land \pvar{TestimonyArguments.SolaFide}{29}\)\gloss{\emph{at issue:} \(\pvar{TestimonyArguments.SolaFide}{27}\) Salvation is by grace, a gift and not wages owed; \(\pvar{TestimonyArguments.SolaFide}{28}\) Salvation is not by works: no work is its ground; \(\pvar{TestimonyArguments.SolaFide}{29}\) Salvation is received through faith}
  \keepwithnext
\end{derivation}
\emph{No premise here rests on scripture alone.}
\par\medskip\noindent\textbf{\texttt{reformed\allowbreak{}Without\allowbreak{}James\allowbreak{}Harmonisation}} --- Reformed, minus the James harmonisation\par\nobreak\noindent
The Reformed case without the premises that answer James.

\begin{derivation}
  \item[(1)] \(\pvar{TestimonyArguments.SolaFide}{8}\)\gloss{Galatians opposes requiring circumcision, besides faith, for justification}
  \item[(2)] \(\pvar{TestimonyArguments.SolaFide}{6}\)\gloss{Paul’s ‘works of the law’ means human works in general, not boundary markers}
  \item[(3)] \(\pvar{TestimonyArguments.SolaFide}{7}\)\gloss{In Galatians 2:16 ‘faith of Christ’ (πίστις Χριστοῦ) means our faith in Christ}
  \breakrow{(4)} \(\pvar{TestimonyArguments.SolaFide}{18}\)\gloss{‘Saved’ (σῴζω) at Luke 7:50 means saved from sin, not healed}
  \breakrow{(5)} \(\pvar{TestimonyArguments.SolaFide}{19}\)\gloss{Luke 7:47 — her love is the evidence of her forgiveness, not its ground}
  \breakrow{(6)} \(\pvar{TestimonyArguments.SolaFide}{10}\)\gloss{The yoke refused at Acts 15:10 is the whole law as a condition of salvation}
  \breakrow{(7)} \(\pvar{TestimonyArguments.SolaFide}{36}\)\gloss{John 6:28–29 — the work God requires is that you believe in him whom he sent}
  \breakrow{(8)} \(\pvar{TestimonyArguments.SolaFide}{37}\)\gloss{John 3:16–18, 3:36, 5:24, 20:31 — eternal life through believing in the Son}
  \breakrow{(9)} \(\pvar{TestimonyArguments.SolaFide}{40}\)\gloss{The believing of John 6:29 is trust, which brings nothing and receives Christ}
  \breakrow{(10)} \(\pvar{TestimonyArguments.SolaFide}{1}\)\gloss{Ephesians 2:8–9 teaches salvation by grace through faith, not of works}
  \breakrow{(11)} \(\pvar{TestimonyArguments.SolaFide}{2}\)\gloss{Romans 3:28 teaches justification by faith apart from works of the law}
  \breakrow{(12)} \(\pvar{TestimonyArguments.SolaFide}{3}\)\gloss{Galatians 2:16 teaches that no one is justified by works of the law}
  \breakrow{(13)} \(\pvar{TestimonyArguments.SolaFide}{4}\)\gloss{Romans 4:4–5 contrasts wages owed with a gift reckoned to the one who believes}
  \breakrow{(14)} \(\pvar{TestimonyArguments.SolaFide}{5}\)\gloss{Titus 3:5 teaches that God saved us not by works done in righteousness}
  \breakrow{(15)} \(\pvar{TestimonyArguments.SolaFide}{17}\)\gloss{Luke 7:50 — Jesus says ‘your faith has saved you’ after declaring sins forgiven}
  \breakrow{(16)} \(\pvar{TestimonyArguments.SolaFide}{9}\)\gloss{Acts 15:9–11 — Peter: hearts cleansed by faith; saved through grace}
  \breakrow{(17)} \(\pvar{TestimonyArguments.SolaFide}{24}\)\gloss{Scripture does not contradict itself}
  \breakrow{(18)} \((\pvar{TestimonyArguments.SolaFide}{2} \land \pvar{TestimonyArguments.SolaFide}{3} \land \pvar{TestimonyArguments.SolaFide}{8} \land \pvar{TestimonyArguments.SolaFide}{7} \land \pvar{TestimonyArguments.SolaFide}{6}) \rightarrow \pvar{TestimonyArguments.SolaFide}{25}\)\gloss{\emph{would give:} Justification is by faith alone}
  \breakrow{(19)} \((\pvar{TestimonyArguments.SolaFide}{17} \land \pvar{TestimonyArguments.SolaFide}{18} \land \pvar{TestimonyArguments.SolaFide}{19}) \rightarrow \pvar{TestimonyArguments.SolaFide}{26}\)\gloss{\emph{would give:} Faith is sufficient: whoever believes is saved}
  \breakrow{(20)} \((\pvar{TestimonyArguments.SolaFide}{9} \land \pvar{TestimonyArguments.SolaFide}{10}) \rightarrow \pvar{TestimonyArguments.SolaFide}{26}\)\gloss{\emph{would give:} Faith is sufficient: whoever believes is saved}
  \breakrow{(21)} \((\pvar{TestimonyArguments.SolaFide}{36} \land \pvar{TestimonyArguments.SolaFide}{37} \land \pvar{TestimonyArguments.SolaFide}{40}) \rightarrow \pvar{TestimonyArguments.SolaFide}{25}\)\gloss{\emph{would give:} Justification is by faith alone}
  \breakrow{(22)} \((\pvar{TestimonyArguments.SolaFide}{1} \land \pvar{TestimonyArguments.SolaFide}{4} \land \pvar{TestimonyArguments.SolaFide}{5}) \rightarrow \pvar{TestimonyArguments.SolaFide}{27}\)\gloss{\emph{would give:} Salvation is by grace, a gift and not wages owed}
  \breakrow{(23)} \((\pvar{TestimonyArguments.SolaFide}{1} \land \pvar{TestimonyArguments.SolaFide}{4} \land \pvar{TestimonyArguments.SolaFide}{5} \land \pvar{TestimonyArguments.SolaFide}{22} \land \pvar{TestimonyArguments.SolaFide}{24}) \rightarrow \pvar{TestimonyArguments.SolaFide}{28}\)\gloss{\emph{would give:} Salvation is not by works: no work is its ground}
  \breakrow{(24)} \(\pvar{TestimonyArguments.SolaFide}{25} \rightarrow \pvar{TestimonyArguments.SolaFide}{26}\)\gloss{\emph{would give:} Faith is sufficient: whoever believes is saved}
  \item[(25)] \((\pvar{TestimonyArguments.SolaFide}{26} \land \pvar{TestimonyArguments.SolaFide}{1}) \rightarrow \pvar{TestimonyArguments.SolaFide}{29}\)\gloss{\emph{would give:} Salvation is received through faith}
  \item[\(\vdash\)] \(\pvar{TestimonyArguments.SolaFide}{27} \land \pvar{TestimonyArguments.SolaFide}{28} \land \pvar{TestimonyArguments.SolaFide}{29}\)\gloss{\emph{at issue:} \(\pvar{TestimonyArguments.SolaFide}{27}\) Salvation is by grace, a gift and not wages owed; \(\pvar{TestimonyArguments.SolaFide}{28}\) Salvation is not by works: no work is its ground; \(\pvar{TestimonyArguments.SolaFide}{29}\) Salvation is received through faith}
  \keepwithnext
\end{derivation}
\emph{Grounded in scripture alone:} John 3:16–18, 3:36, 5:24, 20:31 — eternal life through believing in the Son.
\par\medskip\noindent\textbf{\texttt{grace\allowbreak{}Without\allowbreak{}James\allowbreak{}Harmonisation}} --- Grace alone, minus the James harmonisation\par\nobreak\noindent
The same, asked only about the first part of the conclusion: is salvation by grace, with James 2:24 left unanswered?

\begin{derivation}
  \item[(1)] \(\pvar{TestimonyArguments.SolaFide}{8}\)\gloss{Galatians opposes requiring circumcision, besides faith, for justification}
  \item[(2)] \(\pvar{TestimonyArguments.SolaFide}{6}\)\gloss{Paul’s ‘works of the law’ means human works in general, not boundary markers}
  \item[(3)] \(\pvar{TestimonyArguments.SolaFide}{7}\)\gloss{In Galatians 2:16 ‘faith of Christ’ (πίστις Χριστοῦ) means our faith in Christ}
  \breakrow{(4)} \(\pvar{TestimonyArguments.SolaFide}{18}\)\gloss{‘Saved’ (σῴζω) at Luke 7:50 means saved from sin, not healed}
  \breakrow{(5)} \(\pvar{TestimonyArguments.SolaFide}{19}\)\gloss{Luke 7:47 — her love is the evidence of her forgiveness, not its ground}
  \breakrow{(6)} \(\pvar{TestimonyArguments.SolaFide}{10}\)\gloss{The yoke refused at Acts 15:10 is the whole law as a condition of salvation}
  \breakrow{(7)} \(\pvar{TestimonyArguments.SolaFide}{36}\)\gloss{John 6:28–29 — the work God requires is that you believe in him whom he sent}
  \breakrow{(8)} \(\pvar{TestimonyArguments.SolaFide}{37}\)\gloss{John 3:16–18, 3:36, 5:24, 20:31 — eternal life through believing in the Son}
  \breakrow{(9)} \(\pvar{TestimonyArguments.SolaFide}{40}\)\gloss{The believing of John 6:29 is trust, which brings nothing and receives Christ}
  \breakrow{(10)} \(\pvar{TestimonyArguments.SolaFide}{1}\)\gloss{Ephesians 2:8–9 teaches salvation by grace through faith, not of works}
  \breakrow{(11)} \(\pvar{TestimonyArguments.SolaFide}{2}\)\gloss{Romans 3:28 teaches justification by faith apart from works of the law}
  \breakrow{(12)} \(\pvar{TestimonyArguments.SolaFide}{3}\)\gloss{Galatians 2:16 teaches that no one is justified by works of the law}
  \breakrow{(13)} \(\pvar{TestimonyArguments.SolaFide}{4}\)\gloss{Romans 4:4–5 contrasts wages owed with a gift reckoned to the one who believes}
  \breakrow{(14)} \(\pvar{TestimonyArguments.SolaFide}{5}\)\gloss{Titus 3:5 teaches that God saved us not by works done in righteousness}
  \breakrow{(15)} \(\pvar{TestimonyArguments.SolaFide}{17}\)\gloss{Luke 7:50 — Jesus says ‘your faith has saved you’ after declaring sins forgiven}
  \breakrow{(16)} \(\pvar{TestimonyArguments.SolaFide}{9}\)\gloss{Acts 15:9–11 — Peter: hearts cleansed by faith; saved through grace}
  \breakrow{(17)} \(\pvar{TestimonyArguments.SolaFide}{24}\)\gloss{Scripture does not contradict itself}
  \breakrow{(18)} \((\pvar{TestimonyArguments.SolaFide}{2} \land \pvar{TestimonyArguments.SolaFide}{3} \land \pvar{TestimonyArguments.SolaFide}{8} \land \pvar{TestimonyArguments.SolaFide}{7} \land \pvar{TestimonyArguments.SolaFide}{6}) \rightarrow \pvar{TestimonyArguments.SolaFide}{25}\)\gloss{\emph{would give:} Justification is by faith alone}
  \breakrow{(19)} \((\pvar{TestimonyArguments.SolaFide}{17} \land \pvar{TestimonyArguments.SolaFide}{18} \land \pvar{TestimonyArguments.SolaFide}{19}) \rightarrow \pvar{TestimonyArguments.SolaFide}{26}\)\gloss{\emph{would give:} Faith is sufficient: whoever believes is saved}
  \breakrow{(20)} \((\pvar{TestimonyArguments.SolaFide}{9} \land \pvar{TestimonyArguments.SolaFide}{10}) \rightarrow \pvar{TestimonyArguments.SolaFide}{26}\)\gloss{\emph{would give:} Faith is sufficient: whoever believes is saved}
  \breakrow{(21)} \((\pvar{TestimonyArguments.SolaFide}{36} \land \pvar{TestimonyArguments.SolaFide}{37} \land \pvar{TestimonyArguments.SolaFide}{40}) \rightarrow \pvar{TestimonyArguments.SolaFide}{25}\)\gloss{\emph{would give:} Justification is by faith alone}
  \breakrow{(22)} \((\pvar{TestimonyArguments.SolaFide}{1} \land \pvar{TestimonyArguments.SolaFide}{4} \land \pvar{TestimonyArguments.SolaFide}{5}) \rightarrow \pvar{TestimonyArguments.SolaFide}{27}\)
  \breakrow{(23)} \((\pvar{TestimonyArguments.SolaFide}{1} \land \pvar{TestimonyArguments.SolaFide}{4} \land \pvar{TestimonyArguments.SolaFide}{5} \land \pvar{TestimonyArguments.SolaFide}{22} \land \pvar{TestimonyArguments.SolaFide}{24}) \rightarrow \pvar{TestimonyArguments.SolaFide}{28}\)\gloss{\emph{would give:} Salvation is not by works: no work is its ground}
  \breakrow{(24)} \(\pvar{TestimonyArguments.SolaFide}{25} \rightarrow \pvar{TestimonyArguments.SolaFide}{26}\)\gloss{\emph{would give:} Faith is sufficient: whoever believes is saved}
  \item[(25)] \((\pvar{TestimonyArguments.SolaFide}{26} \land \pvar{TestimonyArguments.SolaFide}{1}) \rightarrow \pvar{TestimonyArguments.SolaFide}{29}\)\gloss{\emph{would give:} Salvation is received through faith}
  \item[\(\vdash\)] \(\pvar{TestimonyArguments.SolaFide}{27}\)\gloss{\emph{at issue:} Salvation is by grace, a gift and not wages owed}
  \keepwithnext
\end{derivation}
\emph{Grounded in scripture alone:} John 3:16–18, 3:36, 5:24, 20:31 — eternal life through believing in the Son.
\par\medskip\noindent\textbf{\texttt{tridentine\allowbreak{}Case}} --- Trent, against ‘not by works’\par\nobreak\noindent
Trent, as the argument it makes when set against the others. Justification is the renewal of the inward person; that renewal grows through good works; so works done in grace merit an increase of justification, and salvation is not apart from works. The premises are those of \texttt{tridentine}. The conclusion is what Trent holds against sola fide.

\begin{derivation}
  \item[(1)] \(\pvar{TestimonyArguments.SolaFide}{1}\)\gloss{Ephesians 2:8–9 teaches salvation by grace through faith, not of works}
  \item[(2)] \(\pvar{TestimonyArguments.SolaFide}{2}\)\gloss{Romans 3:28 teaches justification by faith apart from works of the law}
  \item[(3)] \(\pvar{TestimonyArguments.SolaFide}{3}\)\gloss{Galatians 2:16 teaches that no one is justified by works of the law}
  \breakrow{(4)} \(\pvar{TestimonyArguments.SolaFide}{5}\)\gloss{Titus 3:5 teaches that God saved us not by works done in righteousness}
  \breakrow{(5)} \(\pvar{TestimonyArguments.SolaFide}{24}\)\gloss{Scripture does not contradict itself}
  \breakrow{(6)} \(\pvar{TestimonyArguments.SolaFide}{33}\)\gloss{Justification is not remission of sins only, but renewal of the inward man}
  \breakrow{(7)} \(\pvar{TestimonyArguments.SolaFide}{35}\)\gloss{The inward renewal of the justified grows as they do good works in grace}
  \breakrow{(8)} \((\pvar{TestimonyArguments.SolaFide}{33} \land \pvar{TestimonyArguments.SolaFide}{35}) \rightarrow \pvar{TestimonyArguments.SolaFide}{32}\)\gloss{\emph{would give:} Works performed in grace merit an increase of justification}
  \breakrow{(9)} \(\pvar{TestimonyArguments.SolaFide}{33} \rightarrow \lnot \pvar{TestimonyArguments.SolaFide}{30}\)\gloss{\emph{would deny:} Justification is a verdict only: pardon and righteousness counted, not poured in}
  \breakrow{(10)} \(\pvar{TestimonyArguments.SolaFide}{33} \rightarrow \lnot \pvar{TestimonyArguments.SolaFide}{34}\)\gloss{\emph{would deny:} Justification and sanctification are inseparable but distinct}
  \breakrow{(11)} \(\pvar{TestimonyArguments.SolaFide}{33} \rightarrow \lnot \pvar{TestimonyArguments.SolaFide}{26}\)\gloss{\emph{would deny:} Faith is sufficient: whoever believes is saved}
  \item[(12)] \(\pvar{TestimonyArguments.SolaFide}{32} \rightarrow \lnot \pvar{TestimonyArguments.SolaFide}{28}\)
  \item[\(\vdash\)] \(\lnot \pvar{TestimonyArguments.SolaFide}{28}\)\gloss{\emph{at issue, the denial of:} Salvation is not by works: no work is its ground}
  \keepwithnext
\end{derivation}
\emph{No premise here rests on scripture alone.}
\paragraph*{Paul’s gospel, against Trent’s definition}
\leavevmode\par\nobreak
\par\medskip\noindent\textbf{\texttt{galatian\allowbreak{}Gospel}} --- Paul’s gospel (Galatians 1), against justification as renewal\par\nobreak\noindent
\textbf{Paul’s gospel, against justification as renewal.} Galatians and 1 Corinthians read as one claim, that Christ’s work is the whole ground. With the courtroom sense of “justify” and Romans 8:33–34, it concludes that justification is not the renewal of the inward person. That is the direct denial of Trent’s definition (Session VI, ch. 7).

\begin{derivation}
  \item[(1)] \(\pvar{TestimonyArguments.SolaFide}{44}\)\gloss{Galatians 1:6–9: whoever preaches a gospel contrary to the one received is accursed}
  \item[(2)] \(\pvar{TestimonyArguments.SolaFide}{45}\)\gloss{1 Corinthians 15:3: the gospel of first importance — Christ died for our sins}
  \item[(3)] \(\pvar{TestimonyArguments.SolaFide}{46}\)\gloss{Galatians 2:21: if righteousness were through the law, Christ died for no purpose}
  \breakrow{(4)} \(\pvar{TestimonyArguments.SolaFide}{47}\)\gloss{Galatians 5:2–4: you who would be justified by the law are severed from Christ}
  \breakrow{(5)} \((\pvar{TestimonyArguments.SolaFide}{44} \land \pvar{TestimonyArguments.SolaFide}{45} \land \pvar{TestimonyArguments.SolaFide}{46} \land \pvar{TestimonyArguments.SolaFide}{47}) \rightarrow \pvar{TestimonyArguments.SolaFide}{58}\)\gloss{\emph{would give:} Christ’s death for our sins is the whole ground of justification; to add to it is another gospel}
  \breakrow{(6)} \(\pvar{TestimonyArguments.SolaFide}{48}\)\gloss{Romans 8:33–34: it is God who justifies — who is to condemn?}
  \breakrow{(7)} \(\pvar{TestimonyArguments.SolaFide}{49}\)\gloss{In Paul, ‘justify’ (δικαιόω) is a verdict: to declare righteous}
  \item[(8)] \((\pvar{TestimonyArguments.SolaFide}{58} \land \pvar{TestimonyArguments.SolaFide}{49} \land \pvar{TestimonyArguments.SolaFide}{48}) \rightarrow \lnot \pvar{TestimonyArguments.SolaFide}{33}\)
  \item[\(\vdash\)] \(\lnot \pvar{TestimonyArguments.SolaFide}{33}\)\gloss{\emph{at issue, the denial of:} Justification is not remission of sins only, but renewal of the inward man}
  \keepwithnext
\end{derivation}
\emph{Grounded in scripture alone:} 1 Corinthians 15:3: the gospel of first importance — Christ died for our sins, Romans 8:33–34: it is God who justifies — who is to condemn?.
\par\medskip\noindent\textbf{\texttt{romans\allowbreak{}Four\allowbreak{}On\allowbreak{}The\allowbreak{}Word}} --- Romans 4:3–8, on Paul’s word (by the rule of least meaning)\par\nobreak\noindent
\textbf{Romans 4, on what Paul’s word means.} Paul explains the righteousness God counts to the ungodly as sin not counted. By the rule of least meaning, the verb he explains in that way does not also mean the renewal of the inward person. This reaches from one passage what the argument from the word reaches from usage.

\begin{derivation}
  \item[(1)] \(\pvar{TestimonyArguments.SolaFide}{4}\)\gloss{Romans 4:4–5 contrasts wages owed with a gift reckoned to the one who believes}
  \item[(2)] \(\pvar{TestimonyArguments.SolaFide}{55}\)\gloss{Romans 4:6–8: righteousness counted apart from works — sins forgiven, sin not counted}
  \item[(3)] \(\pvar{TestimonyArguments.SolaFide}{56}\)\gloss{‘Count’ (λογίζομαι) in Romans 4:3–11 is the language of an account: credited to one}
  \breakrow{(4)} \(\pvar{TestimonyArguments.SolaFide}{51}\)\gloss{A word contributes the least meaning its context requires; one occurrence does not carry the whole concept}
  \breakrow{(5)} \((\pvar{TestimonyArguments.SolaFide}{4} \land \pvar{TestimonyArguments.SolaFide}{55} \land \pvar{TestimonyArguments.SolaFide}{56}) \rightarrow \pvar{TestimonyArguments.SolaFide}{57}\)\gloss{\emph{would give:} In Romans 4:6–8 Paul explains righteousness counted apart from works as sin not counted}
  \item[(6)] \((\pvar{TestimonyArguments.SolaFide}{57} \land \pvar{TestimonyArguments.SolaFide}{51}) \rightarrow \lnot \pvar{TestimonyArguments.SolaFide}{52}\)
  \item[\(\vdash\)] \(\lnot \pvar{TestimonyArguments.SolaFide}{52}\)\gloss{\emph{at issue, the denial of:} Paul’s word ‘justify’ means the renewal of the inward man, not only a verdict}
  \keepwithnext
\end{derivation}
\emph{No premise here rests on scripture alone.}
\par\medskip\noindent\textbf{\texttt{romans\allowbreak{}Four\allowbreak{}Case}} --- Romans 4:3–8, against justification as renewal\par\nobreak\noindent
\textbf{Romans 4, against justification as renewal.} Paul explains the righteousness God counts to the ungodly as sin not counted. With the courtroom sense of the verb, justification is not the renewal of the inward person.

\begin{derivation}
  \item[(1)] \(\pvar{TestimonyArguments.SolaFide}{4}\)\gloss{Romans 4:4–5 contrasts wages owed with a gift reckoned to the one who believes}
  \item[(2)] \(\pvar{TestimonyArguments.SolaFide}{55}\)\gloss{Romans 4:6–8: righteousness counted apart from works — sins forgiven, sin not counted}
  \item[(3)] \(\pvar{TestimonyArguments.SolaFide}{56}\)\gloss{‘Count’ (λογίζομαι) in Romans 4:3–11 is the language of an account: credited to one}
  \breakrow{(4)} \(\pvar{TestimonyArguments.SolaFide}{49}\)\gloss{In Paul, ‘justify’ (δικαιόω) is a verdict: to declare righteous}
  \breakrow{(5)} \((\pvar{TestimonyArguments.SolaFide}{4} \land \pvar{TestimonyArguments.SolaFide}{55} \land \pvar{TestimonyArguments.SolaFide}{56}) \rightarrow \pvar{TestimonyArguments.SolaFide}{57}\)\gloss{\emph{would give:} In Romans 4:6–8 Paul explains righteousness counted apart from works as sin not counted}
  \item[(6)] \((\pvar{TestimonyArguments.SolaFide}{57} \land \pvar{TestimonyArguments.SolaFide}{49}) \rightarrow \lnot \pvar{TestimonyArguments.SolaFide}{33}\)
  \item[\(\vdash\)] \(\lnot \pvar{TestimonyArguments.SolaFide}{33}\)\gloss{\emph{at issue, the denial of:} Justification is not remission of sins only, but renewal of the inward man}
  \keepwithnext
\end{derivation}
\emph{No premise here rests on scripture alone.}
\paragraph*{Paul’s word, against justification as renewal}
\leavevmode\par\nobreak
\par\medskip\noindent\textbf{\texttt{lexical\allowbreak{}Case}} --- Paul’s word (δικαιόω, by the rule of least meaning)\par\nobreak\noindent
\textbf{Paul’s word.} “Justify” (δικαιόω) is a courtroom word. Paul names renewal with words of its own. And a word contributes the least meaning its context requires. So Paul’s “justify” does not mean the renewal of the inward person.

\begin{derivation}
  \item[(1)] \(\pvar{TestimonyArguments.SolaFide}{49}\)\gloss{In Paul, ‘justify’ (δικαιόω) is a verdict: to declare righteous}
  \item[(2)] \(\pvar{TestimonyArguments.SolaFide}{50}\)\gloss{Paul names renewal with other words (ἀνακαίνωσις, ἁγιασμός) and sets them beside ‘justify’}
  \item[(3)] \(\pvar{TestimonyArguments.SolaFide}{51}\)\gloss{A word contributes the least meaning its context requires; one occurrence does not carry the whole concept}
  \item[(4)] \((\pvar{TestimonyArguments.SolaFide}{49} \land \pvar{TestimonyArguments.SolaFide}{50} \land \pvar{TestimonyArguments.SolaFide}{51}) \rightarrow \lnot \pvar{TestimonyArguments.SolaFide}{52}\)
  \item[\(\vdash\)] \(\lnot \pvar{TestimonyArguments.SolaFide}{52}\)\gloss{\emph{at issue, the denial of:} Paul’s word ‘justify’ means the renewal of the inward man, not only a verdict}
  \keepwithnext
\end{derivation}
\emph{Grounded in scripture alone:} Paul names renewal with other words (ἀνακαίνωσις, ἁγιασμός) and sets them beside ‘justify’.
\paragraph*{Results}
\leavevmode\par\nobreak
\par\medskip\noindent\textbf{\texttt{reformed\_\allowbreak{}establishes}}\par\nobreak\noindent
\textbf{Grant the Reformed premises, and sola fide follows.} The premises are the texts all four strands read, the disputed reading each strand depends on, and the answer to James 2:24.

The argument is valid. It does not show that the premises are true. Each disputed reading is denied by a scholar cited on this page.

{\small\begin{verbatim}
theorem reformed_establishes : Establishes reformed
-- axioms: propext, Quot.sound
\end{verbatim}}
\par\medskip\noindent\textbf{\texttt{new\allowbreak{}Perspective\_\allowbreak{}establishes}}\par\nobreak\noindent
\textbf{The New Perspective reaches the same conclusion.} Dunn and Wright reject the traditional reading of “works of the law” (ἔργα νόμου). They do not reject justification by faith. Leave Paul’s phrase to them, and sola fide still follows, from Jesus’ words at Luke 7:50.

Winning the argument over “works of the law” does not defeat sola fide. It does not show that Dunn or Wright would put it this way, or that their reading of Paul is right. Their critics dispute the reading, not this result.

{\small\begin{verbatim}
theorem newPerspective_establishes : Establishes newPerspective
-- axioms: propext, Quot.sound
\end{verbatim}}
\par\medskip\noindent\textbf{\texttt{tridentine\allowbreak{}Reading}}\par\nobreak\noindent
The Council of Trent’s reading, written out. Works done in grace merit an increase of justification, so salvation is not apart from works. Justification renews the inward person, so it is not only a verdict and is not distinct from sanctification. And it is not had by faith alone: faith without charity does not suffice. Every text holds.

{\small\begin{verbatim}
def tridentineReading : Valuation Claim :=
  fun a =>
    match a with
    | Claim.salvationNotByWorks => False
    | Claim.justificationIsForensicOnly => False
    | Claim.justificationDistinctFromSanctification => False
    | Claim.faithIsSufficient => False
    | Claim.justificationByFaithAlone => False
    | x => True
\end{verbatim}}
\par\medskip\noindent\textbf{\texttt{tridentine\_\allowbreak{}not\_\allowbreak{}establishes}}\par\nobreak\noindent
\textbf{Trent’s premises do not lead to sola fide.} They lead to the denial of its second part: on Trent’s account salvation is not apart from works.

A reader who holds what Trent holds is not committed to sola fide, and the reading above shows a consistent way to hold it. It does not show that Trent is wrong, or right.

{\small\begin{verbatim}
theorem tridentine_not_establishes : ¬Establishes tridentine
-- axioms: propext, Quot.sound
\end{verbatim}}
\par\medskip\noindent\textbf{\texttt{works\allowbreak{}Of\allowbreak{}Law\_\allowbreak{}not\_\allowbreak{}load\_\allowbreak{}bearing}}\par\nobreak\noindent
\textbf{The argument does not depend on the reading of “works of the law”.} Remove the premise that Paul’s phrase (ἔργα νόμου) means works in general, keep everything else, and sola fide still follows by the other strands.

{\small\begin{verbatim}
theorem worksOfLaw_not_load_bearing : Establishes reformedWithoutWorksOfLaw
-- axioms: propext, Quot.sound
\end{verbatim}}
\par\medskip\noindent\textbf{\texttt{sozo\_\allowbreak{}not\_\allowbreak{}load\_\allowbreak{}bearing}}\par\nobreak\noindent
\textbf{Nor on the reading of “saved” at Luke 7:50.} Remove the premise that the word there (σῴζω) means saved and not healed, keep everything else, and sola fide still follows by the other strands.

{\small\begin{verbatim}
theorem sozo_not_load_bearing : Establishes reformedWithoutSozo
-- axioms: propext, Quot.sound
\end{verbatim}}
\par\medskip\noindent\textbf{\texttt{apostolic\_\allowbreak{}strand\_\allowbreak{}survives\_\allowbreak{}paul\_\allowbreak{}and\_\allowbreak{}luke}}\par\nobreak\noindent
\textbf{Acts supplies “through faith” with the other strands’ disputed readings left out.} Leave out the reading of “works of the law” in Paul, of “saved” in Luke 7:50, and of believing in John 6:29; that is, assume none of them either way. Keep what every strand shares: the verses on grace and works, and the answer to James. Sola fide still follows, with Peter’s speech at Jerusalem (Acts 15:7–11) as the only route to faith.

Peter’s strand does not lean on any of the other three disputes. It does not show that Peter’s strand is beyond dispute, or that Acts 15 is the only text used. The strand rests on its own disputed reading, of the yoke at 15:10. The result’s name is older than the strand from John; John’s reading is left out as well.

{\small\begin{verbatim}
theorem apostolic_strand_survives_paul_and_luke : Establishes
    reformedRestingOnActs
-- axioms: propext, Quot.sound
\end{verbatim}}
\par\medskip\noindent\textbf{\texttt{johannine\_\allowbreak{}strand\_\allowbreak{}survives\_\allowbreak{}the\_\allowbreak{}other\_\allowbreak{}three}}\par\nobreak\noindent
\textbf{John supplies “through faith” with the other strands’ disputed readings left out.} Leave out the reading of “works of the law” in Paul, of “saved” in Luke 7:50, and of the yoke in Acts 15:10; that is, assume none of them either way. Keep what every strand shares: the verses on grace and works, and the answer to James. Sola fide still follows, with Jesus’ answer in John 6:29 as the only route to faith: “This is the work of God, that you believe in him whom he has sent.”

A reader who wins the disputes over Paul, Luke and Acts has not yet answered John. It does not show that John’s strand is beyond dispute, or that John 6:29 is the only text used. Grace and “not by works” still come from the shared verses and the answer to James. The strand rests on reading that believing as trust, which Aquinas denies (\texttt{johannine\_\allowbreak{}strand\_\allowbreak{}rests\_\allowbreak{}on\_\allowbreak{}believing\_\allowbreak{}as\_\allowbreak{}trust}), and on a step from “eternal life through believing” to “justified by faith alone” that is rated \texttt{plausible}.

{\small\begin{verbatim}
theorem johannine_strand_survives_the_other_three : Establishes
    reformedRestingOnJohn
-- axioms: propext, Quot.sound
\end{verbatim}}
\par\medskip\noindent\textbf{\texttt{pistis\allowbreak{}Christou\_\allowbreak{}not\_\allowbreak{}load\_\allowbreak{}bearing}}\par\nobreak\noindent
\textbf{Nor on “faith in Christ”.} Grant Hays that the phrase in Galatians 2:16 (πίστις Χριστοῦ) means Christ’s own faithfulness and not our faith in him. Keep everything else. Sola fide still follows.

A reader who wins this point for Hays has cost the Reformed case its route through Paul, and not its conclusion. It does not show that the point is unimportant to Paul’s strand. Within that strand everything depends on it; see \texttt{pistis\allowbreak{}Christou\_\allowbreak{}jointly\_\allowbreak{}load\_\allowbreak{}bearing}.

{\small\begin{verbatim}
theorem pistisChristou_not_load_bearing : Establishes subjectiveGenitive
-- axioms: propext, Quot.sound
\end{verbatim}}
\par\medskip\noindent\textbf{\texttt{acts15\allowbreak{}Yoke\_\allowbreak{}not\_\allowbreak{}load\_\allowbreak{}bearing}}\par\nobreak\noindent
\textbf{Nor on the yoke.} Grant Jervell that the yoke Peter refuses at Acts 15:10 is Israel’s law laid on gentiles, and not the law as a condition of salvation. Keep everything else. Sola fide still follows, by Paul, by Luke 7:50 and by John.

{\small\begin{verbatim}
theorem acts15Yoke_not_load_bearing : Establishes lawObservantLuke
-- axioms: propext, Quot.sound
\end{verbatim}}
\par\medskip\noindent\textbf{\texttt{neither\allowbreak{}Lexical\allowbreak{}Reading}}\par\nobreak\noindent
A reading on which no strand’s disputed premise holds. Paul’s “works of the law” are the boundary markers. Jesus’ “your faith has saved you” at Luke 7:50 is about healing. Peter’s yoke is Israel’s law laid on gentiles. And the believing of John 6:29 is not bare trust. Every text stands; faith is not said to be enough.

{\small\begin{verbatim}
def neitherLexicalReading : Valuation Claim :=
  fun a =>
    match a with
    | Claim.worksOfLawMeansWorksGenerally => False
    | Claim.sozoIsSoteriological => False
    | Claim.acts15YokeIsLawAsCondition => False
    | Claim.johannineBelievingIsTrust => False
    | Claim.justificationByFaithAlone => False
    | Claim.faithIsSufficient => False
    | Claim.salvationThroughFaith => False
    | x => True
\end{verbatim}}
\par\medskip\noindent\textbf{\texttt{lexical\_\allowbreak{}premises\_\allowbreak{}jointly\_\allowbreak{}load\_\allowbreak{}bearing}}\par\nobreak\noindent
\textbf{Sola fide fails to follow only when every strand loses its disputed reading.} Remove one disputed reading from each strand — “works of the law” as works in general from Paul, “saved” as salvation from Luke 7:50, the yoke as the law as a condition of salvation from Acts 15, believing as trust from John 6:29 — and keep everything else. Sola fide no longer follows. Remove fewer, and a strand that still has its reading reaches the conclusion.

So the argument, as encoded here, does not hang on the meaning of Paul’s “works of the law”. It hangs on that, \emph{or} on Luke’s “saved”, \emph{or} on Peter’s yoke, \emph{or} on John’s “believe”. Someone who denies sola fide has to answer all four. The same holds with “faith of Christ” (πίστις Χριστοῦ) in place of “works of the law”; see \texttt{pistis\allowbreak{}Christou\_\allowbreak{}jointly\_\allowbreak{}load\_\allowbreak{}bearing}.

It does not show that the four are independent in practice. A reader who takes “works of the law” as boundary markers will usually read Peter’s yoke the same way. The result shows that four separate premises would each have to be denied. It does not show that four separate groups of scholars deny them.

{\small\begin{verbatim}
theorem lexical_premises_jointly_load_bearing : ¬Establishes
    reformedWithoutEveryStrandsPremise
-- axioms: propext, Quot.sound
\end{verbatim}}
\par\medskip\noindent\textbf{\texttt{neither\allowbreak{}Pistis\allowbreak{}Nor\allowbreak{}Sozo\allowbreak{}Reading}}\par\nobreak\noindent
A reading on which Galatians 2:16 names Christ’s own faithfulness and not faith in Christ, Jesus’ “your faith has saved you” is about healing, Peter’s yoke is Israel’s law laid on gentiles, and the believing of John 6:29 is not bare trust.

{\small\begin{verbatim}
def neitherPistisNorSozoReading : Valuation Claim :=
  fun a =>
    match a with
    | Claim.pistisChristouObjective => False
    | Claim.sozoIsSoteriological => False
    | Claim.acts15YokeIsLawAsCondition => False
    | Claim.johannineBelievingIsTrust => False
    | Claim.justificationByFaithAlone => False
    | Claim.faithIsSufficient => False
    | Claim.salvationThroughFaith => False
    | x => True
\end{verbatim}}
\par\medskip\noindent\textbf{\texttt{pistis\allowbreak{}Christou\_\allowbreak{}jointly\_\allowbreak{}load\_\allowbreak{}bearing}}\par\nobreak\noindent
\textbf{Within Paul’s strand, “faith in Christ” is needed.} Deny that Galatians 2:16 speaks of our faith in Christ, and leave out the disputed reading of each of the other three strands. Keep the reading of “works of the law” as works in general, and keep Galatians as a letter against requiring circumcision. Sola fide no longer follows.

Excluding works is not yet \emph{faith alone}. The verse must also name the believer’s faith as the means, and on Hays’s reading it does not. So Paul’s strand rests on two disputed readings of one verse, and losing either costs it the conclusion.

{\small\begin{verbatim}
theorem pistisChristou_jointly_load_bearing : ¬Establishes
    reformedWithoutPistisChristouSozoOrYoke
-- axioms: propext, Quot.sound
\end{verbatim}}
\par\medskip\noindent\textbf{\texttt{authorship\_\allowbreak{}not\_\allowbreak{}load\_\allowbreak{}bearing}}\par\nobreak\noindent
\textbf{The argument does not depend on who wrote the disputed letters.} Grant every critical conclusion: that Ephesians and Titus are not by Paul, and 1 and 2 Peter not by Peter. Sola fide still follows.

The argument reads these letters as Scripture, not as the testimony of a named apostle, so their authorship is not one of its premises. The proof shows this directly: it adds the four denials to the Reformed premises and uses none of them.

What this costs is stated at \texttt{critical\allowbreak{}Authorship}. The letters cannot then be used as evidence of what \emph{Paul} meant by “works”, which is why that premise rests on Romans.

{\small\begin{verbatim}
theorem authorship_not_load_bearing : Establishes criticalAuthorship
-- axioms: propext, Quot.sound
\end{verbatim}}
\par\medskip\noindent\textbf{\texttt{apocalyptic\allowbreak{}Reading}}\par\nobreak\noindent
The apocalyptic reading of Paul (Martyn, Campbell), written out. “Faith of Christ” is Christ’s own faithfulness. “The righteousness of God” is God’s act of rescue. Faith is not the condition on which a person is justified. Everything else it grants.

{\small\begin{verbatim}
def apocalypticReading : Valuation Claim :=
  fun a =>
    match a with
    | Claim.pistisChristouObjective => False
    | Claim.justificationByFaithAlone => False
    | Claim.faithIsSufficient => False
    | Claim.salvationThroughFaith => False
    | x => True
\end{verbatim}}
\par\medskip\noindent\textbf{\texttt{apocalyptic\_\allowbreak{}not\_\allowbreak{}establishes}}\par\nobreak\noindent
\textbf{The apocalyptic reading does not lead to sola fide.} It does not reach the third part of the conclusion, “through faith”. Nothing it holds says that faith is enough, and it denies the faith alone from which Paul’s strand would derive that.

What it does reach is stated as a result of its own, \texttt{apocalyptic\_\allowbreak{}establishes\_\allowbreak{}grace\_\allowbreak{}not\_\allowbreak{}works}.

{\small\begin{verbatim}
theorem apocalyptic_not_establishes : ¬Establishes apocalyptic
-- axioms: propext, Quot.sound
\end{verbatim}}
\par\medskip\noindent\textbf{\texttt{apocalyptic\_\allowbreak{}establishes\_\allowbreak{}grace\_\allowbreak{}not\_\allowbreak{}works}}\par\nobreak\noindent
\textbf{The apocalyptic reading keeps grace, and “not by works”.} If “the righteousness of God” is God’s act of rescue, then salvation is a gift that depends on nothing a person does. So the reading reaches the first two parts of the conclusion by its own route.

This places its disagreement with the Reformed reading exactly. They differ over \emph{through faith}. They do not differ over grace or over works.

{\small\begin{verbatim}
theorem apocalyptic_establishes_grace_not_works : Establishes
    apocalypticOnGraceAndWorks
-- axioms: propext, Quot.sound
\end{verbatim}}
\par\medskip\noindent\textbf{\texttt{grace\_\allowbreak{}and\_\allowbreak{}works\_\allowbreak{}rest\_\allowbreak{}on\_\allowbreak{}no\_\allowbreak{}lexical\_\allowbreak{}premise}}\par\nobreak\noindent
\textbf{Grace, and “not by works”, depend on none of the disputed readings.} Remove all of them: both of Paul’s, Luke’s “saved”, the yoke in Acts and believing as trust in John. The first two parts of the conclusion still follow, from the verses that say “not of works” in as many words and from the answer to James.

So the disputed readings decide one thing: whether the texts that exclude works also name faith as the means. They do not decide whether works are excluded.

{\small\begin{verbatim}
theorem grace_and_works_rest_on_no_lexical_premise : Establishes
    graceAndWorksWithoutAnyLexicalPremise
-- axioms: propext, Quot.sound
\end{verbatim}}
\par\medskip\noindent\textbf{\texttt{forensic\_\allowbreak{}justification\_\allowbreak{}not\_\allowbreak{}load\_\allowbreak{}bearing}}\par\nobreak\noindent
\textbf{The argument does not depend on justification being only a verdict.} Grant the Finnish reading of Luther (Mannermaa): justification is not a declaration and nothing more, because Christ himself is present in faith. Sola fide still follows. The proof adds that one denial to the Reformed premises and uses nothing else.

On what justification \emph{is}, Mannermaa stands with Trent against the Reformed account. On sola fide he stands with the Reformed against Trent. So denying the verdict-only account does not by itself deny sola fide. What Trent’s objection needs is narrower: that justification includes the sanctification that grows through works (\texttt{trent\_\allowbreak{}objection\_\allowbreak{}rests\_\allowbreak{}on\_\allowbreak{}its\_\allowbreak{}definition}).

{\small\begin{verbatim}
theorem forensic_justification_not_load_bearing : Establishes finnish
-- axioms: propext, Quot.sound
\end{verbatim}}
\par\medskip\noindent\textbf{\texttt{critics\_\allowbreak{}carry\_\allowbreak{}the\_\allowbreak{}pauline\_\allowbreak{}strand}}\par\nobreak\noindent
\textbf{On the critics’ account of Judaism, Paul alone reaches sola fide.} Gathercole and others deny that Judaism in Paul’s day was simply “in by grace, staying in by obedience”. Grant them that, keep Galatians as a letter against requiring circumcision, and leave Luke, Acts and John out. Sola fide follows from Paul, and the reading of “works of the law” as works in general is derived and not assumed.

{\small\begin{verbatim}
theorem critics_carry_the_pauline_strand : Establishes
    paulineStrandOnTheCritics
-- axioms: propext, Quot.sound
\end{verbatim}}
\par\medskip\noindent\textbf{\texttt{sanders\allowbreak{}Reading}}\par\nobreak\noindent
Sanders’ account, with only Paul to go on. Judaism taught entry by grace and staying in by obedience, so “works of the law” are not works in general, and Paul’s strand does not reach faith alone.

{\small\begin{verbatim}
def sandersReading : Valuation Claim :=
  fun a =>
    match a with
    | Claim.worksOfLawMeansWorksGenerally => False
    | Claim.sozoIsSoteriological => False
    | Claim.acts15YokeIsLawAsCondition => False
    | Claim.justificationByFaithAlone => False
    | Claim.faithIsSufficient => False
    | Claim.salvationThroughFaith => False
    | x => True
\end{verbatim}}
\par\medskip\noindent\textbf{\texttt{sanders\_\allowbreak{}costs\_\allowbreak{}the\_\allowbreak{}pauline\_\allowbreak{}strand}}\par\nobreak\noindent
\textbf{On Sanders’ account, Paul alone does not.} Grant Sanders his account of Judaism and Dunn his inference from it, and “works of the law” does not mean works in general. Leave Luke, Acts and John out, and sola fide does not follow from Paul.

Taken with \texttt{critics\_\allowbreak{}carry\_\allowbreak{}the\_\allowbreak{}pauline\_\allowbreak{}strand}, this shows where the dispute over “works of the law” lies: in a question about first-century Judaism, which has to be answered before the question about Paul’s Greek.

{\small\begin{verbatim}
theorem sanders_costs_the_pauline_strand : ¬Establishes paulineStrandOnSanders
-- axioms: propext, Quot.sound
\end{verbatim}}
\par\medskip\noindent\textbf{\texttt{james\allowbreak{}Unharmonised\allowbreak{}Reading}}\par\nobreak\noindent
A reading on which James 2:24, “justified by works and not by faith alone”, stands against Paul with no answer given.

{\small\begin{verbatim}
def jamesUnharmonisedReading : Valuation Claim :=
  fun a =>
    match a with
    | Claim.james2_24Compatible => False
    | Claim.salvationNotByWorks => False
    | x => True
\end{verbatim}}
\par\medskip\noindent\textbf{\texttt{james\_\allowbreak{}harmonisation\_\allowbreak{}is\_\allowbreak{}load\_\allowbreak{}bearing}}\par\nobreak\noindent
\textbf{The argument depends on answering James.} Remove the two premises that reconcile James 2:24 with Paul, keep everything else, and sola fide no longer follows.

So a case for sola fide that does not deal with James 2:24 is incomplete: its conclusion does not follow from what it has said.

{\small\begin{verbatim}
theorem james_harmonisation_is_load_bearing : ¬Establishes
    reformedWithoutJamesHarmonisation
-- axioms: propext, Quot.sound
\end{verbatim}}
\par\medskip\noindent\textbf{\texttt{grace\_\allowbreak{}needs\_\allowbreak{}no\_\allowbreak{}answer\_\allowbreak{}to\_\allowbreak{}james}}\par\nobreak\noindent
\textbf{Grace does not depend on answering James.} Remove the answer to James, and “salvation is by grace” still follows, from Ephesians 2:8, Romans 4:4–5 and Titus 3:5.

James 2:24 bears on \emph{not by works}, the part of the conclusion it appears to contradict, and on nothing else.

{\small\begin{verbatim}
theorem grace_needs_no_answer_to_james : Establishes
    graceWithoutJamesHarmonisation
-- axioms: propext, Quot.sound
\end{verbatim}}
\paragraph*{Can the premises be true together?}
\leavevmode\par\nobreak
Anything follows from premises that contradict each other. So a position that “establishes” its conclusion only because its premises could never all be true would show nothing, and every other check would still pass. For each position with a result saying its conclusion follows, a reading is named below on which all its premises hold.

A position with a result saying its conclusion does \emph{not} follow needs no such check. The reading that shows the conclusion failing already makes every premise true.

These are checks on the encoding, not claims about justification, so they are not listed among the main results.

\par\medskip\noindent\textbf{\texttt{everything\allowbreak{}Holds\allowbreak{}Reading}}\par\nobreak\noindent
The reading on which every claim in this argument holds at once. It shows that the positions built only from affirmations do not contradict themselves.

{\small\begin{verbatim}
def everythingHoldsReading : Valuation Claim :=
  fun x => True
\end{verbatim}}
\par\medskip\noindent\textbf{\texttt{reformed\_\allowbreak{}is\_\allowbreak{}satisfiable}}\par\nobreak\noindent
The Reformed premises can all be true together. Without this, “sola fide follows from them” would be empty: anything follows from premises that contradict each other.

{\small\begin{verbatim}
theorem reformed_is_satisfiable : Satisfiable reformed.premises
-- axioms: propext, Quot.sound
\end{verbatim}}
\par\medskip\noindent\textbf{\texttt{new\allowbreak{}Perspective\allowbreak{}Own\allowbreak{}Reading}}\par\nobreak\noindent
The New Perspective’s own reading. Paul’s “works of the law” are the boundary markers and not works in general, and justification is by faith all the same. It also holds that faith is the only means and that justification is a verdict, so it denies Trent’s account of baptism as the instrument and of justification as including sanctification.

{\small\begin{verbatim}
def newPerspectiveOwnReading : Valuation Claim :=
  fun a =>
    match a with
    | Claim.paulsJustifyDenotesRenewal => False
    | Claim.rev22_11BeJustifiedStill => False
    | Claim.sir18_22BeJustifiedToDeath => False
    | Claim.jamesUsesDikaioAsPaul => False
    | Claim.worksOfLawMeansWorksGenerally => False
    | Claim.baptismIsInstrumentalCause => False
    | Claim.justificationIncludesSanctification => False
    | x => True
\end{verbatim}}
\par\medskip\noindent\textbf{\texttt{new\allowbreak{}Perspective\_\allowbreak{}is\_\allowbreak{}satisfiable}}\par\nobreak\noindent
The New Perspective’s premises can all be true together.

{\small\begin{verbatim}
theorem newPerspective_is_satisfiable : Satisfiable newPerspective.premises
-- axioms: propext, Quot.sound
\end{verbatim}}
\par\medskip\noindent\textbf{\texttt{subjective\allowbreak{}Genitive\allowbreak{}Reading}}\par\nobreak\noindent
Hays’s reading: “faith of Christ” is Christ’s own faithfulness, and everything else in the Reformed case holds.

{\small\begin{verbatim}
def subjectiveGenitiveReading : Valuation Claim :=
  fun a =>
    match a with
    | Claim.pistisChristouObjective => False
    | x => True
\end{verbatim}}
\par\medskip\noindent\textbf{\texttt{subjective\allowbreak{}Genitive\_\allowbreak{}is\_\allowbreak{}satisfiable}}\par\nobreak\noindent
That position can be held without contradiction, so \texttt{pistis\allowbreak{}Christou\_\allowbreak{}not\_\allowbreak{}load\_\allowbreak{}bearing} is not empty.

{\small\begin{verbatim}
theorem subjectiveGenitive_is_satisfiable : Satisfiable
    subjectiveGenitive.premises
-- axioms: propext, Quot.sound
\end{verbatim}}
\par\medskip\noindent\textbf{\texttt{law\allowbreak{}Observant\allowbreak{}Luke\allowbreak{}Reading}}\par\nobreak\noindent
Jervell’s reading: the yoke is Israel’s law for gentiles, and everything else in the Reformed case holds.

{\small\begin{verbatim}
def lawObservantLukeReading : Valuation Claim :=
  fun a =>
    match a with
    | Claim.acts15YokeIsLawAsCondition => False
    | x => True
\end{verbatim}}
\par\medskip\noindent\textbf{\texttt{law\allowbreak{}Observant\allowbreak{}Luke\_\allowbreak{}is\_\allowbreak{}satisfiable}}\par\nobreak\noindent
That position can be held without contradiction, so \texttt{acts15\allowbreak{}Yoke\_\allowbreak{}not\_\allowbreak{}load\_\allowbreak{}bearing} is not empty.

{\small\begin{verbatim}
theorem lawObservantLuke_is_satisfiable : Satisfiable
    lawObservantLuke.premises
-- axioms: propext, Quot.sound
\end{verbatim}}
\par\medskip\noindent\textbf{\texttt{critical\allowbreak{}Authorship\allowbreak{}Reading}}\par\nobreak\noindent
The critical view of authorship: none of the four disputed letters is by the apostle whose name it bears, and everything else holds.

{\small\begin{verbatim}
def criticalAuthorshipReading : Valuation Claim :=
  fun a =>
    match a with
    | Claim.ephesiansIsPauline => False
    | Claim.titusIsPauline => False
    | Claim.firstPeterIsPetrine => False
    | Claim.secondPeterIsPetrine => False
    | x => True
\end{verbatim}}
\par\medskip\noindent\textbf{\texttt{critical\allowbreak{}Authorship\_\allowbreak{}is\_\allowbreak{}satisfiable}}\par\nobreak\noindent
That position can be held without contradiction, so \texttt{authorship\_\allowbreak{}not\_\allowbreak{}load\_\allowbreak{}bearing} is not empty.

{\small\begin{verbatim}
theorem criticalAuthorship_is_satisfiable : Satisfiable
    criticalAuthorship.premises
-- axioms: propext, Quot.sound
\end{verbatim}}
\par\medskip\noindent\textbf{\texttt{reformed\allowbreak{}Without\allowbreak{}Works\allowbreak{}Of\allowbreak{}Law\_\allowbreak{}is\_\allowbreak{}satisfiable}}\par\nobreak\noindent
The Reformed premises without the reading of “works of the law” can all be true together, so \texttt{works\allowbreak{}Of\allowbreak{}Law\_\allowbreak{}not\_\allowbreak{}load\_\allowbreak{}bearing} is not empty.

{\small\begin{verbatim}
theorem reformedWithoutWorksOfLaw_is_satisfiable : Satisfiable
    reformedWithoutWorksOfLaw.premises
-- axioms: propext, Quot.sound
\end{verbatim}}
\par\medskip\noindent\textbf{\texttt{reformed\allowbreak{}Without\allowbreak{}Sozo\_\allowbreak{}is\_\allowbreak{}satisfiable}}\par\nobreak\noindent
And the same without the reading of “saved” at Luke 7:50.

{\small\begin{verbatim}
theorem reformedWithoutSozo_is_satisfiable : Satisfiable
    reformedWithoutSozo.premises
-- axioms: propext, Quot.sound
\end{verbatim}}
\par\medskip\noindent\textbf{\texttt{reformed\allowbreak{}Resting\allowbreak{}On\allowbreak{}Acts\_\allowbreak{}is\_\allowbreak{}satisfiable}}\par\nobreak\noindent
The package resting on Acts can be held without contradiction, so \texttt{apostolic\_\allowbreak{}strand\_\allowbreak{}survives\_\allowbreak{}paul\_\allowbreak{}and\_\allowbreak{}luke} is not true merely because its premises could never all hold.

{\small\begin{verbatim}
theorem reformedRestingOnActs_is_satisfiable : Satisfiable
    reformedRestingOnActs.premises
-- axioms: propext, Quot.sound
\end{verbatim}}
\par\medskip\noindent\textbf{\texttt{reformed\allowbreak{}Resting\allowbreak{}On\allowbreak{}John\_\allowbreak{}is\_\allowbreak{}satisfiable}}\par\nobreak\noindent
And so can the package resting on John, so \texttt{johannine\_\allowbreak{}strand\_\allowbreak{}survives\_\allowbreak{}the\_\allowbreak{}other\_\allowbreak{}three} is not empty either.

{\small\begin{verbatim}
theorem reformedRestingOnJohn_is_satisfiable : Satisfiable
    reformedRestingOnJohn.premises
-- axioms: propext, Quot.sound
\end{verbatim}}
\par\medskip\noindent\textbf{\texttt{apocalyptic\allowbreak{}On\allowbreak{}Grace\allowbreak{}And\allowbreak{}Works\_\allowbreak{}is\_\allowbreak{}satisfiable}}\par\nobreak\noindent
The apocalyptic reading, asked about grace and works, can be held without contradiction.

{\small\begin{verbatim}
theorem apocalypticOnGraceAndWorks_is_satisfiable : Satisfiable
    apocalypticOnGraceAndWorks.premises
-- axioms: propext, Quot.sound
\end{verbatim}}
\par\medskip\noindent\textbf{\texttt{grace\allowbreak{}And\allowbreak{}Works\allowbreak{}Without\allowbreak{}Any\allowbreak{}Lexical\allowbreak{}Premise\_\allowbreak{}is\_\allowbreak{}satisfiable}}\par\nobreak\noindent
With every disputed reading removed, the premises can still all be true together.

{\small\begin{verbatim}
theorem graceAndWorksWithoutAnyLexicalPremise_is_satisfiable : Satisfiable
    graceAndWorksWithoutAnyLexicalPremise.premises
-- axioms: propext, Quot.sound
\end{verbatim}}
\par\medskip\noindent\textbf{\texttt{finnish\allowbreak{}Reading}}\par\nobreak\noindent
Mannermaa’s reading: justification is not only a verdict, and everything else holds.

{\small\begin{verbatim}
def finnishReading : Valuation Claim :=
  fun a =>
    match a with
    | Claim.justificationIsForensicOnly => False
    | x => True
\end{verbatim}}
\par\medskip\noindent\textbf{\texttt{finnish\_\allowbreak{}is\_\allowbreak{}satisfiable}}\par\nobreak\noindent
That position can be held without contradiction, so \texttt{forensic\_\allowbreak{}justification\_\allowbreak{}not\_\allowbreak{}load\_\allowbreak{}bearing} is not empty.

{\small\begin{verbatim}
theorem finnish_is_satisfiable : Satisfiable finnish.premises
-- axioms: propext, Quot.sound
\end{verbatim}}
\par\medskip\noindent\textbf{\texttt{critics\allowbreak{}Reading}}\par\nobreak\noindent
The critics’ reading (Gathercole): Judaism in Paul’s day was not simply “in by grace, staying in by obedience”, and everything else holds. Faith is the only means and justification is a verdict, so this reading denies Trent’s account of baptism as the instrument and of justification as including sanctification.

{\small\begin{verbatim}
def criticsReading : Valuation Claim :=
  fun a =>
    match a with
    | Claim.paulsJustifyDenotesRenewal => False
    | Claim.rev22_11BeJustifiedStill => False
    | Claim.sir18_22BeJustifiedToDeath => False
    | Claim.jamesUsesDikaioAsPaul => False
    | Claim.secondTempleCovenantalNomism => False
    | Claim.baptismIsInstrumentalCause => False
    | Claim.justificationIncludesSanctification => False
    | x => True
\end{verbatim}}
\par\medskip\noindent\textbf{\texttt{pauline\allowbreak{}Strand\allowbreak{}On\allowbreak{}The\allowbreak{}Critics\_\allowbreak{}is\_\allowbreak{}satisfiable}}\par\nobreak\noindent
That position can be held without contradiction, so \texttt{critics\_\allowbreak{}carry\_\allowbreak{}the\_\allowbreak{}pauline\_\allowbreak{}strand} is not empty.

{\small\begin{verbatim}
theorem paulineStrandOnTheCritics_is_satisfiable : Satisfiable
    paulineStrandOnTheCritics.premises
-- axioms: propext, Quot.sound
\end{verbatim}}
\par\medskip\noindent\textbf{\texttt{grace\allowbreak{}Without\allowbreak{}James\allowbreak{}Harmonisation\_\allowbreak{}is\_\allowbreak{}satisfiable}}\par\nobreak\noindent
With the answer to James removed, the premises asked about grace can still all be true together, so \texttt{grace\_\allowbreak{}needs\_\allowbreak{}no\_\allowbreak{}answer\_\allowbreak{}to\_\allowbreak{}james} is not empty.

{\small\begin{verbatim}
theorem graceWithoutJamesHarmonisation_is_satisfiable : Satisfiable
    graceWithoutJamesHarmonisation.premises
-- axioms: propext, Quot.sound
\end{verbatim}}
\subsection*{John: “this is the work of God, that you believe”}
\addcontentsline{toc}{subsection}{John: “this is the work of God, that you believe”}
The day after the feeding of the five thousand, the crowd finds Jesus across the lake and asks him a question about works:

\begin{quote}
“What must we do, to be doing the works of God?” Jesus answered them, “This is the work of God, that you believe in him whom he has sent” (John 6:28–29).

\end{quote}
The question is about works, and the answer is believing. This is the fourth strand of sola fide, and John says the same thing elsewhere. Whoever believes in the Son has eternal life and does not come into judgement (3:16–18, 3:36, 5:24), and the Gospel gives this as its own purpose: “that by believing you may have life in his name” (20:31).

\subsubsection*{What both sides grant, and where they part}
Nobody in this dispute denies the verse. They differ on what its believing is.

\begin{itemize}
  \item \textbf{Calvin} reads it as trust. Christ calls faith a work only loosely: faith “brings nothing to God, but, on the contrary, places man before God as empty and poor, that he may be filled with Christ”. It is “a passive work, to which no reward can be paid”.
  \item \textbf{Aquinas} reads it as faith formed by charity, that is, faith made alive by love. He takes Paul to distinguish faith “only from external works”, so that believing is itself an inward work. To believe \emph{in} him, as the goal one moves toward, “is proper to faith living through the love of charity”, and such faith “is the principle of all our good works” (\emph{Super Ioannem} 6, lect. 3, n. 901). The Council of Trent describes justifying faith the same way: without hope and charity, faith “neither unites man perfectly with Christ, nor makes him a living member of His body” (Session VI, ch. 7).
\end{itemize}
So this strand has the shape of the other three: texts everyone accepts, and one reading on which the conclusion turns. On Calvin’s reading John supplies faith alone without the readings the other strands depend on (\texttt{johannine\_\allowbreak{}strand\_\allowbreak{}establishes}). On Aquinas’s it does not (\texttt{johannine\_\allowbreak{}strand\_\allowbreak{}rests\_\allowbreak{}on\_\allowbreak{}believing\_\allowbreak{}as\_\allowbreak{}trust}).

\subsubsection*{John has no word for “justify”}
Paul’s verb for “justify” (δικαιόω) does not occur in John’s Gospel. So the strand has to step from “eternal life through believing” to “justified by faith alone”, and that step is an inference, not the meaning of a word. It is cited to Calvin, whose comment on 6:29 speaks of the righteousness faith receives from Christ. It is rated \texttt{plausible}. No cited source accepts the strand’s premises and rejects the step, but no cited source argues the step at length either.

\subsubsection*{Luther’s answer to Aquinas}
Luther, in his own lectures on Galatians, answers the reading Aquinas gives of John 6:29. He names “the scholastics”, not Aquinas; the faith formed by charity that he attacks is the reading Aquinas gives. On 3:11: “The scholastics misconstrue this passage by saying: ‘The just shall live by faith, if it is a working faith, or a faith formed and performed by charitable works.’ … To speak of formed or unformed faith, a sort of double faith, is contrary to the Scriptures.” On 3:12: “Does not the Law command charity? … If the law requires charity, charity is part of the Law and not of faith.”

His argument runs: love is what the law commands, and “the law is not of faith” (Galatians 3:12), so a faith that justifies because love forms it would justify by the law after all.

The argument is encoded as a line of its own (\texttt{luther\allowbreak{}Line}). Its two starting points are ones Aquinas accepts: the text of Galatians 3:11–12, and that love of God and neighbour is what the law commands (Deuteronomy 6:5). The step from them is Luther’s. Grant the step, and Aquinas’s reading of John 6:29 falls (\texttt{luther\_\allowbreak{}answers\_\allowbreak{}aquinas\_\allowbreak{}from\_\allowbreak{}galatians}). Without the step the texts do not exclude his reading (\texttt{luther\_\allowbreak{}answer\_\allowbreak{}rests\_\allowbreak{}on\_\allowbreak{}his\_\allowbreak{}step}). The step is rated \texttt{disputed}, because Aquinas accepts both starting points and rejects the conclusion. His reason for rejecting it is not yet cited here.

Two limits. Removing Aquinas’s reading does not establish Calvin’s: “not formed by charity” is a weaker claim than “trust”. And the Catholic reply rests on Galatians 5:6, “faith working through love”, which Trent cites (Session VI, ch. 7). How Luther and Calvin read that verse is not yet encoded.

Calvin makes a neighbouring argument in the \emph{Institutes} (III.xi.19), against those who “pretend that the works excluded are ceremonial, not moral works”. He cites Galatians 3:10 and 3:12 and concludes that “the whole Law is spoken of when the power of justifying is denied to it”. That answers those who say Paul excludes only ceremonial works. Aquinas does not say that. He says believing is an inward work, formed by love. So Luther’s argument, and not Calvin’s, is the one that addresses him.

\subsubsection*{Where the strand stands in the whole dispute}
The strand is one of the four strands of the Reformed case, and John’s case, Aquinas’s and Luther’s are parties to the sola fide dispute. Three things are checked.

\begin{itemize}
  \item \textbf{John supplies “through faith” without the other three strands.} Leave out the disputed readings in Paul, Luke and Acts. Keep what every strand shares: the verses on grace and works, and the answer to James. Sola fide still follows, with John 6:28–29 as the only route to faith (\texttt{johannine\_\allowbreak{}strand\_\allowbreak{}survives\_\allowbreak{}the\_\allowbreak{}other\_\allowbreak{}three}).
  \item \textbf{John’s disputed reading is not one the whole case depends on.} Grant Aquinas his reading of 6:29, keep the other three strands, and sola fide still follows (\texttt{believing\_\allowbreak{}as\_\allowbreak{}trust\_\allowbreak{}not\_\allowbreak{}load\_\allowbreak{}bearing}).
  \item \textbf{John’s case is not forced, and neither is any rival.} John’s case and Aquinas’s each defeat the other, that is, each contradicts the other and neither is better supported. So do Aquinas’s and Luther’s, John’s and Trent’s, and John’s and the apocalyptic reading’s. Each pair is a tie, because each side rests on a premise rated \texttt{disputed}.
\end{itemize}
\subsubsection*{John 3:36: “whoever does not obey the Son”}
Three chapters on, John sets two words side by side that a reader of this debate cannot pass over:

\begin{quote}
“Whoever believes in the Son has eternal life; whoever does not obey the Son shall not see life, but the wrath of God remains on him” (John 3:36).

\end{quote}
The verse sets \emph{not obeying} against \emph{believing}. Does believing, in John, include obeying? And if it does, can faith alone be what saves? The record points two ways, one on the word and one on the doctrine.

The Greek for “does not obey” (ὁ ἀπειθῶν) is read by the commentators cited here as unbelief seen in what it does. Westcott: “believeth not] disobeyeth. Disbelief is regarded in its activity.” Calvin’s comment takes the clause of “all who do not believe in Christ”. Among modern commentators, Ridderbos: “‘Disobeys’ serves as a description of ‘unbelief’”. Carson takes the verb as “disobeys”, and adds that “failure to trust him is as much disobedience as unbelief”. No source has been found that argues from this word that believing includes obeying.

The doctrine is another matter. Chrysostom, preaching on this verse, does deny that believing is enough: “‘Is it then enough,’ saith one, ‘to believe on the Son, that one may have eternal life?’ By no means.” He does not argue from the word. He argues from Jesus’ saying in Matthew 7:21, “Not everyone who says to me, ‘Lord, Lord,’ will enter the kingdom of heaven”, and concludes: “though a man believe rightly on the Father, the Son, and the Holy Ghost, yet if he lead not a right life, his faith will avail nothing towards his salvation” (\emph{Homilies on John} 31).

Chrysostom’s argument is written out as starting points, a step and a conclusion (\texttt{chrysostom\allowbreak{}Line}). Its starting points are two texts every reader here accepts. The step from them to “faith does not suffice” is his. Three things are checked.

\begin{itemize}
  \item Granted his step, his conclusion follows (\texttt{chrysostom\_\allowbreak{}reads\_\allowbreak{}john3\_\allowbreak{}36\_\allowbreak{}against\_\allowbreak{}faith\_\allowbreak{}alone}).
  \item The wording of the two texts does not settle it. One consistent reading holds both texts and that faith suffices; another holds both and that it does not (\texttt{john3\_\allowbreak{}36\_\allowbreak{}leaves\_\allowbreak{}faiths\_\allowbreak{}sufficiency\_\allowbreak{}open}).
  \item His conclusion, written out as “faith does not suffice”, and John’s strand cannot both be held (\texttt{chrysostom\_\allowbreak{}cannot\_\allowbreak{}be\_\allowbreak{}held\_\allowbreak{}with\_\allowbreak{}johns\_\allowbreak{}case}). Whether the faith he means is the faith the Reformers mean is a further question, noted at that result.
\end{itemize}
The step is rated \texttt{disputed}. Calvin accepts both texts and denies the conclusion: he reads the doing of Matthew 7:21 as including faith, and good works as what follows faith. Among modern readers Chrysostom has company: Stanley holds that “Endurance is a condition for final salvation”, and reads Matthew 7:21 so.

Chrysostom’s case is heard in the main sola fide debate, and so is the Reformed answer to it (below). His conclusion is general, so there he meets every case that concludes that faith suffices: Paul’s, Luke’s, Peter’s and John’s, and the Reformed answer. Each of those meets him in return, and every one of the ties is between positions rated \texttt{disputed}. He does not meet the arguments about Paul’s word or Paul’s gospel, and he contradicts no verse of Paul’s (\texttt{chrysostom\_\allowbreak{}grants\_\allowbreak{}the\_\allowbreak{}pauline\_\allowbreak{}texts}).

The standard lexicon (BDAG) on the word has not been read, and neither have the commentaries of Brown, Schnackenburg and Barrett.

\subsubsection*{The Reformed answer: obedience is sanctification}
The Reformed do not deny what Chrysostom’s texts say. Whoever does not obey the Son shall not see life (John 3:36), and only the one who does the will of the Father enters the kingdom (Matthew 7:21). They deny what he concludes from them.

On their reading the obedience these texts require is sanctification: the life of good works, inward and outward, that follows faith as its fruit and its evidence. It is necessary, for a faith that bears no such fruit is dead. But it is not the ground of justification, and it is not a second thing added to faith as its means. Modern Reformed and evangelical scholars put each part of this.

\begin{itemize}
  \item On John 3:36, Ridderbos: “‘Disobeys’ serves as a description of ‘unbelief’”. Carson: “failure to trust him is as much disobedience as unbelief”.
  \item On Matthew 7:21, Morris: “This is not salvation by works: the contrast is not between merit and grace, but between profession and way of life.”
  \item On the place of works, Schreiner: “the works necessary for justification aren’t the basis for salvation and in that sense don’t contribute to salvation. They are the necessary evidence for and fruit of salvation.” And: “the necessity of good works doesn’t contradict sola fide” (p. 55). Hoekema: good works “must not be thought of as meritorious ... Yet they are necessary”. Murray: “It is living faith that justifies”.
\end{itemize}
Westminster says the same: good works “are the fruits and evidences of a true and lively faith” (XVI.2), and the sanctification of the justified is distinct from their justification (XIII.1, XI.1). So, read whole, John 3:36 says what its first half says in terms: “Whoever believes in the Son has eternal life.”

The reading is written out as a line of its own (\texttt{sanctification\allowbreak{}Line}), and heard in the main debate as the Reformed answer to Chrysostom. It is rated \texttt{disputed}. Chrysostom accepts both texts and reads them the other way, and so does a modern reader: Stanley holds that “Endurance is a condition for final salvation”, and reads Matthew 7:21 as making doing the Father’s will a condition and not only its evidence. The commentaries above were read through their e-book editions, so they are cited by verse and not by page.

\subsubsection*{How this is checked}
\paragraph*{The lines}
\leavevmode\par\nobreak
\par\medskip\noindent\textbf{\texttt{aquinas\allowbreak{}On\allowbreak{}Formed\allowbreak{}Faith}}\par\nobreak\noindent
Aquinas’s step: faith made alive by love is not bare trust. Rated \texttt{consensus}: Calvin accepts it, since the trust he means is exactly a faith that brings nothing. The two divide over the premise, not the step.

{\small\begin{verbatim}
def aquinasOnFormedFaith : Source :=
  {
    primary :=
      Reference.work aquinasJohn
        (Bib.Locus.sectionRef "cap. 6, lect. 3, n. 901"),
    tradition := Tradition.romanCatholic, confidence := Confidence.consensus }
\end{verbatim}}
\par\medskip\noindent\textbf{\texttt{thomist\allowbreak{}Line}} --- Faith formed by charity (Aquinas on John 6:29)\par\nobreak\noindent
\textbf{Aquinas’s line of argument.} The believing of 6:29 is faith made alive by love, and so not the bare trust that John’s strand needs. The denial is derived from Aquinas’s own reading of the verse, not assumed.

\begin{derivation}
  \item[(1)] \(\pvar{TestimonyArguments.SolaFide}{41}\)\gloss{The believing of John 6:29 is faith living through charity, the source of works}
  \item[(2)] \(\pvar{TestimonyArguments.SolaFide}{41} \rightarrow \lnot \pvar{TestimonyArguments.SolaFide}{40}\)
  \item[\(\vdash\)] \(\lnot \pvar{TestimonyArguments.SolaFide}{40}\)\gloss{\emph{at issue, the denial of:} The believing of John 6:29 is trust, which brings nothing and receives Christ}
\end{derivation}
\paragraph*{The packages}
\leavevmode\par\nobreak
\par\medskip\noindent\textbf{\texttt{johannine\allowbreak{}Case}} --- Sola fide from John (6:28–29)\par\nobreak\noindent
Sola fide from John alone: John’s line, with the premises and closing steps every Reformed strand shares.

\begin{derivation}
  \item[(1)] \(\pvar{TestimonyArguments.SolaFide}{36}\)\gloss{John 6:28–29 — the work God requires is that you believe in him whom he sent}
  \item[(2)] \(\pvar{TestimonyArguments.SolaFide}{37}\)\gloss{John 3:16–18, 3:36, 5:24, 20:31 — eternal life through believing in the Son}
  \item[(3)] \(\pvar{TestimonyArguments.SolaFide}{40}\)\gloss{The believing of John 6:29 is trust, which brings nothing and receives Christ}
  \breakrow{(4)} \(\pvar{TestimonyArguments.SolaFide}{1}\)\gloss{Ephesians 2:8–9 teaches salvation by grace through faith, not of works}
  \breakrow{(5)} \(\pvar{TestimonyArguments.SolaFide}{2}\)\gloss{Romans 3:28 teaches justification by faith apart from works of the law}
  \breakrow{(6)} \(\pvar{TestimonyArguments.SolaFide}{3}\)\gloss{Galatians 2:16 teaches that no one is justified by works of the law}
  \breakrow{(7)} \(\pvar{TestimonyArguments.SolaFide}{4}\)\gloss{Romans 4:4–5 contrasts wages owed with a gift reckoned to the one who believes}
  \breakrow{(8)} \(\pvar{TestimonyArguments.SolaFide}{5}\)\gloss{Titus 3:5 teaches that God saved us not by works done in righteousness}
  \breakrow{(9)} \(\pvar{TestimonyArguments.SolaFide}{17}\)\gloss{Luke 7:50 — Jesus says ‘your faith has saved you’ after declaring sins forgiven}
  \breakrow{(10)} \(\pvar{TestimonyArguments.SolaFide}{9}\)\gloss{Acts 15:9–11 — Peter: hearts cleansed by faith; saved through grace}
  \breakrow{(11)} \(\pvar{TestimonyArguments.SolaFide}{20}\)\gloss{James 2:14–26 targets a barren faith — mere assent — not Paul’s doctrine}
  \breakrow{(12)} \(\pvar{TestimonyArguments.SolaFide}{21}\)\gloss{Good works are the fruit and evidence of saving faith, not its ground}
  \breakrow{(13)} \(\pvar{TestimonyArguments.SolaFide}{24}\)\gloss{Scripture does not contradict itself}
  \breakrow{(14)} \((\pvar{TestimonyArguments.SolaFide}{36} \land \pvar{TestimonyArguments.SolaFide}{37} \land \pvar{TestimonyArguments.SolaFide}{40}) \rightarrow \pvar{TestimonyArguments.SolaFide}{25}\)\gloss{\emph{would give:} Justification is by faith alone}
  \breakrow{(15)} \((\pvar{TestimonyArguments.SolaFide}{20} \land \pvar{TestimonyArguments.SolaFide}{21}) \rightarrow \pvar{TestimonyArguments.SolaFide}{22}\)\gloss{\emph{would give:} James 2:24 is compatible with Paul, using ‘justify’ in a different sense}
  \breakrow{(16)} \((\pvar{TestimonyArguments.SolaFide}{1} \land \pvar{TestimonyArguments.SolaFide}{4} \land \pvar{TestimonyArguments.SolaFide}{5}) \rightarrow \pvar{TestimonyArguments.SolaFide}{27}\)\gloss{\emph{would give:} Salvation is by grace, a gift and not wages owed}
  \breakrow{(17)} \((\pvar{TestimonyArguments.SolaFide}{1} \land \pvar{TestimonyArguments.SolaFide}{4} \land \pvar{TestimonyArguments.SolaFide}{5} \land \pvar{TestimonyArguments.SolaFide}{22} \land \pvar{TestimonyArguments.SolaFide}{24}) \rightarrow \pvar{TestimonyArguments.SolaFide}{28}\)\gloss{\emph{would give:} Salvation is not by works: no work is its ground}
  \breakrow{(18)} \(\pvar{TestimonyArguments.SolaFide}{25} \rightarrow \pvar{TestimonyArguments.SolaFide}{26}\)\gloss{\emph{would give:} Faith is sufficient: whoever believes is saved}
  \item[(19)] \((\pvar{TestimonyArguments.SolaFide}{26} \land \pvar{TestimonyArguments.SolaFide}{1}) \rightarrow \pvar{TestimonyArguments.SolaFide}{29}\)\gloss{\emph{would give:} Salvation is received through faith}
  \item[\(\vdash\)] \(\pvar{TestimonyArguments.SolaFide}{27} \land \pvar{TestimonyArguments.SolaFide}{28} \land \pvar{TestimonyArguments.SolaFide}{29}\)\gloss{\emph{at issue:} \(\pvar{TestimonyArguments.SolaFide}{27}\) Salvation is by grace, a gift and not wages owed; \(\pvar{TestimonyArguments.SolaFide}{28}\) Salvation is not by works: no work is its ground; \(\pvar{TestimonyArguments.SolaFide}{29}\) Salvation is received through faith}
  \keepwithnext
\end{derivation}
\emph{Grounded in scripture alone:} John 3:16–18, 3:36, 5:24, 20:31 — eternal life through believing in the Son.
\par\medskip\noindent\textbf{\texttt{thomist\allowbreak{}On\allowbreak{}John}} --- Aquinas on John 6:29 (faith formed by charity)\par\nobreak\noindent
Aquinas on John: every text John’s strand reads, and every shared premise, with the believing read as faith formed by charity. This is the rival John’s strand was written against.

\begin{derivation}
  \item[(1)] \(\pvar{TestimonyArguments.SolaFide}{36}\)\gloss{John 6:28–29 — the work God requires is that you believe in him whom he sent}
  \item[(2)] \(\pvar{TestimonyArguments.SolaFide}{37}\)\gloss{John 3:16–18, 3:36, 5:24, 20:31 — eternal life through believing in the Son}
  \item[(3)] \(\pvar{TestimonyArguments.SolaFide}{41}\)\gloss{The believing of John 6:29 is faith living through charity, the source of works}
  \breakrow{(4)} \(\pvar{TestimonyArguments.SolaFide}{1}\)\gloss{Ephesians 2:8–9 teaches salvation by grace through faith, not of works}
  \breakrow{(5)} \(\pvar{TestimonyArguments.SolaFide}{2}\)\gloss{Romans 3:28 teaches justification by faith apart from works of the law}
  \breakrow{(6)} \(\pvar{TestimonyArguments.SolaFide}{3}\)\gloss{Galatians 2:16 teaches that no one is justified by works of the law}
  \breakrow{(7)} \(\pvar{TestimonyArguments.SolaFide}{4}\)\gloss{Romans 4:4–5 contrasts wages owed with a gift reckoned to the one who believes}
  \breakrow{(8)} \(\pvar{TestimonyArguments.SolaFide}{5}\)\gloss{Titus 3:5 teaches that God saved us not by works done in righteousness}
  \breakrow{(9)} \(\pvar{TestimonyArguments.SolaFide}{17}\)\gloss{Luke 7:50 — Jesus says ‘your faith has saved you’ after declaring sins forgiven}
  \breakrow{(10)} \(\pvar{TestimonyArguments.SolaFide}{9}\)\gloss{Acts 15:9–11 — Peter: hearts cleansed by faith; saved through grace}
  \breakrow{(11)} \(\pvar{TestimonyArguments.SolaFide}{20}\)\gloss{James 2:14–26 targets a barren faith — mere assent — not Paul’s doctrine}
  \breakrow{(12)} \(\pvar{TestimonyArguments.SolaFide}{21}\)\gloss{Good works are the fruit and evidence of saving faith, not its ground}
  \breakrow{(13)} \(\pvar{TestimonyArguments.SolaFide}{24}\)\gloss{Scripture does not contradict itself}
  \breakrow{(14)} \((\pvar{TestimonyArguments.SolaFide}{36} \land \pvar{TestimonyArguments.SolaFide}{37} \land \pvar{TestimonyArguments.SolaFide}{40}) \rightarrow \pvar{TestimonyArguments.SolaFide}{25}\)\gloss{\emph{would give:} Justification is by faith alone}
  \breakrow{(15)} \(\pvar{TestimonyArguments.SolaFide}{41} \rightarrow \lnot \pvar{TestimonyArguments.SolaFide}{40}\)\gloss{\emph{would deny:} The believing of John 6:29 is trust, which brings nothing and receives Christ}
  \breakrow{(16)} \((\pvar{TestimonyArguments.SolaFide}{20} \land \pvar{TestimonyArguments.SolaFide}{21}) \rightarrow \pvar{TestimonyArguments.SolaFide}{22}\)\gloss{\emph{would give:} James 2:24 is compatible with Paul, using ‘justify’ in a different sense}
  \breakrow{(17)} \((\pvar{TestimonyArguments.SolaFide}{1} \land \pvar{TestimonyArguments.SolaFide}{4} \land \pvar{TestimonyArguments.SolaFide}{5}) \rightarrow \pvar{TestimonyArguments.SolaFide}{27}\)\gloss{\emph{would give:} Salvation is by grace, a gift and not wages owed}
  \breakrow{(18)} \((\pvar{TestimonyArguments.SolaFide}{1} \land \pvar{TestimonyArguments.SolaFide}{4} \land \pvar{TestimonyArguments.SolaFide}{5} \land \pvar{TestimonyArguments.SolaFide}{22} \land \pvar{TestimonyArguments.SolaFide}{24}) \rightarrow \pvar{TestimonyArguments.SolaFide}{28}\)\gloss{\emph{would give:} Salvation is not by works: no work is its ground}
  \breakrow{(19)} \(\pvar{TestimonyArguments.SolaFide}{25} \rightarrow \pvar{TestimonyArguments.SolaFide}{26}\)\gloss{\emph{would give:} Faith is sufficient: whoever believes is saved}
  \item[(20)] \((\pvar{TestimonyArguments.SolaFide}{26} \land \pvar{TestimonyArguments.SolaFide}{1}) \rightarrow \pvar{TestimonyArguments.SolaFide}{29}\)\gloss{\emph{would give:} Salvation is received through faith}
  \item[\(\vdash\)] \(\pvar{TestimonyArguments.SolaFide}{27} \land \pvar{TestimonyArguments.SolaFide}{28} \land \pvar{TestimonyArguments.SolaFide}{29}\)\gloss{\emph{at issue:} \(\pvar{TestimonyArguments.SolaFide}{27}\) Salvation is by grace, a gift and not wages owed; \(\pvar{TestimonyArguments.SolaFide}{28}\) Salvation is not by works: no work is its ground; \(\pvar{TestimonyArguments.SolaFide}{29}\) Salvation is received through faith}
  \keepwithnext
\end{derivation}
\emph{Grounded in scripture alone:} John 3:16–18, 3:36, 5:24, 20:31 — eternal life through believing in the Son.
\paragraph*{Results}
\leavevmode\par\nobreak
\par\medskip\noindent\textbf{\texttt{johannine\_\allowbreak{}strand\_\allowbreak{}establishes}}\par\nobreak\noindent
\textbf{John’s strand, argued alone, delivers sola fide.} Grant that the believing of 6:29 is trust, and the work God requires is believing, which earns nothing. Sola fide then follows without the readings Paul’s, Luke’s and Peter’s strands depend on.

It does not show that John 6:29 is the only text used. Grace and “not by works” come from the verses every strand shares and from the answer to James.

{\small\begin{verbatim}
theorem johannine_strand_establishes : Establishes johannineCase
-- axioms: propext, Quot.sound
\end{verbatim}}
\par\medskip\noindent\textbf{\texttt{thomist\allowbreak{}Reading}}\par\nobreak\noindent
Aquinas’s reading, written out. Every text holds. The believing of 6:29 is faith made alive by love, so it is not bare trust; justification is not by faith alone; and faith without charity does not suffice. Grace and “not by works” still hold: Aquinas excludes outward works, and grants that faith is God’s gift.

{\small\begin{verbatim}
def thomistReading : Valuation Claim :=
  fun a =>
    match a with
    | Claim.johannineBelievingIsTrust => False
    | Claim.justificationByFaithAlone => False
    | Claim.faithIsSufficient => False
    | Claim.salvationThroughFaith => False
    | x => True
\end{verbatim}}
\par\medskip\noindent\textbf{\texttt{johannine\_\allowbreak{}strand\_\allowbreak{}rests\_\allowbreak{}on\_\allowbreak{}believing\_\allowbreak{}as\_\allowbreak{}trust}}\par\nobreak\noindent
\textbf{John’s strand depends on reading believing as trust.} Grant Aquinas every text the strand reads, and read their believing as he does, faith made alive by love, and sola fide no longer follows.

So John’s Gospel, like Paul’s Galatians and Luke 7:50, reaches the conclusion only through a reading: here, of what it is to believe \emph{in} him.

{\small\begin{verbatim}
theorem johannine_strand_rests_on_believing_as_trust : ¬Establishes
    thomistOnJohn
-- axioms: propext, Quot.sound
\end{verbatim}}
\par\medskip\noindent\textbf{\texttt{reformed\allowbreak{}Granting\allowbreak{}Aquinas\allowbreak{}On\allowbreak{}John}} --- Reformed, granting Aquinas on John 6:29\par\nobreak\noindent
The Reformed case with Aquinas granted his reading of John 6:29: the believing is faith formed by charity, and so not bare trust. The other three strands are kept whole.

\begin{derivation}
  \item[(1)] \(\pvar{TestimonyArguments.SolaFide}{8}\)\gloss{Galatians opposes requiring circumcision, besides faith, for justification}
  \item[(2)] \(\pvar{TestimonyArguments.SolaFide}{6}\)\gloss{Paul’s ‘works of the law’ means human works in general, not boundary markers}
  \item[(3)] \(\pvar{TestimonyArguments.SolaFide}{7}\)\gloss{In Galatians 2:16 ‘faith of Christ’ (πίστις Χριστοῦ) means our faith in Christ}
  \breakrow{(4)} \(\pvar{TestimonyArguments.SolaFide}{18}\)\gloss{‘Saved’ (σῴζω) at Luke 7:50 means saved from sin, not healed}
  \breakrow{(5)} \(\pvar{TestimonyArguments.SolaFide}{19}\)\gloss{Luke 7:47 — her love is the evidence of her forgiveness, not its ground}
  \breakrow{(6)} \(\pvar{TestimonyArguments.SolaFide}{10}\)\gloss{The yoke refused at Acts 15:10 is the whole law as a condition of salvation}
  \breakrow{(7)} \(\pvar{TestimonyArguments.SolaFide}{36}\)\gloss{John 6:28–29 — the work God requires is that you believe in him whom he sent}
  \breakrow{(8)} \(\pvar{TestimonyArguments.SolaFide}{37}\)\gloss{John 3:16–18, 3:36, 5:24, 20:31 — eternal life through believing in the Son}
  \breakrow{(9)} \(\pvar{TestimonyArguments.SolaFide}{41}\)\gloss{The believing of John 6:29 is faith living through charity, the source of works}
  \breakrow{(10)} \(\pvar{TestimonyArguments.SolaFide}{1}\)\gloss{Ephesians 2:8–9 teaches salvation by grace through faith, not of works}
  \breakrow{(11)} \(\pvar{TestimonyArguments.SolaFide}{2}\)\gloss{Romans 3:28 teaches justification by faith apart from works of the law}
  \breakrow{(12)} \(\pvar{TestimonyArguments.SolaFide}{3}\)\gloss{Galatians 2:16 teaches that no one is justified by works of the law}
  \breakrow{(13)} \(\pvar{TestimonyArguments.SolaFide}{4}\)\gloss{Romans 4:4–5 contrasts wages owed with a gift reckoned to the one who believes}
  \breakrow{(14)} \(\pvar{TestimonyArguments.SolaFide}{5}\)\gloss{Titus 3:5 teaches that God saved us not by works done in righteousness}
  \breakrow{(15)} \(\pvar{TestimonyArguments.SolaFide}{17}\)\gloss{Luke 7:50 — Jesus says ‘your faith has saved you’ after declaring sins forgiven}
  \breakrow{(16)} \(\pvar{TestimonyArguments.SolaFide}{9}\)\gloss{Acts 15:9–11 — Peter: hearts cleansed by faith; saved through grace}
  \breakrow{(17)} \(\pvar{TestimonyArguments.SolaFide}{20}\)\gloss{James 2:14–26 targets a barren faith — mere assent — not Paul’s doctrine}
  \breakrow{(18)} \(\pvar{TestimonyArguments.SolaFide}{21}\)\gloss{Good works are the fruit and evidence of saving faith, not its ground}
  \breakrow{(19)} \(\pvar{TestimonyArguments.SolaFide}{24}\)\gloss{Scripture does not contradict itself}
  \breakrow{(20)} \((\pvar{TestimonyArguments.SolaFide}{2} \land \pvar{TestimonyArguments.SolaFide}{3} \land \pvar{TestimonyArguments.SolaFide}{8} \land \pvar{TestimonyArguments.SolaFide}{7} \land \pvar{TestimonyArguments.SolaFide}{6}) \rightarrow \pvar{TestimonyArguments.SolaFide}{25}\)\gloss{\emph{would give:} Justification is by faith alone}
  \breakrow{(21)} \((\pvar{TestimonyArguments.SolaFide}{17} \land \pvar{TestimonyArguments.SolaFide}{18} \land \pvar{TestimonyArguments.SolaFide}{19}) \rightarrow \pvar{TestimonyArguments.SolaFide}{26}\)\gloss{\emph{would give:} Faith is sufficient: whoever believes is saved}
  \breakrow{(22)} \((\pvar{TestimonyArguments.SolaFide}{9} \land \pvar{TestimonyArguments.SolaFide}{10}) \rightarrow \pvar{TestimonyArguments.SolaFide}{26}\)\gloss{\emph{would give:} Faith is sufficient: whoever believes is saved}
  \breakrow{(23)} \((\pvar{TestimonyArguments.SolaFide}{36} \land \pvar{TestimonyArguments.SolaFide}{37} \land \pvar{TestimonyArguments.SolaFide}{40}) \rightarrow \pvar{TestimonyArguments.SolaFide}{25}\)\gloss{\emph{would give:} Justification is by faith alone}
  \breakrow{(24)} \(\pvar{TestimonyArguments.SolaFide}{41} \rightarrow \lnot \pvar{TestimonyArguments.SolaFide}{40}\)\gloss{\emph{would deny:} The believing of John 6:29 is trust, which brings nothing and receives Christ}
  \breakrow{(25)} \((\pvar{TestimonyArguments.SolaFide}{20} \land \pvar{TestimonyArguments.SolaFide}{21}) \rightarrow \pvar{TestimonyArguments.SolaFide}{22}\)\gloss{\emph{would give:} James 2:24 is compatible with Paul, using ‘justify’ in a different sense}
  \breakrow{(26)} \((\pvar{TestimonyArguments.SolaFide}{1} \land \pvar{TestimonyArguments.SolaFide}{4} \land \pvar{TestimonyArguments.SolaFide}{5}) \rightarrow \pvar{TestimonyArguments.SolaFide}{27}\)\gloss{\emph{would give:} Salvation is by grace, a gift and not wages owed}
  \breakrow{(27)} \((\pvar{TestimonyArguments.SolaFide}{1} \land \pvar{TestimonyArguments.SolaFide}{4} \land \pvar{TestimonyArguments.SolaFide}{5} \land \pvar{TestimonyArguments.SolaFide}{22} \land \pvar{TestimonyArguments.SolaFide}{24}) \rightarrow \pvar{TestimonyArguments.SolaFide}{28}\)\gloss{\emph{would give:} Salvation is not by works: no work is its ground}
  \breakrow{(28)} \(\pvar{TestimonyArguments.SolaFide}{25} \rightarrow \pvar{TestimonyArguments.SolaFide}{26}\)\gloss{\emph{would give:} Faith is sufficient: whoever believes is saved}
  \item[(29)] \((\pvar{TestimonyArguments.SolaFide}{26} \land \pvar{TestimonyArguments.SolaFide}{1}) \rightarrow \pvar{TestimonyArguments.SolaFide}{29}\)\gloss{\emph{would give:} Salvation is received through faith}
  \item[\(\vdash\)] \(\pvar{TestimonyArguments.SolaFide}{27} \land \pvar{TestimonyArguments.SolaFide}{28} \land \pvar{TestimonyArguments.SolaFide}{29}\)\gloss{\emph{at issue:} \(\pvar{TestimonyArguments.SolaFide}{27}\) Salvation is by grace, a gift and not wages owed; \(\pvar{TestimonyArguments.SolaFide}{28}\) Salvation is not by works: no work is its ground; \(\pvar{TestimonyArguments.SolaFide}{29}\) Salvation is received through faith}
  \keepwithnext
\end{derivation}
\emph{Grounded in scripture alone:} John 3:16–18, 3:36, 5:24, 20:31 — eternal life through believing in the Son.
\par\medskip\noindent\textbf{\texttt{believing\_\allowbreak{}as\_\allowbreak{}trust\_\allowbreak{}not\_\allowbreak{}load\_\allowbreak{}bearing}}\par\nobreak\noindent
\textbf{The reading of believing as trust is not one the whole case depends on.} Grant Aquinas that the believing of John 6:29 is faith formed by charity, keep Paul, Luke and Acts as the Reformed read them, and sola fide still follows.

Losing the argument over John 6:29 costs the Reformed case one strand and not its conclusion. It does not show that Aquinas would accept the rest. He would not, and the package is not his position; it is the Reformed case with one concession made.

{\small\begin{verbatim}
theorem believing_as_trust_not_load_bearing : Establishes
    reformedGrantingAquinasOnJohn
-- axioms: propext, Quot.sound
\end{verbatim}}
\par\medskip\noindent\textbf{\texttt{conceding\allowbreak{}John\allowbreak{}Reading}}\par\nobreak\noindent
The Reformed reading with one concession: the believing of John 6:29 is not bare trust, and everything else holds.

{\small\begin{verbatim}
def concedingJohnReading : Valuation Claim :=
  fun a =>
    match a with
    | Claim.johannineBelievingIsTrust => False
    | x => True
\end{verbatim}}
\par\medskip\noindent\textbf{\texttt{reformed\allowbreak{}Granting\allowbreak{}Aquinas\allowbreak{}On\allowbreak{}John\_\allowbreak{}is\_\allowbreak{}satisfiable}}\par\nobreak\noindent
That package can be held without contradiction, so \texttt{believing\_\allowbreak{}as\_\allowbreak{}trust\_\allowbreak{}not\_\allowbreak{}load\_\allowbreak{}bearing} is not true merely because its premises could never all hold.

{\small\begin{verbatim}
theorem reformedGrantingAquinasOnJohn_is_satisfiable : Satisfiable
    reformedGrantingAquinasOnJohn.premises
-- axioms: propext, Quot.sound
\end{verbatim}}
\par\medskip\noindent\textbf{\texttt{johannine\allowbreak{}Case\_\allowbreak{}is\_\allowbreak{}satisfiable}}\par\nobreak\noindent
John’s case can be held without contradiction, so \texttt{johannine\_\allowbreak{}strand\_\allowbreak{}establishes} is not true merely because its premises could never all hold.

{\small\begin{verbatim}
theorem johannineCase_is_satisfiable : Satisfiable johannineCase.premises
-- axioms: propext, Quot.sound
\end{verbatim}}
\paragraph*{Luther’s answer}
\leavevmode\par\nobreak
\par\medskip\noindent\textbf{\texttt{luther\allowbreak{}On\allowbreak{}Galatians\allowbreak{}Three}}\par\nobreak\noindent
Luther’s step, on Galatians 3:11–12: the law commands love, and the law is not of faith, so faith formed by charity would justify by the law. Rated \texttt{disputed}: Aquinas accepts both starting points, the text and that the law commands love, and denies the conclusion, holding that the faith which justifies is formed by charity.

{\small\begin{verbatim}
def lutherOnGalatiansThree : Source :=
  {
    primary :=
      Reference.work lutherGalatians
        (Bib.Locus.adLoc { book := Book.galatians, chapter := 3, verse := 12 }),
    supporting :=
      [Reference.work lutherGalatians
          (Bib.Locus.adLoc
            { book := Book.galatians, chapter := 3, verse := 11 })],
    tradition := Tradition.reformedProtestant,
    confidence := Confidence.disputed }
\end{verbatim}}
\par\medskip\noindent\textbf{\texttt{luther\allowbreak{}Line}} --- Luther on Galatians 3:11–12 (against faith formed by charity)\par\nobreak\noindent
\textbf{Luther’s line of argument.} From Galatians 3:11–12 and the law’s command of love to the denial of Aquinas’s reading: the believing that justifies is not faith formed by charity, since love is what the law commands and the law is not of faith.

\begin{derivation}
  \item[(1)] \(\pvar{TestimonyArguments.SolaFide}{42}\)\gloss{Galatians 3:11–12 — none is justified by the law; the law is not of faith}
  \item[(2)] \(\pvar{TestimonyArguments.SolaFide}{43}\)\gloss{Love of God and neighbour is what the law commands}
  \item[(3)] \((\pvar{TestimonyArguments.SolaFide}{42} \land \pvar{TestimonyArguments.SolaFide}{43}) \rightarrow \lnot \pvar{TestimonyArguments.SolaFide}{41}\)
  \item[\(\vdash\)] \(\lnot \pvar{TestimonyArguments.SolaFide}{41}\)\gloss{\emph{at issue, the denial of:} The believing of John 6:29 is faith living through charity, the source of works}
\end{derivation}
\par\medskip\noindent\textbf{\texttt{luther\allowbreak{}On\allowbreak{}Galatians}} --- Luther on Galatians 3:11–12 (against faith formed by charity)\par\nobreak\noindent
Luther’s answer to Aquinas as a position: what Aquinas grants, with Luther’s step, against reading the believing of 6:29 as faith formed by charity.

\begin{derivation}
  \item[(1)] \(\pvar{TestimonyArguments.SolaFide}{42}\)\gloss{Galatians 3:11–12 — none is justified by the law; the law is not of faith}
  \item[(2)] \(\pvar{TestimonyArguments.SolaFide}{43}\)\gloss{Love of God and neighbour is what the law commands}
  \item[(3)] \((\pvar{TestimonyArguments.SolaFide}{42} \land \pvar{TestimonyArguments.SolaFide}{43}) \rightarrow \lnot \pvar{TestimonyArguments.SolaFide}{41}\)
  \item[\(\vdash\)] \(\lnot \pvar{TestimonyArguments.SolaFide}{41}\)\gloss{\emph{at issue, the denial of:} The believing of John 6:29 is faith living through charity, the source of works}
  \keepwithnext
\end{derivation}
\emph{No premise here rests on scripture alone.}
\par\medskip\noindent\textbf{\texttt{galatians\allowbreak{}Texts\allowbreak{}Alone}} --- Galatians 3:11–12 and the law’s command of love, without Luther’s step\par\nobreak\noindent
The same starting points without Luther’s step: Galatians 3:11–12, and that the law commands love.

\begin{derivation}
  \item[(1)] \(\pvar{TestimonyArguments.SolaFide}{42}\)\gloss{Galatians 3:11–12 — none is justified by the law; the law is not of faith}
  \item[(2)] \(\pvar{TestimonyArguments.SolaFide}{43}\)\gloss{Love of God and neighbour is what the law commands}
  \item[\(\vdash\)] \(\lnot \pvar{TestimonyArguments.SolaFide}{41}\)\gloss{\emph{at issue, the denial of:} The believing of John 6:29 is faith living through charity, the source of works}
  \keepwithnext
\end{derivation}
\emph{No premise here rests on scripture alone.}
\par\medskip\noindent\textbf{\texttt{luther\_\allowbreak{}answers\_\allowbreak{}aquinas\_\allowbreak{}from\_\allowbreak{}galatians}}\par\nobreak\noindent
\textbf{Luther answers Aquinas from Galatians 3.} Grant what Aquinas grants, Galatians 3:11–12 and that the law commands love, and grant Luther’s step. Then the believing of John 6:29 is not faith formed by charity: a faith that justified because love formed it would justify by the law, and the law is not of faith.

It does not show that Calvin’s reading is thereby established. It removes Aquinas’s reading, and only if the step is granted.

{\small\begin{verbatim}
theorem luther_answers_aquinas_from_galatians : Establishes lutherOnGalatians
-- axioms: propext, Quot.sound
\end{verbatim}}
\par\medskip\noindent\textbf{\texttt{luther\_\allowbreak{}answer\_\allowbreak{}rests\_\allowbreak{}on\_\allowbreak{}his\_\allowbreak{}step}}\par\nobreak\noindent
\textbf{Luther’s answer depends on his step.} The texts alone do not exclude Aquinas’s reading. On Aquinas’s reading Galatians 3:11–12 holds, the law commands love, and the believing of 6:29 is faith formed by charity. What decides it is whether faith formed by charity counts as justification “by the law”. That is Luther’s step, and Aquinas denies it.

{\small\begin{verbatim}
theorem luther_answer_rests_on_his_step : ¬Establishes galatiansTextsAlone
-- axioms: propext, Quot.sound
\end{verbatim}}
\par\medskip\noindent\textbf{\texttt{luther\allowbreak{}Reading}}\par\nobreak\noindent
The Reformed reading, on which the believing of John 6:29 is not faith formed by charity, and everything else holds.

{\small\begin{verbatim}
def lutherReading : Valuation Claim :=
  fun a =>
    match a with
    | Claim.johannineBelievingIsFormedByCharity => False
    | x => True
\end{verbatim}}
\par\medskip\noindent\textbf{\texttt{luther\allowbreak{}On\allowbreak{}Galatians\_\allowbreak{}is\_\allowbreak{}satisfiable}}\par\nobreak\noindent
Luther’s position can be held without contradiction, so \texttt{luther\_\allowbreak{}answers\_\allowbreak{}aquinas\_\allowbreak{}from\_\allowbreak{}galatians} is not true merely because its premises could never all hold.

{\small\begin{verbatim}
theorem lutherOnGalatians_is_satisfiable : Satisfiable
    lutherOnGalatians.premises
-- axioms: propext, Quot.sound
\end{verbatim}}
\par\medskip\noindent\textbf{\texttt{why\allowbreak{}Luther\allowbreak{}Answers\allowbreak{}Aquinas}}\par\nobreak\noindent
\textbf{Why Luther’s answer stands against Aquinas.} Everything turns on Luther’s step: from Galatians 3:11–12 and the law’s command of love, the believing that justifies is not faith formed by charity. With the two texts accepted, Aquinas’s reading of John 6:29 cannot keep its one distinctive premise, that the believing is formed by charity.

Unlike the Reformed distinction against Trent, this step is part of Luther’s own case: without it, the texts do not deliver his conclusion. So the answer is exactly as strong as the step, which is rated \texttt{disputed}. The step is Luther’s, and the texts are accepted by both sides.

\paragraph{Why \emph{Luther on Galatians 3:11–12 (against faith formed by charity)} stands against \emph{Aquinas on John 6:29 (faith formed by charity)}.}
\begin{itemize}
\item \textbf{The crux:} \((\pvar{TestimonyArguments.SolaFide}{42} \land \pvar{TestimonyArguments.SolaFide}{43}) \rightarrow \lnot \pvar{TestimonyArguments.SolaFide}{41}\), a premise of \emph{Luther on Galatians 3:11–12 (against faith formed by charity)}.
\item \textbf{What it does:} its conclusion does not follow without it.
\item \textbf{Granted:} \(\pvar{TestimonyArguments.SolaFide}{42}\), \(\pvar{TestimonyArguments.SolaFide}{43}\), which the break also needs.
\item \textbf{Where the rival breaks:} \(\pvar{TestimonyArguments.SolaFide}{41}\) cannot be held together with \((\pvar{TestimonyArguments.SolaFide}{42} \land \pvar{TestimonyArguments.SolaFide}{43}) \rightarrow \lnot \pvar{TestimonyArguments.SolaFide}{41}\) and \(\pvar{TestimonyArguments.SolaFide}{42}\), \(\pvar{TestimonyArguments.SolaFide}{43}\); each is needed for the break, and without the crux they stand.
\item \textbf{What this rests on:}
\begin{itemize}
\item \(\pvar{TestimonyArguments.SolaFide}{42}\) Galatians 3:11–12 — none is justified by the law; the law is not of faith --- \emph{consensus}: Gal 3:11-12; \cite[ad loc. Gal 3:11]{luther-commentary-galatians-1998}
\item \(\pvar{TestimonyArguments.SolaFide}{43}\) Love of God and neighbour is what the law commands --- \emph{consensus}: Deut 6:5; Matt 22:37-40; \cite[ad loc. Gal 3:12]{luther-commentary-galatians-1998}
\item the step itself --- \emph{disputed}: \cite[ad loc. Gal 3:12]{luther-commentary-galatians-1998}; \cite[ad loc. Gal 3:11]{luther-commentary-galatians-1998}
\end{itemize}
\end{itemize}
\par\medskip\noindent\textbf{\texttt{chrysostom\allowbreak{}On\allowbreak{}John\allowbreak{}Three\allowbreak{}Thirty\allowbreak{}Six}}\par\nobreak\noindent
Chrysostom’s step, in Homily 31 on John: from John 3:36 and Matthew 7:21 to “faith alone is not sufficient to salvation”. Stanley, a modern reader, holds works “a condition to final salvation” in the Synoptic Gospels, Matthew 7:21 among the texts. Rated \texttt{disputed}: Calvin accepts both texts and denies the conclusion, and so do the modern Reformed readers cited for the answer.

{\small\begin{verbatim}
def chrysostomOnJohnThreeThirtySix : Source :=
  {
    primary :=
      Reference.work chrysostomHomiliesJohn
        (Bib.Locus.sectionRef "Homily 31, §1"),
    supporting := [Reference.work stanleyDidJesusTeach],
    tradition := Tradition.easternOrthodox, confidence := Confidence.disputed }
\end{verbatim}}
\par\medskip\noindent\textbf{\texttt{chrysostom\allowbreak{}Line}} --- Chrysostom on John 3:36 (believing is not enough without a right life)\par\nobreak\noindent
\textbf{Chrysostom’s line.} From “whoever does not obey the Son shall not see life” (John 3:36) and “not everyone who says to me, ‘Lord, Lord’” (Matthew 7:21) to the conclusion that believing does not suffice without a right life.

\begin{derivation}
  \item[(1)] \(\pvar{TestimonyArguments.SolaFide}{38}\)\gloss{John 3:36 — whoever believes has life; whoever does not obey shall not see life}
  \item[(2)] \(\pvar{TestimonyArguments.SolaFide}{62}\)\gloss{Matthew 7:21–23 — the criterion at the judgment is doing the Father’s will, not profession or mighty works}
  \item[(3)] \((\pvar{TestimonyArguments.SolaFide}{38} \land \pvar{TestimonyArguments.SolaFide}{62}) \rightarrow \lnot \pvar{TestimonyArguments.SolaFide}{26}\)
  \item[\(\vdash\)] \(\lnot \pvar{TestimonyArguments.SolaFide}{26}\)\gloss{\emph{at issue, the denial of:} Faith is sufficient: whoever believes is saved}
\end{derivation}
\par\medskip\noindent\textbf{\texttt{chrysostom\allowbreak{}On\allowbreak{}John}} --- Chrysostom on John 3:36 (believing is not enough without a right life)\par\nobreak\noindent
Chrysostom’s argument as a position: the two texts, with his step.

\begin{derivation}
  \item[(1)] \(\pvar{TestimonyArguments.SolaFide}{38}\)\gloss{John 3:36 — whoever believes has life; whoever does not obey shall not see life}
  \item[(2)] \(\pvar{TestimonyArguments.SolaFide}{62}\)\gloss{Matthew 7:21–23 — the criterion at the judgment is doing the Father’s will, not profession or mighty works}
  \item[(3)] \((\pvar{TestimonyArguments.SolaFide}{38} \land \pvar{TestimonyArguments.SolaFide}{62}) \rightarrow \lnot \pvar{TestimonyArguments.SolaFide}{26}\)
  \item[\(\vdash\)] \(\lnot \pvar{TestimonyArguments.SolaFide}{26}\)\gloss{\emph{at issue, the denial of:} Faith is sufficient: whoever believes is saved}
  \keepwithnext
\end{derivation}
\emph{No premise here rests on scripture alone.}
\par\medskip\noindent\textbf{\texttt{chrysostom\_\allowbreak{}reads\_\allowbreak{}john3\_\allowbreak{}36\_\allowbreak{}against\_\allowbreak{}faith\_\allowbreak{}alone}}\par\nobreak\noindent
\textbf{Granted his step, Chrysostom’s conclusion follows.} From John 3:36 and Matthew 7:21, read as he reads them, believing does not suffice for salvation without a right life.

His argument is valid, and its premises are two texts nobody here disputes. It does not show that the step is sound. Calvin accepts both texts and denies it, so the step is rated \texttt{disputed}.

{\small\begin{verbatim}
theorem chrysostom_reads_john3_36_against_faith_alone : Establishes
    chrysostomOnJohn
-- axioms: propext, Quot.sound
\end{verbatim}}
\par\medskip\noindent\textbf{\texttt{chrysostom\allowbreak{}On\allowbreak{}John\_\allowbreak{}is\_\allowbreak{}satisfiable}}\par\nobreak\noindent
Chrysostom’s position can be held without contradiction, on Aquinas’s reading, where every text holds and faith without charity does not suffice. So \texttt{chrysostom\_\allowbreak{}reads\_\allowbreak{}john3\_\allowbreak{}36\_\allowbreak{}against\_\allowbreak{}faith\_\allowbreak{}alone} is not true merely because its premises could never all hold.

{\small\begin{verbatim}
theorem chrysostomOnJohn_is_satisfiable : Satisfiable
    chrysostomOnJohn.premises
-- axioms: propext, Quot.sound
\end{verbatim}}
\par\medskip\noindent\textbf{\texttt{john3\_\allowbreak{}36\_\allowbreak{}leaves\_\allowbreak{}faiths\_\allowbreak{}sufficiency\_\allowbreak{}open}}\par\nobreak\noindent
\textbf{The wording of John 3:36 and Matthew 7:21 leaves open whether faith suffices.} One consistent reading holds both texts and that faith suffices: the Reformed one, on which not obeying the Son is refusing to believe him. Another holds both texts and that faith does not suffice: Chrysostom’s and Aquinas’s.

Neither side can settle the question by pointing at the words. It does not show that the two readings are equally good. That is a matter for the sources, and the page reports what they say.

{\small\begin{verbatim}
theorem john3_36_leaves_faiths_sufficiency_open : Independent [p
    Claim.john3_36WhoeverDoesNotObey, p Claim.matthew7_21_23DoingTheWill] (p
    Claim.faithIsSufficient)
-- axioms: propext, Quot.sound
\end{verbatim}}
\par\medskip\noindent\textbf{\texttt{john\allowbreak{}With\allowbreak{}Chrysostom}}\par\nobreak\noindent
John’s strand and Chrysostom’s argument, held together: every premise of each.

\begin{derivation}
  \item[(1)] \(\pvar{TestimonyArguments.SolaFide}{36}\)\gloss{John 6:28–29 — the work God requires is that you believe in him whom he sent}
  \item[(2)] \(\pvar{TestimonyArguments.SolaFide}{37}\)\gloss{John 3:16–18, 3:36, 5:24, 20:31 — eternal life through believing in the Son}
  \item[(3)] \(\pvar{TestimonyArguments.SolaFide}{40}\)\gloss{The believing of John 6:29 is trust, which brings nothing and receives Christ}
  \breakrow{(4)} \(\pvar{TestimonyArguments.SolaFide}{1}\)\gloss{Ephesians 2:8–9 teaches salvation by grace through faith, not of works}
  \breakrow{(5)} \(\pvar{TestimonyArguments.SolaFide}{2}\)\gloss{Romans 3:28 teaches justification by faith apart from works of the law}
  \breakrow{(6)} \(\pvar{TestimonyArguments.SolaFide}{3}\)\gloss{Galatians 2:16 teaches that no one is justified by works of the law}
  \breakrow{(7)} \(\pvar{TestimonyArguments.SolaFide}{4}\)\gloss{Romans 4:4–5 contrasts wages owed with a gift reckoned to the one who believes}
  \breakrow{(8)} \(\pvar{TestimonyArguments.SolaFide}{5}\)\gloss{Titus 3:5 teaches that God saved us not by works done in righteousness}
  \breakrow{(9)} \(\pvar{TestimonyArguments.SolaFide}{17}\)\gloss{Luke 7:50 — Jesus says ‘your faith has saved you’ after declaring sins forgiven}
  \breakrow{(10)} \(\pvar{TestimonyArguments.SolaFide}{9}\)\gloss{Acts 15:9–11 — Peter: hearts cleansed by faith; saved through grace}
  \breakrow{(11)} \(\pvar{TestimonyArguments.SolaFide}{20}\)\gloss{James 2:14–26 targets a barren faith — mere assent — not Paul’s doctrine}
  \breakrow{(12)} \(\pvar{TestimonyArguments.SolaFide}{21}\)\gloss{Good works are the fruit and evidence of saving faith, not its ground}
  \breakrow{(13)} \(\pvar{TestimonyArguments.SolaFide}{24}\)\gloss{Scripture does not contradict itself}
  \breakrow{(14)} \((\pvar{TestimonyArguments.SolaFide}{36} \land \pvar{TestimonyArguments.SolaFide}{37} \land \pvar{TestimonyArguments.SolaFide}{40}) \rightarrow \pvar{TestimonyArguments.SolaFide}{25}\)\gloss{\emph{would give:} Justification is by faith alone}
  \breakrow{(15)} \((\pvar{TestimonyArguments.SolaFide}{20} \land \pvar{TestimonyArguments.SolaFide}{21}) \rightarrow \pvar{TestimonyArguments.SolaFide}{22}\)\gloss{\emph{would give:} James 2:24 is compatible with Paul, using ‘justify’ in a different sense}
  \breakrow{(16)} \((\pvar{TestimonyArguments.SolaFide}{1} \land \pvar{TestimonyArguments.SolaFide}{4} \land \pvar{TestimonyArguments.SolaFide}{5}) \rightarrow \pvar{TestimonyArguments.SolaFide}{27}\)\gloss{\emph{would give:} Salvation is by grace, a gift and not wages owed}
  \breakrow{(17)} \((\pvar{TestimonyArguments.SolaFide}{1} \land \pvar{TestimonyArguments.SolaFide}{4} \land \pvar{TestimonyArguments.SolaFide}{5} \land \pvar{TestimonyArguments.SolaFide}{22} \land \pvar{TestimonyArguments.SolaFide}{24}) \rightarrow \pvar{TestimonyArguments.SolaFide}{28}\)\gloss{\emph{would give:} Salvation is not by works: no work is its ground}
  \breakrow{(18)} \(\pvar{TestimonyArguments.SolaFide}{25} \rightarrow \pvar{TestimonyArguments.SolaFide}{26}\)\gloss{\emph{would give:} Faith is sufficient: whoever believes is saved}
  \breakrow{(19)} \((\pvar{TestimonyArguments.SolaFide}{26} \land \pvar{TestimonyArguments.SolaFide}{1}) \rightarrow \pvar{TestimonyArguments.SolaFide}{29}\)\gloss{\emph{would give:} Salvation is received through faith}
  \breakrow{(20)} \(\pvar{TestimonyArguments.SolaFide}{38}\)\gloss{John 3:36 — whoever believes has life; whoever does not obey shall not see life}
  \item[(21)] \(\pvar{TestimonyArguments.SolaFide}{62}\)\gloss{Matthew 7:21–23 — the criterion at the judgment is doing the Father’s will, not profession or mighty works}
  \item[(22)] \((\pvar{TestimonyArguments.SolaFide}{38} \land \pvar{TestimonyArguments.SolaFide}{62}) \rightarrow \lnot \pvar{TestimonyArguments.SolaFide}{26}\)\gloss{\emph{would deny:} Faith is sufficient: whoever believes is saved}
\end{derivation}
\par\medskip\noindent\textbf{\texttt{chrysostom\_\allowbreak{}cannot\_\allowbreak{}be\_\allowbreak{}held\_\allowbreak{}with\_\allowbreak{}johns\_\allowbreak{}case}}\par\nobreak\noindent
\textbf{Chrysostom’s conclusion and John’s strand cannot both be held.} John’s strand concludes that justification is by faith alone, and so that faith suffices. Chrysostom concludes that it does not. No consistent reading holds every premise of both.

A reader has to choose between Calvin’s reading of John and Chrysostom’s conclusion as written out here. It does not claim which to choose. Each rests on a step rated \texttt{disputed}.

Nor does it claim that the faith Chrysostom speaks of, right belief without a right life, is the faith the Reformers say justifies. His conclusion is written out here as “faith does not suffice”. A reader who takes him to deny only what James calls dead faith (see the section on James) need not choose.

{\small\begin{verbatim}
theorem chrysostom_cannot_be_held_with_johns_case : ¬Satisfiable
    johnWithChrysostom
-- axioms: propext, Classical.choice, Quot.sound
\end{verbatim}}
\par\medskip\noindent\textbf{\texttt{obedience\allowbreak{}As\allowbreak{}Fruit\allowbreak{}Source}}\par\nobreak\noindent
The step of the Reformed answer: from John 3:36 and Matthew 7:21, with the obedience they require read as sanctification, to “faith suffices”. Schreiner: “the necessity of good works doesn’t contradict sola fide” (p. 55). Calvin on 3:36: “He shows that enjoyment consists in faith.” Rated \texttt{disputed}: Chrysostom accepts both texts and denies that believing suffices, and Stanley holds works a condition of final salvation.

{\small\begin{verbatim}
def obedienceAsFruitSource : Source :=
  {
    primary :=
      Reference.work schreinerJustificationByWorks (Bib.Locus.pages 55 55),
    supporting :=
      [Reference.work calvinJohn (Bib.Locus.adLoc Scripture.john3_36),
        Reference.work westminsterConfession (Bib.Locus.sectionRef "XVI.2"),
        Reference.work westminsterConfession (Bib.Locus.sectionRef "XIII.1")],
    tradition := Tradition.reformedProtestant,
    confidence := Confidence.disputed }
\end{verbatim}}
\par\medskip\noindent\textbf{\texttt{sanctification\allowbreak{}Line}} --- Obedience as sanctification (the Reformed answer on John 3:36)\par\nobreak\noindent
\textbf{The Reformed answer to Chrysostom, as a line.} John 3:36 and Matthew 7:21 are accepted as they stand. The obedience they require is read as sanctification, the fruit of faith. So they add no condition to faith, and the first half of John 3:36 holds as it is written: whoever believes has eternal life.

\begin{derivation}
  \item[(1)] \(\pvar{TestimonyArguments.SolaFide}{38}\)\gloss{John 3:36 — whoever believes has life; whoever does not obey shall not see life}
  \item[(2)] \(\pvar{TestimonyArguments.SolaFide}{62}\)\gloss{Matthew 7:21–23 — the criterion at the judgment is doing the Father’s will, not profession or mighty works}
  \item[(3)] \(\pvar{TestimonyArguments.SolaFide}{39}\)\gloss{The obedience of John 3:36 and Matthew 7:21 is sanctification, the fruit of faith}
  \item[(4)] \((\pvar{TestimonyArguments.SolaFide}{38} \land \pvar{TestimonyArguments.SolaFide}{62} \land \pvar{TestimonyArguments.SolaFide}{39}) \rightarrow \pvar{TestimonyArguments.SolaFide}{26}\)
  \item[\(\vdash\)] \(\pvar{TestimonyArguments.SolaFide}{26}\)\gloss{\emph{at issue:} Faith is sufficient: whoever believes is saved}
\end{derivation}
\subsection*{Where Trent parts from Paul}
\addcontentsline{toc}{subsection}{Where Trent parts from Paul}
When teachers in Galatia added circumcision to the gospel, Paul answered with a whole letter. A different gospel is no gospel, he wrote, and whoever preaches it is accursed (1:6–9). Then, at the heart of the letter:

\begin{quote}
“If righteousness were through the law, then Christ died for no purpose” (Galatians 2:21).

\end{quote}
Whoever would be justified by the law is “severed from Christ” (5:2–4). The gospel itself is the one Paul received and handed on, that “Christ died for our sins” (1 Corinthians 15:3).

Fifteen centuries later the Council of Trent defined justification this way:

\begin{quote}
“not remission of sins merely, but also the sanctification and renewal of the inward man” (Session VI, ch. 7).

\end{quote}
Where exactly do Trent and Paul part? First, what Trent can accept; then the one step it denies; then what the checking found.

\subsubsection*{What Trent can accept}
Trent can accept every text above. It can also accept that Paul’s verb “justify” (δικαιόω) is a courtroom word, a verdict: scholars across the traditions read it so, the standard lexicon of New Testament Greek (BDAG), Morris, Moo, Irons, Wright, and the Catholic scholar Fitzmyer among them (\texttt{dikaio\allowbreak{}Is\allowbreak{}Forensic}, rated \texttt{well\allowbreak{}Supported}). The \emph{Joint Declaration on the Doctrine of Justification} (1999), signed by the Catholic Church, confesses that God “no longer imputes to them their sin” (§22). It does not speak to the sense of the verb.

\subsubsection*{What Trent denies}
Trent denies one step, which the Westminster Confession states (XI.1): that a verdict given on the ground of Christ’s finished work leaves no room for the renewal worked in us to be part of justification. A modern Catholic text says the same as Trent. The same paragraph of the \emph{Joint Declaration} holds that forgiveness and the love the Spirit brings about are “not to be separated”, and its Annex defines justification as “forgiveness of sins and being made righteous” (Annex 2 A). Gorman reads the verdict as a word that does what it says and changes the person.

\subsubsection*{Where they part}
The checking finds that the whole disagreement is that one step, and that what the step contradicts in Trent is its definition, and nothing else (\texttt{where\allowbreak{}Trent\allowbreak{}Parts\allowbreak{}From\allowbreak{}Paul}).

This does not show that the step is true. It shows that the step is the whole of the disagreement. Grant it, and Trent’s definition cannot stand. Deny it, and Paul’s gospel, as written out here, does not reach Trent’s definition at all.

\subsubsection*{How this is checked}
\paragraph*{The gospel line holds together}
\leavevmode\par\nobreak
\par\medskip\noindent\textbf{\texttt{galatian\allowbreak{}Gospel\allowbreak{}Reading}}\par\nobreak\noindent
Paul’s gospel as this line reads it: every text true, Christ’s work the whole ground, the verdict a courtroom one, and justification not the renewal of the inward person.

{\small\begin{verbatim}
def galatianGospelReading : Valuation Claim :=
  fun a =>
    match a with
    | Claim.justificationIncludesSanctification => False
    | x => True
\end{verbatim}}
\par\medskip\noindent\textbf{\texttt{galatian\allowbreak{}Gospel\_\allowbreak{}establishes}}\par\nobreak\noindent
\textbf{Paul’s gospel, read this way, denies Trent’s definition.} From Galatians 1:6–9, 2:21, 5:2–4 and 1 Corinthians 15:3, read as one claim, with “justify” (δικαιόω) as a courtroom word: justification is not the renewal of the inward person.

It does not show that the reading is beyond dispute. Two of its steps are rated \texttt{disputed}, and Trent accepts every text.

{\small\begin{verbatim}
theorem galatianGospel_establishes : Establishes galatianGospel
-- axioms: propext, Quot.sound
\end{verbatim}}
\par\medskip\noindent\textbf{\texttt{galatian\allowbreak{}Gospel\_\allowbreak{}is\_\allowbreak{}satisfiable}}\par\nobreak\noindent
The gospel line’s premises can all be true together.

{\small\begin{verbatim}
theorem galatianGospel_is_satisfiable : Satisfiable galatianGospel.premises
-- axioms: propext, Quot.sound
\end{verbatim}}
\paragraph*{What Trent can accept}
\leavevmode\par\nobreak
\par\medskip\noindent\textbf{\texttt{joint\allowbreak{}Declaration\allowbreak{}Reading}}\par\nobreak\noindent
Trent’s position with two things added: what the \emph{Joint Declaration} confesses, that God no longer counts sin against the justified (§22); and that Paul’s verb is a courtroom word, as the scholars cited on it read it. Every text of the gospel line holds. Justification still includes renewal, faith without charity does not suffice, and justification is not by faith alone.

Having granted that the verb is a courtroom word, this reading does not also hold that Paul’s verb means renewal: the two exclude each other (\texttt{Testimony\allowbreak{}Arguments.Meanings.Sola\allowbreak{}Fide.joins}). Of the other pairs of claims that exclude each other, it holds the Greek text of Revelation 22:11 and Sirach 18:22, and that James uses “justify” as Paul does.

{\small\begin{verbatim}
def jointDeclarationReading : Valuation Claim :=
  fun a =>
    match a with
    | Claim.paulsJustifyDenotesRenewal => False
    | Claim.rev22_11BeJustifiedStill => False
    | Claim.sir18_22BeJustifiedToDeath => False
    | Claim.james2_24Compatible => False
    | Claim.salvationNotByWorks => False
    | Claim.justificationIsForensicOnly => False
    | Claim.justificationDistinctFromSanctification => False
    | Claim.faithIsSufficient => False
    | Claim.justificationByFaithAlone => False
    | x => True
\end{verbatim}}
\par\medskip\noindent\textbf{\texttt{trent\_\allowbreak{}grants\_\allowbreak{}the\_\allowbreak{}gospel\_\allowbreak{}texts}}\par\nobreak\noindent
\textbf{Trent can accept every text of the gospel line.} Nothing Trent holds denies that Christ died for our sins, that righteousness through the law would make his death needless, or that another gospel is accursed.

{\small\begin{verbatim}
theorem trent_grants_the_gospel_texts : Grants tridentineCase (⋀[p
    Claim.galatians1_6_9, p Claim.firstCorinthians15_3, p Claim.galatians2_21,
    p Claim.galatians5_2_4, p Claim.romans8_33_34])
-- axioms: propext, Quot.sound
\end{verbatim}}
\par\medskip\noindent\textbf{\texttt{trent\_\allowbreak{}grants\_\allowbreak{}the\_\allowbreak{}forensic\_\allowbreak{}sense}}\par\nobreak\noindent
\textbf{Trent’s case can be held together with the verdict sense of Paul’s verb.} Nothing in Trent’s case denies that “justify” is a courtroom word. The meaning of the word is not where they part.

It does not show that Trent or the \emph{Joint Declaration} teaches the verdict sense. The Declaration confesses that God no longer counts sin (§22), and does not speak to the sense of the verb. This result says only that the two can be held together.

{\small\begin{verbatim}
theorem trent_grants_the_forensic_sense : Grants tridentineCase (p
    Claim.dikaioIsForensic)
-- axioms: propext, Quot.sound
\end{verbatim}}
\paragraph*{Where Trent parts from Paul}
\leavevmode\par\nobreak
\par\medskip\noindent\textbf{\texttt{where\allowbreak{}Trent\allowbreak{}Parts\allowbreak{}From\allowbreak{}Paul}}\par\nobreak\noindent
\textbf{Where Trent parts from Paul.} At one step: if justification is God’s verdict on the ground of Christ’s finished work, then the renewal of the inward person is no part of it (\texttt{verdict\allowbreak{}Excludes\allowbreak{}Renewal}; Westminster XI.1).

What that step contradicts in Trent is its definition, justification “not remission of sins merely, but also the sanctification and renewal of the inward man” (Session VI, ch. 7), and nothing else: not Trent’s account of merit, and not its reading of works.

The contradiction needs the gospel line’s starting points as well: the four texts, the reading of them, Romans 8:33–34 and the courtroom sense of the verb. Trent can accept every one of those (\texttt{trent\_\allowbreak{}grants\_\allowbreak{}the\_\allowbreak{}gospel\_\allowbreak{}texts}, \texttt{trent\_\allowbreak{}grants\_\allowbreak{}the\_\allowbreak{}forensic\_\allowbreak{}sense}). It is the step that Trent denies, with the \emph{Joint Declaration}: forgiveness and renewal are “not to be separated”.

\paragraph{Why \emph{Paul’s gospel (Galatians 1), against justification as renewal} stands against \emph{Trent, against ‘not by works’}.}
\begin{itemize}
\item \textbf{The crux:} \((\pvar{TestimonyArguments.SolaFide}{58} \land \pvar{TestimonyArguments.SolaFide}{49} \land \pvar{TestimonyArguments.SolaFide}{48}) \rightarrow \lnot \pvar{TestimonyArguments.SolaFide}{33}\), a premise of \emph{Paul’s gospel (Galatians 1), against justification as renewal}.
\item \textbf{What it does:} its conclusion does not follow without it.
\item \textbf{Granted:} \(\pvar{TestimonyArguments.SolaFide}{44}\), \(\pvar{TestimonyArguments.SolaFide}{45}\), \(\pvar{TestimonyArguments.SolaFide}{46}\), \(\pvar{TestimonyArguments.SolaFide}{47}\), \((\pvar{TestimonyArguments.SolaFide}{44} \land \pvar{TestimonyArguments.SolaFide}{45} \land \pvar{TestimonyArguments.SolaFide}{46} \land \pvar{TestimonyArguments.SolaFide}{47}) \rightarrow \pvar{TestimonyArguments.SolaFide}{58}\), \(\pvar{TestimonyArguments.SolaFide}{48}\), \(\pvar{TestimonyArguments.SolaFide}{49}\), which the break also needs.
\item \textbf{Where the rival breaks:} \(\pvar{TestimonyArguments.SolaFide}{33}\) cannot be held together with \((\pvar{TestimonyArguments.SolaFide}{58} \land \pvar{TestimonyArguments.SolaFide}{49} \land \pvar{TestimonyArguments.SolaFide}{48}) \rightarrow \lnot \pvar{TestimonyArguments.SolaFide}{33}\) and \(\pvar{TestimonyArguments.SolaFide}{44}\), \(\pvar{TestimonyArguments.SolaFide}{45}\), \(\pvar{TestimonyArguments.SolaFide}{46}\), \(\pvar{TestimonyArguments.SolaFide}{47}\), \((\pvar{TestimonyArguments.SolaFide}{44} \land \pvar{TestimonyArguments.SolaFide}{45} \land \pvar{TestimonyArguments.SolaFide}{46} \land \pvar{TestimonyArguments.SolaFide}{47}) \rightarrow \pvar{TestimonyArguments.SolaFide}{58}\), \(\pvar{TestimonyArguments.SolaFide}{48}\), \(\pvar{TestimonyArguments.SolaFide}{49}\); each is needed for the break, and without the crux they stand.
\item \textbf{What this rests on:}
\begin{itemize}
\item \(\pvar{TestimonyArguments.SolaFide}{44}\) Galatians 1:6–9: whoever preaches a gospel contrary to the one received is accursed --- \emph{consensus}: Gal 1:6-9; \cite[ad loc. Gal 1:6]{moo-galatians-2013}
\item \(\pvar{TestimonyArguments.SolaFide}{45}\) 1 Corinthians 15:3: the gospel of first importance — Christ died for our sins --- \emph{consensus}: 1 Cor 15:3
\item \(\pvar{TestimonyArguments.SolaFide}{46}\) Galatians 2:21: if righteousness were through the law, Christ died for no purpose --- \emph{consensus}: Gal 2:21; \cite[ad loc. Gal 2:21]{moo-galatians-2013}
\item \(\pvar{TestimonyArguments.SolaFide}{47}\) Galatians 5:2–4: you who would be justified by the law are severed from Christ --- \emph{consensus}: Gal 5:2-4; \cite[ad loc. Gal 5:2]{moo-galatians-2013}
\item \(\pvar{TestimonyArguments.SolaFide}{58}\) Christ’s death for our sins is the whole ground of justification; to add to it is another gospel --- \emph{disputed}: \cite[XI.1]{westminster-confession-1647}; Gal 2:21; \cite[ad loc. Gal 2:21]{moo-galatians-2013}; \cite[Epistle to Diognetus 9]{roberts-ante-nicene-fathers-1-1885}
\item \(\pvar{TestimonyArguments.SolaFide}{48}\) Romans 8:33–34: it is God who justifies — who is to condemn? --- \emph{consensus}: Rom 8:33-34
\item \(\pvar{TestimonyArguments.SolaFide}{49}\) In Paul, ‘justify’ (δικαιόω) is a verdict: to declare righteous --- \emph{well supported}: \cite[s.v. δικαιόω]{bauer-danker-lexicon-2000}; \cite{morris-apostolic-preaching-1965}; \cite[p. 80]{moo-romans-1996}; \cite{irons-righteousness-of-god-2015}; \cite{wright-justification-2009}; \cite[ad loc. Rom 3:24]{fitzmyer-romans-1993}
\item the step itself --- \emph{disputed}: \cite[ad loc. Gal 1:6]{moo-galatians-2013}; \cite[XI.1]{westminster-confession-1647}; \emph{disputed}: \cite[XI.1]{westminster-confession-1647}; \cite[III.xi.2]{calvin-institutes-1960}; \cite{morris-apostolic-preaching-1965}
\end{itemize}
\end{itemize}
\subsection*{What Trent means by justification}
\addcontentsline{toc}{subsection}{What Trent means by justification}
The last chapter found that Trent and Paul part at one step (\texttt{where\allowbreak{}Trent\allowbreak{}Parts\allowbreak{}From\allowbreak{}Paul}). What lies on Trent’s side of that step is its definition:

\begin{quote}
justification is “not remission of sins merely, but also the sanctification and renewal of the inward man” (Session VI, ch. 7).

\end{quote}
The Westminster Confession denies it (XI.1). But the definition can be taken in two ways, as a claim about a word or as a claim about what God does, and each way meets a different objection.

\subsubsection*{Taken as a claim about Paul’s word}
On this reading the definition says that “justify” (δικαιόω), in Paul, means the renewal of the inward person. Trent invites it. Chapter 8 of the decree sets out how “the Apostle’s” words, justified “by faith and freely”, are to be understood, and chapter 7 of the decree says that the justified “are not only reputed, but are truly called, and are, just”. So read, the definition is a claim about a word, and it meets an answer about the word. The verb is a courtroom word. Paul has other words for renewal and sets them beside it (Titus 3:5–7; 1 Corinthians 6:11). And a word contributes the least meaning its context requires (Joos; Silva). To read renewal into the verdict from the doctrine is what Barr called “illegitimate totality transfer”. Romans 4 gives the same answer from one passage: Paul explains the righteousness God counts to the ungodly as sin not counted, and by the same rule the verb he explains in that way does not also mean renewal. Every step of that answer is rated above \texttt{disputed}.

\subsubsection*{Why Trent would take it so}
The council read Paul in the Latin it had declared “authentic” (Session IV), and in Latin “justify” (\emph{iustificare}) had meant “to make righteous” since Augustine: “what else does the phrase ‘being justified’ signify than being made righteous” (\emph{On the Spirit and the Letter} 26.45). McGrath traces that inheritance through the medieval West. So this way of taking the definition has a starting point of its own, that Augustine explained the word so, which is rated \texttt{consensus} because he did. It then needs a step from the Latin explanation to the Greek word, and that step is where it fails. An explanation in one language does not settle what a word means in another, and the argument from usage says what the Greek word means in Paul (\texttt{where\allowbreak{}The\allowbreak{}Latin\allowbreak{}Reading\allowbreak{}Falls}).

\subsubsection*{Taken as a claim about what God does}
On this reading the definition says that when God justifies he also renews: pardon and renewal given together, “not to be separated”, as the \emph{Joint Declaration} (§22) confesses. No argument about the word reaches it. It can be held with everything the argument from usage says about “justify”, and Calvin himself grants that justification and sanctification cannot be separated. What still divides Trent from Paul is the step found in the last chapter: that a verdict on Christ’s finished work excludes renewal from justification (\texttt{where\allowbreak{}Trent\allowbreak{}Parts\allowbreak{}From\allowbreak{}Paul}). That step is \texttt{disputed}.

Romans 4 leads to the same place by another way, and does not raise the rating. God “justifies the ungodly” (4:5), and the blessing of righteousness counted apart from works is sin not counted (4:6–8). Read with the courtroom sense of the verb, what is counted to the ungodly is not a righteousness worked in them. Trent can accept every word (\texttt{trent\_\allowbreak{}grants\_\allowbreak{}romans\_\allowbreak{}four}), and Augustine reads the same verse the other way: God justifies the ungodly “that he may become a godly one” (\emph{On the Spirit and the Letter} 26.45). So the step is \texttt{disputed}, as the step from Galatians is, and Romans 4 adds a second argument at the same rating. So read, Trent is met by Paul from Galatians and from Romans 4, and each has to be answered (\texttt{where\allowbreak{}Romans\allowbreak{}Four\allowbreak{}Meets\allowbreak{}What\allowbreak{}God\allowbreak{}Does}). The two arguments are not wholly separate: both start from the courtroom sense of the verb, so a reader who denies that sense denies both.

The two are put together and checked in one result (\texttt{what\allowbreak{}Trents\allowbreak{}Definition\allowbreak{}Claims}). Each way of taking the definition is fair, since Trent read either way can still be held without contradiction, and each is answered. The first cannot be held with what Paul’s word means: it fails three times, each at a step rated \texttt{well\allowbreak{}Supported}. The second is untouched by anything about the word, and conflicts with Paul’s gospel and with Romans 4, each time at a step rated \texttt{disputed}. So a reader who grants that step rejects the definition taken that way, and a reader who denies it need not. Which way Trent means it is not something this chapter decides.

One premise, that Paul has his own words for renewal (\texttt{paul\allowbreak{}Names\allowbreak{}Renewal\allowbreak{}Otherwise}), is cited to Scripture alone: it says what the texts say and no more.

\subsubsection*{How this is checked}
\paragraph*{The argument from Paul’s word holds together}
\leavevmode\par\nobreak
\par\medskip\noindent\textbf{\texttt{lexical\allowbreak{}Reading}}\par\nobreak\noindent
The reading on which the argument from Paul’s word holds: every word as that argument takes it, and “justify” meaning the verdict and not renewal.

{\small\begin{verbatim}
def lexicalReading : Valuation Claim :=
  fun a =>
    match a with
    | Claim.paulsJustifyDenotesRenewal => False
    | x => True
\end{verbatim}}
\par\medskip\noindent\textbf{\texttt{lexical\allowbreak{}Case\_\allowbreak{}establishes}}\par\nobreak\noindent
\textbf{Paul’s “justify” (δικαιόω) does not mean the renewal of the inward person.} This follows from three things: the verb is a courtroom word; Paul has other words for renewal; and a word contributes the least meaning its context requires.

It does not claim anything about what God does when he justifies. It is a result about a word.

{\small\begin{verbatim}
theorem lexicalCase_establishes : Establishes lexicalCase
-- axioms: propext, Quot.sound
\end{verbatim}}
\par\medskip\noindent\textbf{\texttt{lexical\allowbreak{}Case\_\allowbreak{}is\_\allowbreak{}satisfiable}}\par\nobreak\noindent
The premises of the argument from Paul’s word can all be true together.

{\small\begin{verbatim}
theorem lexicalCase_is_satisfiable : Satisfiable lexicalCase.premises
-- axioms: propext, Quot.sound
\end{verbatim}}
\paragraph*{Two ways of taking Trent’s definition}
\leavevmode\par\nobreak
\par\medskip\noindent\textbf{\texttt{trent\allowbreak{}Reads\allowbreak{}Pauls\allowbreak{}Word}}\par\nobreak\noindent
Taking Trent’s definition as a claim about Paul’s word, and taking Paul’s word as the Latin West did. Trent’s chapter 8 sets out how “the Apostle’s” words are to be understood, and chapter 7 says the justified are “not only reputed, but are truly called, and are, just”. The council read Paul in the Latin it had declared “authentic” (Session IV), where “justify” (\emph{iustificare}) had meant “to make righteous” since Augustine. McGrath traces that inheritance. Rated \texttt{disputed}: whether the council defines the word or the reality is itself argued, and the step from the Latin explanation to the Greek word is what the argument from Paul’s word denies.

{\small\begin{verbatim}
def trentReadsPaulsWord : Source :=
  {
    primary :=
      Reference.work tannerDecrees
        (Bib.Locus.sectionRef
          "Trent, Session VI (1547), Decree on Justification, ch. 8"),
    supporting :=
      [Reference.work tannerDecrees
          (Bib.Locus.sectionRef
            "Trent, Session VI (1547), Decree on Justification, ch. 7"),
        Reference.work npnfAugustineAntiPelagian
          (Bib.Locus.sectionRef "On the Spirit and the Letter, ch. 45 (26.45)"),
        Reference.work mcgrathIustitiaDei],
    tradition := Tradition.romanCatholic, confidence := Confidence.disputed }
\end{verbatim}}
\par\medskip\noindent\textbf{\texttt{trent\allowbreak{}Reads\allowbreak{}Grace}}\par\nobreak\noindent
Taking Trent’s definition as a claim about what God does in justifying. In the \emph{Joint Declaration}, forgiveness and the renewal the Spirit brings about are two sides of God’s gracious action, “not to be separated” (§22). Rated \texttt{disputed}, for the same reason as the other way of taking it.

{\small\begin{verbatim}
def trentReadsGrace : Source :=
  {
    primary :=
      Reference.work jointDeclarationJustification (Bib.Locus.sectionRef "§22"),
    supporting :=
      [Reference.work tannerDecrees
          (Bib.Locus.sectionRef
            "Trent, Session VI (1547), Decree on Justification, ch. 7")],
    tradition := Tradition.romanCatholic, confidence := Confidence.disputed }
\end{verbatim}}
\par\medskip\noindent\textbf{\texttt{trent\allowbreak{}On\allowbreak{}Pauls\allowbreak{}Word}}\par\nobreak\noindent
Trent’s definition, \textbf{taken as a claim about what Paul’s word means}, and resting, as the Latin West did, on Augustine’s explanation: “being justified” is “being made righteous”.

{\small\begin{verbatim}
def trentOnPaulsWord : Reading Claim :=
  { name := "as what Paul's word means",
    commits := p Claim.paulsJustifyDenotesRenewal,
    grounds :=
      [p Claim.augustineReadsJustifyAsMakeRighteous,
        p Claim.augustineReadsJustifyAsMakeRighteous 🡒
          p Claim.paulsJustifyDenotesRenewal],
    source := trentReadsPaulsWord }
\end{verbatim}}
\par\medskip\noindent\textbf{\texttt{trent\allowbreak{}On\allowbreak{}What\allowbreak{}God\allowbreak{}Does}}\par\nobreak\noindent
Trent’s definition, \textbf{taken as a claim about what God does in justifying}.

{\small\begin{verbatim}
def trentOnWhatGodDoes : Reading Claim :=
  { name := "as what God does in justifying",
    commits := p Claim.justifyingGraceRenews,
    source := trentReadsGrace }
\end{verbatim}}
\par\medskip\noindent\textbf{\texttt{tridentine\allowbreak{}On\allowbreak{}Pauls\allowbreak{}Word}} --- Trent, against ‘not by works’, read as what Paul’s word means\par\nobreak\noindent
Trent, with its definition taken as a claim about Paul’s word.

\begin{derivation}
  \item[(1)] \(\pvar{TestimonyArguments.SolaFide}{1}\)\gloss{Ephesians 2:8–9 teaches salvation by grace through faith, not of works}
  \item[(2)] \(\pvar{TestimonyArguments.SolaFide}{2}\)\gloss{Romans 3:28 teaches justification by faith apart from works of the law}
  \item[(3)] \(\pvar{TestimonyArguments.SolaFide}{3}\)\gloss{Galatians 2:16 teaches that no one is justified by works of the law}
  \breakrow{(4)} \(\pvar{TestimonyArguments.SolaFide}{5}\)\gloss{Titus 3:5 teaches that God saved us not by works done in righteousness}
  \breakrow{(5)} \(\pvar{TestimonyArguments.SolaFide}{24}\)\gloss{Scripture does not contradict itself}
  \breakrow{(6)} \(\pvar{TestimonyArguments.SolaFide}{33}\)\gloss{Justification is not remission of sins only, but renewal of the inward man}
  \breakrow{(7)} \(\pvar{TestimonyArguments.SolaFide}{35}\)\gloss{The inward renewal of the justified grows as they do good works in grace}
  \breakrow{(8)} \((\pvar{TestimonyArguments.SolaFide}{33} \land \pvar{TestimonyArguments.SolaFide}{35}) \rightarrow \pvar{TestimonyArguments.SolaFide}{32}\)\gloss{\emph{would give:} Works performed in grace merit an increase of justification}
  \breakrow{(9)} \(\pvar{TestimonyArguments.SolaFide}{33} \rightarrow \lnot \pvar{TestimonyArguments.SolaFide}{30}\)\gloss{\emph{would deny:} Justification is a verdict only: pardon and righteousness counted, not poured in}
  \breakrow{(10)} \(\pvar{TestimonyArguments.SolaFide}{33} \rightarrow \lnot \pvar{TestimonyArguments.SolaFide}{34}\)\gloss{\emph{would deny:} Justification and sanctification are inseparable but distinct}
  \breakrow{(11)} \(\pvar{TestimonyArguments.SolaFide}{33} \rightarrow \lnot \pvar{TestimonyArguments.SolaFide}{26}\)\gloss{\emph{would deny:} Faith is sufficient: whoever believes is saved}
  \breakrow{(12)} \(\pvar{TestimonyArguments.SolaFide}{32} \rightarrow \lnot \pvar{TestimonyArguments.SolaFide}{28}\)
  \breakrow{(13)} \(\pvar{TestimonyArguments.SolaFide}{54}\)\gloss{Augustine explains “being justified” as “being made righteous”, and the Latin West read iustificare so}
  \breakrow{(14)} \(\pvar{TestimonyArguments.SolaFide}{54} \rightarrow \pvar{TestimonyArguments.SolaFide}{52}\)\gloss{\emph{would give:} Paul’s word ‘justify’ means the renewal of the inward man, not only a verdict}
  \item[(15)] \(\pvar{TestimonyArguments.SolaFide}{33} \rightarrow \pvar{TestimonyArguments.SolaFide}{52}\)\gloss{\emph{would give:} Paul’s word ‘justify’ means the renewal of the inward man, not only a verdict}
  \item[\(\vdash\)] \(\lnot \pvar{TestimonyArguments.SolaFide}{28}\)\gloss{\emph{at issue, the denial of:} Salvation is not by works: no work is its ground}
  \keepwithnext
\end{derivation}
\emph{No premise here rests on scripture alone.}
\par\medskip\noindent\textbf{\texttt{tridentine\allowbreak{}On\allowbreak{}What\allowbreak{}God\allowbreak{}Does}} --- Trent, against ‘not by works’, read as what God does in justifying\par\nobreak\noindent
Trent, with its definition taken as a claim about what God does.

\begin{derivation}
  \item[(1)] \(\pvar{TestimonyArguments.SolaFide}{1}\)\gloss{Ephesians 2:8–9 teaches salvation by grace through faith, not of works}
  \item[(2)] \(\pvar{TestimonyArguments.SolaFide}{2}\)\gloss{Romans 3:28 teaches justification by faith apart from works of the law}
  \item[(3)] \(\pvar{TestimonyArguments.SolaFide}{3}\)\gloss{Galatians 2:16 teaches that no one is justified by works of the law}
  \breakrow{(4)} \(\pvar{TestimonyArguments.SolaFide}{5}\)\gloss{Titus 3:5 teaches that God saved us not by works done in righteousness}
  \breakrow{(5)} \(\pvar{TestimonyArguments.SolaFide}{24}\)\gloss{Scripture does not contradict itself}
  \breakrow{(6)} \(\pvar{TestimonyArguments.SolaFide}{33}\)\gloss{Justification is not remission of sins only, but renewal of the inward man}
  \breakrow{(7)} \(\pvar{TestimonyArguments.SolaFide}{35}\)\gloss{The inward renewal of the justified grows as they do good works in grace}
  \breakrow{(8)} \((\pvar{TestimonyArguments.SolaFide}{33} \land \pvar{TestimonyArguments.SolaFide}{35}) \rightarrow \pvar{TestimonyArguments.SolaFide}{32}\)\gloss{\emph{would give:} Works performed in grace merit an increase of justification}
  \breakrow{(9)} \(\pvar{TestimonyArguments.SolaFide}{33} \rightarrow \lnot \pvar{TestimonyArguments.SolaFide}{30}\)\gloss{\emph{would deny:} Justification is a verdict only: pardon and righteousness counted, not poured in}
  \breakrow{(10)} \(\pvar{TestimonyArguments.SolaFide}{33} \rightarrow \lnot \pvar{TestimonyArguments.SolaFide}{34}\)\gloss{\emph{would deny:} Justification and sanctification are inseparable but distinct}
  \breakrow{(11)} \(\pvar{TestimonyArguments.SolaFide}{33} \rightarrow \lnot \pvar{TestimonyArguments.SolaFide}{26}\)\gloss{\emph{would deny:} Faith is sufficient: whoever believes is saved}
  \breakrow{(12)} \(\pvar{TestimonyArguments.SolaFide}{32} \rightarrow \lnot \pvar{TestimonyArguments.SolaFide}{28}\)
  \item[(13)] \(\pvar{TestimonyArguments.SolaFide}{33} \rightarrow \pvar{TestimonyArguments.SolaFide}{53}\)\gloss{\emph{would give:} When God justifies, he also renews: pardon and renewal are given together, not to be separated}
  \item[\(\vdash\)] \(\lnot \pvar{TestimonyArguments.SolaFide}{28}\)\gloss{\emph{at issue, the denial of:} Salvation is not by works: no work is its ground}
  \keepwithnext
\end{derivation}
\emph{No premise here rests on scripture alone.}
\par\medskip\noindent\textbf{\texttt{tridentine\allowbreak{}On\allowbreak{}Pauls\allowbreak{}Word\_\allowbreak{}establishes}}\par\nobreak\noindent
Trent, taken as a claim about Paul’s word, still concludes what Trent concludes: taking the definition one way only adds premises.

{\small\begin{verbatim}
theorem tridentineOnPaulsWord_establishes : Establishes tridentineOnPaulsWord
-- axioms: propext, Quot.sound
\end{verbatim}}
\par\medskip\noindent\textbf{\texttt{trent\allowbreak{}On\allowbreak{}Pauls\allowbreak{}Word\allowbreak{}Reading}}\par\nobreak\noindent
Trent’s position, taking Paul’s word as it does: “justify” means making just, and not a verdict only; justification includes renewal; faith without charity does not suffice; and justification is not by faith alone.

{\small\begin{verbatim}
def trentOnPaulsWordReading : Valuation Claim :=
  fun a =>
    match a with
    | Claim.salvationNotByWorks => False
    | Claim.justificationIsForensicOnly => False
    | Claim.justificationDistinctFromSanctification => False
    | Claim.faithIsSufficient => False
    | Claim.justificationByFaithAlone => False
    | Claim.dikaioIsForensic => False
    | x => True
\end{verbatim}}
\par\medskip\noindent\textbf{\texttt{tridentine\allowbreak{}On\allowbreak{}Pauls\allowbreak{}Word\_\allowbreak{}is\_\allowbreak{}satisfiable}}\par\nobreak\noindent
\textbf{Trent, taken as a claim about Paul’s word, can be held without contradiction.} So this branch of the question is a fair one: the position does not refute itself.

{\small\begin{verbatim}
theorem tridentineOnPaulsWord_is_satisfiable : Satisfiable
    tridentineOnPaulsWord.premises
-- axioms: propext, Quot.sound
\end{verbatim}}
\par\medskip\noindent\textbf{\texttt{trent\allowbreak{}Grants\allowbreak{}The\allowbreak{}Word\allowbreak{}Reading}}\par\nobreak\noindent
Trent taken as a claim about what God does, with everything the argument from Paul’s word says accepted: the verb a courtroom word, Paul’s own words for renewal, the rule of least meaning, and “justify” not meaning renewal. God, in justifying, also renews.

{\small\begin{verbatim}
def trentGrantsTheWordReading : Valuation Claim :=
  fun a =>
    match a with
    | Claim.salvationNotByWorks => False
    | Claim.justificationIsForensicOnly => False
    | Claim.justificationDistinctFromSanctification => False
    | Claim.faithIsSufficient => False
    | Claim.justificationByFaithAlone => False
    | Claim.paulsJustifyDenotesRenewal => False
    | x => True
\end{verbatim}}
\par\medskip\noindent\textbf{\texttt{tridentine\allowbreak{}On\allowbreak{}What\allowbreak{}God\allowbreak{}Does\_\allowbreak{}is\_\allowbreak{}satisfiable}}\par\nobreak\noindent
\textbf{Trent, taken as a claim about what God does, can be held without contradiction.}

{\small\begin{verbatim}
theorem tridentineOnWhatGodDoes_is_satisfiable : Satisfiable
    tridentineOnWhatGodDoes.premises
-- axioms: propext, Quot.sound
\end{verbatim}}
\par\medskip\noindent\textbf{\texttt{lexical\_\allowbreak{}case\_\allowbreak{}does\_\allowbreak{}not\_\allowbreak{}reach\_\allowbreak{}what\_\allowbreak{}god\_\allowbreak{}does}}\par\nobreak\noindent
\textbf{The argument from Paul’s word does not reach Trent’s definition taken as a claim about what God does.} Taken that way, Trent’s position can be held together with every premise of the argument. The rule of least meaning governs what Paul’s word says, and this way of taking the definition says nothing about the word.

{\small\begin{verbatim}
theorem lexical_case_does_not_reach_what_god_does : Satisfiable
    (lexicalCase.premises ++ tridentineOnWhatGodDoes.premises)
-- axioms: propext, Quot.sound
\end{verbatim}}
\paragraph*{Where each fails}
\leavevmode\par\nobreak
\par\medskip\noindent\textbf{\texttt{where\allowbreak{}Trents\allowbreak{}Word\allowbreak{}Reading\allowbreak{}Falls}}\par\nobreak\noindent
\textbf{Why Paul’s word stands against Trent taken as a claim about it.} Everything turns on one step: with the courtroom sense, Paul’s own words for renewal and the rule of least meaning, “justify” does not mean renewal. What that contradicts is Trent’s definition together with this way of taking it, and nothing else in Trent. Every link is rated above \texttt{disputed}.

\paragraph{Why \emph{Paul’s word (δικαιόω, by the rule of least meaning)} stands against \emph{Trent, against ‘not by works’, read as what Paul’s word means}.}
\begin{itemize}
\item \textbf{The crux:} \((\pvar{TestimonyArguments.SolaFide}{49} \land \pvar{TestimonyArguments.SolaFide}{50} \land \pvar{TestimonyArguments.SolaFide}{51}) \rightarrow \lnot \pvar{TestimonyArguments.SolaFide}{52}\), a premise of \emph{Paul’s word (δικαιόω, by the rule of least meaning)}.
\item \textbf{What it does:} its conclusion does not follow without it.
\item \textbf{Granted:} \(\pvar{TestimonyArguments.SolaFide}{49}\), \(\pvar{TestimonyArguments.SolaFide}{50}\), \(\pvar{TestimonyArguments.SolaFide}{51}\), which the break also needs.
\item \textbf{Where the rival breaks:} \(\pvar{TestimonyArguments.SolaFide}{33}\), \(\pvar{TestimonyArguments.SolaFide}{33} \rightarrow \pvar{TestimonyArguments.SolaFide}{52}\) cannot be held together with \((\pvar{TestimonyArguments.SolaFide}{49} \land \pvar{TestimonyArguments.SolaFide}{50} \land \pvar{TestimonyArguments.SolaFide}{51}) \rightarrow \lnot \pvar{TestimonyArguments.SolaFide}{52}\) and \(\pvar{TestimonyArguments.SolaFide}{49}\), \(\pvar{TestimonyArguments.SolaFide}{50}\), \(\pvar{TestimonyArguments.SolaFide}{51}\); each is needed for the break, and without the crux they stand.
\item \textbf{What this rests on:}
\begin{itemize}
\item \(\pvar{TestimonyArguments.SolaFide}{49}\) In Paul, ‘justify’ (δικαιόω) is a verdict: to declare righteous --- \emph{well supported}: \cite[s.v. δικαιόω]{bauer-danker-lexicon-2000}; \cite{morris-apostolic-preaching-1965}; \cite[p. 80]{moo-romans-1996}; \cite{irons-righteousness-of-god-2015}; \cite{wright-justification-2009}; \cite[ad loc. Rom 3:24]{fitzmyer-romans-1993}
\item \(\pvar{TestimonyArguments.SolaFide}{50}\) Paul names renewal with other words (ἀνακαίνωσις, ἁγιασμός) and sets them beside ‘justify’ --- \emph{consensus}: Titus 3:5-7; Rom 12:2; Rom 6:19-22; 1 Cor 6:11
\item \(\pvar{TestimonyArguments.SolaFide}{51}\) A word contributes the least meaning its context requires; one occurrence does not carry the whole concept --- \emph{well supported}: \cite{silva-biblical-words-1994}; \cite{joos-semantic-axiom-1972}; \cite{barr-semantics-2004}
\item the step itself --- \emph{well supported}: \cite{silva-biblical-words-1994}; \cite[p. 80]{moo-romans-1996}; \cite{barr-semantics-2004}
\end{itemize}
\end{itemize}
\par\medskip\noindent\textbf{\texttt{where\allowbreak{}The\allowbreak{}Latin\allowbreak{}Reading\allowbreak{}Falls}}\par\nobreak\noindent
\textbf{Why Paul’s word stands against the Latin reading of it.} The same step, and a second contradiction in the same branch. Trent rests its reading of Paul’s word on Augustine’s explanation, and what fails is the step from that explanation to what the Greek word means. That Augustine explained it so is not in question; it is rated \texttt{consensus}. But a Latin explanation does not settle what the Greek word means in Paul, and the courtroom sense, Paul’s own words for renewal and the rule of least meaning say it does not mean renewal.

\paragraph{Why \emph{Paul’s word (δικαιόω, by the rule of least meaning)} stands against \emph{Trent, against ‘not by works’, read as what Paul’s word means}.}
\begin{itemize}
\item \textbf{The crux:} \((\pvar{TestimonyArguments.SolaFide}{49} \land \pvar{TestimonyArguments.SolaFide}{50} \land \pvar{TestimonyArguments.SolaFide}{51}) \rightarrow \lnot \pvar{TestimonyArguments.SolaFide}{52}\), a premise of \emph{Paul’s word (δικαιόω, by the rule of least meaning)}.
\item \textbf{What it does:} its conclusion does not follow without it.
\item \textbf{Granted:} \(\pvar{TestimonyArguments.SolaFide}{49}\), \(\pvar{TestimonyArguments.SolaFide}{50}\), \(\pvar{TestimonyArguments.SolaFide}{51}\), which the break also needs.
\item \textbf{Where the rival breaks:} \(\pvar{TestimonyArguments.SolaFide}{54}\), \(\pvar{TestimonyArguments.SolaFide}{54} \rightarrow \pvar{TestimonyArguments.SolaFide}{52}\) cannot be held together with \((\pvar{TestimonyArguments.SolaFide}{49} \land \pvar{TestimonyArguments.SolaFide}{50} \land \pvar{TestimonyArguments.SolaFide}{51}) \rightarrow \lnot \pvar{TestimonyArguments.SolaFide}{52}\) and \(\pvar{TestimonyArguments.SolaFide}{49}\), \(\pvar{TestimonyArguments.SolaFide}{50}\), \(\pvar{TestimonyArguments.SolaFide}{51}\); each is needed for the break, and without the crux they stand.
\item \textbf{What this rests on:}
\begin{itemize}
\item \(\pvar{TestimonyArguments.SolaFide}{49}\) In Paul, ‘justify’ (δικαιόω) is a verdict: to declare righteous --- \emph{well supported}: \cite[s.v. δικαιόω]{bauer-danker-lexicon-2000}; \cite{morris-apostolic-preaching-1965}; \cite[p. 80]{moo-romans-1996}; \cite{irons-righteousness-of-god-2015}; \cite{wright-justification-2009}; \cite[ad loc. Rom 3:24]{fitzmyer-romans-1993}
\item \(\pvar{TestimonyArguments.SolaFide}{50}\) Paul names renewal with other words (ἀνακαίνωσις, ἁγιασμός) and sets them beside ‘justify’ --- \emph{consensus}: Titus 3:5-7; Rom 12:2; Rom 6:19-22; 1 Cor 6:11
\item \(\pvar{TestimonyArguments.SolaFide}{51}\) A word contributes the least meaning its context requires; one occurrence does not carry the whole concept --- \emph{well supported}: \cite{silva-biblical-words-1994}; \cite{joos-semantic-axiom-1972}; \cite{barr-semantics-2004}
\item the step itself --- \emph{well supported}: \cite{silva-biblical-words-1994}; \cite[p. 80]{moo-romans-1996}; \cite{barr-semantics-2004}
\end{itemize}
\end{itemize}
\par\medskip\noindent\textbf{\texttt{where\allowbreak{}Trents\allowbreak{}Grace\allowbreak{}Reading\allowbreak{}Falls}}\par\nobreak\noindent
\textbf{Why Paul’s gospel stands against Trent taken as a claim about what God does.} The same step as \texttt{where\allowbreak{}Trent\allowbreak{}Parts\allowbreak{}From\allowbreak{}Paul}, and the same contradiction: from a verdict on Christ’s finished work to “renewal is no part of justification”, against Trent’s definition. Taken this way, that step is all that divides them, and it is rated \texttt{disputed}.

\paragraph{Why \emph{Paul’s gospel (Galatians 1), against justification as renewal} stands against \emph{Trent, against ‘not by works’, read as what God does in justifying}.}
\begin{itemize}
\item \textbf{The crux:} \((\pvar{TestimonyArguments.SolaFide}{58} \land \pvar{TestimonyArguments.SolaFide}{49} \land \pvar{TestimonyArguments.SolaFide}{48}) \rightarrow \lnot \pvar{TestimonyArguments.SolaFide}{33}\), a premise of \emph{Paul’s gospel (Galatians 1), against justification as renewal}.
\item \textbf{What it does:} its conclusion does not follow without it.
\item \textbf{Granted:} \(\pvar{TestimonyArguments.SolaFide}{44}\), \(\pvar{TestimonyArguments.SolaFide}{45}\), \(\pvar{TestimonyArguments.SolaFide}{46}\), \(\pvar{TestimonyArguments.SolaFide}{47}\), \((\pvar{TestimonyArguments.SolaFide}{44} \land \pvar{TestimonyArguments.SolaFide}{45} \land \pvar{TestimonyArguments.SolaFide}{46} \land \pvar{TestimonyArguments.SolaFide}{47}) \rightarrow \pvar{TestimonyArguments.SolaFide}{58}\), \(\pvar{TestimonyArguments.SolaFide}{48}\), \(\pvar{TestimonyArguments.SolaFide}{49}\), which the break also needs.
\item \textbf{Where the rival breaks:} \(\pvar{TestimonyArguments.SolaFide}{33}\) cannot be held together with \((\pvar{TestimonyArguments.SolaFide}{58} \land \pvar{TestimonyArguments.SolaFide}{49} \land \pvar{TestimonyArguments.SolaFide}{48}) \rightarrow \lnot \pvar{TestimonyArguments.SolaFide}{33}\) and \(\pvar{TestimonyArguments.SolaFide}{44}\), \(\pvar{TestimonyArguments.SolaFide}{45}\), \(\pvar{TestimonyArguments.SolaFide}{46}\), \(\pvar{TestimonyArguments.SolaFide}{47}\), \((\pvar{TestimonyArguments.SolaFide}{44} \land \pvar{TestimonyArguments.SolaFide}{45} \land \pvar{TestimonyArguments.SolaFide}{46} \land \pvar{TestimonyArguments.SolaFide}{47}) \rightarrow \pvar{TestimonyArguments.SolaFide}{58}\), \(\pvar{TestimonyArguments.SolaFide}{48}\), \(\pvar{TestimonyArguments.SolaFide}{49}\); each is needed for the break, and without the crux they stand.
\item \textbf{What this rests on:}
\begin{itemize}
\item \(\pvar{TestimonyArguments.SolaFide}{44}\) Galatians 1:6–9: whoever preaches a gospel contrary to the one received is accursed --- \emph{consensus}: Gal 1:6-9; \cite[ad loc. Gal 1:6]{moo-galatians-2013}
\item \(\pvar{TestimonyArguments.SolaFide}{45}\) 1 Corinthians 15:3: the gospel of first importance — Christ died for our sins --- \emph{consensus}: 1 Cor 15:3
\item \(\pvar{TestimonyArguments.SolaFide}{46}\) Galatians 2:21: if righteousness were through the law, Christ died for no purpose --- \emph{consensus}: Gal 2:21; \cite[ad loc. Gal 2:21]{moo-galatians-2013}
\item \(\pvar{TestimonyArguments.SolaFide}{47}\) Galatians 5:2–4: you who would be justified by the law are severed from Christ --- \emph{consensus}: Gal 5:2-4; \cite[ad loc. Gal 5:2]{moo-galatians-2013}
\item \(\pvar{TestimonyArguments.SolaFide}{58}\) Christ’s death for our sins is the whole ground of justification; to add to it is another gospel --- \emph{disputed}: \cite[XI.1]{westminster-confession-1647}; Gal 2:21; \cite[ad loc. Gal 2:21]{moo-galatians-2013}; \cite[Epistle to Diognetus 9]{roberts-ante-nicene-fathers-1-1885}
\item \(\pvar{TestimonyArguments.SolaFide}{48}\) Romans 8:33–34: it is God who justifies — who is to condemn? --- \emph{consensus}: Rom 8:33-34
\item \(\pvar{TestimonyArguments.SolaFide}{49}\) In Paul, ‘justify’ (δικαιόω) is a verdict: to declare righteous --- \emph{well supported}: \cite[s.v. δικαιόω]{bauer-danker-lexicon-2000}; \cite{morris-apostolic-preaching-1965}; \cite[p. 80]{moo-romans-1996}; \cite{irons-righteousness-of-god-2015}; \cite{wright-justification-2009}; \cite[ad loc. Rom 3:24]{fitzmyer-romans-1993}
\item the step itself --- \emph{disputed}: \cite[ad loc. Gal 1:6]{moo-galatians-2013}; \cite[XI.1]{westminster-confession-1647}; \emph{disputed}: \cite[XI.1]{westminster-confession-1647}; \cite[III.xi.2]{calvin-institutes-1960}; \cite{morris-apostolic-preaching-1965}
\end{itemize}
\end{itemize}
\paragraph*{Romans 4}
\leavevmode\par\nobreak
\par\medskip\noindent\textbf{\texttt{romans\allowbreak{}Four\allowbreak{}Case\_\allowbreak{}establishes}}\par\nobreak\noindent
\textbf{Romans 4, read this way, denies Trent’s definition.} Paul explains the righteousness God counts to the ungodly as sin not counted. With the courtroom sense of the verb, renewal is no part of justification.

It does not show that the reading is beyond dispute. Augustine accepts the same verses and reads them the other way, so the step is rated \texttt{disputed}.

{\small\begin{verbatim}
theorem romansFourCase_establishes : Establishes romansFourCase
-- axioms: propext, Quot.sound
\end{verbatim}}
\par\medskip\noindent\textbf{\texttt{romans\allowbreak{}Four\allowbreak{}Case\_\allowbreak{}is\_\allowbreak{}satisfiable}}\par\nobreak\noindent
The premises of the argument from Romans 4 can all be true together, on the reading given for Paul’s gospel.

{\small\begin{verbatim}
theorem romansFourCase_is_satisfiable : Satisfiable romansFourCase.premises
-- axioms: propext, Quot.sound
\end{verbatim}}
\par\medskip\noindent\textbf{\texttt{romans\allowbreak{}Four\allowbreak{}On\allowbreak{}The\allowbreak{}Word\_\allowbreak{}establishes}}\par\nobreak\noindent
\textbf{Romans 4, read by the rule of least meaning, shows what Paul’s word means.} The verb Paul explains as sin not counted does not also mean the renewal of the inward person.

{\small\begin{verbatim}
theorem romansFourOnTheWord_establishes : Establishes romansFourOnTheWord
-- axioms: propext, Quot.sound
\end{verbatim}}
\par\medskip\noindent\textbf{\texttt{romans\allowbreak{}Four\allowbreak{}On\allowbreak{}The\allowbreak{}Word\_\allowbreak{}is\_\allowbreak{}satisfiable}}\par\nobreak\noindent
The premises of that argument can all be true together, on the reading given for the argument from Paul’s word.

{\small\begin{verbatim}
theorem romansFourOnTheWord_is_satisfiable : Satisfiable
    romansFourOnTheWord.premises
-- axioms: propext, Quot.sound
\end{verbatim}}
\par\medskip\noindent\textbf{\texttt{trent\_\allowbreak{}grants\_\allowbreak{}romans\_\allowbreak{}four}}\par\nobreak\noindent
\textbf{Trent can accept every word of Romans 4:5–8}, as Augustine does: God justifies the ungodly, and the blessing is sin not counted. The texts are accepted on both sides. The reading of them is not.

{\small\begin{verbatim}
theorem trent_grants_romans_four : Grants tridentineCase (⋀[p
    Claim.romans4_4_5, p Claim.romans4_6_8])
-- axioms: propext, Quot.sound
\end{verbatim}}
\par\medskip\noindent\textbf{\texttt{romans\_\allowbreak{}four\_\allowbreak{}on\_\allowbreak{}the\_\allowbreak{}word\_\allowbreak{}does\_\allowbreak{}not\_\allowbreak{}reach\_\allowbreak{}what\_\allowbreak{}god\_\allowbreak{}does}}\par\nobreak\noindent
\textbf{Romans 4, read for Paul’s word, does not reach Trent’s definition taken as a claim about what God does.} Taken that way, Trent’s position can be held with every premise of the argument. The argument says what the counting is, and this way of taking the definition says nothing about the word.

{\small\begin{verbatim}
theorem romans_four_on_the_word_does_not_reach_what_god_does : Satisfiable
    (romansFourOnTheWord.premises ++ tridentineOnWhatGodDoes.premises)
-- axioms: propext, Quot.sound
\end{verbatim}}
\par\medskip\noindent\textbf{\texttt{where\allowbreak{}Romans\allowbreak{}Four\allowbreak{}Meets\allowbreak{}The\allowbreak{}Word\allowbreak{}Reading}}\par\nobreak\noindent
\textbf{Why Romans 4 stands against Trent taken as a claim about Paul’s word.} Everything turns on the step from Paul’s explanation, that the righteousness counted to the ungodly is sin not counted, by the rule of least meaning, to “the verb does not also mean renewal”. What it contradicts is Trent’s definition taken this way. Every link is rated \texttt{well\allowbreak{}Supported} or better. It arrives from one passage where the argument from usage arrives from the word.

\paragraph{Why \emph{Romans 4:3–8, on Paul’s word (by the rule of least meaning)} stands against \emph{Trent, against ‘not by works’, read as what Paul’s word means}.}
\begin{itemize}
\item \textbf{The crux:} \((\pvar{TestimonyArguments.SolaFide}{57} \land \pvar{TestimonyArguments.SolaFide}{51}) \rightarrow \lnot \pvar{TestimonyArguments.SolaFide}{52}\), a premise of \emph{Romans 4:3–8, on Paul’s word (by the rule of least meaning)}.
\item \textbf{What it does:} its conclusion does not follow without it.
\item \textbf{Granted:} \(\pvar{TestimonyArguments.SolaFide}{4}\), \(\pvar{TestimonyArguments.SolaFide}{55}\), \(\pvar{TestimonyArguments.SolaFide}{56}\), \(\pvar{TestimonyArguments.SolaFide}{51}\), \((\pvar{TestimonyArguments.SolaFide}{4} \land \pvar{TestimonyArguments.SolaFide}{55} \land \pvar{TestimonyArguments.SolaFide}{56}) \rightarrow \pvar{TestimonyArguments.SolaFide}{57}\), which the break also needs.
\item \textbf{Where the rival breaks:} \(\pvar{TestimonyArguments.SolaFide}{33}\), \(\pvar{TestimonyArguments.SolaFide}{33} \rightarrow \pvar{TestimonyArguments.SolaFide}{52}\) cannot be held together with \((\pvar{TestimonyArguments.SolaFide}{57} \land \pvar{TestimonyArguments.SolaFide}{51}) \rightarrow \lnot \pvar{TestimonyArguments.SolaFide}{52}\) and \(\pvar{TestimonyArguments.SolaFide}{4}\), \(\pvar{TestimonyArguments.SolaFide}{55}\), \(\pvar{TestimonyArguments.SolaFide}{56}\), \(\pvar{TestimonyArguments.SolaFide}{51}\), \((\pvar{TestimonyArguments.SolaFide}{4} \land \pvar{TestimonyArguments.SolaFide}{55} \land \pvar{TestimonyArguments.SolaFide}{56}) \rightarrow \pvar{TestimonyArguments.SolaFide}{57}\); each is needed for the break, and without the crux they stand.
\item \textbf{What this rests on:}
\begin{itemize}
\item \(\pvar{TestimonyArguments.SolaFide}{4}\) Romans 4:4–5 contrasts wages owed with a gift reckoned to the one who believes --- \emph{well supported}: Rom 4:4-5; \cite[III.xi.18]{calvin-institutes-1960}
\item \(\pvar{TestimonyArguments.SolaFide}{55}\) Romans 4:6–8: righteousness counted apart from works — sins forgiven, sin not counted --- \emph{consensus}: Rom 4:6-8; \cite[ad loc. Rom 4:6]{moo-romans-1996}
\item \(\pvar{TestimonyArguments.SolaFide}{56}\) ‘Count’ (λογίζομαι) in Romans 4:3–11 is the language of an account: credited to one --- \emph{well supported}: \cite[s.v. λογίζομαι]{bauer-danker-lexicon-2000}; \cite[ad loc. Rom 4:3]{moo-romans-1996}
\item \(\pvar{TestimonyArguments.SolaFide}{51}\) A word contributes the least meaning its context requires; one occurrence does not carry the whole concept --- \emph{well supported}: \cite{silva-biblical-words-1994}; \cite{joos-semantic-axiom-1972}; \cite{barr-semantics-2004}
\item \(\pvar{TestimonyArguments.SolaFide}{57}\) In Romans 4:6–8 Paul explains righteousness counted apart from works as sin not counted --- \emph{well supported}: \cite[ad loc. Rom 4:6]{moo-romans-1996}; Rom 4:6-8
\item the step itself --- \emph{well supported}: \cite[ad loc. Rom 4:6]{moo-romans-1996}; \cite[s.v. λογίζομαι]{bauer-danker-lexicon-2000}; \emph{well supported}: \cite{silva-biblical-words-1994}; \cite[p. 80]{moo-romans-1996}; \cite{barr-semantics-2004}
\end{itemize}
\end{itemize}
\par\medskip\noindent\textbf{\texttt{where\allowbreak{}Romans\allowbreak{}Four\allowbreak{}Meets\allowbreak{}What\allowbreak{}God\allowbreak{}Does}}\par\nobreak\noindent
\textbf{Why Romans 4 stands against Trent taken as a claim about what God does.} Everything turns on the step from Paul’s explanation, with the courtroom sense of the verb, to “renewal is no part of justification”. It contradicts Trent’s definition and nothing else.

The step is rated \texttt{disputed}, as the step from Galatians is. Augustine accepts the same verses and reads them the other way: God justifies the ungodly “that he may become a godly one”. So taken this way, Trent is met by Paul twice, from Galatians and from Romans 4, each time at a \texttt{disputed} step. The two arguments share one starting point, the courtroom sense of the verb.

\paragraph{Why \emph{Romans 4:3–8, against justification as renewal} stands against \emph{Trent, against ‘not by works’, read as what God does in justifying}.}
\begin{itemize}
\item \textbf{The crux:} \((\pvar{TestimonyArguments.SolaFide}{57} \land \pvar{TestimonyArguments.SolaFide}{49}) \rightarrow \lnot \pvar{TestimonyArguments.SolaFide}{33}\), a premise of \emph{Romans 4:3–8, against justification as renewal}.
\item \textbf{What it does:} its conclusion does not follow without it.
\item \textbf{Granted:} \(\pvar{TestimonyArguments.SolaFide}{4}\), \(\pvar{TestimonyArguments.SolaFide}{55}\), \(\pvar{TestimonyArguments.SolaFide}{56}\), \(\pvar{TestimonyArguments.SolaFide}{49}\), \((\pvar{TestimonyArguments.SolaFide}{4} \land \pvar{TestimonyArguments.SolaFide}{55} \land \pvar{TestimonyArguments.SolaFide}{56}) \rightarrow \pvar{TestimonyArguments.SolaFide}{57}\), which the break also needs.
\item \textbf{Where the rival breaks:} \(\pvar{TestimonyArguments.SolaFide}{33}\) cannot be held together with \((\pvar{TestimonyArguments.SolaFide}{57} \land \pvar{TestimonyArguments.SolaFide}{49}) \rightarrow \lnot \pvar{TestimonyArguments.SolaFide}{33}\) and \(\pvar{TestimonyArguments.SolaFide}{4}\), \(\pvar{TestimonyArguments.SolaFide}{55}\), \(\pvar{TestimonyArguments.SolaFide}{56}\), \(\pvar{TestimonyArguments.SolaFide}{49}\), \((\pvar{TestimonyArguments.SolaFide}{4} \land \pvar{TestimonyArguments.SolaFide}{55} \land \pvar{TestimonyArguments.SolaFide}{56}) \rightarrow \pvar{TestimonyArguments.SolaFide}{57}\); each is needed for the break, and without the crux they stand.
\item \textbf{What this rests on:}
\begin{itemize}
\item \(\pvar{TestimonyArguments.SolaFide}{4}\) Romans 4:4–5 contrasts wages owed with a gift reckoned to the one who believes --- \emph{well supported}: Rom 4:4-5; \cite[III.xi.18]{calvin-institutes-1960}
\item \(\pvar{TestimonyArguments.SolaFide}{55}\) Romans 4:6–8: righteousness counted apart from works — sins forgiven, sin not counted --- \emph{consensus}: Rom 4:6-8; \cite[ad loc. Rom 4:6]{moo-romans-1996}
\item \(\pvar{TestimonyArguments.SolaFide}{56}\) ‘Count’ (λογίζομαι) in Romans 4:3–11 is the language of an account: credited to one --- \emph{well supported}: \cite[s.v. λογίζομαι]{bauer-danker-lexicon-2000}; \cite[ad loc. Rom 4:3]{moo-romans-1996}
\item \(\pvar{TestimonyArguments.SolaFide}{49}\) In Paul, ‘justify’ (δικαιόω) is a verdict: to declare righteous --- \emph{well supported}: \cite[s.v. δικαιόω]{bauer-danker-lexicon-2000}; \cite{morris-apostolic-preaching-1965}; \cite[p. 80]{moo-romans-1996}; \cite{irons-righteousness-of-god-2015}; \cite{wright-justification-2009}; \cite[ad loc. Rom 3:24]{fitzmyer-romans-1993}
\item \(\pvar{TestimonyArguments.SolaFide}{57}\) In Romans 4:6–8 Paul explains righteousness counted apart from works as sin not counted --- \emph{well supported}: \cite[ad loc. Rom 4:6]{moo-romans-1996}; Rom 4:6-8
\item the step itself --- \emph{well supported}: \cite[ad loc. Rom 4:6]{moo-romans-1996}; \cite[s.v. λογίζομαι]{bauer-danker-lexicon-2000}; \emph{disputed}: \cite[ad loc. Rom 4:5]{moo-romans-1996}; \cite[XI.1]{westminster-confession-1647}
\end{itemize}
\end{itemize}
\paragraph*{The two put together}
\leavevmode\par\nobreak
\par\medskip\noindent\textbf{\texttt{what\allowbreak{}Trents\allowbreak{}Definition\allowbreak{}Claims}}\par\nobreak\noindent
\textbf{What Trent’s definition claims, taken both ways.} The definition is that justification is the sanctification and renewal of the inward person.

\emph{Taken as a claim about what Paul’s word means}, namely that “justify” (δικαιόω) means that renewal, it cannot be held with the evidence for what the word means. It fails three times, each at a step rated \texttt{well\allowbreak{}Supported}: against the argument from usage, as Trent states the definition; against the same argument, at the Latin explanation this reading rests on (“being justified” as “being made righteous”); and against Romans 4 read for the word.

\emph{Taken as a claim about what God does in justifying}, namely that he also renews, nothing about the word reaches it: neither the argument from usage nor Romans 4 read for the word. It conflicts with Paul twice, each time at a step rated \texttt{disputed}: with his gospel in Galatians, at the step \texttt{where\allowbreak{}Trent\allowbreak{}Parts\allowbreak{}From\allowbreak{}Paul} names, and with Romans 4, at the reading of “justifies the ungodly”. A reader who grants those steps rejects the definition taken this way. A reader who denies them need not. Augustine does, and so does the \emph{Joint Declaration}‘s Annex, which defines justification as “forgiveness of sins and being made righteous” (Annex 2 A).

Each way of taking it is fair: Trent’s position, taken either way, can still be held without contradiction.

\paragraph{Every reading of \(P_{33}\), answered.} \emph{Trent, against ‘not by works’} holds \(\pvar{TestimonyArguments.SolaFide}{33}\). It is read 2 ways here, and each reading is checked; none can be left out.

\paragraph{1. Read as what Paul’s word means.} The claim commits \emph{Trent, against ‘not by works’} to \(\pvar{TestimonyArguments.SolaFide}{52}\) --- so read by \cite[Trent, Session VI (1547), Decree on Justification, ch. 8]{tanner-decrees-1990}; \cite[Trent, Session VI (1547), Decree on Justification, ch. 7]{tanner-decrees-1990}; \cite[On the Spirit and the Letter, ch. 45 (26.45)]{npnf1-05-1887}; \cite{mcgrath-iustitia-dei-2005} (\emph{disputed}). It rests the reading on \(\pvar{TestimonyArguments.SolaFide}{54}\), \(\pvar{TestimonyArguments.SolaFide}{54} \rightarrow \pvar{TestimonyArguments.SolaFide}{52}\).

Against \emph{Paul’s word (δικαιόω, by the rule of least meaning)}, it cannot be held. Where it breaks:
\paragraph{Why \emph{Paul’s word (δικαιόω, by the rule of least meaning)} stands against \emph{Trent, against ‘not by works’, read as what Paul’s word means}.}
\begin{itemize}
\item \textbf{The crux:} \((\pvar{TestimonyArguments.SolaFide}{49} \land \pvar{TestimonyArguments.SolaFide}{50} \land \pvar{TestimonyArguments.SolaFide}{51}) \rightarrow \lnot \pvar{TestimonyArguments.SolaFide}{52}\), a premise of \emph{Paul’s word (δικαιόω, by the rule of least meaning)}.
\item \textbf{What it does:} its conclusion does not follow without it.
\item \textbf{Granted:} \(\pvar{TestimonyArguments.SolaFide}{49}\), \(\pvar{TestimonyArguments.SolaFide}{50}\), \(\pvar{TestimonyArguments.SolaFide}{51}\), which the break also needs.
\item \textbf{Where the rival breaks:} \(\pvar{TestimonyArguments.SolaFide}{33}\), \(\pvar{TestimonyArguments.SolaFide}{33} \rightarrow \pvar{TestimonyArguments.SolaFide}{52}\) cannot be held together with \((\pvar{TestimonyArguments.SolaFide}{49} \land \pvar{TestimonyArguments.SolaFide}{50} \land \pvar{TestimonyArguments.SolaFide}{51}) \rightarrow \lnot \pvar{TestimonyArguments.SolaFide}{52}\) and \(\pvar{TestimonyArguments.SolaFide}{49}\), \(\pvar{TestimonyArguments.SolaFide}{50}\), \(\pvar{TestimonyArguments.SolaFide}{51}\); each is needed for the break, and without the crux they stand.
\item \textbf{What this rests on:}
\begin{itemize}
\item \(\pvar{TestimonyArguments.SolaFide}{49}\) In Paul, ‘justify’ (δικαιόω) is a verdict: to declare righteous --- \emph{well supported}: \cite[s.v. δικαιόω]{bauer-danker-lexicon-2000}; \cite{morris-apostolic-preaching-1965}; \cite[p. 80]{moo-romans-1996}; \cite{irons-righteousness-of-god-2015}; \cite{wright-justification-2009}; \cite[ad loc. Rom 3:24]{fitzmyer-romans-1993}
\item \(\pvar{TestimonyArguments.SolaFide}{50}\) Paul names renewal with other words (ἀνακαίνωσις, ἁγιασμός) and sets them beside ‘justify’ --- \emph{consensus}: Titus 3:5-7; Rom 12:2; Rom 6:19-22; 1 Cor 6:11
\item \(\pvar{TestimonyArguments.SolaFide}{51}\) A word contributes the least meaning its context requires; one occurrence does not carry the whole concept --- \emph{well supported}: \cite{silva-biblical-words-1994}; \cite{joos-semantic-axiom-1972}; \cite{barr-semantics-2004}
\item the step itself --- \emph{well supported}: \cite{silva-biblical-words-1994}; \cite[p. 80]{moo-romans-1996}; \cite{barr-semantics-2004}
\end{itemize}
\end{itemize}
Against \emph{Paul’s word (δικαιόω, by the rule of least meaning)}, it cannot be held. Where it breaks:
\paragraph{Why \emph{Paul’s word (δικαιόω, by the rule of least meaning)} stands against \emph{Trent, against ‘not by works’, read as what Paul’s word means}.}
\begin{itemize}
\item \textbf{The crux:} \((\pvar{TestimonyArguments.SolaFide}{49} \land \pvar{TestimonyArguments.SolaFide}{50} \land \pvar{TestimonyArguments.SolaFide}{51}) \rightarrow \lnot \pvar{TestimonyArguments.SolaFide}{52}\), a premise of \emph{Paul’s word (δικαιόω, by the rule of least meaning)}.
\item \textbf{What it does:} its conclusion does not follow without it.
\item \textbf{Granted:} \(\pvar{TestimonyArguments.SolaFide}{49}\), \(\pvar{TestimonyArguments.SolaFide}{50}\), \(\pvar{TestimonyArguments.SolaFide}{51}\), which the break also needs.
\item \textbf{Where the rival breaks:} \(\pvar{TestimonyArguments.SolaFide}{54}\), \(\pvar{TestimonyArguments.SolaFide}{54} \rightarrow \pvar{TestimonyArguments.SolaFide}{52}\) cannot be held together with \((\pvar{TestimonyArguments.SolaFide}{49} \land \pvar{TestimonyArguments.SolaFide}{50} \land \pvar{TestimonyArguments.SolaFide}{51}) \rightarrow \lnot \pvar{TestimonyArguments.SolaFide}{52}\) and \(\pvar{TestimonyArguments.SolaFide}{49}\), \(\pvar{TestimonyArguments.SolaFide}{50}\), \(\pvar{TestimonyArguments.SolaFide}{51}\); each is needed for the break, and without the crux they stand.
\item \textbf{What this rests on:}
\begin{itemize}
\item \(\pvar{TestimonyArguments.SolaFide}{49}\) In Paul, ‘justify’ (δικαιόω) is a verdict: to declare righteous --- \emph{well supported}: \cite[s.v. δικαιόω]{bauer-danker-lexicon-2000}; \cite{morris-apostolic-preaching-1965}; \cite[p. 80]{moo-romans-1996}; \cite{irons-righteousness-of-god-2015}; \cite{wright-justification-2009}; \cite[ad loc. Rom 3:24]{fitzmyer-romans-1993}
\item \(\pvar{TestimonyArguments.SolaFide}{50}\) Paul names renewal with other words (ἀνακαίνωσις, ἁγιασμός) and sets them beside ‘justify’ --- \emph{consensus}: Titus 3:5-7; Rom 12:2; Rom 6:19-22; 1 Cor 6:11
\item \(\pvar{TestimonyArguments.SolaFide}{51}\) A word contributes the least meaning its context requires; one occurrence does not carry the whole concept --- \emph{well supported}: \cite{silva-biblical-words-1994}; \cite{joos-semantic-axiom-1972}; \cite{barr-semantics-2004}
\item the step itself --- \emph{well supported}: \cite{silva-biblical-words-1994}; \cite[p. 80]{moo-romans-1996}; \cite{barr-semantics-2004}
\end{itemize}
\end{itemize}
Against \emph{Romans 4:3–8, on Paul’s word (by the rule of least meaning)}, it cannot be held. Where it breaks:
\paragraph{Why \emph{Romans 4:3–8, on Paul’s word (by the rule of least meaning)} stands against \emph{Trent, against ‘not by works’, read as what Paul’s word means}.}
\begin{itemize}
\item \textbf{The crux:} \((\pvar{TestimonyArguments.SolaFide}{57} \land \pvar{TestimonyArguments.SolaFide}{51}) \rightarrow \lnot \pvar{TestimonyArguments.SolaFide}{52}\), a premise of \emph{Romans 4:3–8, on Paul’s word (by the rule of least meaning)}.
\item \textbf{What it does:} its conclusion does not follow without it.
\item \textbf{Granted:} \(\pvar{TestimonyArguments.SolaFide}{4}\), \(\pvar{TestimonyArguments.SolaFide}{55}\), \(\pvar{TestimonyArguments.SolaFide}{56}\), \(\pvar{TestimonyArguments.SolaFide}{51}\), \((\pvar{TestimonyArguments.SolaFide}{4} \land \pvar{TestimonyArguments.SolaFide}{55} \land \pvar{TestimonyArguments.SolaFide}{56}) \rightarrow \pvar{TestimonyArguments.SolaFide}{57}\), which the break also needs.
\item \textbf{Where the rival breaks:} \(\pvar{TestimonyArguments.SolaFide}{33}\), \(\pvar{TestimonyArguments.SolaFide}{33} \rightarrow \pvar{TestimonyArguments.SolaFide}{52}\) cannot be held together with \((\pvar{TestimonyArguments.SolaFide}{57} \land \pvar{TestimonyArguments.SolaFide}{51}) \rightarrow \lnot \pvar{TestimonyArguments.SolaFide}{52}\) and \(\pvar{TestimonyArguments.SolaFide}{4}\), \(\pvar{TestimonyArguments.SolaFide}{55}\), \(\pvar{TestimonyArguments.SolaFide}{56}\), \(\pvar{TestimonyArguments.SolaFide}{51}\), \((\pvar{TestimonyArguments.SolaFide}{4} \land \pvar{TestimonyArguments.SolaFide}{55} \land \pvar{TestimonyArguments.SolaFide}{56}) \rightarrow \pvar{TestimonyArguments.SolaFide}{57}\); each is needed for the break, and without the crux they stand.
\item \textbf{What this rests on:}
\begin{itemize}
\item \(\pvar{TestimonyArguments.SolaFide}{4}\) Romans 4:4–5 contrasts wages owed with a gift reckoned to the one who believes --- \emph{well supported}: Rom 4:4-5; \cite[III.xi.18]{calvin-institutes-1960}
\item \(\pvar{TestimonyArguments.SolaFide}{55}\) Romans 4:6–8: righteousness counted apart from works — sins forgiven, sin not counted --- \emph{consensus}: Rom 4:6-8; \cite[ad loc. Rom 4:6]{moo-romans-1996}
\item \(\pvar{TestimonyArguments.SolaFide}{56}\) ‘Count’ (λογίζομαι) in Romans 4:3–11 is the language of an account: credited to one --- \emph{well supported}: \cite[s.v. λογίζομαι]{bauer-danker-lexicon-2000}; \cite[ad loc. Rom 4:3]{moo-romans-1996}
\item \(\pvar{TestimonyArguments.SolaFide}{51}\) A word contributes the least meaning its context requires; one occurrence does not carry the whole concept --- \emph{well supported}: \cite{silva-biblical-words-1994}; \cite{joos-semantic-axiom-1972}; \cite{barr-semantics-2004}
\item \(\pvar{TestimonyArguments.SolaFide}{57}\) In Romans 4:6–8 Paul explains righteousness counted apart from works as sin not counted --- \emph{well supported}: \cite[ad loc. Rom 4:6]{moo-romans-1996}; Rom 4:6-8
\item the step itself --- \emph{well supported}: \cite[ad loc. Rom 4:6]{moo-romans-1996}; \cite[s.v. λογίζομαι]{bauer-danker-lexicon-2000}; \emph{well supported}: \cite{silva-biblical-words-1994}; \cite[p. 80]{moo-romans-1996}; \cite{barr-semantics-2004}
\end{itemize}
\end{itemize}
\paragraph{2. Read as what God does in justifying.} The claim commits \emph{Trent, against ‘not by works’} to \(\pvar{TestimonyArguments.SolaFide}{53}\) --- so read by \cite[§22]{lwf-catholic-joint-declaration-2000}; \cite[Trent, Session VI (1547), Decree on Justification, ch. 7]{tanner-decrees-1990} (\emph{disputed}).

Against \emph{Paul’s word (δικαιόω, by the rule of least meaning)}, it is not reached: \emph{Paul’s word (δικαιόω, by the rule of least meaning)} holds, and can be held together with it.

Against \emph{Romans 4:3–8, on Paul’s word (by the rule of least meaning)}, it is not reached: \emph{Romans 4:3–8, on Paul’s word (by the rule of least meaning)} holds, and can be held together with it.

Against \emph{Paul’s gospel (Galatians 1), against justification as renewal}, it cannot be held. Where it breaks:
\paragraph{Why \emph{Paul’s gospel (Galatians 1), against justification as renewal} stands against \emph{Trent, against ‘not by works’, read as what God does in justifying}.}
\begin{itemize}
\item \textbf{The crux:} \((\pvar{TestimonyArguments.SolaFide}{58} \land \pvar{TestimonyArguments.SolaFide}{49} \land \pvar{TestimonyArguments.SolaFide}{48}) \rightarrow \lnot \pvar{TestimonyArguments.SolaFide}{33}\), a premise of \emph{Paul’s gospel (Galatians 1), against justification as renewal}.
\item \textbf{What it does:} its conclusion does not follow without it.
\item \textbf{Granted:} \(\pvar{TestimonyArguments.SolaFide}{44}\), \(\pvar{TestimonyArguments.SolaFide}{45}\), \(\pvar{TestimonyArguments.SolaFide}{46}\), \(\pvar{TestimonyArguments.SolaFide}{47}\), \((\pvar{TestimonyArguments.SolaFide}{44} \land \pvar{TestimonyArguments.SolaFide}{45} \land \pvar{TestimonyArguments.SolaFide}{46} \land \pvar{TestimonyArguments.SolaFide}{47}) \rightarrow \pvar{TestimonyArguments.SolaFide}{58}\), \(\pvar{TestimonyArguments.SolaFide}{48}\), \(\pvar{TestimonyArguments.SolaFide}{49}\), which the break also needs.
\item \textbf{Where the rival breaks:} \(\pvar{TestimonyArguments.SolaFide}{33}\) cannot be held together with \((\pvar{TestimonyArguments.SolaFide}{58} \land \pvar{TestimonyArguments.SolaFide}{49} \land \pvar{TestimonyArguments.SolaFide}{48}) \rightarrow \lnot \pvar{TestimonyArguments.SolaFide}{33}\) and \(\pvar{TestimonyArguments.SolaFide}{44}\), \(\pvar{TestimonyArguments.SolaFide}{45}\), \(\pvar{TestimonyArguments.SolaFide}{46}\), \(\pvar{TestimonyArguments.SolaFide}{47}\), \((\pvar{TestimonyArguments.SolaFide}{44} \land \pvar{TestimonyArguments.SolaFide}{45} \land \pvar{TestimonyArguments.SolaFide}{46} \land \pvar{TestimonyArguments.SolaFide}{47}) \rightarrow \pvar{TestimonyArguments.SolaFide}{58}\), \(\pvar{TestimonyArguments.SolaFide}{48}\), \(\pvar{TestimonyArguments.SolaFide}{49}\); each is needed for the break, and without the crux they stand.
\item \textbf{What this rests on:}
\begin{itemize}
\item \(\pvar{TestimonyArguments.SolaFide}{44}\) Galatians 1:6–9: whoever preaches a gospel contrary to the one received is accursed --- \emph{consensus}: Gal 1:6-9; \cite[ad loc. Gal 1:6]{moo-galatians-2013}
\item \(\pvar{TestimonyArguments.SolaFide}{45}\) 1 Corinthians 15:3: the gospel of first importance — Christ died for our sins --- \emph{consensus}: 1 Cor 15:3
\item \(\pvar{TestimonyArguments.SolaFide}{46}\) Galatians 2:21: if righteousness were through the law, Christ died for no purpose --- \emph{consensus}: Gal 2:21; \cite[ad loc. Gal 2:21]{moo-galatians-2013}
\item \(\pvar{TestimonyArguments.SolaFide}{47}\) Galatians 5:2–4: you who would be justified by the law are severed from Christ --- \emph{consensus}: Gal 5:2-4; \cite[ad loc. Gal 5:2]{moo-galatians-2013}
\item \(\pvar{TestimonyArguments.SolaFide}{58}\) Christ’s death for our sins is the whole ground of justification; to add to it is another gospel --- \emph{disputed}: \cite[XI.1]{westminster-confession-1647}; Gal 2:21; \cite[ad loc. Gal 2:21]{moo-galatians-2013}; \cite[Epistle to Diognetus 9]{roberts-ante-nicene-fathers-1-1885}
\item \(\pvar{TestimonyArguments.SolaFide}{48}\) Romans 8:33–34: it is God who justifies — who is to condemn? --- \emph{consensus}: Rom 8:33-34
\item \(\pvar{TestimonyArguments.SolaFide}{49}\) In Paul, ‘justify’ (δικαιόω) is a verdict: to declare righteous --- \emph{well supported}: \cite[s.v. δικαιόω]{bauer-danker-lexicon-2000}; \cite{morris-apostolic-preaching-1965}; \cite[p. 80]{moo-romans-1996}; \cite{irons-righteousness-of-god-2015}; \cite{wright-justification-2009}; \cite[ad loc. Rom 3:24]{fitzmyer-romans-1993}
\item the step itself --- \emph{disputed}: \cite[ad loc. Gal 1:6]{moo-galatians-2013}; \cite[XI.1]{westminster-confession-1647}; \emph{disputed}: \cite[XI.1]{westminster-confession-1647}; \cite[III.xi.2]{calvin-institutes-1960}; \cite{morris-apostolic-preaching-1965}
\end{itemize}
\end{itemize}
Against \emph{Romans 4:3–8, against justification as renewal}, it cannot be held. Where it breaks:
\paragraph{Why \emph{Romans 4:3–8, against justification as renewal} stands against \emph{Trent, against ‘not by works’, read as what God does in justifying}.}
\begin{itemize}
\item \textbf{The crux:} \((\pvar{TestimonyArguments.SolaFide}{57} \land \pvar{TestimonyArguments.SolaFide}{49}) \rightarrow \lnot \pvar{TestimonyArguments.SolaFide}{33}\), a premise of \emph{Romans 4:3–8, against justification as renewal}.
\item \textbf{What it does:} its conclusion does not follow without it.
\item \textbf{Granted:} \(\pvar{TestimonyArguments.SolaFide}{4}\), \(\pvar{TestimonyArguments.SolaFide}{55}\), \(\pvar{TestimonyArguments.SolaFide}{56}\), \(\pvar{TestimonyArguments.SolaFide}{49}\), \((\pvar{TestimonyArguments.SolaFide}{4} \land \pvar{TestimonyArguments.SolaFide}{55} \land \pvar{TestimonyArguments.SolaFide}{56}) \rightarrow \pvar{TestimonyArguments.SolaFide}{57}\), which the break also needs.
\item \textbf{Where the rival breaks:} \(\pvar{TestimonyArguments.SolaFide}{33}\) cannot be held together with \((\pvar{TestimonyArguments.SolaFide}{57} \land \pvar{TestimonyArguments.SolaFide}{49}) \rightarrow \lnot \pvar{TestimonyArguments.SolaFide}{33}\) and \(\pvar{TestimonyArguments.SolaFide}{4}\), \(\pvar{TestimonyArguments.SolaFide}{55}\), \(\pvar{TestimonyArguments.SolaFide}{56}\), \(\pvar{TestimonyArguments.SolaFide}{49}\), \((\pvar{TestimonyArguments.SolaFide}{4} \land \pvar{TestimonyArguments.SolaFide}{55} \land \pvar{TestimonyArguments.SolaFide}{56}) \rightarrow \pvar{TestimonyArguments.SolaFide}{57}\); each is needed for the break, and without the crux they stand.
\item \textbf{What this rests on:}
\begin{itemize}
\item \(\pvar{TestimonyArguments.SolaFide}{4}\) Romans 4:4–5 contrasts wages owed with a gift reckoned to the one who believes --- \emph{well supported}: Rom 4:4-5; \cite[III.xi.18]{calvin-institutes-1960}
\item \(\pvar{TestimonyArguments.SolaFide}{55}\) Romans 4:6–8: righteousness counted apart from works — sins forgiven, sin not counted --- \emph{consensus}: Rom 4:6-8; \cite[ad loc. Rom 4:6]{moo-romans-1996}
\item \(\pvar{TestimonyArguments.SolaFide}{56}\) ‘Count’ (λογίζομαι) in Romans 4:3–11 is the language of an account: credited to one --- \emph{well supported}: \cite[s.v. λογίζομαι]{bauer-danker-lexicon-2000}; \cite[ad loc. Rom 4:3]{moo-romans-1996}
\item \(\pvar{TestimonyArguments.SolaFide}{49}\) In Paul, ‘justify’ (δικαιόω) is a verdict: to declare righteous --- \emph{well supported}: \cite[s.v. δικαιόω]{bauer-danker-lexicon-2000}; \cite{morris-apostolic-preaching-1965}; \cite[p. 80]{moo-romans-1996}; \cite{irons-righteousness-of-god-2015}; \cite{wright-justification-2009}; \cite[ad loc. Rom 3:24]{fitzmyer-romans-1993}
\item \(\pvar{TestimonyArguments.SolaFide}{57}\) In Romans 4:6–8 Paul explains righteousness counted apart from works as sin not counted --- \emph{well supported}: \cite[ad loc. Rom 4:6]{moo-romans-1996}; Rom 4:6-8
\item the step itself --- \emph{well supported}: \cite[ad loc. Rom 4:6]{moo-romans-1996}; \cite[s.v. λογίζομαι]{bauer-danker-lexicon-2000}; \emph{disputed}: \cite[ad loc. Rom 4:5]{moo-romans-1996}; \cite[XI.1]{westminster-confession-1647}
\end{itemize}
\end{itemize}
\subsection*{Does justification have degrees?}
\addcontentsline{toc}{subsection}{Does justification have degrees?}
Trent’s decree on justification has a chapter “On the increase of Justification received”. The justified, it says,

\begin{quote}
“through the observance of the commandments of God and of the Church, faith co-operating with good works, increase in that justice which they have received through the grace of Christ, and are still further justified” (Session VI, ch. 10).

\end{quote}
Canon 24 condemns anyone who says that good works are “merely the fruits and signs of Justification obtained, but not a cause of the increase thereof”.

The Reformed answer is short, and it starts in a courtroom.

\begin{quote}
“If there is a dispute between men and they come into court and the judges decide between them, acquitting the innocent and condemning the guilty” (Deuteronomy 25:1).

\end{quote}
\begin{quote}
“He who justifies the wicked and he who condemns the righteous are both alike an abomination to the LORD” (Proverbs 17:15).

\end{quote}
Paul uses the same pair of words of God: “It is God who justifies. Who is to condemn?” (Romans 8:33–34). A judge’s verdict acquits or condemns. Nobody leaves a courtroom more acquitted than the man beside him, and nobody is acquitted a little now and more later. So if justification is a verdict, it has no degrees, and there is no “increase of justification” for works to cause.

\subsubsection*{Who has said so}
Calvin answered Trent’s decree within the year, in his \emph{Antidote}. Of the council’s account of justification he wrote: “I maintain that it is one, and simple, and is wholly included in the gratuitous acceptance of God” (p. 116). And of canon 24: “whosoever asserts that works have the effect of increasing justification, understands neither what is the meaning of justification nor its cause” (p. 158). Turretin, more than a century later, gave a whole question of his \emph{Institutes} to the point (Topic XVI, question 10): justification is perfect, so that in itself it neither grows nor shrinks. James Buchanan put it in English in 1867: justification is “an act, too, which is completed at once, and not a work which is gradually accomplished by successive acts”, and “there is no second Justification, properly so called” (p. 251).

Two limits on the courtroom picture. Wright and Dunn speak of a present and a future justification, a verdict given now and affirmed at the last judgment. That is a verdict given twice, not a verdict that grows, and the step here claims only that a verdict has no degrees, not that it is given once. And not every picture of a divine court has two outcomes: the Talmud’s three books, for the wholly righteous, the wholly wicked and those in between (Rosh Hashanah 16b), has three. Those readers do not grant Romans 8:33–34 as Paul reads it, and are not counted against the step.

\subsubsection*{Where the argument has to go}
The courtroom texts settle what a verdict is. They do not settle whether justification is \emph{only} a verdict, and that is where Trent stands. For Trent, justification is “not remission of sins merely, but also the sanctification and renewal of the inward man” (ch. 7), and a renewal can grow. So the argument from the courtroom reaches Trent only with the Reformed definition beside it, that justification is a verdict and nothing more (Westminster XI.1).

The checking bears this out in three results.

\begin{itemize}
  \item \textbf{Granted the Reformed definition, justification has no degrees}, and works cannot increase it (\texttt{verdict\allowbreak{}On\allowbreak{}Degrees\_\allowbreak{}establishes}, \texttt{no\_\allowbreak{}degrees\_\allowbreak{}denies\_\allowbreak{}the\_\allowbreak{}increase}).
  \item \textbf{Trent can grant every courtroom text, and that a verdict has no degrees.} Its increase of justification still stands, because for Trent justification is more than a verdict (\texttt{trent\_\allowbreak{}grants\_\allowbreak{}that\_\allowbreak{}verdicts\_\allowbreak{}have\_\allowbreak{}no\_\allowbreak{}degrees}).
  \item \textbf{So the two part at one claim: whether justification is only a verdict.} Add that claim to Trent’s position, and it can no longer be held (\texttt{degrees\_\allowbreak{}turn\_\allowbreak{}on\_\allowbreak{}the\_\allowbreak{}definition}). The courtroom texts alone leave the question open (\texttt{courtroom\_\allowbreak{}texts\_\allowbreak{}leave\_\allowbreak{}degrees\_\allowbreak{}open}).
\end{itemize}
So the argument that justification has no degrees is sound, and it is no stronger than the definition it rests on. The chapter “What Trent means by justification” weighs that definition. What this chapter adds is the reason the Reformers cared about it: on their definition, Trent’s increase has nothing to increase.

\subsubsection*{What this does not show}
It does not show that the justified do not grow. Both sides hold that the renewal of the justified grows through good works (\texttt{renewal\allowbreak{}Grows\allowbreak{}Through\allowbreak{}Good\allowbreak{}Works}). The Reformed call that growth sanctification, as the Westminster Larger Catechism does when it sets the two side by side (Q. 77): justification “doth equally free all believers from the revenging wrath of God, and that perfectly in this life”; sanctification “is neither equal in all, nor in this life perfect in any, but growing up to perfection”. What the courtroom argument denies is that the growth is an increase of \emph{justification}.

\subsubsection*{How this is checked}
\paragraph*{Trent’s increase, written out first}
\leavevmode\par\nobreak
\par\medskip\noindent\textbf{\texttt{increase\allowbreak{}Means\allowbreak{}Degrees\allowbreak{}Source}}\par\nobreak\noindent
The step from an increase of justification to degrees of it. An increase is more of something, so if works merit an increase of justification, justification admits of more and less. Rated \texttt{consensus}: it is what the word “increase” says, and both sides read Trent’s chapter 10 so.

{\small\begin{verbatim}
def increaseMeansDegreesSource : Source :=
  { primary := Reference.work waterworthTrent (Bib.Locus.page 37),
    supporting := [Reference.work calvinAntidote (Bib.Locus.page 158)],
    tradition := Tradition.romanCatholic, confidence := Confidence.consensus }
\end{verbatim}}
\par\medskip\noindent\textbf{\texttt{increase\allowbreak{}Means\allowbreak{}Degrees\allowbreak{}Line}} --- Trent’s increase, as degrees\par\nobreak\noindent
\textbf{An increase of justification means degrees of it.} If works merit an increase of justification, then justification has degrees.

\begin{derivation}
  \item[(1)] \(\pvar{TestimonyArguments.SolaFide}{32}\)\gloss{Works performed in grace merit an increase of justification}
  \item[(2)] \(\pvar{TestimonyArguments.SolaFide}{32} \rightarrow \pvar{TestimonyArguments.SolaFide}{109}\)
  \item[\(\vdash\)] \(\pvar{TestimonyArguments.SolaFide}{109}\)\gloss{\emph{at issue:} Justification has degrees: the justified are still further justified}
\end{derivation}
\par\medskip\noindent\textbf{\texttt{trent\allowbreak{}On\allowbreak{}Degrees}} --- Trent, on the increase of justification\par\nobreak\noindent
\textbf{Trent, on degrees of justification.} Justification is the renewal of the inward man, the renewal grows through good works, so good works increase justification (ch. 10, canon 24), and justification has degrees.

\begin{derivation}
  \item[(1)] \(\pvar{TestimonyArguments.SolaFide}{33}\)\gloss{Justification is not remission of sins only, but renewal of the inward man}
  \item[(2)] \(\pvar{TestimonyArguments.SolaFide}{35}\)\gloss{The inward renewal of the justified grows as they do good works in grace}
  \item[(3)] \((\pvar{TestimonyArguments.SolaFide}{33} \land \pvar{TestimonyArguments.SolaFide}{35}) \rightarrow \pvar{TestimonyArguments.SolaFide}{32}\)\gloss{\emph{would give:} Works performed in grace merit an increase of justification}
  \item[(4)] \(\pvar{TestimonyArguments.SolaFide}{32} \rightarrow \pvar{TestimonyArguments.SolaFide}{109}\)
  \item[\(\vdash\)] \(\pvar{TestimonyArguments.SolaFide}{109}\)\gloss{\emph{at issue:} Justification has degrees: the justified are still further justified}
  \keepwithnext
\end{derivation}
\emph{No premise here rests on scripture alone.}
\paragraph*{The courtroom answer}
\leavevmode\par\nobreak
\par\medskip\noindent\textbf{\texttt{courtroom\allowbreak{}Source}}\par\nobreak\noindent
The step from the courtroom texts to “a verdict has no degrees”. Judges acquit or condemn (Deuteronomy 25:1; Proverbs 17:15), and God’s justifying is set against condemning in the same way (Romans 8:33–34). Cited to Buchanan (p. 251), with Calvin’s \emph{Antidote} (p. 116) and Turretin (XVI, q. 10).

{\small\begin{verbatim}
def courtroomSource : Source :=
  { primary := Reference.work buchananJustification (Bib.Locus.page 251),
    supporting :=
      [Reference.work calvinAntidote (Bib.Locus.page 116),
        Reference.work turretinInstitutes2
          (Bib.Locus.sectionRef "Topic XVI, q. 10")],
    tradition := Tradition.reformedProtestant,
    confidence := Confidence.wellSupported }
\end{verbatim}}
\par\medskip\noindent\textbf{\texttt{courtroom\allowbreak{}Line}} --- The courtroom (Deuteronomy 25:1; Proverbs 17:15; Romans 8:33–34)\par\nobreak\noindent
\textbf{The courtroom line.} From Deuteronomy 25:1, Proverbs 17:15 and Romans 8:33–34, a verdict acquits or condemns and has no degrees.

\begin{derivation}
  \item[(1)] \(\pvar{TestimonyArguments.SolaFide}{110}\)\gloss{Deuteronomy 25:1, Proverbs 17:15: judges acquit the innocent and condemn the guilty}
  \item[(2)] \(\pvar{TestimonyArguments.SolaFide}{48}\)\gloss{Romans 8:33–34: it is God who justifies — who is to condemn?}
  \item[(3)] \((\pvar{TestimonyArguments.SolaFide}{110} \land \pvar{TestimonyArguments.SolaFide}{48}) \rightarrow \pvar{TestimonyArguments.SolaFide}{111}\)
  \item[\(\vdash\)] \(\pvar{TestimonyArguments.SolaFide}{111}\)\gloss{\emph{at issue:} A verdict acquits or condemns, and has no degrees}
\end{derivation}
\par\medskip\noindent\textbf{\texttt{no\allowbreak{}Degrees\allowbreak{}Source}}\par\nobreak\noindent
The step from “justification is a verdict only” and “a verdict has no degrees” to “justification has no degrees”. Calvin against canon 24 (\emph{Antidote}, p. 158), and Turretin (XVI, q. 10).

{\small\begin{verbatim}
def noDegreesSource : Source :=
  { primary := Reference.work calvinAntidote (Bib.Locus.page 158),
    supporting :=
      [Reference.work turretinInstitutes2
          (Bib.Locus.sectionRef "Topic XVI, q. 10")],
    tradition := Tradition.reformedProtestant,
    confidence := Confidence.wellSupported }
\end{verbatim}}
\par\medskip\noindent\textbf{\texttt{no\allowbreak{}Degrees\allowbreak{}Line}} --- A verdict only, so no degrees\par\nobreak\noindent
\textbf{No degrees.} If justification is a verdict and nothing more, and a verdict has no degrees, then justification has none.

\begin{derivation}
  \item[(1)] \(\pvar{TestimonyArguments.SolaFide}{30}\)\gloss{Justification is a verdict only: pardon and righteousness counted, not poured in}
  \item[(2)] \(\pvar{TestimonyArguments.SolaFide}{111}\)\gloss{A verdict acquits or condemns, and has no degrees}
  \item[(3)] \((\pvar{TestimonyArguments.SolaFide}{30} \land \pvar{TestimonyArguments.SolaFide}{111}) \rightarrow \lnot \pvar{TestimonyArguments.SolaFide}{109}\)
  \item[\(\vdash\)] \(\lnot \pvar{TestimonyArguments.SolaFide}{109}\)\gloss{\emph{at issue, the denial of:} Justification has degrees: the justified are still further justified}
\end{derivation}
\par\medskip\noindent\textbf{\texttt{verdict\allowbreak{}On\allowbreak{}Degrees}} --- Justification as a verdict, which has no degrees\par\nobreak\noindent
\textbf{The Reformed position on degrees.} The courtroom texts, that a verdict has no degrees, and the Reformed definition, that justification is a verdict only. So justification has no degrees. It also holds what an increase means, so that it can say what it denies of Trent.

\begin{derivation}
  \item[(1)] \(\pvar{TestimonyArguments.SolaFide}{110}\)\gloss{Deuteronomy 25:1, Proverbs 17:15: judges acquit the innocent and condemn the guilty}
  \item[(2)] \(\pvar{TestimonyArguments.SolaFide}{48}\)\gloss{Romans 8:33–34: it is God who justifies — who is to condemn?}
  \item[(3)] \((\pvar{TestimonyArguments.SolaFide}{110} \land \pvar{TestimonyArguments.SolaFide}{48}) \rightarrow \pvar{TestimonyArguments.SolaFide}{111}\)\gloss{\emph{would give:} A verdict acquits or condemns, and has no degrees}
  \breakrow{(4)} \(\pvar{TestimonyArguments.SolaFide}{30}\)\gloss{Justification is a verdict only: pardon and righteousness counted, not poured in}
  \breakrow{(5)} \((\pvar{TestimonyArguments.SolaFide}{30} \land \pvar{TestimonyArguments.SolaFide}{111}) \rightarrow \lnot \pvar{TestimonyArguments.SolaFide}{109}\)
  \item[(6)] \(\pvar{TestimonyArguments.SolaFide}{32} \rightarrow \pvar{TestimonyArguments.SolaFide}{109}\)\gloss{\emph{would give:} Justification has degrees: the justified are still further justified}
  \item[\(\vdash\)] \(\lnot \pvar{TestimonyArguments.SolaFide}{109}\)\gloss{\emph{at issue, the denial of:} Justification has degrees: the justified are still further justified}
  \keepwithnext
\end{derivation}
\emph{Grounded in scripture alone:} Romans 8:33–34: it is God who justifies — who is to condemn?.
\par\medskip\noindent\textbf{\texttt{verdict\allowbreak{}Against\allowbreak{}Increase}} --- Justification as a verdict, against the increase\par\nobreak\noindent
\textbf{The same position, against Trent’s increase.} It concludes what canon 24 condemns: good works do not merit an increase of justification.

\begin{derivation}
  \item[(1)] \(\pvar{TestimonyArguments.SolaFide}{110}\)\gloss{Deuteronomy 25:1, Proverbs 17:15: judges acquit the innocent and condemn the guilty}
  \item[(2)] \(\pvar{TestimonyArguments.SolaFide}{48}\)\gloss{Romans 8:33–34: it is God who justifies — who is to condemn?}
  \item[(3)] \((\pvar{TestimonyArguments.SolaFide}{110} \land \pvar{TestimonyArguments.SolaFide}{48}) \rightarrow \pvar{TestimonyArguments.SolaFide}{111}\)\gloss{\emph{would give:} A verdict acquits or condemns, and has no degrees}
  \breakrow{(4)} \(\pvar{TestimonyArguments.SolaFide}{30}\)\gloss{Justification is a verdict only: pardon and righteousness counted, not poured in}
  \breakrow{(5)} \((\pvar{TestimonyArguments.SolaFide}{30} \land \pvar{TestimonyArguments.SolaFide}{111}) \rightarrow \lnot \pvar{TestimonyArguments.SolaFide}{109}\)\gloss{\emph{would deny:} Justification has degrees: the justified are still further justified}
  \item[(6)] \(\pvar{TestimonyArguments.SolaFide}{32} \rightarrow \pvar{TestimonyArguments.SolaFide}{109}\)\gloss{\emph{would give:} Justification has degrees: the justified are still further justified}
  \item[\(\vdash\)] \(\lnot \pvar{TestimonyArguments.SolaFide}{32}\)\gloss{\emph{at issue, the denial of:} Works performed in grace merit an increase of justification}
  \keepwithnext
\end{derivation}
\emph{Grounded in scripture alone:} Romans 8:33–34: it is God who justifies — who is to condemn?.
\paragraph*{The readings}
\leavevmode\par\nobreak
\par\medskip\noindent\textbf{\texttt{verdict\allowbreak{}On\allowbreak{}Degrees\allowbreak{}Reading}}\par\nobreak\noindent
The Reformed reading: every courtroom text holds, a verdict has no degrees, justification is a verdict only, and so it has no degrees and nothing increases it. Trent’s definition does not hold.

{\small\begin{verbatim}
def verdictOnDegreesReading : Valuation Claim :=
  fun a =>
    match a with
    | Claim.justificationHasDegrees => False
    | Claim.worksMeritIncreaseOfJustification => False
    | Claim.justificationIncludesSanctification => False
    | x => True
\end{verbatim}}
\par\medskip\noindent\textbf{\texttt{trent\allowbreak{}On\allowbreak{}Degrees\allowbreak{}Reading}}\par\nobreak\noindent
Trent’s reading, granting the courtroom texts and that a verdict has no degrees: justification is more than a verdict, so its increase stands.

{\small\begin{verbatim}
def trentOnDegreesReading : Valuation Claim :=
  fun a =>
    match a with
    | Claim.justificationIsForensicOnly => False
    | x => True
\end{verbatim}}
\paragraph*{Results}
\leavevmode\par\nobreak
\par\medskip\noindent\textbf{\texttt{verdict\allowbreak{}On\allowbreak{}Degrees\_\allowbreak{}establishes}}\par\nobreak\noindent
\textbf{Granted the Reformed definition, justification has no degrees.} From the courtroom texts, a verdict acquits or condemns and has no degrees; justification is a verdict and nothing more; so it has no degrees.

It does not show that justification is only a verdict. That is the Reformed definition, and Trent denies it.

{\small\begin{verbatim}
theorem verdictOnDegrees_establishes : Establishes verdictOnDegrees
-- axioms: propext, Quot.sound
\end{verbatim}}
\par\medskip\noindent\textbf{\texttt{verdict\allowbreak{}On\allowbreak{}Degrees\_\allowbreak{}is\_\allowbreak{}satisfiable}}\par\nobreak\noindent
The Reformed position on degrees can be held without contradiction.

{\small\begin{verbatim}
theorem verdictOnDegrees_is_satisfiable : Satisfiable
    verdictOnDegrees.premises
-- axioms: propext, Quot.sound
\end{verbatim}}
\par\medskip\noindent\textbf{\texttt{no\_\allowbreak{}degrees\_\allowbreak{}denies\_\allowbreak{}the\_\allowbreak{}increase}}\par\nobreak\noindent
\textbf{On the same premises, works do not merit an increase of justification.} This is what canon 24 condemns, reached from the courtroom and the Reformed definition.

It does not show that the justified do not grow. Both sides hold that the renewal of the justified grows; the Reformed call it sanctification.

{\small\begin{verbatim}
theorem no_degrees_denies_the_increase : Establishes verdictAgainstIncrease
-- axioms: propext, Quot.sound
\end{verbatim}}
\par\medskip\noindent\textbf{\texttt{verdict\allowbreak{}Against\allowbreak{}Increase\_\allowbreak{}is\_\allowbreak{}satisfiable}}\par\nobreak\noindent
The same premises, so the same reading shows they can be held together.

{\small\begin{verbatim}
theorem verdictAgainstIncrease_is_satisfiable : Satisfiable
    verdictAgainstIncrease.premises
-- axioms: propext, Quot.sound
\end{verbatim}}
\par\medskip\noindent\textbf{\texttt{trent\allowbreak{}On\allowbreak{}Degrees\_\allowbreak{}establishes}}\par\nobreak\noindent
\textbf{Trent’s increase follows from its own premises.} Justification is the renewal of the inward man, the renewal grows through good works, and so justification increases and has degrees.

It does not show that Trent’s definition is right. The Westminster Confession denies it (XI.1).

{\small\begin{verbatim}
theorem trentOnDegrees_establishes : Establishes trentOnDegrees
-- axioms: propext, Quot.sound
\end{verbatim}}
\par\medskip\noindent\textbf{\texttt{trent\allowbreak{}On\allowbreak{}Degrees\_\allowbreak{}is\_\allowbreak{}satisfiable}}\par\nobreak\noindent
Trent’s position on degrees can be held without contradiction.

{\small\begin{verbatim}
theorem trentOnDegrees_is_satisfiable : Satisfiable trentOnDegrees.premises
-- axioms: propext, Quot.sound
\end{verbatim}}
\par\medskip\noindent\textbf{\texttt{trent\_\allowbreak{}grants\_\allowbreak{}that\_\allowbreak{}verdicts\_\allowbreak{}have\_\allowbreak{}no\_\allowbreak{}degrees}}\par\nobreak\noindent
\textbf{Trent can grant every courtroom text, and that a verdict has no degrees.} Its position on the increase of justification can be held together with Deuteronomy 25:1, Proverbs 17:15, Romans 8:33–34 and the step from them to “a verdict has no degrees”. For Trent justification is more than a verdict, so a verdict’s having no degrees does not touch it.

It does not show that Trent says verdicts have no degrees. The result says only that nothing in its position denies it.

{\small\begin{verbatim}
theorem trent_grants_that_verdicts_have_no_degrees : Satisfiable
    (trentOnDegrees.premises ++ courtroomLine.premises)
-- axioms: propext, Quot.sound
\end{verbatim}}
\par\medskip\noindent\textbf{\texttt{degrees\_\allowbreak{}turn\_\allowbreak{}on\_\allowbreak{}the\_\allowbreak{}definition}}\par\nobreak\noindent
\textbf{The two part at one claim: whether justification is only a verdict.} Add to Trent’s position that a verdict has no degrees, that justification is a verdict only, and the step from those two to “no degrees”, and it can no longer be held. Together with \texttt{trent\_\allowbreak{}grants\_\allowbreak{}that\_\allowbreak{}verdicts\_\allowbreak{}have\_\allowbreak{}no\_\allowbreak{}degrees}, this puts the whole disagreement about degrees at the Reformed definition.

It does not claim which side is right about the definition. That is weighed in the chapter “What Trent means by justification”.

{\small\begin{verbatim}
theorem degrees_turn_on_the_definition : ¬Satisfiable (trentOnDegrees.premises
    ++ [p Claim.justificationIsForensicOnly, p Claim.verdictAdmitsNoDegrees,
    noDegreesLine.step])
-- axioms: propext, Classical.choice, Quot.sound
\end{verbatim}}
\par\medskip\noindent\textbf{\texttt{courtroom\_\allowbreak{}texts\_\allowbreak{}leave\_\allowbreak{}degrees\_\allowbreak{}open}}\par\nobreak\noindent
\textbf{The courtroom texts alone leave the question open.} Deuteronomy 25:1, Proverbs 17:15 and Romans 8:33–34 are kept alike by the Reformed reading, on which justification has no degrees, and by Trent’s, on which it has. What decides between them is the definition, not the texts.

{\small\begin{verbatim}
theorem courtroom_texts_leave_degrees_open : Independent [p
    Claim.deut25_1JudgesAcquitOrCondemn, p Claim.romans8_33_34] (p
    Claim.justificationHasDegrees)
-- axioms: propext, Quot.sound
\end{verbatim}}
\subsection*{James 2:24: “not by faith alone”}
\addcontentsline{toc}{subsection}{James 2:24: “not by faith alone”}
\begin{quote}
“You see that a person is justified by works and not by faith alone” (James 2:24).

\end{quote}
Few verses have been argued over more in this debate. Trent read it as teaching the increase of justification. The justified, says the decree,

\begin{quote}
“increase in that justice which they have received ... and are still further justified, as it is written; He that is just, let him be justified still; and again, Be not afraid to be justified even to death; and also, Do you see that by works a man is justified, and not by faith only” (Session VI, ch. 10).

\end{quote}
Canon 24 condemns anyone who says good works “are merely the fruits and signs of Justification obtained, but not a cause of the increase thereof”. This chapter weighs that reading, the three texts the decree cites for it, and the Reformed reply.

\subsubsection*{Two ways of taking Trent’s reading}
Trent’s claim about James 2:24 can be taken two ways, and each meets different evidence.

\begin{itemize}
  \item \textbf{As a claim about James’s word.} “Justified” (δικαιόω) in James \emph{means} the increase of the justice already received: being made more just. Then it meets the evidence about the word. Outside Paul, the word declares, acquits or vindicates: “wisdom is justified by her deeds” (Matthew 11:19, ESV), where “deeds” translates ἔργα, the word James uses for “works”. James uses it of Abraham offering Isaac (Genesis 22), long after “Abraham believed God, and it was counted to him as righteousness” (James 2:23, quoting Genesis 15:6), and of Rahab. And a word contributes the least meaning its context requires.
  \item \textbf{As a claim about what works do.} James’s “justified” is a verdict, and the works it is “by” are a cause of that verdict’s increase, not only its fruit (canon 24). Nothing about the word bears on this. It meets the Reformed reading, on which the works are evidence of a living faith, and at that point each side rests on its own confession.
\end{itemize}
\subsubsection*{The other two texts Trent cites}
Two of Trent’s three texts read differently in the Latin Trent quotes and in the Greek.

“He that is just, let him be justified still” is the Latin Vulgate’s wording (\emph{iustificetur adhuc}) at Revelation 22:11. The Greek, as modern critical editions print it, says “let the righteous still do righteousness”, and the Catholic Church’s own English translation gives “The righteous must still do right”. The printed Textus Receptus of the sixteenth century reads “let him be justified still” (δικαιωθήτω), as the Latin does; the critical editions and the Byzantine Majority Text do not.

“Be not afraid to be justified even to death” is the Vulgate’s Sirach 18:22. The Greek does have the verb “to be justified” (δικαιωθῆναι), but in a sentence about vows: “let nothing hinder you from paying a vow in due time, and do not wait until death to be justified”. The Catholic Church’s own translation renders the clause “do not wait until death to fulfill them”. So the Latin and the Greek share “until death ... to be justified”. They differ in the rest: the Latin has “pray always” and “be not afraid”, where the Greek has “pay a vow” and “do not wait”.

Pius XII teaches that the original text, “written by the inspired author himself, has more authority and greater weight than any even the very best translation” (\emph{Divino Afflante Spiritu} §16). Westminster says the same (I.8). So taken as claims about what the inspired authors wrote, the two Latin readings cannot be defended.

Pius XII does not say that a Latin reading absent from the original is an error. Of the Vulgate as a whole he says the opposite. As the Church understands it, it is “free from any error whatsoever in matters of faith and morals”, and its authority is “juridical”, that is, a matter of the Church’s approval, and not a claim about the original wording (§21). Taken that way, Trent’s “as it is written” at these two verses quotes the Church’s approved Latin and not the inspired author’s words, and is consistent with everything Pius XII says (the chapter “Trent’s two Latin texts”). What it then rests on is the Church’s authority. Pius XII himself asks that a doctrine taught in this way be confirmed from the original texts (§22), and at these two verses the originals do not confirm it.

\subsubsection*{The phrase “faith alone”}
James 2:24 is the one place in the New Testament where the phrase occurs, and it is often cited as denying \emph{sola fide} in so many words. The faith James says does not justify is the faith he calls dead (2:17, 2:26). Trent reads it so: faith “unless hope and charity be added thereto ... For which reason it is most truly said, that Faith without works is dead” (Session VI, ch. 7). The Westminster Confession, a Reformed confession of 1647, says the faith that alone justifies “is no dead faith, but worketh by love” (Westminster XI.2). So the phrase James denies is not the formula the Reformers affirm (\texttt{james\allowbreak{}Faith\allowbreak{}Line}). This answers an argument from the phrase. It does not answer James’s “by works”, which is the question above.

\subsubsection*{The words used here}
These are the words of the first chapter, which “Weighing the positions” uses too. One position \textbf{defeats} another when it contradicts something the other rests on, concludes or derives, and is not the weaker. A position is \textbf{forced} when it must be accepted however every tie is settled. It \textbf{can be defended} when some consistent set of positions includes it and answers every attack. It \textbf{cannot be defended} when none does.

\subsubsection*{What the debate decides}
\begin{itemize}
  \item \textbf{Forced} (\texttt{james\_\allowbreak{}words\_\allowbreak{}prevail}). James’s “justified” does not mean the increase of justice. The faith James 2:24 denies is not the Reformers’ faith. Revelation 22:11, in the critical Greek text, does not say “be justified still”. Sirach 18:22 in Greek does not say “be not afraid to be justified even to death”.
  \item \textbf{Cannot be defended.} Trent’s reading of James’s word (\texttt{trent\_\allowbreak{}on\_\allowbreak{}james\_\allowbreak{}word\_\allowbreak{}indefensible}), and its use of Revelation 22:11 (\texttt{trent\_\allowbreak{}on\_\allowbreak{}revelation\_\allowbreak{}indefensible}) and Sirach 18:22 (\texttt{trent\_\allowbreak{}on\_\allowbreak{}sirach\_\allowbreak{}indefensible}) as texts for the increase of justification.
  \item \textbf{Can be defended, and not forced.} Trent’s reading of what works do (\texttt{trent\_\allowbreak{}on\_\allowbreak{}works\_\allowbreak{}defensible}, \texttt{trent\_\allowbreak{}on\_\allowbreak{}works\_\allowbreak{}not\_\allowbreak{}forced}), and, against it, the Reformed reconciliation of James with Paul. Whether works are only the fruit of justification is contested at the ratings the sources were given, before they are recomputed under the three standards (\texttt{works\allowbreak{}Contested}). Trent’s reading of the texts that promise a reward (chapter 16, \texttt{trent\allowbreak{}Reward\allowbreak{}Line}) meets the Reformed reconciliation in the same way. It also meets Westminster’s reading of the same texts (\texttt{westminster\allowbreak{}Reward\allowbreak{}Line}), which denies that the reward is given to merits. When only evidence anyone can check is counted, neither of Trent’s readings of what works do stands as an objection, and the Reformed reconciliation is forced (the chapter “Why each step is rated as it is”).
  \item \textbf{Cannot be defended, at the ratings the sources were given.} The reading on which James’s word means renewal because Paul’s does (\texttt{renewal\allowbreak{}Word\allowbreak{}Line}). The argument from James’s word defeats it. When a tradition’s own teaching is allowed to count as an objection, it does stand as one, and then the question is contested (the chapter “Why each step is rated as it is”). It also meets the Reformed reconciliation, and the two defeat each other: the reconciliation says James’s word does not have Paul’s sense, and this reading says it does.
\end{itemize}
So of the three texts chapter 10 cites, one is left: James 2:24, taken as a claim about what works do before God. Each reading of the three that rests on what a word or a verse \emph{says} cannot be defended. Trent’s reading of the reward texts (chapter 16) can also be defended. No other text was weighed here. What remains is a claim about God’s verdict, and there these pages cannot choose.

\subsubsection*{A rating that was lowered, against the Reformed side}
The premise that works are the fruit of faith and not its ground (\texttt{works\allowbreak{}Are\allowbreak{}Fruit\allowbreak{}Not\allowbreak{}Ground}) was rated \texttt{well\allowbreak{}Supported} while no position here denied it. In this debate Trent’s reading of what works do denies it, and the rating decides who prevails. So the rating was re-examined. Canon 24 condemns exactly that claim. As a premise it is now \texttt{disputed}. Left at \texttt{well\allowbreak{}Supported}, it would have made Trent’s reading impossible to defend, and that verdict would have come from the rating and not from the evidence.

As a step (\texttt{fruit\allowbreak{}Line}), its rating is now worked out and not simply assigned (the chapter “Why each step is rated as it is”). Its support is \texttt{well\allowbreak{}Supported}: the \emph{Joint Declaration} confesses works to be fruits (§37) and not the basis of justification (§25). On whether works contribute anything further, it records the sides apart: Catholics hold that good works “contribute to growth in grace” (§38). Each of Trent’s two objections to it rests, at one claim, on the council’s own word. So the step stands when only checkable evidence counts, and is \texttt{disputed} when a tradition’s teaching counts. This debate weighs the Reformed reconciliation at its weakest link in any case: that James is attacking a barren faith is rated \texttt{disputed} (\texttt{james\allowbreak{}Harmony\allowbreak{}Source}).

\subsubsection*{What each position rests on}
\texttt{james\_\allowbreak{}warrants} lists, for each position, the claims it makes on someone’s word alone. Each of Trent’s five readings rests on Trent’s word: four at the reading itself, and the reading of the reward texts at the claim that the reward is given to merits. The reading through Paul rests on Trent’s word at one claim, that James uses the word as Paul does. The two readings of the Greek rest, at “the original text outweighs any translation”, on the word of Pius XII and of Westminster. Both sides accept it, so it decides the matter between them. It is still a claim taken on authority. The argument from James’s word, the argument about James’s faith, and the Reformed reconciliation rest on nothing but Scripture and evidence. This is reported beside the verdicts and does not affect them.

\subsubsection*{Westminster’s texts for “fruit, not ground”}
The Reformed reconciliation does not assume that works are the fruit of faith and not its ground. It derives it (\texttt{fruit\allowbreak{}Line}) from the texts the Westminster Assembly gave as proofs. Good works “are the fruits and evidences of a true and lively faith” (XVI.2: James 2:18, 22; Ephesians 2:10; Romans 6:22). And “we cannot by our best works merit” (XVI.5: Ephesians 2:8–9; Luke 17:10).

Trent accepts the texts, and the word “fruit”: its chapter 16 is headed “On the fruit of Justification, that is, on the merit of good works”, and canon 24 condemns only those who call works “\emph{merely} the fruits and signs”. The disagreement is whether the fruit is also merit. Trent’s own answer, that the reward the texts promise is given to merits, is a claim it makes on its own word (“we must believe”, chapter 16).

So the Reformed reconciliation now rests on Scripture and not on the Confession’s word. In this debate, at the ratings the sources were given, whether works are only fruit is still contested (\texttt{works\allowbreak{}Contested}). When only checkable evidence counts, it is not (the chapter “Why each step is rated as it is”).

\subsubsection*{What has been re-examined, and what has not}
The steps in the argument from James’s word and the argument about James’s faith are cited at \texttt{well\allowbreak{}Supported}.

The rating of the step about the word is now worked out as well (the chapter “Why each step is rated as it is”). The strongest known objection to it is the reading that takes James’s word as Paul’s and Paul’s as renewal (\texttt{renewal\allowbreak{}Word\allowbreak{}Line}). That reading accepts the step’s starting points and argues for the opposite conclusion. But one of its own starting points, that James uses the word as Paul does, rests on Trent’s word alone. So the step stands at \texttt{well\allowbreak{}Supported} when only checkable evidence counts, and is \texttt{disputed} when a tradition’s teaching counts. This debate weighs it at its cited rating, the first of these. A scholar, verified, who holds that James’s word is Paul’s would make that objection stand either way, and the step \texttt{disputed} either way.

The step about James’s faith has no objection written down here beyond Howell’s, which states the opposite without arguing it (\texttt{Testimony\allowbreak{}Arguments.Articles.Howell2003}). A fuller search across the Catholic commentaries on James is still to be done.

\subsubsection*{How this is checked}
\paragraph*{How each step is rated}
\leavevmode\par\nobreak
\par\medskip\noindent\textbf{\texttt{james\allowbreak{}Lexical\allowbreak{}Source}}\par\nobreak\noindent
The step from three things to “James’s ‘justified’ does not mean the increase of justice”: how the word (δικαιόω) is used outside Paul, how James himself uses it, and the rule of least meaning. Its rating has yet to be re-examined.

{\small\begin{verbatim}
def jamesLexicalSource : Source :=
  {
    primary :=
      Reference.work mooJames
        (Bib.Locus.adLoc { book := Book.james, chapter := 2, verse := 21 }),
    supporting :=
      [Reference.work silvaBiblicalWords, Reference.work barrSemantics],
    tradition := Tradition.christianHistoricalGrammatical,
    confidence := Confidence.wellSupported }
\end{verbatim}}
\par\medskip\noindent\textbf{\texttt{james\allowbreak{}Faith\allowbreak{}Source}}\par\nobreak\noindent
The step from James 2:14–26, the dead faith it names, and what the Reformers’ formula means, to “the faith James 2:24 denies is not the Reformers’ faith”. Its rating has yet to be re-examined.

{\small\begin{verbatim}
def jamesFaithSource : Source :=
  {
    primary :=
      Reference.work mooJames
        (Bib.Locus.adLoc { book := Book.james, chapter := 2, verse := 24 }),
    supporting :=
      [Reference.work westminsterConfession (Bib.Locus.sectionRef "XI.2")],
    tradition := Tradition.christianHistoricalGrammatical,
    confidence := Confidence.wellSupported }
\end{verbatim}}
\par\medskip\noindent\textbf{\texttt{trent\allowbreak{}Increase\allowbreak{}Source}}\par\nobreak\noindent
Trent’s reading of its three texts as teaching the increase of justification (Session VI, ch. 10). Rated \texttt{disputed}: the Reformed accept the texts and deny that they teach it.

{\small\begin{verbatim}
def trentIncreaseSource : Source :=
  {
    primary :=
      Reference.work tannerDecrees
        (Bib.Locus.sectionRef "Trent, Session VI (1547), ch. 10"),
    tradition := Tradition.romanCatholic, confidence := Confidence.disputed }
\end{verbatim}}
\par\medskip\noindent\textbf{\texttt{trent\allowbreak{}Works\allowbreak{}Source}}\par\nobreak\noindent
Trent’s reading of what works do (Session VI, canon 24), applied to “by works” in James 2:24. Rated \texttt{disputed}: the Reformed accept James’s words and read the works as evidence, not cause.

{\small\begin{verbatim}
def trentWorksSource : Source :=
  {
    primary :=
      Reference.work tannerDecrees
        (Bib.Locus.sectionRef "Trent, Session VI (1547), canon 24"),
    tradition := Tradition.romanCatholic, confidence := Confidence.disputed }
\end{verbatim}}
\par\medskip\noindent\textbf{\texttt{textual\allowbreak{}Source}}\par\nobreak\noindent
The step from what the original text says, and the principle that it outweighs any translation, to “the inspired text does not say what the Latin says”. Rated \texttt{well\allowbreak{}Supported}: Pius XII (\emph{Divino Afflante Spiritu} §16) and Westminster (I.8) both make the original the measure of a translation, and no reader was found who accepts both and holds that the inspired author wrote what only the Latin has.

{\small\begin{verbatim}
def textualSource : Source :=
  {
    primary :=
      Reference.work divinoAfflanteSpiritu (Bib.Locus.sectionRef "§16"),
    supporting := [Reference.work Bib.sblgnt, Reference.work nabre],
    tradition := Tradition.christianHistoricalGrammatical,
    confidence := Confidence.wellSupported }
\end{verbatim}}
\par\medskip\noindent\textbf{\texttt{james\allowbreak{}Harmony\allowbreak{}Source}}\par\nobreak\noindent
The step in the Reformed reconciliation: from “James attacks a barren faith” and “works are the fruit of saving faith” to “James 2:24 does not contradict Paul”. Rated \texttt{disputed}: Trent accepts James’s words and reads the works as a cause of the increase of justification (Session VI, ch. 10, canon 24).

{\small\begin{verbatim}
def jamesHarmonySource : Source :=
  {
    primary :=
      Reference.work mooJames
        (Bib.Locus.adLoc { book := Book.james, chapter := 2, verse := 24 }),
    supporting :=
      [Reference.work Bib.calvinInstitutes
          (Bib.Locus.sectionRef "III.xvii.11")],
    tradition := Tradition.reformedProtestant,
    confidence := Confidence.disputed }
\end{verbatim}}
\par\medskip\noindent\textbf{\texttt{trent\allowbreak{}Reward\allowbreak{}Source}}\par\nobreak\noindent
Trent’s step from the texts that promise a reward, and the reward as given to merits, to “works are not merely fruit”. Eternal life is “a reward to be faithfully rendered to their good works and merits” (Session VI, ch. 16), and the good works of the justified truly merit it (canon 32). Rated \texttt{disputed}: Westminster accepts that God rewards the good works of believers, and holds the reward to be of grace, given to works accepted in Christ and not earned by them (XVI.5, XVI.6).

{\small\begin{verbatim}
def trentRewardSource : Source :=
  {
    primary :=
      Reference.work tannerDecrees
        (Bib.Locus.sectionRef "Trent, Session VI (1547), ch. 16, canon 32"),
    supporting :=
      [Reference.work westminsterConfession
          (Bib.Locus.sectionRef "XVI.5, XVI.6")],
    tradition := Tradition.romanCatholic, confidence := Confidence.disputed }
\end{verbatim}}
\par\medskip\noindent\textbf{\texttt{renewal\allowbreak{}Word\allowbreak{}Source}}\par\nobreak\noindent
The step in the reading that takes James’s word through Paul’s: from Paul’s “justify” as renewal, James’s word as Paul’s, and James’s “justified by works” coming after Genesis 15:6, to “James’s ‘justified’ means the increase of justice” (Trent, Session VI, ch. 10). Rated \texttt{disputed}: the argument from James’s word accepts James’s words and denies the conclusion.

{\small\begin{verbatim}
def renewalWordSource : Source :=
  {
    primary :=
      Reference.work tannerDecrees
        (Bib.Locus.sectionRef "Trent, Session VI (1547), ch. 10"),
    supporting := [Reference.work vanLandinghamJudgment],
    tradition := Tradition.romanCatholic, confidence := Confidence.disputed }
\end{verbatim}}
\par\medskip\noindent\textbf{\texttt{westminster\allowbreak{}Reward\allowbreak{}Source}}\par\nobreak\noindent
Westminster’s step from the texts it gives as proofs to “the reward is not given to merits”. God does not forget the work of believers (Hebrews 6:10, a proof of XVI.6, and one of Trent’s own four texts). When they have done all, they are unworthy servants who have done only their duty (Luke 17:10, XVI.5). To one who works, wages are counted as due, and to one who believes, as a gift (Romans 4:4–5, XVI.5). And their sacrifices are acceptable to God through Jesus Christ (1 Peter 2:5, XVI.6). So God “is pleased to accept and reward” works accepted in his Son, and the reward is of grace, not merit (XVI.5, XVI.6). Rated \texttt{disputed}: Trent accepts these texts and holds the reward to be given to merits (Session VI, ch. 16, canon 32).

{\small\begin{verbatim}
def westminsterRewardSource : Source :=
  {
    primary :=
      Reference.work westminsterConfession
        (Bib.Locus.sectionRef "XVI.5, XVI.6"),
    tradition := Tradition.reformedProtestant,
    confidence := Confidence.disputed }
\end{verbatim}}
\paragraph*{The lines of argument}
\leavevmode\par\nobreak
\par\medskip\noindent\textbf{\texttt{james\allowbreak{}Lexical\allowbreak{}Line}} --- James’s word (δικαιόω, by the rule of least meaning)\par\nobreak\noindent
\textbf{James’s word}, read by the rule of least meaning. Outside Paul, “justify” (δικαιόω) declares, acquits or vindicates. James uses it of Abraham at Genesis 22, after Genesis 15:6 had been spoken, and of Rahab. And a word contributes the least meaning its context requires. So James’s “justified” does not mean the increase of justice.

\begin{derivation}
  \item[(1)] \(\pvar{TestimonyArguments.SolaFide}{85}\)\gloss{Outside Paul, ‘justify’ (δικαιόω) declares or vindicates; it does not mean making righteous}
  \item[(2)] \(\pvar{TestimonyArguments.SolaFide}{84}\)\gloss{James 2:21–25: Abraham justified by works at Genesis 22, fulfilling Genesis 15:6; Rahab justified by works}
  \item[(3)] \(\pvar{TestimonyArguments.SolaFide}{51}\)\gloss{A word contributes the least meaning its context requires; one occurrence does not carry the whole concept}
  \item[(4)] \((\pvar{TestimonyArguments.SolaFide}{85} \land \pvar{TestimonyArguments.SolaFide}{84} \land \pvar{TestimonyArguments.SolaFide}{51}) \rightarrow \lnot \pvar{TestimonyArguments.SolaFide}{86}\)
  \item[\(\vdash\)] \(\lnot \pvar{TestimonyArguments.SolaFide}{86}\)\gloss{\emph{at issue, the denial of:} James’s ‘justified’ (δικαιόω) means the increase of the justice received}
\end{derivation}
\par\medskip\noindent\textbf{\texttt{james\allowbreak{}Faith\allowbreak{}Line}} --- James’s faith (the faith James calls dead)\par\nobreak\noindent
\textbf{James’s faith.} The faith that James 2:24 says does not justify is the faith apart from works that he calls dead. The faith the Reformers say alone justifies is no dead faith. So the phrase James denies is not the formula they affirm.

\begin{derivation}
  \item[(1)] \(\pvar{TestimonyArguments.SolaFide}{83}\)\gloss{James 2:14–26: faith without works cannot save, and is dead (2:17, 2:26)}
  \item[(2)] \(\pvar{TestimonyArguments.SolaFide}{87}\)\gloss{The faith alone that James 2:24 denies justifies is the dead faith of 2:14–26}
  \item[(3)] \(\pvar{TestimonyArguments.SolaFide}{88}\)\gloss{The faith the Reformers say alone justifies is no dead faith, but works by love}
  \item[(4)] \((\pvar{TestimonyArguments.SolaFide}{83} \land \pvar{TestimonyArguments.SolaFide}{87} \land \pvar{TestimonyArguments.SolaFide}{88}) \rightarrow \pvar{TestimonyArguments.SolaFide}{89}\)
  \item[\(\vdash\)] \(\pvar{TestimonyArguments.SolaFide}{89}\)\gloss{\emph{at issue:} The faith James 2:24 denies justifies is not the faith of the Reformers’ formula}
\end{derivation}
\par\medskip\noindent\textbf{\texttt{catholic\allowbreak{}James\allowbreak{}Word\allowbreak{}Line}} --- James 2:24, read for the increase of justification\par\nobreak\noindent
\textbf{Trent’s reading of James’s word}: “justified by works” is the increase of the justice received, so James 2:24 teaches the increase of justification.

\begin{derivation}
  \item[(1)] \(\pvar{TestimonyArguments.SolaFide}{23}\)\gloss{James 2:24 says a person is justified by works and not by faith alone}
  \item[(2)] \(\pvar{TestimonyArguments.SolaFide}{84}\)\gloss{James 2:21–25: Abraham justified by works at Genesis 22, fulfilling Genesis 15:6; Rahab justified by works}
  \item[(3)] \(\pvar{TestimonyArguments.SolaFide}{86}\)\gloss{James’s ‘justified’ (δικαιόω) means the increase of the justice received}
  \item[(4)] \((\pvar{TestimonyArguments.SolaFide}{23} \land \pvar{TestimonyArguments.SolaFide}{84} \land \pvar{TestimonyArguments.SolaFide}{86}) \rightarrow \pvar{TestimonyArguments.SolaFide}{90}\)
  \item[\(\vdash\)] \(\pvar{TestimonyArguments.SolaFide}{90}\)\gloss{\emph{at issue:} Scripture teaches an increase of justification by good works}
\end{derivation}
\par\medskip\noindent\textbf{\texttt{catholic\allowbreak{}James\allowbreak{}Works\allowbreak{}Line}} --- James 2:24, read by canon 24\par\nobreak\noindent
\textbf{Trent’s reading of what works do}, applied to James 2:24: the works by which a person is justified are a cause of the increase, so they are not merely fruit and evidence.

\begin{derivation}
  \item[(1)] \(\pvar{TestimonyArguments.SolaFide}{23}\)\gloss{James 2:24 says a person is justified by works and not by faith alone}
  \item[(2)] \(\pvar{TestimonyArguments.SolaFide}{84}\)\gloss{James 2:21–25: Abraham justified by works at Genesis 22, fulfilling Genesis 15:6; Rahab justified by works}
  \item[(3)] \(\pvar{TestimonyArguments.SolaFide}{91}\)\gloss{Good works are a cause of the increase of justification, not merely its fruits}
  \item[(4)] \((\pvar{TestimonyArguments.SolaFide}{23} \land \pvar{TestimonyArguments.SolaFide}{84} \land \pvar{TestimonyArguments.SolaFide}{91}) \rightarrow \lnot \pvar{TestimonyArguments.SolaFide}{21}\)
  \item[\(\vdash\)] \(\lnot \pvar{TestimonyArguments.SolaFide}{21}\)\gloss{\emph{at issue, the denial of:} Good works are the fruit and evidence of saving faith, not its ground}
\end{derivation}
\par\medskip\noindent\textbf{\texttt{trent\allowbreak{}Revelation\allowbreak{}Line}} --- Revelation 22:11, as Trent quotes it\par\nobreak\noindent
\textbf{Trent’s Revelation 22:11}: “He that is just, let him be justified still”.

\begin{derivation}
  \item[(1)] \(\pvar{TestimonyArguments.SolaFide}{92}\)\gloss{Revelation 22:11 says: he that is just, let him be justified still}
  \item[(2)] \(\pvar{TestimonyArguments.SolaFide}{92} \rightarrow \pvar{TestimonyArguments.SolaFide}{90}\)
  \item[\(\vdash\)] \(\pvar{TestimonyArguments.SolaFide}{90}\)\gloss{\emph{at issue:} Scripture teaches an increase of justification by good works}
\end{derivation}
\par\medskip\noindent\textbf{\texttt{trent\allowbreak{}Sirach\allowbreak{}Line}} --- Sirach 18:22, as Trent quotes it\par\nobreak\noindent
\textbf{Trent’s Sirach 18:22}: “Be not afraid to be justified even to death”.

\begin{derivation}
  \item[(1)] \(\pvar{TestimonyArguments.SolaFide}{94}\)\gloss{Sirach 18:22 says: be not afraid to be justified even to death}
  \item[(2)] \(\pvar{TestimonyArguments.SolaFide}{94} \rightarrow \pvar{TestimonyArguments.SolaFide}{90}\)
  \item[\(\vdash\)] \(\pvar{TestimonyArguments.SolaFide}{90}\)\gloss{\emph{at issue:} Scripture teaches an increase of justification by good works}
\end{derivation}
\par\medskip\noindent\textbf{\texttt{revelation\allowbreak{}Text\allowbreak{}Line}} --- Revelation 22:11, in the Greek\par\nobreak\noindent
\textbf{Revelation 22:11 in Greek.} The verse says “let the righteous still do righteousness”. The original outweighs any translation. So its inspired text does not say “be justified still”.

\begin{derivation}
  \item[(1)] \(\pvar{TestimonyArguments.SolaFide}{93}\)\gloss{The Greek of Revelation 22:11 says: let the righteous still do righteousness}
  \item[(2)] \(\pvar{TestimonyArguments.SolaFide}{96}\)\gloss{The original text, written by the inspired author, outweighs any translation: a reading it lacks is not his word}
  \item[(3)] \((\pvar{TestimonyArguments.SolaFide}{93} \land \pvar{TestimonyArguments.SolaFide}{96}) \rightarrow \lnot \pvar{TestimonyArguments.SolaFide}{92}\)
  \item[\(\vdash\)] \(\lnot \pvar{TestimonyArguments.SolaFide}{92}\)\gloss{\emph{at issue, the denial of:} Revelation 22:11 says: he that is just, let him be justified still}
\end{derivation}
\par\medskip\noindent\textbf{\texttt{sirach\allowbreak{}Text\allowbreak{}Line}} --- Sirach 18:22, in the Greek\par\nobreak\noindent
\textbf{Sirach 18:22 in Greek.} The verse is about paying a vow promptly, and says “do not wait until death to be justified”. The original outweighs any translation. So its inspired text does not say what the Latin says, “be not afraid to be justified even to death”.

\begin{derivation}
  \item[(1)] \(\pvar{TestimonyArguments.SolaFide}{95}\)\gloss{The Greek of Sirach 18:22 is about paying a vow: do not wait until death to be justified}
  \item[(2)] \(\pvar{TestimonyArguments.SolaFide}{96}\)\gloss{The original text, written by the inspired author, outweighs any translation: a reading it lacks is not his word}
  \item[(3)] \((\pvar{TestimonyArguments.SolaFide}{95} \land \pvar{TestimonyArguments.SolaFide}{96}) \rightarrow \lnot \pvar{TestimonyArguments.SolaFide}{94}\)
  \item[\(\vdash\)] \(\lnot \pvar{TestimonyArguments.SolaFide}{94}\)\gloss{\emph{at issue, the denial of:} Sirach 18:22 says: be not afraid to be justified even to death}
\end{derivation}
\par\medskip\noindent\textbf{\texttt{trent\allowbreak{}Reward\allowbreak{}Line}} --- The reward of works (Trent, Session VI, ch. 16)\par\nobreak\noindent
\textbf{The reward of works, as Trent reads it.} God rewards the labour of believers (1 Corinthians 15:58, Hebrews 6:10, Hebrews 10:35, 2 Timothy 4:8). The reward is given to those works as merits (chapter 16, canon 32). So works are not merely the fruit of justification. The second premise is stated as Trent states it. The texts say “reward”; that the reward is given to merits is the council’s own claim, introduced with “we must believe”.

\begin{derivation}
  \item[(1)] \(\pvar{TestimonyArguments.SolaFide}{105}\)\gloss{1 Cor 15:58; Heb 6:10; 10:35; 2 Tim 4:8: God rewards the labor of believers}
  \item[(2)] \(\pvar{TestimonyArguments.SolaFide}{107}\)\gloss{The reward God renders to the works of the justified is rendered to them as merits}
  \item[(3)] \((\pvar{TestimonyArguments.SolaFide}{105} \land \pvar{TestimonyArguments.SolaFide}{107}) \rightarrow \lnot \pvar{TestimonyArguments.SolaFide}{21}\)
  \item[\(\vdash\)] \(\lnot \pvar{TestimonyArguments.SolaFide}{21}\)\gloss{\emph{at issue, the denial of:} Good works are the fruit and evidence of saving faith, not its ground}
\end{derivation}
\par\medskip\noindent\textbf{\texttt{renewal\allowbreak{}Word\allowbreak{}Line}} --- James’s word, read through Paul’s (the renewal reading)\par\nobreak\noindent
\textbf{James’s word, read through Paul’s.} Paul’s “justify” (δικαιόω) means the renewal of the inward person, as Augustine explained it and as VanLandingham argues from the lexicon. James uses the word as Paul does. And James says Abraham was justified by works at Genesis 22, after Genesis 15:6 had been spoken. So James’s “justified” means the increase of the justice received.

This is the strongest reading of the word against the argument from James’s usage that has been found: it argues for the increase from Paul’s word, and does not merely state it.

\begin{derivation}
  \item[(1)] \(\pvar{TestimonyArguments.SolaFide}{52}\)\gloss{Paul’s word ‘justify’ means the renewal of the inward man, not only a verdict}
  \item[(2)] \(\pvar{TestimonyArguments.SolaFide}{106}\)\gloss{James’s ‘justify’ (δικαιόω) has the sense Paul’s has: the justification of Romans and Galatians}
  \item[(3)] \(\pvar{TestimonyArguments.SolaFide}{84}\)\gloss{James 2:21–25: Abraham justified by works at Genesis 22, fulfilling Genesis 15:6; Rahab justified by works}
  \item[(4)] \((\pvar{TestimonyArguments.SolaFide}{52} \land \pvar{TestimonyArguments.SolaFide}{106} \land \pvar{TestimonyArguments.SolaFide}{84}) \rightarrow \pvar{TestimonyArguments.SolaFide}{86}\)
  \item[\(\vdash\)] \(\pvar{TestimonyArguments.SolaFide}{86}\)\gloss{\emph{at issue:} James’s ‘justified’ (δικαιόω) means the increase of the justice received}
\end{derivation}
\par\medskip\noindent\textbf{\texttt{westminster\allowbreak{}Reward\allowbreak{}Line}} --- The reward of works, of grace (Westminster XVI.5, XVI.6)\par\nobreak\noindent
\textbf{The reward of works, as Westminster reads it} (XVI.5, XVI.6). God rewards the labour of believers, Hebrews 6:10 among the texts. Having done all, they are unworthy servants who have done only their duty (Luke 17:10). Wages are counted as due to one who works, and as a gift to one who believes (Romans 4:4–5). And their sacrifices are acceptable to God through Jesus Christ (1 Peter 2:5). So the reward is not given to their works as merits. Every premise is a text, and Westminster gives each as a proof of these sections. What is Westminster’s own is the step, a reading of the texts, rated \texttt{disputed}.

\begin{derivation}
  \item[(1)] \(\pvar{TestimonyArguments.SolaFide}{105}\)\gloss{1 Cor 15:58; Heb 6:10; 10:35; 2 Tim 4:8: God rewards the labor of believers}
  \item[(2)] \(\pvar{TestimonyArguments.SolaFide}{104}\)\gloss{Luke 17:10: having done all that was commanded, we are unworthy servants}
  \item[(3)] \(\pvar{TestimonyArguments.SolaFide}{4}\)\gloss{Romans 4:4–5 contrasts wages owed with a gift reckoned to the one who believes}
  \breakrow{(4)} \(\pvar{TestimonyArguments.SolaFide}{108}\)\gloss{1 Peter 2:5: spiritual sacrifices acceptable to God through Jesus Christ}
  \item[(5)] \((\pvar{TestimonyArguments.SolaFide}{105} \land \pvar{TestimonyArguments.SolaFide}{104} \land \pvar{TestimonyArguments.SolaFide}{4} \land \pvar{TestimonyArguments.SolaFide}{108}) \rightarrow \lnot \pvar{TestimonyArguments.SolaFide}{107}\)
  \item[\(\vdash\)] \(\lnot \pvar{TestimonyArguments.SolaFide}{107}\)\gloss{\emph{at issue, the denial of:} The reward God renders to the works of the justified is rendered to them as merits}
\end{derivation}
\paragraph*{The positions}
\leavevmode\par\nobreak
\par\medskip\noindent\textbf{\texttt{james\allowbreak{}Lexical\allowbreak{}Case}} --- James’s word (δικαιόω, by the rule of least meaning)\par\nobreak\noindent
James’s word: “justified” in James does not mean the increase of justice.

\begin{derivation}
  \item[(1)] \(\pvar{TestimonyArguments.SolaFide}{85}\)\gloss{Outside Paul, ‘justify’ (δικαιόω) declares or vindicates; it does not mean making righteous}
  \item[(2)] \(\pvar{TestimonyArguments.SolaFide}{84}\)\gloss{James 2:21–25: Abraham justified by works at Genesis 22, fulfilling Genesis 15:6; Rahab justified by works}
  \item[(3)] \(\pvar{TestimonyArguments.SolaFide}{51}\)\gloss{A word contributes the least meaning its context requires; one occurrence does not carry the whole concept}
  \item[(4)] \((\pvar{TestimonyArguments.SolaFide}{85} \land \pvar{TestimonyArguments.SolaFide}{84} \land \pvar{TestimonyArguments.SolaFide}{51}) \rightarrow \lnot \pvar{TestimonyArguments.SolaFide}{86}\)
  \item[\(\vdash\)] \(\lnot \pvar{TestimonyArguments.SolaFide}{86}\)\gloss{\emph{at issue, the denial of:} James’s ‘justified’ (δικαιόω) means the increase of the justice received}
  \keepwithnext
\end{derivation}
\emph{Grounded in scripture alone:} James 2:21–25: Abraham justified by works at Genesis 22, fulfilling Genesis 15:6; Rahab justified by works.
\par\medskip\noindent\textbf{\texttt{james\allowbreak{}Faith\allowbreak{}Case}} --- James’s faith (the faith James calls dead)\par\nobreak\noindent
James’s faith: the phrase James denies is not the Reformers’ formula.

\begin{derivation}
  \item[(1)] \(\pvar{TestimonyArguments.SolaFide}{83}\)\gloss{James 2:14–26: faith without works cannot save, and is dead (2:17, 2:26)}
  \item[(2)] \(\pvar{TestimonyArguments.SolaFide}{87}\)\gloss{The faith alone that James 2:24 denies justifies is the dead faith of 2:14–26}
  \item[(3)] \(\pvar{TestimonyArguments.SolaFide}{88}\)\gloss{The faith the Reformers say alone justifies is no dead faith, but works by love}
  \item[(4)] \((\pvar{TestimonyArguments.SolaFide}{83} \land \pvar{TestimonyArguments.SolaFide}{87} \land \pvar{TestimonyArguments.SolaFide}{88}) \rightarrow \pvar{TestimonyArguments.SolaFide}{89}\)
  \item[\(\vdash\)] \(\pvar{TestimonyArguments.SolaFide}{89}\)\gloss{\emph{at issue:} The faith James 2:24 denies justifies is not the faith of the Reformers’ formula}
  \keepwithnext
\end{derivation}
\emph{No premise here rests on scripture alone.}
\par\medskip\noindent\textbf{\texttt{fruit\allowbreak{}Source}}\par\nobreak\noindent
Westminster’s reading of its proof texts: from Ephesians 2:8–10, James 2:18–22, Romans 6:22 and Luke 17:10 to “works are the fruit and evidence of saving faith, not its ground” (XVI.2, XVI.5).

This rates how well the step is \textbf{supported}, and it is \texttt{well\allowbreak{}Supported}. The Reformed confess it. The Lutheran World Federation and the Catholic Church confess two parts of it together: good works “follow justification and are its fruits” (§37), and “whatever in the justified precedes or follows the free gift of faith is neither the basis of justification nor merits it” (§25, \emph{Joint Declaration}).

They do not confess that works are \emph{only} fruit. The Declaration records the Catholic understanding that good works “contribute to growth in grace, so that the righteousness that comes from God is preserved and communion with Christ is deepened” (§38). That is the point canon 24 makes, and it is weighed as Trent’s objection, not as part of this step’s support.

Whether an objection brings the step down is worked out separately (the chapter “Why each step is rated as it is”). Trent objects twice: in canon 24, and in chapter 16’s reading of the reward texts. When a tradition’s teaching counts, both stand, and the step is \texttt{disputed}. When only checkable evidence counts, neither does, because each rests at one claim on the council’s word.

{\small\begin{verbatim}
def fruitSource : Source :=
  {
    primary :=
      Reference.work westminsterConfession
        (Bib.Locus.sectionRef "XVI.2, XVI.5"),
    supporting :=
      [Reference.work jointDeclarationJustification
          (Bib.Locus.sectionRef "§§25, 37; see §38")],
    tradition := Tradition.reformedProtestant,
    confidence := Confidence.wellSupported }
\end{verbatim}}
\par\medskip\noindent\textbf{\texttt{fruit\allowbreak{}Line}} --- Works as fruit (Westminster XVI.2, XVI.5, and its proofs)\par\nobreak\noindent
\textbf{Works as fruit, from Westminster’s proof texts.} “Not a result of works”, and “created in Christ Jesus for good works” (Ephesians 2:8–10). Faith shown by works (James 2:18, 22). “The fruit you get leads to sanctification” (Romans 6:22). “We are unworthy servants; we have only done what was our duty” (Luke 17:10). So works are the fruit and evidence of saving faith, not its ground.

\begin{derivation}
  \item[(1)] \(\pvar{TestimonyArguments.SolaFide}{1}\)\gloss{Ephesians 2:8–9 teaches salvation by grace through faith, not of works}
  \item[(2)] \(\pvar{TestimonyArguments.SolaFide}{101}\)\gloss{Ephesians 2:10: created in Christ Jesus for good works, after ‘not a result of works’}
  \item[(3)] \(\pvar{TestimonyArguments.SolaFide}{102}\)\gloss{James 2:18, 22: I will show you my faith by my works; faith completed by works}
  \breakrow{(4)} \(\pvar{TestimonyArguments.SolaFide}{103}\)\gloss{Romans 6:22: the fruit you get leads to sanctification, and its end is eternal life}
  \breakrow{(5)} \(\pvar{TestimonyArguments.SolaFide}{104}\)\gloss{Luke 17:10: having done all that was commanded, we are unworthy servants}
  \item[(6)] \((\pvar{TestimonyArguments.SolaFide}{1} \land \pvar{TestimonyArguments.SolaFide}{101} \land \pvar{TestimonyArguments.SolaFide}{102} \land \pvar{TestimonyArguments.SolaFide}{103} \land \pvar{TestimonyArguments.SolaFide}{104}) \rightarrow \pvar{TestimonyArguments.SolaFide}{21}\)
  \item[\(\vdash\)] \(\pvar{TestimonyArguments.SolaFide}{21}\)\gloss{\emph{at issue:} Good works are the fruit and evidence of saving faith, not its ground}
\end{derivation}
\par\medskip\noindent\textbf{\texttt{james\allowbreak{}Harmony\allowbreak{}Case}} --- Harmonisation of James 2:24\par\nobreak\noindent
The Reformed reconciliation of James with Paul: James is attacking a barren faith, and works are the fruit of saving faith and not its ground. The second point is derived from Westminster’s proof texts, not assumed. The line here is \texttt{james\allowbreak{}Line} with “fruit, not ground” taken from \texttt{fruit\allowbreak{}Line}.

\begin{derivation}
  \item[(1)] \(\pvar{TestimonyArguments.SolaFide}{1}\)\gloss{Ephesians 2:8–9 teaches salvation by grace through faith, not of works}
  \item[(2)] \(\pvar{TestimonyArguments.SolaFide}{101}\)\gloss{Ephesians 2:10: created in Christ Jesus for good works, after ‘not a result of works’}
  \item[(3)] \(\pvar{TestimonyArguments.SolaFide}{102}\)\gloss{James 2:18, 22: I will show you my faith by my works; faith completed by works}
  \breakrow{(4)} \(\pvar{TestimonyArguments.SolaFide}{103}\)\gloss{Romans 6:22: the fruit you get leads to sanctification, and its end is eternal life}
  \breakrow{(5)} \(\pvar{TestimonyArguments.SolaFide}{104}\)\gloss{Luke 17:10: having done all that was commanded, we are unworthy servants}
  \breakrow{(6)} \(\pvar{TestimonyArguments.SolaFide}{20}\)\gloss{James 2:14–26 targets a barren faith — mere assent — not Paul’s doctrine}
  \breakrow{(7)} \((\pvar{TestimonyArguments.SolaFide}{1} \land \pvar{TestimonyArguments.SolaFide}{101} \land \pvar{TestimonyArguments.SolaFide}{102} \land \pvar{TestimonyArguments.SolaFide}{103} \land \pvar{TestimonyArguments.SolaFide}{104}) \rightarrow \pvar{TestimonyArguments.SolaFide}{21}\)\gloss{\emph{would give:} Good works are the fruit and evidence of saving faith, not its ground}
  \item[(8)] \((\pvar{TestimonyArguments.SolaFide}{20} \land \pvar{TestimonyArguments.SolaFide}{21}) \rightarrow \pvar{TestimonyArguments.SolaFide}{22}\)
  \item[\(\vdash\)] \(\pvar{TestimonyArguments.SolaFide}{22}\)\gloss{\emph{at issue:} James 2:24 is compatible with Paul, using ‘justify’ in a different sense}
  \keepwithnext
\end{derivation}
\emph{No premise here rests on scripture alone.}
\par\medskip\noindent\textbf{\texttt{catholic\allowbreak{}James\allowbreak{}On\allowbreak{}The\allowbreak{}Word}} --- James 2:24, read for the increase of justification\par\nobreak\noindent
Trent on James’s word: James 2:24 teaches the increase of justification, because that is what James’s “justified” means.

\begin{derivation}
  \item[(1)] \(\pvar{TestimonyArguments.SolaFide}{23}\)\gloss{James 2:24 says a person is justified by works and not by faith alone}
  \item[(2)] \(\pvar{TestimonyArguments.SolaFide}{84}\)\gloss{James 2:21–25: Abraham justified by works at Genesis 22, fulfilling Genesis 15:6; Rahab justified by works}
  \item[(3)] \(\pvar{TestimonyArguments.SolaFide}{86}\)\gloss{James’s ‘justified’ (δικαιόω) means the increase of the justice received}
  \item[(4)] \((\pvar{TestimonyArguments.SolaFide}{23} \land \pvar{TestimonyArguments.SolaFide}{84} \land \pvar{TestimonyArguments.SolaFide}{86}) \rightarrow \pvar{TestimonyArguments.SolaFide}{90}\)
  \item[\(\vdash\)] \(\pvar{TestimonyArguments.SolaFide}{90}\)\gloss{\emph{at issue:} Scripture teaches an increase of justification by good works}
  \keepwithnext
\end{derivation}
\emph{Grounded in scripture alone:} James 2:21–25: Abraham justified by works at Genesis 22, fulfilling Genesis 15:6; Rahab justified by works.
\par\medskip\noindent\textbf{\texttt{catholic\allowbreak{}James\allowbreak{}On\allowbreak{}Works}} --- James 2:24, read by canon 24\par\nobreak\noindent
Trent on what works do: the works of James 2:24 are a cause of the increase of justification, not merely its fruit.

\begin{derivation}
  \item[(1)] \(\pvar{TestimonyArguments.SolaFide}{23}\)\gloss{James 2:24 says a person is justified by works and not by faith alone}
  \item[(2)] \(\pvar{TestimonyArguments.SolaFide}{84}\)\gloss{James 2:21–25: Abraham justified by works at Genesis 22, fulfilling Genesis 15:6; Rahab justified by works}
  \item[(3)] \(\pvar{TestimonyArguments.SolaFide}{91}\)\gloss{Good works are a cause of the increase of justification, not merely its fruits}
  \item[(4)] \((\pvar{TestimonyArguments.SolaFide}{23} \land \pvar{TestimonyArguments.SolaFide}{84} \land \pvar{TestimonyArguments.SolaFide}{91}) \rightarrow \lnot \pvar{TestimonyArguments.SolaFide}{21}\)
  \item[\(\vdash\)] \(\lnot \pvar{TestimonyArguments.SolaFide}{21}\)\gloss{\emph{at issue, the denial of:} Good works are the fruit and evidence of saving faith, not its ground}
  \keepwithnext
\end{derivation}
\emph{Grounded in scripture alone:} James 2:21–25: Abraham justified by works at Genesis 22, fulfilling Genesis 15:6; Rahab justified by works.
\par\medskip\noindent\textbf{\texttt{trent\allowbreak{}On\allowbreak{}Revelation}} --- Revelation 22:11, as Trent quotes it\par\nobreak\noindent
Trent on Revelation 22:11.

\begin{derivation}
  \item[(1)] \(\pvar{TestimonyArguments.SolaFide}{92}\)\gloss{Revelation 22:11 says: he that is just, let him be justified still}
  \item[(2)] \(\pvar{TestimonyArguments.SolaFide}{92} \rightarrow \pvar{TestimonyArguments.SolaFide}{90}\)
  \item[\(\vdash\)] \(\pvar{TestimonyArguments.SolaFide}{90}\)\gloss{\emph{at issue:} Scripture teaches an increase of justification by good works}
  \keepwithnext
\end{derivation}
\emph{No premise here rests on scripture alone.}
\par\medskip\noindent\textbf{\texttt{trent\allowbreak{}On\allowbreak{}Sirach}} --- Sirach 18:22, as Trent quotes it\par\nobreak\noindent
Trent on Sirach 18:22.

\begin{derivation}
  \item[(1)] \(\pvar{TestimonyArguments.SolaFide}{94}\)\gloss{Sirach 18:22 says: be not afraid to be justified even to death}
  \item[(2)] \(\pvar{TestimonyArguments.SolaFide}{94} \rightarrow \pvar{TestimonyArguments.SolaFide}{90}\)
  \item[\(\vdash\)] \(\pvar{TestimonyArguments.SolaFide}{90}\)\gloss{\emph{at issue:} Scripture teaches an increase of justification by good works}
  \keepwithnext
\end{derivation}
\emph{No premise here rests on scripture alone.}
\par\medskip\noindent\textbf{\texttt{revelation\allowbreak{}Text\allowbreak{}Case}} --- Revelation 22:11, in the Greek\par\nobreak\noindent
What Revelation 22:11 says, in the Greek.

\begin{derivation}
  \item[(1)] \(\pvar{TestimonyArguments.SolaFide}{93}\)\gloss{The Greek of Revelation 22:11 says: let the righteous still do righteousness}
  \item[(2)] \(\pvar{TestimonyArguments.SolaFide}{96}\)\gloss{The original text, written by the inspired author, outweighs any translation: a reading it lacks is not his word}
  \item[(3)] \((\pvar{TestimonyArguments.SolaFide}{93} \land \pvar{TestimonyArguments.SolaFide}{96}) \rightarrow \lnot \pvar{TestimonyArguments.SolaFide}{92}\)
  \item[\(\vdash\)] \(\lnot \pvar{TestimonyArguments.SolaFide}{92}\)\gloss{\emph{at issue, the denial of:} Revelation 22:11 says: he that is just, let him be justified still}
  \keepwithnext
\end{derivation}
\emph{No premise here rests on scripture alone.}
\par\medskip\noindent\textbf{\texttt{sirach\allowbreak{}Text\allowbreak{}Case}} --- Sirach 18:22, in the Greek\par\nobreak\noindent
What Sirach 18:22 says, in the Greek.

\begin{derivation}
  \item[(1)] \(\pvar{TestimonyArguments.SolaFide}{95}\)\gloss{The Greek of Sirach 18:22 is about paying a vow: do not wait until death to be justified}
  \item[(2)] \(\pvar{TestimonyArguments.SolaFide}{96}\)\gloss{The original text, written by the inspired author, outweighs any translation: a reading it lacks is not his word}
  \item[(3)] \((\pvar{TestimonyArguments.SolaFide}{95} \land \pvar{TestimonyArguments.SolaFide}{96}) \rightarrow \lnot \pvar{TestimonyArguments.SolaFide}{94}\)
  \item[\(\vdash\)] \(\lnot \pvar{TestimonyArguments.SolaFide}{94}\)\gloss{\emph{at issue, the denial of:} Sirach 18:22 says: be not afraid to be justified even to death}
  \keepwithnext
\end{derivation}
\emph{No premise here rests on scripture alone.}
\par\medskip\noindent\textbf{\texttt{trent\allowbreak{}On\allowbreak{}Reward}} --- The reward of works (Trent, Session VI, ch. 16)\par\nobreak\noindent
Trent on the reward of works: what God rewards is not merely fruit.

\begin{derivation}
  \item[(1)] \(\pvar{TestimonyArguments.SolaFide}{105}\)\gloss{1 Cor 15:58; Heb 6:10; 10:35; 2 Tim 4:8: God rewards the labor of believers}
  \item[(2)] \(\pvar{TestimonyArguments.SolaFide}{107}\)\gloss{The reward God renders to the works of the justified is rendered to them as merits}
  \item[(3)] \((\pvar{TestimonyArguments.SolaFide}{105} \land \pvar{TestimonyArguments.SolaFide}{107}) \rightarrow \lnot \pvar{TestimonyArguments.SolaFide}{21}\)
  \item[\(\vdash\)] \(\lnot \pvar{TestimonyArguments.SolaFide}{21}\)\gloss{\emph{at issue, the denial of:} Good works are the fruit and evidence of saving faith, not its ground}
  \keepwithnext
\end{derivation}
\emph{No premise here rests on scripture alone.}
\par\medskip\noindent\textbf{\texttt{westminster\allowbreak{}On\allowbreak{}Reward}} --- The reward of works, of grace (Westminster XVI.5, XVI.6)\par\nobreak\noindent
Westminster on the reward of works: the reward is of grace, and is not given to merits.

\begin{derivation}
  \item[(1)] \(\pvar{TestimonyArguments.SolaFide}{105}\)\gloss{1 Cor 15:58; Heb 6:10; 10:35; 2 Tim 4:8: God rewards the labor of believers}
  \item[(2)] \(\pvar{TestimonyArguments.SolaFide}{104}\)\gloss{Luke 17:10: having done all that was commanded, we are unworthy servants}
  \item[(3)] \(\pvar{TestimonyArguments.SolaFide}{4}\)\gloss{Romans 4:4–5 contrasts wages owed with a gift reckoned to the one who believes}
  \breakrow{(4)} \(\pvar{TestimonyArguments.SolaFide}{108}\)\gloss{1 Peter 2:5: spiritual sacrifices acceptable to God through Jesus Christ}
  \item[(5)] \((\pvar{TestimonyArguments.SolaFide}{105} \land \pvar{TestimonyArguments.SolaFide}{104} \land \pvar{TestimonyArguments.SolaFide}{4} \land \pvar{TestimonyArguments.SolaFide}{108}) \rightarrow \lnot \pvar{TestimonyArguments.SolaFide}{107}\)
  \item[\(\vdash\)] \(\lnot \pvar{TestimonyArguments.SolaFide}{107}\)\gloss{\emph{at issue, the denial of:} The reward God renders to the works of the justified is rendered to them as merits}
  \keepwithnext
\end{derivation}
\emph{No premise here rests on scripture alone.}
\par\medskip\noindent\textbf{\texttt{renewal\allowbreak{}On\allowbreak{}The\allowbreak{}Word}} --- James’s word, read through Paul’s (the renewal reading)\par\nobreak\noindent
The reading through Paul: James’s “justified” means the increase of justice, because Paul’s word means renewal and James’s is Paul’s.

\begin{derivation}
  \item[(1)] \(\pvar{TestimonyArguments.SolaFide}{52}\)\gloss{Paul’s word ‘justify’ means the renewal of the inward man, not only a verdict}
  \item[(2)] \(\pvar{TestimonyArguments.SolaFide}{106}\)\gloss{James’s ‘justify’ (δικαιόω) has the sense Paul’s has: the justification of Romans and Galatians}
  \item[(3)] \(\pvar{TestimonyArguments.SolaFide}{84}\)\gloss{James 2:21–25: Abraham justified by works at Genesis 22, fulfilling Genesis 15:6; Rahab justified by works}
  \item[(4)] \((\pvar{TestimonyArguments.SolaFide}{52} \land \pvar{TestimonyArguments.SolaFide}{106} \land \pvar{TestimonyArguments.SolaFide}{84}) \rightarrow \pvar{TestimonyArguments.SolaFide}{86}\)
  \item[\(\vdash\)] \(\pvar{TestimonyArguments.SolaFide}{86}\)\gloss{\emph{at issue:} James’s ‘justified’ (δικαιόω) means the increase of the justice received}
  \keepwithnext
\end{derivation}
\emph{Grounded in scripture alone:} James 2:21–25: Abraham justified by works at Genesis 22, fulfilling Genesis 15:6; Rahab justified by works.
\paragraph*{Each position holds together and reaches its conclusion}
\leavevmode\par\nobreak
\par\medskip\noindent\textbf{\texttt{reformed\allowbreak{}James\allowbreak{}Reading}}\par\nobreak\noindent
\textbf{The Reformed reading of James.} Every text as it stands. James’s “justified” a declaration and not an increase. The faith James denies not the Reformers’ faith. Trent’s Latin readings of Revelation 22:11 and Sirach 18:22 not what the verses say. And the reward of works not given to them as merits.

{\small\begin{verbatim}
def reformedJamesReading : Valuation Claim :=
  fun a =>
    match a with
    | Claim.jamesJustifyDenotesIncrease => False
    | Claim.worksCauseIncreaseOfJustification => False
    | Claim.rev22_11BeJustifiedStill => False
    | Claim.sir18_22BeJustifiedToDeath => False
    | Claim.scriptureTeachesIncreaseOfJustification => False
    | Claim.rewardRenderedToMerits => False
    | x => True
\end{verbatim}}
\par\medskip\noindent\textbf{\texttt{tridentine\allowbreak{}James\allowbreak{}Reading}}\par\nobreak\noindent
\textbf{Trent’s reading of James.} Every text as Trent reads it: James’s “justified” the increase of justice, its works a cause of that increase and not merely fruit, and Revelation 22:11 and Sirach 18:22 as the Latin has them.

{\small\begin{verbatim}
def tridentineJamesReading : Valuation Claim :=
  fun a =>
    match a with
    | Claim.worksAreFruitNotGround => False
    | x => True
\end{verbatim}}
\par\medskip\noindent\textbf{\texttt{james\allowbreak{}Lexical\allowbreak{}Case\_\allowbreak{}establishes}}\par\nobreak\noindent
The argument from James’s word reaches its conclusion.

{\small\begin{verbatim}
theorem jamesLexicalCase_establishes : Establishes jamesLexicalCase
-- axioms: propext, Quot.sound
\end{verbatim}}
\par\medskip\noindent\textbf{\texttt{james\allowbreak{}Lexical\allowbreak{}Case\_\allowbreak{}is\_\allowbreak{}satisfiable}}\par\nobreak\noindent
The argument from James’s word can be held without contradiction.

{\small\begin{verbatim}
theorem jamesLexicalCase_is_satisfiable : Satisfiable
    jamesLexicalCase.premises
-- axioms: propext, Quot.sound
\end{verbatim}}
\par\medskip\noindent\textbf{\texttt{james\allowbreak{}Faith\allowbreak{}Case\_\allowbreak{}establishes}}\par\nobreak\noindent
The argument about James’s faith reaches its conclusion.

{\small\begin{verbatim}
theorem jamesFaithCase_establishes : Establishes jamesFaithCase
-- axioms: propext, Quot.sound
\end{verbatim}}
\par\medskip\noindent\textbf{\texttt{james\allowbreak{}Faith\allowbreak{}Case\_\allowbreak{}is\_\allowbreak{}satisfiable}}\par\nobreak\noindent
The argument about James’s faith can be held without contradiction.

{\small\begin{verbatim}
theorem jamesFaithCase_is_satisfiable : Satisfiable jamesFaithCase.premises
-- axioms: propext, Quot.sound
\end{verbatim}}
\par\medskip\noindent\textbf{\texttt{james\allowbreak{}Harmony\allowbreak{}Case\_\allowbreak{}establishes}}\par\nobreak\noindent
The Reformed reconciliation reaches its conclusion.

{\small\begin{verbatim}
theorem jamesHarmonyCase_establishes : Establishes jamesHarmonyCase
-- axioms: propext, Quot.sound
\end{verbatim}}
\par\medskip\noindent\textbf{\texttt{james\allowbreak{}Harmony\allowbreak{}Case\_\allowbreak{}is\_\allowbreak{}satisfiable}}\par\nobreak\noindent
The Reformed reconciliation can be held without contradiction.

{\small\begin{verbatim}
theorem jamesHarmonyCase_is_satisfiable : Satisfiable
    jamesHarmonyCase.premises
-- axioms: propext, Quot.sound
\end{verbatim}}
\par\medskip\noindent\textbf{\texttt{catholic\allowbreak{}James\allowbreak{}On\allowbreak{}The\allowbreak{}Word\_\allowbreak{}establishes}}\par\nobreak\noindent
Trent on James’s word reaches its conclusion.

{\small\begin{verbatim}
theorem catholicJamesOnTheWord_establishes : Establishes
    catholicJamesOnTheWord
-- axioms: propext, Quot.sound
\end{verbatim}}
\par\medskip\noindent\textbf{\texttt{catholic\allowbreak{}James\allowbreak{}On\allowbreak{}The\allowbreak{}Word\_\allowbreak{}is\_\allowbreak{}satisfiable}}\par\nobreak\noindent
Trent on James’s word can be held without contradiction.

{\small\begin{verbatim}
theorem catholicJamesOnTheWord_is_satisfiable : Satisfiable
    catholicJamesOnTheWord.premises
-- axioms: propext, Quot.sound
\end{verbatim}}
\par\medskip\noindent\textbf{\texttt{catholic\allowbreak{}James\allowbreak{}On\allowbreak{}Works\_\allowbreak{}establishes}}\par\nobreak\noindent
Trent on what works do reaches its conclusion.

{\small\begin{verbatim}
theorem catholicJamesOnWorks_establishes : Establishes catholicJamesOnWorks
-- axioms: propext, Quot.sound
\end{verbatim}}
\par\medskip\noindent\textbf{\texttt{catholic\allowbreak{}James\allowbreak{}On\allowbreak{}Works\_\allowbreak{}is\_\allowbreak{}satisfiable}}\par\nobreak\noindent
Trent on what works do can be held without contradiction.

{\small\begin{verbatim}
theorem catholicJamesOnWorks_is_satisfiable : Satisfiable
    catholicJamesOnWorks.premises
-- axioms: propext, Quot.sound
\end{verbatim}}
\par\medskip\noindent\textbf{\texttt{trent\allowbreak{}On\allowbreak{}Revelation\_\allowbreak{}establishes}}\par\nobreak\noindent
Trent on Revelation 22:11 reaches its conclusion.

{\small\begin{verbatim}
theorem trentOnRevelation_establishes : Establishes trentOnRevelation
-- axioms: propext, Quot.sound
\end{verbatim}}
\par\medskip\noindent\textbf{\texttt{trent\allowbreak{}On\allowbreak{}Revelation\_\allowbreak{}is\_\allowbreak{}satisfiable}}\par\nobreak\noindent
Trent on Revelation 22:11 can be held without contradiction.

{\small\begin{verbatim}
theorem trentOnRevelation_is_satisfiable : Satisfiable
    trentOnRevelation.premises
-- axioms: propext, Quot.sound
\end{verbatim}}
\par\medskip\noindent\textbf{\texttt{trent\allowbreak{}On\allowbreak{}Sirach\_\allowbreak{}establishes}}\par\nobreak\noindent
Trent on Sirach 18:22 reaches its conclusion.

{\small\begin{verbatim}
theorem trentOnSirach_establishes : Establishes trentOnSirach
-- axioms: propext, Quot.sound
\end{verbatim}}
\par\medskip\noindent\textbf{\texttt{trent\allowbreak{}On\allowbreak{}Sirach\_\allowbreak{}is\_\allowbreak{}satisfiable}}\par\nobreak\noindent
Trent on Sirach 18:22 can be held without contradiction.

{\small\begin{verbatim}
theorem trentOnSirach_is_satisfiable : Satisfiable trentOnSirach.premises
-- axioms: propext, Quot.sound
\end{verbatim}}
\par\medskip\noindent\textbf{\texttt{revelation\allowbreak{}Text\allowbreak{}Case\_\allowbreak{}establishes}}\par\nobreak\noindent
The Greek of Revelation 22:11 reaches its conclusion.

{\small\begin{verbatim}
theorem revelationTextCase_establishes : Establishes revelationTextCase
-- axioms: propext, Quot.sound
\end{verbatim}}
\par\medskip\noindent\textbf{\texttt{revelation\allowbreak{}Text\allowbreak{}Case\_\allowbreak{}is\_\allowbreak{}satisfiable}}\par\nobreak\noindent
The Greek of Revelation 22:11 can be held without contradiction.

{\small\begin{verbatim}
theorem revelationTextCase_is_satisfiable : Satisfiable
    revelationTextCase.premises
-- axioms: propext, Quot.sound
\end{verbatim}}
\par\medskip\noindent\textbf{\texttt{sirach\allowbreak{}Text\allowbreak{}Case\_\allowbreak{}establishes}}\par\nobreak\noindent
The Greek of Sirach 18:22 reaches its conclusion.

{\small\begin{verbatim}
theorem sirachTextCase_establishes : Establishes sirachTextCase
-- axioms: propext, Quot.sound
\end{verbatim}}
\par\medskip\noindent\textbf{\texttt{sirach\allowbreak{}Text\allowbreak{}Case\_\allowbreak{}is\_\allowbreak{}satisfiable}}\par\nobreak\noindent
The Greek of Sirach 18:22 can be held without contradiction.

{\small\begin{verbatim}
theorem sirachTextCase_is_satisfiable : Satisfiable sirachTextCase.premises
-- axioms: propext, Quot.sound
\end{verbatim}}
\par\medskip\noindent\textbf{\texttt{trent\allowbreak{}On\allowbreak{}Reward\_\allowbreak{}establishes}}\par\nobreak\noindent
Trent on the reward of works reaches its conclusion.

{\small\begin{verbatim}
theorem trentOnReward_establishes : Establishes trentOnReward
-- axioms: propext, Quot.sound
\end{verbatim}}
\par\medskip\noindent\textbf{\texttt{trent\allowbreak{}On\allowbreak{}Reward\_\allowbreak{}is\_\allowbreak{}satisfiable}}\par\nobreak\noindent
Trent on the reward of works can be held without contradiction.

{\small\begin{verbatim}
theorem trentOnReward_is_satisfiable : Satisfiable trentOnReward.premises
-- axioms: propext, Quot.sound
\end{verbatim}}
\par\medskip\noindent\textbf{\texttt{westminster\allowbreak{}On\allowbreak{}Reward\_\allowbreak{}establishes}}\par\nobreak\noindent
Westminster on the reward of works reaches its conclusion.

{\small\begin{verbatim}
theorem westminsterOnReward_establishes : Establishes westminsterOnReward
-- axioms: propext, Quot.sound
\end{verbatim}}
\par\medskip\noindent\textbf{\texttt{westminster\allowbreak{}On\allowbreak{}Reward\_\allowbreak{}is\_\allowbreak{}satisfiable}}\par\nobreak\noindent
Westminster on the reward of works can be held without contradiction.

{\small\begin{verbatim}
theorem westminsterOnReward_is_satisfiable : Satisfiable
    westminsterOnReward.premises
-- axioms: propext, Quot.sound
\end{verbatim}}
\par\medskip\noindent\textbf{\texttt{renewal\allowbreak{}On\allowbreak{}The\allowbreak{}Word\_\allowbreak{}establishes}}\par\nobreak\noindent
The reading through Paul reaches its conclusion.

{\small\begin{verbatim}
theorem renewalOnTheWord_establishes : Establishes renewalOnTheWord
-- axioms: propext, Quot.sound
\end{verbatim}}
\par\medskip\noindent\textbf{\texttt{renewal\allowbreak{}On\allowbreak{}The\allowbreak{}Word\_\allowbreak{}is\_\allowbreak{}satisfiable}}\par\nobreak\noindent
The reading through Paul can be held without contradiction.

{\small\begin{verbatim}
theorem renewalOnTheWord_is_satisfiable : Satisfiable
    renewalOnTheWord.premises
-- axioms: propext, Quot.sound
\end{verbatim}}
\paragraph*{How strong each position is}
\leavevmode\par\nobreak
\par\medskip\noindent\textbf{\texttt{james\allowbreak{}Lexical\allowbreak{}Case\_\allowbreak{}strength}}\par\nobreak\noindent
The argument from James’s word rests on nothing rated below \texttt{well\allowbreak{}Supported}.

{\small\begin{verbatim}
theorem jamesLexicalCase_strength : jamesLexicalCase.strength = 2
-- axioms: propext
\end{verbatim}}
\par\medskip\noindent\textbf{\texttt{james\allowbreak{}Faith\allowbreak{}Case\_\allowbreak{}strength}}\par\nobreak\noindent
The argument about James’s faith rests on nothing rated below \texttt{well\allowbreak{}Supported}.

{\small\begin{verbatim}
theorem jamesFaithCase_strength : jamesFaithCase.strength = 2
-- axioms: propext
\end{verbatim}}
\par\medskip\noindent\textbf{\texttt{james\allowbreak{}Harmony\allowbreak{}Case\_\allowbreak{}strength}}\par\nobreak\noindent
The Reformed reconciliation rests on its reading of what James is attacking, and on works as fruit and not ground. Both are \texttt{disputed}.

{\small\begin{verbatim}
theorem jamesHarmonyCase_strength : jamesHarmonyCase.strength = 0
-- axioms: propext
\end{verbatim}}
\par\medskip\noindent\textbf{\texttt{catholic\allowbreak{}James\allowbreak{}On\allowbreak{}The\allowbreak{}Word\_\allowbreak{}strength}}\par\nobreak\noindent
Trent on James’s word rests on its reading of the word, rated \texttt{disputed}.

{\small\begin{verbatim}
theorem catholicJamesOnTheWord_strength : catholicJamesOnTheWord.strength = 0
-- axioms: propext
\end{verbatim}}
\par\medskip\noindent\textbf{\texttt{catholic\allowbreak{}James\allowbreak{}On\allowbreak{}Works\_\allowbreak{}strength}}\par\nobreak\noindent
Trent on what works do rests on canon 24, rated \texttt{disputed}.

{\small\begin{verbatim}
theorem catholicJamesOnWorks_strength : catholicJamesOnWorks.strength = 0
-- axioms: propext
\end{verbatim}}
\par\medskip\noindent\textbf{\texttt{trent\allowbreak{}On\allowbreak{}Revelation\_\allowbreak{}strength}}\par\nobreak\noindent
Trent on Revelation 22:11 rests on the Latin reading, rated \texttt{disputed}.

{\small\begin{verbatim}
theorem trentOnRevelation_strength : trentOnRevelation.strength = 0
-- axioms: propext
\end{verbatim}}
\par\medskip\noindent\textbf{\texttt{trent\allowbreak{}On\allowbreak{}Sirach\_\allowbreak{}strength}}\par\nobreak\noindent
Trent on Sirach 18:22 rests on the Latin reading, rated \texttt{disputed}.

{\small\begin{verbatim}
theorem trentOnSirach_strength : trentOnSirach.strength = 0
-- axioms: propext
\end{verbatim}}
\par\medskip\noindent\textbf{\texttt{revelation\allowbreak{}Text\allowbreak{}Case\_\allowbreak{}strength}}\par\nobreak\noindent
The Greek of Revelation 22:11 rests on nothing below \texttt{well\allowbreak{}Supported}.

{\small\begin{verbatim}
theorem revelationTextCase_strength : revelationTextCase.strength = 2
-- axioms: propext
\end{verbatim}}
\par\medskip\noindent\textbf{\texttt{sirach\allowbreak{}Text\allowbreak{}Case\_\allowbreak{}strength}}\par\nobreak\noindent
The Greek of Sirach 18:22 rests on nothing below \texttt{well\allowbreak{}Supported}.

{\small\begin{verbatim}
theorem sirachTextCase_strength : sirachTextCase.strength = 2
-- axioms: propext
\end{verbatim}}
\par\medskip\noindent\textbf{\texttt{trent\allowbreak{}On\allowbreak{}Reward\_\allowbreak{}strength}}\par\nobreak\noindent
Trent on the reward of works rests on its reading of the texts, rated \texttt{disputed}.

{\small\begin{verbatim}
theorem trentOnReward_strength : trentOnReward.strength = 0
-- axioms: propext
\end{verbatim}}
\par\medskip\noindent\textbf{\texttt{westminster\allowbreak{}On\allowbreak{}Reward\_\allowbreak{}strength}}\par\nobreak\noindent
Westminster on the reward rests on its reading of the texts, rated \texttt{disputed}.

{\small\begin{verbatim}
theorem westminsterOnReward_strength : westminsterOnReward.strength = 0
-- axioms: propext
\end{verbatim}}
\par\medskip\noindent\textbf{\texttt{renewal\allowbreak{}On\allowbreak{}The\allowbreak{}Word\_\allowbreak{}strength}}\par\nobreak\noindent
The reading through Paul rests on Paul’s word as renewal, rated \texttt{disputed}.

{\small\begin{verbatim}
theorem renewalOnTheWord_strength : renewalOnTheWord.strength = 0
-- axioms: propext
\end{verbatim}}
\paragraph*{The debate}
\leavevmode\par\nobreak
\par\medskip\noindent\textbf{\texttt{james\allowbreak{}Party\allowbreak{}Node}}\par\nobreak\noindent
The premises each position argues from.

{\small\begin{verbatim}
def jamesPartyNode : JamesParty → ArgumentPackage Claim
\end{verbatim}}
\par\medskip\noindent\textbf{\texttt{james\allowbreak{}Dispute}}\par\nobreak\noindent
The debate over James 2:24 and the increase of justification. For every position heard: its premises can be held without contradiction, its conclusion follows from them, and every step it uses is rated. Each position’s weakest link is \texttt{well\allowbreak{}Supported} for the readings of the words and the texts, and the bottom for the rest.

{\small\begin{verbatim}
def jamesDispute : Dispute Claim JamesParty :=
  { node := jamesPartyNode,
    consistent := jamesDispute._proof_1,
    sound := jamesDispute._proof_2,
    rated := jamesDispute._proof_15 }
\end{verbatim}}
\par\medskip\noindent\textbf{\texttt{james\allowbreak{}Dispute\_\allowbreak{}defeats}}\par\nobreak\noindent
\textbf{Who defeats whom}, for all 144 pairs.

The argument from James’s word defeats Trent’s reading of that word, and the reply fails. It defeats the reading through Paul too, whose reply fails for the same reason: it rests on Paul’s word as renewal, rated \texttt{disputed}, against an argument rated \texttt{well\allowbreak{}Supported}.

The Greek of Revelation 22:11 and of Sirach 18:22 defeat Trent’s Latin readings of them, and the replies fail.

The Reformed reconciliation and each of Trent’s readings of what works do defeat each other: canon 24’s, and chapter 16’s from the reward texts. Westminster’s reading of the reward texts and chapter 16’s defeat each other: each denies what the other holds about merit, and both rest on a reading rated \texttt{disputed}.

The Reformed reconciliation and the reading through Paul defeat each other through what their claims mean. The reconciliation holds that James 2:24 does not contradict Paul, which says James’s word does not have Paul’s sense; the other holds that it does (\texttt{Testimony\allowbreak{}Arguments.Meanings.Sola\allowbreak{}Fide.joins}). Both rest on a reading rated \texttt{disputed}.

Nothing else. In particular, nothing here contradicts the argument about James’s faith. And the argument from James’s word does not bear on Trent’s reading of what works do, which accepts the word and makes a claim about God’s verdict.

Every answer is worked out from the positions’ premises and checked by machine.

{\small\begin{verbatim}
theorem jamesDispute_defeats : ∀ (i j : JamesParty), jamesDispute.defeats i j
    ↔ (i, j) ∈ [(JamesParty.lexical, JamesParty.trentOnTheWord),
    (JamesParty.harmony, JamesParty.trentOnWorks), (JamesParty.trentOnWorks,
    JamesParty.harmony), (JamesParty.revelationText,
    JamesParty.trentOnRevelation), (JamesParty.sirachText,
    JamesParty.trentOnSirach), (JamesParty.harmony, JamesParty.trentOnReward),
    (JamesParty.trentOnReward, JamesParty.harmony), (JamesParty.lexical,
    JamesParty.renewalOnTheWord), (JamesParty.westminsterOnReward,
    JamesParty.trentOnReward), (JamesParty.trentOnReward,
    JamesParty.westminsterOnReward), (JamesParty.harmony,
    JamesParty.renewalOnTheWord), (JamesParty.renewalOnTheWord,
    JamesParty.harmony)]
-- axioms: propext, Classical.choice, Quot.sound
\end{verbatim}}
\paragraph*{What the debate decides}
\leavevmode\par\nobreak
\par\medskip\noindent\textbf{\texttt{james\allowbreak{}Words\allowbreak{}Unanswered}}\par\nobreak\noindent
How far the debate can be settled. Nothing defeats the argument from James’s word, the argument about James’s faith, or the Greek of the two verses, so those are accepted first, and nothing joins them. Trent’s readings of James’s word, of Revelation 22:11 and of Sirach 18:22 each have an attacker that remains: the one that defeats it. The Reformed reconciliation and Trent’s reading of what works do have each other.

\paragraph{What the dispute forces is exactly \emph{James’s word (δικαιόω, by the rule of least meaning)}, \emph{James’s faith (the faith James calls dead)}, \emph{Revelation 22:11, in the Greek}, \emph{Sirach 18:22, in the Greek}.}
\begin{itemize}
\item Stage 1: \emph{James’s word (δικαιόω, by the rule of least meaning)}, \emph{James’s faith (the faith James calls dead)}, \emph{Revelation 22:11, in the Greek}, \emph{Sirach 18:22, in the Greek}.
\begin{itemize}
\item \emph{James’s word (δικαιόω, by the rule of least meaning)} is defeated by nothing.
\item \emph{James’s faith (the faith James calls dead)} is defeated by nothing.
\item \emph{Revelation 22:11, in the Greek} is defeated by nothing.
\item \emph{Sirach 18:22, in the Greek} is defeated by nothing.
\end{itemize}
\item Nothing else is forced:
\begin{itemize}
\item \emph{James 2:24, read by canon 24} defeats \emph{Harmonisation of James 2:24}, and nothing forced defeats \emph{James 2:24, read by canon 24}.
\item \emph{James’s word (δικαιόω, by the rule of least meaning)} defeats \emph{James 2:24, read for the increase of justification}, and nothing forced defeats \emph{James’s word (δικαιόω, by the rule of least meaning)}.
\item \emph{Harmonisation of James 2:24} defeats \emph{James 2:24, read by canon 24}, and nothing forced defeats \emph{Harmonisation of James 2:24}.
\item \emph{Revelation 22:11, in the Greek} defeats \emph{Revelation 22:11, as Trent quotes it}, and nothing forced defeats \emph{Revelation 22:11, in the Greek}.
\item \emph{Sirach 18:22, in the Greek} defeats \emph{Sirach 18:22, as Trent quotes it}, and nothing forced defeats \emph{Sirach 18:22, in the Greek}.
\item \emph{Harmonisation of James 2:24} defeats \emph{The reward of works (Trent, Session VI, ch. 16)}, and nothing forced defeats \emph{Harmonisation of James 2:24}.
\item \emph{James’s word (δικαιόω, by the rule of least meaning)} defeats \emph{James’s word, read through Paul’s (the renewal reading)}, and nothing forced defeats \emph{James’s word (δικαιόω, by the rule of least meaning)}.
\item \emph{The reward of works (Trent, Session VI, ch. 16)} defeats \emph{The reward of works, of grace (Westminster XVI.5, XVI.6)}, and nothing forced defeats \emph{The reward of works (Trent, Session VI, ch. 16)}.
\end{itemize}
\end{itemize}
\par\noindent\textbf{What this rests on.} The reasons state 8 defeats and 60 absences of defeat, each a cell of the defeat table, \texttt{james\allowbreak{}Dispute\_\allowbreak{}defeats}, computed from the two parties' premises and checked by the kernel. Whether an attack survives turns on the attacker's weakest link and on the rating of what it attacks, so a changed rating can change the verdict. The weakest links:
\begin{itemize}
\item \emph{James’s word (δικαιόω, by the rule of least meaning)}, weakest at \emph{well supported}:
\begin{itemize}
\item Outside Paul, ‘justify’ (δικαιόω) declares or vindicates; it does not mean making righteous --- \cite[s.v. δικαιόω]{bauer-danker-lexicon-2000}; Deut 25:1; Matt 11:19; Luke 7:35; \cite[ad loc. Matt 11:19]{crossway-esv-2001}
\item A word contributes the least meaning its context requires; one occurrence does not carry the whole concept --- \cite{silva-biblical-words-1994}; \cite{joos-semantic-axiom-1972}; \cite{barr-semantics-2004}
\item an inference step --- \cite[ad loc. Jas 2:21]{moo-james-2000}; \cite{silva-biblical-words-1994}; \cite{barr-semantics-2004}
\end{itemize}
\item \emph{James’s faith (the faith James calls dead)}, weakest at \emph{well supported}:
\begin{itemize}
\item The faith alone that James 2:24 denies justifies is the dead faith of 2:14–26 --- Jas 2:14-26; \cite[Trent, Session VI (1547), ch. 7]{tanner-decrees-1990}; \cite[ad loc. Jas 2:24]{moo-james-2000}
\item an inference step --- \cite[ad loc. Jas 2:24]{moo-james-2000}; \cite[XI.2]{westminster-confession-1647}
\end{itemize}
\item \emph{Harmonisation of James 2:24}, weakest at \emph{disputed}:
\begin{itemize}
\item an inference step --- \cite[ad loc. Jas 2:24]{moo-james-2000}; \cite[III.xvii.11]{calvin-institutes-1960}
\end{itemize}
\item \emph{James 2:24, read for the increase of justification}, weakest at \emph{disputed}:
\begin{itemize}
\item James’s ‘justified’ (δικαιόω) means the increase of the justice received --- \cite[Trent, Session VI (1547), ch. 10]{tanner-decrees-1990}
\item an inference step --- \cite[Trent, Session VI (1547), ch. 10]{tanner-decrees-1990}
\end{itemize}
\item \emph{James 2:24, read by canon 24}, weakest at \emph{disputed}:
\begin{itemize}
\item Good works are a cause of the increase of justification, not merely its fruits --- \cite[Trent, Session VI (1547), canon 24]{tanner-decrees-1990}
\item an inference step --- \cite[Trent, Session VI (1547), canon 24]{tanner-decrees-1990}
\end{itemize}
\item \emph{Revelation 22:11, as Trent quotes it}, weakest at \emph{disputed}:
\begin{itemize}
\item Revelation 22:11 says: he that is just, let him be justified still --- \cite[Trent, Session VI (1547), ch. 10]{tanner-decrees-1990}
\item an inference step --- \cite[Trent, Session VI (1547), ch. 10]{tanner-decrees-1990}
\end{itemize}
\item \emph{Sirach 18:22, as Trent quotes it}, weakest at \emph{disputed}:
\begin{itemize}
\item Sirach 18:22 says: be not afraid to be justified even to death --- \cite[Trent, Session VI (1547), ch. 10]{tanner-decrees-1990}
\item an inference step --- \cite[Trent, Session VI (1547), ch. 10]{tanner-decrees-1990}
\end{itemize}
\item \emph{Revelation 22:11, in the Greek}, weakest at \emph{well supported}:
\begin{itemize}
\item an inference step --- \cite[§16]{pius-divino-afflante-spiritu-1943}; \cite{holmes-sblgnt-2010}; \cite{usccb-nabre-2011}
\end{itemize}
\item \emph{Sirach 18:22, in the Greek}, weakest at \emph{well supported}:
\begin{itemize}
\item The Greek of Sirach 18:22 is about paying a vow: do not wait until death to be justified --- \cite[ad loc. Sir 18:22]{swete-old-testament-in-greek}; \cite[Sirach 18:22]{usccb-nabre-2011}
\item an inference step --- \cite[§16]{pius-divino-afflante-spiritu-1943}; \cite{holmes-sblgnt-2010}; \cite{usccb-nabre-2011}
\end{itemize}
\item \emph{The reward of works (Trent, Session VI, ch. 16)}, weakest at \emph{disputed}:
\begin{itemize}
\item The reward God renders to the works of the justified is rendered to them as merits --- \cite[Trent, Session VI (1547), ch. 16; canon 32]{tanner-decrees-1990}; \cite{dulles-two-languages-1999}
\item an inference step --- \cite[Trent, Session VI (1547), ch. 16, canon 32]{tanner-decrees-1990}; \cite[XVI.5, XVI.6]{westminster-confession-1647}
\end{itemize}
\item \emph{James’s word, read through Paul’s (the renewal reading)}, weakest at \emph{disputed}:
\begin{itemize}
\item Paul’s word ‘justify’ means the renewal of the inward man, not only a verdict --- \cite[Trent, Session VI (1547), Decree on Justification, chs. 7–8]{tanner-decrees-1990}; \cite{vanlandingham-judgment-justification-2006}
\item James’s ‘justify’ (δικαιόω) has the sense Paul’s has: the justification of Romans and Galatians --- \cite[Trent, Session VI (1547), ch. 10]{tanner-decrees-1990}
\item an inference step --- \cite[Trent, Session VI (1547), ch. 10]{tanner-decrees-1990}; \cite{vanlandingham-judgment-justification-2006}
\end{itemize}
\item \emph{The reward of works, of grace (Westminster XVI.5, XVI.6)}, weakest at \emph{disputed}:
\begin{itemize}
\item an inference step --- \cite[XVI.5, XVI.6]{westminster-confession-1647}
\end{itemize}
\end{itemize}
\par\medskip\noindent\textbf{\texttt{james\_\allowbreak{}words\_\allowbreak{}prevail}}\par\nobreak\noindent
\textbf{What James’s words say, and what the Greek of Trent’s other two texts says, prevails outright.} Exactly four positions are forced. James’s “justified” (δικαιόω) does not mean the increase of justice. The faith James 2:24 denies is not the faith of the Reformers’ formula. Revelation 22:11, in the critical Greek text, does not say “be justified still”. And Sirach 18:22, in Greek, does not say “be not afraid to be justified even to death”. Nothing in the debate answers any of them.

It does not show that James 2:24 agrees with Paul, or that works are not a cause of the increase of justification. Neither is forced. Nor that the Greek of Sirach 18:22 lacks the word “justified”: it has it, in a sentence about vows.

{\small\begin{verbatim}
theorem james_words_prevail : Framework.grounded jamesDispute.defeats =
    {JamesParty.lexical, JamesParty.faith, JamesParty.revelationText,
    JamesParty.sirachText}
-- axioms: propext, Classical.choice, Quot.sound
\end{verbatim}}
\par\medskip\noindent\textbf{\texttt{trent\allowbreak{}On\allowbreak{}The\allowbreak{}Word\allowbreak{}Answered}}\par\nobreak\noindent
Why Trent’s reading of James’s word cannot be defended: the argument from James’s word defeats it, and nothing defeats that argument.

\paragraph{\emph{James 2:24, read for the increase of justification} cannot be defended: no admissible position holds it.}
\begin{itemize}
\item \emph{James’s word (δικαιόω, by the rule of least meaning)} defeats \emph{James 2:24, read for the increase of justification}, and nothing defeats \emph{James’s word (δικαιόω, by the rule of least meaning)}.
\end{itemize}
\par\noindent\textbf{What this rests on.} The reasons state 1 defeat and 12 absences of defeat, each a cell of the defeat table, \texttt{james\allowbreak{}Dispute\_\allowbreak{}defeats}, computed from the two parties' premises and checked by the kernel. Whether an attack survives turns on the attacker's weakest link and on the rating of what it attacks, so a changed rating can change the verdict. The weakest links:
\begin{itemize}
\item \emph{James’s word (δικαιόω, by the rule of least meaning)}, weakest at \emph{well supported}:
\begin{itemize}
\item Outside Paul, ‘justify’ (δικαιόω) declares or vindicates; it does not mean making righteous --- \cite[s.v. δικαιόω]{bauer-danker-lexicon-2000}; Deut 25:1; Matt 11:19; Luke 7:35; \cite[ad loc. Matt 11:19]{crossway-esv-2001}
\item A word contributes the least meaning its context requires; one occurrence does not carry the whole concept --- \cite{silva-biblical-words-1994}; \cite{joos-semantic-axiom-1972}; \cite{barr-semantics-2004}
\item an inference step --- \cite[ad loc. Jas 2:21]{moo-james-2000}; \cite{silva-biblical-words-1994}; \cite{barr-semantics-2004}
\end{itemize}
\item \emph{James 2:24, read for the increase of justification}, weakest at \emph{disputed}:
\begin{itemize}
\item James’s ‘justified’ (δικαιόω) means the increase of the justice received --- \cite[Trent, Session VI (1547), ch. 10]{tanner-decrees-1990}
\item an inference step --- \cite[Trent, Session VI (1547), ch. 10]{tanner-decrees-1990}
\end{itemize}
\item \emph{James’s faith (the faith James calls dead)}, weakest at \emph{well supported}:
\begin{itemize}
\item The faith alone that James 2:24 denies justifies is the dead faith of 2:14–26 --- Jas 2:14-26; \cite[Trent, Session VI (1547), ch. 7]{tanner-decrees-1990}; \cite[ad loc. Jas 2:24]{moo-james-2000}
\item an inference step --- \cite[ad loc. Jas 2:24]{moo-james-2000}; \cite[XI.2]{westminster-confession-1647}
\end{itemize}
\item \emph{Harmonisation of James 2:24}, weakest at \emph{disputed}:
\begin{itemize}
\item an inference step --- \cite[ad loc. Jas 2:24]{moo-james-2000}; \cite[III.xvii.11]{calvin-institutes-1960}
\end{itemize}
\item \emph{James 2:24, read by canon 24}, weakest at \emph{disputed}:
\begin{itemize}
\item Good works are a cause of the increase of justification, not merely its fruits --- \cite[Trent, Session VI (1547), canon 24]{tanner-decrees-1990}
\item an inference step --- \cite[Trent, Session VI (1547), canon 24]{tanner-decrees-1990}
\end{itemize}
\item \emph{Revelation 22:11, as Trent quotes it}, weakest at \emph{disputed}:
\begin{itemize}
\item Revelation 22:11 says: he that is just, let him be justified still --- \cite[Trent, Session VI (1547), ch. 10]{tanner-decrees-1990}
\item an inference step --- \cite[Trent, Session VI (1547), ch. 10]{tanner-decrees-1990}
\end{itemize}
\item \emph{Sirach 18:22, as Trent quotes it}, weakest at \emph{disputed}:
\begin{itemize}
\item Sirach 18:22 says: be not afraid to be justified even to death --- \cite[Trent, Session VI (1547), ch. 10]{tanner-decrees-1990}
\item an inference step --- \cite[Trent, Session VI (1547), ch. 10]{tanner-decrees-1990}
\end{itemize}
\item \emph{Revelation 22:11, in the Greek}, weakest at \emph{well supported}:
\begin{itemize}
\item an inference step --- \cite[§16]{pius-divino-afflante-spiritu-1943}; \cite{holmes-sblgnt-2010}; \cite{usccb-nabre-2011}
\end{itemize}
\item \emph{Sirach 18:22, in the Greek}, weakest at \emph{well supported}:
\begin{itemize}
\item The Greek of Sirach 18:22 is about paying a vow: do not wait until death to be justified --- \cite[ad loc. Sir 18:22]{swete-old-testament-in-greek}; \cite[Sirach 18:22]{usccb-nabre-2011}
\item an inference step --- \cite[§16]{pius-divino-afflante-spiritu-1943}; \cite{holmes-sblgnt-2010}; \cite{usccb-nabre-2011}
\end{itemize}
\item \emph{The reward of works (Trent, Session VI, ch. 16)}, weakest at \emph{disputed}:
\begin{itemize}
\item The reward God renders to the works of the justified is rendered to them as merits --- \cite[Trent, Session VI (1547), ch. 16; canon 32]{tanner-decrees-1990}; \cite{dulles-two-languages-1999}
\item an inference step --- \cite[Trent, Session VI (1547), ch. 16, canon 32]{tanner-decrees-1990}; \cite[XVI.5, XVI.6]{westminster-confession-1647}
\end{itemize}
\item \emph{James’s word, read through Paul’s (the renewal reading)}, weakest at \emph{disputed}:
\begin{itemize}
\item Paul’s word ‘justify’ means the renewal of the inward man, not only a verdict --- \cite[Trent, Session VI (1547), Decree on Justification, chs. 7–8]{tanner-decrees-1990}; \cite{vanlandingham-judgment-justification-2006}
\item James’s ‘justify’ (δικαιόω) has the sense Paul’s has: the justification of Romans and Galatians --- \cite[Trent, Session VI (1547), ch. 10]{tanner-decrees-1990}
\item an inference step --- \cite[Trent, Session VI (1547), ch. 10]{tanner-decrees-1990}; \cite{vanlandingham-judgment-justification-2006}
\end{itemize}
\item \emph{The reward of works, of grace (Westminster XVI.5, XVI.6)}, weakest at \emph{disputed}:
\begin{itemize}
\item an inference step --- \cite[XVI.5, XVI.6]{westminster-confession-1647}
\end{itemize}
\end{itemize}
\par\medskip\noindent\textbf{\texttt{trent\_\allowbreak{}on\_\allowbreak{}james\_\allowbreak{}word\_\allowbreak{}indefensible}}\par\nobreak\noindent
\textbf{Trent’s reading of James’s word cannot be defended.} If “justified by works” in James 2:24 is taken to mean the increase of the justice received, being made more just, it cannot be defended in this debate. Outside Paul the word declares, acquits or vindicates. James uses it of Abraham at Genesis 22, after Genesis 15:6 had been spoken, and of Rahab. And nothing in James’s setting requires a second, inward sense. The step from that evidence to “the word does not mean the increase of justice” is rated \texttt{well\allowbreak{}Supported}. The reading rests on Trent’s word alone.

It does not show that works are no cause of the increase of justification. That is a claim about God’s verdict, not about James’s word, and it is weighed separately (\texttt{trent\_\allowbreak{}on\_\allowbreak{}works\_\allowbreak{}defensible}).

{\small\begin{verbatim}
theorem trent_on_james_word_indefensible : ∀ (S : Set JamesParty),
    Framework.Admissible jamesDispute.defeats S → JamesParty.trentOnTheWord ∉
    S
-- axioms: propext, Classical.choice, Quot.sound
\end{verbatim}}
\par\medskip\noindent\textbf{\texttt{trent\allowbreak{}On\allowbreak{}Revelation\allowbreak{}Answered}}\par\nobreak\noindent
Why Trent’s Revelation 22:11 cannot be defended: the Greek defeats it, and nothing defeats the Greek.

\paragraph{\emph{Revelation 22:11, as Trent quotes it} cannot be defended: no admissible position holds it.}
\begin{itemize}
\item \emph{Revelation 22:11, in the Greek} defeats \emph{Revelation 22:11, as Trent quotes it}, and nothing defeats \emph{Revelation 22:11, in the Greek}.
\end{itemize}
\par\noindent\textbf{What this rests on.} The reasons state 1 defeat and 12 absences of defeat, each a cell of the defeat table, \texttt{james\allowbreak{}Dispute\_\allowbreak{}defeats}, computed from the two parties' premises and checked by the kernel. Whether an attack survives turns on the attacker's weakest link and on the rating of what it attacks, so a changed rating can change the verdict. The weakest links:
\begin{itemize}
\item \emph{Revelation 22:11, in the Greek}, weakest at \emph{well supported}:
\begin{itemize}
\item an inference step --- \cite[§16]{pius-divino-afflante-spiritu-1943}; \cite{holmes-sblgnt-2010}; \cite{usccb-nabre-2011}
\end{itemize}
\item \emph{Revelation 22:11, as Trent quotes it}, weakest at \emph{disputed}:
\begin{itemize}
\item Revelation 22:11 says: he that is just, let him be justified still --- \cite[Trent, Session VI (1547), ch. 10]{tanner-decrees-1990}
\item an inference step --- \cite[Trent, Session VI (1547), ch. 10]{tanner-decrees-1990}
\end{itemize}
\item \emph{James’s word (δικαιόω, by the rule of least meaning)}, weakest at \emph{well supported}:
\begin{itemize}
\item Outside Paul, ‘justify’ (δικαιόω) declares or vindicates; it does not mean making righteous --- \cite[s.v. δικαιόω]{bauer-danker-lexicon-2000}; Deut 25:1; Matt 11:19; Luke 7:35; \cite[ad loc. Matt 11:19]{crossway-esv-2001}
\item A word contributes the least meaning its context requires; one occurrence does not carry the whole concept --- \cite{silva-biblical-words-1994}; \cite{joos-semantic-axiom-1972}; \cite{barr-semantics-2004}
\item an inference step --- \cite[ad loc. Jas 2:21]{moo-james-2000}; \cite{silva-biblical-words-1994}; \cite{barr-semantics-2004}
\end{itemize}
\item \emph{James’s faith (the faith James calls dead)}, weakest at \emph{well supported}:
\begin{itemize}
\item The faith alone that James 2:24 denies justifies is the dead faith of 2:14–26 --- Jas 2:14-26; \cite[Trent, Session VI (1547), ch. 7]{tanner-decrees-1990}; \cite[ad loc. Jas 2:24]{moo-james-2000}
\item an inference step --- \cite[ad loc. Jas 2:24]{moo-james-2000}; \cite[XI.2]{westminster-confession-1647}
\end{itemize}
\item \emph{Harmonisation of James 2:24}, weakest at \emph{disputed}:
\begin{itemize}
\item an inference step --- \cite[ad loc. Jas 2:24]{moo-james-2000}; \cite[III.xvii.11]{calvin-institutes-1960}
\end{itemize}
\item \emph{James 2:24, read for the increase of justification}, weakest at \emph{disputed}:
\begin{itemize}
\item James’s ‘justified’ (δικαιόω) means the increase of the justice received --- \cite[Trent, Session VI (1547), ch. 10]{tanner-decrees-1990}
\item an inference step --- \cite[Trent, Session VI (1547), ch. 10]{tanner-decrees-1990}
\end{itemize}
\item \emph{James 2:24, read by canon 24}, weakest at \emph{disputed}:
\begin{itemize}
\item Good works are a cause of the increase of justification, not merely its fruits --- \cite[Trent, Session VI (1547), canon 24]{tanner-decrees-1990}
\item an inference step --- \cite[Trent, Session VI (1547), canon 24]{tanner-decrees-1990}
\end{itemize}
\item \emph{Sirach 18:22, as Trent quotes it}, weakest at \emph{disputed}:
\begin{itemize}
\item Sirach 18:22 says: be not afraid to be justified even to death --- \cite[Trent, Session VI (1547), ch. 10]{tanner-decrees-1990}
\item an inference step --- \cite[Trent, Session VI (1547), ch. 10]{tanner-decrees-1990}
\end{itemize}
\item \emph{Sirach 18:22, in the Greek}, weakest at \emph{well supported}:
\begin{itemize}
\item The Greek of Sirach 18:22 is about paying a vow: do not wait until death to be justified --- \cite[ad loc. Sir 18:22]{swete-old-testament-in-greek}; \cite[Sirach 18:22]{usccb-nabre-2011}
\item an inference step --- \cite[§16]{pius-divino-afflante-spiritu-1943}; \cite{holmes-sblgnt-2010}; \cite{usccb-nabre-2011}
\end{itemize}
\item \emph{The reward of works (Trent, Session VI, ch. 16)}, weakest at \emph{disputed}:
\begin{itemize}
\item The reward God renders to the works of the justified is rendered to them as merits --- \cite[Trent, Session VI (1547), ch. 16; canon 32]{tanner-decrees-1990}; \cite{dulles-two-languages-1999}
\item an inference step --- \cite[Trent, Session VI (1547), ch. 16, canon 32]{tanner-decrees-1990}; \cite[XVI.5, XVI.6]{westminster-confession-1647}
\end{itemize}
\item \emph{James’s word, read through Paul’s (the renewal reading)}, weakest at \emph{disputed}:
\begin{itemize}
\item Paul’s word ‘justify’ means the renewal of the inward man, not only a verdict --- \cite[Trent, Session VI (1547), Decree on Justification, chs. 7–8]{tanner-decrees-1990}; \cite{vanlandingham-judgment-justification-2006}
\item James’s ‘justify’ (δικαιόω) has the sense Paul’s has: the justification of Romans and Galatians --- \cite[Trent, Session VI (1547), ch. 10]{tanner-decrees-1990}
\item an inference step --- \cite[Trent, Session VI (1547), ch. 10]{tanner-decrees-1990}; \cite{vanlandingham-judgment-justification-2006}
\end{itemize}
\item \emph{The reward of works, of grace (Westminster XVI.5, XVI.6)}, weakest at \emph{disputed}:
\begin{itemize}
\item an inference step --- \cite[XVI.5, XVI.6]{westminster-confession-1647}
\end{itemize}
\end{itemize}
\par\medskip\noindent\textbf{\texttt{trent\_\allowbreak{}on\_\allowbreak{}revelation\_\allowbreak{}indefensible}}\par\nobreak\noindent
\textbf{Trent’s Revelation 22:11 cannot be defended as a text for the increase of justification.} “He that is just, let him be justified still” is the Latin Vulgate’s wording (\emph{iustificetur adhuc}). The Greek, as modern editions print it and as the Catholic Church’s own translation gives it, says “the righteous must still do right”. And the original, “written by the inspired author himself”, outweighs any translation (\emph{Divino Afflante Spiritu} §16).

It does not show that Trent’s citation is an error of doctrine. Taken as Pius XII takes Trent’s decree on the Vulgate, as a reading the Church approves and holds free from error in faith, it is consistent with all he says, and rests on the Church’s authority and not on the inspired text (the chapter “Trent’s two Latin texts”).

{\small\begin{verbatim}
theorem trent_on_revelation_indefensible : ∀ (S : Set JamesParty),
    Framework.Admissible jamesDispute.defeats S → JamesParty.trentOnRevelation
    ∉ S
-- axioms: propext, Classical.choice, Quot.sound
\end{verbatim}}
\par\medskip\noindent\textbf{\texttt{trent\allowbreak{}On\allowbreak{}Sirach\allowbreak{}Answered}}\par\nobreak\noindent
Why Trent’s Sirach 18:22 cannot be defended: the Greek defeats it, and nothing defeats the Greek.

\paragraph{\emph{Sirach 18:22, as Trent quotes it} cannot be defended: no admissible position holds it.}
\begin{itemize}
\item \emph{Sirach 18:22, in the Greek} defeats \emph{Sirach 18:22, as Trent quotes it}, and nothing defeats \emph{Sirach 18:22, in the Greek}.
\end{itemize}
\par\noindent\textbf{What this rests on.} The reasons state 1 defeat and 12 absences of defeat, each a cell of the defeat table, \texttt{james\allowbreak{}Dispute\_\allowbreak{}defeats}, computed from the two parties' premises and checked by the kernel. Whether an attack survives turns on the attacker's weakest link and on the rating of what it attacks, so a changed rating can change the verdict. The weakest links:
\begin{itemize}
\item \emph{Sirach 18:22, in the Greek}, weakest at \emph{well supported}:
\begin{itemize}
\item The Greek of Sirach 18:22 is about paying a vow: do not wait until death to be justified --- \cite[ad loc. Sir 18:22]{swete-old-testament-in-greek}; \cite[Sirach 18:22]{usccb-nabre-2011}
\item an inference step --- \cite[§16]{pius-divino-afflante-spiritu-1943}; \cite{holmes-sblgnt-2010}; \cite{usccb-nabre-2011}
\end{itemize}
\item \emph{Sirach 18:22, as Trent quotes it}, weakest at \emph{disputed}:
\begin{itemize}
\item Sirach 18:22 says: be not afraid to be justified even to death --- \cite[Trent, Session VI (1547), ch. 10]{tanner-decrees-1990}
\item an inference step --- \cite[Trent, Session VI (1547), ch. 10]{tanner-decrees-1990}
\end{itemize}
\item \emph{James’s word (δικαιόω, by the rule of least meaning)}, weakest at \emph{well supported}:
\begin{itemize}
\item Outside Paul, ‘justify’ (δικαιόω) declares or vindicates; it does not mean making righteous --- \cite[s.v. δικαιόω]{bauer-danker-lexicon-2000}; Deut 25:1; Matt 11:19; Luke 7:35; \cite[ad loc. Matt 11:19]{crossway-esv-2001}
\item A word contributes the least meaning its context requires; one occurrence does not carry the whole concept --- \cite{silva-biblical-words-1994}; \cite{joos-semantic-axiom-1972}; \cite{barr-semantics-2004}
\item an inference step --- \cite[ad loc. Jas 2:21]{moo-james-2000}; \cite{silva-biblical-words-1994}; \cite{barr-semantics-2004}
\end{itemize}
\item \emph{James’s faith (the faith James calls dead)}, weakest at \emph{well supported}:
\begin{itemize}
\item The faith alone that James 2:24 denies justifies is the dead faith of 2:14–26 --- Jas 2:14-26; \cite[Trent, Session VI (1547), ch. 7]{tanner-decrees-1990}; \cite[ad loc. Jas 2:24]{moo-james-2000}
\item an inference step --- \cite[ad loc. Jas 2:24]{moo-james-2000}; \cite[XI.2]{westminster-confession-1647}
\end{itemize}
\item \emph{Harmonisation of James 2:24}, weakest at \emph{disputed}:
\begin{itemize}
\item an inference step --- \cite[ad loc. Jas 2:24]{moo-james-2000}; \cite[III.xvii.11]{calvin-institutes-1960}
\end{itemize}
\item \emph{James 2:24, read for the increase of justification}, weakest at \emph{disputed}:
\begin{itemize}
\item James’s ‘justified’ (δικαιόω) means the increase of the justice received --- \cite[Trent, Session VI (1547), ch. 10]{tanner-decrees-1990}
\item an inference step --- \cite[Trent, Session VI (1547), ch. 10]{tanner-decrees-1990}
\end{itemize}
\item \emph{James 2:24, read by canon 24}, weakest at \emph{disputed}:
\begin{itemize}
\item Good works are a cause of the increase of justification, not merely its fruits --- \cite[Trent, Session VI (1547), canon 24]{tanner-decrees-1990}
\item an inference step --- \cite[Trent, Session VI (1547), canon 24]{tanner-decrees-1990}
\end{itemize}
\item \emph{Revelation 22:11, as Trent quotes it}, weakest at \emph{disputed}:
\begin{itemize}
\item Revelation 22:11 says: he that is just, let him be justified still --- \cite[Trent, Session VI (1547), ch. 10]{tanner-decrees-1990}
\item an inference step --- \cite[Trent, Session VI (1547), ch. 10]{tanner-decrees-1990}
\end{itemize}
\item \emph{Revelation 22:11, in the Greek}, weakest at \emph{well supported}:
\begin{itemize}
\item an inference step --- \cite[§16]{pius-divino-afflante-spiritu-1943}; \cite{holmes-sblgnt-2010}; \cite{usccb-nabre-2011}
\end{itemize}
\item \emph{The reward of works (Trent, Session VI, ch. 16)}, weakest at \emph{disputed}:
\begin{itemize}
\item The reward God renders to the works of the justified is rendered to them as merits --- \cite[Trent, Session VI (1547), ch. 16; canon 32]{tanner-decrees-1990}; \cite{dulles-two-languages-1999}
\item an inference step --- \cite[Trent, Session VI (1547), ch. 16, canon 32]{tanner-decrees-1990}; \cite[XVI.5, XVI.6]{westminster-confession-1647}
\end{itemize}
\item \emph{James’s word, read through Paul’s (the renewal reading)}, weakest at \emph{disputed}:
\begin{itemize}
\item Paul’s word ‘justify’ means the renewal of the inward man, not only a verdict --- \cite[Trent, Session VI (1547), Decree on Justification, chs. 7–8]{tanner-decrees-1990}; \cite{vanlandingham-judgment-justification-2006}
\item James’s ‘justify’ (δικαιόω) has the sense Paul’s has: the justification of Romans and Galatians --- \cite[Trent, Session VI (1547), ch. 10]{tanner-decrees-1990}
\item an inference step --- \cite[Trent, Session VI (1547), ch. 10]{tanner-decrees-1990}; \cite{vanlandingham-judgment-justification-2006}
\end{itemize}
\item \emph{The reward of works, of grace (Westminster XVI.5, XVI.6)}, weakest at \emph{disputed}:
\begin{itemize}
\item an inference step --- \cite[XVI.5, XVI.6]{westminster-confession-1647}
\end{itemize}
\end{itemize}
\par\medskip\noindent\textbf{\texttt{trent\_\allowbreak{}on\_\allowbreak{}sirach\_\allowbreak{}indefensible}}\par\nobreak\noindent
\textbf{Trent’s Sirach 18:22 cannot be defended as what the verse says.} “Be not afraid to be justified even to death” is the Vulgate’s verse. The Greek, in a sentence about paying vows, says “do not wait until death to be justified”, and the Catholic Church’s own translation renders that “do not wait until death to fulfill them”. Whether Sirach is Scripture is not in question here: the verdict holds on Trent’s own list of the biblical books.

It does not show that the Greek has no “be justified”. It has the verb. What it lacks is the Latin’s “be not afraid ... even to death”, the words in which Trent reads continued growth in justice.

{\small\begin{verbatim}
theorem trent_on_sirach_indefensible : ∀ (S : Set JamesParty),
    Framework.Admissible jamesDispute.defeats S → JamesParty.trentOnSirach ∉ S
-- axioms: propext, Classical.choice, Quot.sound
\end{verbatim}}
\par\medskip\noindent\textbf{\texttt{trent\allowbreak{}On\allowbreak{}Works\allowbreak{}Stands}}\par\nobreak\noindent
Why Trent’s reading of what works do can be defended: it answers the one position that attacks it, the Reformed reconciliation.

\paragraph{\emph{James 2:24, read by canon 24} is accepted on some resolution.}
\begin{itemize}
\item A position holding \emph{James’s word (δικαιόω, by the rule of least meaning)}, \emph{James’s faith (the faith James calls dead)}, \emph{James 2:24, read by canon 24}, \emph{Revelation 22:11, in the Greek}, \emph{Sirach 18:22, in the Greek} can be held.
\item None of \emph{James’s word (δικαιόω, by the rule of least meaning)}, \emph{James’s faith (the faith James calls dead)}, \emph{James 2:24, read by canon 24}, \emph{Revelation 22:11, in the Greek}, \emph{Sirach 18:22, in the Greek} defeats another, and each attack on them is answered from among them:
\begin{itemize}
\item \emph{Harmonisation of James 2:24} defeats \emph{James 2:24, read by canon 24}, and \emph{James 2:24, read by canon 24} defeats \emph{Harmonisation of James 2:24}.
\end{itemize}
\end{itemize}
\par\noindent\textbf{What this rests on.} The reasons state 2 defeats and 59 absences of defeat, each a cell of the defeat table, \texttt{james\allowbreak{}Dispute\_\allowbreak{}defeats}, computed from the two parties' premises and checked by the kernel. Whether an attack survives turns on the attacker's weakest link and on the rating of what it attacks, so a changed rating can change the verdict. The weakest links:
\begin{itemize}
\item \emph{James’s word (δικαιόω, by the rule of least meaning)}, weakest at \emph{well supported}:
\begin{itemize}
\item Outside Paul, ‘justify’ (δικαιόω) declares or vindicates; it does not mean making righteous --- \cite[s.v. δικαιόω]{bauer-danker-lexicon-2000}; Deut 25:1; Matt 11:19; Luke 7:35; \cite[ad loc. Matt 11:19]{crossway-esv-2001}
\item A word contributes the least meaning its context requires; one occurrence does not carry the whole concept --- \cite{silva-biblical-words-1994}; \cite{joos-semantic-axiom-1972}; \cite{barr-semantics-2004}
\item an inference step --- \cite[ad loc. Jas 2:21]{moo-james-2000}; \cite{silva-biblical-words-1994}; \cite{barr-semantics-2004}
\end{itemize}
\item \emph{James’s faith (the faith James calls dead)}, weakest at \emph{well supported}:
\begin{itemize}
\item The faith alone that James 2:24 denies justifies is the dead faith of 2:14–26 --- Jas 2:14-26; \cite[Trent, Session VI (1547), ch. 7]{tanner-decrees-1990}; \cite[ad loc. Jas 2:24]{moo-james-2000}
\item an inference step --- \cite[ad loc. Jas 2:24]{moo-james-2000}; \cite[XI.2]{westminster-confession-1647}
\end{itemize}
\item \emph{James 2:24, read by canon 24}, weakest at \emph{disputed}:
\begin{itemize}
\item Good works are a cause of the increase of justification, not merely its fruits --- \cite[Trent, Session VI (1547), canon 24]{tanner-decrees-1990}
\item an inference step --- \cite[Trent, Session VI (1547), canon 24]{tanner-decrees-1990}
\end{itemize}
\item \emph{Revelation 22:11, in the Greek}, weakest at \emph{well supported}:
\begin{itemize}
\item an inference step --- \cite[§16]{pius-divino-afflante-spiritu-1943}; \cite{holmes-sblgnt-2010}; \cite{usccb-nabre-2011}
\end{itemize}
\item \emph{Sirach 18:22, in the Greek}, weakest at \emph{well supported}:
\begin{itemize}
\item The Greek of Sirach 18:22 is about paying a vow: do not wait until death to be justified --- \cite[ad loc. Sir 18:22]{swete-old-testament-in-greek}; \cite[Sirach 18:22]{usccb-nabre-2011}
\item an inference step --- \cite[§16]{pius-divino-afflante-spiritu-1943}; \cite{holmes-sblgnt-2010}; \cite{usccb-nabre-2011}
\end{itemize}
\item \emph{Harmonisation of James 2:24}, weakest at \emph{disputed}:
\begin{itemize}
\item an inference step --- \cite[ad loc. Jas 2:24]{moo-james-2000}; \cite[III.xvii.11]{calvin-institutes-1960}
\end{itemize}
\item \emph{James 2:24, read for the increase of justification}, weakest at \emph{disputed}:
\begin{itemize}
\item James’s ‘justified’ (δικαιόω) means the increase of the justice received --- \cite[Trent, Session VI (1547), ch. 10]{tanner-decrees-1990}
\item an inference step --- \cite[Trent, Session VI (1547), ch. 10]{tanner-decrees-1990}
\end{itemize}
\item \emph{Revelation 22:11, as Trent quotes it}, weakest at \emph{disputed}:
\begin{itemize}
\item Revelation 22:11 says: he that is just, let him be justified still --- \cite[Trent, Session VI (1547), ch. 10]{tanner-decrees-1990}
\item an inference step --- \cite[Trent, Session VI (1547), ch. 10]{tanner-decrees-1990}
\end{itemize}
\item \emph{Sirach 18:22, as Trent quotes it}, weakest at \emph{disputed}:
\begin{itemize}
\item Sirach 18:22 says: be not afraid to be justified even to death --- \cite[Trent, Session VI (1547), ch. 10]{tanner-decrees-1990}
\item an inference step --- \cite[Trent, Session VI (1547), ch. 10]{tanner-decrees-1990}
\end{itemize}
\item \emph{The reward of works (Trent, Session VI, ch. 16)}, weakest at \emph{disputed}:
\begin{itemize}
\item The reward God renders to the works of the justified is rendered to them as merits --- \cite[Trent, Session VI (1547), ch. 16; canon 32]{tanner-decrees-1990}; \cite{dulles-two-languages-1999}
\item an inference step --- \cite[Trent, Session VI (1547), ch. 16, canon 32]{tanner-decrees-1990}; \cite[XVI.5, XVI.6]{westminster-confession-1647}
\end{itemize}
\item \emph{James’s word, read through Paul’s (the renewal reading)}, weakest at \emph{disputed}:
\begin{itemize}
\item Paul’s word ‘justify’ means the renewal of the inward man, not only a verdict --- \cite[Trent, Session VI (1547), Decree on Justification, chs. 7–8]{tanner-decrees-1990}; \cite{vanlandingham-judgment-justification-2006}
\item James’s ‘justify’ (δικαιόω) has the sense Paul’s has: the justification of Romans and Galatians --- \cite[Trent, Session VI (1547), ch. 10]{tanner-decrees-1990}
\item an inference step --- \cite[Trent, Session VI (1547), ch. 10]{tanner-decrees-1990}; \cite{vanlandingham-judgment-justification-2006}
\end{itemize}
\item \emph{The reward of works, of grace (Westminster XVI.5, XVI.6)}, weakest at \emph{disputed}:
\begin{itemize}
\item an inference step --- \cite[XVI.5, XVI.6]{westminster-confession-1647}
\end{itemize}
\end{itemize}
\par\medskip\noindent\textbf{\texttt{trent\_\allowbreak{}on\_\allowbreak{}works\_\allowbreak{}defensible}}\par\nobreak\noindent
\textbf{Trent’s reading of what works do can be defended}, together with the argument from James’s word, the argument about James’s faith, and the Greek of both verses. It accepts that James’s “justified” is a verdict, and holds that the works by which James says a person is justified are a cause of that verdict’s increase and not merely its fruit. The Reformed reconciliation denies this, and the debate cannot choose: at that point each side rests on its own confession, rated \texttt{disputed}.

It does not show that works are a cause of the increase. The reading is not forced (\texttt{trent\_\allowbreak{}on\_\allowbreak{}works\_\allowbreak{}not\_\allowbreak{}forced}).

{\small\begin{verbatim}
theorem trent_on_works_defensible : Framework.CredulouslyAccepted
    jamesDispute.defeats JamesParty.trentOnWorks
-- axioms: propext, Classical.choice, Quot.sound
\end{verbatim}}
\par\medskip\noindent\textbf{\texttt{trent\allowbreak{}On\allowbreak{}Works\allowbreak{}Answered}}\par\nobreak\noindent
Why it is not forced: a set of positions that can be defended includes the Reformed reconciliation, which defeats it.

\paragraph{\emph{James 2:24, read by canon 24} is not accepted on every resolution.}
\begin{itemize}
\item A position holding \emph{James’s word (δικαιόω, by the rule of least meaning)}, \emph{James’s faith (the faith James calls dead)}, \emph{Harmonisation of James 2:24}, \emph{Revelation 22:11, in the Greek}, \emph{Sirach 18:22, in the Greek} can be held, and it holds \emph{Harmonisation of James 2:24}; \emph{Harmonisation of James 2:24} conflicts with \emph{James 2:24, read by canon 24}: \emph{Harmonisation of James 2:24} defeats \emph{James 2:24, read by canon 24}, and \emph{James 2:24, read by canon 24} defeats \emph{Harmonisation of James 2:24}.
\item None of \emph{James’s word (δικαιόω, by the rule of least meaning)}, \emph{James’s faith (the faith James calls dead)}, \emph{Harmonisation of James 2:24}, \emph{Revelation 22:11, in the Greek}, \emph{Sirach 18:22, in the Greek} defeats another, and each attack on them is answered from among them:
\begin{itemize}
\item \emph{James 2:24, read by canon 24} defeats \emph{Harmonisation of James 2:24}, and \emph{Harmonisation of James 2:24} defeats \emph{James 2:24, read by canon 24}.
\item \emph{The reward of works (Trent, Session VI, ch. 16)} defeats \emph{Harmonisation of James 2:24}, and \emph{Harmonisation of James 2:24} defeats \emph{The reward of works (Trent, Session VI, ch. 16)}.
\item \emph{James’s word, read through Paul’s (the renewal reading)} defeats \emph{Harmonisation of James 2:24}, and \emph{James’s word (δικαιόω, by the rule of least meaning)} defeats \emph{James’s word, read through Paul’s (the renewal reading)}.
\end{itemize}
\end{itemize}
\par\noindent\textbf{What this rests on.} The reasons state 6 defeats and 57 absences of defeat, each a cell of the defeat table, \texttt{james\allowbreak{}Dispute\_\allowbreak{}defeats}, computed from the two parties' premises and checked by the kernel. Whether an attack survives turns on the attacker's weakest link and on the rating of what it attacks, so a changed rating can change the verdict. The weakest links:
\begin{itemize}
\item \emph{Harmonisation of James 2:24}, weakest at \emph{disputed}:
\begin{itemize}
\item an inference step --- \cite[ad loc. Jas 2:24]{moo-james-2000}; \cite[III.xvii.11]{calvin-institutes-1960}
\end{itemize}
\item \emph{James 2:24, read by canon 24}, weakest at \emph{disputed}:
\begin{itemize}
\item Good works are a cause of the increase of justification, not merely its fruits --- \cite[Trent, Session VI (1547), canon 24]{tanner-decrees-1990}
\item an inference step --- \cite[Trent, Session VI (1547), canon 24]{tanner-decrees-1990}
\end{itemize}
\item \emph{James’s word (δικαιόω, by the rule of least meaning)}, weakest at \emph{well supported}:
\begin{itemize}
\item Outside Paul, ‘justify’ (δικαιόω) declares or vindicates; it does not mean making righteous --- \cite[s.v. δικαιόω]{bauer-danker-lexicon-2000}; Deut 25:1; Matt 11:19; Luke 7:35; \cite[ad loc. Matt 11:19]{crossway-esv-2001}
\item A word contributes the least meaning its context requires; one occurrence does not carry the whole concept --- \cite{silva-biblical-words-1994}; \cite{joos-semantic-axiom-1972}; \cite{barr-semantics-2004}
\item an inference step --- \cite[ad loc. Jas 2:21]{moo-james-2000}; \cite{silva-biblical-words-1994}; \cite{barr-semantics-2004}
\end{itemize}
\item \emph{James’s faith (the faith James calls dead)}, weakest at \emph{well supported}:
\begin{itemize}
\item The faith alone that James 2:24 denies justifies is the dead faith of 2:14–26 --- Jas 2:14-26; \cite[Trent, Session VI (1547), ch. 7]{tanner-decrees-1990}; \cite[ad loc. Jas 2:24]{moo-james-2000}
\item an inference step --- \cite[ad loc. Jas 2:24]{moo-james-2000}; \cite[XI.2]{westminster-confession-1647}
\end{itemize}
\item \emph{Revelation 22:11, in the Greek}, weakest at \emph{well supported}:
\begin{itemize}
\item an inference step --- \cite[§16]{pius-divino-afflante-spiritu-1943}; \cite{holmes-sblgnt-2010}; \cite{usccb-nabre-2011}
\end{itemize}
\item \emph{Sirach 18:22, in the Greek}, weakest at \emph{well supported}:
\begin{itemize}
\item The Greek of Sirach 18:22 is about paying a vow: do not wait until death to be justified --- \cite[ad loc. Sir 18:22]{swete-old-testament-in-greek}; \cite[Sirach 18:22]{usccb-nabre-2011}
\item an inference step --- \cite[§16]{pius-divino-afflante-spiritu-1943}; \cite{holmes-sblgnt-2010}; \cite{usccb-nabre-2011}
\end{itemize}
\item \emph{James 2:24, read for the increase of justification}, weakest at \emph{disputed}:
\begin{itemize}
\item James’s ‘justified’ (δικαιόω) means the increase of the justice received --- \cite[Trent, Session VI (1547), ch. 10]{tanner-decrees-1990}
\item an inference step --- \cite[Trent, Session VI (1547), ch. 10]{tanner-decrees-1990}
\end{itemize}
\item \emph{Revelation 22:11, as Trent quotes it}, weakest at \emph{disputed}:
\begin{itemize}
\item Revelation 22:11 says: he that is just, let him be justified still --- \cite[Trent, Session VI (1547), ch. 10]{tanner-decrees-1990}
\item an inference step --- \cite[Trent, Session VI (1547), ch. 10]{tanner-decrees-1990}
\end{itemize}
\item \emph{Sirach 18:22, as Trent quotes it}, weakest at \emph{disputed}:
\begin{itemize}
\item Sirach 18:22 says: be not afraid to be justified even to death --- \cite[Trent, Session VI (1547), ch. 10]{tanner-decrees-1990}
\item an inference step --- \cite[Trent, Session VI (1547), ch. 10]{tanner-decrees-1990}
\end{itemize}
\item \emph{The reward of works (Trent, Session VI, ch. 16)}, weakest at \emph{disputed}:
\begin{itemize}
\item The reward God renders to the works of the justified is rendered to them as merits --- \cite[Trent, Session VI (1547), ch. 16; canon 32]{tanner-decrees-1990}; \cite{dulles-two-languages-1999}
\item an inference step --- \cite[Trent, Session VI (1547), ch. 16, canon 32]{tanner-decrees-1990}; \cite[XVI.5, XVI.6]{westminster-confession-1647}
\end{itemize}
\item \emph{James’s word, read through Paul’s (the renewal reading)}, weakest at \emph{disputed}:
\begin{itemize}
\item Paul’s word ‘justify’ means the renewal of the inward man, not only a verdict --- \cite[Trent, Session VI (1547), Decree on Justification, chs. 7–8]{tanner-decrees-1990}; \cite{vanlandingham-judgment-justification-2006}
\item James’s ‘justify’ (δικαιόω) has the sense Paul’s has: the justification of Romans and Galatians --- \cite[Trent, Session VI (1547), ch. 10]{tanner-decrees-1990}
\item an inference step --- \cite[Trent, Session VI (1547), ch. 10]{tanner-decrees-1990}; \cite{vanlandingham-judgment-justification-2006}
\end{itemize}
\item \emph{The reward of works, of grace (Westminster XVI.5, XVI.6)}, weakest at \emph{disputed}:
\begin{itemize}
\item an inference step --- \cite[XVI.5, XVI.6]{westminster-confession-1647}
\end{itemize}
\end{itemize}
\par\medskip\noindent\textbf{\texttt{trent\_\allowbreak{}on\_\allowbreak{}works\_\allowbreak{}not\_\allowbreak{}forced}}\par\nobreak\noindent
\textbf{Nor is it forced.} A consistent set of positions that answers every attack holds the Reformed reconciliation, with the arguments from James’s word and James’s faith, and cannot hold Trent’s reading of what works do with it.

{\small\begin{verbatim}
theorem trent_on_works_not_forced : ¬Framework.SkepticallyAccepted
    jamesDispute.defeats JamesParty.trentOnWorks
-- axioms: propext, Classical.choice, Quot.sound
\end{verbatim}}
\par\medskip\noindent\textbf{\texttt{harmony\allowbreak{}Stands}}\par\nobreak\noindent
Why the Reformed reconciliation can be defended: it answers Trent’s reading of what works do.

\paragraph{\emph{Harmonisation of James 2:24} is accepted on some resolution.}
\begin{itemize}
\item A position holding \emph{James’s word (δικαιόω, by the rule of least meaning)}, \emph{James’s faith (the faith James calls dead)}, \emph{Harmonisation of James 2:24}, \emph{Revelation 22:11, in the Greek}, \emph{Sirach 18:22, in the Greek} can be held.
\item None of \emph{James’s word (δικαιόω, by the rule of least meaning)}, \emph{James’s faith (the faith James calls dead)}, \emph{Harmonisation of James 2:24}, \emph{Revelation 22:11, in the Greek}, \emph{Sirach 18:22, in the Greek} defeats another, and each attack on them is answered from among them:
\begin{itemize}
\item \emph{James 2:24, read by canon 24} defeats \emph{Harmonisation of James 2:24}, and \emph{Harmonisation of James 2:24} defeats \emph{James 2:24, read by canon 24}.
\item \emph{The reward of works (Trent, Session VI, ch. 16)} defeats \emph{Harmonisation of James 2:24}, and \emph{Harmonisation of James 2:24} defeats \emph{The reward of works (Trent, Session VI, ch. 16)}.
\item \emph{James’s word, read through Paul’s (the renewal reading)} defeats \emph{Harmonisation of James 2:24}, and \emph{James’s word (δικαιόω, by the rule of least meaning)} defeats \emph{James’s word, read through Paul’s (the renewal reading)}.
\end{itemize}
\end{itemize}
\par\noindent\textbf{What this rests on.} The reasons state 6 defeats and 57 absences of defeat, each a cell of the defeat table, \texttt{james\allowbreak{}Dispute\_\allowbreak{}defeats}, computed from the two parties' premises and checked by the kernel. Whether an attack survives turns on the attacker's weakest link and on the rating of what it attacks, so a changed rating can change the verdict. The weakest links:
\begin{itemize}
\item \emph{James’s word (δικαιόω, by the rule of least meaning)}, weakest at \emph{well supported}:
\begin{itemize}
\item Outside Paul, ‘justify’ (δικαιόω) declares or vindicates; it does not mean making righteous --- \cite[s.v. δικαιόω]{bauer-danker-lexicon-2000}; Deut 25:1; Matt 11:19; Luke 7:35; \cite[ad loc. Matt 11:19]{crossway-esv-2001}
\item A word contributes the least meaning its context requires; one occurrence does not carry the whole concept --- \cite{silva-biblical-words-1994}; \cite{joos-semantic-axiom-1972}; \cite{barr-semantics-2004}
\item an inference step --- \cite[ad loc. Jas 2:21]{moo-james-2000}; \cite{silva-biblical-words-1994}; \cite{barr-semantics-2004}
\end{itemize}
\item \emph{James’s faith (the faith James calls dead)}, weakest at \emph{well supported}:
\begin{itemize}
\item The faith alone that James 2:24 denies justifies is the dead faith of 2:14–26 --- Jas 2:14-26; \cite[Trent, Session VI (1547), ch. 7]{tanner-decrees-1990}; \cite[ad loc. Jas 2:24]{moo-james-2000}
\item an inference step --- \cite[ad loc. Jas 2:24]{moo-james-2000}; \cite[XI.2]{westminster-confession-1647}
\end{itemize}
\item \emph{Harmonisation of James 2:24}, weakest at \emph{disputed}:
\begin{itemize}
\item an inference step --- \cite[ad loc. Jas 2:24]{moo-james-2000}; \cite[III.xvii.11]{calvin-institutes-1960}
\end{itemize}
\item \emph{Revelation 22:11, in the Greek}, weakest at \emph{well supported}:
\begin{itemize}
\item an inference step --- \cite[§16]{pius-divino-afflante-spiritu-1943}; \cite{holmes-sblgnt-2010}; \cite{usccb-nabre-2011}
\end{itemize}
\item \emph{Sirach 18:22, in the Greek}, weakest at \emph{well supported}:
\begin{itemize}
\item The Greek of Sirach 18:22 is about paying a vow: do not wait until death to be justified --- \cite[ad loc. Sir 18:22]{swete-old-testament-in-greek}; \cite[Sirach 18:22]{usccb-nabre-2011}
\item an inference step --- \cite[§16]{pius-divino-afflante-spiritu-1943}; \cite{holmes-sblgnt-2010}; \cite{usccb-nabre-2011}
\end{itemize}
\item \emph{James 2:24, read for the increase of justification}, weakest at \emph{disputed}:
\begin{itemize}
\item James’s ‘justified’ (δικαιόω) means the increase of the justice received --- \cite[Trent, Session VI (1547), ch. 10]{tanner-decrees-1990}
\item an inference step --- \cite[Trent, Session VI (1547), ch. 10]{tanner-decrees-1990}
\end{itemize}
\item \emph{James 2:24, read by canon 24}, weakest at \emph{disputed}:
\begin{itemize}
\item Good works are a cause of the increase of justification, not merely its fruits --- \cite[Trent, Session VI (1547), canon 24]{tanner-decrees-1990}
\item an inference step --- \cite[Trent, Session VI (1547), canon 24]{tanner-decrees-1990}
\end{itemize}
\item \emph{Revelation 22:11, as Trent quotes it}, weakest at \emph{disputed}:
\begin{itemize}
\item Revelation 22:11 says: he that is just, let him be justified still --- \cite[Trent, Session VI (1547), ch. 10]{tanner-decrees-1990}
\item an inference step --- \cite[Trent, Session VI (1547), ch. 10]{tanner-decrees-1990}
\end{itemize}
\item \emph{Sirach 18:22, as Trent quotes it}, weakest at \emph{disputed}:
\begin{itemize}
\item Sirach 18:22 says: be not afraid to be justified even to death --- \cite[Trent, Session VI (1547), ch. 10]{tanner-decrees-1990}
\item an inference step --- \cite[Trent, Session VI (1547), ch. 10]{tanner-decrees-1990}
\end{itemize}
\item \emph{The reward of works (Trent, Session VI, ch. 16)}, weakest at \emph{disputed}:
\begin{itemize}
\item The reward God renders to the works of the justified is rendered to them as merits --- \cite[Trent, Session VI (1547), ch. 16; canon 32]{tanner-decrees-1990}; \cite{dulles-two-languages-1999}
\item an inference step --- \cite[Trent, Session VI (1547), ch. 16, canon 32]{tanner-decrees-1990}; \cite[XVI.5, XVI.6]{westminster-confession-1647}
\end{itemize}
\item \emph{James’s word, read through Paul’s (the renewal reading)}, weakest at \emph{disputed}:
\begin{itemize}
\item Paul’s word ‘justify’ means the renewal of the inward man, not only a verdict --- \cite[Trent, Session VI (1547), Decree on Justification, chs. 7–8]{tanner-decrees-1990}; \cite{vanlandingham-judgment-justification-2006}
\item James’s ‘justify’ (δικαιόω) has the sense Paul’s has: the justification of Romans and Galatians --- \cite[Trent, Session VI (1547), ch. 10]{tanner-decrees-1990}
\item an inference step --- \cite[Trent, Session VI (1547), ch. 10]{tanner-decrees-1990}; \cite{vanlandingham-judgment-justification-2006}
\end{itemize}
\item \emph{The reward of works, of grace (Westminster XVI.5, XVI.6)}, weakest at \emph{disputed}:
\begin{itemize}
\item an inference step --- \cite[XVI.5, XVI.6]{westminster-confession-1647}
\end{itemize}
\end{itemize}
\par\medskip\noindent\textbf{\texttt{harmony\_\allowbreak{}holds\_\allowbreak{}fruit}}\par\nobreak\noindent
The Reformed reconciliation derives, from Westminster’s proof texts, that works are the fruit of faith and not its ground.

{\small\begin{verbatim}
theorem harmony_holds_fruit : Entails (jamesDispute.node
    JamesParty.harmony).premises (p Claim.worksAreFruitNotGround)
-- axioms: propext, Classical.choice, Quot.sound
\end{verbatim}}
\par\medskip\noindent\textbf{\texttt{trent\allowbreak{}On\allowbreak{}Works\_\allowbreak{}denies\_\allowbreak{}fruit}}\par\nobreak\noindent
Trent’s reading of what works do denies that.

{\small\begin{verbatim}
theorem trentOnWorks_denies_fruit : Entails (jamesDispute.node
    JamesParty.trentOnWorks).premises (∼p Claim.worksAreFruitNotGround)
-- axioms: propext, Classical.choice, Quot.sound
\end{verbatim}}
\par\medskip\noindent\textbf{\texttt{works\allowbreak{}Contested}}\par\nobreak\noindent
\textbf{Whether works are only the fruit of justification is contested, and the weighing confirms it.} The Reformed reconciliation holds it. Trent’s reading of James 2:24 by canon 24 denies it. Each can be defended, and neither is forced. So the rating \texttt{disputed} on that premise (\texttt{works\allowbreak{}Are\allowbreak{}Fruit\allowbreak{}Not\allowbreak{}Ground}) agrees with what the weighing finds.

{\small\begin{verbatim}
def worksContested : Contested jamesDispute (p Claim.worksAreFruitNotGround) :=
  { holder := JamesParty.harmony,
    denier := JamesParty.trentOnWorks,
    holds := harmony_holds_fruit,
    denies := trentOnWorks_denies_fruit,
    holderDefensible := worksContested._proof_1,
    denierDefensible := trent_on_works_defensible }
\end{verbatim}}
\paragraph*{What each position rests on}
\leavevmode\par\nobreak
\par\medskip\noindent\textbf{\texttt{james\_\allowbreak{}warrants}}\par\nobreak\noindent
\textbf{Where each position rests on someone’s word alone.}

Trent’s reading of the reward texts rests on the council’s word at one claim: that the reward is given to merits. The reading through Paul rests on Trent’s word at one claim: that James’s “justified” is Paul’s. For Paul’s word as renewal it has VanLandingham’s argument from the lexicon.

Each of Trent’s other four readings rests on Trent’s word at the reading itself. The two readings of the Greek rest on the word of Pius XII and of Westminster alike, at the principle that the original outweighs any translation.

The argument from James’s word, the argument about James’s faith, and the Reformed reconciliation, which now derives “fruit, not ground” from Westminster’s proof texts, rest on nothing but Scripture and evidence.

{\small\begin{verbatim}
theorem james_warrants : jamesLexicalCase.restingOnAuthority = [] ∧
    jamesFaithCase.restingOnAuthority = [] ∧
    jamesHarmonyCase.restingOnAuthority = [] ∧
    catholicJamesOnTheWord.restingOnAuthority =
    [Claim.jamesJustifyDenotesIncrease] ∧
    catholicJamesOnWorks.restingOnAuthority =
    [Claim.worksCauseIncreaseOfJustification] ∧
    trentOnRevelation.restingOnAuthority = [Claim.rev22_11BeJustifiedStill] ∧
    trentOnSirach.restingOnAuthority = [Claim.sir18_22BeJustifiedToDeath] ∧
    revelationTextCase.restingOnAuthority = [Claim.originalTextDecides] ∧
    sirachTextCase.restingOnAuthority = [Claim.originalTextDecides] ∧
    trentOnReward.restingOnAuthority = [Claim.rewardRenderedToMerits] ∧
    renewalOnTheWord.restingOnAuthority = [Claim.jamesUsesDikaioAsPaul] ∧
    westminsterOnReward.restingOnAuthority = []
-- axioms: propext
\end{verbatim}}
\subsection*{Abraham, counted righteous}
\addcontentsline{toc}{subsection}{Abraham, counted righteous}
Paul’s chief witness for justification by faith is not a verse of his own but a sentence from the story of Abraham:

\begin{quote}
“And he believed the LORD, and he counted it to him as righteousness” (Genesis 15:6).

\end{quote}
Paul quotes it in Romans 4 and in Galatians 3. And he presses a point of timing:

\begin{quote}
“How then was it counted to him? Was it before or after he had been circumcised? It was not after, but before he was circumcised” (Romans 4:10).

\end{quote}
\subsubsection*{What happens next in Genesis 15}
The rest of the chapter is a covenant ceremony, and it is a strange one. God tells Abram to cut animals in two and lay the halves opposite each other. Then:

\begin{quote}
“As the sun was going down, a deep sleep fell on Abram. And behold, dreadful and great darkness fell upon him” (Genesis 15:12).

\end{quote}
\begin{quote}
“When the sun had gone down and it was dark, behold, a smoking fire pot and a flaming torch passed between these pieces” (Genesis 15:17).

\end{quote}
Jeremiah shows what walking between the pieces meant. Those who broke a covenant would be made “like the calf that they cut in two and passed between its parts” (Jeremiah 34:18): whoever passed through called that fate down on himself if he broke his word. In Genesis 15 only the fire and the torch pass through, and Abram is asleep. The letter to the Hebrews draws the conclusion: “when God made a promise to Abraham, since he had no one greater by whom to swear, he swore by himself” (Hebrews 6:13).

On this reading the covenant binds God alone, and nothing Abram does is a condition of the promise that was counted to him. Reformed readers of the covenants read it so: O. Palmer Robertson and Meredith Kline, and Gordon Wenham in his commentary on Genesis. Moshe Weinfeld, comparing the covenants of the ancient Near East, set the covenant with Abraham beside the royal grants that a king gave and did not take back. None of the four has been read for this chapter, so each is cited as a whole.

Weinfeld also gives the strongest reason to doubt the step. For him, as Gary Knoppers quotes him, a grant is “a reward for loyalty and good deeds already performed”, and he ties Abraham’s to Genesis 26:5. Knoppers himself argues that the division is overdrawn, and in the other direction: “Land grants are not predominantly unconditional.” So his critique does not help the step; but he adds that “Continuing or future loyalty can be assumed or stipulated, even though the basic promise may be unaffected by (dis)loyalty”. The next sections test whether Weinfeld’s reason reaches the step.

\subsubsection*{The other way of reading Abraham}
Genesis does not stop at chapter 15. When Abraham has offered Isaac, the angel of the LORD says:

\begin{quote}
“By myself I have sworn, declares the LORD, because you have done this and have not withheld your son, your only son, I will surely bless you” (Genesis 22:16–17).

\end{quote}
And to Isaac: the blessing comes “because Abraham obeyed my voice and kept my charge, my commandments, my statutes, and my laws” (Genesis 26:5).

Jewish writers before Paul read Genesis 15:6 in the light of Genesis 22. 1 Maccabees asks, in the Catholic Church’s own translation (NABRE):

\begin{quote}
“Was not Abraham found faithful in trial, and it was credited to him as righteousness?” (1 Maccabees 2:52).

\end{quote}
Sirach says that “when tested [he] was found loyal” (44:20). The same translation cross-refers the first to Genesis 15:6 and to Genesis 22, and glosses the second as “Abraham’s willingness to sacrifice his son Isaac”. On this reading the righteousness counted to Abraham is tied to his obedience under testing. James draws on the same story: Abraham was “justified by works when he offered up his son Isaac on the altar” (James 2:21).

These are not books of the Protestant canon, and nothing here treats them as Scripture. They are witnesses to how Abraham was read in Paul’s day, which is why Paul argues as he does in Romans 4.

\subsubsection*{Does Weinfeld’s grant reach what was counted?}
Set Weinfeld’s argument out step by step:

\begin{itemize}
  \item the covenant of Genesis 15 is a royal grant, one-sided;
  \item a grant rewarded loyalty already shown;
  \item the blessing came “because you have obeyed my voice” (Genesis 22:18; 26:5);
  \item so the grant to Abraham rewarded his obedience.
\end{itemize}
Every line may be true, and the step from Genesis 15 still stands. The step is about something else: what was counted to Abraham as righteousness at 15:6, and whether anything he did was its condition. A grant is not a counting. Weinfeld’s premises are kept alike by a reading on which Abraham was counted righteous without works and by one on which he was not, so they decide nothing about it (\texttt{weinfeld\_\allowbreak{}leaves\_\allowbreak{}the\_\allowbreak{}counting\_\allowbreak{}open}). To reach the step, the argument needs one more premise: that what was counted at 15:6 was credited for Abraham’s obedience. Weinfeld, as he has been read here, does not state it, so his objection does not dissent from the step at all (\texttt{weinfeld\_\allowbreak{}does\_\allowbreak{}not\_\allowbreak{}reach\_\allowbreak{}the\_\allowbreak{}counting}). Treating “the grant rewarded his obedience” as if it said “his righteousness was credited for his obedience” runs two different claims together.

Abraham Kuruvilla, writing in 2012, takes the same “because” further. He grants that “faithlessness characterized Abraham’s response to God” in the Hagar episode, and holds that in Genesis 22, “Contingent upon his obedience, every aspect of the earlier promises to Abraham is now ‘augmented and guaranteed’”. He answers Calvin, who read the oath as “nothing more than confirmation”. But his contingency, too, is of the blessing sworn in Genesis 22, and he calls it “more like an unexpected bonus, which, of course, is what grace is all about”. It does not reach what was counted at 15:6 (\texttt{kuruvilla\_\allowbreak{}does\_\allowbreak{}not\_\allowbreak{}reach\_\allowbreak{}the\_\allowbreak{}counting}). R. W. L. Moberly, as Kuruvilla quotes him, puts the point exactly: “It is not that divine promise has become contingent upon Abraham’s obedience, but that Abraham’s obedience has been incorporated into the divine promise.”

\subsubsection*{The missing premise, where it is stated}
1 Maccabees states it. “Was not Abraham found faithful in trial, and it was credited to him as righteousness?” reads the counting of Genesis 15:6 as credited for the testing of Genesis 22. With that premise, the argument does deny the step, and it is an argument and not a bare disagreement (\texttt{maccabees\_\allowbreak{}unstruck\_\allowbreak{}counts}).

But Genesis reads otherwise, and in plain words. It counts Abram’s believing at 15:6:

\begin{quote}
“And he believed the LORD, and he counted it to him as righteousness” (Genesis 15:6).

\end{quote}
And it places the testing later: “After these things God tested Abraham” (Genesis 22:1). The counting came first, decades before Mount Moriah. In between, Abraham and Sarah took matters into their own hands with Hagar (Genesis 16), and God still confirmed the covenant to him (Genesis 17). Paul argues from the same kind of order about circumcision (Romans 4:10), and Calvin made this very argument against reading James 2 as the ground of Genesis 15:6: “The imputation of righteousness which Moses mentions, preceded more than thirty years the work by which they would have Abraham to have been justified.” No published reader was found who argues from the Hagar episode in particular; that part of the argument is made here, from the texts. So the reading 1 Maccabees gives Genesis 15:6 is refuted by what Genesis reads, and with it struck, 1 Maccabees no longer dissents from the step (\texttt{genesis\_\allowbreak{}strikes\_\allowbreak{}the\_\allowbreak{}maccabees\_\allowbreak{}reading}). What remains true in it is that Abraham was found faithful when tested. James says the same, and calls that testing the fulfilment of the earlier verse: “and the Scripture was fulfilled that says, ‘Abraham believed God, and it was counted to him as righteousness’” (James 2:23).

\subsubsection*{What the checking finds}
\begin{itemize}
  \item \textbf{Read from the ceremony, Abraham was counted righteous with no work as a condition} (\texttt{genesis\allowbreak{}Fifteen\_\allowbreak{}establishes}).
  \item \textbf{Read from his testing, he was not} (\texttt{second\allowbreak{}Temple\allowbreak{}Abraham\_\allowbreak{}establishes}), but only through the reading of 15:6 that Genesis refutes.
  \item \textbf{The texts of Genesis 15 and 22, quoted alone, do not decide} (\texttt{genesis\_\allowbreak{}leaves\_\allowbreak{}abraham\_\allowbreak{}open}): it takes the order of the narrative to strike the other reading.
  \item \textbf{No dissent found survives, so the step stands at \texttt{well\allowbreak{}Supported} under every standard} (\texttt{abraham\_\allowbreak{}step\_\allowbreak{}ratings}). The rating is worked out, not assigned: it would fall again if a verified reader argued that the counting of 15:6 was conditioned on Abraham’s obedience in a way the order of Genesis does not refute, for example on his leaving Haran (Genesis 12:4). Hebrews calls that obedience itself an act of faith: “By faith Abraham obeyed when he was called” (Hebrews 11:8).
\end{itemize}
The Reformed answer to Genesis 22 is the answer this library already gives to James: Abraham’s offering of Isaac shows the faith that was counted to him, years before. And the chapter “Does justification have degrees?” adds a reason of its own. If justification is a verdict, which has no degrees, then James’s “justified” at Genesis 22 cannot mean the \emph{increase} of the justice Abraham already had (\texttt{no\_\allowbreak{}degrees\_\allowbreak{}answers\_\allowbreak{}james\_\allowbreak{}on\_\allowbreak{}abraham}). That reason rests, as before, on the Reformed definition of justification, which Trent denies.

\subsubsection*{What this does not show}
Some readers press a further question. Scott Hahn and Paul Williamson read Genesis 15, 17 and 22 as stages of one relationship, so that “the promise” may name the one-sided oath of Genesis 15 or the promise as Genesis 22 and Hebrews 6 give it, sworn after the test. Hahn takes the second: the covenant Paul means in Galatians 3:17 is “the Abrahamic covenant in its final form, as ratified most solemnly by God’s oath after the Aqedah”, and “the binding of Isaac ... merits from God a covenant oath to bless the ἔθνη [the nations] through Abraham’s ‘seed’”. But what he says the binding of Isaac merited is the oath and the blessing of the nations, not the righteousness counted at 15:6, which he does not discuss in what was read of him; he even calls that oath “unilaterally” sworn. So on either reading of “the promise”, nothing reaches what was counted (\texttt{hahn\_\allowbreak{}does\_\allowbreak{}not\_\allowbreak{}reach\_\allowbreak{}the\_\allowbreak{}counting}). Jeffrey Niehaus, who reads one covenant cut in Genesis 15, says so in terms: Abraham is the father of “those who, like him, would be justified by faith without circumcision (as in Genesis 15)”, and the works that followed “showed that his faith was real”. Williamson’s own text has not been read, and he is not encoded.

It does not show that Abraham’s obedience did not matter. Every reader cited here holds that it did. The question is whether it was the condition of the righteousness counted to him at Genesis 15:6, or its fruit. Nor does it show that Trent denies the reading of the ceremony: Trent holds that “none of those things which precede justification, whether faith or works, merit the grace itself of justification” (Session VI, ch. 8). Its disagreement is over what follows.

\subsubsection*{How this is checked}
\paragraph*{The reading tied to Genesis 22, written out first}
\leavevmode\par\nobreak
\par\medskip\noindent\textbf{\texttt{second\allowbreak{}Temple\allowbreak{}On\allowbreak{}Abraham\allowbreak{}Source}}\par\nobreak\noindent
The Second Temple reading’s step: Abraham was “found faithful in trial, and it was credited to him as righteousness” (1 Maccabees 2:52, NABRE), so the counting of Genesis 15:6 is tied to his obedience in Genesis 22, and was not without works. Rated \texttt{disputed}: Paul accepts Genesis 15:6 and argues the opposite from its timing (Romans 4:10).

{\small\begin{verbatim}
def secondTempleOnAbrahamSource : Source :=
  {
    primary :=
      Reference.work nabre
        (Bib.Locus.adLoc
          { book := Book.firstMaccabees, chapter := 2, verse := 52 }),
    supporting :=
      [Reference.work nabre
          (Bib.Locus.adLoc
            { book := Book.sirach, chapter := 44, verse := 20 })],
    tradition := Tradition.secondTempleJewish,
    confidence := Confidence.disputed }
\end{verbatim}}
\par\medskip\noindent\textbf{\texttt{second\allowbreak{}Temple\allowbreak{}Abraham\allowbreak{}Line}} --- Abraham found faithful when tested (1 Maccabees 2:52; Sirach 44:20)\par\nobreak\noindent
\textbf{Abraham, read from his testing.} From Genesis 15:6 and 22:16–18, with 1 Maccabees 2:52 and Sirach 44:20, and the reading 1 Maccabees gives Genesis 15:6, that what was counted was credited for his faithfulness when tested, the righteousness counted to Abraham is tied to his obedience.

\begin{derivation}
  \item[(1)] \(\pvar{TestimonyArguments.SolaFide}{112}\)\gloss{Genesis 15:6: Abram believed the LORD, and he counted it to him as righteousness}
  \item[(2)] \(\pvar{TestimonyArguments.SolaFide}{117}\)\gloss{Genesis 22:16–18, 26:5: blessed because Abraham obeyed}
  \item[(3)] \(\pvar{TestimonyArguments.SolaFide}{118}\)\gloss{1 Maccabees 2:52, Sirach 44:20: Abraham found faithful in trial, and it was credited to him as righteousness}
  \breakrow{(4)} \(\pvar{TestimonyArguments.SolaFide}{121}\)\gloss{The righteousness counted at Genesis 15:6 was credited for Abraham’s faithfulness when tested}
  \item[(5)] \((\pvar{TestimonyArguments.SolaFide}{112} \land \pvar{TestimonyArguments.SolaFide}{117} \land \pvar{TestimonyArguments.SolaFide}{118} \land \pvar{TestimonyArguments.SolaFide}{121}) \rightarrow \lnot \pvar{TestimonyArguments.SolaFide}{116}\)
  \item[\(\vdash\)] \(\lnot \pvar{TestimonyArguments.SolaFide}{116}\)\gloss{\emph{at issue, the denial of:} Abraham was counted righteous on believing, with no work as a condition}
\end{derivation}
\par\medskip\noindent\textbf{\texttt{second\allowbreak{}Temple\allowbreak{}Abraham}} --- Abraham, counted righteous when tested\par\nobreak\noindent
\textbf{The reading tied to Genesis 22}, as a position.

\begin{derivation}
  \item[(1)] \(\pvar{TestimonyArguments.SolaFide}{112}\)\gloss{Genesis 15:6: Abram believed the LORD, and he counted it to him as righteousness}
  \item[(2)] \(\pvar{TestimonyArguments.SolaFide}{117}\)\gloss{Genesis 22:16–18, 26:5: blessed because Abraham obeyed}
  \item[(3)] \(\pvar{TestimonyArguments.SolaFide}{118}\)\gloss{1 Maccabees 2:52, Sirach 44:20: Abraham found faithful in trial, and it was credited to him as righteousness}
  \breakrow{(4)} \(\pvar{TestimonyArguments.SolaFide}{121}\)\gloss{The righteousness counted at Genesis 15:6 was credited for Abraham’s faithfulness when tested}
  \item[(5)] \((\pvar{TestimonyArguments.SolaFide}{112} \land \pvar{TestimonyArguments.SolaFide}{117} \land \pvar{TestimonyArguments.SolaFide}{118} \land \pvar{TestimonyArguments.SolaFide}{121}) \rightarrow \lnot \pvar{TestimonyArguments.SolaFide}{116}\)
  \item[\(\vdash\)] \(\lnot \pvar{TestimonyArguments.SolaFide}{116}\)\gloss{\emph{at issue, the denial of:} Abraham was counted righteous on believing, with no work as a condition}
  \keepwithnext
\end{derivation}
\emph{No premise here rests on scripture alone.}
\paragraph*{The reading of the ceremony}
\leavevmode\par\nobreak
\par\medskip\noindent\textbf{\texttt{genesis\allowbreak{}Fifteen\allowbreak{}Source}}\par\nobreak\noindent
The Reformed step: Abraham believed and it was counted to him (Genesis 15:6), before he was circumcised (Romans 4:10), and the covenant that followed bound God alone (15:12–17). So nothing he did was a condition of what was counted to him. Its support is Paul’s own argument from the same verse and its timing (Romans 4:3–12), rated \texttt{well\allowbreak{}Supported}. The rating the page gives is worked out from the dissents found against it (\texttt{abraham\_\allowbreak{}step\_\allowbreak{}ratings}).

{\small\begin{verbatim}
def genesisFifteenSource : Source
\end{verbatim}}
\par\medskip\noindent\textbf{\texttt{genesis\allowbreak{}Fifteen\allowbreak{}Line}} --- Abraham, counted righteous before the covenant’s sign (Genesis 15; Romans 4)\par\nobreak\noindent
\textbf{Abraham, read from the ceremony.} From Genesis 15:6, Romans 4:10, the ceremony of Genesis 15:12–17 and the reading of it as God’s oath alone, Abraham was counted righteous with no work as a condition.

\begin{derivation}
  \item[(1)] \(\pvar{TestimonyArguments.SolaFide}{112}\)\gloss{Genesis 15:6: Abram believed the LORD, and he counted it to him as righteousness}
  \item[(2)] \(\pvar{TestimonyArguments.SolaFide}{114}\)\gloss{Romans 4:10: righteousness was counted to Abraham before he was circumcised}
  \item[(3)] \(\pvar{TestimonyArguments.SolaFide}{113}\)\gloss{Genesis 15:12, 17: Abram asleep; a fire pot and a torch pass between the pieces}
  \breakrow{(4)} \(\pvar{TestimonyArguments.SolaFide}{115}\)\gloss{The covenant of Genesis 15 binds God alone, by an oath on himself}
  \item[(5)] \((\pvar{TestimonyArguments.SolaFide}{112} \land \pvar{TestimonyArguments.SolaFide}{114} \land \pvar{TestimonyArguments.SolaFide}{113} \land \pvar{TestimonyArguments.SolaFide}{115}) \rightarrow \pvar{TestimonyArguments.SolaFide}{116}\)
  \item[\(\vdash\)] \(\pvar{TestimonyArguments.SolaFide}{116}\)\gloss{\emph{at issue:} Abraham was counted righteous on believing, with no work as a condition}
\end{derivation}
\par\medskip\noindent\textbf{\texttt{genesis\allowbreak{}Fifteen}} --- Abraham, counted righteous by God’s oath alone\par\nobreak\noindent
\textbf{The reading of the ceremony}, as a position.

\begin{derivation}
  \item[(1)] \(\pvar{TestimonyArguments.SolaFide}{112}\)\gloss{Genesis 15:6: Abram believed the LORD, and he counted it to him as righteousness}
  \item[(2)] \(\pvar{TestimonyArguments.SolaFide}{114}\)\gloss{Romans 4:10: righteousness was counted to Abraham before he was circumcised}
  \item[(3)] \(\pvar{TestimonyArguments.SolaFide}{113}\)\gloss{Genesis 15:12, 17: Abram asleep; a fire pot and a torch pass between the pieces}
  \breakrow{(4)} \(\pvar{TestimonyArguments.SolaFide}{115}\)\gloss{The covenant of Genesis 15 binds God alone, by an oath on himself}
  \item[(5)] \((\pvar{TestimonyArguments.SolaFide}{112} \land \pvar{TestimonyArguments.SolaFide}{114} \land \pvar{TestimonyArguments.SolaFide}{113} \land \pvar{TestimonyArguments.SolaFide}{115}) \rightarrow \pvar{TestimonyArguments.SolaFide}{116}\)
  \item[\(\vdash\)] \(\pvar{TestimonyArguments.SolaFide}{116}\)\gloss{\emph{at issue:} Abraham was counted righteous on believing, with no work as a condition}
  \keepwithnext
\end{derivation}
\emph{No premise here rests on scripture alone.}
\paragraph*{Genesis 22, and degrees}
\leavevmode\par\nobreak
\par\medskip\noindent\textbf{\texttt{james\allowbreak{}Increase\allowbreak{}Means\allowbreak{}Degrees\allowbreak{}Source}}\par\nobreak\noindent
The step from James’s word to degrees: if James’s “justified” at Genesis 22 means the increase of the justice Abraham already had, then justification has degrees. Rated \texttt{consensus}: it is what “increase” says.

{\small\begin{verbatim}
def jamesIncreaseMeansDegreesSource : Source :=
  { primary := Reference.work waterworthTrent (Bib.Locus.page 37),
    tradition := Tradition.romanCatholic, confidence := Confidence.consensus }
\end{verbatim}}
\par\medskip\noindent\textbf{\texttt{james\allowbreak{}Increase\allowbreak{}Means\allowbreak{}Degrees\allowbreak{}Line}} --- James’s increase, as degrees\par\nobreak\noindent
\textbf{James’s increase, as degrees.} If James’s “justified” means an increase of justice, justification has degrees.

\begin{derivation}
  \item[(1)] \(\pvar{TestimonyArguments.SolaFide}{86}\)\gloss{James’s ‘justified’ (δικαιόω) means the increase of the justice received}
  \item[(2)] \(\pvar{TestimonyArguments.SolaFide}{86} \rightarrow \pvar{TestimonyArguments.SolaFide}{109}\)
  \item[\(\vdash\)] \(\pvar{TestimonyArguments.SolaFide}{109}\)\gloss{\emph{at issue:} Justification has degrees: the justified are still further justified}
\end{derivation}
\par\medskip\noindent\textbf{\texttt{verdict\allowbreak{}On\allowbreak{}James\allowbreak{}Increase}} --- Justification as a verdict, on James at Genesis 22\par\nobreak\noindent
\textbf{The verdict position, on James at Genesis 22.} The Reformed position on degrees, with what an increase in James’s word would mean. It concludes that James’s “justified” does not mean an increase of justice.

\begin{derivation}
  \item[(1)] \(\pvar{TestimonyArguments.SolaFide}{110}\)\gloss{Deuteronomy 25:1, Proverbs 17:15: judges acquit the innocent and condemn the guilty}
  \item[(2)] \(\pvar{TestimonyArguments.SolaFide}{48}\)\gloss{Romans 8:33–34: it is God who justifies — who is to condemn?}
  \item[(3)] \((\pvar{TestimonyArguments.SolaFide}{110} \land \pvar{TestimonyArguments.SolaFide}{48}) \rightarrow \pvar{TestimonyArguments.SolaFide}{111}\)\gloss{\emph{would give:} A verdict acquits or condemns, and has no degrees}
  \breakrow{(4)} \(\pvar{TestimonyArguments.SolaFide}{30}\)\gloss{Justification is a verdict only: pardon and righteousness counted, not poured in}
  \breakrow{(5)} \((\pvar{TestimonyArguments.SolaFide}{30} \land \pvar{TestimonyArguments.SolaFide}{111}) \rightarrow \lnot \pvar{TestimonyArguments.SolaFide}{109}\)\gloss{\emph{would deny:} Justification has degrees: the justified are still further justified}
  \breakrow{(6)} \(\pvar{TestimonyArguments.SolaFide}{32} \rightarrow \pvar{TestimonyArguments.SolaFide}{109}\)\gloss{\emph{would give:} Justification has degrees: the justified are still further justified}
  \item[(7)] \(\pvar{TestimonyArguments.SolaFide}{86} \rightarrow \pvar{TestimonyArguments.SolaFide}{109}\)\gloss{\emph{would give:} Justification has degrees: the justified are still further justified}
  \item[\(\vdash\)] \(\lnot \pvar{TestimonyArguments.SolaFide}{86}\)\gloss{\emph{at issue, the denial of:} James’s ‘justified’ (δικαιόω) means the increase of the justice received}
  \keepwithnext
\end{derivation}
\emph{Grounded in scripture alone:} Romans 8:33–34: it is God who justifies — who is to condemn?.
\paragraph*{The readings}
\leavevmode\par\nobreak
\par\medskip\noindent\textbf{\texttt{genesis\allowbreak{}Fifteen\allowbreak{}Reading}}\par\nobreak\noindent
The reading of the ceremony: every text holds, the covenant binds God alone, and Abraham was counted righteous with no work as a condition. The reading tied to his testing, from 1 Maccabees and Sirach, is not held.

{\small\begin{verbatim}
def genesisFifteenReading : Valuation Claim :=
  fun a =>
    match a with
    | Claim.abrahamFoundFaithfulWhenTested => False
    | Claim.countingWasForTheTesting => False
    | Claim.justificationHasDegrees => False
    | Claim.worksMeritIncreaseOfJustification => False
    | Claim.jamesJustifyDenotesIncrease => False
    | Claim.justificationIncludesSanctification => False
    | x => True
\end{verbatim}}
\par\medskip\noindent\textbf{\texttt{second\allowbreak{}Temple\allowbreak{}Abraham\allowbreak{}Reading}}\par\nobreak\noindent
The reading tied to Genesis 22: every text holds, and the righteousness counted to Abraham is tied to his testing. The reading of the ceremony as God’s oath alone is not held.

{\small\begin{verbatim}
def secondTempleAbrahamReading : Valuation Claim :=
  fun a =>
    match a with
    | Claim.abrahamCountedRighteousWithoutWorks => False
    | Claim.gen15CovenantIsGodsOathAlone => False
    | x => True
\end{verbatim}}
\paragraph*{Results}
\leavevmode\par\nobreak
\par\medskip\noindent\textbf{\texttt{genesis\allowbreak{}Fifteen\_\allowbreak{}establishes}}\par\nobreak\noindent
\textbf{Read from the ceremony, Abraham was counted righteous with no work as a condition.} Genesis 15:6, Romans 4:10, the ceremony of Genesis 15:12–17, and the reading of it as God’s oath alone.

It does not show that the reading of the ceremony is beyond dispute. It is rated \texttt{disputed}, and Genesis 22 is read the other way.

{\small\begin{verbatim}
theorem genesisFifteen_establishes : Establishes genesisFifteen
-- axioms: propext, Quot.sound
\end{verbatim}}
\par\medskip\noindent\textbf{\texttt{genesis\allowbreak{}Fifteen\_\allowbreak{}is\_\allowbreak{}satisfiable}}\par\nobreak\noindent
The reading of the ceremony can be held without contradiction.

{\small\begin{verbatim}
theorem genesisFifteen_is_satisfiable : Satisfiable genesisFifteen.premises
-- axioms: propext, Quot.sound
\end{verbatim}}
\par\medskip\noindent\textbf{\texttt{second\allowbreak{}Temple\allowbreak{}Abraham\_\allowbreak{}establishes}}\par\nobreak\noindent
\textbf{Read from his testing, Abraham was not counted righteous without works.} Genesis 15:6 and 22:16–18, with 1 Maccabees 2:52 and Sirach 44:20.

It does not show that 1 Maccabees or Sirach is Scripture. They are witnesses to how Abraham was read before Paul.

{\small\begin{verbatim}
theorem secondTempleAbraham_establishes : Establishes secondTempleAbraham
-- axioms: propext, Quot.sound
\end{verbatim}}
\par\medskip\noindent\textbf{\texttt{second\allowbreak{}Temple\allowbreak{}Abraham\_\allowbreak{}is\_\allowbreak{}satisfiable}}\par\nobreak\noindent
The reading tied to Genesis 22 can be held without contradiction.

{\small\begin{verbatim}
theorem secondTempleAbraham_is_satisfiable : Satisfiable
    secondTempleAbraham.premises
-- axioms: propext, Quot.sound
\end{verbatim}}
\par\medskip\noindent\textbf{\texttt{genesis\_\allowbreak{}leaves\_\allowbreak{}abraham\_\allowbreak{}open}}\par\nobreak\noindent
\textbf{The texts of Genesis alone do not decide.} Genesis 15:6, the ceremony of 15:12–17, and Genesis 22:16–18 are kept alike by the reading of the ceremony, on which Abraham was counted righteous with no work as a condition, and by the reading tied to his testing, on which he was not.

{\small\begin{verbatim}
theorem genesis_leaves_abraham_open : Independent [p
    Claim.gen15_6AbramBelieved, p Claim.gen15_17GodAlonePasses, p
    Claim.gen22_16BecauseYouObeyed] (p
    Claim.abrahamCountedRighteousWithoutWorks)
-- axioms: propext, Quot.sound
\end{verbatim}}
\par\medskip\noindent\textbf{\texttt{no\_\allowbreak{}degrees\_\allowbreak{}answers\_\allowbreak{}james\_\allowbreak{}on\_\allowbreak{}abraham}}\par\nobreak\noindent
\textbf{If justification has no degrees, James’s “justified” at Genesis 22 is not an increase.} On the Reformed position on degrees, Abraham’s being “justified by works when he offered up his son Isaac” (James 2:21) cannot mean the increase of the justice he already had. That leaves the Reformed reading of James’s word, “shown to be righteous”.

It does not show that justification has no degrees. That rests on the Reformed definition, which Trent denies (\texttt{degrees\_\allowbreak{}turn\_\allowbreak{}on\_\allowbreak{}the\_\allowbreak{}definition}).

{\small\begin{verbatim}
theorem no_degrees_answers_james_on_abraham : Establishes
    verdictOnJamesIncrease
-- axioms: propext, Quot.sound
\end{verbatim}}
\par\medskip\noindent\textbf{\texttt{verdict\allowbreak{}On\allowbreak{}James\allowbreak{}Increase\_\allowbreak{}is\_\allowbreak{}satisfiable}}\par\nobreak\noindent
The verdict position on James at Genesis 22 can be held without contradiction.

{\small\begin{verbatim}
theorem verdictOnJamesIncrease_is_satisfiable : Satisfiable
    verdictOnJamesIncrease.premises
-- axioms: propext, Quot.sound
\end{verbatim}}
\paragraph*{Weinfeld’s grant, and the bridge it would need}
\leavevmode\par\nobreak
\par\medskip\noindent\textbf{\texttt{weinfeld\allowbreak{}Grant\allowbreak{}Source}}\par\nobreak\noindent
Weinfeld’s step: the covenant of Genesis 15 is a grant, a grant rewards loyalty already shown, and Genesis 22:18 and 26:5 give the blessing “because” Abraham obeyed. So the grant rewarded his obedience. Rated \texttt{disputed}: Knoppers argues that Weinfeld’s division of grants from treaties is overdrawn.

{\small\begin{verbatim}
def weinfeldGrantSource : Source :=
  { primary := Reference.work weinfeldCovenantOfGrant (Bib.Locus.page 185),
    tradition := Tradition.criticalScholarship,
    confidence := Confidence.disputed }
\end{verbatim}}
\par\medskip\noindent\textbf{\texttt{weinfeld\allowbreak{}Grant\allowbreak{}Line}} --- Abraham’s covenant as a royal grant (Weinfeld)\par\nobreak\noindent
\textbf{Weinfeld’s grant.} From the reading of Genesis 15 as one-sided, what a grant was, and Genesis 22:18 and 26:5, the grant to Abraham rewarded his obedience.

\begin{derivation}
  \item[(1)] \(\pvar{TestimonyArguments.SolaFide}{115}\)\gloss{The covenant of Genesis 15 binds God alone, by an oath on himself}
  \item[(2)] \(\pvar{TestimonyArguments.SolaFide}{119}\)\gloss{A royal grant rewarded loyalty already shown}
  \item[(3)] \(\pvar{TestimonyArguments.SolaFide}{117}\)\gloss{Genesis 22:16–18, 26:5: blessed because Abraham obeyed}
  \item[(4)] \((\pvar{TestimonyArguments.SolaFide}{115} \land \pvar{TestimonyArguments.SolaFide}{119} \land \pvar{TestimonyArguments.SolaFide}{117}) \rightarrow \pvar{TestimonyArguments.SolaFide}{120}\)
  \item[\(\vdash\)] \(\pvar{TestimonyArguments.SolaFide}{120}\)\gloss{\emph{at issue:} The grant to Abraham rewarded his obedience (Genesis 22:18; 26:5)}
\end{derivation}
\par\medskip\noindent\textbf{\texttt{weinfeld\allowbreak{}On\allowbreak{}Abraham}} --- Weinfeld, on Abraham’s grant\par\nobreak\noindent
\textbf{Weinfeld, on Abraham}, as his own claims carry him: the grant to Abraham rewarded his obedience.

\begin{derivation}
  \item[(1)] \(\pvar{TestimonyArguments.SolaFide}{115}\)\gloss{The covenant of Genesis 15 binds God alone, by an oath on himself}
  \item[(2)] \(\pvar{TestimonyArguments.SolaFide}{119}\)\gloss{A royal grant rewarded loyalty already shown}
  \item[(3)] \(\pvar{TestimonyArguments.SolaFide}{117}\)\gloss{Genesis 22:16–18, 26:5: blessed because Abraham obeyed}
  \item[(4)] \((\pvar{TestimonyArguments.SolaFide}{115} \land \pvar{TestimonyArguments.SolaFide}{119} \land \pvar{TestimonyArguments.SolaFide}{117}) \rightarrow \pvar{TestimonyArguments.SolaFide}{120}\)
  \item[\(\vdash\)] \(\pvar{TestimonyArguments.SolaFide}{120}\)\gloss{\emph{at issue:} The grant to Abraham rewarded his obedience (Genesis 22:18; 26:5)}
  \keepwithnext
\end{derivation}
\emph{No premise here rests on scripture alone.}
\par\medskip\noindent\textbf{\texttt{weinfeld\allowbreak{}With\allowbreak{}Works\allowbreak{}Reading}}\par\nobreak\noindent
Weinfeld’s reading, with the counting of Genesis 15:6 tied to Abraham’s works: the grant rewarded his obedience, and he was not counted righteous without works.

{\small\begin{verbatim}
def weinfeldWithWorksReading : Valuation Claim :=
  fun a =>
    match a with
    | Claim.abrahamCountedRighteousWithoutWorks => False
    | x => True
\end{verbatim}}
\par\medskip\noindent\textbf{\texttt{weinfeld\allowbreak{}With\allowbreak{}Faith\allowbreak{}Reading}}\par\nobreak\noindent
Weinfeld’s reading, with the counting of Genesis 15:6 left as the text gives it: the grant rewarded his obedience, and what was counted at 15:6 was his believing, with no work as its condition.

{\small\begin{verbatim}
def weinfeldWithFaithReading : Valuation Claim :=
  fun a =>
    match a with
    | Claim.countingWasForTheTesting => False
    | x => True
\end{verbatim}}
\par\medskip\noindent\textbf{\texttt{genesis\allowbreak{}Reads\allowbreak{}The\allowbreak{}Order}}\par\nobreak\noindent
\textbf{What Genesis reads}: Genesis 15:6 counts Abram’s believing as righteousness, and the testing comes afterwards, “after these things” (Genesis 22:1). So what was counted at 15:6 was not credited for the faithfulness of Genesis 22. Paul argues from the same order in Romans 4:10, of circumcision.

{\small\begin{verbatim}
def genesisReadsTheOrder : Refutation Claim
\end{verbatim}}
\paragraph*{Kuruvilla on the Aqedah}
\leavevmode\par\nobreak
\par\medskip\noindent\textbf{\texttt{kuruvilla\allowbreak{}Source}}\par\nobreak\noindent
Kuruvilla’s step: the oath of Genesis 22:16–18 is sworn “because you have done this” and “because you have obeyed my voice”, so the blessing it swears is augmented because of Abraham’s obedience. Rated \texttt{disputed}: Calvin accepts the text and reads the oath as “nothing more than confirmation”.

{\small\begin{verbatim}
def kuruvillaSource : Source :=
  { primary := Reference.work kuruvillaAqedah (Bib.Locus.page 507),
    tradition := Tradition.christianHistoricalGrammatical,
    confidence := Confidence.disputed }
\end{verbatim}}
\par\medskip\noindent\textbf{\texttt{kuruvilla\allowbreak{}Line}} --- The blessing of Genesis 22, augmented for obedience (Kuruvilla)\par\nobreak\noindent
\textbf{The Aqedah’s blessing, augmented for obedience} (Kuruvilla).

\begin{derivation}
  \item[(1)] \(\pvar{TestimonyArguments.SolaFide}{112}\)\gloss{Genesis 15:6: Abram believed the LORD, and he counted it to him as righteousness}
  \item[(2)] \(\pvar{TestimonyArguments.SolaFide}{117}\)\gloss{Genesis 22:16–18, 26:5: blessed because Abraham obeyed}
  \item[(3)] \((\pvar{TestimonyArguments.SolaFide}{112} \land \pvar{TestimonyArguments.SolaFide}{117}) \rightarrow \pvar{TestimonyArguments.SolaFide}{122}\)
  \item[\(\vdash\)] \(\pvar{TestimonyArguments.SolaFide}{122}\)\gloss{\emph{at issue:} The blessing of Genesis 22:16–18 was augmented because of Abraham’s obedience}
\end{derivation}
\par\medskip\noindent\textbf{\texttt{kuruvilla\allowbreak{}On\allowbreak{}The\allowbreak{}Aqedah}} --- Kuruvilla, on the Aqedah\par\nobreak\noindent
\textbf{Kuruvilla, on the Aqedah}: the blessing of Genesis 22 was augmented because of Abraham’s obedience.

\begin{derivation}
  \item[(1)] \(\pvar{TestimonyArguments.SolaFide}{112}\)\gloss{Genesis 15:6: Abram believed the LORD, and he counted it to him as righteousness}
  \item[(2)] \(\pvar{TestimonyArguments.SolaFide}{117}\)\gloss{Genesis 22:16–18, 26:5: blessed because Abraham obeyed}
  \item[(3)] \((\pvar{TestimonyArguments.SolaFide}{112} \land \pvar{TestimonyArguments.SolaFide}{117}) \rightarrow \pvar{TestimonyArguments.SolaFide}{122}\)
  \item[\(\vdash\)] \(\pvar{TestimonyArguments.SolaFide}{122}\)\gloss{\emph{at issue:} The blessing of Genesis 22:16–18 was augmented because of Abraham’s obedience}
  \keepwithnext
\end{derivation}
\emph{No premise here rests on scripture alone.}
\par\medskip\noindent\textbf{\texttt{kuruvilla\allowbreak{}Assessment}}\par\nobreak\noindent
\textbf{Kuruvilla, against the step from Genesis 15.} He grants Genesis 15:6 and reads the oath of Genesis 22 as a reward for obedience. His claims reach the augmented blessing, and say nothing of the condition of what was counted at 15:6.

\paragraph{Dissent.} \emph{Kuruvilla, The Aqedah}, against \emph{Abraham, counted righteous before the covenant’s sign (Genesis 15; Romans 4)}. \textbf{Not a critique.} Not a critique of this step: it does not deny what the step concludes.
\begin{itemize}
\item Under the evidence standard: not credible.
\item Under the tradition standard: not credible.
\item Under the Scripture standard: not credible.
\end{itemize}
\paragraph*{Hahn on the oath of Genesis 22}
\leavevmode\par\nobreak
\par\medskip\noindent\textbf{\texttt{hahn\allowbreak{}Source}}\par\nobreak\noindent
Hahn’s step: the oath of Genesis 22:16–18 is sworn “because you have done this” and “because you have obeyed my voice”, so the oath to bless the nations was merited by the binding of Isaac. Rated \texttt{disputed}: Calvin accepts the text and reads the oath as “nothing more than confirmation”.

{\small\begin{verbatim}
def hahnSource : Source :=
  { primary := Reference.work hahnCovenantOathAqedah (Bib.Locus.page 93),
    tradition := Tradition.romanCatholic, confidence := Confidence.disputed }
\end{verbatim}}
\par\medskip\noindent\textbf{\texttt{hahn\allowbreak{}Line}} --- The oath of Genesis 22, merited by the binding of Isaac (Hahn)\par\nobreak\noindent
\textbf{The oath of the Aqedah, merited} (Hahn).

\begin{derivation}
  \item[(1)] \(\pvar{TestimonyArguments.SolaFide}{112}\)\gloss{Genesis 15:6: Abram believed the LORD, and he counted it to him as righteousness}
  \item[(2)] \(\pvar{TestimonyArguments.SolaFide}{117}\)\gloss{Genesis 22:16–18, 26:5: blessed because Abraham obeyed}
  \item[(3)] \((\pvar{TestimonyArguments.SolaFide}{112} \land \pvar{TestimonyArguments.SolaFide}{117}) \rightarrow \pvar{TestimonyArguments.SolaFide}{123}\)
  \item[\(\vdash\)] \(\pvar{TestimonyArguments.SolaFide}{123}\)\gloss{\emph{at issue:} The oath of Genesis 22:16–18 was merited by the binding of Isaac}
\end{derivation}
\par\medskip\noindent\textbf{\texttt{hahn\allowbreak{}On\allowbreak{}The\allowbreak{}Oath}} --- Hahn, on the oath of Genesis 22\par\nobreak\noindent
\textbf{Hahn, on the oath of Genesis 22}: the binding of Isaac merited the oath to bless the nations.

\begin{derivation}
  \item[(1)] \(\pvar{TestimonyArguments.SolaFide}{112}\)\gloss{Genesis 15:6: Abram believed the LORD, and he counted it to him as righteousness}
  \item[(2)] \(\pvar{TestimonyArguments.SolaFide}{117}\)\gloss{Genesis 22:16–18, 26:5: blessed because Abraham obeyed}
  \item[(3)] \((\pvar{TestimonyArguments.SolaFide}{112} \land \pvar{TestimonyArguments.SolaFide}{117}) \rightarrow \pvar{TestimonyArguments.SolaFide}{123}\)
  \item[\(\vdash\)] \(\pvar{TestimonyArguments.SolaFide}{123}\)\gloss{\emph{at issue:} The oath of Genesis 22:16–18 was merited by the binding of Isaac}
  \keepwithnext
\end{derivation}
\emph{No premise here rests on scripture alone.}
\par\medskip\noindent\textbf{\texttt{hahn\allowbreak{}Assessment}}\par\nobreak\noindent
\textbf{Hahn, against the step from Genesis 15.} He reads the covenant Paul means as the oath sworn after the binding of Isaac, which the binding of Isaac merited. His claims reach the oath of Genesis 22, and say nothing of the condition of what was counted at 15:6.

\paragraph{Dissent.} \emph{Hahn, Covenant, Oath, and the Aqedah}, against \emph{Abraham, counted righteous before the covenant’s sign (Genesis 15; Romans 4)}. \textbf{Not a critique.} Not a critique of this step: it does not deny what the step concludes.
\begin{itemize}
\item Under the evidence standard: not credible.
\item Under the tradition standard: not credible.
\item Under the Scripture standard: not credible.
\end{itemize}
\par\medskip\noindent\textbf{\texttt{hahn\_\allowbreak{}does\_\allowbreak{}not\_\allowbreak{}reach\_\allowbreak{}the\_\allowbreak{}counting}}\par\nobreak\noindent
\textbf{Hahn does not dissent from the step from Genesis 15}, under any standard. What he says the binding of Isaac merited is the oath of Genesis 22 and the blessing of the nations it swears, not the righteousness counted at 15:6. It does not show that Hahn would accept the step: in what was read of him he does not discuss Genesis 15:6.

{\small\begin{verbatim}
theorem hahn_does_not_reach_the_counting : ∀ (s : Standard), ¬Dissent.Credible
    s hahnAssessment.dissent
-- axioms: propext, Classical.choice, Quot.sound
\end{verbatim}}
\par\medskip\noindent\textbf{\texttt{weinfeld\allowbreak{}Assessment}}\par\nobreak\noindent
\textbf{Weinfeld, against the step from Genesis 15.} As his own claims carry him, he does not deny the step. He reaches the grant, and says nothing of what was counted at 15:6. Every claim he asserts is Scripture or the evidence of the ancient grants, so he would meet the evidence and tradition standards, if he denied the step.

\paragraph{Dissent.} \emph{Weinfeld, The Covenant of Grant}, against \emph{Abraham, counted righteous before the covenant’s sign (Genesis 15; Romans 4)}. \textbf{Not a critique.} Not a critique of this step: it does not deny what the step concludes.
\begin{itemize}
\item Under the evidence standard: not credible.
\item Under the tradition standard: not credible.
\item Under the Scripture standard: not credible; it rests on what the standard does not accept: A royal grant rewarded loyalty already shown.
\end{itemize}
\par\medskip\noindent\textbf{\texttt{maccabees\allowbreak{}Unstruck\allowbreak{}Assessment}}\par\nobreak\noindent
\textbf{1 Maccabees, against the step from Genesis 15, before anything is struck.} It denies the step, through its reading of Genesis 15:6: that what was counted was credited for Abraham’s faithfulness when tested. Every claim it asserts is a text or a reading of one, so it is an argument under the evidence and tradition standards. Under the Scripture standard it is not: 1 Maccabees is not Scripture in the canon these chapters read, and neither is its reading of Genesis 15:6.

\paragraph{Dissent.} \emph{1 Maccabees 2:52, unstruck}, against \emph{Abraham, counted righteous before the covenant’s sign (Genesis 15; Romans 4)}. \textbf{A critique.} A critique: it can be held with the step’s grounds, and it argues its denial by a step of its own.
\begin{itemize}
\item Under the evidence standard: credible.
\item Under the tradition standard: credible.
\item Under the Scripture standard: not credible; it rests on what the standard does not accept: 1 Maccabees 2:52, Sirach 44:20: Abraham found faithful in trial, and it was credited to him as righteousness; The righteousness counted at Genesis 15:6 was credited for Abraham’s faithfulness when tested.
\end{itemize}
\par\medskip\noindent\textbf{\texttt{maccabees\allowbreak{}Assessment}}\par\nobreak\noindent
\textbf{1 Maccabees, against the step from Genesis 15}, with what Genesis reads struck: the reading that the counting of 15:6 was credited for the testing, which Genesis places “after these things” (22:1).

\paragraph{Dissent.} \emph{1 Maccabees 2:52}, against \emph{Abraham, counted righteous before the covenant’s sign (Genesis 15; Romans 4)}. \textbf{Not a critique.} Not a critique of this step: with the claims the texts refute struck, what remains does not deny what the step concludes.
\begin{itemize}
\item Struck, because the text reads otherwise: The righteousness counted at Genesis 15:6 was credited for Abraham’s faithfulness when tested. Genesis 15:6 counts Abram’s believing as righteousness; the testing comes “after these things” (Genesis 22:1).
\item Under the evidence standard: not credible.
\item Under the tradition standard: not credible.
\item Under the Scripture standard: not credible; it rests on what the standard does not accept: 1 Maccabees 2:52, Sirach 44:20: Abraham found faithful in trial, and it was credited to him as righteousness.
\end{itemize}
\par\medskip\noindent\textbf{\texttt{weinfeld\_\allowbreak{}leaves\_\allowbreak{}the\_\allowbreak{}counting\_\allowbreak{}open}}\par\nobreak\noindent
\textbf{Weinfeld’s grant leaves the counting open.} His premises, that Genesis 15 is a one-sided grant, that a grant rewarded loyalty already shown, and that the blessing came “because” Abraham obeyed, are kept alike by a reading on which Abraham was counted righteous without works and by one on which he was not. What they establish is the grant; what was counted at 15:6 is another claim.

It does not show that Weinfeld is wrong about the grant. It may well have rewarded Abraham’s obedience. The result says that, even so, nothing follows about the condition of what was counted to him.

{\small\begin{verbatim}
theorem weinfeld_leaves_the_counting_open : Independent
    weinfeldOnAbraham.premises (p Claim.abrahamCountedRighteousWithoutWorks)
-- axioms: propext, Quot.sound
\end{verbatim}}
\par\medskip\noindent\textbf{\texttt{weinfeld\_\allowbreak{}reaches\_\allowbreak{}the\_\allowbreak{}grant}}\par\nobreak\noindent
\textbf{Weinfeld’s grant reaches its own conclusion.} On his premises, the grant to Abraham rewarded his obedience.

{\small\begin{verbatim}
theorem weinfeld_reaches_the_grant : Establishes weinfeldOnAbraham
-- axioms: propext, Quot.sound
\end{verbatim}}
\par\medskip\noindent\textbf{\texttt{weinfeld\allowbreak{}On\allowbreak{}Abraham\_\allowbreak{}is\_\allowbreak{}satisfiable}}\par\nobreak\noindent
Weinfeld’s premises can be held without contradiction.

{\small\begin{verbatim}
theorem weinfeldOnAbraham_is_satisfiable : Satisfiable
    weinfeldOnAbraham.premises
-- axioms: propext, Quot.sound
\end{verbatim}}
\par\medskip\noindent\textbf{\texttt{weinfeld\_\allowbreak{}does\_\allowbreak{}not\_\allowbreak{}reach\_\allowbreak{}the\_\allowbreak{}counting}}\par\nobreak\noindent
\textbf{Weinfeld does not dissent from the step from Genesis 15}, under any standard. His claims do not deny it: they reach the grant, not the counting.

It does not show that Weinfeld would agree with the step. A reader who adds that what was counted at 15:6 was credited for Abraham’s obedience would deny it; that is what 1 Maccabees does.

{\small\begin{verbatim}
theorem weinfeld_does_not_reach_the_counting : ∀ (s : Standard),
    ¬Dissent.Credible s weinfeldAssessment.dissent
-- axioms: propext, Classical.choice, Quot.sound
\end{verbatim}}
\par\medskip\noindent\textbf{\texttt{maccabees\_\allowbreak{}unstruck\_\allowbreak{}counts}}\par\nobreak\noindent
\textbf{Unstruck, 1 Maccabees is an argument against the step}, under the evidence and tradition standards: it reads Genesis 15:6 as credited for the testing, and from that denies that Abraham was counted righteous without works.

{\small\begin{verbatim}
theorem maccabees_unstruck_counts : Dissent.Credible Standard.evidence
    maccabeesUnstruckAssessment.dissent ∧ Dissent.Credible Standard.tradition
    maccabeesUnstruckAssessment.dissent
-- axioms: propext, Classical.choice, Quot.sound
\end{verbatim}}
\par\medskip\noindent\textbf{\texttt{genesis\_\allowbreak{}strikes\_\allowbreak{}the\_\allowbreak{}maccabees\_\allowbreak{}reading}}\par\nobreak\noindent
\textbf{Genesis strikes the reading 1 Maccabees needs, and then it no longer dissents.} Genesis counts Abram’s believing at 15:6 and places the testing “after these things” (22:1). With the reading that the counting was credited for the testing struck, nothing in 1 Maccabees’ position denies the step, under any standard.

It does not show that 1 Maccabees says nothing true of Abraham. He was found faithful when tested. What the text refutes is only that the counting of 15:6 was credited for it.

{\small\begin{verbatim}
theorem genesis_strikes_the_maccabees_reading :
    maccabeesAssessment.findings.struck = [Claim.countingWasForTheTesting] ∧ ∀
    (s : Standard), ¬Dissent.Credible s maccabeesAssessment.dissent
-- axioms: propext, Classical.choice, Quot.sound
\end{verbatim}}
\par\medskip\noindent\textbf{\texttt{kuruvilla\_\allowbreak{}does\_\allowbreak{}not\_\allowbreak{}reach\_\allowbreak{}the\_\allowbreak{}counting}}\par\nobreak\noindent
\textbf{Kuruvilla does not dissent from the step from Genesis 15}, under any standard. His claims reach the blessing of Genesis 22, augmented for obedience, and say nothing of the condition of what was counted at 15:6.

It does not show that Kuruvilla is wrong about Genesis 22. Calvin reads the oath as “nothing more than confirmation”, and Kuruvilla answers him; that disagreement is about the blessing of Genesis 22, not the counting of 15:6.

{\small\begin{verbatim}
theorem kuruvilla_does_not_reach_the_counting : ∀ (s : Standard),
    ¬Dissent.Credible s kuruvillaAssessment.dissent
-- axioms: propext, Classical.choice, Quot.sound
\end{verbatim}}
\par\medskip\noindent\textbf{\texttt{genesis\allowbreak{}Fifteen\allowbreak{}Step}}\par\nobreak\noindent
\textbf{The step from Genesis 15, with the dissents found against it}: Weinfeld’s grant, Kuruvilla on the Aqedah, Hahn on its oath, and 1 Maccabees, the strongest known, which alone denies the step.

\paragraph{The rating, computed.} \emph{Abraham, counted righteous before the covenant’s sign (Genesis 15; Romans 4)} is cited at \emph{well supported} (Rom 4:3-12; \cite[ad loc. Gen 22:15]{calvin-commentary-genesis-1847}; \cite[ad loc. Jas 2:23]{calvin-commentaries-catholic-epistles}; \cite[260--263]{niehaus-gods-covenant-abraham-2013}; \cite{robertson-christ-covenants-1980}; \cite{kline-kingdom-prologue-2006}). Dissents encoded: \emph{Weinfeld, The Covenant of Grant}, \emph{1 Maccabees 2:52}, \emph{Kuruvilla, The Aqedah}, \emph{Hahn, Covenant, Oath, and the Aqedah}. The strongest known: \emph{1 Maccabees 2:52}.
\begin{itemize}
\item Under the evidence standard: \textbf{well supported} --- every encoded dissent fails a check, the strongest known among them, so the step’s support stands.
\item Under the tradition standard: \textbf{well supported} --- every encoded dissent fails a check, the strongest known among them, so the step’s support stands.
\item Under the Scripture standard: \textbf{well supported} --- every encoded dissent fails a check, the strongest known among them, so the step’s support stands.
\end{itemize}
\par\medskip\noindent\textbf{\texttt{abraham\_\allowbreak{}step\_\allowbreak{}ratings}}\par\nobreak\noindent
\textbf{The step from Genesis 15 stands at \texttt{well\allowbreak{}Supported} under every standard.} No dissent found survives its checks. Weinfeld’s, Kuruvilla’s and Hahn’s do not deny the step: they reach the grant, and the blessing and the oath of Genesis 22, not the counting of 15:6. 1 Maccabees denies it only through its reading of Genesis 15:6, which Genesis itself refutes by its order.

It does not show that no reader could dissent. A verified reader who argued that the counting of 15:6 was conditioned on Abraham’s obedience without tying it to the later testing, for example on his leaving Haran (Genesis 12:4), would be a new dissent, and the rating would be worked out again. Hahn’s \emph{Kinship by Covenant} and Williamson’s books have not been read.

{\small\begin{verbatim}
theorem abraham_step_ratings : genesisFifteenStep.ratings =
    [(Standard.evidence, Computed.stands Confidence.wellSupported),
    (Standard.tradition, Computed.stands Confidence.wellSupported),
    (Standard.scriptural, Computed.stands Confidence.wellSupported)]
-- axioms: propext
\end{verbatim}}
\subsection*{Trent’s two Latin texts}
\addcontentsline{toc}{subsection}{Trent’s two Latin texts}
Two of the three texts Trent cites for the increase of justification read differently in the Latin Trent quotes and in the Greek (the chapter “James 2:24”). “He that is just, let him be justified still” is Revelation 22:11 in the Latin Vulgate; the Greek, as the critical editions print it, says “let the righteous still do righteousness”. “Be not afraid to be justified even to death” is Sirach 18:22 in the Vulgate; the Greek has “to be justified” too, but in a sentence about paying vows: “do not wait until death to be justified”. Trent introduces both with “as it is written” (Session VI, ch. 10), having declared the Vulgate “authentic” for “public lectures, disputations, sermons and expositions” (Session IV).

In 1943 Pius XII urged Catholic scholars back to the Hebrew and the Greek. The original text, he wrote in \emph{Divino Afflante Spiritu},

\begin{quote}
“written by the inspired author himself, has more authority and greater weight than any even the very best translation” (§16).

\end{quote}
Four centuries earlier, Trent had taught the increase of justification from two verses whose Latin reads differently from the Greek. Do the pope and the council contradict each other? The answer here is drawn from Rome’s own documents and nothing else: the council’s decree, the pope’s letter, and the Catholic Church’s own translation of the Greek.

\subsubsection*{The answer depends on what “as it is written” claims}
\textbf{Taken as a claim about what the inspired authors wrote}, it cannot be held together with Pius XII. He teaches that the original text, “written by the inspired author himself, has more authority and greater weight than any even the very best translation” (§16). The original of Revelation 22:11, as the Church’s own translation gives it, says “the righteous must still do right”, not “be justified still”. Those commitments cannot all be true together (\texttt{trent\_\allowbreak{}as\_\allowbreak{}the\_\allowbreak{}inspired\_\allowbreak{}text\_\allowbreak{}contradicts\_\allowbreak{}pius}, which is checked for Revelation 22:11). Every one of them is Rome’s.

\textbf{Taken as Pius XII takes Trent’s decree on the Vulgate}, there is no contradiction. Pius XII denies that going back to the originals “in any way derogates” from Trent (§20). The Vulgate’s authority is “juridical”, a matter of the Church’s approval, and not a claim about the original wording. And the Vulgate, “in the sense in which the Church has understood and understands it”, is “free from any error whatsoever in matters of faith and morals” and “may be quoted safely ... in disputations” (§21). Taken that way, Trent quotes the Church’s approved Latin, not the inspired authors’ words, and everything Rome says can be held together (\texttt{pius\_\allowbreak{}reading\_\allowbreak{}is\_\allowbreak{}consistent}).

\textbf{What the consistent way of taking it costs.} It moves the support for the two texts from the inspired authors to the Church: the claims that carry it rest on Rome’s word alone (\texttt{pius\_\allowbreak{}reading\_\allowbreak{}rests\_\allowbreak{}on\_\allowbreak{}authority}). And Pius XII asks that a doctrine taught from the Vulgate be confirmed from the original texts (§22), which at these two verses do not confirm it (\texttt{pius\_\allowbreak{}finds\_\allowbreak{}no\_\allowbreak{}confirmation}). So of the three texts chapter 10 cites, only James 2:24 is left to confirm the doctrine from the originals. It is weighed in the chapter “James 2:24”.

\subsubsection*{What this does not show}
It does not show that Rome contradicts itself. Trent does not say which way its “as it is written” is meant, and Pius XII supplies the way on which there is no contradiction. The first contradicts Pius XII. The second makes the two texts an appeal to the Church’s authority, which Pius XII himself asks to be confirmed from the originals, and the originals do not confirm it.

Nor does it show that the doctrine is false. Pius XII holds the Vulgate free from error in faith because Rome holds the doctrine true on other grounds. This section is about what these two texts can bear.

\subsubsection*{How this is checked}
\paragraph*{Taken as the inspired text}
\leavevmode\par\nobreak
\par\medskip\noindent\textbf{\texttt{rome\allowbreak{}On\allowbreak{}The\allowbreak{}Inspired\allowbreak{}Text}}\par\nobreak\noindent
\textbf{Rome’s own commitments, with Trent’s citation taken as the inspired text.} Trent’s “as it is written” at Revelation 22:11, as a claim about what its author wrote. The Greek, as the Catholic Church’s own translation gives it. Pius XII’s principle that the original outweighs any translation. And the step from the last two to “the inspired text does not say ‘be justified still’”.

\begin{derivation}
  \item[(1)] \(\pvar{TestimonyArguments.SolaFide}{92}\)\gloss{Revelation 22:11 says: he that is just, let him be justified still}
  \item[(2)] \(\pvar{TestimonyArguments.SolaFide}{93}\)\gloss{The Greek of Revelation 22:11 says: let the righteous still do righteousness}
  \item[(3)] \(\pvar{TestimonyArguments.SolaFide}{96}\)\gloss{The original text, written by the inspired author, outweighs any translation: a reading it lacks is not his word}
  \item[(4)] \((\pvar{TestimonyArguments.SolaFide}{93} \land \pvar{TestimonyArguments.SolaFide}{96}) \rightarrow \lnot \pvar{TestimonyArguments.SolaFide}{92}\)\gloss{\emph{would deny:} Revelation 22:11 says: he that is just, let him be justified still}
\end{derivation}
\par\medskip\noindent\textbf{\texttt{trent\_\allowbreak{}as\_\allowbreak{}the\_\allowbreak{}inspired\_\allowbreak{}text\_\allowbreak{}contradicts\_\allowbreak{}pius}}\par\nobreak\noindent
\textbf{Trent’s Revelation 22:11, taken as what its inspired author wrote, cannot be held together with Pius XII.} The original “has more authority and greater weight than any even the very best translation” (\emph{Divino Afflante Spiritu} §16). The original of Revelation 22:11, as the Catholic Church’s own translation gives it, says “the righteous must still do right”. With those two, the claim that the inspired author wrote “let him be justified still” cannot be true. Every premise is Rome’s own.

It does not show that Trent meant its citation that way. Taken as Pius XII takes Trent’s use of the Vulgate, there is no contradiction (\texttt{pius\_\allowbreak{}reading\_\allowbreak{}is\_\allowbreak{}consistent}).

{\small\begin{verbatim}
theorem trent_as_the_inspired_text_contradicts_pius : ¬Satisfiable
    romeOnTheInspiredText
-- axioms: propext, Classical.choice, Quot.sound
\end{verbatim}}
\paragraph*{Taken as Pius XII takes the Vulgate}
\leavevmode\par\nobreak
\par\medskip\noindent\textbf{\texttt{pius\allowbreak{}Confirmation\allowbreak{}Source}}\par\nobreak\noindent
Pius XII’s step. A doctrine taught from the Vulgate is to be confirmed from the original texts (§22). The original of Revelation 22:11 has no “be justified”. The Greek of Sirach 18:22 has it, but says “do not wait until death to be justified” of paying a vow, not “be not afraid to be justified even to death”. So these two texts do not confirm the doctrine from the originals. Rated \texttt{well\allowbreak{}Supported}: it is Pius XII’s own requirement, applied to texts whose originals the Church’s own translation gives.

{\small\begin{verbatim}
def piusConfirmationSource : Source :=
  {
    primary :=
      Reference.work divinoAfflanteSpiritu (Bib.Locus.sectionRef "§22"),
    supporting := [Reference.work nabre], tradition := Tradition.romanCatholic,
    confidence := Confidence.wellSupported }
\end{verbatim}}
\par\medskip\noindent\textbf{\texttt{pius\allowbreak{}On\allowbreak{}Trents\allowbreak{}Vulgate}} --- Pius XII, on Trent’s Vulgate\par\nobreak\noindent
\textbf{Rome’s own position on the two Latin texts, taken as Pius XII takes Trent.} Trent quotes the Vulgate as the Church’s approved text. The Vulgate is free from error in faith and morals. The original outweighs any translation, and the originals of the two verses do not say what the Latin says. A doctrine taught from the Vulgate is to be confirmed from the originals. The conclusion: the two texts do not confirm the increase of justification from the originals.

\begin{derivation}
  \item[(1)] \(\pvar{TestimonyArguments.SolaFide}{98}\)\gloss{Trent quotes the Vulgate’s Revelation 22:11 and Sirach 18:22 as juridically authentic, not as the inspired authors’ words}
  \item[(2)] \(\pvar{TestimonyArguments.SolaFide}{97}\)\gloss{The Vulgate, as the Church understands it, is free from error in faith and morals: its authenticity is juridical}
  \item[(3)] \(\pvar{TestimonyArguments.SolaFide}{96}\)\gloss{The original text, written by the inspired author, outweighs any translation: a reading it lacks is not his word}
  \breakrow{(4)} \(\pvar{TestimonyArguments.SolaFide}{93}\)\gloss{The Greek of Revelation 22:11 says: let the righteous still do righteousness}
  \breakrow{(5)} \(\pvar{TestimonyArguments.SolaFide}{95}\)\gloss{The Greek of Sirach 18:22 is about paying a vow: do not wait until death to be justified}
  \breakrow{(6)} \(\pvar{TestimonyArguments.SolaFide}{99}\)\gloss{A doctrine taught from the Vulgate is to be confirmed from the original texts}
  \breakrow{(7)} \((\pvar{TestimonyArguments.SolaFide}{93} \land \pvar{TestimonyArguments.SolaFide}{96}) \rightarrow \lnot \pvar{TestimonyArguments.SolaFide}{92}\)\gloss{\emph{would deny:} Revelation 22:11 says: he that is just, let him be justified still}
  \breakrow{(8)} \((\pvar{TestimonyArguments.SolaFide}{95} \land \pvar{TestimonyArguments.SolaFide}{96}) \rightarrow \lnot \pvar{TestimonyArguments.SolaFide}{94}\)\gloss{\emph{would deny:} Sirach 18:22 says: be not afraid to be justified even to death}
  \item[(9)] \((\pvar{TestimonyArguments.SolaFide}{99} \land \lnot \pvar{TestimonyArguments.SolaFide}{92} \land \lnot \pvar{TestimonyArguments.SolaFide}{94}) \rightarrow \lnot \pvar{TestimonyArguments.SolaFide}{100}\)
  \item[\(\vdash\)] \(\lnot \pvar{TestimonyArguments.SolaFide}{100}\)\gloss{\emph{at issue, the denial of:} Revelation 22:11 and Sirach 18:22 confirm the increase of justification from the original texts}
  \keepwithnext
\end{derivation}
\emph{No premise here rests on scripture alone.}
\par\medskip\noindent\textbf{\texttt{pius\allowbreak{}Reading}}\par\nobreak\noindent
\textbf{Rome’s position, as Pius XII takes Trent.} The Vulgate is the Church’s approved text and free from error in faith. Trent quotes it as such. The originals of the two verses do not say what the Latin says, and outweigh any translation. So the two texts do not confirm the doctrine from the originals. Everything else is as Trent holds it.

{\small\begin{verbatim}
def piusReading : Valuation Claim :=
  fun a =>
    match a with
    | Claim.rev22_11BeJustifiedStill => False
    | Claim.sir18_22BeJustifiedToDeath => False
    | Claim.latinTextsConfirmedFromOriginals => False
    | Claim.worksAreFruitNotGround => False
    | x => True
\end{verbatim}}
\par\medskip\noindent\textbf{\texttt{pius\_\allowbreak{}reading\_\allowbreak{}is\_\allowbreak{}consistent}}\par\nobreak\noindent
\textbf{Taken as Pius XII takes Trent’s decree on the Vulgate, Rome is consistent.} All of this can be held together: Trent quoting the Vulgate as the Church’s approved text; the Vulgate free from error in faith and morals; the original outweighing any translation; the Greek of the two verses as the Church’s own translation gives it; and the requirement that a doctrine taught from the Vulgate be confirmed from the originals.

{\small\begin{verbatim}
theorem pius_reading_is_consistent : Satisfiable piusOnTrentsVulgate.premises
-- axioms: propext, Quot.sound
\end{verbatim}}
\par\medskip\noindent\textbf{\texttt{pius\_\allowbreak{}finds\_\allowbreak{}no\_\allowbreak{}confirmation}}\par\nobreak\noindent
\textbf{By Pius XII’s own teaching, the two Latin texts do not confirm the increase of justification from the originals.} He asks that a doctrine taught from the Vulgate be confirmed from the original texts (§22). The original of Revelation 22:11 has no “be justified”. The Greek of Sirach 18:22 has the verb, in a sentence about paying a vow before death, and not the Latin’s “be not afraid ... even to death”. So of the three texts chapter 10 cites, only James 2:24 is left to confirm the doctrine from the originals.

It does not show that the Greek of Sirach 18:22 could not be read in Trent’s favour. That would be a further argument about what its “be justified” means, and none is written down here.

{\small\begin{verbatim}
theorem pius_finds_no_confirmation : Establishes piusOnTrentsVulgate
-- axioms: propext, Quot.sound
\end{verbatim}}
\par\medskip\noindent\textbf{\texttt{pius\_\allowbreak{}reading\_\allowbreak{}rests\_\allowbreak{}on\_\allowbreak{}authority}}\par\nobreak\noindent
\textbf{Taken that way, the two texts rest on the Church’s word.} The claims that carry Trent’s citation are cited to Pius XII alone: that it quotes the Vulgate as the Church’s approved text, and that the Vulgate is free from error in faith. So are the principle that the original outweighs any translation and the requirement of confirmation from the originals, though Westminster holds the principle too.

{\small\begin{verbatim}
theorem pius_reading_rests_on_authority :
    piusOnTrentsVulgate.restingOnAuthority =
    [Claim.trentQuotesTheVulgateJuridically,
    Claim.vulgateFreeFromDoctrinalError, Claim.originalTextDecides,
    Claim.doctrineToBeConfirmedFromOriginals]
-- axioms: propext
\end{verbatim}}
\subsection*{What Jesus’ own words say about faith and works}
\addcontentsline{toc}{subsection}{What Jesus’ own words say about faith and works}
The four witnesses include two sayings of Jesus, Luke 7:50 and John 6:29. But a reader who has followed Paul and Peter this far will want to hear more from Jesus himself. Three more passages are the ones each side cites.

\subsubsection*{Luke 18:9–14: the Pharisee and the tax collector}
The Pharisee lists his fasting and his tithes. The tax collector stands far off, beats his breast, and prays, “God, be merciful to me, a sinner!” Then Jesus gives the verdict:

\begin{quote}
“I tell you, this man went down to his house justified, rather than the other” (Luke 18:14).

\end{quote}
It is Jesus’ own teaching on justification and works, and it uses Paul’s word, “justified” (δεδικαιωμένος).

Two things here are accepted on both sides. First, the word gives a verdict. No reader cited here takes “justified” at 18:14 as the inner renewal itself, and Calvin reads it as “to stand before God as if we were righteous”. That is a verdict; it is not said to be \emph{only} a verdict. Trent holds both that the justified are counted just and that they are made just (ch. 7), and that is where the parable and Trent part. Second, the tax collector brings no works that earn anything. Trent itself says that nothing that comes before justification, “whether faith or works”, merits it (Session VI, ch. 8). So Trent’s case is consistent with every word of the parable (\texttt{trent\_\allowbreak{}grants\_\allowbreak{}the\_\allowbreak{}parable}), though the Council never refers to it.

What divides them is the step from the plea to “faith suffices”. Trent never cites Luke 18. But it requires the dispositions it lists, namely fear, hope, beginning to love God, and penitence (ch. 6), and it condemns whoever says that nothing else is required to cooperate (canon 9). Read by those chapters, the tax collector’s humility and contrition are dispositions his justification required, and not a plea that was enough alone. So the step is rated \texttt{disputed}, and Luke 18 meets Trent where Luke 7:50 does: at Trent’s definition of justification (\texttt{why\allowbreak{}Luke\allowbreak{}Eighteen\allowbreak{}Stands\allowbreak{}Against\allowbreak{}Trent}).

The verdict sense at 18:14 is what the argument from Paul’s word claims of “justify” in Paul, found in Jesus’ own parable. It is not added to that argument. What a word means in Luke is evidence for what it means in Paul, not a premise of it, and that argument stands on its own sources.

\subsubsection*{Matthew 7:21–23: “Lord, Lord”}
\begin{quote}
“Not everyone who says to me, ‘Lord, Lord,’ will enter the kingdom of heaven, but the one who does the will of my Father who is in heaven.”

\end{quote}
And to those who prophesied in his name: “I never knew you.” The passage does not mention faith, and gives doing the Father’s will as the test. Rome reads it as condemning profession without obedience, with obedience a condition of entry in its own right. Calvin reads the doing as including faith, and “I never knew you” as meaning they were never united to Christ. Both are written down here as readings, each cited and each rated \texttt{disputed}. Neither is given as what the passage plainly says (\texttt{matthew7\_\allowbreak{}leaves\_\allowbreak{}the\_\allowbreak{}ground\_\allowbreak{}open}).

\subsubsection*{Matthew 19:16–22: the rich young man}
\begin{quote}
“If you would enter life, keep the commandments” (Matthew 19:17).

\end{quote}
Trent quotes these words (Session VI, ch. 7). Rome reads them as the way to life, possible by grace (\emph{Veritatis Splendor} 12). Calvin reads them as Christ answering on the law’s own terms, so that the young man may see he has not kept it and turn to faith. Here too each reading is cited, and each is rated \texttt{disputed}.

\subsubsection*{What the debate decides}
Heard together, with Luke 7:50, Luke 18, Trent, and each reading of Matthew 7 and 19, nothing is forced (\texttt{jesus\_\allowbreak{}words\_\allowbreak{}force\_\allowbreak{}nothing}). That is, nothing has to be accepted however the ties are settled. Every position is in a tie the debate cannot settle: each saying in Luke with Trent, and each reading of Matthew with its rival.

Heard without Trent, Luke 7:50 and Luke 18 are forced (\texttt{without\_\allowbreak{}trent\_\allowbreak{}luke\_\allowbreak{}is\_\allowbreak{}forced}). What is forced is that faith sufficed, not that faith alone saves. Matthew is still not forced. The rival readings of Matthew 7 and 19 conflict with each other and not with Luke.

Trent here is the same position as in the main debate, concluding that salvation is not apart from works. It meets Luke through its definition of justification, not through that conclusion. Rome’s readings of Matthew reach the same conclusion without the definition, and so do not conflict with Luke (\texttt{luke\_\allowbreak{}stands\_\allowbreak{}with\_\allowbreak{}rome\_\allowbreak{}on\_\allowbreak{}matthew}).

So Jesus’ words do not settle the Reformation debate alone. Without Trent, Luke’s two sayings stand. Matthew’s are argued both ways, and the ratings cannot choose.

\subsubsection*{What it depends on: two accounts of justification}
Every tie between Luke and Trent here sets Luke’s step, from the saying to “faith sufficed”, rated \texttt{disputed}, against Trent’s definition of justification together with the step from it to canon 9. Nothing else in Trent’s case meets Luke (\texttt{why\allowbreak{}Luke\allowbreak{}Eighteen\allowbreak{}Stands\allowbreak{}Against\allowbreak{}Trent}, \texttt{why\allowbreak{}The\allowbreak{}Dominical\allowbreak{}Case\allowbreak{}Stands\allowbreak{}Against\allowbreak{}Trent}). Trent accepts every word of the parable (\texttt{trent\_\allowbreak{}grants\_\allowbreak{}the\_\allowbreak{}parable}). So the words do not divide them. What justification is does.

For the Reformers justification is a verdict: God declares the ungodly righteous, as a judge acquits. For Trent it is “not remission of sins merely, but also the sanctification and renewal of the inward man” (Session VI, ch. 7): the justified are not only counted just but made just.

What can be shown against that definition from the meaning of words in their setting is checked in \texttt{what\allowbreak{}Trents\allowbreak{}Definition\allowbreak{}Claims}, and it depends on how the definition is taken.

\begin{itemize}
  \item \textbf{As a claim about the word}, that Paul’s verb “justify” (δικαιόω) \emph{means} renewal, it fails three times, each at a step rated \texttt{well\allowbreak{}Supported}. It fails against the argument from usage: the verb is a courtroom word; Paul has other words for renewal and sets them beside it (Titus 3:5–7; 1 Corinthians 6:11); and a word contributes the least meaning its context requires. It fails at the Latin explanation it rests on, since “justify” (\emph{iustificare}) as “make righteous” does not settle what the Greek word means. And it fails against Romans 4 read for the word. In the main debate, Trent taken this way cannot be defended (\texttt{trent\_\allowbreak{}on\_\allowbreak{}pauls\_\allowbreak{}word\_\allowbreak{}indefensible}). Luke 18 agrees with this without adding a fourth failure: no reader cited here takes “justified” at 18:14 as the renewal itself.
  \item \textbf{As a claim about what God does}, that when he justifies he also renews, nothing about the word reaches it. Taken this way, Trent can be held with every premise of the argument from usage (\texttt{lexical\_\allowbreak{}case\_\allowbreak{}does\_\allowbreak{}not\_\allowbreak{}reach\_\allowbreak{}what\_\allowbreak{}god\_\allowbreak{}does}). It meets Paul (Galatians, Romans 4) and Luke only at steps rated \texttt{disputed}. A reader who grants those steps rejects it, and one who denies them, as Augustine and the \emph{Joint Declaration} do, need not.
\end{itemize}
Which of the two Trent means is itself rated \texttt{disputed}, and the check does not decide it. In this debate Trent is heard as it states its definition, tied to neither, so the check changes no verdict above.

It does not show that logic refutes Rome. Trent’s case can be held without contradiction taken either way. What the check shows is narrower, and it is as far as an argument from the meaning of words can go. The evidence about the word, rated \texttt{well\allowbreak{}Supported}, tells against taking the definition as what the word means. Whether God’s justifying act also includes renewal is something the word’s meaning leaves open, and there the matter stands as the main debate leaves it: a tie at the \texttt{disputed} step that \texttt{where\allowbreak{}Trent\allowbreak{}Parts\allowbreak{}From\allowbreak{}Paul} names (\texttt{gospel\_\allowbreak{}not\_\allowbreak{}forced}).

This debate is kept apart from the main one. Adding its positions there would change what that section says about its own seventeen, and that change should be made and reported by itself.

\subsubsection*{How strong each position is}
A position’s weakest link is the lowest-rated premise or step it rests on. An attack fails only when the position attacked is rated above its attacker at that link. The words of Luke 18 and its verdict are rated above \texttt{disputed}, but its step is not, and each reading of Matthew rests on a reading rated \texttt{disputed}. So here every weakest link is \texttt{disputed}, and every attack stands.

\subsubsection*{What the debate decides}
A few words recur. \textbf{Forced}: accepted however every tie is settled. \textbf{Can be defended}: included in some consistent set of positions that answers every attack on its members. \textbf{Not forced}: some such set leaves it out, which is not the same as false.

\subsubsection*{How this is checked}
\paragraph*{Luke 18:9–14}
\leavevmode\par\nobreak
\par\medskip\noindent\textbf{\texttt{luke18\allowbreak{}To\allowbreak{}Sufficiency}}\par\nobreak\noindent
\textbf{From the parable to “faith suffices”.} The tax collector brings no works, only a plea for mercy, and goes home justified, which is God’s verdict. So what justified him was not his works but his trust in God’s mercy: faith sufficed. Rated \texttt{disputed}: Trent accepts the text and the absence of works, and denies that the plea was enough alone (Session VI, ch. 6 and canon 9).

\begin{derivation}
  \item[] \((\pvar{TestimonyArguments.SolaFide}{59} \land \pvar{TestimonyArguments.SolaFide}{60} \land \pvar{TestimonyArguments.SolaFide}{61}) \rightarrow \pvar{TestimonyArguments.SolaFide}{26}\)\gloss{\emph{would give:} Faith is sufficient: whoever believes is saved}
\end{derivation}
\par\medskip\noindent\textbf{\texttt{luke18\allowbreak{}Line}} --- Luke 18:9–14 (the tax collector justified)\par\nobreak\noindent
\textbf{Luke 18, as a line of argument}: from the parable, the verdict, and the plea without works, to “faith suffices”.

\begin{derivation}
  \item[(1)] \(\pvar{TestimonyArguments.SolaFide}{59}\)\gloss{Luke 18:9–14 — the tax collector pleads for mercy and goes home justified, rather than the Pharisee who listed his works}
  \item[(2)] \(\pvar{TestimonyArguments.SolaFide}{60}\)\gloss{‘Justified’ at Luke 18:14 is God’s verdict: accepted as righteous before God}
  \item[(3)] \(\pvar{TestimonyArguments.SolaFide}{61}\)\gloss{The tax collector brings no works that merit, only a humble and contrite plea for mercy}
  \item[(4)] \((\pvar{TestimonyArguments.SolaFide}{59} \land \pvar{TestimonyArguments.SolaFide}{60} \land \pvar{TestimonyArguments.SolaFide}{61}) \rightarrow \pvar{TestimonyArguments.SolaFide}{26}\)
  \item[\(\vdash\)] \(\pvar{TestimonyArguments.SolaFide}{26}\)\gloss{\emph{at issue:} Faith is sufficient: whoever believes is saved}
\end{derivation}
\par\medskip\noindent\textbf{\texttt{luke18\allowbreak{}Case}} --- Luke 18:14: “this man went down to his house justified”\par\nobreak\noindent
\textbf{Luke 18’s case}: what the parable says, and no more. It claims that the tax collector’s faith sufficed, not that faith \emph{alone} saves. Luke 7:50 makes the same claim from a second saying.

\begin{derivation}
  \item[(1)] \(\pvar{TestimonyArguments.SolaFide}{59}\)\gloss{Luke 18:9–14 — the tax collector pleads for mercy and goes home justified, rather than the Pharisee who listed his works}
  \item[(2)] \(\pvar{TestimonyArguments.SolaFide}{60}\)\gloss{‘Justified’ at Luke 18:14 is God’s verdict: accepted as righteous before God}
  \item[(3)] \(\pvar{TestimonyArguments.SolaFide}{61}\)\gloss{The tax collector brings no works that merit, only a humble and contrite plea for mercy}
  \item[(4)] \((\pvar{TestimonyArguments.SolaFide}{59} \land \pvar{TestimonyArguments.SolaFide}{60} \land \pvar{TestimonyArguments.SolaFide}{61}) \rightarrow \pvar{TestimonyArguments.SolaFide}{26}\)
  \item[\(\vdash\)] \(\pvar{TestimonyArguments.SolaFide}{26}\)\gloss{\emph{at issue:} Faith is sufficient: whoever believes is saved}
  \keepwithnext
\end{derivation}
\emph{No premise here rests on scripture alone.}
\paragraph*{Matthew 7:21–23}
\leavevmode\par\nobreak
\par\medskip\noindent\textbf{\texttt{matthew7\allowbreak{}To\allowbreak{}Works}}\par\nobreak\noindent
The Catholic step from Matthew 7:21. If doing the Father’s will is a condition of entering the kingdom in its own right, then salvation is not apart from works. Rated \texttt{disputed}: Calvin accepts the text and denies that the doing is more than the fruit of faith.

\begin{derivation}
  \item[] \((\pvar{TestimonyArguments.SolaFide}{62} \land \pvar{TestimonyArguments.SolaFide}{63}) \rightarrow \lnot \pvar{TestimonyArguments.SolaFide}{28}\)\gloss{\emph{would deny:} Salvation is not by works: no work is its ground}
\end{derivation}
\par\medskip\noindent\textbf{\texttt{matthew7\allowbreak{}To\allowbreak{}Fruit}}\par\nobreak\noindent
The Reformed step from Matthew 7:21. If doing the Father’s will includes believing, and those rejected were never known, then the doing is the fruit of faith and not a condition of entry in its own right. Rated \texttt{disputed}: the \emph{Catechism} accepts the text and reads it the other way.

\begin{derivation}
  \item[] \((\pvar{TestimonyArguments.SolaFide}{62} \land \pvar{TestimonyArguments.SolaFide}{64}) \rightarrow \lnot \pvar{TestimonyArguments.SolaFide}{63}\)\gloss{\emph{would deny:} In Matthew 7:21–23, doing the Father’s will is a condition of entering the kingdom in its own right, not only the fruit of faith}
\end{derivation}
\par\medskip\noindent\textbf{\texttt{matthew7\allowbreak{}Obedience\allowbreak{}Line}} --- Matthew 7:21–23, read as obedience a condition of entry\par\nobreak\noindent
\textbf{Matthew 7:21–23, as Rome reads it}: doing the Father’s will is a condition of entry in its own right, so salvation is not apart from works.

\begin{derivation}
  \item[(1)] \(\pvar{TestimonyArguments.SolaFide}{62}\)\gloss{Matthew 7:21–23 — the criterion at the judgment is doing the Father’s will, not profession or mighty works}
  \item[(2)] \(\pvar{TestimonyArguments.SolaFide}{63}\)\gloss{In Matthew 7:21–23, doing the Father’s will is a condition of entering the kingdom in its own right, not only the fruit of faith}
  \item[(3)] \((\pvar{TestimonyArguments.SolaFide}{62} \land \pvar{TestimonyArguments.SolaFide}{63}) \rightarrow \lnot \pvar{TestimonyArguments.SolaFide}{28}\)
  \item[\(\vdash\)] \(\lnot \pvar{TestimonyArguments.SolaFide}{28}\)\gloss{\emph{at issue, the denial of:} Salvation is not by works: no work is its ground}
\end{derivation}
\par\medskip\noindent\textbf{\texttt{matthew7\allowbreak{}Faith\allowbreak{}Line}} --- Matthew 7:21–23, read as the doing including faith\par\nobreak\noindent
\textbf{Matthew 7:21–23, as Calvin reads it}: the doing includes believing, so obedience is not a condition of entry in its own right.

\begin{derivation}
  \item[(1)] \(\pvar{TestimonyArguments.SolaFide}{62}\)\gloss{Matthew 7:21–23 — the criterion at the judgment is doing the Father’s will, not profession or mighty works}
  \item[(2)] \(\pvar{TestimonyArguments.SolaFide}{64}\)\gloss{In Matthew 7:21–23, doing the Father’s will includes believing in Christ, and “I never knew you” means no union with him by faith}
  \item[(3)] \((\pvar{TestimonyArguments.SolaFide}{62} \land \pvar{TestimonyArguments.SolaFide}{64}) \rightarrow \lnot \pvar{TestimonyArguments.SolaFide}{63}\)
  \item[\(\vdash\)] \(\lnot \pvar{TestimonyArguments.SolaFide}{63}\)\gloss{\emph{at issue, the denial of:} In Matthew 7:21–23, doing the Father’s will is a condition of entering the kingdom in its own right, not only the fruit of faith}
\end{derivation}
\par\medskip\noindent\textbf{\texttt{matthew7\allowbreak{}Obedience\allowbreak{}Case}} --- Matthew 7:21–23, read as obedience a condition of entry\par\nobreak\noindent
Rome’s reading of Matthew 7:21–23, as a position in the debate.

\begin{derivation}
  \item[(1)] \(\pvar{TestimonyArguments.SolaFide}{62}\)\gloss{Matthew 7:21–23 — the criterion at the judgment is doing the Father’s will, not profession or mighty works}
  \item[(2)] \(\pvar{TestimonyArguments.SolaFide}{63}\)\gloss{In Matthew 7:21–23, doing the Father’s will is a condition of entering the kingdom in its own right, not only the fruit of faith}
  \item[(3)] \((\pvar{TestimonyArguments.SolaFide}{62} \land \pvar{TestimonyArguments.SolaFide}{63}) \rightarrow \lnot \pvar{TestimonyArguments.SolaFide}{28}\)
  \item[\(\vdash\)] \(\lnot \pvar{TestimonyArguments.SolaFide}{28}\)\gloss{\emph{at issue, the denial of:} Salvation is not by works: no work is its ground}
  \keepwithnext
\end{derivation}
\emph{No premise here rests on scripture alone.}
\par\medskip\noindent\textbf{\texttt{matthew7\allowbreak{}Faith\allowbreak{}Case}} --- Matthew 7:21–23, read as the doing including faith\par\nobreak\noindent
Calvin’s reading of Matthew 7:21–23, as a position in the debate.

\begin{derivation}
  \item[(1)] \(\pvar{TestimonyArguments.SolaFide}{62}\)\gloss{Matthew 7:21–23 — the criterion at the judgment is doing the Father’s will, not profession or mighty works}
  \item[(2)] \(\pvar{TestimonyArguments.SolaFide}{64}\)\gloss{In Matthew 7:21–23, doing the Father’s will includes believing in Christ, and “I never knew you” means no union with him by faith}
  \item[(3)] \((\pvar{TestimonyArguments.SolaFide}{62} \land \pvar{TestimonyArguments.SolaFide}{64}) \rightarrow \lnot \pvar{TestimonyArguments.SolaFide}{63}\)
  \item[\(\vdash\)] \(\lnot \pvar{TestimonyArguments.SolaFide}{63}\)\gloss{\emph{at issue, the denial of:} In Matthew 7:21–23, doing the Father’s will is a condition of entering the kingdom in its own right, not only the fruit of faith}
  \keepwithnext
\end{derivation}
\emph{No premise here rests on scripture alone.}
\paragraph*{Matthew 19:16–22}
\leavevmode\par\nobreak
\par\medskip\noindent\textbf{\texttt{matthew19\allowbreak{}To\allowbreak{}Works}}\par\nobreak\noindent
The Catholic step from Matthew 19:17. If keeping the commandments is the way to life, then salvation is not apart from works. Rated \texttt{disputed}: Calvin accepts the text and denies that it is the way for any sinner.

\begin{derivation}
  \item[] \((\pvar{TestimonyArguments.SolaFide}{65} \land \pvar{TestimonyArguments.SolaFide}{66}) \rightarrow \lnot \pvar{TestimonyArguments.SolaFide}{28}\)\gloss{\emph{would deny:} Salvation is not by works: no work is its ground}
\end{derivation}
\par\medskip\noindent\textbf{\texttt{matthew19\allowbreak{}To\allowbreak{}Legal\allowbreak{}Reply}}\par\nobreak\noindent
The Reformed step from Matthew 19:17. If Christ answers on the law’s own terms, to show the young man what he has not kept, then the commandments are not offered as the way to life. Rated \texttt{disputed}: \emph{Veritatis Splendor} accepts the text and reads it as the way.

\begin{derivation}
  \item[] \((\pvar{TestimonyArguments.SolaFide}{65} \land \pvar{TestimonyArguments.SolaFide}{67}) \rightarrow \lnot \pvar{TestimonyArguments.SolaFide}{66}\)\gloss{\emph{would deny:} In Matthew 19:17, keeping the commandments is the way to life, made possible by grace}
\end{derivation}
\par\medskip\noindent\textbf{\texttt{matthew19\allowbreak{}Way\allowbreak{}Line}} --- Matthew 19:17, read as the commandments the way to life\par\nobreak\noindent
\textbf{Matthew 19:17, as Rome reads it}: the commandments are the way to life, so salvation is not apart from works.

\begin{derivation}
  \item[(1)] \(\pvar{TestimonyArguments.SolaFide}{65}\)\gloss{Matthew 19:17 — “If you would enter life, keep the commandments”}
  \item[(2)] \(\pvar{TestimonyArguments.SolaFide}{66}\)\gloss{In Matthew 19:17, keeping the commandments is the way to life, made possible by grace}
  \item[(3)] \((\pvar{TestimonyArguments.SolaFide}{65} \land \pvar{TestimonyArguments.SolaFide}{66}) \rightarrow \lnot \pvar{TestimonyArguments.SolaFide}{28}\)
  \item[\(\vdash\)] \(\lnot \pvar{TestimonyArguments.SolaFide}{28}\)\gloss{\emph{at issue, the denial of:} Salvation is not by works: no work is its ground}
\end{derivation}
\par\medskip\noindent\textbf{\texttt{matthew19\allowbreak{}Legal\allowbreak{}Line}} --- Matthew 19:17, read as the law exposing what he has not kept\par\nobreak\noindent
\textbf{Matthew 19:17, as Calvin reads it}: a reply on the law’s own terms, showing the young man what he has not kept.

\begin{derivation}
  \item[(1)] \(\pvar{TestimonyArguments.SolaFide}{65}\)\gloss{Matthew 19:17 — “If you would enter life, keep the commandments”}
  \item[(2)] \(\pvar{TestimonyArguments.SolaFide}{67}\)\gloss{In Matthew 19:17, Christ answers on the law’s own terms, so that the young man may see he has not kept it and turn to faith}
  \item[(3)] \((\pvar{TestimonyArguments.SolaFide}{65} \land \pvar{TestimonyArguments.SolaFide}{67}) \rightarrow \lnot \pvar{TestimonyArguments.SolaFide}{66}\)
  \item[\(\vdash\)] \(\lnot \pvar{TestimonyArguments.SolaFide}{66}\)\gloss{\emph{at issue, the denial of:} In Matthew 19:17, keeping the commandments is the way to life, made possible by grace}
\end{derivation}
\par\medskip\noindent\textbf{\texttt{matthew19\allowbreak{}Way\allowbreak{}Case}} --- Matthew 19:17, read as the commandments the way to life\par\nobreak\noindent
Rome’s reading of Matthew 19:17, as a position in the debate.

\begin{derivation}
  \item[(1)] \(\pvar{TestimonyArguments.SolaFide}{65}\)\gloss{Matthew 19:17 — “If you would enter life, keep the commandments”}
  \item[(2)] \(\pvar{TestimonyArguments.SolaFide}{66}\)\gloss{In Matthew 19:17, keeping the commandments is the way to life, made possible by grace}
  \item[(3)] \((\pvar{TestimonyArguments.SolaFide}{65} \land \pvar{TestimonyArguments.SolaFide}{66}) \rightarrow \lnot \pvar{TestimonyArguments.SolaFide}{28}\)
  \item[\(\vdash\)] \(\lnot \pvar{TestimonyArguments.SolaFide}{28}\)\gloss{\emph{at issue, the denial of:} Salvation is not by works: no work is its ground}
  \keepwithnext
\end{derivation}
\emph{No premise here rests on scripture alone.}
\par\medskip\noindent\textbf{\texttt{matthew19\allowbreak{}Legal\allowbreak{}Case}} --- Matthew 19:17, read as the law exposing what he has not kept\par\nobreak\noindent
Calvin’s reading of Matthew 19:17, as a position in the debate.

\begin{derivation}
  \item[(1)] \(\pvar{TestimonyArguments.SolaFide}{65}\)\gloss{Matthew 19:17 — “If you would enter life, keep the commandments”}
  \item[(2)] \(\pvar{TestimonyArguments.SolaFide}{67}\)\gloss{In Matthew 19:17, Christ answers on the law’s own terms, so that the young man may see he has not kept it and turn to faith}
  \item[(3)] \((\pvar{TestimonyArguments.SolaFide}{65} \land \pvar{TestimonyArguments.SolaFide}{67}) \rightarrow \lnot \pvar{TestimonyArguments.SolaFide}{66}\)
  \item[\(\vdash\)] \(\lnot \pvar{TestimonyArguments.SolaFide}{66}\)\gloss{\emph{at issue, the denial of:} In Matthew 19:17, keeping the commandments is the way to life, made possible by grace}
  \keepwithnext
\end{derivation}
\emph{No premise here rests on scripture alone.}
\paragraph*{The two readings, written out}
\leavevmode\par\nobreak
\par\medskip\noindent\textbf{\texttt{rome\allowbreak{}On\allowbreak{}Jesus\allowbreak{}Words\allowbreak{}Reading}}\par\nobreak\noindent
Rome’s reading of Jesus’ words. In Matthew 7 obedience is a condition of entry in its own right. In Matthew 19 the commandments are the way to life. So salvation is not apart from works, and faith, without the dispositions and the charity Trent requires, does not suffice.

{\small\begin{verbatim}
def romeOnJesusWordsReading : Valuation Claim :=
  fun a =>
    match a with
    | Claim.paulsJustifyDenotesRenewal => False
    | Claim.rev22_11BeJustifiedStill => False
    | Claim.sir18_22BeJustifiedToDeath => False
    | Claim.jamesUsesDikaioAsPaul => False
    | Claim.salvationNotByWorks => False
    | Claim.justificationIsForensicOnly => False
    | Claim.justificationDistinctFromSanctification => False
    | Claim.faithIsSufficient => False
    | Claim.justificationByFaithAlone => False
    | Claim.matthew7DoingIncludesBelieving => False
    | Claim.matthew19LawExposesInability => False
    | x => True
\end{verbatim}}
\par\medskip\noindent\textbf{\texttt{calvin\allowbreak{}On\allowbreak{}Jesus\allowbreak{}Words\allowbreak{}Reading}}\par\nobreak\noindent
Calvin’s reading of Jesus’ words. In Matthew 7 the doing includes believing and is not a condition of entry in its own right. In Matthew 19 Christ replies on the law’s own terms and does not give the way to life. In Luke the plea and the faith sufficed.

{\small\begin{verbatim}
def calvinOnJesusWordsReading : Valuation Claim :=
  fun a =>
    match a with
    | Claim.matthew7ObedienceIsAGround => False
    | Claim.matthew19CommandmentsAreTheWayToLife => False
    | x => True
\end{verbatim}}
\paragraph*{What each position holds}
\leavevmode\par\nobreak
\par\medskip\noindent\textbf{\texttt{luke18\allowbreak{}Case\_\allowbreak{}establishes}}\par\nobreak\noindent
\textbf{Luke 18 reaches “faith suffices”}, given its step.

{\small\begin{verbatim}
theorem luke18Case_establishes : Establishes luke18Case
-- axioms: propext, Quot.sound
\end{verbatim}}
\par\medskip\noindent\textbf{\texttt{luke18\allowbreak{}Case\_\allowbreak{}is\_\allowbreak{}satisfiable}}\par\nobreak\noindent
Luke 18’s case can be held without contradiction.

{\small\begin{verbatim}
theorem luke18Case_is_satisfiable : Satisfiable luke18Case.premises
-- axioms: propext, Quot.sound
\end{verbatim}}
\par\medskip\noindent\textbf{\texttt{trent\_\allowbreak{}grants\_\allowbreak{}the\_\allowbreak{}parable}}\par\nobreak\noindent
\textbf{Trent’s case is consistent with every word of the parable.} Trent never cites Luke 18. But nothing in its decree denies that the tax collector pleaded for mercy, brought no works that earn anything, and went home justified by God’s verdict. Its case can be held with all of that, on Rome’s reading, where faith without the dispositions does not suffice.

What Trent’s case denies, when applied to the parable, is the step from it, not the text.

{\small\begin{verbatim}
theorem trent_grants_the_parable : Grants tridentineCase (⋀[p
    Claim.luke18_9_14TaxCollectorJustified, p Claim.luke18JustifiedIsVerdict,
    p Claim.taxCollectorBringsNoWorks])
-- axioms: propext, Quot.sound
\end{verbatim}}
\par\medskip\noindent\textbf{\texttt{matthew7\allowbreak{}Obedience\allowbreak{}Case\_\allowbreak{}establishes}}\par\nobreak\noindent
Rome’s reading of Matthew 7 reaches its conclusion.

{\small\begin{verbatim}
theorem matthew7ObedienceCase_establishes : Establishes matthew7ObedienceCase
-- axioms: propext, Quot.sound
\end{verbatim}}
\par\medskip\noindent\textbf{\texttt{matthew7\allowbreak{}Obedience\allowbreak{}Case\_\allowbreak{}is\_\allowbreak{}satisfiable}}\par\nobreak\noindent
Rome’s reading of Matthew 7 can be held without contradiction.

{\small\begin{verbatim}
theorem matthew7ObedienceCase_is_satisfiable : Satisfiable
    matthew7ObedienceCase.premises
-- axioms: propext, Quot.sound
\end{verbatim}}
\par\medskip\noindent\textbf{\texttt{matthew7\allowbreak{}Faith\allowbreak{}Case\_\allowbreak{}establishes}}\par\nobreak\noindent
Calvin’s reading of Matthew 7 reaches its conclusion.

{\small\begin{verbatim}
theorem matthew7FaithCase_establishes : Establishes matthew7FaithCase
-- axioms: propext, Quot.sound
\end{verbatim}}
\par\medskip\noindent\textbf{\texttt{matthew7\allowbreak{}Faith\allowbreak{}Case\_\allowbreak{}is\_\allowbreak{}satisfiable}}\par\nobreak\noindent
Calvin’s reading of Matthew 7 can be held without contradiction.

{\small\begin{verbatim}
theorem matthew7FaithCase_is_satisfiable : Satisfiable
    matthew7FaithCase.premises
-- axioms: propext, Quot.sound
\end{verbatim}}
\par\medskip\noindent\textbf{\texttt{matthew19\allowbreak{}Way\allowbreak{}Case\_\allowbreak{}establishes}}\par\nobreak\noindent
Rome’s reading of Matthew 19 reaches its conclusion.

{\small\begin{verbatim}
theorem matthew19WayCase_establishes : Establishes matthew19WayCase
-- axioms: propext, Quot.sound
\end{verbatim}}
\par\medskip\noindent\textbf{\texttt{matthew19\allowbreak{}Way\allowbreak{}Case\_\allowbreak{}is\_\allowbreak{}satisfiable}}\par\nobreak\noindent
Rome’s reading of Matthew 19 can be held without contradiction.

{\small\begin{verbatim}
theorem matthew19WayCase_is_satisfiable : Satisfiable
    matthew19WayCase.premises
-- axioms: propext, Quot.sound
\end{verbatim}}
\par\medskip\noindent\textbf{\texttt{matthew19\allowbreak{}Legal\allowbreak{}Case\_\allowbreak{}establishes}}\par\nobreak\noindent
Calvin’s reading of Matthew 19 reaches its conclusion.

{\small\begin{verbatim}
theorem matthew19LegalCase_establishes : Establishes matthew19LegalCase
-- axioms: propext, Quot.sound
\end{verbatim}}
\par\medskip\noindent\textbf{\texttt{matthew19\allowbreak{}Legal\allowbreak{}Case\_\allowbreak{}is\_\allowbreak{}satisfiable}}\par\nobreak\noindent
Calvin’s reading of Matthew 19 can be held without contradiction.

{\small\begin{verbatim}
theorem matthew19LegalCase_is_satisfiable : Satisfiable
    matthew19LegalCase.premises
-- axioms: propext, Quot.sound
\end{verbatim}}
\paragraph*{What the texts leave open}
\leavevmode\par\nobreak
\par\medskip\noindent\textbf{\texttt{matthew7\_\allowbreak{}leaves\_\allowbreak{}the\_\allowbreak{}ground\_\allowbreak{}open}}\par\nobreak\noindent
\textbf{Matthew 7:21–23 leaves open whether obedience is a condition of entry in its own right.} Its words say that the test is doing the Father’s will. They do not say that the doing is such a condition, and they do not say that it is only the fruit of faith. Both readings keep the words, and differ only in what they add to them.

{\small\begin{verbatim}
theorem matthew7_leaves_the_ground_open : Independent [p
    Claim.matthew7_21_23DoingTheWill] (p Claim.matthew7ObedienceIsAGround)
-- axioms: propext, Quot.sound
\end{verbatim}}
\par\medskip\noindent\textbf{\texttt{matthew19\_\allowbreak{}leaves\_\allowbreak{}the\_\allowbreak{}way\_\allowbreak{}open}}\par\nobreak\noindent
\textbf{Matthew 19:17 leaves open whether the commandments are the way to life.} “If you would enter life, keep the commandments” is kept alike by the reading on which it is the way, and by the reading on which it is a reply on the law’s own terms.

{\small\begin{verbatim}
theorem matthew19_leaves_the_way_open : Independent [p
    Claim.matthew19_17KeepTheCommandments] (p
    Claim.matthew19CommandmentsAreTheWayToLife)
-- axioms: propext, Quot.sound
\end{verbatim}}
\paragraph*{Why Luke 18 stands against Trent}
\leavevmode\par\nobreak
\par\medskip\noindent\textbf{\texttt{why\allowbreak{}Luke\allowbreak{}Eighteen\allowbreak{}Stands\allowbreak{}Against\allowbreak{}Trent}}\par\nobreak\noindent
\textbf{Why Luke 18 stands against Trent.} Everything turns on Luke 18’s own step: from the parable, the verdict, and the plea without works, to “faith suffices”.

Trent cannot hold that step. Two things in its case contradict it, and both are needed: its definition of justification as the renewal of the inward person (Session VI, ch. 7), and the step from that definition to denying that faith suffices (canon 9). Nothing else in Trent’s case does. So Luke 18 meets Trent exactly where Luke 7:50 does (\texttt{why\allowbreak{}The\allowbreak{}Dominical\allowbreak{}Case\allowbreak{}Stands\allowbreak{}Against\allowbreak{}Trent}): not over what the parable says, which Trent’s case is consistent with, but over what justification is.

The dispositions of chapter 6 are why the step is rated \texttt{disputed}. They are not a premise of Trent’s case here. What can be shown against Trent’s definition from the meaning of the word, and what cannot, is checked in \texttt{what\allowbreak{}Trents\allowbreak{}Definition\allowbreak{}Claims}.

\paragraph{Why \emph{Luke 18:14: “this man went down to his house justified”} stands against \emph{Trent, against ‘not by works’}.}
\begin{itemize}
\item \textbf{The crux:} \((\pvar{TestimonyArguments.SolaFide}{59} \land \pvar{TestimonyArguments.SolaFide}{60} \land \pvar{TestimonyArguments.SolaFide}{61}) \rightarrow \pvar{TestimonyArguments.SolaFide}{26}\), a premise of \emph{Luke 18:14: “this man went down to his house justified”}.
\item \textbf{What it does:} its conclusion does not follow without it.
\item \textbf{Granted:} \(\pvar{TestimonyArguments.SolaFide}{59}\), \(\pvar{TestimonyArguments.SolaFide}{60}\), \(\pvar{TestimonyArguments.SolaFide}{61}\), which the break also needs.
\item \textbf{Where the rival breaks:} \(\pvar{TestimonyArguments.SolaFide}{33}\), \(\pvar{TestimonyArguments.SolaFide}{33} \rightarrow \lnot \pvar{TestimonyArguments.SolaFide}{26}\) cannot be held together with \((\pvar{TestimonyArguments.SolaFide}{59} \land \pvar{TestimonyArguments.SolaFide}{60} \land \pvar{TestimonyArguments.SolaFide}{61}) \rightarrow \pvar{TestimonyArguments.SolaFide}{26}\) and \(\pvar{TestimonyArguments.SolaFide}{59}\), \(\pvar{TestimonyArguments.SolaFide}{60}\), \(\pvar{TestimonyArguments.SolaFide}{61}\); each is needed for the break, and without the crux they stand.
\item \textbf{What this rests on:}
\begin{itemize}
\item \(\pvar{TestimonyArguments.SolaFide}{59}\) Luke 18:9–14 — the tax collector pleads for mercy and goes home justified, rather than the Pharisee who listed his works --- \emph{consensus}: Luke 18:9-14; \cite[201--207]{calvin-harmony-evangelists-2-1845}
\item \(\pvar{TestimonyArguments.SolaFide}{60}\) ‘Justified’ at Luke 18:14 is God’s verdict: accepted as righteous before God --- \emph{well supported}: \cite[206--207]{calvin-harmony-evangelists-2-1845}; Luke 18:14
\item \(\pvar{TestimonyArguments.SolaFide}{61}\) The tax collector brings no works that merit, only a humble and contrite plea for mercy --- \emph{consensus}: Luke 18:9-14; \cite[206--207]{calvin-harmony-evangelists-2-1845}; \cite[36]{waterworth-trent-1848}
\item the step itself --- \emph{disputed}: \cite[45]{waterworth-trent-1848}; \cite[33--34]{waterworth-trent-1848}
\end{itemize}
\end{itemize}
\par\medskip\noindent\textbf{\texttt{luke18\allowbreak{}Case\_\allowbreak{}strength}}\par\nobreak\noindent
Luke 18’s case rests on a step Trent contests, so its weakest link is \texttt{disputed}.

{\small\begin{verbatim}
theorem luke18Case_strength : luke18Case.strength = 0
-- axioms: propext
\end{verbatim}}
\par\medskip\noindent\textbf{\texttt{matthew7\allowbreak{}Obedience\allowbreak{}Case\_\allowbreak{}strength}}\par\nobreak\noindent
Rome’s reading of Matthew 7 rests on that reading, so its weakest link is \texttt{disputed}.

{\small\begin{verbatim}
theorem matthew7ObedienceCase_strength : matthew7ObedienceCase.strength = 0
-- axioms: propext
\end{verbatim}}
\par\medskip\noindent\textbf{\texttt{matthew7\allowbreak{}Faith\allowbreak{}Case\_\allowbreak{}strength}}\par\nobreak\noindent
Calvin’s reading of Matthew 7 rests on his, so its weakest link is \texttt{disputed}.

{\small\begin{verbatim}
theorem matthew7FaithCase_strength : matthew7FaithCase.strength = 0
-- axioms: propext
\end{verbatim}}
\par\medskip\noindent\textbf{\texttt{matthew19\allowbreak{}Way\allowbreak{}Case\_\allowbreak{}strength}}\par\nobreak\noindent
Rome’s reading of Matthew 19 rests on that reading, so its weakest link is \texttt{disputed}.

{\small\begin{verbatim}
theorem matthew19WayCase_strength : matthew19WayCase.strength = 0
-- axioms: propext
\end{verbatim}}
\par\medskip\noindent\textbf{\texttt{matthew19\allowbreak{}Legal\allowbreak{}Case\_\allowbreak{}strength}}\par\nobreak\noindent
Calvin’s reading of Matthew 19 rests on his, so its weakest link is \texttt{disputed}.

{\small\begin{verbatim}
theorem matthew19LegalCase_strength : matthew19LegalCase.strength = 0
-- axioms: propext
\end{verbatim}}
\paragraph*{The debate over Jesus’ words}
\leavevmode\par\nobreak
\par\medskip\noindent\textbf{\texttt{words\allowbreak{}Party\allowbreak{}Node}}\par\nobreak\noindent
The premises each position argues from.

{\small\begin{verbatim}
def wordsPartyNode : WordsParty → ArgumentPackage Claim :=
  fun x =>
    match x with
    | WordsParty.luke7 => dominicalCase
    | WordsParty.luke18 => luke18Case
    | WordsParty.trent => tridentineCase
    | WordsParty.matthew7Obedience => matthew7ObedienceCase
    | WordsParty.matthew7Faith => matthew7FaithCase
    | WordsParty.matthew19Way => matthew19WayCase
    | WordsParty.matthew19Legal => matthew19LegalCase
\end{verbatim}}
\par\medskip\noindent\textbf{\texttt{jesus\allowbreak{}Words\allowbreak{}Dispute}}\par\nobreak\noindent
The debate over Jesus’ words. For every position heard: its premises can be held without contradiction, its conclusion follows from them, and every step it uses is rated. Each position’s weakest link is \texttt{disputed}.

{\small\begin{verbatim}
def jesusWordsDispute : Dispute Claim WordsParty :=
  { node := wordsPartyNode,
    consistent := jesusWordsDispute._proof_1,
    sound := jesusWordsDispute._proof_2,
    rated := jesusWordsDispute._proof_10 }
\end{verbatim}}
\par\medskip\noindent\textbf{\texttt{jesus\allowbreak{}Words\allowbreak{}Dispute\_\allowbreak{}defeats}}\par\nobreak\noindent
\textbf{Who defeats whom}, for all 49 pairs: four ties, and nothing else. Luke 7:50 and Trent defeat each other, and so do Luke 18 and Trent, over Trent’s definition of justification. Each reading of Matthew defeats its rival and is defeated in return. Every answer is worked out from the positions’ premises and checked by machine.

{\small\begin{verbatim}
theorem jesusWordsDispute_defeats : ∀ (i j : WordsParty),
    jesusWordsDispute.defeats i j ↔ (i, j) ∈ [(WordsParty.luke7,
    WordsParty.trent), (WordsParty.trent, WordsParty.luke7),
    (WordsParty.luke18, WordsParty.trent), (WordsParty.trent,
    WordsParty.luke18), (WordsParty.matthew7Obedience,
    WordsParty.matthew7Faith), (WordsParty.matthew7Faith,
    WordsParty.matthew7Obedience), (WordsParty.matthew19Way,
    WordsParty.matthew19Legal), (WordsParty.matthew19Legal,
    WordsParty.matthew19Way)]
-- axioms: propext, Classical.choice, Quot.sound
\end{verbatim}}
\par\medskip\noindent\textbf{\texttt{jesus\allowbreak{}Words\allowbreak{}Dispute\_\allowbreak{}supports}}\par\nobreak\noindent
\textbf{No position supports another, and no position’s case is contained in another’s}, in any of the 49 pairs. Every answer is worked out and checked by machine.

{\small\begin{verbatim}
theorem jesusWordsDispute_supports : ∀ (i j : WordsParty),
    ¬jesusWordsDispute.supports i j
-- axioms: propext, Classical.choice, Quot.sound
\end{verbatim}}
\par\medskip\noindent\textbf{\texttt{jesus\allowbreak{}Words\allowbreak{}Dispute\_\allowbreak{}part\allowbreak{}Of}}\par\nobreak\noindent
Each position’s case contains itself and no other. Every answer is worked out and checked by machine.

{\small\begin{verbatim}
theorem jesusWordsDispute_partOf : ∀ (i j : WordsParty),
    jesusWordsDispute.partOf i j ↔ i = j ∨ (i, j) ∈ []
-- axioms: propext, Classical.choice, Quot.sound
\end{verbatim}}
\par\medskip\noindent\textbf{\texttt{jesus\allowbreak{}Words\allowbreak{}Map}}\par\nobreak\noindent
The debate over Jesus’ words drawn as a map: who defeats whom.

\begin{center}
\begin{tikzpicture}[>={Stealth[length=2mm]}]
\node[draw,circle,minimum size=6mm,inner sep=0pt] (n0) at (0.00,3.20) {1};
\node[draw,circle,minimum size=6mm,inner sep=0pt] (n1) at (2.50,2.00) {2};
\node[draw,circle,minimum size=6mm,inner sep=0pt] (n2) at (3.12,-0.71) {3};
\node[draw,circle,minimum size=6mm,inner sep=0pt] (n3) at (1.39,-2.88) {4};
\node[draw,circle,minimum size=6mm,inner sep=0pt] (n4) at (-1.39,-2.88) {5};
\node[draw,circle,minimum size=6mm,inner sep=0pt] (n5) at (-3.12,-0.71) {6};
\node[draw,circle,minimum size=6mm,inner sep=0pt] (n6) at (-2.50,2.00) {7};
\draw[->,red!70!black,thick] (n0) to[bend left=10] (n2);
\draw[->,red!70!black,thick] (n1) to[bend left=10] (n2);
\draw[->,red!70!black,thick] (n2) to[bend left=10] (n0);
\draw[->,red!70!black,thick] (n2) to[bend left=10] (n1);
\draw[->,red!70!black,thick] (n3) to[bend left=10] (n4);
\draw[->,red!70!black,thick] (n4) to[bend left=10] (n3);
\draw[->,red!70!black,thick] (n5) to[bend left=10] (n6);
\draw[->,red!70!black,thick] (n6) to[bend left=10] (n5);
\end{tikzpicture}
\end{center}
\noindent Solid: defeats. Dashed: supports. Dotted: one party's case is part of another's.
\begin{enumerate}
\item Luke 7:50: “your faith has saved you”
\item Luke 18:14: “this man went down to his house justified”
\item Trent, against ‘not by works’
\item Matthew 7:21–23, read as obedience a condition of entry
\item Matthew 7:21–23, read as the doing including faith
\item Matthew 19:17, read as the commandments the way to life
\item Matthew 19:17, read as the law exposing what he has not kept
\end{enumerate}
\begin{longtable}{@{}>{\raggedright\arraybackslash}p{0.3\textwidth} >{\raggedright\arraybackslash}p{0.3\textwidth} >{\raggedright\arraybackslash}p{0.3\textwidth}@{}}
\toprule
\textbf{From} & \textbf{To} & \textbf{Edge} \\
\midrule
\endhead
1 \emph{Luke 7:50: “your faith has saved you”} & 3 \emph{Trent, against ‘not by works’} & defeats \\
2 \emph{Luke 18:14: “this man went down to his house justified”} & 3 \emph{Trent, against ‘not by works’} & defeats \\
3 \emph{Trent, against ‘not by works’} & 1 \emph{Luke 7:50: “your faith has saved you”} & defeats \\
3 \emph{Trent, against ‘not by works’} & 2 \emph{Luke 18:14: “this man went down to his house justified”} & defeats \\
4 \emph{Matthew 7:21–23, read as obedience a condition of entry} & 5 \emph{Matthew 7:21–23, read as the doing including faith} & defeats \\
5 \emph{Matthew 7:21–23, read as the doing including faith} & 4 \emph{Matthew 7:21–23, read as obedience a condition of entry} & defeats \\
6 \emph{Matthew 19:17, read as the commandments the way to life} & 7 \emph{Matthew 19:17, read as the law exposing what he has not kept} & defeats \\
7 \emph{Matthew 19:17, read as the law exposing what he has not kept} & 6 \emph{Matthew 19:17, read as the commandments the way to life} & defeats \\
\bottomrule
\end{longtable}
\noindent The defeats are the cells of \texttt{jesus\allowbreak{}Words\allowbreak{}Dispute\_\allowbreak{}defeats}, the supports the cells of \texttt{jesus\allowbreak{}Words\allowbreak{}Dispute\_\allowbreak{}supports} and the parts the cells of \texttt{jesus\allowbreak{}Words\allowbreak{}Dispute\_\allowbreak{}part\allowbreak{}Of}, each computed from the parties' premises and checked by the kernel. The attacks derived through support and through parts are reported, not counted: every verdict is computed from the defeats alone. Every derived attack is already a defeat.
\par\medskip\noindent\textbf{\texttt{luke\allowbreak{}With\allowbreak{}Rome\allowbreak{}On\allowbreak{}Matthew\allowbreak{}Reading}}\par\nobreak\noindent
\textbf{Rome’s readings of Matthew can be held together with Luke’s sayings.} One consistent reading holds Luke 7:50 and Luke 18, that faith sufficed, and Rome’s readings of Matthew 7 and 19, that salvation is not apart from works. As written out here, “faith suffices” and “not apart from works” are different claims, and Luke’s sayings make only the first. So whoever holds both is not contradicted by these texts. This reading takes neither of Trent’s further claims: baptism as the instrument of justification, and justification as including sanctification.

{\small\begin{verbatim}
def lukeWithRomeOnMatthewReading : Valuation Claim :=
  fun a =>
    match a with
    | Claim.paulsJustifyDenotesRenewal => False
    | Claim.rev22_11BeJustifiedStill => False
    | Claim.sir18_22BeJustifiedToDeath => False
    | Claim.jamesUsesDikaioAsPaul => False
    | Claim.salvationNotByWorks => False
    | Claim.matthew7DoingIncludesBelieving => False
    | Claim.matthew19LawExposesInability => False
    | Claim.baptismIsInstrumentalCause => False
    | Claim.justificationIncludesSanctification => False
    | x => True
\end{verbatim}}
\par\medskip\noindent\textbf{\texttt{luke\_\allowbreak{}stands\_\allowbreak{}with\_\allowbreak{}rome\_\allowbreak{}on\_\allowbreak{}matthew}}\par\nobreak\noindent
Luke’s two sayings can be held together with Rome’s readings of Matthew.

{\small\begin{verbatim}
theorem luke_stands_with_rome_on_matthew : jesusWordsDispute.StandTogether
    [WordsParty.luke7, WordsParty.luke18, WordsParty.matthew7Obedience,
    WordsParty.matthew19Way]
-- axioms: propext, Classical.choice, Quot.sound
\end{verbatim}}
\par\medskip\noindent\textbf{\texttt{jesus\allowbreak{}Words\allowbreak{}Settle\allowbreak{}Nothing}}\par\nobreak\noindent
Why nothing is forced: every position has an attacker. Luke 7:50 and Luke 18 each have Trent, Trent has Luke 7:50, and each reading of Matthew has its rival.

\paragraph{Nothing prevails outright: every party is defeated by another.}
\begin{itemize}
\item \emph{Trent, against ‘not by works’} defeats \emph{Luke 7:50: “your faith has saved you”}.
\item \emph{Trent, against ‘not by works’} defeats \emph{Luke 18:14: “this man went down to his house justified”}.
\item \emph{Luke 7:50: “your faith has saved you”} defeats \emph{Trent, against ‘not by works’} (\texttt{why\allowbreak{}The\allowbreak{}Dominical\allowbreak{}Case\allowbreak{}Stands\allowbreak{}Against\allowbreak{}Trent}).
\item \emph{Matthew 7:21–23, read as the doing including faith} defeats \emph{Matthew 7:21–23, read as obedience a condition of entry}.
\item \emph{Matthew 7:21–23, read as obedience a condition of entry} defeats \emph{Matthew 7:21–23, read as the doing including faith}.
\item \emph{Matthew 19:17, read as the law exposing what he has not kept} defeats \emph{Matthew 19:17, read as the commandments the way to life}.
\item \emph{Matthew 19:17, read as the commandments the way to life} defeats \emph{Matthew 19:17, read as the law exposing what he has not kept}.
\end{itemize}
\par\noindent\textbf{What this rests on.} The reasons state 7 defeats and 0 absences of defeat, each a cell of the defeat table, \texttt{jesus\allowbreak{}Words\allowbreak{}Dispute\_\allowbreak{}defeats}, computed from the two parties' premises and checked by the kernel. Whether an attack survives turns on the attacker's weakest link and on the rating of what it attacks, so a changed rating can change the verdict. The weakest links:
\begin{itemize}
\item \emph{Trent, against ‘not by works’}, weakest at \emph{disputed}:
\begin{itemize}
\item Justification is not remission of sins only, but renewal of the inward man --- \cite[Trent, Session VI (1547), Decree on Justification, ch. 7]{tanner-decrees-1990}; \cite[Trent, Session VI (1547), canon 11]{tanner-decrees-1990}
\end{itemize}
\item \emph{Luke 7:50: “your faith has saved you”}, weakest at \emph{disputed}:
\begin{itemize}
\item Luke 7:47 — her love is the evidence of her forgiveness, not its ground --- \cite[ad loc. Luke 7:47]{marshall-luke-1978}; \cite[54--54]{padilla-narrative-criticism-2021}; \cite[Apology IV, on love and the fulfilling of the law: Luke 7:47–50]{kolb-wengert-book-of-concord-2000}; Luke 7:41-43; Luke 7:47; Luke 7:47-50
\item an inference step --- \cite[Trent, Session VI (1547), canon 9]{tanner-decrees-1990}
\end{itemize}
\item \emph{Luke 18:14: “this man went down to his house justified”}, weakest at \emph{disputed}:
\begin{itemize}
\item an inference step --- \cite[45]{waterworth-trent-1848}; \cite[33--34]{waterworth-trent-1848}
\end{itemize}
\item \emph{Matthew 7:21–23, read as the doing including faith}, weakest at \emph{disputed}:
\begin{itemize}
\item In Matthew 7:21–23, doing the Father’s will includes believing in Christ, and “I never knew you” means no union with him by faith --- \cite[367--368]{calvin-harmony-evangelists-1845}
\item an inference step --- \cite[2826]{catholic-church-catechism-2000}
\end{itemize}
\item \emph{Matthew 7:21–23, read as obedience a condition of entry}, weakest at \emph{disputed}:
\begin{itemize}
\item In Matthew 7:21–23, doing the Father’s will is a condition of entering the kingdom in its own right, not only the fruit of faith --- \cite[2826]{catholic-church-catechism-2000}; \cite[1821]{catholic-church-catechism-2000}
\item an inference step --- \cite[367--368]{calvin-harmony-evangelists-1845}
\end{itemize}
\item \emph{Matthew 19:17, read as the law exposing what he has not kept}, weakest at \emph{disputed}:
\begin{itemize}
\item In Matthew 19:17, Christ answers on the law’s own terms, so that the young man may see he has not kept it and turn to faith --- \cite[393--394]{calvin-harmony-evangelists-2-1845}
\item an inference step --- \cite[12]{john-paul-ii-veritatis-splendor-1993}
\end{itemize}
\item \emph{Matthew 19:17, read as the commandments the way to life}, weakest at \emph{disputed}:
\begin{itemize}
\item In Matthew 19:17, keeping the commandments is the way to life, made possible by grace --- \cite[12]{john-paul-ii-veritatis-splendor-1993}; \cite[2052–2053]{catholic-church-catechism-2000}; \cite[35]{waterworth-trent-1848}
\item an inference step --- \cite[394]{calvin-harmony-evangelists-2-1845}
\end{itemize}
\end{itemize}
\noindent Why \emph{Luke 7:50: “your faith has saved you”} defeats \emph{Trent, against ‘not by works’}, at the level of the claims: \texttt{why\allowbreak{}The\allowbreak{}Dominical\allowbreak{}Case\allowbreak{}Stands\allowbreak{}Against\allowbreak{}Trent}.
\par\medskip\noindent\textbf{\texttt{jesus\_\allowbreak{}words\_\allowbreak{}force\_\allowbreak{}nothing}}\par\nobreak\noindent
\textbf{Heard together, Jesus’ words force nothing.} Every position is in a tie the debate cannot settle: each saying in Luke with Trent, over what justification is, and each reading of Matthew with its rival. Every weakest link is \texttt{disputed}, so no rating breaks a tie.

It does not show that Jesus’ words say nothing. Each position can be defended, and Trent accepts every word of Luke 18 (\texttt{trent\_\allowbreak{}grants\_\allowbreak{}the\_\allowbreak{}parable}). What the debate cannot do is choose between readings rated alike.

{\small\begin{verbatim}
theorem jesus_words_force_nothing : Framework.grounded
    jesusWordsDispute.defeats = ∅
-- axioms: propext, Classical.choice, Quot.sound
\end{verbatim}}
\par\medskip\noindent\textbf{\texttt{luke18\allowbreak{}Stands\allowbreak{}With\allowbreak{}Luke7}}\par\nobreak\noindent
Why Luke 18 can be defended: it can be held with Luke 7:50 and with Rome’s readings of Matthew, and Trent, its one attacker, is defeated from within that set, by Luke 7:50 and by Luke 18 itself.

\paragraph{\emph{Luke 18:14: “this man went down to his house justified”} is accepted on some resolution.}
\begin{itemize}
\item A position holding \emph{Luke 7:50: “your faith has saved you”}, \emph{Luke 18:14: “this man went down to his house justified”}, \emph{Matthew 7:21–23, read as obedience a condition of entry}, \emph{Matthew 19:17, read as the commandments the way to life} can be held.
\item None of \emph{Luke 7:50: “your faith has saved you”}, \emph{Luke 18:14: “this man went down to his house justified”}, \emph{Matthew 7:21–23, read as obedience a condition of entry}, \emph{Matthew 19:17, read as the commandments the way to life} defeats another, and each attack on them is answered from among them:
\begin{itemize}
\item \emph{Trent, against ‘not by works’} defeats \emph{Luke 7:50: “your faith has saved you”}, and \emph{Luke 7:50: “your faith has saved you”} defeats \emph{Trent, against ‘not by works’} (\texttt{why\allowbreak{}The\allowbreak{}Dominical\allowbreak{}Case\allowbreak{}Stands\allowbreak{}Against\allowbreak{}Trent}).
\item \emph{Trent, against ‘not by works’} defeats \emph{Luke 18:14: “this man went down to his house justified”}, and \emph{Luke 7:50: “your faith has saved you”} defeats \emph{Trent, against ‘not by works’} (\texttt{why\allowbreak{}The\allowbreak{}Dominical\allowbreak{}Case\allowbreak{}Stands\allowbreak{}Against\allowbreak{}Trent}).
\item \emph{Matthew 7:21–23, read as the doing including faith} defeats \emph{Matthew 7:21–23, read as obedience a condition of entry}, and \emph{Matthew 7:21–23, read as obedience a condition of entry} defeats \emph{Matthew 7:21–23, read as the doing including faith}.
\item \emph{Matthew 19:17, read as the law exposing what he has not kept} defeats \emph{Matthew 19:17, read as the commandments the way to life}, and \emph{Matthew 19:17, read as the commandments the way to life} defeats \emph{Matthew 19:17, read as the law exposing what he has not kept}.
\end{itemize}
\end{itemize}
\par\noindent\textbf{What this rests on.} The reasons state 7 defeats and 24 absences of defeat, each a cell of the defeat table, \texttt{jesus\allowbreak{}Words\allowbreak{}Dispute\_\allowbreak{}defeats}, computed from the two parties' premises and checked by the kernel. Whether an attack survives turns on the attacker's weakest link and on the rating of what it attacks, so a changed rating can change the verdict. The weakest links:
\begin{itemize}
\item \emph{Luke 7:50: “your faith has saved you”}, weakest at \emph{disputed}:
\begin{itemize}
\item Luke 7:47 — her love is the evidence of her forgiveness, not its ground --- \cite[ad loc. Luke 7:47]{marshall-luke-1978}; \cite[54--54]{padilla-narrative-criticism-2021}; \cite[Apology IV, on love and the fulfilling of the law: Luke 7:47–50]{kolb-wengert-book-of-concord-2000}; Luke 7:41-43; Luke 7:47; Luke 7:47-50
\item an inference step --- \cite[Trent, Session VI (1547), canon 9]{tanner-decrees-1990}
\end{itemize}
\item \emph{Luke 18:14: “this man went down to his house justified”}, weakest at \emph{disputed}:
\begin{itemize}
\item an inference step --- \cite[45]{waterworth-trent-1848}; \cite[33--34]{waterworth-trent-1848}
\end{itemize}
\item \emph{Matthew 7:21–23, read as obedience a condition of entry}, weakest at \emph{disputed}:
\begin{itemize}
\item In Matthew 7:21–23, doing the Father’s will is a condition of entering the kingdom in its own right, not only the fruit of faith --- \cite[2826]{catholic-church-catechism-2000}; \cite[1821]{catholic-church-catechism-2000}
\item an inference step --- \cite[367--368]{calvin-harmony-evangelists-1845}
\end{itemize}
\item \emph{Matthew 19:17, read as the commandments the way to life}, weakest at \emph{disputed}:
\begin{itemize}
\item In Matthew 19:17, keeping the commandments is the way to life, made possible by grace --- \cite[12]{john-paul-ii-veritatis-splendor-1993}; \cite[2052–2053]{catholic-church-catechism-2000}; \cite[35]{waterworth-trent-1848}
\item an inference step --- \cite[394]{calvin-harmony-evangelists-2-1845}
\end{itemize}
\item \emph{Trent, against ‘not by works’}, weakest at \emph{disputed}:
\begin{itemize}
\item Justification is not remission of sins only, but renewal of the inward man --- \cite[Trent, Session VI (1547), Decree on Justification, ch. 7]{tanner-decrees-1990}; \cite[Trent, Session VI (1547), canon 11]{tanner-decrees-1990}
\end{itemize}
\item \emph{Matthew 7:21–23, read as the doing including faith}, weakest at \emph{disputed}:
\begin{itemize}
\item In Matthew 7:21–23, doing the Father’s will includes believing in Christ, and “I never knew you” means no union with him by faith --- \cite[367--368]{calvin-harmony-evangelists-1845}
\item an inference step --- \cite[2826]{catholic-church-catechism-2000}
\end{itemize}
\item \emph{Matthew 19:17, read as the law exposing what he has not kept}, weakest at \emph{disputed}:
\begin{itemize}
\item In Matthew 19:17, Christ answers on the law’s own terms, so that the young man may see he has not kept it and turn to faith --- \cite[393--394]{calvin-harmony-evangelists-2-1845}
\item an inference step --- \cite[12]{john-paul-ii-veritatis-splendor-1993}
\end{itemize}
\end{itemize}
\noindent Why \emph{Luke 7:50: “your faith has saved you”} defeats \emph{Trent, against ‘not by works’}, at the level of the claims: \texttt{why\allowbreak{}The\allowbreak{}Dominical\allowbreak{}Case\allowbreak{}Stands\allowbreak{}Against\allowbreak{}Trent}.
\par\medskip\noindent\textbf{\texttt{luke18\_\allowbreak{}defensible}}\par\nobreak\noindent
\textbf{Luke 18 can be defended.} A consistent set of positions that answers every attack holds it, with Luke 7:50. Its only attacker is Trent, over Trent’s definition of justification, and it defeats Trent in return. It does not conflict with either reading of Matthew. The set that holds it here holds Rome’s readings of Matthew 7 and 19 as well (\texttt{luke\_\allowbreak{}stands\_\allowbreak{}with\_\allowbreak{}rome\_\allowbreak{}on\_\allowbreak{}matthew}), and the set that holds Calvin’s readings holds it too (\texttt{matthew7\allowbreak{}Rome\allowbreak{}Answered\allowbreak{}By\allowbreak{}Calvin}).

{\small\begin{verbatim}
theorem luke18_defensible : Framework.CredulouslyAccepted
    jesusWordsDispute.defeats WordsParty.luke18
-- axioms: propext, Classical.choice, Quot.sound
\end{verbatim}}
\par\medskip\noindent\textbf{\texttt{luke18\allowbreak{}Answered\allowbreak{}By\allowbreak{}Trent}}\par\nobreak\noindent
Why Luke 18 is not forced: a set of positions that can be defended includes Trent, and Trent defeats it.

\paragraph{\emph{Luke 18:14: “this man went down to his house justified”} is not accepted on every resolution.}
\begin{itemize}
\item A position holding \emph{Trent, against ‘not by works’}, \emph{Matthew 7:21–23, read as obedience a condition of entry}, \emph{Matthew 19:17, read as the commandments the way to life} can be held, and it holds \emph{Trent, against ‘not by works’}; \emph{Trent, against ‘not by works’} conflicts with \emph{Luke 18:14: “this man went down to his house justified”}: \emph{Trent, against ‘not by works’} defeats \emph{Luke 18:14: “this man went down to his house justified”}, and \emph{Luke 18:14: “this man went down to his house justified”} defeats \emph{Trent, against ‘not by works’} (\texttt{why\allowbreak{}Luke\allowbreak{}Eighteen\allowbreak{}Stands\allowbreak{}Against\allowbreak{}Trent}).
\item None of \emph{Trent, against ‘not by works’}, \emph{Matthew 7:21–23, read as obedience a condition of entry}, \emph{Matthew 19:17, read as the commandments the way to life} defeats another, and each attack on them is answered from among them:
\begin{itemize}
\item \emph{Luke 7:50: “your faith has saved you”} defeats \emph{Trent, against ‘not by works’} (\texttt{why\allowbreak{}The\allowbreak{}Dominical\allowbreak{}Case\allowbreak{}Stands\allowbreak{}Against\allowbreak{}Trent}), and \emph{Trent, against ‘not by works’} defeats \emph{Luke 7:50: “your faith has saved you”}.
\item \emph{Luke 18:14: “this man went down to his house justified”} defeats \emph{Trent, against ‘not by works’} (\texttt{why\allowbreak{}Luke\allowbreak{}Eighteen\allowbreak{}Stands\allowbreak{}Against\allowbreak{}Trent}), and \emph{Trent, against ‘not by works’} defeats \emph{Luke 18:14: “this man went down to his house justified”}.
\item \emph{Matthew 7:21–23, read as the doing including faith} defeats \emph{Matthew 7:21–23, read as obedience a condition of entry}, and \emph{Matthew 7:21–23, read as obedience a condition of entry} defeats \emph{Matthew 7:21–23, read as the doing including faith}.
\item \emph{Matthew 19:17, read as the law exposing what he has not kept} defeats \emph{Matthew 19:17, read as the commandments the way to life}, and \emph{Matthew 19:17, read as the commandments the way to life} defeats \emph{Matthew 19:17, read as the law exposing what he has not kept}.
\end{itemize}
\end{itemize}
\par\noindent\textbf{What this rests on.} The reasons state 8 defeats and 17 absences of defeat, each a cell of the defeat table, \texttt{jesus\allowbreak{}Words\allowbreak{}Dispute\_\allowbreak{}defeats}, computed from the two parties' premises and checked by the kernel. Whether an attack survives turns on the attacker's weakest link and on the rating of what it attacks, so a changed rating can change the verdict. The weakest links:
\begin{itemize}
\item \emph{Trent, against ‘not by works’}, weakest at \emph{disputed}:
\begin{itemize}
\item Justification is not remission of sins only, but renewal of the inward man --- \cite[Trent, Session VI (1547), Decree on Justification, ch. 7]{tanner-decrees-1990}; \cite[Trent, Session VI (1547), canon 11]{tanner-decrees-1990}
\end{itemize}
\item \emph{Luke 18:14: “this man went down to his house justified”}, weakest at \emph{disputed}:
\begin{itemize}
\item an inference step --- \cite[45]{waterworth-trent-1848}; \cite[33--34]{waterworth-trent-1848}
\end{itemize}
\item \emph{Matthew 7:21–23, read as obedience a condition of entry}, weakest at \emph{disputed}:
\begin{itemize}
\item In Matthew 7:21–23, doing the Father’s will is a condition of entering the kingdom in its own right, not only the fruit of faith --- \cite[2826]{catholic-church-catechism-2000}; \cite[1821]{catholic-church-catechism-2000}
\item an inference step --- \cite[367--368]{calvin-harmony-evangelists-1845}
\end{itemize}
\item \emph{Matthew 19:17, read as the commandments the way to life}, weakest at \emph{disputed}:
\begin{itemize}
\item In Matthew 19:17, keeping the commandments is the way to life, made possible by grace --- \cite[12]{john-paul-ii-veritatis-splendor-1993}; \cite[2052–2053]{catholic-church-catechism-2000}; \cite[35]{waterworth-trent-1848}
\item an inference step --- \cite[394]{calvin-harmony-evangelists-2-1845}
\end{itemize}
\item \emph{Luke 7:50: “your faith has saved you”}, weakest at \emph{disputed}:
\begin{itemize}
\item Luke 7:47 — her love is the evidence of her forgiveness, not its ground --- \cite[ad loc. Luke 7:47]{marshall-luke-1978}; \cite[54--54]{padilla-narrative-criticism-2021}; \cite[Apology IV, on love and the fulfilling of the law: Luke 7:47–50]{kolb-wengert-book-of-concord-2000}; Luke 7:41-43; Luke 7:47; Luke 7:47-50
\item an inference step --- \cite[Trent, Session VI (1547), canon 9]{tanner-decrees-1990}
\end{itemize}
\item \emph{Matthew 7:21–23, read as the doing including faith}, weakest at \emph{disputed}:
\begin{itemize}
\item In Matthew 7:21–23, doing the Father’s will includes believing in Christ, and “I never knew you” means no union with him by faith --- \cite[367--368]{calvin-harmony-evangelists-1845}
\item an inference step --- \cite[2826]{catholic-church-catechism-2000}
\end{itemize}
\item \emph{Matthew 19:17, read as the law exposing what he has not kept}, weakest at \emph{disputed}:
\begin{itemize}
\item In Matthew 19:17, Christ answers on the law’s own terms, so that the young man may see he has not kept it and turn to faith --- \cite[393--394]{calvin-harmony-evangelists-2-1845}
\item an inference step --- \cite[12]{john-paul-ii-veritatis-splendor-1993}
\end{itemize}
\end{itemize}
\noindent Why \emph{Luke 18:14: “this man went down to his house justified”} defeats \emph{Trent, against ‘not by works’}, at the level of the claims: \texttt{why\allowbreak{}Luke\allowbreak{}Eighteen\allowbreak{}Stands\allowbreak{}Against\allowbreak{}Trent}.
\noindent Why \emph{Luke 7:50: “your faith has saved you”} defeats \emph{Trent, against ‘not by works’}, at the level of the claims: \texttt{why\allowbreak{}The\allowbreak{}Dominical\allowbreak{}Case\allowbreak{}Stands\allowbreak{}Against\allowbreak{}Trent}.
\par\medskip\noindent\textbf{\texttt{luke18\_\allowbreak{}not\_\allowbreak{}forced}}\par\nobreak\noindent
\textbf{Nor is Luke 18 forced.} A consistent set of positions that answers every attack holds Trent, and cannot hold Luke 18 with it. What would decide it is the step from the plea to “faith suffices”. That step is rated \texttt{disputed}, because Trent requires the dispositions of chapter 6 and condemns “nothing else is required” (canon 9). Behind the step is Trent’s definition of justification (\texttt{why\allowbreak{}Luke\allowbreak{}Eighteen\allowbreak{}Stands\allowbreak{}Against\allowbreak{}Trent}).

{\small\begin{verbatim}
theorem luke18_not_forced : ¬Framework.SkepticallyAccepted
    jesusWordsDispute.defeats WordsParty.luke18
-- axioms: propext, Classical.choice, Quot.sound
\end{verbatim}}
\par\medskip\noindent\textbf{\texttt{matthew7\allowbreak{}Rome\allowbreak{}Answered\allowbreak{}By\allowbreak{}Calvin}}\par\nobreak\noindent
Why Rome’s reading of Matthew 7 is not forced: a set of positions that can be defended includes Calvin’s reading, which defeats Rome’s.

\paragraph{\emph{Matthew 7:21–23, read as obedience a condition of entry} is not accepted on every resolution.}
\begin{itemize}
\item A position holding \emph{Luke 7:50: “your faith has saved you”}, \emph{Luke 18:14: “this man went down to his house justified”}, \emph{Matthew 7:21–23, read as the doing including faith}, \emph{Matthew 19:17, read as the law exposing what he has not kept} can be held, and it holds \emph{Matthew 7:21–23, read as the doing including faith}; \emph{Matthew 7:21–23, read as the doing including faith} conflicts with \emph{Matthew 7:21–23, read as obedience a condition of entry}: \emph{Matthew 7:21–23, read as the doing including faith} defeats \emph{Matthew 7:21–23, read as obedience a condition of entry}, and \emph{Matthew 7:21–23, read as obedience a condition of entry} defeats \emph{Matthew 7:21–23, read as the doing including faith}.
\item None of \emph{Luke 7:50: “your faith has saved you”}, \emph{Luke 18:14: “this man went down to his house justified”}, \emph{Matthew 7:21–23, read as the doing including faith}, \emph{Matthew 19:17, read as the law exposing what he has not kept} defeats another, and each attack on them is answered from among them:
\begin{itemize}
\item \emph{Trent, against ‘not by works’} defeats \emph{Luke 7:50: “your faith has saved you”}, and \emph{Luke 7:50: “your faith has saved you”} defeats \emph{Trent, against ‘not by works’} (\texttt{why\allowbreak{}The\allowbreak{}Dominical\allowbreak{}Case\allowbreak{}Stands\allowbreak{}Against\allowbreak{}Trent}).
\item \emph{Trent, against ‘not by works’} defeats \emph{Luke 18:14: “this man went down to his house justified”}, and \emph{Luke 7:50: “your faith has saved you”} defeats \emph{Trent, against ‘not by works’} (\texttt{why\allowbreak{}The\allowbreak{}Dominical\allowbreak{}Case\allowbreak{}Stands\allowbreak{}Against\allowbreak{}Trent}).
\item \emph{Matthew 7:21–23, read as obedience a condition of entry} defeats \emph{Matthew 7:21–23, read as the doing including faith}, and \emph{Matthew 7:21–23, read as the doing including faith} defeats \emph{Matthew 7:21–23, read as obedience a condition of entry}.
\item \emph{Matthew 19:17, read as the commandments the way to life} defeats \emph{Matthew 19:17, read as the law exposing what he has not kept}, and \emph{Matthew 19:17, read as the law exposing what he has not kept} defeats \emph{Matthew 19:17, read as the commandments the way to life}.
\end{itemize}
\end{itemize}
\par\noindent\textbf{What this rests on.} The reasons state 7 defeats and 24 absences of defeat, each a cell of the defeat table, \texttt{jesus\allowbreak{}Words\allowbreak{}Dispute\_\allowbreak{}defeats}, computed from the two parties' premises and checked by the kernel. Whether an attack survives turns on the attacker's weakest link and on the rating of what it attacks, so a changed rating can change the verdict. The weakest links:
\begin{itemize}
\item \emph{Matthew 7:21–23, read as the doing including faith}, weakest at \emph{disputed}:
\begin{itemize}
\item In Matthew 7:21–23, doing the Father’s will includes believing in Christ, and “I never knew you” means no union with him by faith --- \cite[367--368]{calvin-harmony-evangelists-1845}
\item an inference step --- \cite[2826]{catholic-church-catechism-2000}
\end{itemize}
\item \emph{Matthew 7:21–23, read as obedience a condition of entry}, weakest at \emph{disputed}:
\begin{itemize}
\item In Matthew 7:21–23, doing the Father’s will is a condition of entering the kingdom in its own right, not only the fruit of faith --- \cite[2826]{catholic-church-catechism-2000}; \cite[1821]{catholic-church-catechism-2000}
\item an inference step --- \cite[367--368]{calvin-harmony-evangelists-1845}
\end{itemize}
\item \emph{Luke 7:50: “your faith has saved you”}, weakest at \emph{disputed}:
\begin{itemize}
\item Luke 7:47 — her love is the evidence of her forgiveness, not its ground --- \cite[ad loc. Luke 7:47]{marshall-luke-1978}; \cite[54--54]{padilla-narrative-criticism-2021}; \cite[Apology IV, on love and the fulfilling of the law: Luke 7:47–50]{kolb-wengert-book-of-concord-2000}; Luke 7:41-43; Luke 7:47; Luke 7:47-50
\item an inference step --- \cite[Trent, Session VI (1547), canon 9]{tanner-decrees-1990}
\end{itemize}
\item \emph{Luke 18:14: “this man went down to his house justified”}, weakest at \emph{disputed}:
\begin{itemize}
\item an inference step --- \cite[45]{waterworth-trent-1848}; \cite[33--34]{waterworth-trent-1848}
\end{itemize}
\item \emph{Matthew 19:17, read as the law exposing what he has not kept}, weakest at \emph{disputed}:
\begin{itemize}
\item In Matthew 19:17, Christ answers on the law’s own terms, so that the young man may see he has not kept it and turn to faith --- \cite[393--394]{calvin-harmony-evangelists-2-1845}
\item an inference step --- \cite[12]{john-paul-ii-veritatis-splendor-1993}
\end{itemize}
\item \emph{Trent, against ‘not by works’}, weakest at \emph{disputed}:
\begin{itemize}
\item Justification is not remission of sins only, but renewal of the inward man --- \cite[Trent, Session VI (1547), Decree on Justification, ch. 7]{tanner-decrees-1990}; \cite[Trent, Session VI (1547), canon 11]{tanner-decrees-1990}
\end{itemize}
\item \emph{Matthew 19:17, read as the commandments the way to life}, weakest at \emph{disputed}:
\begin{itemize}
\item In Matthew 19:17, keeping the commandments is the way to life, made possible by grace --- \cite[12]{john-paul-ii-veritatis-splendor-1993}; \cite[2052–2053]{catholic-church-catechism-2000}; \cite[35]{waterworth-trent-1848}
\item an inference step --- \cite[394]{calvin-harmony-evangelists-2-1845}
\end{itemize}
\end{itemize}
\noindent Why \emph{Luke 7:50: “your faith has saved you”} defeats \emph{Trent, against ‘not by works’}, at the level of the claims: \texttt{why\allowbreak{}The\allowbreak{}Dominical\allowbreak{}Case\allowbreak{}Stands\allowbreak{}Against\allowbreak{}Trent}.
\par\medskip\noindent\textbf{\texttt{matthew7\_\allowbreak{}rome\_\allowbreak{}reading\_\allowbreak{}not\_\allowbreak{}forced}}\par\nobreak\noindent
\textbf{Rome’s reading of Matthew 7 is not forced.} A consistent set of positions that answers every attack holds Calvin’s reading, on which the doing includes believing and is not a condition of entry in its own right, and cannot hold Rome’s with it.

{\small\begin{verbatim}
theorem matthew7_rome_reading_not_forced : ¬Framework.SkepticallyAccepted
    jesusWordsDispute.defeats WordsParty.matthew7Obedience
-- axioms: propext, Classical.choice, Quot.sound
\end{verbatim}}
\par\medskip\noindent\textbf{\texttt{matthew7\allowbreak{}Calvin\allowbreak{}Answered\allowbreak{}By\allowbreak{}Rome}}\par\nobreak\noindent
Why Calvin’s reading of Matthew 7 is not forced: a set of positions that can be defended includes Rome’s reading, which defeats Calvin’s.

\paragraph{\emph{Matthew 7:21–23, read as the doing including faith} is not accepted on every resolution.}
\begin{itemize}
\item A position holding \emph{Luke 7:50: “your faith has saved you”}, \emph{Luke 18:14: “this man went down to his house justified”}, \emph{Matthew 7:21–23, read as obedience a condition of entry}, \emph{Matthew 19:17, read as the commandments the way to life} can be held, and it holds \emph{Matthew 7:21–23, read as obedience a condition of entry}; \emph{Matthew 7:21–23, read as obedience a condition of entry} conflicts with \emph{Matthew 7:21–23, read as the doing including faith}: \emph{Matthew 7:21–23, read as obedience a condition of entry} defeats \emph{Matthew 7:21–23, read as the doing including faith}, and \emph{Matthew 7:21–23, read as the doing including faith} defeats \emph{Matthew 7:21–23, read as obedience a condition of entry}.
\item None of \emph{Luke 7:50: “your faith has saved you”}, \emph{Luke 18:14: “this man went down to his house justified”}, \emph{Matthew 7:21–23, read as obedience a condition of entry}, \emph{Matthew 19:17, read as the commandments the way to life} defeats another, and each attack on them is answered from among them:
\begin{itemize}
\item \emph{Trent, against ‘not by works’} defeats \emph{Luke 7:50: “your faith has saved you”}, and \emph{Luke 7:50: “your faith has saved you”} defeats \emph{Trent, against ‘not by works’} (\texttt{why\allowbreak{}The\allowbreak{}Dominical\allowbreak{}Case\allowbreak{}Stands\allowbreak{}Against\allowbreak{}Trent}).
\item \emph{Trent, against ‘not by works’} defeats \emph{Luke 18:14: “this man went down to his house justified”}, and \emph{Luke 7:50: “your faith has saved you”} defeats \emph{Trent, against ‘not by works’} (\texttt{why\allowbreak{}The\allowbreak{}Dominical\allowbreak{}Case\allowbreak{}Stands\allowbreak{}Against\allowbreak{}Trent}).
\item \emph{Matthew 7:21–23, read as the doing including faith} defeats \emph{Matthew 7:21–23, read as obedience a condition of entry}, and \emph{Matthew 7:21–23, read as obedience a condition of entry} defeats \emph{Matthew 7:21–23, read as the doing including faith}.
\item \emph{Matthew 19:17, read as the law exposing what he has not kept} defeats \emph{Matthew 19:17, read as the commandments the way to life}, and \emph{Matthew 19:17, read as the commandments the way to life} defeats \emph{Matthew 19:17, read as the law exposing what he has not kept}.
\end{itemize}
\end{itemize}
\par\noindent\textbf{What this rests on.} The reasons state 7 defeats and 24 absences of defeat, each a cell of the defeat table, \texttt{jesus\allowbreak{}Words\allowbreak{}Dispute\_\allowbreak{}defeats}, computed from the two parties' premises and checked by the kernel. Whether an attack survives turns on the attacker's weakest link and on the rating of what it attacks, so a changed rating can change the verdict. The weakest links:
\begin{itemize}
\item \emph{Matthew 7:21–23, read as obedience a condition of entry}, weakest at \emph{disputed}:
\begin{itemize}
\item In Matthew 7:21–23, doing the Father’s will is a condition of entering the kingdom in its own right, not only the fruit of faith --- \cite[2826]{catholic-church-catechism-2000}; \cite[1821]{catholic-church-catechism-2000}
\item an inference step --- \cite[367--368]{calvin-harmony-evangelists-1845}
\end{itemize}
\item \emph{Matthew 7:21–23, read as the doing including faith}, weakest at \emph{disputed}:
\begin{itemize}
\item In Matthew 7:21–23, doing the Father’s will includes believing in Christ, and “I never knew you” means no union with him by faith --- \cite[367--368]{calvin-harmony-evangelists-1845}
\item an inference step --- \cite[2826]{catholic-church-catechism-2000}
\end{itemize}
\item \emph{Luke 7:50: “your faith has saved you”}, weakest at \emph{disputed}:
\begin{itemize}
\item Luke 7:47 — her love is the evidence of her forgiveness, not its ground --- \cite[ad loc. Luke 7:47]{marshall-luke-1978}; \cite[54--54]{padilla-narrative-criticism-2021}; \cite[Apology IV, on love and the fulfilling of the law: Luke 7:47–50]{kolb-wengert-book-of-concord-2000}; Luke 7:41-43; Luke 7:47; Luke 7:47-50
\item an inference step --- \cite[Trent, Session VI (1547), canon 9]{tanner-decrees-1990}
\end{itemize}
\item \emph{Luke 18:14: “this man went down to his house justified”}, weakest at \emph{disputed}:
\begin{itemize}
\item an inference step --- \cite[45]{waterworth-trent-1848}; \cite[33--34]{waterworth-trent-1848}
\end{itemize}
\item \emph{Matthew 19:17, read as the commandments the way to life}, weakest at \emph{disputed}:
\begin{itemize}
\item In Matthew 19:17, keeping the commandments is the way to life, made possible by grace --- \cite[12]{john-paul-ii-veritatis-splendor-1993}; \cite[2052–2053]{catholic-church-catechism-2000}; \cite[35]{waterworth-trent-1848}
\item an inference step --- \cite[394]{calvin-harmony-evangelists-2-1845}
\end{itemize}
\item \emph{Trent, against ‘not by works’}, weakest at \emph{disputed}:
\begin{itemize}
\item Justification is not remission of sins only, but renewal of the inward man --- \cite[Trent, Session VI (1547), Decree on Justification, ch. 7]{tanner-decrees-1990}; \cite[Trent, Session VI (1547), canon 11]{tanner-decrees-1990}
\end{itemize}
\item \emph{Matthew 19:17, read as the law exposing what he has not kept}, weakest at \emph{disputed}:
\begin{itemize}
\item In Matthew 19:17, Christ answers on the law’s own terms, so that the young man may see he has not kept it and turn to faith --- \cite[393--394]{calvin-harmony-evangelists-2-1845}
\item an inference step --- \cite[12]{john-paul-ii-veritatis-splendor-1993}
\end{itemize}
\end{itemize}
\noindent Why \emph{Luke 7:50: “your faith has saved you”} defeats \emph{Trent, against ‘not by works’}, at the level of the claims: \texttt{why\allowbreak{}The\allowbreak{}Dominical\allowbreak{}Case\allowbreak{}Stands\allowbreak{}Against\allowbreak{}Trent}.
\par\medskip\noindent\textbf{\texttt{matthew7\_\allowbreak{}calvin\_\allowbreak{}reading\_\allowbreak{}not\_\allowbreak{}forced}}\par\nobreak\noindent
\textbf{Nor is Calvin’s.} A consistent set of positions that answers every attack holds Rome’s reading, on which obedience is a condition of entry, and cannot hold Calvin’s with it. Between the two readings of Matthew 7 the debate chooses neither, as the passage itself does not (\texttt{matthew7\_\allowbreak{}leaves\_\allowbreak{}the\_\allowbreak{}ground\_\allowbreak{}open}).

{\small\begin{verbatim}
theorem matthew7_calvin_reading_not_forced : ¬Framework.SkepticallyAccepted
    jesusWordsDispute.defeats WordsParty.matthew7Faith
-- axioms: propext, Classical.choice, Quot.sound
\end{verbatim}}
\par\medskip\noindent\textbf{\texttt{matthew19\allowbreak{}Rome\allowbreak{}Answered\allowbreak{}By\allowbreak{}Calvin}}\par\nobreak\noindent
Why Rome’s reading of Matthew 19 is not forced: a set of positions that can be defended includes Calvin’s reading, which defeats Rome’s.

\paragraph{\emph{Matthew 19:17, read as the commandments the way to life} is not accepted on every resolution.}
\begin{itemize}
\item A position holding \emph{Luke 7:50: “your faith has saved you”}, \emph{Luke 18:14: “this man went down to his house justified”}, \emph{Matthew 7:21–23, read as the doing including faith}, \emph{Matthew 19:17, read as the law exposing what he has not kept} can be held, and it holds \emph{Matthew 19:17, read as the law exposing what he has not kept}; \emph{Matthew 19:17, read as the law exposing what he has not kept} conflicts with \emph{Matthew 19:17, read as the commandments the way to life}: \emph{Matthew 19:17, read as the law exposing what he has not kept} defeats \emph{Matthew 19:17, read as the commandments the way to life}, and \emph{Matthew 19:17, read as the commandments the way to life} defeats \emph{Matthew 19:17, read as the law exposing what he has not kept}.
\item None of \emph{Luke 7:50: “your faith has saved you”}, \emph{Luke 18:14: “this man went down to his house justified”}, \emph{Matthew 7:21–23, read as the doing including faith}, \emph{Matthew 19:17, read as the law exposing what he has not kept} defeats another, and each attack on them is answered from among them:
\begin{itemize}
\item \emph{Trent, against ‘not by works’} defeats \emph{Luke 7:50: “your faith has saved you”}, and \emph{Luke 7:50: “your faith has saved you”} defeats \emph{Trent, against ‘not by works’} (\texttt{why\allowbreak{}The\allowbreak{}Dominical\allowbreak{}Case\allowbreak{}Stands\allowbreak{}Against\allowbreak{}Trent}).
\item \emph{Trent, against ‘not by works’} defeats \emph{Luke 18:14: “this man went down to his house justified”}, and \emph{Luke 7:50: “your faith has saved you”} defeats \emph{Trent, against ‘not by works’} (\texttt{why\allowbreak{}The\allowbreak{}Dominical\allowbreak{}Case\allowbreak{}Stands\allowbreak{}Against\allowbreak{}Trent}).
\item \emph{Matthew 7:21–23, read as obedience a condition of entry} defeats \emph{Matthew 7:21–23, read as the doing including faith}, and \emph{Matthew 7:21–23, read as the doing including faith} defeats \emph{Matthew 7:21–23, read as obedience a condition of entry}.
\item \emph{Matthew 19:17, read as the commandments the way to life} defeats \emph{Matthew 19:17, read as the law exposing what he has not kept}, and \emph{Matthew 19:17, read as the law exposing what he has not kept} defeats \emph{Matthew 19:17, read as the commandments the way to life}.
\end{itemize}
\end{itemize}
\par\noindent\textbf{What this rests on.} The reasons state 7 defeats and 24 absences of defeat, each a cell of the defeat table, \texttt{jesus\allowbreak{}Words\allowbreak{}Dispute\_\allowbreak{}defeats}, computed from the two parties' premises and checked by the kernel. Whether an attack survives turns on the attacker's weakest link and on the rating of what it attacks, so a changed rating can change the verdict. The weakest links:
\begin{itemize}
\item \emph{Matthew 19:17, read as the law exposing what he has not kept}, weakest at \emph{disputed}:
\begin{itemize}
\item In Matthew 19:17, Christ answers on the law’s own terms, so that the young man may see he has not kept it and turn to faith --- \cite[393--394]{calvin-harmony-evangelists-2-1845}
\item an inference step --- \cite[12]{john-paul-ii-veritatis-splendor-1993}
\end{itemize}
\item \emph{Matthew 19:17, read as the commandments the way to life}, weakest at \emph{disputed}:
\begin{itemize}
\item In Matthew 19:17, keeping the commandments is the way to life, made possible by grace --- \cite[12]{john-paul-ii-veritatis-splendor-1993}; \cite[2052–2053]{catholic-church-catechism-2000}; \cite[35]{waterworth-trent-1848}
\item an inference step --- \cite[394]{calvin-harmony-evangelists-2-1845}
\end{itemize}
\item \emph{Luke 7:50: “your faith has saved you”}, weakest at \emph{disputed}:
\begin{itemize}
\item Luke 7:47 — her love is the evidence of her forgiveness, not its ground --- \cite[ad loc. Luke 7:47]{marshall-luke-1978}; \cite[54--54]{padilla-narrative-criticism-2021}; \cite[Apology IV, on love and the fulfilling of the law: Luke 7:47–50]{kolb-wengert-book-of-concord-2000}; Luke 7:41-43; Luke 7:47; Luke 7:47-50
\item an inference step --- \cite[Trent, Session VI (1547), canon 9]{tanner-decrees-1990}
\end{itemize}
\item \emph{Luke 18:14: “this man went down to his house justified”}, weakest at \emph{disputed}:
\begin{itemize}
\item an inference step --- \cite[45]{waterworth-trent-1848}; \cite[33--34]{waterworth-trent-1848}
\end{itemize}
\item \emph{Matthew 7:21–23, read as the doing including faith}, weakest at \emph{disputed}:
\begin{itemize}
\item In Matthew 7:21–23, doing the Father’s will includes believing in Christ, and “I never knew you” means no union with him by faith --- \cite[367--368]{calvin-harmony-evangelists-1845}
\item an inference step --- \cite[2826]{catholic-church-catechism-2000}
\end{itemize}
\item \emph{Trent, against ‘not by works’}, weakest at \emph{disputed}:
\begin{itemize}
\item Justification is not remission of sins only, but renewal of the inward man --- \cite[Trent, Session VI (1547), Decree on Justification, ch. 7]{tanner-decrees-1990}; \cite[Trent, Session VI (1547), canon 11]{tanner-decrees-1990}
\end{itemize}
\item \emph{Matthew 7:21–23, read as obedience a condition of entry}, weakest at \emph{disputed}:
\begin{itemize}
\item In Matthew 7:21–23, doing the Father’s will is a condition of entering the kingdom in its own right, not only the fruit of faith --- \cite[2826]{catholic-church-catechism-2000}; \cite[1821]{catholic-church-catechism-2000}
\item an inference step --- \cite[367--368]{calvin-harmony-evangelists-1845}
\end{itemize}
\end{itemize}
\noindent Why \emph{Luke 7:50: “your faith has saved you”} defeats \emph{Trent, against ‘not by works’}, at the level of the claims: \texttt{why\allowbreak{}The\allowbreak{}Dominical\allowbreak{}Case\allowbreak{}Stands\allowbreak{}Against\allowbreak{}Trent}.
\par\medskip\noindent\textbf{\texttt{matthew19\_\allowbreak{}rome\_\allowbreak{}reading\_\allowbreak{}not\_\allowbreak{}forced}}\par\nobreak\noindent
\textbf{Rome’s reading of Matthew 19 is not forced.} A consistent set of positions that answers every attack holds Calvin’s reading, a reply on the law’s own terms that shows the young man what he has not kept, and cannot hold Rome’s with it.

{\small\begin{verbatim}
theorem matthew19_rome_reading_not_forced : ¬Framework.SkepticallyAccepted
    jesusWordsDispute.defeats WordsParty.matthew19Way
-- axioms: propext, Classical.choice, Quot.sound
\end{verbatim}}
\par\medskip\noindent\textbf{\texttt{matthew19\allowbreak{}Calvin\allowbreak{}Answered\allowbreak{}By\allowbreak{}Rome}}\par\nobreak\noindent
Why Calvin’s reading of Matthew 19 is not forced: a set of positions that can be defended includes Rome’s reading, which defeats Calvin’s.

\paragraph{\emph{Matthew 19:17, read as the law exposing what he has not kept} is not accepted on every resolution.}
\begin{itemize}
\item A position holding \emph{Luke 7:50: “your faith has saved you”}, \emph{Luke 18:14: “this man went down to his house justified”}, \emph{Matthew 7:21–23, read as obedience a condition of entry}, \emph{Matthew 19:17, read as the commandments the way to life} can be held, and it holds \emph{Matthew 19:17, read as the commandments the way to life}; \emph{Matthew 19:17, read as the commandments the way to life} conflicts with \emph{Matthew 19:17, read as the law exposing what he has not kept}: \emph{Matthew 19:17, read as the commandments the way to life} defeats \emph{Matthew 19:17, read as the law exposing what he has not kept}, and \emph{Matthew 19:17, read as the law exposing what he has not kept} defeats \emph{Matthew 19:17, read as the commandments the way to life}.
\item None of \emph{Luke 7:50: “your faith has saved you”}, \emph{Luke 18:14: “this man went down to his house justified”}, \emph{Matthew 7:21–23, read as obedience a condition of entry}, \emph{Matthew 19:17, read as the commandments the way to life} defeats another, and each attack on them is answered from among them:
\begin{itemize}
\item \emph{Trent, against ‘not by works’} defeats \emph{Luke 7:50: “your faith has saved you”}, and \emph{Luke 7:50: “your faith has saved you”} defeats \emph{Trent, against ‘not by works’} (\texttt{why\allowbreak{}The\allowbreak{}Dominical\allowbreak{}Case\allowbreak{}Stands\allowbreak{}Against\allowbreak{}Trent}).
\item \emph{Trent, against ‘not by works’} defeats \emph{Luke 18:14: “this man went down to his house justified”}, and \emph{Luke 7:50: “your faith has saved you”} defeats \emph{Trent, against ‘not by works’} (\texttt{why\allowbreak{}The\allowbreak{}Dominical\allowbreak{}Case\allowbreak{}Stands\allowbreak{}Against\allowbreak{}Trent}).
\item \emph{Matthew 7:21–23, read as the doing including faith} defeats \emph{Matthew 7:21–23, read as obedience a condition of entry}, and \emph{Matthew 7:21–23, read as obedience a condition of entry} defeats \emph{Matthew 7:21–23, read as the doing including faith}.
\item \emph{Matthew 19:17, read as the law exposing what he has not kept} defeats \emph{Matthew 19:17, read as the commandments the way to life}, and \emph{Matthew 19:17, read as the commandments the way to life} defeats \emph{Matthew 19:17, read as the law exposing what he has not kept}.
\end{itemize}
\end{itemize}
\par\noindent\textbf{What this rests on.} The reasons state 7 defeats and 24 absences of defeat, each a cell of the defeat table, \texttt{jesus\allowbreak{}Words\allowbreak{}Dispute\_\allowbreak{}defeats}, computed from the two parties' premises and checked by the kernel. Whether an attack survives turns on the attacker's weakest link and on the rating of what it attacks, so a changed rating can change the verdict. The weakest links:
\begin{itemize}
\item \emph{Matthew 19:17, read as the commandments the way to life}, weakest at \emph{disputed}:
\begin{itemize}
\item In Matthew 19:17, keeping the commandments is the way to life, made possible by grace --- \cite[12]{john-paul-ii-veritatis-splendor-1993}; \cite[2052–2053]{catholic-church-catechism-2000}; \cite[35]{waterworth-trent-1848}
\item an inference step --- \cite[394]{calvin-harmony-evangelists-2-1845}
\end{itemize}
\item \emph{Matthew 19:17, read as the law exposing what he has not kept}, weakest at \emph{disputed}:
\begin{itemize}
\item In Matthew 19:17, Christ answers on the law’s own terms, so that the young man may see he has not kept it and turn to faith --- \cite[393--394]{calvin-harmony-evangelists-2-1845}
\item an inference step --- \cite[12]{john-paul-ii-veritatis-splendor-1993}
\end{itemize}
\item \emph{Luke 7:50: “your faith has saved you”}, weakest at \emph{disputed}:
\begin{itemize}
\item Luke 7:47 — her love is the evidence of her forgiveness, not its ground --- \cite[ad loc. Luke 7:47]{marshall-luke-1978}; \cite[54--54]{padilla-narrative-criticism-2021}; \cite[Apology IV, on love and the fulfilling of the law: Luke 7:47–50]{kolb-wengert-book-of-concord-2000}; Luke 7:41-43; Luke 7:47; Luke 7:47-50
\item an inference step --- \cite[Trent, Session VI (1547), canon 9]{tanner-decrees-1990}
\end{itemize}
\item \emph{Luke 18:14: “this man went down to his house justified”}, weakest at \emph{disputed}:
\begin{itemize}
\item an inference step --- \cite[45]{waterworth-trent-1848}; \cite[33--34]{waterworth-trent-1848}
\end{itemize}
\item \emph{Matthew 7:21–23, read as obedience a condition of entry}, weakest at \emph{disputed}:
\begin{itemize}
\item In Matthew 7:21–23, doing the Father’s will is a condition of entering the kingdom in its own right, not only the fruit of faith --- \cite[2826]{catholic-church-catechism-2000}; \cite[1821]{catholic-church-catechism-2000}
\item an inference step --- \cite[367--368]{calvin-harmony-evangelists-1845}
\end{itemize}
\item \emph{Trent, against ‘not by works’}, weakest at \emph{disputed}:
\begin{itemize}
\item Justification is not remission of sins only, but renewal of the inward man --- \cite[Trent, Session VI (1547), Decree on Justification, ch. 7]{tanner-decrees-1990}; \cite[Trent, Session VI (1547), canon 11]{tanner-decrees-1990}
\end{itemize}
\item \emph{Matthew 7:21–23, read as the doing including faith}, weakest at \emph{disputed}:
\begin{itemize}
\item In Matthew 7:21–23, doing the Father’s will includes believing in Christ, and “I never knew you” means no union with him by faith --- \cite[367--368]{calvin-harmony-evangelists-1845}
\item an inference step --- \cite[2826]{catholic-church-catechism-2000}
\end{itemize}
\end{itemize}
\noindent Why \emph{Luke 7:50: “your faith has saved you”} defeats \emph{Trent, against ‘not by works’}, at the level of the claims: \texttt{why\allowbreak{}The\allowbreak{}Dominical\allowbreak{}Case\allowbreak{}Stands\allowbreak{}Against\allowbreak{}Trent}.
\par\medskip\noindent\textbf{\texttt{matthew19\_\allowbreak{}calvin\_\allowbreak{}reading\_\allowbreak{}not\_\allowbreak{}forced}}\par\nobreak\noindent
\textbf{Nor is Calvin’s.} A consistent set of positions that answers every attack holds Rome’s reading, on which the commandments are the way to life, and cannot hold Calvin’s with it. Between the two readings of Matthew 19 the debate chooses neither (\texttt{matthew19\_\allowbreak{}leaves\_\allowbreak{}the\_\allowbreak{}way\_\allowbreak{}open}).

{\small\begin{verbatim}
theorem matthew19_calvin_reading_not_forced : ¬Framework.SkepticallyAccepted
    jesusWordsDispute.defeats WordsParty.matthew19Legal
-- axioms: propext, Classical.choice, Quot.sound
\end{verbatim}}
\paragraph*{Heard without Trent}
\leavevmode\par\nobreak
\par\medskip\noindent\textbf{\texttt{jesus\allowbreak{}Words\allowbreak{}Without\allowbreak{}Trent}}\par\nobreak\noindent
Every position but Trent: Luke’s two sayings, and both readings of each passage in Matthew.

{\small\begin{verbatim}
def jesusWordsWithoutTrent : Dispute Claim { i // i ∉ [WordsParty.trent] } :=
  jesusWordsDispute.restrict fun x => x ∉ [WordsParty.trent]
\end{verbatim}}
\par\medskip\noindent\textbf{\texttt{without\allowbreak{}Trent\allowbreak{}Luke\allowbreak{}Settles}}\par\nobreak\noindent
How the hearing is settled. With Trent not heard, nothing attacks Luke 7:50 or Luke 18, so both are accepted first. Nothing joins them, because each reading of Matthew still has its rival.

\paragraph{What the dispute forces is exactly \emph{Luke 7:50: “your faith has saved you”}, \emph{Luke 18:14: “this man went down to his house justified”}.}
\begin{itemize}
\item Stage 1: \emph{Luke 7:50: “your faith has saved you”}, \emph{Luke 18:14: “this man went down to his house justified”}.
\begin{itemize}
\item \emph{Luke 7:50: “your faith has saved you”} is defeated by nothing.
\item \emph{Luke 18:14: “this man went down to his house justified”} is defeated by nothing.
\end{itemize}
\item Nothing else is forced:
\begin{itemize}
\item \emph{Matthew 7:21–23, read as the doing including faith} defeats \emph{Matthew 7:21–23, read as obedience a condition of entry}, and nothing forced defeats \emph{Matthew 7:21–23, read as the doing including faith}.
\item \emph{Matthew 7:21–23, read as obedience a condition of entry} defeats \emph{Matthew 7:21–23, read as the doing including faith}, and nothing forced defeats \emph{Matthew 7:21–23, read as obedience a condition of entry}.
\item \emph{Matthew 19:17, read as the law exposing what he has not kept} defeats \emph{Matthew 19:17, read as the commandments the way to life}, and nothing forced defeats \emph{Matthew 19:17, read as the law exposing what he has not kept}.
\item \emph{Matthew 19:17, read as the commandments the way to life} defeats \emph{Matthew 19:17, read as the law exposing what he has not kept}, and nothing forced defeats \emph{Matthew 19:17, read as the commandments the way to life}.
\end{itemize}
\end{itemize}
\par\noindent\textbf{What this rests on.} The reasons state 4 defeats and 20 absences of defeat, each a cell of the defeat table, \texttt{jesus\allowbreak{}Words\allowbreak{}Dispute\_\allowbreak{}defeats}, computed from the two parties' premises and checked by the kernel. Whether an attack survives turns on the attacker's weakest link and on the rating of what it attacks, so a changed rating can change the verdict. The weakest links:
\begin{itemize}
\item \emph{Luke 7:50: “your faith has saved you”}, weakest at \emph{disputed}:
\begin{itemize}
\item Luke 7:47 — her love is the evidence of her forgiveness, not its ground --- \cite[ad loc. Luke 7:47]{marshall-luke-1978}; \cite[54--54]{padilla-narrative-criticism-2021}; \cite[Apology IV, on love and the fulfilling of the law: Luke 7:47–50]{kolb-wengert-book-of-concord-2000}; Luke 7:41-43; Luke 7:47; Luke 7:47-50
\item an inference step --- \cite[Trent, Session VI (1547), canon 9]{tanner-decrees-1990}
\end{itemize}
\item \emph{Luke 18:14: “this man went down to his house justified”}, weakest at \emph{disputed}:
\begin{itemize}
\item an inference step --- \cite[45]{waterworth-trent-1848}; \cite[33--34]{waterworth-trent-1848}
\end{itemize}
\item \emph{Matthew 7:21–23, read as obedience a condition of entry}, weakest at \emph{disputed}:
\begin{itemize}
\item In Matthew 7:21–23, doing the Father’s will is a condition of entering the kingdom in its own right, not only the fruit of faith --- \cite[2826]{catholic-church-catechism-2000}; \cite[1821]{catholic-church-catechism-2000}
\item an inference step --- \cite[367--368]{calvin-harmony-evangelists-1845}
\end{itemize}
\item \emph{Matthew 7:21–23, read as the doing including faith}, weakest at \emph{disputed}:
\begin{itemize}
\item In Matthew 7:21–23, doing the Father’s will includes believing in Christ, and “I never knew you” means no union with him by faith --- \cite[367--368]{calvin-harmony-evangelists-1845}
\item an inference step --- \cite[2826]{catholic-church-catechism-2000}
\end{itemize}
\item \emph{Matthew 19:17, read as the commandments the way to life}, weakest at \emph{disputed}:
\begin{itemize}
\item In Matthew 19:17, keeping the commandments is the way to life, made possible by grace --- \cite[12]{john-paul-ii-veritatis-splendor-1993}; \cite[2052–2053]{catholic-church-catechism-2000}; \cite[35]{waterworth-trent-1848}
\item an inference step --- \cite[394]{calvin-harmony-evangelists-2-1845}
\end{itemize}
\item \emph{Matthew 19:17, read as the law exposing what he has not kept}, weakest at \emph{disputed}:
\begin{itemize}
\item In Matthew 19:17, Christ answers on the law’s own terms, so that the young man may see he has not kept it and turn to faith --- \cite[393--394]{calvin-harmony-evangelists-2-1845}
\item an inference step --- \cite[12]{john-paul-ii-veritatis-splendor-1993}
\end{itemize}
\end{itemize}
\par\medskip\noindent\textbf{\texttt{without\_\allowbreak{}trent\_\allowbreak{}luke\_\allowbreak{}is\_\allowbreak{}forced}}\par\nobreak\noindent
\textbf{Without Trent, Luke’s two sayings are forced, and Matthew is not.} With Trent not heard, nothing attacks “your faith has saved you” or the tax collector justified. That faith suffices is forced twice, from Jesus’ own words. The readings of Matthew 7 and 19 still tie, each with its rival, because they conflict with each other and not with Luke.

It does not show that faith \emph{alone} is forced. Luke’s sayings claim that faith sufficed, not that nothing else is a condition. Rome’s readings of Matthew, which say that salvation is not apart from works, can be held with them (\texttt{luke\_\allowbreak{}stands\_\allowbreak{}with\_\allowbreak{}rome\_\allowbreak{}on\_\allowbreak{}matthew}).

{\small\begin{verbatim}
theorem without_trent_luke_is_forced : Framework.grounded
    jesusWordsWithoutTrent.defeats = Solver.toSet (Witness.listSub (fun x => x
    ∉ [WordsParty.trent]) [WordsParty.luke7, WordsParty.luke18])
-- axioms: propext, Classical.choice, Quot.sound
\end{verbatim}}
\subsection*{Is baptism how we are justified?}
\addcontentsline{toc}{subsection}{Is baptism how we are justified?}
When Trent lists the causes of justification, it names baptism among them:

\begin{quote}
“the instrumental cause is the sacrament of baptism, which is the sacrament of faith, without which (faith) no man was ever justified” (Session VI, ch. 7).

\end{quote}
The “without which” refers to faith, not to baptism, as the translator Waterworth’s bracket shows. The Reformed hold baptism to be Christ’s own ordinance, a sign and seal of the covenant of grace (Westminster XXVIII.1), and faith to be “the alone instrument of justification” (XI.2). So the question this chapter takes is a plain one: does the New Testament make water baptism the instrument that justifies?

A reminder of the ratings. A step is \texttt{disputed} when a cited reader accepts its starting points and denies its conclusion. It is \texttt{well\allowbreak{}Supported} when no such reader was found. It is \texttt{consensus} when every side accepts it.

\subsubsection*{Three things called baptism}
“Baptism” names three things in the New Testament. Each is a separate claim here, so that no argument slides from one to another unnoticed.

\begin{itemize}
  \item \textbf{Water.} John’s baptism, and the one Christ commands (Matthew 28:19; Acts 2:38). Trent’s instrumental cause is \emph{this} one. It is not a “commandment of men”, since Jesus commands it, so the argument from Mark 7 against human tradition (in the sola scriptura argument) does not apply to it as such. The disagreement is over the role Trent gives it.
  \item \textbf{The Spirit.} “He will baptize you with the Holy Spirit.” In Acts the Spirit comes after the water (Samaria, 8:14–17) and before it (Cornelius, 10:44–48).
  \item \textbf{“And fire”} (Matthew 3:11; Luke 3:16). Gregory Nazianzen reads the fire as judgement, the chaff burned. Calvin reads it as purification. Cyril of Jerusalem reads it as Pentecost’s tongues “as of fire”. All three readings are written down, and none is chosen: the text does not settle between them (\texttt{the\_\allowbreak{}baptist\_\allowbreak{}leaves\_\allowbreak{}the\_\allowbreak{}fire\_\allowbreak{}open} and the two results after it).
\end{itemize}
\subsubsection*{What Trent’s claim says}
Trent qualifies its own claim. Justification “cannot be effected, without the laver of regeneration, \textbf{or the desire thereof}” (ch. 4). So the claim can be taken two ways, and \texttt{what\allowbreak{}Trents\allowbreak{}Instrumental\allowbreak{}Cause\allowbreak{}Claims} answers both.

\begin{itemize}
  \item \textbf{Taken so that the washing itself is necessary for salvation.} This is Feeney’s reading: desire may justify, but “It is now: Baptism of Water, or damnation!” No reader was found who holds the water itself necessary for \emph{justification}; Feeney too grants that desire justifies. Taken Feeney’s way, the claim cannot be held together with the thief promised paradise unbaptised (Luke 23:43), read as Augustine reads him. Nor can it be held with Rome’s own teaching that the desire “assures them the salvation” (Catechism 1259). Both steps are \texttt{well\allowbreak{}Supported}.
  \item \textbf{Taken as “baptism or the desire for it”}, Trent’s own words. The thief does not count against this. What it meets is a question the desire raises. Where only the desire is present, what justifies is not the washing, since the desire “brings about the fruits of Baptism without being a sacrament” (Catechism 1258). Then what is the instrument, and how does it differ from faith? The Reformed answer is that the instrument is faith, and the sacrament is not. Aquinas accepts the premises and denies that: the desire is a desire \emph{for} the sacrament, and when the sacrament follows it gives “a yet greater fulness of grace”. So the step is \texttt{disputed}.
\end{itemize}
Neither way of taking the claim is ruled out by what is argued here. The first is answered by texts and teaching that Rome itself holds. The second stands at a disputed step, where the question is whether the desire is faith under another name.

\subsubsection*{John 3:5}
\begin{quote}
“Unless one is born of water and the Spirit, he cannot enter the kingdom of God” (John 3:5).

\end{quote}
Trent cites this verse in chapter 4, and Session VII condemns whoever “wrests, to some sort of metaphor” its water. \texttt{what\allowbreak{}John\allowbreak{}Three\allowbreak{}Five\allowbreak{}Means} takes the four readings of “water” that are on record.

\begin{itemize}
  \item Baptismal water: Trent, and most early Christian writers, as Calvin reports.
  \item Natural birth: Oliver.
  \item The cleansing promised in Ezekiel 36:25–27: Carson.
  \item The Spirit’s own cleansing, with “water and Spirit” naming one thing: Calvin.
\end{itemize}
What can be taken more than one way here is the verse, not Trent. Trent argues from the verse to water baptism, and each reading is a reading of the verse’s “water”. Only the baptismal reading leaves Trent’s step standing. The other three each rule it out, at a step rated \texttt{well\allowbreak{}Supported}.

The baptismal reading cannot be held together with Calvin’s. The step from Calvin’s reading is \texttt{well\allowbreak{}Supported}, but his reading itself is \texttt{disputed}, and Trent condemns it. Even on the baptismal reading Trent’s step stays \texttt{disputed}: Calvin accepts the verse and denies that it requires water baptism. Which reading is right is not decided here.

\subsubsection*{Acts gives both orders}
In Samaria the water comes first and the Spirit after. At Caesarea the Spirit comes first and the water after. The Reformed argue from this that the Spirit is not tied to the water; Calvin on Acts 10:47: “the Spirit is not included in baptism”. Rome says as much in its own words: “God has bound salvation to the sacrament of Baptism, but he himself is not bound by his sacraments” (Catechism 1257).

So that God is not bound to the water is accepted on both sides (\texttt{trent\_\allowbreak{}grants\_\allowbreak{}god\_\allowbreak{}is\_\allowbreak{}not\_\allowbreak{}bound}). It is not a point won by the Reformed. Calvin’s further claim, that the Spirit is not included in the sign, is not what Rome accepts. What remains in dispute is the first half of the Catechism’s sentence, that God has bound salvation to the sacrament, and Trent’s further claim that the sacrament is the instrument of justification, which is the question above.

\subsubsection*{How this is checked}
\paragraph*{Trent’s claim, and two ways of taking it}
\leavevmode\par\nobreak
\par\medskip\noindent\textbf{\texttt{trent\allowbreak{}On\allowbreak{}Baptism}} --- Trent: baptism the instrumental cause of justification\par\nobreak\noindent
\textbf{Trent’s instrumental cause}, as Trent states it: water baptism is the instrumental cause of justification.

\begin{derivation}
  \item[(1)] \(\pvar{TestimonyArguments.SolaFide}{68}\)\gloss{Water baptism is the instrumental cause of justification}
  \item[\(\vdash\)] \(\pvar{TestimonyArguments.SolaFide}{68}\)\gloss{\emph{at issue:} Water baptism is the instrumental cause of justification}
  \keepwithnext
\end{derivation}
\emph{No premise here rests on scripture alone.}
\par\medskip\noindent\textbf{\texttt{feeney\allowbreak{}Reads\allowbreak{}Trent}}\par\nobreak\noindent
Taking the instrumental cause so that the washing itself is necessary for salvation. This is Feeney’s reading. Rated \texttt{disputed}: Rome’s own teaching denies it (Catechism 1258–1260).

{\small\begin{verbatim}
def feeneyReadsTrent : Source :=
  { primary := Reference.work feeneyBreadOfLife,
    tradition := Tradition.romanCatholic, confidence := Confidence.disputed }
\end{verbatim}}
\par\medskip\noindent\textbf{\texttt{trent\allowbreak{}Reads\allowbreak{}Its\allowbreak{}Own\allowbreak{}Claim}}\par\nobreak\noindent
Taking the instrumental cause as Trent’s chapter 4 qualifies it: the laver of regeneration, or the desire for it. Rated \texttt{consensus}, since it is what Trent says.

{\small\begin{verbatim}
def trentReadsItsOwnClaim : Source :=
  { primary := Reference.work waterworthTrent (Bib.Locus.page 32),
    tradition := Tradition.romanCatholic, confidence := Confidence.consensus }
\end{verbatim}}
\par\medskip\noindent\textbf{\texttt{as\allowbreak{}The\allowbreak{}Water\allowbreak{}Itself}}\par\nobreak\noindent
The instrumental cause, \textbf{taken so that the washing itself is necessary for salvation}.

{\small\begin{verbatim}
def asTheWaterItself : Reading Claim :=
  { name := "so that the washing itself is necessary for salvation",
    commits := p Claim.waterItselfNecessaryForSalvation,
    source := feeneyReadsTrent }
\end{verbatim}}
\par\medskip\noindent\textbf{\texttt{as\allowbreak{}Baptism\allowbreak{}Or\allowbreak{}Its\allowbreak{}Desire}}\par\nobreak\noindent
The instrumental cause, \textbf{taken as baptism or the desire for it}.

{\small\begin{verbatim}
def asBaptismOrItsDesire : Reading Claim :=
  { name := "as baptism or the desire of it",
    commits := p Claim.baptismOrItsDesire,
    source := trentReadsItsOwnClaim }
\end{verbatim}}
\par\medskip\noindent\textbf{\texttt{trent\allowbreak{}On\allowbreak{}The\allowbreak{}Water\allowbreak{}Itself}} --- Trent: baptism the instrumental cause of justification, read so that the washing itself is necessary for salvation\par\nobreak\noindent
Trent’s instrumental cause, taken so that the washing itself is necessary.

\begin{derivation}
  \item[(1)] \(\pvar{TestimonyArguments.SolaFide}{68}\)\gloss{Water baptism is the instrumental cause of justification}
  \item[(2)] \(\pvar{TestimonyArguments.SolaFide}{68} \rightarrow \pvar{TestimonyArguments.SolaFide}{69}\)\gloss{\emph{would give:} The washing itself is necessary for salvation, though desire may justify}
  \item[\(\vdash\)] \(\pvar{TestimonyArguments.SolaFide}{68}\)\gloss{\emph{at issue:} Water baptism is the instrumental cause of justification}
  \keepwithnext
\end{derivation}
\emph{No premise here rests on scripture alone.}
\par\medskip\noindent\textbf{\texttt{trent\allowbreak{}On\allowbreak{}Baptism\allowbreak{}Or\allowbreak{}Desire}} --- Trent: baptism the instrumental cause of justification, read as baptism or the desire of it\par\nobreak\noindent
Trent’s instrumental cause, taken as baptism or the desire for it.

\begin{derivation}
  \item[(1)] \(\pvar{TestimonyArguments.SolaFide}{68}\)\gloss{Water baptism is the instrumental cause of justification}
  \item[(2)] \(\pvar{TestimonyArguments.SolaFide}{68} \rightarrow \pvar{TestimonyArguments.SolaFide}{70}\)\gloss{\emph{would give:} Justification cannot be effected without the laver of regeneration, or the desire thereof}
  \item[\(\vdash\)] \(\pvar{TestimonyArguments.SolaFide}{68}\)\gloss{\emph{at issue:} Water baptism is the instrumental cause of justification}
  \keepwithnext
\end{derivation}
\emph{No premise here rests on scripture alone.}
\paragraph*{What each is set against}
\leavevmode\par\nobreak
\par\medskip\noindent\textbf{\texttt{thief\allowbreak{}To\allowbreak{}Not\allowbreak{}Necessary}}\par\nobreak\noindent
From the thief to “the washing is not necessary”: he was promised paradise unbaptised, and his case bears on salvation under the Gospel.

\begin{derivation}
  \item[] \((\pvar{TestimonyArguments.SolaFide}{72} \land \pvar{TestimonyArguments.SolaFide}{73}) \rightarrow \lnot \pvar{TestimonyArguments.SolaFide}{69}\)\gloss{\emph{would deny:} The washing itself is necessary for salvation, though desire may justify}
\end{derivation}
\par\medskip\noindent\textbf{\texttt{thief\allowbreak{}Case}} --- The thief on the cross (Luke 23:43), read with Augustine\par\nobreak\noindent
\textbf{The thief on the cross}, as Augustine reads him: the washing itself is not necessary for salvation.

\begin{derivation}
  \item[(1)] \(\pvar{TestimonyArguments.SolaFide}{72}\)\gloss{Luke 23:43 — the thief, unbaptised, is promised paradise}
  \item[(2)] \(\pvar{TestimonyArguments.SolaFide}{73}\)\gloss{The thief’s case bears on salvation under the Gospel}
  \item[(3)] \((\pvar{TestimonyArguments.SolaFide}{72} \land \pvar{TestimonyArguments.SolaFide}{73}) \rightarrow \lnot \pvar{TestimonyArguments.SolaFide}{69}\)
  \item[\(\vdash\)] \(\lnot \pvar{TestimonyArguments.SolaFide}{69}\)\gloss{\emph{at issue, the denial of:} The washing itself is necessary for salvation, though desire may justify}
  \keepwithnext
\end{derivation}
\emph{No premise here rests on scripture alone.}
\par\medskip\noindent\textbf{\texttt{desire\allowbreak{}To\allowbreak{}Not\allowbreak{}Necessary}}\par\nobreak\noindent
From what the desire brings to “the washing is not necessary”.

\begin{derivation}
  \item[] \(\pvar{TestimonyArguments.SolaFide}{71} \rightarrow \lnot \pvar{TestimonyArguments.SolaFide}{69}\)\gloss{\emph{would deny:} The washing itself is necessary for salvation, though desire may justify}
\end{derivation}
\par\medskip\noindent\textbf{\texttt{desire\allowbreak{}Case}} --- The desire for baptism (Catechism 1258–1259)\par\nobreak\noindent
\textbf{Rome’s own teaching on the desire}: it brings about the fruits of baptism, so the washing itself is not necessary for salvation.

\begin{derivation}
  \item[(1)] \(\pvar{TestimonyArguments.SolaFide}{71}\)\gloss{The desire for baptism brings about its fruits without being a sacrament}
  \item[(2)] \(\pvar{TestimonyArguments.SolaFide}{71} \rightarrow \lnot \pvar{TestimonyArguments.SolaFide}{69}\)
  \item[\(\vdash\)] \(\lnot \pvar{TestimonyArguments.SolaFide}{69}\)\gloss{\emph{at issue, the denial of:} The washing itself is necessary for salvation, though desire may justify}
  \keepwithnext
\end{derivation}
\emph{No premise here rests on scripture alone.}
\par\medskip\noindent\textbf{\texttt{desire\allowbreak{}To\allowbreak{}Not\allowbreak{}The\allowbreak{}Instrument}}\par\nobreak\noindent
The Reformed step from the desire. If justification needs only the laver or the desire for it, and the desire alone brings the fruits of baptism without the sacrament, then the sacrament is not the instrument of justification. Held by Westminster (XXVIII.5; XI.2). Rated \texttt{disputed}, because Aquinas accepts the starting points and denies the conclusion.

\begin{derivation}
  \item[] \((\pvar{TestimonyArguments.SolaFide}{70} \land \pvar{TestimonyArguments.SolaFide}{71}) \rightarrow \lnot \pvar{TestimonyArguments.SolaFide}{68}\)\gloss{\emph{would deny:} Water baptism is the instrumental cause of justification}
\end{derivation}
\par\medskip\noindent\textbf{\texttt{reformed\allowbreak{}On\allowbreak{}The\allowbreak{}Desire}} --- What justifies where only the desire is present\par\nobreak\noindent
\textbf{The question the desire raises.} Where only the desire is present, what justifies is not the washing, so the washing is not the instrument. This is the Reformed answer to Trent taken in its own words.

\begin{derivation}
  \item[(1)] \(\pvar{TestimonyArguments.SolaFide}{70}\)\gloss{Justification cannot be effected without the laver of regeneration, or the desire thereof}
  \item[(2)] \(\pvar{TestimonyArguments.SolaFide}{71}\)\gloss{The desire for baptism brings about its fruits without being a sacrament}
  \item[(3)] \((\pvar{TestimonyArguments.SolaFide}{70} \land \pvar{TestimonyArguments.SolaFide}{71}) \rightarrow \lnot \pvar{TestimonyArguments.SolaFide}{68}\)
  \item[\(\vdash\)] \(\lnot \pvar{TestimonyArguments.SolaFide}{68}\)\gloss{\emph{at issue, the denial of:} Water baptism is the instrumental cause of justification}
  \keepwithnext
\end{derivation}
\emph{No premise here rests on scripture alone.}
\paragraph*{The positions, written out}
\leavevmode\par\nobreak
\par\medskip\noindent\textbf{\texttt{rome\allowbreak{}On\allowbreak{}Baptism\allowbreak{}Reading}}\par\nobreak\noindent
Rome’s position. Baptism is the instrumental cause. The desire for it brings its fruits. The thief is saved. God is not bound to his sacraments. The water of John 3:5 is baptismal. And the washing itself is not necessary for salvation. Since baptism is a means of justification, faith is not the only means; and since justification includes sanctification, it is not a verdict only.

{\small\begin{verbatim}
def romeOnBaptismReading : Valuation Claim :=
  fun a =>
    match a with
    | Claim.paulsJustifyDenotesRenewal => False
    | Claim.rev22_11BeJustifiedStill => False
    | Claim.sir18_22BeJustifiedToDeath => False
    | Claim.jamesUsesDikaioAsPaul => False
    | Claim.waterItselfNecessaryForSalvation => False
    | Claim.john3_5WaterIsNaturalBirth => False
    | Claim.john3_5WaterIsEzekielsCleansing => False
    | Claim.john3_5WaterIsTheSpiritsCleansing => False
    | Claim.justificationByFaithAlone => False
    | Claim.justificationIsForensicOnly => False
    | x => True
\end{verbatim}}
\par\medskip\noindent\textbf{\texttt{feeney\allowbreak{}Reading}}\par\nobreak\noindent
A consistent way of holding Feeney’s reading: baptism is the instrumental cause, and the washing itself is necessary for salvation. It has to deny that the thief’s case bears on the Gospel, and that the desire brings the fruits that save. The first denial is what his reading requires; it is not a sentence quoted from him. With Trent, it denies that faith is the only means of justification, and that justification is a verdict only.

{\small\begin{verbatim}
def feeneyReading : Valuation Claim :=
  fun a =>
    match a with
    | Claim.thiefBearsOnTheGospel => False
    | Claim.desireBringsFruitsWithoutTheSacrament => False
    | Claim.john3_5WaterIsNaturalBirth => False
    | Claim.john3_5WaterIsEzekielsCleansing => False
    | Claim.john3_5WaterIsTheSpiritsCleansing => False
    | Claim.justificationByFaithAlone => False
    | Claim.justificationIsForensicOnly => False
    | x => True
\end{verbatim}}
\par\medskip\noindent\textbf{\texttt{reformed\allowbreak{}On\allowbreak{}Baptism\allowbreak{}Reading}}\par\nobreak\noindent
The Reformed position. Faith is the instrument, and baptism is not. Justification is the remission of sins, not sanctification. The desire brings the fruits, the thief is saved, and God is not bound to water. The water of John 3:5 is the Spirit’s own cleansing, and the verse does not require water baptism.

{\small\begin{verbatim}
def reformedOnBaptismReading : Valuation Claim :=
  fun a =>
    match a with
    | Claim.baptismIsInstrumentalCause => False
    | Claim.waterItselfNecessaryForSalvation => False
    | Claim.john3_5RequiresWaterBaptism => False
    | Claim.john3_5WaterIsBaptism => False
    | Claim.justificationIncludesSanctification => False
    | x => True
\end{verbatim}}
\paragraph*{Each position holds together}
\leavevmode\par\nobreak
\par\medskip\noindent\textbf{\texttt{trent\allowbreak{}On\allowbreak{}Baptism\_\allowbreak{}establishes}}\par\nobreak\noindent
Trent’s claim follows from its premises, trivially, since its one premise is the claim itself. This records the position. It does not argue for it.

{\small\begin{verbatim}
theorem trentOnBaptism_establishes : Establishes trentOnBaptism
-- axioms: propext
\end{verbatim}}
\par\medskip\noindent\textbf{\texttt{trent\allowbreak{}On\allowbreak{}Baptism\_\allowbreak{}is\_\allowbreak{}satisfiable}}\par\nobreak\noindent
Trent’s instrumental cause, as stated, can be held without contradiction.

{\small\begin{verbatim}
theorem trentOnBaptism_is_satisfiable : Satisfiable trentOnBaptism.premises
-- axioms: propext, Quot.sound
\end{verbatim}}
\par\medskip\noindent\textbf{\texttt{trent\allowbreak{}On\allowbreak{}The\allowbreak{}Water\allowbreak{}Itself\_\allowbreak{}is\_\allowbreak{}satisfiable}}\par\nobreak\noindent
\textbf{The first way of taking it is fair.} By itself, Trent taken Feeney’s way contradicts nothing, so it is not a position invented to be knocked down.

{\small\begin{verbatim}
theorem trentOnTheWaterItself_is_satisfiable : Satisfiable
    trentOnTheWaterItself.premises
-- axioms: propext, Quot.sound
\end{verbatim}}
\par\medskip\noindent\textbf{\texttt{trent\allowbreak{}On\allowbreak{}Baptism\allowbreak{}Or\allowbreak{}Desire\_\allowbreak{}is\_\allowbreak{}satisfiable}}\par\nobreak\noindent
\textbf{The second is fair too.} Trent taken in its own words can be held without contradiction.

{\small\begin{verbatim}
theorem trentOnBaptismOrDesire_is_satisfiable : Satisfiable
    trentOnBaptismOrDesire.premises
-- axioms: propext, Quot.sound
\end{verbatim}}
\par\medskip\noindent\textbf{\texttt{thief\allowbreak{}Case\_\allowbreak{}establishes}}\par\nobreak\noindent
\textbf{Grant that the thief was promised paradise unbaptised, and that his case bears on the Gospel. It follows that the washing itself is not necessary for salvation.}

{\small\begin{verbatim}
theorem thiefCase_establishes : Establishes thiefCase
-- axioms: propext, Quot.sound
\end{verbatim}}
\par\medskip\noindent\textbf{\texttt{thief\allowbreak{}Case\_\allowbreak{}is\_\allowbreak{}satisfiable}}\par\nobreak\noindent
The thief’s case can be held without contradiction.

{\small\begin{verbatim}
theorem thiefCase_is_satisfiable : Satisfiable thiefCase.premises
-- axioms: propext, Quot.sound
\end{verbatim}}
\par\medskip\noindent\textbf{\texttt{desire\allowbreak{}Case\_\allowbreak{}establishes}}\par\nobreak\noindent
\textbf{Grant Rome’s own teaching that the desire brings the fruits of baptism, and the same follows.}

{\small\begin{verbatim}
theorem desireCase_establishes : Establishes desireCase
-- axioms: propext, Quot.sound
\end{verbatim}}
\par\medskip\noindent\textbf{\texttt{desire\allowbreak{}Case\_\allowbreak{}is\_\allowbreak{}satisfiable}}\par\nobreak\noindent
The desire’s case can be held without contradiction.

{\small\begin{verbatim}
theorem desireCase_is_satisfiable : Satisfiable desireCase.premises
-- axioms: propext, Quot.sound
\end{verbatim}}
\par\medskip\noindent\textbf{\texttt{reformed\allowbreak{}On\allowbreak{}The\allowbreak{}Desire\_\allowbreak{}establishes}}\par\nobreak\noindent
\textbf{The Reformed answer reaches its conclusion}, given its step: the sacrament is not the instrument of justification.

{\small\begin{verbatim}
theorem reformedOnTheDesire_establishes : Establishes reformedOnTheDesire
-- axioms: propext, Quot.sound
\end{verbatim}}
\par\medskip\noindent\textbf{\texttt{reformed\allowbreak{}On\allowbreak{}The\allowbreak{}Desire\_\allowbreak{}is\_\allowbreak{}satisfiable}}\par\nobreak\noindent
The Reformed answer can be held without contradiction.

{\small\begin{verbatim}
theorem reformedOnTheDesire_is_satisfiable : Satisfiable
    reformedOnTheDesire.premises
-- axioms: propext, Quot.sound
\end{verbatim}}
\par\medskip\noindent\textbf{\texttt{the\_\allowbreak{}thief\_\allowbreak{}does\_\allowbreak{}not\_\allowbreak{}reach\_\allowbreak{}the\_\allowbreak{}desire}}\par\nobreak\noindent
\textbf{The thief does not count against Trent taken in its own words.} Trent’s instrumental cause, as baptism or the desire for it, can be held with everything the thief’s case says: the thief had the desire, and the desire saves.

{\small\begin{verbatim}
theorem the_thief_does_not_reach_the_desire : Satisfiable (thiefCase.premises
    ++ trentOnBaptismOrDesire.premises)
-- axioms: propext, Quot.sound
\end{verbatim}}
\paragraph*{Where each fails}
\leavevmode\par\nobreak
\par\medskip\noindent\textbf{\texttt{where\allowbreak{}The\allowbreak{}Thief\allowbreak{}Meets\allowbreak{}Feeney}}\par\nobreak\noindent
\textbf{Why the thief stands against Feeney’s reading.} Everything turns on the step from the thief: he was promised paradise unbaptised, and his case bears on the Gospel, so the washing itself is not necessary for salvation. What that contradicts is Trent’s claim taken Feeney’s way, and nothing else.

\paragraph{Why \emph{The thief on the cross (Luke 23:43), read with Augustine} stands against \emph{Trent: baptism the instrumental cause of justification, read so that the washing itself is necessary for salvation}.}
\begin{itemize}
\item \textbf{The crux:} \((\pvar{TestimonyArguments.SolaFide}{72} \land \pvar{TestimonyArguments.SolaFide}{73}) \rightarrow \lnot \pvar{TestimonyArguments.SolaFide}{69}\), a premise of \emph{The thief on the cross (Luke 23:43), read with Augustine}.
\item \textbf{What it does:} its conclusion does not follow without it.
\item \textbf{Granted:} \(\pvar{TestimonyArguments.SolaFide}{72}\), \(\pvar{TestimonyArguments.SolaFide}{73}\), which the break also needs.
\item \textbf{Where the rival breaks:} \(\pvar{TestimonyArguments.SolaFide}{68}\), \(\pvar{TestimonyArguments.SolaFide}{68} \rightarrow \pvar{TestimonyArguments.SolaFide}{69}\) cannot be held together with \((\pvar{TestimonyArguments.SolaFide}{72} \land \pvar{TestimonyArguments.SolaFide}{73}) \rightarrow \lnot \pvar{TestimonyArguments.SolaFide}{69}\) and \(\pvar{TestimonyArguments.SolaFide}{72}\), \(\pvar{TestimonyArguments.SolaFide}{73}\); each is needed for the break, and without the crux they stand.
\item \textbf{What this rests on:}
\begin{itemize}
\item \(\pvar{TestimonyArguments.SolaFide}{72}\) Luke 23:43 — the thief, unbaptised, is promised paradise --- \emph{consensus}: Luke 23:43; \cite[On Baptism IV.23.31]{npnf1-04-1887}
\item \(\pvar{TestimonyArguments.SolaFide}{73}\) The thief’s case bears on salvation under the Gospel --- \emph{well supported}: \cite[On Baptism IV.22.30–23.31]{npnf1-04-1887}; \cite[III q.66 a.11]{aquinas-summa-1920}
\item the step itself --- \emph{well supported}: \cite[On Baptism IV.23.31]{npnf1-04-1887}; \cite[III q.66 a.11]{aquinas-summa-1920}
\end{itemize}
\end{itemize}
\par\medskip\noindent\textbf{\texttt{where\allowbreak{}The\allowbreak{}Desire\allowbreak{}Meets\allowbreak{}Feeney}}\par\nobreak\noindent
\textbf{Why Rome’s own teaching stands against Feeney’s reading.} Everything turns on the step from what the desire brings: if the desire brings about the fruits of baptism, the washing itself is not necessary for salvation. Feeney denies the starting point. Whoever accepts it, as the Catechism does, cannot hold his reading.

\paragraph{Why \emph{The desire for baptism (Catechism 1258–1259)} stands against \emph{Trent: baptism the instrumental cause of justification, read so that the washing itself is necessary for salvation}.}
\begin{itemize}
\item \textbf{The crux:} \(\pvar{TestimonyArguments.SolaFide}{71} \rightarrow \lnot \pvar{TestimonyArguments.SolaFide}{69}\), a premise of \emph{The desire for baptism (Catechism 1258–1259)}.
\item \textbf{What it does:} its conclusion does not follow without it.
\item \textbf{Granted:} \(\pvar{TestimonyArguments.SolaFide}{71}\), which the break also needs.
\item \textbf{Where the rival breaks:} \(\pvar{TestimonyArguments.SolaFide}{68}\), \(\pvar{TestimonyArguments.SolaFide}{68} \rightarrow \pvar{TestimonyArguments.SolaFide}{69}\) cannot be held together with \(\pvar{TestimonyArguments.SolaFide}{71} \rightarrow \lnot \pvar{TestimonyArguments.SolaFide}{69}\) and \(\pvar{TestimonyArguments.SolaFide}{71}\); each is needed for the break, and without the crux they stand.
\item \textbf{What this rests on:}
\begin{itemize}
\item \(\pvar{TestimonyArguments.SolaFide}{71}\) The desire for baptism brings about its fruits without being a sacrament --- \emph{consensus}: \cite[1258]{catholic-church-catechism-2000}; \cite[III q.69 a.4 ad 2]{aquinas-summa-1920}
\item the step itself --- \emph{well supported}: \cite[1259]{catholic-church-catechism-2000}
\end{itemize}
\end{itemize}
\par\medskip\noindent\textbf{\texttt{where\allowbreak{}The\allowbreak{}Desire\allowbreak{}Meets\allowbreak{}Trent}}\par\nobreak\noindent
\textbf{Why the Reformed answer stands against Trent taken in its own words.} Everything turns on the step from the desire: where the desire alone brings the fruits, what justifies is not the washing, so the washing is not the instrument. What that contradicts is the instrumental cause itself. The step is rated \texttt{disputed}: Aquinas accepts that the desire brings grace, and holds that it does so as a desire for the sacrament, which then gives a fuller grace.

\paragraph{Why \emph{What justifies where only the desire is present} stands against \emph{Trent: baptism the instrumental cause of justification, read as baptism or the desire of it}.}
\begin{itemize}
\item \textbf{The crux:} \((\pvar{TestimonyArguments.SolaFide}{70} \land \pvar{TestimonyArguments.SolaFide}{71}) \rightarrow \lnot \pvar{TestimonyArguments.SolaFide}{68}\), a premise of \emph{What justifies where only the desire is present}.
\item \textbf{What it does:} its conclusion does not follow without it.
\item \textbf{Granted:} \(\pvar{TestimonyArguments.SolaFide}{70}\), \(\pvar{TestimonyArguments.SolaFide}{71}\), which the break also needs.
\item \textbf{Where the rival breaks:} \(\pvar{TestimonyArguments.SolaFide}{68}\) cannot be held together with \((\pvar{TestimonyArguments.SolaFide}{70} \land \pvar{TestimonyArguments.SolaFide}{71}) \rightarrow \lnot \pvar{TestimonyArguments.SolaFide}{68}\) and \(\pvar{TestimonyArguments.SolaFide}{70}\), \(\pvar{TestimonyArguments.SolaFide}{71}\); each is needed for the break, and without the crux they stand.
\item \textbf{What this rests on:}
\begin{itemize}
\item \(\pvar{TestimonyArguments.SolaFide}{70}\) Justification cannot be effected without the laver of regeneration, or the desire thereof --- \emph{consensus}: \cite[32]{waterworth-trent-1848}; \cite[III q.68 a.2]{aquinas-summa-1920}
\item \(\pvar{TestimonyArguments.SolaFide}{71}\) The desire for baptism brings about its fruits without being a sacrament --- \emph{consensus}: \cite[1258]{catholic-church-catechism-2000}; \cite[III q.69 a.4 ad 2]{aquinas-summa-1920}
\item the step itself --- \emph{disputed}: \cite[XXVIII.5]{westminster-confession-1647}; \cite[XI.2]{westminster-confession-1647}; \emph{disputed}: \cite[III q.69 a.4 ad 2]{aquinas-summa-1920}
\end{itemize}
\end{itemize}
\par\medskip\noindent\textbf{\texttt{what\allowbreak{}Trents\allowbreak{}Instrumental\allowbreak{}Cause\allowbreak{}Claims}}\par\nobreak\noindent
\textbf{What Trent’s instrumental cause claims, taken both ways.}

\emph{Taken so that the washing itself is necessary for salvation}, Feeney’s reading, it cannot be held with the thief as Augustine reads him, nor with Rome’s own teaching that the desire brings the fruits of baptism. Each fails at a step rated \texttt{well\allowbreak{}Supported}. Both answers are Catholic ones.

\emph{Taken in Trent’s own words}, baptism or the desire for it, the thief does not count against it. It cannot be held with the Reformed answer, that where only the desire is present what justifies is not the washing. But that step is \texttt{disputed}, by Aquinas.

Each way of taking it is fair. By itself it contradicts nothing, so it is not a position invented to be knocked down. It fails only when joined to what the positions set against it hold.

\paragraph{Every reading of \(P_{68}\), answered.} \emph{Trent: baptism the instrumental cause of justification} holds \(\pvar{TestimonyArguments.SolaFide}{68}\). It is read 2 ways here, and each reading is checked; none can be left out.

\paragraph{1. Read so that the washing itself is necessary for salvation.} The claim commits \emph{Trent: baptism the instrumental cause of justification} to \(\pvar{TestimonyArguments.SolaFide}{69}\) --- so read by \cite{feeney-bread-of-life-1952} (\emph{disputed}).

Against \emph{The thief on the cross (Luke 23:43), read with Augustine}, it cannot be held. Where it breaks:
\paragraph{Why \emph{The thief on the cross (Luke 23:43), read with Augustine} stands against \emph{Trent: baptism the instrumental cause of justification, read so that the washing itself is necessary for salvation}.}
\begin{itemize}
\item \textbf{The crux:} \((\pvar{TestimonyArguments.SolaFide}{72} \land \pvar{TestimonyArguments.SolaFide}{73}) \rightarrow \lnot \pvar{TestimonyArguments.SolaFide}{69}\), a premise of \emph{The thief on the cross (Luke 23:43), read with Augustine}.
\item \textbf{What it does:} its conclusion does not follow without it.
\item \textbf{Granted:} \(\pvar{TestimonyArguments.SolaFide}{72}\), \(\pvar{TestimonyArguments.SolaFide}{73}\), which the break also needs.
\item \textbf{Where the rival breaks:} \(\pvar{TestimonyArguments.SolaFide}{68}\), \(\pvar{TestimonyArguments.SolaFide}{68} \rightarrow \pvar{TestimonyArguments.SolaFide}{69}\) cannot be held together with \((\pvar{TestimonyArguments.SolaFide}{72} \land \pvar{TestimonyArguments.SolaFide}{73}) \rightarrow \lnot \pvar{TestimonyArguments.SolaFide}{69}\) and \(\pvar{TestimonyArguments.SolaFide}{72}\), \(\pvar{TestimonyArguments.SolaFide}{73}\); each is needed for the break, and without the crux they stand.
\item \textbf{What this rests on:}
\begin{itemize}
\item \(\pvar{TestimonyArguments.SolaFide}{72}\) Luke 23:43 — the thief, unbaptised, is promised paradise --- \emph{consensus}: Luke 23:43; \cite[On Baptism IV.23.31]{npnf1-04-1887}
\item \(\pvar{TestimonyArguments.SolaFide}{73}\) The thief’s case bears on salvation under the Gospel --- \emph{well supported}: \cite[On Baptism IV.22.30–23.31]{npnf1-04-1887}; \cite[III q.66 a.11]{aquinas-summa-1920}
\item the step itself --- \emph{well supported}: \cite[On Baptism IV.23.31]{npnf1-04-1887}; \cite[III q.66 a.11]{aquinas-summa-1920}
\end{itemize}
\end{itemize}
Against \emph{The desire for baptism (Catechism 1258–1259)}, it cannot be held. Where it breaks:
\paragraph{Why \emph{The desire for baptism (Catechism 1258–1259)} stands against \emph{Trent: baptism the instrumental cause of justification, read so that the washing itself is necessary for salvation}.}
\begin{itemize}
\item \textbf{The crux:} \(\pvar{TestimonyArguments.SolaFide}{71} \rightarrow \lnot \pvar{TestimonyArguments.SolaFide}{69}\), a premise of \emph{The desire for baptism (Catechism 1258–1259)}.
\item \textbf{What it does:} its conclusion does not follow without it.
\item \textbf{Granted:} \(\pvar{TestimonyArguments.SolaFide}{71}\), which the break also needs.
\item \textbf{Where the rival breaks:} \(\pvar{TestimonyArguments.SolaFide}{68}\), \(\pvar{TestimonyArguments.SolaFide}{68} \rightarrow \pvar{TestimonyArguments.SolaFide}{69}\) cannot be held together with \(\pvar{TestimonyArguments.SolaFide}{71} \rightarrow \lnot \pvar{TestimonyArguments.SolaFide}{69}\) and \(\pvar{TestimonyArguments.SolaFide}{71}\); each is needed for the break, and without the crux they stand.
\item \textbf{What this rests on:}
\begin{itemize}
\item \(\pvar{TestimonyArguments.SolaFide}{71}\) The desire for baptism brings about its fruits without being a sacrament --- \emph{consensus}: \cite[1258]{catholic-church-catechism-2000}; \cite[III q.69 a.4 ad 2]{aquinas-summa-1920}
\item the step itself --- \emph{well supported}: \cite[1259]{catholic-church-catechism-2000}
\end{itemize}
\end{itemize}
\paragraph{2. Read as baptism or the desire of it.} The claim commits \emph{Trent: baptism the instrumental cause of justification} to \(\pvar{TestimonyArguments.SolaFide}{70}\) --- so read by \cite[32]{waterworth-trent-1848} (\emph{consensus}).

Against \emph{The thief on the cross (Luke 23:43), read with Augustine}, it is not reached: \emph{The thief on the cross (Luke 23:43), read with Augustine} holds, and can be held together with it.

Against \emph{What justifies where only the desire is present}, it cannot be held. Where it breaks:
\paragraph{Why \emph{What justifies where only the desire is present} stands against \emph{Trent: baptism the instrumental cause of justification, read as baptism or the desire of it}.}
\begin{itemize}
\item \textbf{The crux:} \((\pvar{TestimonyArguments.SolaFide}{70} \land \pvar{TestimonyArguments.SolaFide}{71}) \rightarrow \lnot \pvar{TestimonyArguments.SolaFide}{68}\), a premise of \emph{What justifies where only the desire is present}.
\item \textbf{What it does:} its conclusion does not follow without it.
\item \textbf{Granted:} \(\pvar{TestimonyArguments.SolaFide}{70}\), \(\pvar{TestimonyArguments.SolaFide}{71}\), which the break also needs.
\item \textbf{Where the rival breaks:} \(\pvar{TestimonyArguments.SolaFide}{68}\) cannot be held together with \((\pvar{TestimonyArguments.SolaFide}{70} \land \pvar{TestimonyArguments.SolaFide}{71}) \rightarrow \lnot \pvar{TestimonyArguments.SolaFide}{68}\) and \(\pvar{TestimonyArguments.SolaFide}{70}\), \(\pvar{TestimonyArguments.SolaFide}{71}\); each is needed for the break, and without the crux they stand.
\item \textbf{What this rests on:}
\begin{itemize}
\item \(\pvar{TestimonyArguments.SolaFide}{70}\) Justification cannot be effected without the laver of regeneration, or the desire thereof --- \emph{consensus}: \cite[32]{waterworth-trent-1848}; \cite[III q.68 a.2]{aquinas-summa-1920}
\item \(\pvar{TestimonyArguments.SolaFide}{71}\) The desire for baptism brings about its fruits without being a sacrament --- \emph{consensus}: \cite[1258]{catholic-church-catechism-2000}; \cite[III q.69 a.4 ad 2]{aquinas-summa-1920}
\item the step itself --- \emph{disputed}: \cite[XXVIII.5]{westminster-confession-1647}; \cite[XI.2]{westminster-confession-1647}; \emph{disputed}: \cite[III q.69 a.4 ad 2]{aquinas-summa-1920}
\end{itemize}
\end{itemize}
\paragraph*{John 3:5}
\leavevmode\par\nobreak
\par\medskip\noindent\textbf{\texttt{john\allowbreak{}Three\allowbreak{}Five\allowbreak{}To\allowbreak{}Water}}\par\nobreak\noindent
Trent’s step from John 3:5: the verse requires water baptism. Rated \texttt{disputed}: Calvin accepts the verse and denies the step.

\begin{derivation}
  \item[] \(\pvar{TestimonyArguments.SolaFide}{77} \rightarrow \pvar{TestimonyArguments.SolaFide}{78}\)\gloss{\emph{would give:} John 3:5 requires water baptism for entry into the kingdom}
\end{derivation}
\par\medskip\noindent\textbf{\texttt{trent\allowbreak{}On\allowbreak{}John\allowbreak{}Three\allowbreak{}Five}} --- Trent on John 3:5: born of water, so water baptism required\par\nobreak\noindent
\textbf{Trent on John 3:5}: born of water and the Spirit, so the verse requires water baptism for entry into the kingdom.

\begin{derivation}
  \item[(1)] \(\pvar{TestimonyArguments.SolaFide}{77}\)\gloss{John 3:5 — unless one is born of water and the Spirit, he cannot enter the kingdom}
  \item[(2)] \(\pvar{TestimonyArguments.SolaFide}{77} \rightarrow \pvar{TestimonyArguments.SolaFide}{78}\)
  \item[\(\vdash\)] \(\pvar{TestimonyArguments.SolaFide}{78}\)\gloss{\emph{at issue:} John 3:5 requires water baptism for entry into the kingdom}
  \keepwithnext
\end{derivation}
\emph{No premise here rests on scripture alone.}
\par\medskip\noindent\textbf{\texttt{water\allowbreak{}Read\allowbreak{}As\allowbreak{}Baptism}}\par\nobreak\noindent
Who reads the water of John 3:5 as baptismal water: Trent (Session VII, On Baptism, canon 2), and “Chrysostom, with whom the greater part of expounders agree”, as Calvin reports.

{\small\begin{verbatim}
def waterReadAsBaptism : Source :=
  { primary := Reference.work waterworthTrent (Bib.Locus.page 56),
    supporting := [Reference.work calvinJohn (Bib.Locus.page 110)],
    tradition := Tradition.romanCatholic, confidence := Confidence.disputed }
\end{verbatim}}
\par\medskip\noindent\textbf{\texttt{water\allowbreak{}Read\allowbreak{}As\allowbreak{}Birth}}\par\nobreak\noindent
Who reads it as natural birth: Oliver, who argues for it as his own reading. The water “could be a reference to the amniotic fluid that surrounds the baby in the womb”.

{\small\begin{verbatim}
def waterReadAsBirth : Source :=
  { primary := Reference.work oliverWaterJohn,
    tradition := Tradition.criticalScholarship,
    confidence := Confidence.disputed }
\end{verbatim}}
\par\medskip\noindent\textbf{\texttt{water\allowbreak{}Read\allowbreak{}As\allowbreak{}Ezekiel}}\par\nobreak\noindent
Who reads it as the cleansing of Ezekiel 36: Carson, who finds there “a transformative new beginning, characterized by spectacular cleansing symbolized by water”.

{\small\begin{verbatim}
def waterReadAsEzekiel : Source :=
  { primary := Reference.work carsonBornOfWater,
    supporting := [Reference.work oliverWaterJohn],
    tradition := Tradition.christianHistoricalGrammatical,
    confidence := Confidence.disputed }
\end{verbatim}}
\par\medskip\noindent\textbf{\texttt{water\allowbreak{}Read\allowbreak{}As\allowbreak{}The\allowbreak{}Spirit}}\par\nobreak\noindent
Who reads it as the Spirit’s own cleansing: Calvin, “By ‘water and the Spirit,’ therefore, I simply understand the Spirit, which is water” (\emph{Institutes} IV.xvi.25).

{\small\begin{verbatim}
def waterReadAsTheSpirit : Source :=
  {
    primary :=
      Reference.work Bib.calvinInstitutes (Bib.Locus.sectionRef "IV.xvi.25"),
    supporting := [Reference.work calvinJohn (Bib.Locus.pages 110 111)],
    tradition := Tradition.reformedProtestant,
    confidence := Confidence.disputed }
\end{verbatim}}
\par\medskip\noindent\textbf{\texttt{water\allowbreak{}As\allowbreak{}Baptism}}\par\nobreak\noindent
The water of John 3:5, \textbf{read as baptismal water}.

{\small\begin{verbatim}
def waterAsBaptism : Reading Claim :=
  { name := "with its water as baptismal water",
    commits := p Claim.john3_5WaterIsBaptism,
    source := waterReadAsBaptism }
\end{verbatim}}
\par\medskip\noindent\textbf{\texttt{water\allowbreak{}As\allowbreak{}Birth}}\par\nobreak\noindent
The water of John 3:5, \textbf{read as natural birth}.

{\small\begin{verbatim}
def waterAsBirth : Reading Claim :=
  { name := "with its water as natural birth",
    commits := p Claim.john3_5WaterIsNaturalBirth,
    source := waterReadAsBirth }
\end{verbatim}}
\par\medskip\noindent\textbf{\texttt{water\allowbreak{}As\allowbreak{}Ezekiel}}\par\nobreak\noindent
The water of John 3:5, \textbf{read as the cleansing of Ezekiel 36}.

{\small\begin{verbatim}
def waterAsEzekiel : Reading Claim :=
  { name := "with its water as Ezekiel's cleansing",
    commits := p Claim.john3_5WaterIsEzekielsCleansing,
    source := waterReadAsEzekiel }
\end{verbatim}}
\par\medskip\noindent\textbf{\texttt{water\allowbreak{}As\allowbreak{}The\allowbreak{}Spirit}}\par\nobreak\noindent
The water of John 3:5, \textbf{read as the Spirit’s own cleansing}.

{\small\begin{verbatim}
def waterAsTheSpirit : Reading Claim :=
  { name := "with its water as the Spirit's own cleansing",
    commits := p Claim.john3_5WaterIsTheSpiritsCleansing,
    source := waterReadAsTheSpirit }
\end{verbatim}}
\par\medskip\noindent\textbf{\texttt{birth\allowbreak{}To\allowbreak{}No\allowbreak{}Water}}\par\nobreak\noindent
If the water is natural birth, the verse does not require water baptism.

\begin{derivation}
  \item[] \(\pvar{TestimonyArguments.SolaFide}{80} \rightarrow \lnot \pvar{TestimonyArguments.SolaFide}{78}\)\gloss{\emph{would deny:} John 3:5 requires water baptism for entry into the kingdom}
\end{derivation}
\par\medskip\noindent\textbf{\texttt{ezekiel\allowbreak{}To\allowbreak{}No\allowbreak{}Water}}\par\nobreak\noindent
If the water is Ezekiel’s cleansing and not baptism, the verse does not require water baptism.

\begin{derivation}
  \item[] \(\pvar{TestimonyArguments.SolaFide}{81} \rightarrow \lnot \pvar{TestimonyArguments.SolaFide}{78}\)\gloss{\emph{would deny:} John 3:5 requires water baptism for entry into the kingdom}
\end{derivation}
\par\medskip\noindent\textbf{\texttt{spirit\allowbreak{}To\allowbreak{}No\allowbreak{}Water}}\par\nobreak\noindent
If the water is the Spirit’s own cleansing, the verse does not require water baptism.

\begin{derivation}
  \item[] \(\pvar{TestimonyArguments.SolaFide}{82} \rightarrow \lnot \pvar{TestimonyArguments.SolaFide}{78}\)\gloss{\emph{would deny:} John 3:5 requires water baptism for entry into the kingdom}
\end{derivation}
\par\medskip\noindent\textbf{\texttt{spirit\allowbreak{}To\allowbreak{}Not\allowbreak{}Baptism}}\par\nobreak\noindent
And if the water is the Spirit’s own cleansing, it is not baptismal water.

\begin{derivation}
  \item[] \((\pvar{TestimonyArguments.SolaFide}{77} \land \pvar{TestimonyArguments.SolaFide}{82}) \rightarrow \lnot \pvar{TestimonyArguments.SolaFide}{79}\)\gloss{\emph{would deny:} The water of John 3:5 is baptismal water}
\end{derivation}
\par\medskip\noindent\textbf{\texttt{birth\allowbreak{}On\allowbreak{}John\allowbreak{}Three\allowbreak{}Five}} --- John 3:5, the water as natural birth (Oliver)\par\nobreak\noindent
\textbf{The natural-birth reading, as a position}: the water is natural birth, so the verse does not require water baptism.

\begin{derivation}
  \item[(1)] \(\pvar{TestimonyArguments.SolaFide}{80}\)\gloss{The water of John 3:5 is natural birth}
  \item[(2)] \(\pvar{TestimonyArguments.SolaFide}{80} \rightarrow \lnot \pvar{TestimonyArguments.SolaFide}{78}\)
  \item[\(\vdash\)] \(\lnot \pvar{TestimonyArguments.SolaFide}{78}\)\gloss{\emph{at issue, the denial of:} John 3:5 requires water baptism for entry into the kingdom}
  \keepwithnext
\end{derivation}
\emph{No premise here rests on scripture alone.}
\par\medskip\noindent\textbf{\texttt{ezekiel\allowbreak{}On\allowbreak{}John\allowbreak{}Three\allowbreak{}Five}} --- John 3:5, the water as Ezekiel’s cleansing (Carson)\par\nobreak\noindent
\textbf{The Ezekiel reading, as a position.}

\begin{derivation}
  \item[(1)] \(\pvar{TestimonyArguments.SolaFide}{81}\)\gloss{The water of John 3:5 is the cleansing promised in Ezekiel 36:25–27, not baptism}
  \item[(2)] \(\pvar{TestimonyArguments.SolaFide}{81} \rightarrow \lnot \pvar{TestimonyArguments.SolaFide}{78}\)
  \item[\(\vdash\)] \(\lnot \pvar{TestimonyArguments.SolaFide}{78}\)\gloss{\emph{at issue, the denial of:} John 3:5 requires water baptism for entry into the kingdom}
  \keepwithnext
\end{derivation}
\emph{No premise here rests on scripture alone.}
\par\medskip\noindent\textbf{\texttt{calvin\allowbreak{}On\allowbreak{}John\allowbreak{}Three\allowbreak{}Five}} --- John 3:5, the water as the Spirit’s own cleansing (Calvin)\par\nobreak\noindent
\textbf{Calvin’s reading, as a position}: water and Spirit are one, so the water is not baptismal water, and the verse does not require water baptism.

\begin{derivation}
  \item[(1)] \(\pvar{TestimonyArguments.SolaFide}{77}\)\gloss{John 3:5 — unless one is born of water and the Spirit, he cannot enter the kingdom}
  \item[(2)] \(\pvar{TestimonyArguments.SolaFide}{82}\)\gloss{The water of John 3:5 is the Spirit’s own cleansing: water and Spirit are one}
  \item[(3)] \(\pvar{TestimonyArguments.SolaFide}{82} \rightarrow \lnot \pvar{TestimonyArguments.SolaFide}{78}\)
  \item[(4)] \((\pvar{TestimonyArguments.SolaFide}{77} \land \pvar{TestimonyArguments.SolaFide}{82}) \rightarrow \lnot \pvar{TestimonyArguments.SolaFide}{79}\)\gloss{\emph{would deny:} The water of John 3:5 is baptismal water}
  \item[\(\vdash\)] \(\lnot \pvar{TestimonyArguments.SolaFide}{78}\)\gloss{\emph{at issue, the denial of:} John 3:5 requires water baptism for entry into the kingdom}
  \keepwithnext
\end{derivation}
\emph{No premise here rests on scripture alone.}
\par\medskip\noindent\textbf{\texttt{no\allowbreak{}Water\allowbreak{}In\allowbreak{}John\allowbreak{}Three\allowbreak{}Five\allowbreak{}Reading}}\par\nobreak\noindent
A reading on which all three non-baptismal readings hold at once: the water is natural birth, Ezekiel’s cleansing and the Spirit’s own cleansing, and not baptism, and the verse does not require water baptism. Nobody holds this, since the three readings are rivals. It is used only to show that each position’s premises can be true together.

{\small\begin{verbatim}
def noWaterInJohnThreeFiveReading : Valuation Claim :=
  fun a =>
    match a with
    | Claim.john3_5RequiresWaterBaptism => False
    | Claim.john3_5WaterIsBaptism => False
    | x => True
\end{verbatim}}
\par\medskip\noindent\textbf{\texttt{trent\allowbreak{}On\allowbreak{}John\allowbreak{}Three\allowbreak{}Five\_\allowbreak{}establishes}}\par\nobreak\noindent
Grant the verse and Trent’s step from it, and the requirement of water baptism follows. The step is \texttt{disputed}.

{\small\begin{verbatim}
theorem trentOnJohnThreeFive_establishes : Establishes trentOnJohnThreeFive
-- axioms: propext, Quot.sound
\end{verbatim}}
\par\medskip\noindent\textbf{\texttt{trent\allowbreak{}On\allowbreak{}John\allowbreak{}Three\allowbreak{}Five\_\allowbreak{}is\_\allowbreak{}satisfiable}}\par\nobreak\noindent
Trent on John 3:5 can be held without contradiction, on Rome’s position.

{\small\begin{verbatim}
theorem trentOnJohnThreeFive_is_satisfiable : Satisfiable
    trentOnJohnThreeFive.premises
-- axioms: propext, Quot.sound
\end{verbatim}}
\par\medskip\noindent\textbf{\texttt{birth\allowbreak{}On\allowbreak{}John\allowbreak{}Three\allowbreak{}Five\_\allowbreak{}establishes}}\par\nobreak\noindent
The natural-birth reading reaches its conclusion.

{\small\begin{verbatim}
theorem birthOnJohnThreeFive_establishes : Establishes birthOnJohnThreeFive
-- axioms: propext, Quot.sound
\end{verbatim}}
\par\medskip\noindent\textbf{\texttt{birth\allowbreak{}On\allowbreak{}John\allowbreak{}Three\allowbreak{}Five\_\allowbreak{}is\_\allowbreak{}satisfiable}}\par\nobreak\noindent
The natural-birth reading can be held without contradiction.

{\small\begin{verbatim}
theorem birthOnJohnThreeFive_is_satisfiable : Satisfiable
    birthOnJohnThreeFive.premises
-- axioms: propext, Quot.sound
\end{verbatim}}
\par\medskip\noindent\textbf{\texttt{ezekiel\allowbreak{}On\allowbreak{}John\allowbreak{}Three\allowbreak{}Five\_\allowbreak{}establishes}}\par\nobreak\noindent
The Ezekiel reading reaches its conclusion.

{\small\begin{verbatim}
theorem ezekielOnJohnThreeFive_establishes : Establishes
    ezekielOnJohnThreeFive
-- axioms: propext, Quot.sound
\end{verbatim}}
\par\medskip\noindent\textbf{\texttt{ezekiel\allowbreak{}On\allowbreak{}John\allowbreak{}Three\allowbreak{}Five\_\allowbreak{}is\_\allowbreak{}satisfiable}}\par\nobreak\noindent
The Ezekiel reading can be held without contradiction.

{\small\begin{verbatim}
theorem ezekielOnJohnThreeFive_is_satisfiable : Satisfiable
    ezekielOnJohnThreeFive.premises
-- axioms: propext, Quot.sound
\end{verbatim}}
\par\medskip\noindent\textbf{\texttt{calvin\allowbreak{}On\allowbreak{}John\allowbreak{}Three\allowbreak{}Five\_\allowbreak{}establishes}}\par\nobreak\noindent
Calvin’s reading reaches its conclusion.

{\small\begin{verbatim}
theorem calvinOnJohnThreeFive_establishes : Establishes calvinOnJohnThreeFive
-- axioms: propext, Quot.sound
\end{verbatim}}
\par\medskip\noindent\textbf{\texttt{calvin\allowbreak{}On\allowbreak{}John\allowbreak{}Three\allowbreak{}Five\_\allowbreak{}is\_\allowbreak{}satisfiable}}\par\nobreak\noindent
Calvin’s reading can be held without contradiction.

{\small\begin{verbatim}
theorem calvinOnJohnThreeFive_is_satisfiable : Satisfiable
    calvinOnJohnThreeFive.premises
-- axioms: propext, Quot.sound
\end{verbatim}}
\par\medskip\noindent\textbf{\texttt{trent\allowbreak{}With\allowbreak{}Baptismal\allowbreak{}Water}} --- Trent on John 3:5: born of water, so water baptism required, read with its water as baptismal water\par\nobreak\noindent
Trent on John 3:5, with the water read as baptism.

\begin{derivation}
  \item[(1)] \(\pvar{TestimonyArguments.SolaFide}{77}\)\gloss{John 3:5 — unless one is born of water and the Spirit, he cannot enter the kingdom}
  \item[(2)] \(\pvar{TestimonyArguments.SolaFide}{77} \rightarrow \pvar{TestimonyArguments.SolaFide}{78}\)
  \item[(3)] \(\pvar{TestimonyArguments.SolaFide}{77} \rightarrow \pvar{TestimonyArguments.SolaFide}{79}\)\gloss{\emph{would give:} The water of John 3:5 is baptismal water}
  \item[\(\vdash\)] \(\pvar{TestimonyArguments.SolaFide}{78}\)\gloss{\emph{at issue:} John 3:5 requires water baptism for entry into the kingdom}
  \keepwithnext
\end{derivation}
\emph{No premise here rests on scripture alone.}
\par\medskip\noindent\textbf{\texttt{trent\allowbreak{}With\allowbreak{}Natural\allowbreak{}Birth}} --- Trent on John 3:5: born of water, so water baptism required, read with its water as natural birth\par\nobreak\noindent
Trent on John 3:5, with the water read as natural birth.

\begin{derivation}
  \item[(1)] \(\pvar{TestimonyArguments.SolaFide}{77}\)\gloss{John 3:5 — unless one is born of water and the Spirit, he cannot enter the kingdom}
  \item[(2)] \(\pvar{TestimonyArguments.SolaFide}{77} \rightarrow \pvar{TestimonyArguments.SolaFide}{78}\)
  \item[(3)] \(\pvar{TestimonyArguments.SolaFide}{77} \rightarrow \pvar{TestimonyArguments.SolaFide}{80}\)\gloss{\emph{would give:} The water of John 3:5 is natural birth}
  \item[\(\vdash\)] \(\pvar{TestimonyArguments.SolaFide}{78}\)\gloss{\emph{at issue:} John 3:5 requires water baptism for entry into the kingdom}
  \keepwithnext
\end{derivation}
\emph{No premise here rests on scripture alone.}
\par\medskip\noindent\textbf{\texttt{trent\allowbreak{}With\allowbreak{}Ezekiels\allowbreak{}Cleansing}} --- Trent on John 3:5: born of water, so water baptism required, read with its water as Ezekiel’s cleansing\par\nobreak\noindent
Trent on John 3:5, with the water read as Ezekiel’s cleansing.

\begin{derivation}
  \item[(1)] \(\pvar{TestimonyArguments.SolaFide}{77}\)\gloss{John 3:5 — unless one is born of water and the Spirit, he cannot enter the kingdom}
  \item[(2)] \(\pvar{TestimonyArguments.SolaFide}{77} \rightarrow \pvar{TestimonyArguments.SolaFide}{78}\)
  \item[(3)] \(\pvar{TestimonyArguments.SolaFide}{77} \rightarrow \pvar{TestimonyArguments.SolaFide}{81}\)\gloss{\emph{would give:} The water of John 3:5 is the cleansing promised in Ezekiel 36:25–27, not baptism}
  \item[\(\vdash\)] \(\pvar{TestimonyArguments.SolaFide}{78}\)\gloss{\emph{at issue:} John 3:5 requires water baptism for entry into the kingdom}
  \keepwithnext
\end{derivation}
\emph{No premise here rests on scripture alone.}
\par\medskip\noindent\textbf{\texttt{trent\allowbreak{}With\allowbreak{}The\allowbreak{}Spirits\allowbreak{}Cleansing}} --- Trent on John 3:5: born of water, so water baptism required, read with its water as the Spirit’s own cleansing\par\nobreak\noindent
Trent on John 3:5, with the water read as the Spirit’s own cleansing.

\begin{derivation}
  \item[(1)] \(\pvar{TestimonyArguments.SolaFide}{77}\)\gloss{John 3:5 — unless one is born of water and the Spirit, he cannot enter the kingdom}
  \item[(2)] \(\pvar{TestimonyArguments.SolaFide}{77} \rightarrow \pvar{TestimonyArguments.SolaFide}{78}\)
  \item[(3)] \(\pvar{TestimonyArguments.SolaFide}{77} \rightarrow \pvar{TestimonyArguments.SolaFide}{82}\)\gloss{\emph{would give:} The water of John 3:5 is the Spirit’s own cleansing: water and Spirit are one}
  \item[\(\vdash\)] \(\pvar{TestimonyArguments.SolaFide}{78}\)\gloss{\emph{at issue:} John 3:5 requires water baptism for entry into the kingdom}
  \keepwithnext
\end{derivation}
\emph{No premise here rests on scripture alone.}
\par\medskip\noindent\textbf{\texttt{trent\allowbreak{}Reading\allowbreak{}Of\allowbreak{}John\allowbreak{}Three\allowbreak{}Five}}\par\nobreak\noindent
The verse read Trent’s way and no other: its water is baptismal water, and not natural birth, Ezekiel’s cleansing or the Spirit’s own cleansing.

{\small\begin{verbatim}
def trentReadingOfJohnThreeFive : Valuation Claim :=
  fun a =>
    match a with
    | Claim.john3_5WaterIsNaturalBirth => False
    | Claim.john3_5WaterIsEzekielsCleansing => False
    | Claim.john3_5WaterIsTheSpiritsCleansing => False
    | x => True
\end{verbatim}}
\par\medskip\noindent\textbf{\texttt{trent\allowbreak{}With\allowbreak{}Birth\allowbreak{}Reading}}\par\nobreak\noindent
The verse read with its water as natural birth, and Trent’s step kept. By itself this contradicts nothing. It fails only when that reading’s own step is added: then the verse does not require water baptism.

{\small\begin{verbatim}
def trentWithBirthReading : Valuation Claim :=
  fun a =>
    match a with
    | Claim.john3_5WaterIsEzekielsCleansing => False
    | Claim.john3_5WaterIsTheSpiritsCleansing => False
    | x => True
\end{verbatim}}
\par\medskip\noindent\textbf{\texttt{trent\allowbreak{}With\allowbreak{}Ezekiel\allowbreak{}Reading}}\par\nobreak\noindent
Trent’s step, with the water read as Ezekiel’s cleansing.

{\small\begin{verbatim}
def trentWithEzekielReading : Valuation Claim :=
  fun a =>
    match a with
    | Claim.john3_5WaterIsNaturalBirth => False
    | Claim.john3_5WaterIsTheSpiritsCleansing => False
    | x => True
\end{verbatim}}
\par\medskip\noindent\textbf{\texttt{trent\allowbreak{}With\allowbreak{}Spirit\allowbreak{}Reading}}\par\nobreak\noindent
Trent’s step, with the water read as the Spirit’s own cleansing.

{\small\begin{verbatim}
def trentWithSpiritReading : Valuation Claim :=
  fun a =>
    match a with
    | Claim.john3_5WaterIsNaturalBirth => False
    | Claim.john3_5WaterIsEzekielsCleansing => False
    | x => True
\end{verbatim}}
\par\medskip\noindent\textbf{\texttt{trent\allowbreak{}With\allowbreak{}Baptismal\allowbreak{}Water\_\allowbreak{}is\_\allowbreak{}satisfiable}}\par\nobreak\noindent
The baptismal reading of Trent’s step can be held without contradiction.

{\small\begin{verbatim}
theorem trentWithBaptismalWater_is_satisfiable : Satisfiable
    trentWithBaptismalWater.premises
-- axioms: propext, Quot.sound
\end{verbatim}}
\par\medskip\noindent\textbf{\texttt{trent\allowbreak{}With\allowbreak{}Natural\allowbreak{}Birth\_\allowbreak{}is\_\allowbreak{}satisfiable}}\par\nobreak\noindent
The natural-birth reading of Trent’s step can be held without contradiction.

{\small\begin{verbatim}
theorem trentWithNaturalBirth_is_satisfiable : Satisfiable
    trentWithNaturalBirth.premises
-- axioms: propext, Quot.sound
\end{verbatim}}
\par\medskip\noindent\textbf{\texttt{trent\allowbreak{}With\allowbreak{}Ezekiels\allowbreak{}Cleansing\_\allowbreak{}is\_\allowbreak{}satisfiable}}\par\nobreak\noindent
The Ezekiel reading of Trent’s step can be held without contradiction.

{\small\begin{verbatim}
theorem trentWithEzekielsCleansing_is_satisfiable : Satisfiable
    trentWithEzekielsCleansing.premises
-- axioms: propext, Quot.sound
\end{verbatim}}
\par\medskip\noindent\textbf{\texttt{trent\allowbreak{}With\allowbreak{}The\allowbreak{}Spirits\allowbreak{}Cleansing\_\allowbreak{}is\_\allowbreak{}satisfiable}}\par\nobreak\noindent
The Spirit’s-cleansing reading of Trent’s step can be held without contradiction.

{\small\begin{verbatim}
theorem trentWithTheSpiritsCleansing_is_satisfiable : Satisfiable
    trentWithTheSpiritsCleansing.premises
-- axioms: propext, Quot.sound
\end{verbatim}}
\par\medskip\noindent\textbf{\texttt{where\allowbreak{}Calvin\allowbreak{}Meets\allowbreak{}The\allowbreak{}Baptismal\allowbreak{}Reading}}\par\nobreak\noindent
\textbf{Why Calvin’s reading stands against Trent with the water read as baptism.} Everything turns on Calvin’s step: if water and Spirit are one thing, the water is not baptismal water. What that contradicts is the baptismal reading itself. The step is \texttt{well\allowbreak{}Supported}. But Calvin’s reading, which it starts from, is \texttt{disputed}, and Trent condemns it (Session VII, On Baptism, canon 2).

\paragraph{Why \emph{John 3:5, the water as the Spirit’s own cleansing (Calvin)} stands against \emph{Trent on John 3:5: born of water, so water baptism required, read with its water as baptismal water}.}
\begin{itemize}
\item \textbf{The crux:} \((\pvar{TestimonyArguments.SolaFide}{77} \land \pvar{TestimonyArguments.SolaFide}{82}) \rightarrow \lnot \pvar{TestimonyArguments.SolaFide}{79}\), a premise of \emph{John 3:5, the water as the Spirit’s own cleansing (Calvin)}.
\item \textbf{What it does:} its conclusion follows without it --- the crux is its answer to \emph{Trent on John 3:5: born of water, so water baptism required, read with its water as baptismal water}.
\item \textbf{Granted:} \(\pvar{TestimonyArguments.SolaFide}{77}\), \(\pvar{TestimonyArguments.SolaFide}{82}\), which the break also needs.
\item \textbf{Where the rival breaks:} \(\pvar{TestimonyArguments.SolaFide}{77} \rightarrow \pvar{TestimonyArguments.SolaFide}{79}\) cannot be held together with \((\pvar{TestimonyArguments.SolaFide}{77} \land \pvar{TestimonyArguments.SolaFide}{82}) \rightarrow \lnot \pvar{TestimonyArguments.SolaFide}{79}\) and \(\pvar{TestimonyArguments.SolaFide}{77}\), \(\pvar{TestimonyArguments.SolaFide}{82}\); each is needed for the break, and without the crux they stand.
\item \textbf{What this rests on:}
\begin{itemize}
\item \(\pvar{TestimonyArguments.SolaFide}{77}\) John 3:5 — unless one is born of water and the Spirit, he cannot enter the kingdom --- \emph{consensus}: John 3:5; \cite{oliver-water-john-2022}
\item \(\pvar{TestimonyArguments.SolaFide}{82}\) The water of John 3:5 is the Spirit’s own cleansing: water and Spirit are one --- \emph{disputed}: \cite[110--111]{calvin-commentary-john-1847}; \cite[IV.xvi.25]{calvin-institutes-1960}
\item the step itself --- \emph{well supported}: \cite{oliver-water-john-2022}; \cite[IV.xvi.25]{calvin-institutes-1960}
\end{itemize}
\end{itemize}
\par\medskip\noindent\textbf{\texttt{where\allowbreak{}Birth\allowbreak{}Meets\allowbreak{}Trents\allowbreak{}Step}}\par\nobreak\noindent
\textbf{Why the natural-birth reading cannot keep Trent’s step.} If the water is natural birth, the verse does not require water baptism. So Trent’s step from the verse, with the verse read that way, fails, at a step rated \texttt{well\allowbreak{}Supported}.

\paragraph{Why \emph{John 3:5, the water as natural birth (Oliver)} stands against \emph{Trent on John 3:5: born of water, so water baptism required, read with its water as natural birth}.}
\begin{itemize}
\item \textbf{The crux:} \(\pvar{TestimonyArguments.SolaFide}{80} \rightarrow \lnot \pvar{TestimonyArguments.SolaFide}{78}\), a premise of \emph{John 3:5, the water as natural birth (Oliver)}.
\item \textbf{What it does:} its conclusion does not follow without it.
\item \textbf{Where the rival breaks:} \(\pvar{TestimonyArguments.SolaFide}{77}\), \(\pvar{TestimonyArguments.SolaFide}{77} \rightarrow \pvar{TestimonyArguments.SolaFide}{78}\), \(\pvar{TestimonyArguments.SolaFide}{77} \rightarrow \pvar{TestimonyArguments.SolaFide}{80}\) cannot be held together with \(\pvar{TestimonyArguments.SolaFide}{80} \rightarrow \lnot \pvar{TestimonyArguments.SolaFide}{78}\); each is needed for the break, and without the crux they stand.
\item \textbf{What this rests on:}
\begin{itemize}
\item the step itself --- \emph{well supported}: \cite{oliver-water-john-2022}; \cite[IV.xvi.25]{calvin-institutes-1960}
\end{itemize}
\end{itemize}
\par\medskip\noindent\textbf{\texttt{where\allowbreak{}Ezekiel\allowbreak{}Meets\allowbreak{}Trents\allowbreak{}Step}}\par\nobreak\noindent
\textbf{Why the Ezekiel reading cannot keep Trent’s step.} If the water is Ezekiel’s cleansing and not baptism, the verse does not require water baptism.

\paragraph{Why \emph{John 3:5, the water as Ezekiel’s cleansing (Carson)} stands against \emph{Trent on John 3:5: born of water, so water baptism required, read with its water as Ezekiel’s cleansing}.}
\begin{itemize}
\item \textbf{The crux:} \(\pvar{TestimonyArguments.SolaFide}{81} \rightarrow \lnot \pvar{TestimonyArguments.SolaFide}{78}\), a premise of \emph{John 3:5, the water as Ezekiel’s cleansing (Carson)}.
\item \textbf{What it does:} its conclusion does not follow without it.
\item \textbf{Where the rival breaks:} \(\pvar{TestimonyArguments.SolaFide}{77}\), \(\pvar{TestimonyArguments.SolaFide}{77} \rightarrow \pvar{TestimonyArguments.SolaFide}{78}\), \(\pvar{TestimonyArguments.SolaFide}{77} \rightarrow \pvar{TestimonyArguments.SolaFide}{81}\) cannot be held together with \(\pvar{TestimonyArguments.SolaFide}{81} \rightarrow \lnot \pvar{TestimonyArguments.SolaFide}{78}\); each is needed for the break, and without the crux they stand.
\item \textbf{What this rests on:}
\begin{itemize}
\item the step itself --- \emph{well supported}: \cite{oliver-water-john-2022}; \cite[IV.xvi.25]{calvin-institutes-1960}
\end{itemize}
\end{itemize}
\par\medskip\noindent\textbf{\texttt{where\allowbreak{}The\allowbreak{}Spirit\allowbreak{}Meets\allowbreak{}Trents\allowbreak{}Step}}\par\nobreak\noindent
\textbf{Why Calvin’s reading cannot keep Trent’s step.} If the water is the Spirit’s own cleansing, the verse does not require water baptism.

\paragraph{Why \emph{John 3:5, the water as the Spirit’s own cleansing (Calvin)} stands against \emph{Trent on John 3:5: born of water, so water baptism required, read with its water as the Spirit’s own cleansing}.}
\begin{itemize}
\item \textbf{The crux:} \(\pvar{TestimonyArguments.SolaFide}{82} \rightarrow \lnot \pvar{TestimonyArguments.SolaFide}{78}\), a premise of \emph{John 3:5, the water as the Spirit’s own cleansing (Calvin)}.
\item \textbf{What it does:} its conclusion does not follow without it.
\item \textbf{Where the rival breaks:} \(\pvar{TestimonyArguments.SolaFide}{77}\), \(\pvar{TestimonyArguments.SolaFide}{77} \rightarrow \pvar{TestimonyArguments.SolaFide}{78}\), \(\pvar{TestimonyArguments.SolaFide}{77} \rightarrow \pvar{TestimonyArguments.SolaFide}{82}\) cannot be held together with \(\pvar{TestimonyArguments.SolaFide}{82} \rightarrow \lnot \pvar{TestimonyArguments.SolaFide}{78}\); each is needed for the break, and without the crux they stand.
\item \textbf{What this rests on:}
\begin{itemize}
\item the step itself --- \emph{well supported}: \cite{oliver-water-john-2022}; \cite[IV.xvi.25]{calvin-institutes-1960}
\end{itemize}
\end{itemize}
\par\medskip\noindent\textbf{\texttt{what\allowbreak{}John\allowbreak{}Three\allowbreak{}Five\allowbreak{}Means}}\par\nobreak\noindent
\textbf{What John 3:5 means, read four ways.} What can be taken more than one way is the verse Trent cites, not Trent. Trent argues: unless one is born of water and the Spirit, he cannot enter the kingdom of God; therefore water baptism is required. Each of the four readings is a reading of the verse’s “water”, held by the reader named. For each, the question is whether Trent’s step from the verse still stands once the verse is read that way.

Read with its water as baptismal water, Trent’s step stands. That reading cannot be held with Calvin’s, at a \texttt{well\allowbreak{}Supported} step from Calvin’s reading, which is itself \texttt{disputed}. Even so, Trent’s step stays \texttt{disputed}.

Read with its water as natural birth, as Ezekiel’s cleansing, or as the Spirit’s own cleansing, Trent’s step cannot be kept, each time at a \texttt{well\allowbreak{}Supported} step.

So only the baptismal reading leaves Trent’s step standing. Which reading is right is not decided here: each is cited, and each is disputed. Each is fair. By itself it contradicts nothing, and it fails only when joined to the step its own holders draw from it.

\paragraph{Every reading of \(P_{77}\), answered.} \emph{Trent on John 3:5: born of water, so water baptism required} holds \(\pvar{TestimonyArguments.SolaFide}{77}\). It is read 4 ways here, and each reading is checked; none can be left out.

\paragraph{1. Read with its water as baptismal water.} The claim commits \emph{Trent on John 3:5: born of water, so water baptism required} to \(\pvar{TestimonyArguments.SolaFide}{79}\) --- so read by \cite[56]{waterworth-trent-1848}; \cite[110]{calvin-commentary-john-1847} (\emph{disputed}).

Against \emph{John 3:5, the water as the Spirit’s own cleansing (Calvin)}, it cannot be held. Where it breaks:
\paragraph{Why \emph{John 3:5, the water as the Spirit’s own cleansing (Calvin)} stands against \emph{Trent on John 3:5: born of water, so water baptism required, read with its water as baptismal water}.}
\begin{itemize}
\item \textbf{The crux:} \((\pvar{TestimonyArguments.SolaFide}{77} \land \pvar{TestimonyArguments.SolaFide}{82}) \rightarrow \lnot \pvar{TestimonyArguments.SolaFide}{79}\), a premise of \emph{John 3:5, the water as the Spirit’s own cleansing (Calvin)}.
\item \textbf{What it does:} its conclusion follows without it --- the crux is its answer to \emph{Trent on John 3:5: born of water, so water baptism required, read with its water as baptismal water}.
\item \textbf{Granted:} \(\pvar{TestimonyArguments.SolaFide}{77}\), \(\pvar{TestimonyArguments.SolaFide}{82}\), which the break also needs.
\item \textbf{Where the rival breaks:} \(\pvar{TestimonyArguments.SolaFide}{77} \rightarrow \pvar{TestimonyArguments.SolaFide}{79}\) cannot be held together with \((\pvar{TestimonyArguments.SolaFide}{77} \land \pvar{TestimonyArguments.SolaFide}{82}) \rightarrow \lnot \pvar{TestimonyArguments.SolaFide}{79}\) and \(\pvar{TestimonyArguments.SolaFide}{77}\), \(\pvar{TestimonyArguments.SolaFide}{82}\); each is needed for the break, and without the crux they stand.
\item \textbf{What this rests on:}
\begin{itemize}
\item \(\pvar{TestimonyArguments.SolaFide}{77}\) John 3:5 — unless one is born of water and the Spirit, he cannot enter the kingdom --- \emph{consensus}: John 3:5; \cite{oliver-water-john-2022}
\item \(\pvar{TestimonyArguments.SolaFide}{82}\) The water of John 3:5 is the Spirit’s own cleansing: water and Spirit are one --- \emph{disputed}: \cite[110--111]{calvin-commentary-john-1847}; \cite[IV.xvi.25]{calvin-institutes-1960}
\item the step itself --- \emph{well supported}: \cite{oliver-water-john-2022}; \cite[IV.xvi.25]{calvin-institutes-1960}
\end{itemize}
\end{itemize}
\paragraph{2. Read with its water as natural birth.} The claim commits \emph{Trent on John 3:5: born of water, so water baptism required} to \(\pvar{TestimonyArguments.SolaFide}{80}\) --- so read by \cite{oliver-water-john-2022} (\emph{disputed}).

Against \emph{John 3:5, the water as natural birth (Oliver)}, it cannot be held. Where it breaks:
\paragraph{Why \emph{John 3:5, the water as natural birth (Oliver)} stands against \emph{Trent on John 3:5: born of water, so water baptism required, read with its water as natural birth}.}
\begin{itemize}
\item \textbf{The crux:} \(\pvar{TestimonyArguments.SolaFide}{80} \rightarrow \lnot \pvar{TestimonyArguments.SolaFide}{78}\), a premise of \emph{John 3:5, the water as natural birth (Oliver)}.
\item \textbf{What it does:} its conclusion does not follow without it.
\item \textbf{Where the rival breaks:} \(\pvar{TestimonyArguments.SolaFide}{77}\), \(\pvar{TestimonyArguments.SolaFide}{77} \rightarrow \pvar{TestimonyArguments.SolaFide}{78}\), \(\pvar{TestimonyArguments.SolaFide}{77} \rightarrow \pvar{TestimonyArguments.SolaFide}{80}\) cannot be held together with \(\pvar{TestimonyArguments.SolaFide}{80} \rightarrow \lnot \pvar{TestimonyArguments.SolaFide}{78}\); each is needed for the break, and without the crux they stand.
\item \textbf{What this rests on:}
\begin{itemize}
\item the step itself --- \emph{well supported}: \cite{oliver-water-john-2022}; \cite[IV.xvi.25]{calvin-institutes-1960}
\end{itemize}
\end{itemize}
\paragraph{3. Read with its water as Ezekiel’s cleansing.} The claim commits \emph{Trent on John 3:5: born of water, so water baptism required} to \(\pvar{TestimonyArguments.SolaFide}{81}\) --- so read by \cite{carson-born-of-water-2019}; \cite{oliver-water-john-2022} (\emph{disputed}).

Against \emph{John 3:5, the water as Ezekiel’s cleansing (Carson)}, it cannot be held. Where it breaks:
\paragraph{Why \emph{John 3:5, the water as Ezekiel’s cleansing (Carson)} stands against \emph{Trent on John 3:5: born of water, so water baptism required, read with its water as Ezekiel’s cleansing}.}
\begin{itemize}
\item \textbf{The crux:} \(\pvar{TestimonyArguments.SolaFide}{81} \rightarrow \lnot \pvar{TestimonyArguments.SolaFide}{78}\), a premise of \emph{John 3:5, the water as Ezekiel’s cleansing (Carson)}.
\item \textbf{What it does:} its conclusion does not follow without it.
\item \textbf{Where the rival breaks:} \(\pvar{TestimonyArguments.SolaFide}{77}\), \(\pvar{TestimonyArguments.SolaFide}{77} \rightarrow \pvar{TestimonyArguments.SolaFide}{78}\), \(\pvar{TestimonyArguments.SolaFide}{77} \rightarrow \pvar{TestimonyArguments.SolaFide}{81}\) cannot be held together with \(\pvar{TestimonyArguments.SolaFide}{81} \rightarrow \lnot \pvar{TestimonyArguments.SolaFide}{78}\); each is needed for the break, and without the crux they stand.
\item \textbf{What this rests on:}
\begin{itemize}
\item the step itself --- \emph{well supported}: \cite{oliver-water-john-2022}; \cite[IV.xvi.25]{calvin-institutes-1960}
\end{itemize}
\end{itemize}
\paragraph{4. Read with its water as the Spirit’s own cleansing.} The claim commits \emph{Trent on John 3:5: born of water, so water baptism required} to \(\pvar{TestimonyArguments.SolaFide}{82}\) --- so read by \cite[IV.xvi.25]{calvin-institutes-1960}; \cite[110--111]{calvin-commentary-john-1847} (\emph{disputed}).

Against \emph{John 3:5, the water as the Spirit’s own cleansing (Calvin)}, it cannot be held. Where it breaks:
\paragraph{Why \emph{John 3:5, the water as the Spirit’s own cleansing (Calvin)} stands against \emph{Trent on John 3:5: born of water, so water baptism required, read with its water as the Spirit’s own cleansing}.}
\begin{itemize}
\item \textbf{The crux:} \(\pvar{TestimonyArguments.SolaFide}{82} \rightarrow \lnot \pvar{TestimonyArguments.SolaFide}{78}\), a premise of \emph{John 3:5, the water as the Spirit’s own cleansing (Calvin)}.
\item \textbf{What it does:} its conclusion does not follow without it.
\item \textbf{Where the rival breaks:} \(\pvar{TestimonyArguments.SolaFide}{77}\), \(\pvar{TestimonyArguments.SolaFide}{77} \rightarrow \pvar{TestimonyArguments.SolaFide}{78}\), \(\pvar{TestimonyArguments.SolaFide}{77} \rightarrow \pvar{TestimonyArguments.SolaFide}{82}\) cannot be held together with \(\pvar{TestimonyArguments.SolaFide}{82} \rightarrow \lnot \pvar{TestimonyArguments.SolaFide}{78}\); each is needed for the break, and without the crux they stand.
\item \textbf{What this rests on:}
\begin{itemize}
\item the step itself --- \emph{well supported}: \cite{oliver-water-john-2022}; \cite[IV.xvi.25]{calvin-institutes-1960}
\end{itemize}
\end{itemize}
\paragraph*{Acts, both orders}
\leavevmode\par\nobreak
\par\medskip\noindent\textbf{\texttt{acts\allowbreak{}To\allowbreak{}Unbound}}\par\nobreak\noindent
From Samaria and Caesarea: the Spirit comes after the water and before it, so God is not bound to water baptism. Rated \texttt{consensus}: Calvin and the Catechism both accept that much, though Calvin’s own sentence goes further.

\begin{derivation}
  \item[] \((\pvar{TestimonyArguments.SolaFide}{74} \land \pvar{TestimonyArguments.SolaFide}{75}) \rightarrow \pvar{TestimonyArguments.SolaFide}{76}\)\gloss{\emph{would give:} God is not bound to water baptism: he gives the Spirit before it and after it}
\end{derivation}
\par\medskip\noindent\textbf{\texttt{acts\allowbreak{}Case}} --- Acts 8 and 10: the Spirit after the water, and before it\par\nobreak\noindent
\textbf{Acts, both orders}: the Spirit given after the water in Samaria and before it at Caesarea, so God is not bound to water baptism.

\begin{derivation}
  \item[(1)] \(\pvar{TestimonyArguments.SolaFide}{74}\)\gloss{Acts 8:14–17 — the Samaritans, baptized, receive the Spirit afterwards}
  \item[(2)] \(\pvar{TestimonyArguments.SolaFide}{75}\)\gloss{Acts 10:44–48 — the Spirit falls on Cornelius’s household before baptism}
  \item[(3)] \((\pvar{TestimonyArguments.SolaFide}{74} \land \pvar{TestimonyArguments.SolaFide}{75}) \rightarrow \pvar{TestimonyArguments.SolaFide}{76}\)
  \item[\(\vdash\)] \(\pvar{TestimonyArguments.SolaFide}{76}\)\gloss{\emph{at issue:} God is not bound to water baptism: he gives the Spirit before it and after it}
  \keepwithnext
\end{derivation}
\emph{No premise here rests on scripture alone.}
\par\medskip\noindent\textbf{\texttt{acts\allowbreak{}Case\_\allowbreak{}establishes}}\par\nobreak\noindent
\textbf{Grant the two orders in Acts, and it follows that God is not bound to water.}

{\small\begin{verbatim}
theorem actsCase_establishes : Establishes actsCase
-- axioms: propext, Quot.sound
\end{verbatim}}
\par\medskip\noindent\textbf{\texttt{acts\allowbreak{}Case\_\allowbreak{}is\_\allowbreak{}satisfiable}}\par\nobreak\noindent
Acts’ case can be held without contradiction.

{\small\begin{verbatim}
theorem actsCase_is_satisfiable : Satisfiable actsCase.premises
-- axioms: propext, Quot.sound
\end{verbatim}}
\par\medskip\noindent\textbf{\texttt{trent\_\allowbreak{}grants\_\allowbreak{}god\_\allowbreak{}is\_\allowbreak{}not\_\allowbreak{}bound}}\par\nobreak\noindent
\textbf{Trent’s instrumental cause is consistent with God not being bound to it.} Rome holds both: “God has bound salvation to the sacrament of Baptism, but he himself is not bound by his sacraments” (Catechism 1257). So what Acts shows is accepted on both sides, and is not a point won by the Reformed. What remains in dispute is the first half of that sentence, that God has bound salvation to the sacrament, and Trent’s further claim that the sacrament is the instrument of justification.

{\small\begin{verbatim}
theorem trent_grants_god_is_not_bound : Grants trentOnBaptism (p
    Claim.godNotBoundToWater)
-- axioms: propext, Quot.sound
\end{verbatim}}
\paragraph*{“And fire”: three readings, and none chosen}
\leavevmode\par\nobreak
\par\medskip\noindent\textbf{\texttt{fire\allowbreak{}As\allowbreak{}Judgment\allowbreak{}Reading}}\par\nobreak\noindent
A reading on which the fire is judgement, and neither of the other readings holds. This is enough to show that the verse does not require them. Many readers take the fire as both judging and purifying; nothing here says the readings exclude each other.

{\small\begin{verbatim}
def fireAsJudgmentReading : Valuation Claim :=
  fun a =>
    match a with
    | Claim.fireIsPurification => False
    | Claim.fireIsPentecost => False
    | x => True
\end{verbatim}}
\par\medskip\noindent\textbf{\texttt{fire\allowbreak{}As\allowbreak{}Purification\allowbreak{}Reading}}\par\nobreak\noindent
A reading on which the fire is purification, and neither of the other readings holds. This is enough to show that the verse does not require them.

{\small\begin{verbatim}
def fireAsPurificationReading : Valuation Claim :=
  fun a =>
    match a with
    | Claim.fireIsJudgment => False
    | Claim.fireIsPentecost => False
    | x => True
\end{verbatim}}
\par\medskip\noindent\textbf{\texttt{the\_\allowbreak{}baptist\_\allowbreak{}leaves\_\allowbreak{}the\_\allowbreak{}fire\_\allowbreak{}open}}\par\nobreak\noindent
\textbf{John the Baptist’s words leave open whether the fire is judgement.} “He will baptize you with the Holy Spirit and fire” is kept alike by Gregory Nazianzen, who reads the fire as the burning of the chaff, and by Calvin, who reads it as the Spirit’s purifying. The text does not settle it either way.

{\small\begin{verbatim}
theorem the_baptist_leaves_the_fire_open : Independent [p
    Claim.matthew3_11SpiritAndFire] (p Claim.fireIsJudgment)
-- axioms: propext, Quot.sound
\end{verbatim}}
\par\medskip\noindent\textbf{\texttt{the\_\allowbreak{}baptist\_\allowbreak{}leaves\_\allowbreak{}purification\_\allowbreak{}open}}\par\nobreak\noindent
\textbf{Nor whether it is purification.}

{\small\begin{verbatim}
theorem the_baptist_leaves_purification_open : Independent [p
    Claim.matthew3_11SpiritAndFire] (p Claim.fireIsPurification)
-- axioms: propext, Quot.sound
\end{verbatim}}
\par\medskip\noindent\textbf{\texttt{the\_\allowbreak{}baptist\_\allowbreak{}leaves\_\allowbreak{}pentecost\_\allowbreak{}open}}\par\nobreak\noindent
\textbf{Nor whether it is Pentecost}, as Cyril of Jerusalem reads it: “because the descent of the Holy Ghost was in fiery tongues”.

{\small\begin{verbatim}
theorem the_baptist_leaves_pentecost_open : Independent [p
    Claim.matthew3_11SpiritAndFire] (p Claim.fireIsPentecost)
-- axioms: propext, Quot.sound
\end{verbatim}}
\subsection*{Weighing the positions}
\addcontentsline{toc}{subsection}{Weighing the positions}
The earlier chapters each heard one question: Trent against Paul’s gospel, James 2:24, Jesus’ words, baptism. This one hears all the positions at once and asks which are left standing.

Seventeen positions are heard.

\begin{itemize}
  \item The four Reformed strands, each argued alone: Paul’s, Luke’s, Peter’s and John’s.
  \item The Council of Trent.
  \item The apocalyptic reading of Paul (Martyn, Campbell).
  \item Sanders and Dunn on Judaism in Paul’s day, and the critics who answer them.
  \item Jervell on the yoke of Acts 15.
  \item Aquinas on John 6:29, and Luther’s answer to him.
  \item Chrysostom on John 3:36 and Matthew 7:21, and the Reformed answer to him: the obedience those texts require is sanctification, the fruit of faith.
  \item Paul’s gospel in Galatians, set against Trent’s definition of justification.
  \item The meaning of Paul’s word “justify”, argued two ways: from how the word is used, and from Romans 4.
  \item Trent, with its definition taken as a claim about that word.
\end{itemize}
Who defeats whom is not decided by hand. Each defeat, and each absence of one, is worked out from the two positions’ premises and checked by machine.

\subsubsection*{The words used here}
\begin{itemize}
  \item \textbf{Defeats.} One position contradicts another, and is not the weaker of the two. To contradict it is to deny one of its premises, its conclusion, or something it derives on the way.
  \item \textbf{Weakest link.} A position is only as strong as its least-supported premise or step. Ratings run from \texttt{disputed}, the lowest, through \texttt{plausible} and \texttt{well\allowbreak{}Supported}, to \texttt{consensus}. Denying a claim counts as \texttt{disputed}.
  \item \textbf{Tie.} Two positions defeat each other. A reader may hold either.
  \item \textbf{Meets.} One position meets another when they contradict each other, so that one defeats the other unless it is the weaker. Where both defeat, they meet “in both directions”.
  \item \textbf{Forced} (also: \emph{prevails outright}). The position must be accepted however every tie is settled. Either nothing defeats it, or every position that defeats it is itself defeated by a forced position. The technical name for the set of forced positions is the \emph{grounded extension}.
  \item \textbf{Can be defended.} Some consistent set of positions includes it and answers every attack on its members. A set is consistent when no member defeats another.
  \item \textbf{Not forced.} Some set of positions that can be defended, and is as large as it can be, leaves it out. Not forced does not mean false. It means the debate, as argued here, cannot make a reader accept it. The verdicts below that say \emph{not forced} use this test; \emph{forced} means the test above.
  \item \textbf{Cannot be defended.} No consistent set of positions can include it and answer its attackers.
\end{itemize}
\subsubsection*{The ratings decide one thing}
Almost every position here has a weakest link rated \texttt{disputed}.

\begin{itemize}
  \item Each Reformed strand uses a step that Trent’s canon 9 denies: to faith alone in Paul’s case and John’s, and to “faith suffices” in Luke’s and Peter’s.
  \item Paul’s, Peter’s and John’s cases hold the Reformed distinction between justification and sanctification, which Trent condemns.
  \item Trent rests on its definition of justification, which the Westminster Confession denies.
  \item The apocalyptic reading denies a claim, and a premise that denies a claim is rated \texttt{disputed} by rule.
\end{itemize}
Between such positions no rating blocks an attack. The outcome depends only on who contradicts whom.

The two exceptions are the two arguments about Paul’s word. The first is from usage (\texttt{lexical\allowbreak{}Case}): “justify” (δικαιόω) is a courtroom word, Paul has other words for renewal, and a word contributes the least meaning its context requires. The second is from Romans 4 (\texttt{romans\allowbreak{}Four\allowbreak{}On\allowbreak{}The\allowbreak{}Word}): Paul explains the righteousness God counts to the ungodly as sin not counted, and by the same rule the verb he explains in that way does not also mean renewal. Every link of each is rated \texttt{well\allowbreak{}Supported} or better. Each contradicts one position, Trent taken as a claim about Paul’s word, and those are the attacks the ratings decide.

\subsubsection*{Who defeats whom}
\begin{itemize}
  \item \textbf{Trent and each Reformed strand defeat each other.} Paul’s, Peter’s and John’s cases conclude “not by works”, which Trent denies. Trent’s definition of justification, “not remission of sins merely, but also the sanctification and renewal of the inward man” (Session VI, ch. 7), contradicts the distinction between justification and sanctification that they hold. That is where Trent’s objection comes from: put the distinction in place of the definition, and Trent’s case no longer denies “not by works” (\texttt{trent\_\allowbreak{}objection\_\allowbreak{}rests\_\allowbreak{}on\_\allowbreak{}its\_\allowbreak{}definition}). Luke’s case holds only Luke’s words and meets Trent on one point. Trent’s definition, with canon 9, denies that faith without charity suffices. Luke’s case reads “Your faith has saved you” (Luke 7:50) as saying that her faith sufficed.
  \item \textbf{The apocalyptic reading and Paul’s case defeat each other.} It denies that Galatians 2:16 speaks of faith \emph{in} Christ, and Paul’s case concludes faith alone. \textbf{It does not conflict with Luke’s case or Peter’s.} They say that faith saves, which it grants. It says that faith is not the condition, which they do not deny.
  \item \textbf{The apocalyptic reading and Trent defeat each other.} They agree that justification is not by faith alone. But the apocalyptic reading derives “not by works”, since God’s rescue depends on nothing a person does, and Trent denies that.
  \item \textbf{Sanders defeats Paul’s case}, by denying the reading of “works of the law”, and Paul’s case defeats him in return. \textbf{Sanders and the critics defeat each other.}
  \item \textbf{Jervell and Peter’s case defeat each other}, over the yoke.
  \item \textbf{John’s case meets four positions, each in both directions.} Trent, however its definition is taken, over whether faith suffices and over “not by works”, as with Paul and Peter. The apocalyptic reading, because John’s case concludes faith alone, as Paul’s does and Luke’s does not. And Aquinas, over what the believing of John 6:29 is: Aquinas concludes it is not bare trust, and John’s case holds that it is.
  \item \textbf{Chrysostom meets every position that concludes faith suffices, each in both directions}: Paul’s case, Luke’s, Peter’s and John’s, and the Reformed answer to him. He contradicts no text they read. He contradicts the step each takes to “faith suffices”, and each contradicts his step to “faith does not suffice”. Each pair is a tie at \texttt{disputed}. The Reformed answer meets Trent too, however Trent’s definition is taken, because it holds the distinction between justification and sanctification.
  \item \textbf{Aquinas and Luther defeat each other.} Luther concludes that the faith which justifies is not faith formed by charity, which is Aquinas’s premise. Luther meets no one else. Aquinas meets only John’s case besides.
  \item \textbf{Paul’s gospel and Trent defeat each other}, over Trent’s definition. The gospel line concludes that justification is not the renewal of the inward person. Both rest on something rated \texttt{disputed}.
  \item \textbf{Paul’s word defeats Trent taken as a claim about it, and is not defeated in return} (\texttt{lexical\_\allowbreak{}defeats\_\allowbreak{}trent\_\allowbreak{}on\_\allowbreak{}pauls\_\allowbreak{}word}, \texttt{trent\_\allowbreak{}on\_\allowbreak{}pauls\_\allowbreak{}word\_\allowbreak{}does\_\allowbreak{}not\_\allowbreak{}defeat\_\allowbreak{}lexical}). Taken that way, Trent’s case derives that “justify” means renewal. The argument from the word denies it. Trent’s reply rests on a definition rated \texttt{disputed}, and fails against an argument rated \texttt{well\allowbreak{}Supported}. Romans 4, read for the word, does the same from one passage. Neither argument conflicts with anything else. In particular neither conflicts with Trent’s case as Trent states it, whose definition, taken as it stands, says nothing about the word. Trent taken as a claim about Paul’s word also meets every position Trent meets, in both directions: the four strands, the apocalyptic reading, Paul’s gospel and the Reformed answer to Chrysostom.
\end{itemize}
\subsubsection*{What follows}
\subsubsection*{Only the meaning of Paul’s word prevails outright}
The checking finds that the meaning of Paul’s word is the one thing forced (\texttt{only\_\allowbreak{}pauls\_\allowbreak{}word\_\allowbreak{}prevails}). The forced positions are exactly the two arguments about it. Every other position is attacked by one the debate cannot rule out. One loses: Trent taken as a claim about Paul’s word. Every other is in a tie with a rival the debate cannot choose between.

\subsubsection*{Trent, taken as a claim about Paul’s word, cannot be defended}
This is the one position that loses (\texttt{trent\_\allowbreak{}on\_\allowbreak{}pauls\_\allowbreak{}word\_\allowbreak{}indefensible}). Both arguments about the word defeat it, and nothing answers either.

\subsubsection*{Beyond that, nothing wins outright}
Luke’s case, Paul’s gospel and Trent as Trent states it can each be defended, and none is forced (\texttt{dominical\_\allowbreak{}case\_\allowbreak{}defensible}, \texttt{dominical\_\allowbreak{}case\_\allowbreak{}not\_\allowbreak{}forced}, \texttt{gospel\_\allowbreak{}defensible}, \texttt{gospel\_\allowbreak{}not\_\allowbreak{}forced}, \texttt{trent\_\allowbreak{}defensible}, \texttt{trent\_\allowbreak{}not\_\allowbreak{}forced}). Every attack between Trent and a Reformed position runs in both directions, with every rating at the bottom, and the debate has no ground to choose between them.

\subsubsection*{Between Paul and the apocalyptic reading, the debate chooses neither}
Each defeats the other over “faith of Christ” (πίστις Χριστοῦ). Each can be defended, and neither is forced (\texttt{apocalyptic\_\allowbreak{}defensible}, \texttt{apocalyptic\_\allowbreak{}not\_\allowbreak{}forced}). The apocalyptic reading can be held together with Luke’s case and Peter’s, which claim that faith saves and not that faith alone does.

\subsubsection*{What the verdict rests on}
\subsubsection*{What Trent’s definition claims}
The definition can be taken two ways, and the chapter “What Trent means by justification” answers both (\texttt{what\allowbreak{}Trents\allowbreak{}Definition\allowbreak{}Claims}). Taken as a claim about Paul’s word, it cannot be held with the evidence for what the word means, and in this debate it cannot be defended. Taken as a claim about what God does in justifying, the argument from the word does not reach it. What then divides it from Paul is one step: that a verdict on Christ’s finished work excludes the renewal worked in us (\texttt{where\allowbreak{}Trent\allowbreak{}Parts\allowbreak{}From\allowbreak{}Paul}). That step is rated \texttt{disputed}. Trent’s case as Trent states it stands here for that second way of taking it. So the debate does not settle between Trent and the Reformed. It shows that Trent can keep its definition only as a claim about what God does, not about Paul’s word, and it names the one step on which that claim depends.

\subsubsection*{An absence}
No source cited here argues that “saved” (σῴζω) at Luke 7:50 means healed, so no position attacks Luke’s case at that reading. Trent, however its definition is taken, and Chrysostom attack it at the step to “faith suffices”. That absence is also the premise’s rating. It is \texttt{well\allowbreak{}Supported}, not \texttt{disputed}: a search for a scholar arguing the healing sense \emph{at 7:50} found none, and at 7:50 no one is ill. The saying follows “your sins are forgiven”. Marshall’s commentary on the Greek text says the same: the saying refers mainly to healing elsewhere in Luke, and is used here “in a case where physical healing is not involved”. The rating is not \texttt{consensus}, because the same words mean “made you well” at Luke 8:48, 17:19 and 18:42. One older lexicon, Thayer’s, lists 7:50 with the healing uses without arguing it. It is not yet cited, and whether an unargued listing should lower the rating is an open question. A cited argument for the healing reading would attack Luke’s case and lower the rating.

\subsubsection*{An objection built in, not heard}
“Her sins, which are many, are forgiven—for she loved much” (7:47) is the verse a rival would use to make love, not faith, the reason she was forgiven. Luke’s case does not leave that open. It carries, as a premise of its own, that her love is the evidence of her forgiveness and not its cause (\texttt{luke7\_\allowbreak{}47\allowbreak{}Love\allowbreak{}Is\allowbreak{}Evidence}), as the rest of the verse and the parable of 7:41–43 suggest. Marshall argues this from the Greek: “for” (ὅτι) here means “as is evidenced by the fact that”. Padilla reads it so: she “loves much because she has been forgiven much”. So does Melanchthon in the \emph{Apology of the Augsburg Confession}: “Christ interprets Himself when He adds: Thy faith hath saved thee”, so the woman did not merit forgiveness “by that work of love”. The premise is rated \texttt{disputed}. Marshall names those who read the love as the cause: Wellhausen, Creed, and most Roman Catholic commentators, Lagrange and Schürmann among them. Trent counts beginning to love God among the dispositions to justification (Session VI, ch. 6). None of them is heard here as a position of his own. No one in this debate argues that reading of 7:47, and none of those commentators is yet cited from his own work.

\subsubsection*{John}
John’s case, Aquinas’s reading of John 6:29 and Luther’s answer are each in a tie with a position rated no higher than itself. So none of the three is forced, and none of them changes what is forced. Without Trent, John’s case is still not forced, because Aquinas’s reading can be defended (\texttt{sola\_\allowbreak{}fide\_\allowbreak{}from\_\allowbreak{}john\_\allowbreak{}not\_\allowbreak{}forced\_\allowbreak{}without\_\allowbreak{}trent}).

\subsubsection*{Who is heard}
A debate is over the arguments put into it. Hays’s reading of “faith of Christ” is represented through the apocalyptic reading, which holds it. A position for Hays alone, one for the \emph{Joint Declaration}, and one for the Finnish reading of Luther would each be new arguments to weigh. So would a cited argument that Paul’s “justify”, granted to be a courtroom word, still means renewal: none was found, and the rating of the argument from the word says so. Romans 4 as an argument about what God does (\texttt{romans\allowbreak{}Four\allowbreak{}Case}) is set against Trent in the chapter “What Trent means by justification”, and is not added to this debate. It meets Trent exactly where Paul’s gospel does, at Trent’s definition, from a step rated \texttt{disputed}, and a position of its own would be a second argument to weigh.

\subsubsection*{Two narrower questions}
The full debate hears the Reformers, Trent and several schools of modern scholarship at once. The two hearings below each set some positions aside and keep every defeat the full debate established. Nothing is argued again.

\subsubsection*{What did the Reformation itself decide?}
The sixteenth-century debate was between Rome and the Reformers, before anyone had proposed a New Perspective on Paul or an apocalyptic reading of him. Heard that way (\texttt{reformation\allowbreak{}Alone}), the apocalyptic reading, Sanders and his critics, and Jervell are set aside. Aquinas and Luther, and Chrysostom and the Reformed answer to him, are still heard, and modern scholars are still cited for the premises of the positions that remain. One thing is forced: Paul’s “justify” (δικαιόω) does not mean the renewal of the inward person. So Trent cannot claim that its definition is what Paul’s word means. What fails there is not Augustine’s explanation of the word, which he certainly gave, but the step from a Latin explanation to what the Greek word means.

Trent’s definition taken as a claim about what God does is not touched by anything about the word. There Trent and the Reformers tie: each defeats the other, and the ratings, all \texttt{disputed}, cannot choose. The Reformation debate, heard alone, comes down to what justification \emph{is}, which is the question of the chapter “What Trent means by justification” (\texttt{what\allowbreak{}Trents\allowbreak{}Definition\allowbreak{}Claims}).

\subsubsection*{If Trent were not heard, would Scripture establish sola fide?}
Heard without Trent (\texttt{without\allowbreak{}Trent}), two things are forced: the meaning of Paul’s word; and Paul’s gospel, that justification is not the renewal of the inward person. \emph{Faith alone} is not forced, from Paul, from Peter or from John, and neither is Luke’s “your faith has saved you”, that faith suffices. Five disagreements remain with Trent left out, though Aquinas, a Catholic reader, is still heard: three in modern New Testament scholarship, and two older.

\begin{itemize}
  \item Campbell’s apocalyptic reading takes “faith of Christ” (πίστις Χριστοῦ) as Christ’s own faithfulness, and meets Paul’s case there.
  \item Sanders and Dunn read “works of the law” (ἔργα νόμου) as the marks of Israel. The critics answer them, and are answered in return.
  \item Jervell reads the yoke of Acts 15 as something other than the law as a condition of salvation, and meets Peter’s case.
  \item Aquinas reads the believing of John 6:29 as faith formed by charity, and meets John’s case. Luther answers him from Galatians 3, and is answered in return.
  \item Chrysostom reads John 3:36 and Matthew 7:21 as requiring a right life beside believing, and meets every case that concludes faith suffices, Luke’s among them. The Reformed answer, that the obedience those texts require is sanctification, the fruit of faith, meets him and is met in return.
\end{itemize}
Four of these are questions about what a text says, and one, Sanders against his critics, about what Judaism in Paul’s day taught. Leaving Trent out does not make them go away.

\subsubsection*{Augustine, in both hearings}
Augustine stays where the full debate has him. In the Reformation hearing, his explanation is what Trent’s reading of Paul’s word rests on. Without Trent, it goes out with that reading, since nothing else rests on it. In both he remains the cited reader who accepts Romans 4:5–8 and reads it the other way, which is why the step from Romans 4 is rated \texttt{disputed}. A reader is not removed from these pages because he is judged mistaken. The pages show where his reading fails, and where it does not.

\subsubsection*{How this is checked}
\paragraph*{The positions}
\leavevmode\par\nobreak
Each position heard is an argument in its own right: premises that can be held together, a conclusion that follows from them, and steps that someone has rated. The Reformed case is heard as four positions, one for each strand. A debate weighs each argument whole. One position holding all four strands would be defeated by a contradiction of any one of them, and the fact that each strand reaches the conclusion without the others would count for nothing.

\par\medskip\noindent\textbf{\texttt{reformed\allowbreak{}Ontology}}\par\nobreak\noindent
What the Reformed hold about what justification \emph{is}: justification and sanctification are distinct, and sanctification follows. No strand needs this to reach sola fide. But it is the premise Trent’s definition contradicts, so it belongs to Paul’s, Peter’s and John’s cases, and Trent’s attack on each falls on it as well as on “not by works”. It is not added to the Reformed case itself, because the other positions built from that case, the Finnish reading among them, would then hold it on their authors’ behalf.

\begin{derivation}
  \item[(1)] \(\pvar{TestimonyArguments.SolaFide}{34}\)\gloss{Justification and sanctification are inseparable but distinct}
\end{derivation}
\par\medskip\noindent\textbf{\texttt{pauline\allowbreak{}Case}} --- Sola fide from Paul (Galatians 2:16)\par\nobreak\noindent
\textbf{Paul’s case}: sola fide from Galatians 2:16, argued alone.

Its reading of “works of the law” is derived, not assumed. It rests on the critics’ line, which denies Sanders’ account of Judaism and reads Galatians as written against requiring circumcision. That is where the scholarly debate lies (\texttt{critics\_\allowbreak{}carry\_\allowbreak{}the\_\allowbreak{}pauline\_\allowbreak{}strand}, \texttt{sanders\_\allowbreak{}costs\_\allowbreak{}the\_\allowbreak{}pauline\_\allowbreak{}strand}). So Sanders meets Paul’s case here over the history of Judaism in Paul’s day, not over Paul’s Greek. The critics’ inference is rated with the rest, as Gathercole’s.

\begin{derivation}
  \item[(1)] \(\pvar{TestimonyArguments.SolaFide}{8}\)\gloss{Galatians opposes requiring circumcision, besides faith, for justification}
  \item[(2)] \(\pvar{TestimonyArguments.SolaFide}{7}\)\gloss{In Galatians 2:16 ‘faith of Christ’ (πίστις Χριστοῦ) means our faith in Christ}
  \item[(3)] \(\lnot \pvar{TestimonyArguments.SolaFide}{31}\)\gloss{\emph{denied:} Second Temple Judaism: in by grace, staying in by works}
  \breakrow{(4)} \(\pvar{TestimonyArguments.SolaFide}{8}\)\gloss{Galatians opposes requiring circumcision, besides faith, for justification}
  \breakrow{(5)} \(\pvar{TestimonyArguments.SolaFide}{1}\)\gloss{Ephesians 2:8–9 teaches salvation by grace through faith, not of works}
  \breakrow{(6)} \(\pvar{TestimonyArguments.SolaFide}{2}\)\gloss{Romans 3:28 teaches justification by faith apart from works of the law}
  \breakrow{(7)} \(\pvar{TestimonyArguments.SolaFide}{3}\)\gloss{Galatians 2:16 teaches that no one is justified by works of the law}
  \breakrow{(8)} \(\pvar{TestimonyArguments.SolaFide}{4}\)\gloss{Romans 4:4–5 contrasts wages owed with a gift reckoned to the one who believes}
  \breakrow{(9)} \(\pvar{TestimonyArguments.SolaFide}{5}\)\gloss{Titus 3:5 teaches that God saved us not by works done in righteousness}
  \breakrow{(10)} \(\pvar{TestimonyArguments.SolaFide}{17}\)\gloss{Luke 7:50 — Jesus says ‘your faith has saved you’ after declaring sins forgiven}
  \breakrow{(11)} \(\pvar{TestimonyArguments.SolaFide}{9}\)\gloss{Acts 15:9–11 — Peter: hearts cleansed by faith; saved through grace}
  \breakrow{(12)} \(\pvar{TestimonyArguments.SolaFide}{20}\)\gloss{James 2:14–26 targets a barren faith — mere assent — not Paul’s doctrine}
  \breakrow{(13)} \(\pvar{TestimonyArguments.SolaFide}{21}\)\gloss{Good works are the fruit and evidence of saving faith, not its ground}
  \breakrow{(14)} \(\pvar{TestimonyArguments.SolaFide}{24}\)\gloss{Scripture does not contradict itself}
  \breakrow{(15)} \(\pvar{TestimonyArguments.SolaFide}{34}\)\gloss{Justification and sanctification are inseparable but distinct}
  \breakrow{(16)} \((\pvar{TestimonyArguments.SolaFide}{2} \land \pvar{TestimonyArguments.SolaFide}{3} \land \pvar{TestimonyArguments.SolaFide}{8} \land \pvar{TestimonyArguments.SolaFide}{7} \land \pvar{TestimonyArguments.SolaFide}{6}) \rightarrow \pvar{TestimonyArguments.SolaFide}{25}\)\gloss{\emph{would give:} Justification is by faith alone}
  \breakrow{(17)} \((\lnot \pvar{TestimonyArguments.SolaFide}{31} \land \pvar{TestimonyArguments.SolaFide}{8}) \rightarrow \pvar{TestimonyArguments.SolaFide}{6}\)\gloss{\emph{would give:} Paul’s ‘works of the law’ means human works in general, not boundary markers}
  \breakrow{(18)} \((\pvar{TestimonyArguments.SolaFide}{20} \land \pvar{TestimonyArguments.SolaFide}{21}) \rightarrow \pvar{TestimonyArguments.SolaFide}{22}\)\gloss{\emph{would give:} James 2:24 is compatible with Paul, using ‘justify’ in a different sense}
  \breakrow{(19)} \((\pvar{TestimonyArguments.SolaFide}{1} \land \pvar{TestimonyArguments.SolaFide}{4} \land \pvar{TestimonyArguments.SolaFide}{5}) \rightarrow \pvar{TestimonyArguments.SolaFide}{27}\)\gloss{\emph{would give:} Salvation is by grace, a gift and not wages owed}
  \breakrow{(20)} \((\pvar{TestimonyArguments.SolaFide}{1} \land \pvar{TestimonyArguments.SolaFide}{4} \land \pvar{TestimonyArguments.SolaFide}{5} \land \pvar{TestimonyArguments.SolaFide}{22} \land \pvar{TestimonyArguments.SolaFide}{24}) \rightarrow \pvar{TestimonyArguments.SolaFide}{28}\)\gloss{\emph{would give:} Salvation is not by works: no work is its ground}
  \breakrow{(21)} \(\pvar{TestimonyArguments.SolaFide}{25} \rightarrow \pvar{TestimonyArguments.SolaFide}{26}\)\gloss{\emph{would give:} Faith is sufficient: whoever believes is saved}
  \item[(22)] \((\pvar{TestimonyArguments.SolaFide}{26} \land \pvar{TestimonyArguments.SolaFide}{1}) \rightarrow \pvar{TestimonyArguments.SolaFide}{29}\)\gloss{\emph{would give:} Salvation is received through faith}
  \item[\(\vdash\)] \(\pvar{TestimonyArguments.SolaFide}{27} \land \pvar{TestimonyArguments.SolaFide}{28} \land \pvar{TestimonyArguments.SolaFide}{29}\)\gloss{\emph{at issue:} \(\pvar{TestimonyArguments.SolaFide}{27}\) Salvation is by grace, a gift and not wages owed; \(\pvar{TestimonyArguments.SolaFide}{28}\) Salvation is not by works: no work is its ground; \(\pvar{TestimonyArguments.SolaFide}{29}\) Salvation is received through faith}
  \keepwithnext
\end{derivation}
\emph{No premise here rests on scripture alone.}
\par\medskip\noindent\textbf{\texttt{dominical\allowbreak{}Case}} --- Luke 7:50: “your faith has saved you”\par\nobreak\noindent
\textbf{Luke’s case}: what Jesus says at Luke 7:50, and no more. “Your faith has saved you”, with “saved” (σῴζω) read as saved from sin and the woman’s love (7:47) read as the evidence of her forgiveness: her faith sufficed.

It carries nothing Luke does not say. It does not carry Paul’s texts, the answer to James, or the Reformed account of what justification is. And it does not claim what Paul argues, that faith \emph{alone} saves. So it is attacked only by what denies Luke’s own claim.

\begin{derivation}
  \item[(1)] \(\pvar{TestimonyArguments.SolaFide}{17}\)\gloss{Luke 7:50 — Jesus says ‘your faith has saved you’ after declaring sins forgiven}
  \item[(2)] \(\pvar{TestimonyArguments.SolaFide}{18}\)\gloss{‘Saved’ (σῴζω) at Luke 7:50 means saved from sin, not healed}
  \item[(3)] \(\pvar{TestimonyArguments.SolaFide}{19}\)\gloss{Luke 7:47 — her love is the evidence of her forgiveness, not its ground}
  \item[(4)] \((\pvar{TestimonyArguments.SolaFide}{17} \land \pvar{TestimonyArguments.SolaFide}{18} \land \pvar{TestimonyArguments.SolaFide}{19}) \rightarrow \pvar{TestimonyArguments.SolaFide}{26}\)
  \item[\(\vdash\)] \(\pvar{TestimonyArguments.SolaFide}{26}\)\gloss{\emph{at issue:} Faith is sufficient: whoever believes is saved}
  \keepwithnext
\end{derivation}
\emph{No premise here rests on scripture alone.}
\par\medskip\noindent\textbf{\texttt{apostolic\allowbreak{}Case}} --- Sola fide from Peter (Acts 15:9–11)\par\nobreak\noindent
\textbf{Peter’s case}: sola fide from Peter’s speech at Jerusalem (Acts 15:7–11), argued alone.

\begin{derivation}
  \item[(1)] \(\pvar{TestimonyArguments.SolaFide}{10}\)\gloss{The yoke refused at Acts 15:10 is the whole law as a condition of salvation}
  \item[(2)] \(\pvar{TestimonyArguments.SolaFide}{1}\)\gloss{Ephesians 2:8–9 teaches salvation by grace through faith, not of works}
  \item[(3)] \(\pvar{TestimonyArguments.SolaFide}{2}\)\gloss{Romans 3:28 teaches justification by faith apart from works of the law}
  \breakrow{(4)} \(\pvar{TestimonyArguments.SolaFide}{3}\)\gloss{Galatians 2:16 teaches that no one is justified by works of the law}
  \breakrow{(5)} \(\pvar{TestimonyArguments.SolaFide}{4}\)\gloss{Romans 4:4–5 contrasts wages owed with a gift reckoned to the one who believes}
  \breakrow{(6)} \(\pvar{TestimonyArguments.SolaFide}{5}\)\gloss{Titus 3:5 teaches that God saved us not by works done in righteousness}
  \breakrow{(7)} \(\pvar{TestimonyArguments.SolaFide}{17}\)\gloss{Luke 7:50 — Jesus says ‘your faith has saved you’ after declaring sins forgiven}
  \breakrow{(8)} \(\pvar{TestimonyArguments.SolaFide}{9}\)\gloss{Acts 15:9–11 — Peter: hearts cleansed by faith; saved through grace}
  \breakrow{(9)} \(\pvar{TestimonyArguments.SolaFide}{20}\)\gloss{James 2:14–26 targets a barren faith — mere assent — not Paul’s doctrine}
  \breakrow{(10)} \(\pvar{TestimonyArguments.SolaFide}{21}\)\gloss{Good works are the fruit and evidence of saving faith, not its ground}
  \breakrow{(11)} \(\pvar{TestimonyArguments.SolaFide}{24}\)\gloss{Scripture does not contradict itself}
  \breakrow{(12)} \(\pvar{TestimonyArguments.SolaFide}{34}\)\gloss{Justification and sanctification are inseparable but distinct}
  \breakrow{(13)} \((\pvar{TestimonyArguments.SolaFide}{9} \land \pvar{TestimonyArguments.SolaFide}{10}) \rightarrow \pvar{TestimonyArguments.SolaFide}{26}\)\gloss{\emph{would give:} Faith is sufficient: whoever believes is saved}
  \breakrow{(14)} \((\pvar{TestimonyArguments.SolaFide}{20} \land \pvar{TestimonyArguments.SolaFide}{21}) \rightarrow \pvar{TestimonyArguments.SolaFide}{22}\)\gloss{\emph{would give:} James 2:24 is compatible with Paul, using ‘justify’ in a different sense}
  \breakrow{(15)} \((\pvar{TestimonyArguments.SolaFide}{1} \land \pvar{TestimonyArguments.SolaFide}{4} \land \pvar{TestimonyArguments.SolaFide}{5}) \rightarrow \pvar{TestimonyArguments.SolaFide}{27}\)\gloss{\emph{would give:} Salvation is by grace, a gift and not wages owed}
  \breakrow{(16)} \((\pvar{TestimonyArguments.SolaFide}{1} \land \pvar{TestimonyArguments.SolaFide}{4} \land \pvar{TestimonyArguments.SolaFide}{5} \land \pvar{TestimonyArguments.SolaFide}{22} \land \pvar{TestimonyArguments.SolaFide}{24}) \rightarrow \pvar{TestimonyArguments.SolaFide}{28}\)\gloss{\emph{would give:} Salvation is not by works: no work is its ground}
  \breakrow{(17)} \(\pvar{TestimonyArguments.SolaFide}{25} \rightarrow \pvar{TestimonyArguments.SolaFide}{26}\)\gloss{\emph{would give:} Faith is sufficient: whoever believes is saved}
  \item[(18)] \((\pvar{TestimonyArguments.SolaFide}{26} \land \pvar{TestimonyArguments.SolaFide}{1}) \rightarrow \pvar{TestimonyArguments.SolaFide}{29}\)\gloss{\emph{would give:} Salvation is received through faith}
  \item[\(\vdash\)] \(\pvar{TestimonyArguments.SolaFide}{27} \land \pvar{TestimonyArguments.SolaFide}{28} \land \pvar{TestimonyArguments.SolaFide}{29}\)\gloss{\emph{at issue:} \(\pvar{TestimonyArguments.SolaFide}{27}\) Salvation is by grace, a gift and not wages owed; \(\pvar{TestimonyArguments.SolaFide}{28}\) Salvation is not by works: no work is its ground; \(\pvar{TestimonyArguments.SolaFide}{29}\) Salvation is received through faith}
  \keepwithnext
\end{derivation}
\emph{No premise here rests on scripture alone.}
\par\medskip\noindent\textbf{\texttt{johannine\allowbreak{}Party}} --- Sola fide from John (6:28–29)\par\nobreak\noindent
\textbf{John’s case}: sola fide from Jesus’ answer in John 6:28–29, argued alone. It holds the Reformed distinction between justification and sanctification, as Paul’s and Peter’s cases do, so Trent meets it on the same ground. Two inferences are rated: Calvin’s step from John’s “eternal life through believing” to justification, and the step to faith alone that Trent’s canon 9 denies.

\begin{derivation}
  \item[(1)] \(\pvar{TestimonyArguments.SolaFide}{36}\)\gloss{John 6:28–29 — the work God requires is that you believe in him whom he sent}
  \item[(2)] \(\pvar{TestimonyArguments.SolaFide}{37}\)\gloss{John 3:16–18, 3:36, 5:24, 20:31 — eternal life through believing in the Son}
  \item[(3)] \(\pvar{TestimonyArguments.SolaFide}{40}\)\gloss{The believing of John 6:29 is trust, which brings nothing and receives Christ}
  \breakrow{(4)} \(\pvar{TestimonyArguments.SolaFide}{1}\)\gloss{Ephesians 2:8–9 teaches salvation by grace through faith, not of works}
  \breakrow{(5)} \(\pvar{TestimonyArguments.SolaFide}{2}\)\gloss{Romans 3:28 teaches justification by faith apart from works of the law}
  \breakrow{(6)} \(\pvar{TestimonyArguments.SolaFide}{3}\)\gloss{Galatians 2:16 teaches that no one is justified by works of the law}
  \breakrow{(7)} \(\pvar{TestimonyArguments.SolaFide}{4}\)\gloss{Romans 4:4–5 contrasts wages owed with a gift reckoned to the one who believes}
  \breakrow{(8)} \(\pvar{TestimonyArguments.SolaFide}{5}\)\gloss{Titus 3:5 teaches that God saved us not by works done in righteousness}
  \breakrow{(9)} \(\pvar{TestimonyArguments.SolaFide}{17}\)\gloss{Luke 7:50 — Jesus says ‘your faith has saved you’ after declaring sins forgiven}
  \breakrow{(10)} \(\pvar{TestimonyArguments.SolaFide}{9}\)\gloss{Acts 15:9–11 — Peter: hearts cleansed by faith; saved through grace}
  \breakrow{(11)} \(\pvar{TestimonyArguments.SolaFide}{20}\)\gloss{James 2:14–26 targets a barren faith — mere assent — not Paul’s doctrine}
  \breakrow{(12)} \(\pvar{TestimonyArguments.SolaFide}{21}\)\gloss{Good works are the fruit and evidence of saving faith, not its ground}
  \breakrow{(13)} \(\pvar{TestimonyArguments.SolaFide}{24}\)\gloss{Scripture does not contradict itself}
  \breakrow{(14)} \(\pvar{TestimonyArguments.SolaFide}{34}\)\gloss{Justification and sanctification are inseparable but distinct}
  \breakrow{(15)} \((\pvar{TestimonyArguments.SolaFide}{36} \land \pvar{TestimonyArguments.SolaFide}{37} \land \pvar{TestimonyArguments.SolaFide}{40}) \rightarrow \pvar{TestimonyArguments.SolaFide}{25}\)\gloss{\emph{would give:} Justification is by faith alone}
  \breakrow{(16)} \((\pvar{TestimonyArguments.SolaFide}{20} \land \pvar{TestimonyArguments.SolaFide}{21}) \rightarrow \pvar{TestimonyArguments.SolaFide}{22}\)\gloss{\emph{would give:} James 2:24 is compatible with Paul, using ‘justify’ in a different sense}
  \breakrow{(17)} \((\pvar{TestimonyArguments.SolaFide}{1} \land \pvar{TestimonyArguments.SolaFide}{4} \land \pvar{TestimonyArguments.SolaFide}{5}) \rightarrow \pvar{TestimonyArguments.SolaFide}{27}\)\gloss{\emph{would give:} Salvation is by grace, a gift and not wages owed}
  \breakrow{(18)} \((\pvar{TestimonyArguments.SolaFide}{1} \land \pvar{TestimonyArguments.SolaFide}{4} \land \pvar{TestimonyArguments.SolaFide}{5} \land \pvar{TestimonyArguments.SolaFide}{22} \land \pvar{TestimonyArguments.SolaFide}{24}) \rightarrow \pvar{TestimonyArguments.SolaFide}{28}\)\gloss{\emph{would give:} Salvation is not by works: no work is its ground}
  \breakrow{(19)} \(\pvar{TestimonyArguments.SolaFide}{25} \rightarrow \pvar{TestimonyArguments.SolaFide}{26}\)\gloss{\emph{would give:} Faith is sufficient: whoever believes is saved}
  \item[(20)] \((\pvar{TestimonyArguments.SolaFide}{26} \land \pvar{TestimonyArguments.SolaFide}{1}) \rightarrow \pvar{TestimonyArguments.SolaFide}{29}\)\gloss{\emph{would give:} Salvation is received through faith}
  \item[\(\vdash\)] \(\pvar{TestimonyArguments.SolaFide}{27} \land \pvar{TestimonyArguments.SolaFide}{28} \land \pvar{TestimonyArguments.SolaFide}{29}\)\gloss{\emph{at issue:} \(\pvar{TestimonyArguments.SolaFide}{27}\) Salvation is by grace, a gift and not wages owed; \(\pvar{TestimonyArguments.SolaFide}{28}\) Salvation is not by works: no work is its ground; \(\pvar{TestimonyArguments.SolaFide}{29}\) Salvation is received through faith}
  \keepwithnext
\end{derivation}
\emph{Grounded in scripture alone:} John 3:16–18, 3:36, 5:24, 20:31 — eternal life through believing in the Son.
\par\medskip\noindent\textbf{\texttt{aquinas\allowbreak{}Case}} --- Faith formed by charity (Aquinas on John 6:29)\par\nobreak\noindent
\textbf{Aquinas on John 6:29}, as the argument he makes: the believing Jesus names is faith made alive by love, so it is not bare trust.

\begin{derivation}
  \item[(1)] \(\pvar{TestimonyArguments.SolaFide}{41}\)\gloss{The believing of John 6:29 is faith living through charity, the source of works}
  \item[(2)] \(\pvar{TestimonyArguments.SolaFide}{41} \rightarrow \lnot \pvar{TestimonyArguments.SolaFide}{40}\)
  \item[\(\vdash\)] \(\lnot \pvar{TestimonyArguments.SolaFide}{40}\)\gloss{\emph{at issue, the denial of:} The believing of John 6:29 is trust, which brings nothing and receives Christ}
  \keepwithnext
\end{derivation}
\emph{No premise here rests on scripture alone.}
\par\medskip\noindent\textbf{\texttt{sanctification\allowbreak{}Case}} --- Obedience as sanctification (John 3:36, Matthew 7:21)\par\nobreak\noindent
\textbf{The Reformed answer to Chrysostom}, as a position: the obedience John 3:36 and Matthew 7:21 require is sanctification, the fruit of faith, so the texts add no condition to faith, and faith suffices. It holds the Reformed distinction between justification and sanctification, as the other Reformed positions do, so Trent meets it on the same ground.

\begin{derivation}
  \item[(1)] \(\pvar{TestimonyArguments.SolaFide}{38}\)\gloss{John 3:36 — whoever believes has life; whoever does not obey shall not see life}
  \item[(2)] \(\pvar{TestimonyArguments.SolaFide}{62}\)\gloss{Matthew 7:21–23 — the criterion at the judgment is doing the Father’s will, not profession or mighty works}
  \item[(3)] \(\pvar{TestimonyArguments.SolaFide}{39}\)\gloss{The obedience of John 3:36 and Matthew 7:21 is sanctification, the fruit of faith}
  \breakrow{(4)} \(\pvar{TestimonyArguments.SolaFide}{34}\)\gloss{Justification and sanctification are inseparable but distinct}
  \item[(5)] \((\pvar{TestimonyArguments.SolaFide}{38} \land \pvar{TestimonyArguments.SolaFide}{62} \land \pvar{TestimonyArguments.SolaFide}{39}) \rightarrow \pvar{TestimonyArguments.SolaFide}{26}\)
  \item[\(\vdash\)] \(\pvar{TestimonyArguments.SolaFide}{26}\)\gloss{\emph{at issue:} Faith is sufficient: whoever believes is saved}
  \keepwithnext
\end{derivation}
\emph{No premise here rests on scripture alone.}
\par\medskip\noindent\textbf{\texttt{apocalyptic\allowbreak{}Case}} --- Apocalyptic reading, against faith as the condition\par\nobreak\noindent
The apocalyptic reading, as the argument it makes: faith is not the condition of justification. It keeps grace and “not by works”, by a step of its own.

\begin{derivation}
  \item[(1)] \(\lnot \pvar{TestimonyArguments.SolaFide}{7}\)\gloss{\emph{denied:} In Galatians 2:16 ‘faith of Christ’ (πίστις Χριστοῦ) means our faith in Christ}
  \item[(2)] \(\pvar{TestimonyArguments.SolaFide}{16}\)\gloss{‘The righteousness of God’ (δικαιοσύνη θεοῦ) names God’s act of rescue in Christ}
  \item[(3)] \(\pvar{TestimonyArguments.SolaFide}{3}\)\gloss{Galatians 2:16 teaches that no one is justified by works of the law}
  \breakrow{(4)} \(\pvar{TestimonyArguments.SolaFide}{8}\)\gloss{Galatians opposes requiring circumcision, besides faith, for justification}
  \breakrow{(5)} \((\lnot \pvar{TestimonyArguments.SolaFide}{7} \land \pvar{TestimonyArguments.SolaFide}{16}) \rightarrow \lnot \pvar{TestimonyArguments.SolaFide}{25}\)
  \breakrow{(6)} \(\pvar{TestimonyArguments.SolaFide}{16} \rightarrow \pvar{TestimonyArguments.SolaFide}{27} \land \pvar{TestimonyArguments.SolaFide}{28}\)\gloss{\emph{would give:} \(\pvar{TestimonyArguments.SolaFide}{27}\) Salvation is by grace, a gift and not wages owed; \(\pvar{TestimonyArguments.SolaFide}{28}\) Salvation is not by works: no work is its ground}
  \item[(7)] \((\pvar{TestimonyArguments.SolaFide}{26} \land \pvar{TestimonyArguments.SolaFide}{1}) \rightarrow \pvar{TestimonyArguments.SolaFide}{29}\)\gloss{\emph{would give:} Salvation is received through faith}
  \item[\(\vdash\)] \(\lnot \pvar{TestimonyArguments.SolaFide}{25}\)\gloss{\emph{at issue, the denial of:} Justification is by faith alone}
  \keepwithnext
\end{derivation}
\emph{No premise here rests on scripture alone.}
\par\medskip\noindent\textbf{\texttt{sanders\allowbreak{}Case}} --- Covenantal nomism (Sanders, Dunn)\par\nobreak\noindent
Sanders and Dunn: Judaism in Paul’s day taught entry by grace and staying in by obedience, so “works of the law” are the boundary markers.

\begin{derivation}
  \item[(1)] \(\pvar{TestimonyArguments.SolaFide}{31}\)\gloss{Second Temple Judaism: in by grace, staying in by works}
  \item[(2)] \(\pvar{TestimonyArguments.SolaFide}{31} \rightarrow \lnot \pvar{TestimonyArguments.SolaFide}{6}\)
  \item[\(\vdash\)] \(\lnot \pvar{TestimonyArguments.SolaFide}{6}\)\gloss{\emph{at issue, the denial of:} Paul’s ‘works of the law’ means human works in general, not boundary markers}
  \keepwithnext
\end{derivation}
\emph{No premise here rests on scripture alone.}
\par\medskip\noindent\textbf{\texttt{critics\allowbreak{}Case}} --- Variegated nomism (Gathercole, Carson et al.)\par\nobreak\noindent
Gathercole, and the volume edited by Carson, O’Brien and Seifrid: Judaism in Paul’s day was not simply that, so “works of the law” are works in general.

\begin{derivation}
  \item[(1)] \(\lnot \pvar{TestimonyArguments.SolaFide}{31}\)\gloss{\emph{denied:} Second Temple Judaism: in by grace, staying in by works}
  \item[(2)] \(\pvar{TestimonyArguments.SolaFide}{8}\)\gloss{Galatians opposes requiring circumcision, besides faith, for justification}
  \item[(3)] \((\lnot \pvar{TestimonyArguments.SolaFide}{31} \land \pvar{TestimonyArguments.SolaFide}{8}) \rightarrow \pvar{TestimonyArguments.SolaFide}{6}\)
  \item[\(\vdash\)] \(\pvar{TestimonyArguments.SolaFide}{6}\)\gloss{\emph{at issue:} Paul’s ‘works of the law’ means human works in general, not boundary markers}
  \keepwithnext
\end{derivation}
\emph{No premise here rests on scripture alone.}
\par\medskip\noindent\textbf{\texttt{jervell\allowbreak{}Case}} --- Law-observant Luke (Jervell)\par\nobreak\noindent
Jervell: Luke shows the church keeping the law, so the yoke is not the law as a condition of salvation.

\begin{derivation}
  \item[(1)] \(\pvar{TestimonyArguments.SolaFide}{11}\)\gloss{Acts 15:20–21, 21:20–24: Jewish believers keep the law; gentiles keep part}
  \item[(2)] \(\pvar{TestimonyArguments.SolaFide}{11} \rightarrow \lnot \pvar{TestimonyArguments.SolaFide}{10}\)
  \item[\(\vdash\)] \(\lnot \pvar{TestimonyArguments.SolaFide}{10}\)\gloss{\emph{at issue, the denial of:} The yoke refused at Acts 15:10 is the whole law as a condition of salvation}
  \keepwithnext
\end{derivation}
\emph{No premise here rests on scripture alone.}
\paragraph*{Each position holds together and reaches its conclusion}
\leavevmode\par\nobreak
\par\medskip\noindent\textbf{\texttt{pauline\allowbreak{}Case\_\allowbreak{}establishes}}\par\nobreak\noindent
Paul’s case reaches sola fide.

{\small\begin{verbatim}
theorem paulineCase_establishes : Establishes paulineCase
-- axioms: propext, Quot.sound
\end{verbatim}}
\par\medskip\noindent\textbf{\texttt{pauline\allowbreak{}Case\_\allowbreak{}is\_\allowbreak{}satisfiable}}\par\nobreak\noindent
Paul’s case can be held without contradiction.

{\small\begin{verbatim}
theorem paulineCase_is_satisfiable : Satisfiable paulineCase.premises
-- axioms: propext, Quot.sound
\end{verbatim}}
\par\medskip\noindent\textbf{\texttt{dominical\allowbreak{}Case\_\allowbreak{}establishes}}\par\nobreak\noindent
Luke’s case reaches its conclusion: faith suffices.

{\small\begin{verbatim}
theorem dominicalCase_establishes : Establishes dominicalCase
-- axioms: propext, Quot.sound
\end{verbatim}}
\par\medskip\noindent\textbf{\texttt{dominical\allowbreak{}Case\_\allowbreak{}is\_\allowbreak{}satisfiable}}\par\nobreak\noindent
Luke’s case can be held without contradiction.

{\small\begin{verbatim}
theorem dominicalCase_is_satisfiable : Satisfiable dominicalCase.premises
-- axioms: propext, Quot.sound
\end{verbatim}}
\par\medskip\noindent\textbf{\texttt{apostolic\allowbreak{}Case\_\allowbreak{}establishes}}\par\nobreak\noindent
Peter’s case reaches sola fide.

{\small\begin{verbatim}
theorem apostolicCase_establishes : Establishes apostolicCase
-- axioms: propext, Quot.sound
\end{verbatim}}
\par\medskip\noindent\textbf{\texttt{apostolic\allowbreak{}Case\_\allowbreak{}is\_\allowbreak{}satisfiable}}\par\nobreak\noindent
Peter’s case can be held without contradiction.

{\small\begin{verbatim}
theorem apostolicCase_is_satisfiable : Satisfiable apostolicCase.premises
-- axioms: propext, Quot.sound
\end{verbatim}}
\par\medskip\noindent\textbf{\texttt{tridentine\allowbreak{}Case\_\allowbreak{}establishes}}\par\nobreak\noindent
Trent’s case reaches its conclusion: salvation is not apart from works.

{\small\begin{verbatim}
theorem tridentineCase_establishes : Establishes tridentineCase
-- axioms: propext, Quot.sound
\end{verbatim}}
\par\medskip\noindent\textbf{\texttt{tridentine\allowbreak{}Case\_\allowbreak{}is\_\allowbreak{}satisfiable}}\par\nobreak\noindent
Trent’s case can be held without contradiction.

{\small\begin{verbatim}
theorem tridentineCase_is_satisfiable : Satisfiable tridentineCase.premises
-- axioms: propext, Quot.sound
\end{verbatim}}
\paragraph*{Where Trent’s objection comes from}
\leavevmode\par\nobreak
Trent’s objection to “not by works” is not a separate thesis about works. It follows from what Trent says justification \emph{is}: “not remission of sins merely, but also the sanctification and renewal of the inward man” (ch. 7). Grant that. Grant also that the renewal grows as the justified do good works (ch. 10). Then good works increase justification (canon 24), which denies “not by works” (ch. 16). The Reformed grant the growth, which Westminster XIII.1 calls sanctification, and deny the definition.

\par\medskip\noindent\textbf{\texttt{tridentine\allowbreak{}Case\allowbreak{}On\allowbreak{}The\allowbreak{}Reformed\allowbreak{}Distinction}} --- Trent, on the Reformed distinction\par\nobreak\noindent
Trent’s case with one change: its definition of justification is exchanged for the Reformed distinction. Everything else Trent holds is kept, including that the renewal grows through good works.

\begin{derivation}
  \item[(1)] \(\pvar{TestimonyArguments.SolaFide}{1}\)\gloss{Ephesians 2:8–9 teaches salvation by grace through faith, not of works}
  \item[(2)] \(\pvar{TestimonyArguments.SolaFide}{2}\)\gloss{Romans 3:28 teaches justification by faith apart from works of the law}
  \item[(3)] \(\pvar{TestimonyArguments.SolaFide}{3}\)\gloss{Galatians 2:16 teaches that no one is justified by works of the law}
  \breakrow{(4)} \(\pvar{TestimonyArguments.SolaFide}{5}\)\gloss{Titus 3:5 teaches that God saved us not by works done in righteousness}
  \breakrow{(5)} \(\pvar{TestimonyArguments.SolaFide}{24}\)\gloss{Scripture does not contradict itself}
  \breakrow{(6)} \(\pvar{TestimonyArguments.SolaFide}{35}\)\gloss{The inward renewal of the justified grows as they do good works in grace}
  \breakrow{(7)} \((\pvar{TestimonyArguments.SolaFide}{33} \land \pvar{TestimonyArguments.SolaFide}{35}) \rightarrow \pvar{TestimonyArguments.SolaFide}{32}\)\gloss{\emph{would give:} Works performed in grace merit an increase of justification}
  \breakrow{(8)} \(\pvar{TestimonyArguments.SolaFide}{33} \rightarrow \lnot \pvar{TestimonyArguments.SolaFide}{30}\)\gloss{\emph{would deny:} Justification is a verdict only: pardon and righteousness counted, not poured in}
  \breakrow{(9)} \(\pvar{TestimonyArguments.SolaFide}{33} \rightarrow \lnot \pvar{TestimonyArguments.SolaFide}{34}\)\gloss{\emph{would deny:} Justification and sanctification are inseparable but distinct}
  \breakrow{(10)} \(\pvar{TestimonyArguments.SolaFide}{33} \rightarrow \lnot \pvar{TestimonyArguments.SolaFide}{26}\)\gloss{\emph{would deny:} Faith is sufficient: whoever believes is saved}
  \breakrow{(11)} \(\pvar{TestimonyArguments.SolaFide}{32} \rightarrow \lnot \pvar{TestimonyArguments.SolaFide}{28}\)
  \item[(12)] \(\pvar{TestimonyArguments.SolaFide}{34}\)\gloss{Justification and sanctification are inseparable but distinct}
  \item[\(\vdash\)] \(\lnot \pvar{TestimonyArguments.SolaFide}{28}\)\gloss{\emph{at issue, the denial of:} Salvation is not by works: no work is its ground}
  \keepwithnext
\end{derivation}
\emph{No premise here rests on scripture alone.}
\par\medskip\noindent\textbf{\texttt{reformed\allowbreak{}Distinction\allowbreak{}Reading}}\par\nobreak\noindent
Trent’s position with the Reformed distinction in place of its definition. The renewal grows, but it is sanctification, so it increases nothing called justification, and salvation is not by works.

{\small\begin{verbatim}
def reformedDistinctionReading : Valuation Claim :=
  fun a =>
    match a with
    | Claim.justificationIncludesSanctification => False
    | Claim.worksMeritIncreaseOfJustification => False
    | x => True
\end{verbatim}}
\par\medskip\noindent\textbf{\texttt{trent\_\allowbreak{}objection\_\allowbreak{}rests\_\allowbreak{}on\_\allowbreak{}its\_\allowbreak{}definition}}\par\nobreak\noindent
\textbf{Trent’s objection depends on its definition of justification.} Keep all of Trent’s case except the definition, put the Reformed distinction in its place, and “salvation is not apart from works” no longer follows. Both sides still hold that the renewal of the justified grows through good works.

So the disagreement over works comes down to what justification \emph{is}: whether the renewal that grows is part of justification, or is the sanctification that follows it.

{\small\begin{verbatim}
theorem trent_objection_rests_on_its_definition : ¬Establishes
    tridentineCaseOnTheReformedDistinction
-- axioms: propext, Quot.sound
\end{verbatim}}
\par\medskip\noindent\textbf{\texttt{apocalyptic\allowbreak{}Case\_\allowbreak{}establishes}}\par\nobreak\noindent
The apocalyptic case reaches its conclusion.

{\small\begin{verbatim}
theorem apocalypticCase_establishes : Establishes apocalypticCase
-- axioms: propext, Quot.sound
\end{verbatim}}
\par\medskip\noindent\textbf{\texttt{apocalyptic\allowbreak{}Case\_\allowbreak{}is\_\allowbreak{}satisfiable}}\par\nobreak\noindent
The apocalyptic case can be held without contradiction.

{\small\begin{verbatim}
theorem apocalypticCase_is_satisfiable : Satisfiable apocalypticCase.premises
-- axioms: propext, Quot.sound
\end{verbatim}}
\par\medskip\noindent\textbf{\texttt{sanders\allowbreak{}Case\_\allowbreak{}establishes}}\par\nobreak\noindent
Sanders’ case reaches its conclusion.

{\small\begin{verbatim}
theorem sandersCase_establishes : Establishes sandersCase
-- axioms: propext, Quot.sound
\end{verbatim}}
\par\medskip\noindent\textbf{\texttt{sanders\allowbreak{}Case\_\allowbreak{}is\_\allowbreak{}satisfiable}}\par\nobreak\noindent
Sanders’ case can be held without contradiction.

{\small\begin{verbatim}
theorem sandersCase_is_satisfiable : Satisfiable sandersCase.premises
-- axioms: propext, Quot.sound
\end{verbatim}}
\par\medskip\noindent\textbf{\texttt{critics\allowbreak{}Case\_\allowbreak{}establishes}}\par\nobreak\noindent
The critics’ case reaches its conclusion.

{\small\begin{verbatim}
theorem criticsCase_establishes : Establishes criticsCase
-- axioms: propext, Quot.sound
\end{verbatim}}
\par\medskip\noindent\textbf{\texttt{critics\allowbreak{}Case\_\allowbreak{}is\_\allowbreak{}satisfiable}}\par\nobreak\noindent
The critics’ case can be held without contradiction.

{\small\begin{verbatim}
theorem criticsCase_is_satisfiable : Satisfiable criticsCase.premises
-- axioms: propext, Quot.sound
\end{verbatim}}
\par\medskip\noindent\textbf{\texttt{jervell\allowbreak{}Case\_\allowbreak{}establishes}}\par\nobreak\noindent
Jervell’s case reaches its conclusion.

{\small\begin{verbatim}
theorem jervellCase_establishes : Establishes jervellCase
-- axioms: propext, Quot.sound
\end{verbatim}}
\par\medskip\noindent\textbf{\texttt{jervell\allowbreak{}Case\_\allowbreak{}is\_\allowbreak{}satisfiable}}\par\nobreak\noindent
Jervell’s case can be held without contradiction.

{\small\begin{verbatim}
theorem jervellCase_is_satisfiable : Satisfiable jervellCase.premises
-- axioms: propext, Quot.sound
\end{verbatim}}
\par\medskip\noindent\textbf{\texttt{johannine\allowbreak{}Party\_\allowbreak{}establishes}}\par\nobreak\noindent
John’s case reaches sola fide.

{\small\begin{verbatim}
theorem johannineParty_establishes : Establishes johannineParty
-- axioms: propext, Quot.sound
\end{verbatim}}
\par\medskip\noindent\textbf{\texttt{aquinas\allowbreak{}Case\_\allowbreak{}establishes}}\par\nobreak\noindent
Aquinas’s case reaches its conclusion: the believing is not bare trust.

{\small\begin{verbatim}
theorem aquinasCase_establishes : Establishes aquinasCase
-- axioms: propext, Quot.sound
\end{verbatim}}
\par\medskip\noindent\textbf{\texttt{johannine\allowbreak{}Party\_\allowbreak{}is\_\allowbreak{}satisfiable}}\par\nobreak\noindent
John’s case can be held without contradiction: on the Reformed reading, every claim in it holds.

{\small\begin{verbatim}
theorem johannineParty_is_satisfiable : Satisfiable johannineParty.premises
-- axioms: propext, Quot.sound
\end{verbatim}}
\par\medskip\noindent\textbf{\texttt{aquinas\allowbreak{}Case\_\allowbreak{}is\_\allowbreak{}satisfiable}}\par\nobreak\noindent
Aquinas’s case can be held without contradiction: on his own reading, the believing of John 6:29 is faith formed by charity and not bare trust.

{\small\begin{verbatim}
theorem aquinasCase_is_satisfiable : Satisfiable aquinasCase.premises
-- axioms: propext, Quot.sound
\end{verbatim}}
\par\medskip\noindent\textbf{\texttt{chrysostom\allowbreak{}Case\_\allowbreak{}establishes}}\par\nobreak\noindent
Chrysostom’s case reaches its conclusion, as the John section shows.

{\small\begin{verbatim}
theorem chrysostomCase_establishes : Establishes chrysostomOnJohn
-- axioms: propext, Quot.sound
\end{verbatim}}
\par\medskip\noindent\textbf{\texttt{sanctification\allowbreak{}Case\_\allowbreak{}establishes}}\par\nobreak\noindent
The Reformed answer to Chrysostom reaches its conclusion: faith suffices.

{\small\begin{verbatim}
theorem sanctificationCase_establishes : Establishes sanctificationCase
-- axioms: propext, Quot.sound
\end{verbatim}}
\par\medskip\noindent\textbf{\texttt{sanctification\allowbreak{}Case\_\allowbreak{}is\_\allowbreak{}satisfiable}}\par\nobreak\noindent
The Reformed answer to Chrysostom can be held without contradiction: on the Reformed reading, every claim in it holds.

{\small\begin{verbatim}
theorem sanctificationCase_is_satisfiable : Satisfiable
    sanctificationCase.premises
-- axioms: propext, Quot.sound
\end{verbatim}}
\par\medskip\noindent\textbf{\texttt{chrysostom\allowbreak{}Reading}}\par\nobreak\noindent
Chrysostom’s position, with every verse of Paul’s that the Reformed strands read. Every text holds. Believing does not suffice without a right life, so faith is not sufficient and justification is not by faith alone. With Rome, it takes justification to include sanctification, and the readings of Matthew 7 and 19 that make obedience a condition. On the readings the meanings exclude, it takes the Greek of Revelation 22:11 and Sirach 18:22, and that James does not use “justify” as Paul does.

{\small\begin{verbatim}
def chrysostomReading : Valuation Claim :=
  fun a =>
    match a with
    | Claim.paulsJustifyDenotesRenewal => False
    | Claim.rev22_11BeJustifiedStill => False
    | Claim.sir18_22BeJustifiedToDeath => False
    | Claim.jamesUsesDikaioAsPaul => False
    | Claim.salvationNotByWorks => False
    | Claim.justificationIsForensicOnly => False
    | Claim.justificationDistinctFromSanctification => False
    | Claim.faithIsSufficient => False
    | Claim.justificationByFaithAlone => False
    | Claim.matthew7DoingIncludesBelieving => False
    | Claim.matthew19LawExposesInability => False
    | x => True
\end{verbatim}}
\par\medskip\noindent\textbf{\texttt{chrysostom\_\allowbreak{}grants\_\allowbreak{}the\_\allowbreak{}pauline\_\allowbreak{}texts}}\par\nobreak\noindent
\textbf{Chrysostom contradicts no text of Paul’s.} His position can be held together with every verse the Reformed strands read from Paul: Romans 3:28, Galatians 2:16, Romans 4:4–5, Ephesians 2:8–9 and Titus 3:5.

Where Chrysostom stands against the Reformed case, he stands against a reading of those verses, the step from them to faith alone, and not against anything they say. It does not show that his reading of John 3:36 is right.

{\small\begin{verbatim}
theorem chrysostom_grants_the_pauline_texts : Grants chrysostomOnJohn (⋀[p
    Claim.romans3_28, p Claim.galatians2_16, p Claim.romans4_4_5, p
    Claim.ephesians2_8_9, p Claim.titus3_5])
-- axioms: propext, Quot.sound
\end{verbatim}}
\paragraph*{How strong each position is}
\leavevmode\par\nobreak
Every position but two rests on something rated \texttt{disputed}, or denies a claim, which ranks at the bottom. Between two such positions no rating blocks an attack.

The exceptions are the two arguments about Paul’s word, from usage and from Romans 4. Every link of each is rated \texttt{well\allowbreak{}Supported} or better. So an attack on either from a position at the bottom does not defeat it, and its attack on such a position does.

\par\medskip\noindent\textbf{\texttt{pauline\allowbreak{}Case\_\allowbreak{}strength}}\par\nobreak\noindent
Paul’s case rests on its two disputed readings, and on the Reformed distinction. All are rated \texttt{disputed}.

{\small\begin{verbatim}
theorem paulineCase_strength : paulineCase.strength = 0
-- axioms: propext
\end{verbatim}}
\par\medskip\noindent\textbf{\texttt{dominical\allowbreak{}Case\_\allowbreak{}strength}}\par\nobreak\noindent
Luke’s case uses a step that Trent’s canon 9 denies, so its weakest link is \texttt{disputed}. Of its own two premises, the sense of “saved” (σῴζω) is rated \texttt{well\allowbreak{}Supported}. The reading of 7:47 is \texttt{disputed}, because Marshall names the commentators who read her love as the cause of her forgiveness.

{\small\begin{verbatim}
theorem dominicalCase_strength : dominicalCase.strength = 0
-- axioms: propext
\end{verbatim}}
\par\medskip\noindent\textbf{\texttt{apostolic\allowbreak{}Case\_\allowbreak{}strength}}\par\nobreak\noindent
Peter’s case rests on the reading of the yoke, rated \texttt{disputed}.

{\small\begin{verbatim}
theorem apostolicCase_strength : apostolicCase.strength = 0
-- axioms: propext
\end{verbatim}}
\par\medskip\noindent\textbf{\texttt{tridentine\allowbreak{}Case\_\allowbreak{}strength}}\par\nobreak\noindent
Trent rests on its definition of justification, rated \texttt{disputed}.

{\small\begin{verbatim}
theorem tridentineCase_strength : tridentineCase.strength = 0
-- axioms: propext
\end{verbatim}}
\par\medskip\noindent\textbf{\texttt{apocalyptic\allowbreak{}Case\_\allowbreak{}strength}}\par\nobreak\noindent
The apocalyptic case denies that Galatians 2:16 speaks of faith \emph{in} Christ, and a denial ranks at the bottom.

{\small\begin{verbatim}
theorem apocalypticCase_strength : apocalypticCase.strength = 0
-- axioms: propext
\end{verbatim}}
\par\medskip\noindent\textbf{\texttt{sanders\allowbreak{}Case\_\allowbreak{}strength}}\par\nobreak\noindent
Sanders’ case rests on his account of Judaism, rated \texttt{disputed}.

{\small\begin{verbatim}
theorem sandersCase_strength : sandersCase.strength = 0
-- axioms: propext
\end{verbatim}}
\par\medskip\noindent\textbf{\texttt{critics\allowbreak{}Case\_\allowbreak{}strength}}\par\nobreak\noindent
The critics’ case denies that account.

{\small\begin{verbatim}
theorem criticsCase_strength : criticsCase.strength = 0
-- axioms: propext
\end{verbatim}}
\par\medskip\noindent\textbf{\texttt{jervell\allowbreak{}Case\_\allowbreak{}strength}}\par\nobreak\noindent
Jervell’s step is contested, by Bruce.

{\small\begin{verbatim}
theorem jervellCase_strength : jervellCase.strength = 0
-- axioms: propext
\end{verbatim}}
\par\medskip\noindent\textbf{\texttt{galatian\allowbreak{}Gospel\_\allowbreak{}strength}}\par\nobreak\noindent
Paul’s gospel rests on its reading of Galatians, and on the step from a verdict to “renewal is no part of justification”. Both are rated \texttt{disputed}.

{\small\begin{verbatim}
theorem galatianGospel_strength : galatianGospel.strength = 0
-- axioms: propext
\end{verbatim}}
\par\medskip\noindent\textbf{\texttt{johannine\allowbreak{}Party\_\allowbreak{}strength}}\par\nobreak\noindent
John’s case rests on the Reformed distinction and on a step that Trent’s canon 9 denies, both \texttt{disputed}. Its own reading, believing as trust, is \texttt{disputed} too: Aquinas denies it.

{\small\begin{verbatim}
theorem johannineParty_strength : johannineParty.strength = 0
-- axioms: propext
\end{verbatim}}
\par\medskip\noindent\textbf{\texttt{aquinas\allowbreak{}Case\_\allowbreak{}strength}}\par\nobreak\noindent
Aquinas’s case rests on his reading of John 6:29, rated \texttt{disputed}: Calvin denies it.

{\small\begin{verbatim}
theorem aquinasCase_strength : aquinasCase.strength = 0
-- axioms: propext
\end{verbatim}}
\par\medskip\noindent\textbf{\texttt{luther\allowbreak{}On\allowbreak{}Galatians\_\allowbreak{}strength}}\par\nobreak\noindent
Luther’s answer rests on his step, which Aquinas denies.

{\small\begin{verbatim}
theorem lutherOnGalatians_strength : lutherOnGalatians.strength = 0
-- axioms: propext
\end{verbatim}}
\par\medskip\noindent\textbf{\texttt{chrysostom\allowbreak{}Case\_\allowbreak{}strength}}\par\nobreak\noindent
Chrysostom’s case rests on his step, rated \texttt{disputed}: Calvin accepts both texts and denies it.

{\small\begin{verbatim}
theorem chrysostomCase_strength : chrysostomOnJohn.strength = 0
-- axioms: propext
\end{verbatim}}
\par\medskip\noindent\textbf{\texttt{sanctification\allowbreak{}Case\_\allowbreak{}strength}}\par\nobreak\noindent
The Reformed answer rests on its reading of the obedience, rated \texttt{disputed}: Chrysostom denies it.

{\small\begin{verbatim}
theorem sanctificationCase_strength : sanctificationCase.strength = 0
-- axioms: propext
\end{verbatim}}
\par\medskip\noindent\textbf{\texttt{lexical\allowbreak{}Case\_\allowbreak{}strength}}\par\nobreak\noindent
\textbf{The argument from Paul’s word rests on nothing disputed.} That the verb is a courtroom word, and the rule of least meaning, are \texttt{well\allowbreak{}Supported}. That Paul has other words for renewal is \texttt{consensus}. The step from them is \texttt{well\allowbreak{}Supported}. So its weakest link ranks above the bottom.

{\small\begin{verbatim}
theorem lexicalCase_strength : lexicalCase.strength = 2
-- axioms: propext
\end{verbatim}}
\par\medskip\noindent\textbf{\texttt{tridentine\allowbreak{}On\allowbreak{}Pauls\allowbreak{}Word\_\allowbreak{}strength}}\par\nobreak\noindent
Trent, taken as a claim about Paul’s word, rests on everything Trent rests on.

{\small\begin{verbatim}
theorem tridentineOnPaulsWord_strength : tridentineOnPaulsWord.strength = 0
-- axioms: propext
\end{verbatim}}
\par\medskip\noindent\textbf{\texttt{romans\allowbreak{}Four\allowbreak{}On\allowbreak{}The\allowbreak{}Word\_\allowbreak{}strength}}\par\nobreak\noindent
\textbf{Nor does Romans 4, read for Paul’s word.} Its texts are rated \texttt{well\allowbreak{}Supported} or better. So are the sense of “count” (λογίζομαι), Paul’s explanation, the rule of least meaning, and both steps.

{\small\begin{verbatim}
theorem romansFourOnTheWord_strength : romansFourOnTheWord.strength = 2
-- axioms: propext
\end{verbatim}}
\paragraph*{The defeats}
\leavevmode\par\nobreak
Each is worked out from the two positions’ premises and checked by machine: first whether one contradicts the other, then which is the weaker.

\par\medskip\noindent\textbf{\texttt{trent\_\allowbreak{}defeats\_\allowbreak{}pauline}}\par\nobreak\noindent
\textbf{Trent defeats Paul’s case.} Trent’s case leads to the denial of “not by works”, which contradicts the conclusion of Paul’s.

{\small\begin{verbatim}
theorem trent_defeats_pauline : Defeats tridentineCase paulineCase
-- axioms: propext, Classical.choice, Quot.sound
\end{verbatim}}
\par\medskip\noindent\textbf{\texttt{trent\_\allowbreak{}defeats\_\allowbreak{}dominical}}\par\nobreak\noindent
\textbf{Trent defeats Luke’s case.} Trent’s definition, with canon 9, denies that faith suffices, which is what Luke’s case concludes.

{\small\begin{verbatim}
theorem trent_defeats_dominical : Defeats tridentineCase dominicalCase
-- axioms: propext, Classical.choice, Quot.sound
\end{verbatim}}
\par\medskip\noindent\textbf{\texttt{trent\_\allowbreak{}defeats\_\allowbreak{}apostolic}}\par\nobreak\noindent
\textbf{Trent defeats Peter’s case}, as it defeats Paul’s.

{\small\begin{verbatim}
theorem trent_defeats_apostolic : Defeats tridentineCase apostolicCase
-- axioms: propext, Classical.choice, Quot.sound
\end{verbatim}}
\par\medskip\noindent\textbf{\texttt{pauline\_\allowbreak{}defeats\_\allowbreak{}trent}}\par\nobreak\noindent
\textbf{Paul’s case defeats Trent in return.} It leads to “not by works”, which Trent denies.

{\small\begin{verbatim}
theorem pauline_defeats_trent : Defeats paulineCase tridentineCase
-- axioms: propext, Classical.choice, Quot.sound
\end{verbatim}}
\par\medskip\noindent\textbf{\texttt{dominical\_\allowbreak{}defeats\_\allowbreak{}trent}}\par\nobreak\noindent
\textbf{Luke’s case defeats Trent in return.}

{\small\begin{verbatim}
theorem dominical_defeats_trent : Defeats dominicalCase tridentineCase
-- axioms: propext, Classical.choice, Quot.sound
\end{verbatim}}
\par\medskip\noindent\textbf{\texttt{apostolic\_\allowbreak{}defeats\_\allowbreak{}trent}}\par\nobreak\noindent
\textbf{Peter’s case defeats Trent in return.}

{\small\begin{verbatim}
theorem apostolic_defeats_trent : Defeats apostolicCase tridentineCase
-- axioms: propext, Classical.choice, Quot.sound
\end{verbatim}}
\par\medskip\noindent\textbf{\texttt{apocalyptic\_\allowbreak{}defeats\_\allowbreak{}trent}}\par\nobreak\noindent
\textbf{The apocalyptic reading defeats Trent.} It holds “not by works” by a step of its own, since God’s rescue depends on nothing a person does, and Trent denies it.

{\small\begin{verbatim}
theorem apocalyptic_defeats_trent : Defeats apocalypticCase tridentineCase
-- axioms: propext, Classical.choice, Quot.sound
\end{verbatim}}
\par\medskip\noindent\textbf{\texttt{apocalyptic\_\allowbreak{}defeats\_\allowbreak{}pauline}}\par\nobreak\noindent
\textbf{The apocalyptic reading defeats Paul’s case.} It denies that Galatians 2:16 speaks of faith \emph{in} Christ, a premise of Paul’s strand rated \texttt{disputed}.

{\small\begin{verbatim}
theorem apocalyptic_defeats_pauline : Defeats apocalypticCase paulineCase
-- axioms: propext, Classical.choice, Quot.sound
\end{verbatim}}
\par\medskip\noindent\textbf{\texttt{pauline\_\allowbreak{}defeats\_\allowbreak{}apocalyptic}}\par\nobreak\noindent
\textbf{Paul’s case defeats the apocalyptic reading in return.} It leads to justification by faith alone.

{\small\begin{verbatim}
theorem pauline_defeats_apocalyptic : Defeats paulineCase apocalypticCase
-- axioms: propext, Classical.choice, Quot.sound
\end{verbatim}}
\par\medskip\noindent\textbf{\texttt{dominical\_\allowbreak{}does\_\allowbreak{}not\_\allowbreak{}defeat\_\allowbreak{}apocalyptic}}\par\nobreak\noindent
\textbf{Luke’s case does not defeat the apocalyptic reading.} Luke 7:50 says that faith saved her, not that faith alone does. The apocalyptic reading grants that faith saves, and denies only that faith is the condition. So nothing in Luke’s case contradicts it (\texttt{luke\_\allowbreak{}grants\_\allowbreak{}the\_\allowbreak{}apocalyptic\_\allowbreak{}conclusion}).

{\small\begin{verbatim}
theorem dominical_does_not_defeat_apocalyptic : ¬Defeats dominicalCase
    apocalypticCase
-- axioms: propext, Classical.choice, Quot.sound
\end{verbatim}}
\par\medskip\noindent\textbf{\texttt{apostolic\_\allowbreak{}does\_\allowbreak{}not\_\allowbreak{}defeat\_\allowbreak{}apocalyptic}}\par\nobreak\noindent
\textbf{Nor does Peter’s case}, for the same reason: Peter says that faith suffices, not that faith alone does (\texttt{peter\_\allowbreak{}grants\_\allowbreak{}the\_\allowbreak{}apocalyptic\_\allowbreak{}conclusion}).

{\small\begin{verbatim}
theorem apostolic_does_not_defeat_apocalyptic : ¬Defeats apostolicCase
    apocalypticCase
-- axioms: propext, Classical.choice, Quot.sound
\end{verbatim}}
\par\medskip\noindent\textbf{\texttt{sanders\_\allowbreak{}defeats\_\allowbreak{}pauline}}\par\nobreak\noindent
\textbf{Sanders defeats Paul’s case.} His account of Judaism, with Dunn’s inference, denies the reading of “works of the law” as works in general.

{\small\begin{verbatim}
theorem sanders_defeats_pauline : Defeats sandersCase paulineCase
-- axioms: propext, Classical.choice, Quot.sound
\end{verbatim}}
\par\medskip\noindent\textbf{\texttt{pauline\_\allowbreak{}defeats\_\allowbreak{}sanders}}\par\nobreak\noindent
\textbf{Paul’s case defeats Sanders in return.} It holds the premise Sanders’ conclusion denies.

{\small\begin{verbatim}
theorem pauline_defeats_sanders : Defeats paulineCase sandersCase
-- axioms: propext, Classical.choice, Quot.sound
\end{verbatim}}
\par\medskip\noindent\textbf{\texttt{sanders\_\allowbreak{}defeats\_\allowbreak{}critics}}\par\nobreak\noindent
\textbf{Sanders defeats the critics.} Their conclusions contradict.

{\small\begin{verbatim}
theorem sanders_defeats_critics : Defeats sandersCase criticsCase
-- axioms: propext, Classical.choice, Quot.sound
\end{verbatim}}
\par\medskip\noindent\textbf{\texttt{critics\_\allowbreak{}defeat\_\allowbreak{}sanders}}\par\nobreak\noindent
\textbf{The critics defeat Sanders.} They deny the account of Judaism he starts from.

{\small\begin{verbatim}
theorem critics_defeat_sanders : Defeats criticsCase sandersCase
-- axioms: propext, Classical.choice, Quot.sound
\end{verbatim}}
\par\medskip\noindent\textbf{\texttt{jervell\_\allowbreak{}defeats\_\allowbreak{}apostolic}}\par\nobreak\noindent
\textbf{Jervell defeats Peter’s case.} He denies its reading of the yoke.

{\small\begin{verbatim}
theorem jervell_defeats_apostolic : Defeats jervellCase apostolicCase
-- axioms: propext, Classical.choice, Quot.sound
\end{verbatim}}
\par\medskip\noindent\textbf{\texttt{apostolic\_\allowbreak{}defeats\_\allowbreak{}jervell}}\par\nobreak\noindent
\textbf{Peter’s case defeats Jervell in return.} It holds the premise his conclusion denies.

{\small\begin{verbatim}
theorem apostolic_defeats_jervell : Defeats apostolicCase jervellCase
-- axioms: propext, Classical.choice, Quot.sound
\end{verbatim}}
\paragraph*{What does not defeat}
\leavevmode\par\nobreak
Most pairs of positions do not conflict: some consistent reading holds both. Seven such readings are named below, each one a way of holding a whole group of positions at once. Two pairs conflict in one direction only: each argument about Paul’s word defeats Trent taken as a claim about it, and is not defeated in return. All of this, like every answer in the table, is worked out from the premises.

\par\medskip\noindent\textbf{\texttt{critics\allowbreak{}Jervell\allowbreak{}Reading}}\par\nobreak\noindent
The critics’ reading together with Jervell’s: Sanders’ account of Judaism is denied, and the yoke is Israel’s law for gentiles. Paul’s case and Luke’s hold on it. Faith is the only means and justification is a verdict, so it denies Trent’s account of baptism as the instrument and of justification as including sanctification.

{\small\begin{verbatim}
def criticsJervellReading : Valuation Claim :=
  fun a =>
    match a with
    | Claim.paulsJustifyDenotesRenewal => False
    | Claim.rev22_11BeJustifiedStill => False
    | Claim.sir18_22BeJustifiedToDeath => False
    | Claim.jamesUsesDikaioAsPaul => False
    | Claim.secondTempleCovenantalNomism => False
    | Claim.acts15YokeIsLawAsCondition => False
    | Claim.baptismIsInstrumentalCause => False
    | Claim.justificationIncludesSanctification => False
    | x => True
\end{verbatim}}
\par\medskip\noindent\textbf{\texttt{trent\allowbreak{}Sanders\allowbreak{}Jervell\allowbreak{}Reading}}\par\nobreak\noindent
Trent’s position together with Sanders’ and Jervell’s. With Trent, it denies that faith is the only means of justification.

{\small\begin{verbatim}
def trentSandersJervellReading : Valuation Claim :=
  fun a =>
    match a with
    | Claim.paulsJustifyDenotesRenewal => False
    | Claim.rev22_11BeJustifiedStill => False
    | Claim.sir18_22BeJustifiedToDeath => False
    | Claim.jamesUsesDikaioAsPaul => False
    | Claim.salvationNotByWorks => False
    | Claim.justificationIsForensicOnly => False
    | Claim.worksOfLawMeansWorksGenerally => False
    | Claim.acts15YokeIsLawAsCondition => False
    | Claim.justificationDistinctFromSanctification => False
    | Claim.faithIsSufficient => False
    | Claim.justificationByFaithAlone => False
    | x => True
\end{verbatim}}
\par\medskip\noindent\textbf{\texttt{trent\allowbreak{}Critics\allowbreak{}Jervell\allowbreak{}Reading}}\par\nobreak\noindent
Trent’s position together with the critics’ and Jervell’s. With Trent, it denies that faith is the only means of justification.

{\small\begin{verbatim}
def trentCriticsJervellReading : Valuation Claim :=
  fun a =>
    match a with
    | Claim.paulsJustifyDenotesRenewal => False
    | Claim.rev22_11BeJustifiedStill => False
    | Claim.sir18_22BeJustifiedToDeath => False
    | Claim.jamesUsesDikaioAsPaul => False
    | Claim.salvationNotByWorks => False
    | Claim.justificationIsForensicOnly => False
    | Claim.secondTempleCovenantalNomism => False
    | Claim.acts15YokeIsLawAsCondition => False
    | Claim.justificationDistinctFromSanctification => False
    | Claim.faithIsSufficient => False
    | Claim.justificationByFaithAlone => False
    | x => True
\end{verbatim}}
\par\medskip\noindent\textbf{\texttt{apocalyptic\allowbreak{}Sanders\allowbreak{}Jervell\allowbreak{}Reading}}\par\nobreak\noindent
The apocalyptic reading together with Sanders’ and Jervell’s. Justification is a verdict only here, as on every reading but Trent’s, so it does not include sanctification.

{\small\begin{verbatim}
def apocalypticSandersJervellReading : Valuation Claim :=
  fun a =>
    match a with
    | Claim.paulsJustifyDenotesRenewal => False
    | Claim.rev22_11BeJustifiedStill => False
    | Claim.sir18_22BeJustifiedToDeath => False
    | Claim.jamesUsesDikaioAsPaul => False
    | Claim.pistisChristouObjective => False
    | Claim.justificationByFaithAlone => False
    | Claim.worksOfLawMeansWorksGenerally => False
    | Claim.acts15YokeIsLawAsCondition => False
    | Claim.justificationIncludesSanctification => False
    | x => True
\end{verbatim}}
\par\medskip\noindent\textbf{\texttt{apocalyptic\allowbreak{}Critics\allowbreak{}Jervell\allowbreak{}Reading}}\par\nobreak\noindent
The apocalyptic reading together with the critics’ and Jervell’s. Justification is a verdict only here, as on every reading but Trent’s, so it does not include sanctification.

{\small\begin{verbatim}
def apocalypticCriticsJervellReading : Valuation Claim :=
  fun a =>
    match a with
    | Claim.paulsJustifyDenotesRenewal => False
    | Claim.rev22_11BeJustifiedStill => False
    | Claim.sir18_22BeJustifiedToDeath => False
    | Claim.jamesUsesDikaioAsPaul => False
    | Claim.pistisChristouObjective => False
    | Claim.justificationByFaithAlone => False
    | Claim.secondTempleCovenantalNomism => False
    | Claim.acts15YokeIsLawAsCondition => False
    | Claim.justificationIncludesSanctification => False
    | x => True
\end{verbatim}}
\par\medskip\noindent\textbf{\texttt{apocalyptic\_\allowbreak{}does\_\allowbreak{}not\_\allowbreak{}defeat\_\allowbreak{}dominical}}\par\nobreak\noindent
\textbf{The apocalyptic reading does not defeat Luke’s case.} It contradicts nothing Luke’s case rests on: not Luke 7:50, not the sense of “saved”, not the reading of 7:47. And it does not deny the conclusion. It can grant that salvation is received through faith while denying that faith is its condition. The reasons, premise by premise, are in \texttt{apocalyptic\_\allowbreak{}grants\_\allowbreak{}sola\_\allowbreak{}fide\_\allowbreak{}as\_\allowbreak{}stated} and \texttt{apocalyptic\_\allowbreak{}grants\_\allowbreak{}the\_\allowbreak{}dominical\_\allowbreak{}step}.

{\small\begin{verbatim}
theorem apocalyptic_does_not_defeat_dominical : ¬Defeats apocalypticCase
    dominicalCase
-- axioms: propext, Classical.choice, Quot.sound
\end{verbatim}}
\par\medskip\noindent\textbf{\texttt{apocalyptic\_\allowbreak{}does\_\allowbreak{}not\_\allowbreak{}defeat\_\allowbreak{}apostolic}}\par\nobreak\noindent
\textbf{Nor Peter’s case}, for the same reasons (\texttt{apocalyptic\_\allowbreak{}grants\_\allowbreak{}the\_\allowbreak{}apostolic\_\allowbreak{}step}).

{\small\begin{verbatim}
theorem apocalyptic_does_not_defeat_apostolic : ¬Defeats apocalypticCase
    apostolicCase
-- axioms: propext, Classical.choice, Quot.sound
\end{verbatim}}
\par\medskip\noindent\textbf{\texttt{trent\_\allowbreak{}defeats\_\allowbreak{}apocalyptic}}\par\nobreak\noindent
\textbf{Trent defeats the apocalyptic reading.} The two agree on its conclusion: justification is not by faith alone. But the apocalyptic reading \emph{derives} “not by works” on the way, since God’s rescue is a gift that depends on nothing a person does. Trent denies that: works done in grace merit an increase of justification. So Trent contradicts something the apocalyptic reading derives.

{\small\begin{verbatim}
theorem trent_defeats_apocalyptic : Defeats tridentineCase apocalypticCase
-- axioms: propext, Classical.choice, Quot.sound
\end{verbatim}}
\par\medskip\noindent\textbf{\texttt{gospel\_\allowbreak{}defeats\_\allowbreak{}trent}}\par\nobreak\noindent
\textbf{Paul’s gospel defeats Trent.} It leads to the denial of Trent’s definition: justification is not the renewal of the inward person. \texttt{where\allowbreak{}Trent\allowbreak{}Parts\allowbreak{}From\allowbreak{}Paul} says at which step.

{\small\begin{verbatim}
theorem gospel_defeats_trent : Defeats galatianGospel tridentineCase
-- axioms: propext, Classical.choice, Quot.sound
\end{verbatim}}
\par\medskip\noindent\textbf{\texttt{trent\_\allowbreak{}defeats\_\allowbreak{}gospel}}\par\nobreak\noindent
\textbf{Trent defeats Paul’s gospel in return}, by the same contradiction from the other side. Its definition denies the gospel line’s conclusion, and both rest on something rated \texttt{disputed}.

{\small\begin{verbatim}
theorem trent_defeats_gospel : Defeats tridentineCase galatianGospel
-- axioms: propext, Classical.choice, Quot.sound
\end{verbatim}}
\par\medskip\noindent\textbf{\texttt{lexical\_\allowbreak{}defeats\_\allowbreak{}trent\_\allowbreak{}on\_\allowbreak{}pauls\_\allowbreak{}word}}\par\nobreak\noindent
\textbf{Paul’s word defeats Trent taken as a claim about it.} Taken that way, Trent derives that “justify” (δικαιόω) means renewal. The argument from the word concludes that it does not. The attack succeeds because Trent’s weakest link is not above the argument’s.

{\small\begin{verbatim}
theorem lexical_defeats_trent_on_pauls_word : Defeats lexicalCase
    tridentineOnPaulsWord
-- axioms: propext, Classical.choice, Quot.sound
\end{verbatim}}
\par\medskip\noindent\textbf{\texttt{trent\_\allowbreak{}on\_\allowbreak{}pauls\_\allowbreak{}word\_\allowbreak{}does\_\allowbreak{}not\_\allowbreak{}defeat\_\allowbreak{}lexical}}\par\nobreak\noindent
\textbf{And the reply fails.} Trent, taken as a claim about Paul’s word, attacks the argument from the word in turn. But Trent’s weakest link is its definition, rated \texttt{disputed}, and the argument’s is \texttt{well\allowbreak{}Supported}. A weaker argument’s contradiction does not defeat a stronger one.

{\small\begin{verbatim}
theorem trent_on_pauls_word_does_not_defeat_lexical : ¬Defeats
    tridentineOnPaulsWord lexicalCase
-- axioms: propext, Classical.choice, Quot.sound
\end{verbatim}}
\paragraph*{Why these do not conflict}
\leavevmode\par\nobreak
The absences above are worked out by machine, and a machine’s working gives no reason. The readings below do. Each is a way of holding one position that leaves a premise or conclusion of another standing, and each result checks that the reading holds everything the first position holds.

\par\medskip\noindent\textbf{\texttt{apocalyptic\allowbreak{}With\allowbreak{}Faith\allowbreak{}Reading}}\par\nobreak\noindent
The apocalyptic reading with salvation still received through faith: faith is not the condition, but it is the means. Grace and “not by works” hold. Justification is a verdict only here, as on every reading but Trent’s, so it does not include sanctification.

{\small\begin{verbatim}
def apocalypticWithFaithReading : Valuation Claim :=
  fun a =>
    match a with
    | Claim.paulsJustifyDenotesRenewal => False
    | Claim.rev22_11BeJustifiedStill => False
    | Claim.sir18_22BeJustifiedToDeath => False
    | Claim.jamesUsesDikaioAsPaul => False
    | Claim.pistisChristouObjective => False
    | Claim.justificationByFaithAlone => False
    | Claim.justificationIncludesSanctification => False
    | x => True
\end{verbatim}}
\par\medskip\noindent\textbf{\texttt{apocalyptic\_\allowbreak{}grants\_\allowbreak{}sola\_\allowbreak{}fide\_\allowbreak{}as\_\allowbreak{}stated}}\par\nobreak\noindent
\textbf{The apocalyptic reading can grant what the Reformed strands conclude.} Grace, “not by works”, and salvation received through faith can all hold on it. What it denies is that faith is the condition. So it contradicts neither Luke’s case nor Peter’s.

{\small\begin{verbatim}
theorem apocalyptic_grants_sola_fide_as_stated : Grants apocalypticCase
    dominicalCase.conclusion
-- axioms: propext, Quot.sound
\end{verbatim}}
\par\medskip\noindent\textbf{\texttt{apocalyptic\_\allowbreak{}grants\_\allowbreak{}the\_\allowbreak{}dominical\_\allowbreak{}step}}\par\nobreak\noindent
\textbf{The apocalyptic reader can grant Luke’s step.} The step delivers that faith suffices, “your faith has saved you”, and the apocalyptic reading can hold that: faith is the means, though not the condition.

{\small\begin{verbatim}
theorem apocalyptic_grants_the_dominical_step : Grants apocalypticCase
    dominicalToSufficiency
-- axioms: propext, Quot.sound
\end{verbatim}}
\par\medskip\noindent\textbf{\texttt{apocalyptic\_\allowbreak{}grants\_\allowbreak{}the\_\allowbreak{}apostolic\_\allowbreak{}step}}\par\nobreak\noindent
\textbf{And Peter’s step}, for the same reason: Peter says that faith suffices, and the apocalyptic reading does not deny it.

{\small\begin{verbatim}
theorem apocalyptic_grants_the_apostolic_step : Grants apocalypticCase
    apostolicToSufficiency
-- axioms: propext, Quot.sound
\end{verbatim}}
\par\medskip\noindent\textbf{\texttt{luke\allowbreak{}Without\allowbreak{}Alone\allowbreak{}Reading}}\par\nobreak\noindent
Luke’s case without \emph{alone}: every premise of the case, faith sufficient, and justification not by faith alone, because nothing Luke says excludes another condition. Justification is a verdict only here, as on every reading but Trent’s, so it does not include sanctification.

{\small\begin{verbatim}
def lukeWithoutAloneReading : Valuation Claim :=
  fun a =>
    match a with
    | Claim.paulsJustifyDenotesRenewal => False
    | Claim.rev22_11BeJustifiedStill => False
    | Claim.sir18_22BeJustifiedToDeath => False
    | Claim.jamesUsesDikaioAsPaul => False
    | Claim.justificationByFaithAlone => False
    | Claim.justificationIncludesSanctification => False
    | x => True
\end{verbatim}}
\par\medskip\noindent\textbf{\texttt{luke\_\allowbreak{}grants\_\allowbreak{}the\_\allowbreak{}apocalyptic\_\allowbreak{}conclusion}}\par\nobreak\noindent
\textbf{Luke’s case can grant the apocalyptic conclusion.} Everything Luke’s case holds is consistent with justification not being by faith alone, so it does not contradict the apocalyptic reading. Among the cases from Paul, Luke and Acts, only Paul’s, which argues for \emph{alone}, does.

{\small\begin{verbatim}
theorem luke_grants_the_apocalyptic_conclusion : Grants dominicalCase
    apocalypticCase.conclusion
-- axioms: propext, Quot.sound
\end{verbatim}}
\par\medskip\noindent\textbf{\texttt{peter\_\allowbreak{}grants\_\allowbreak{}the\_\allowbreak{}apocalyptic\_\allowbreak{}conclusion}}\par\nobreak\noindent
\textbf{And so can Peter’s case.}

{\small\begin{verbatim}
theorem peter_grants_the_apocalyptic_conclusion : Grants apostolicCase
    apocalypticCase.conclusion
-- axioms: propext, Quot.sound
\end{verbatim}}
\par\medskip\noindent\textbf{\texttt{trent\allowbreak{}Without\allowbreak{}Deliverance\allowbreak{}Reading}}\par\nobreak\noindent
Trent’s position, with both of these denied: that Galatians 2:16 speaks of faith \emph{in} Christ, and that “the righteousness of God” is God’s act of rescue. With Trent, it denies that faith is the only means of justification.

{\small\begin{verbatim}
def trentWithoutDeliveranceReading : Valuation Claim :=
  fun a =>
    match a with
    | Claim.paulsJustifyDenotesRenewal => False
    | Claim.rev22_11BeJustifiedStill => False
    | Claim.sir18_22BeJustifiedToDeath => False
    | Claim.jamesUsesDikaioAsPaul => False
    | Claim.salvationNotByWorks => False
    | Claim.justificationIsForensicOnly => False
    | Claim.pistisChristouObjective => False
    | Claim.righteousnessOfGodIsDeliverance => False
    | Claim.justificationDistinctFromSanctification => False
    | Claim.faithIsSufficient => False
    | Claim.justificationByFaithAlone => False
    | x => True
\end{verbatim}}
\par\medskip\noindent\textbf{\texttt{trent\_\allowbreak{}grants\_\allowbreak{}the\_\allowbreak{}subjective\_\allowbreak{}genitive}}\par\nobreak\noindent
\textbf{Trent can grant Hays’s reading of “faith of Christ”.} Nothing Trent holds settles how the phrase (πίστις Χριστοῦ) is read, so the apocalyptic reading’s denial that it means faith \emph{in} Christ is not something Trent answers.

{\small\begin{verbatim}
theorem trent_grants_the_subjective_genitive : Grants tridentineCase (notP
    Claim.pistisChristouObjective)
-- axioms: propext, Quot.sound
\end{verbatim}}
\par\medskip\noindent\textbf{\texttt{trent\allowbreak{}Without\allowbreak{}Faith\allowbreak{}Alone\allowbreak{}Reading}}\par\nobreak\noindent
Trent’s position, on which justification is not by faith alone. That is the apocalyptic reading’s conclusion, which Trent shares.

{\small\begin{verbatim}
def trentWithoutFaithAloneReading : Valuation Claim :=
  fun a =>
    match a with
    | Claim.paulsJustifyDenotesRenewal => False
    | Claim.rev22_11BeJustifiedStill => False
    | Claim.sir18_22BeJustifiedToDeath => False
    | Claim.jamesUsesDikaioAsPaul => False
    | Claim.salvationNotByWorks => False
    | Claim.justificationIsForensicOnly => False
    | Claim.justificationByFaithAlone => False
    | Claim.justificationDistinctFromSanctification => False
    | Claim.faithIsSufficient => False
    | x => True
\end{verbatim}}
\par\medskip\noindent\textbf{\texttt{trent\_\allowbreak{}grants\_\allowbreak{}the\_\allowbreak{}apocalyptic\_\allowbreak{}conclusion}}\par\nobreak\noindent
\textbf{Trent grants the apocalyptic reading’s conclusion}: justification is not by faith alone. The two hold it for different reasons, so Trent does not contradict that conclusion.

{\small\begin{verbatim}
theorem trent_grants_the_apocalyptic_conclusion : Grants tridentineCase
    apocalypticCase.conclusion
-- axioms: propext, Quot.sound
\end{verbatim}}
\paragraph*{The dispute}
\leavevmode\par\nobreak
\par\medskip\noindent\textbf{\texttt{party\allowbreak{}Node}}\par\nobreak\noindent
The premises each position argues from.

{\small\begin{verbatim}
def partyNode : Party → ArgumentPackage Claim
\end{verbatim}}
\par\medskip\noindent\textbf{\texttt{sola\allowbreak{}Fide\allowbreak{}Dispute}}\par\nobreak\noindent
The debate over sola fide. For every position heard: its premises can be held without contradiction, its conclusion follows from them, and every step it uses is rated.

{\small\begin{verbatim}
def solaFideDispute : Dispute Claim Party :=
  { node := partyNode, consistent := solaFideDispute._proof_1,
    sound := solaFideDispute._proof_2,
    rated := solaFideDispute._proof_20 }
\end{verbatim}}
\par\medskip\noindent\textbf{\texttt{sola\allowbreak{}Fide\allowbreak{}Dispute\_\allowbreak{}defeats}}\par\nobreak\noindent
\textbf{Who defeats whom}, for all 289 pairs of positions.

Trent and each Reformed strand defeat each other. The apocalyptic reading and Paul’s case defeat each other; it does not conflict with Luke’s case or Peter’s; it and Trent defeat each other. Sanders defeats Paul’s case and the critics’, and both defeat him in return. Jervell and Peter’s case defeat each other. Paul’s gospel and Trent defeat each other, over Trent’s definition. John’s case meets Trent, the apocalyptic reading and Aquinas, each in both directions, and Aquinas and Luther defeat each other. Trent taken as a claim about Paul’s word meets every position Trent meets, in both directions, and two more: the two arguments about Paul’s word defeat it, and it defeats neither in return. Nothing else.

Every one of the 225 answers, the defeats and the absences alike, is worked out from the positions’ premises and checked by machine. The table above is what that working is checked against. None of it is assumed.

{\small\begin{verbatim}
theorem solaFideDispute_defeats : ∀ (i j : Party),
 solaFideDispute.defeats i
    j ↔
 (i, j) ∈
 have y :=
 have y :=
 have y := [(Party.johannine,
    Party.chrysostom), (Party.chrysostom, Party.sanctification),
    (Party.sanctification, Party.chrysostom), (Party.trent,
    Party.sanctification), (Party.sanctification, Party.trent),
    (Party.trentOnPaulsWord, Party.sanctification), (Party.sanctification,
    Party.trentOnPaulsWord)];
 (Party.chrysostom, Party.pauline) ::
    (Party.pauline, Party.chrysostom) :: (Party.chrysostom, Party.dominical)
    :: (Party.dominical, Party.chrysostom) :: (Party.chrysostom,
    Party.apostolic) :: (Party.apostolic, Party.chrysostom) ::
    (Party.chrysostom, Party.johannine) :: y;
 have y := (Party.johannine,
    Party.trentOnPaulsWord) :: (Party.apocalyptic, Party.johannine) ::
    (Party.johannine, Party.apocalyptic) :: (Party.aquinas, Party.johannine)
    :: (Party.johannine, Party.aquinas) :: (Party.luther, Party.aquinas) ::
    (Party.aquinas, Party.luther) :: y;
 (Party.trentOnPaulsWord,
    Party.dominical) :: (Party.trentOnPaulsWord, Party.apostolic) ::
    (Party.trentOnPaulsWord, Party.apocalyptic) :: (Party.trent,
    Party.johannine) :: (Party.johannine, Party.trent) ::
    (Party.trentOnPaulsWord, Party.johannine) :: y;
 have y :=
 have y :=
    (Party.lexical, Party.trentOnPaulsWord) :: (Party.romansOnTheWord,
    Party.trentOnPaulsWord) :: (Party.pauline, Party.trentOnPaulsWord) ::
    (Party.dominical, Party.trentOnPaulsWord) :: (Party.apostolic,
    Party.trentOnPaulsWord) :: (Party.apocalyptic, Party.trentOnPaulsWord) ::
    (Party.trentOnPaulsWord, Party.pauline) :: y;
 (Party.critics,
    Party.sanders) :: (Party.jervell, Party.apostolic) :: (Party.apostolic,
    Party.jervell) :: (Party.gospel, Party.trent) :: (Party.trent,
    Party.gospel) :: (Party.gospel, Party.trentOnPaulsWord) ::
    (Party.trentOnPaulsWord, Party.gospel) :: y;
 have y :=
    (Party.apocalyptic, Party.trent) :: (Party.trent, Party.apocalyptic) ::
    (Party.apocalyptic, Party.pauline) :: (Party.pauline, Party.apocalyptic)
    :: (Party.sanders, Party.pauline) :: (Party.pauline, Party.sanders) ::
    (Party.sanders, Party.critics) :: y;
 (Party.trent, Party.pauline) ::
    (Party.trent, Party.dominical) :: (Party.trent, Party.apostolic) ::
    (Party.pauline, Party.trent) :: (Party.dominical, Party.trent) ::
    (Party.apostolic, Party.trent) :: y
-- axioms: propext, Classical.choice, Quot.sound
\end{verbatim}}
\par\medskip\noindent\textbf{\texttt{sola\allowbreak{}Fide\allowbreak{}Dispute\_\allowbreak{}supports}}\par\nobreak\noindent
\textbf{No position supports another} in this debate, in any of the 289 pairs. One position supports another when its conclusion gives the other a claim it rests on. The critics’ conclusion used to be one of Paul’s premises. Now Paul’s case contains the critics’ whole line (\texttt{pauline\allowbreak{}Case}), which is the closer relation recorded below. The Reformed strands share a conclusion; that is agreement, not support (see \texttt{Testimony.Logic.Support}). Every answer is worked out and checked by machine.

{\small\begin{verbatim}
theorem solaFideDispute_supports : ∀ (i j : Party), ¬solaFideDispute.supports
    i j
-- axioms: propext, Classical.choice, Quot.sound
\end{verbatim}}
\par\medskip\noindent\textbf{\texttt{sola\allowbreak{}Fide\allowbreak{}Dispute\_\allowbreak{}part\allowbreak{}Of}}\par\nobreak\noindent
\textbf{Whose case is contained in whose}, for all 289 pairs. The critics’ case is contained in Paul’s, since Paul’s case derives its reading of “works of the law” from the critics’ line. Trent’s case is contained in Trent taken as a claim about Paul’s word, which only adds to it. Every case contains itself. Nothing else. Every answer is worked out and checked by machine.

{\small\begin{verbatim}
theorem solaFideDispute_partOf : ∀ (i j : Party), solaFideDispute.partOf i j ↔
    i = j ∨ (i, j) ∈ [(Party.critics, Party.pauline), (Party.trent,
    Party.trentOnPaulsWord)]
-- axioms: propext, Classical.choice, Quot.sound
\end{verbatim}}
\par\medskip\noindent\textbf{\texttt{sola\allowbreak{}Fide\allowbreak{}Map}}\par\nobreak\noindent
The debate drawn as a map: who defeats whom, and whose case is contained in whose.

\begin{center}
\begin{tikzpicture}[>={Stealth[length=2mm]}]
\node[draw,circle,minimum size=6mm,inner sep=0pt] (n0) at (0.00,3.20) {1};
\node[draw,circle,minimum size=6mm,inner sep=0pt] (n1) at (1.16,2.98) {2};
\node[draw,circle,minimum size=6mm,inner sep=0pt] (n2) at (2.16,2.36) {3};
\node[draw,circle,minimum size=6mm,inner sep=0pt] (n3) at (2.86,1.43) {4};
\node[draw,circle,minimum size=6mm,inner sep=0pt] (n4) at (3.19,0.30) {5};
\node[draw,circle,minimum size=6mm,inner sep=0pt] (n5) at (3.08,-0.88) {6};
\node[draw,circle,minimum size=6mm,inner sep=0pt] (n6) at (2.55,-1.93) {7};
\node[draw,circle,minimum size=6mm,inner sep=0pt] (n7) at (1.68,-2.72) {8};
\node[draw,circle,minimum size=6mm,inner sep=0pt] (n8) at (0.59,-3.15) {9};
\node[draw,circle,minimum size=6mm,inner sep=0pt] (n9) at (-0.59,-3.15) {10};
\node[draw,circle,minimum size=6mm,inner sep=0pt] (n10) at (-1.68,-2.72) {11};
\node[draw,circle,minimum size=6mm,inner sep=0pt] (n11) at (-2.55,-1.93) {12};
\node[draw,circle,minimum size=6mm,inner sep=0pt] (n12) at (-3.08,-0.88) {13};
\node[draw,circle,minimum size=6mm,inner sep=0pt] (n13) at (-3.19,0.30) {14};
\node[draw,circle,minimum size=6mm,inner sep=0pt] (n14) at (-2.86,1.43) {15};
\node[draw,circle,minimum size=6mm,inner sep=0pt] (n15) at (-2.16,2.36) {16};
\node[draw,circle,minimum size=6mm,inner sep=0pt] (n16) at (-1.16,2.98) {17};
\draw[->,red!70!black,thick] (n0) to[bend left=10] (n3);
\draw[->,red!70!black,thick] (n0) to[bend left=10] (n4);
\draw[->,red!70!black,thick] (n0) to[bend left=10] (n5);
\draw[->,red!70!black,thick] (n0) to[bend left=10] (n10);
\draw[->,red!70!black,thick] (n0) to[bend left=10] (n15);
\draw[->,red!70!black,thick] (n1) to[bend left=10] (n3);
\draw[->,red!70!black,thick] (n1) to[bend left=10] (n10);
\draw[->,red!70!black,thick] (n1) to[bend left=10] (n15);
\draw[->,red!70!black,thick] (n2) to[bend left=10] (n3);
\draw[->,red!70!black,thick] (n2) to[bend left=10] (n7);
\draw[->,red!70!black,thick] (n2) to[bend left=10] (n10);
\draw[->,red!70!black,thick] (n2) to[bend left=10] (n15);
\draw[->,red!70!black,thick] (n3) to[bend left=10] (n0);
\draw[->,red!70!black,thick] (n3) to[bend left=10] (n1);
\draw[->,red!70!black,thick] (n3) to[bend left=10] (n2);
\draw[->,red!70!black,thick] (n3) to[bend left=10] (n4);
\draw[->,red!70!black,thick] (n3) to[bend left=10] (n8);
\draw[->,red!70!black,thick] (n3) to[bend left=10] (n12);
\draw[->,red!70!black,thick] (n3) to[bend left=10] (n16);
\draw[->,red!70!black,thick] (n4) to[bend left=10] (n0);
\draw[->,red!70!black,thick] (n4) to[bend left=10] (n3);
\draw[->,red!70!black,thick] (n4) to[bend left=10] (n10);
\draw[->,red!70!black,thick] (n4) to[bend left=10] (n12);
\draw[->,red!70!black,thick] (n5) to[bend left=10] (n0);
\draw[->,red!70!black,thick] (n5) to[bend left=10] (n6);
\draw[->,red!70!black,thick] (n6) to[bend left=10] (n5);
\draw[->,red!70!black,thick] (n7) to[bend left=10] (n2);
\draw[->,red!70!black,thick] (n8) to[bend left=10] (n3);
\draw[->,red!70!black,thick] (n8) to[bend left=10] (n10);
\draw[->,red!70!black,thick] (n9) to[bend left=10] (n10);
\draw[->,red!70!black,thick] (n10) to[bend left=10] (n0);
\draw[->,red!70!black,thick] (n10) to[bend left=10] (n1);
\draw[->,red!70!black,thick] (n10) to[bend left=10] (n2);
\draw[->,red!70!black,thick] (n10) to[bend left=10] (n4);
\draw[->,red!70!black,thick] (n10) to[bend left=10] (n8);
\draw[->,red!70!black,thick] (n10) to[bend left=10] (n12);
\draw[->,red!70!black,thick] (n10) to[bend left=10] (n16);
\draw[->,red!70!black,thick] (n11) to[bend left=10] (n10);
\draw[->,red!70!black,thick] (n12) to[bend left=10] (n3);
\draw[->,red!70!black,thick] (n12) to[bend left=10] (n4);
\draw[->,red!70!black,thick] (n12) to[bend left=10] (n10);
\draw[->,red!70!black,thick] (n12) to[bend left=10] (n13);
\draw[->,red!70!black,thick] (n12) to[bend left=10] (n15);
\draw[->,red!70!black,thick] (n13) to[bend left=10] (n12);
\draw[->,red!70!black,thick] (n13) to[bend left=10] (n14);
\draw[->,red!70!black,thick] (n14) to[bend left=10] (n13);
\draw[->,red!70!black,thick] (n15) to[bend left=10] (n0);
\draw[->,red!70!black,thick] (n15) to[bend left=10] (n1);
\draw[->,red!70!black,thick] (n15) to[bend left=10] (n2);
\draw[->,red!70!black,thick] (n15) to[bend left=10] (n12);
\draw[->,red!70!black,thick] (n15) to[bend left=10] (n16);
\draw[->,red!70!black,thick] (n16) to[bend left=10] (n3);
\draw[->,red!70!black,thick] (n16) to[bend left=10] (n10);
\draw[->,red!70!black,thick] (n16) to[bend left=10] (n15);
\draw[->,blue!60!black,thick,dotted] (n3) to[bend left=10] (n10);
\draw[->,blue!60!black,thick,dotted] (n6) to[bend left=10] (n0);
\end{tikzpicture}
\end{center}
\noindent Solid: defeats. Dashed: supports. Dotted: one party's case is part of another's.
\begin{enumerate}
\item Sola fide from Paul (Galatians 2:16)
\item Luke 7:50: “your faith has saved you”
\item Sola fide from Peter (Acts 15:9–11)
\item Trent, against ‘not by works’
\item Apocalyptic reading, against faith as the condition
\item Covenantal nomism (Sanders, Dunn)
\item Variegated nomism (Gathercole, Carson et al.)
\item Law-observant Luke (Jervell)
\item Paul’s gospel (Galatians 1), against justification as renewal
\item Paul’s word (δικαιόω, by the rule of least meaning)
\item Trent, against ‘not by works’, read as what Paul’s word means
\item Romans 4:3–8, on Paul’s word (by the rule of least meaning)
\item Sola fide from John (6:28–29)
\item Faith formed by charity (Aquinas on John 6:29)
\item Luther on Galatians 3:11–12 (against faith formed by charity)
\item Chrysostom on John 3:36 (believing is not enough without a right life)
\item Obedience as sanctification (John 3:36, Matthew 7:21)
\end{enumerate}
\begin{longtable}{@{}>{\raggedright\arraybackslash}p{0.3\textwidth} >{\raggedright\arraybackslash}p{0.3\textwidth} >{\raggedright\arraybackslash}p{0.3\textwidth}@{}}
\toprule
\textbf{From} & \textbf{To} & \textbf{Edge} \\
\midrule
\endhead
1 \emph{Sola fide from Paul (Galatians 2:16)} & 4 \emph{Trent, against ‘not by works’} & defeats \\
1 \emph{Sola fide from Paul (Galatians 2:16)} & 5 \emph{Apocalyptic reading, against faith as the condition} & defeats \\
1 \emph{Sola fide from Paul (Galatians 2:16)} & 6 \emph{Covenantal nomism (Sanders, Dunn)} & defeats \\
1 \emph{Sola fide from Paul (Galatians 2:16)} & 11 \emph{Trent, against ‘not by works’, read as what Paul’s word means} & defeats \\
1 \emph{Sola fide from Paul (Galatians 2:16)} & 16 \emph{Chrysostom on John 3:36 (believing is not enough without a right life)} & defeats \\
2 \emph{Luke 7:50: “your faith has saved you”} & 4 \emph{Trent, against ‘not by works’} & defeats \\
2 \emph{Luke 7:50: “your faith has saved you”} & 11 \emph{Trent, against ‘not by works’, read as what Paul’s word means} & defeats \\
2 \emph{Luke 7:50: “your faith has saved you”} & 16 \emph{Chrysostom on John 3:36 (believing is not enough without a right life)} & defeats \\
3 \emph{Sola fide from Peter (Acts 15:9–11)} & 4 \emph{Trent, against ‘not by works’} & defeats \\
3 \emph{Sola fide from Peter (Acts 15:9–11)} & 8 \emph{Law-observant Luke (Jervell)} & defeats \\
3 \emph{Sola fide from Peter (Acts 15:9–11)} & 11 \emph{Trent, against ‘not by works’, read as what Paul’s word means} & defeats \\
3 \emph{Sola fide from Peter (Acts 15:9–11)} & 16 \emph{Chrysostom on John 3:36 (believing is not enough without a right life)} & defeats \\
4 \emph{Trent, against ‘not by works’} & 1 \emph{Sola fide from Paul (Galatians 2:16)} & defeats \\
4 \emph{Trent, against ‘not by works’} & 2 \emph{Luke 7:50: “your faith has saved you”} & defeats \\
4 \emph{Trent, against ‘not by works’} & 3 \emph{Sola fide from Peter (Acts 15:9–11)} & defeats \\
4 \emph{Trent, against ‘not by works’} & 5 \emph{Apocalyptic reading, against faith as the condition} & defeats \\
4 \emph{Trent, against ‘not by works’} & 9 \emph{Paul’s gospel (Galatians 1), against justification as renewal} & defeats \\
4 \emph{Trent, against ‘not by works’} & 13 \emph{Sola fide from John (6:28–29)} & defeats \\
4 \emph{Trent, against ‘not by works’} & 17 \emph{Obedience as sanctification (John 3:36, Matthew 7:21)} & defeats \\
5 \emph{Apocalyptic reading, against faith as the condition} & 1 \emph{Sola fide from Paul (Galatians 2:16)} & defeats \\
5 \emph{Apocalyptic reading, against faith as the condition} & 4 \emph{Trent, against ‘not by works’} & defeats \\
5 \emph{Apocalyptic reading, against faith as the condition} & 11 \emph{Trent, against ‘not by works’, read as what Paul’s word means} & defeats \\
5 \emph{Apocalyptic reading, against faith as the condition} & 13 \emph{Sola fide from John (6:28–29)} & defeats \\
6 \emph{Covenantal nomism (Sanders, Dunn)} & 1 \emph{Sola fide from Paul (Galatians 2:16)} & defeats \\
6 \emph{Covenantal nomism (Sanders, Dunn)} & 7 \emph{Variegated nomism (Gathercole, Carson et al.)} & defeats \\
7 \emph{Variegated nomism (Gathercole, Carson et al.)} & 6 \emph{Covenantal nomism (Sanders, Dunn)} & defeats \\
8 \emph{Law-observant Luke (Jervell)} & 3 \emph{Sola fide from Peter (Acts 15:9–11)} & defeats \\
9 \emph{Paul’s gospel (Galatians 1), against justification as renewal} & 4 \emph{Trent, against ‘not by works’} & defeats \\
9 \emph{Paul’s gospel (Galatians 1), against justification as renewal} & 11 \emph{Trent, against ‘not by works’, read as what Paul’s word means} & defeats \\
10 \emph{Paul’s word (δικαιόω, by the rule of least meaning)} & 11 \emph{Trent, against ‘not by works’, read as what Paul’s word means} & defeats \\
11 \emph{Trent, against ‘not by works’, read as what Paul’s word means} & 1 \emph{Sola fide from Paul (Galatians 2:16)} & defeats \\
11 \emph{Trent, against ‘not by works’, read as what Paul’s word means} & 2 \emph{Luke 7:50: “your faith has saved you”} & defeats \\
11 \emph{Trent, against ‘not by works’, read as what Paul’s word means} & 3 \emph{Sola fide from Peter (Acts 15:9–11)} & defeats \\
11 \emph{Trent, against ‘not by works’, read as what Paul’s word means} & 5 \emph{Apocalyptic reading, against faith as the condition} & defeats \\
11 \emph{Trent, against ‘not by works’, read as what Paul’s word means} & 9 \emph{Paul’s gospel (Galatians 1), against justification as renewal} & defeats \\
11 \emph{Trent, against ‘not by works’, read as what Paul’s word means} & 13 \emph{Sola fide from John (6:28–29)} & defeats \\
11 \emph{Trent, against ‘not by works’, read as what Paul’s word means} & 17 \emph{Obedience as sanctification (John 3:36, Matthew 7:21)} & defeats \\
12 \emph{Romans 4:3–8, on Paul’s word (by the rule of least meaning)} & 11 \emph{Trent, against ‘not by works’, read as what Paul’s word means} & defeats \\
13 \emph{Sola fide from John (6:28–29)} & 4 \emph{Trent, against ‘not by works’} & defeats \\
13 \emph{Sola fide from John (6:28–29)} & 5 \emph{Apocalyptic reading, against faith as the condition} & defeats \\
13 \emph{Sola fide from John (6:28–29)} & 11 \emph{Trent, against ‘not by works’, read as what Paul’s word means} & defeats \\
13 \emph{Sola fide from John (6:28–29)} & 14 \emph{Faith formed by charity (Aquinas on John 6:29)} & defeats \\
13 \emph{Sola fide from John (6:28–29)} & 16 \emph{Chrysostom on John 3:36 (believing is not enough without a right life)} & defeats \\
14 \emph{Faith formed by charity (Aquinas on John 6:29)} & 13 \emph{Sola fide from John (6:28–29)} & defeats \\
14 \emph{Faith formed by charity (Aquinas on John 6:29)} & 15 \emph{Luther on Galatians 3:11–12 (against faith formed by charity)} & defeats \\
15 \emph{Luther on Galatians 3:11–12 (against faith formed by charity)} & 14 \emph{Faith formed by charity (Aquinas on John 6:29)} & defeats \\
16 \emph{Chrysostom on John 3:36 (believing is not enough without a right life)} & 1 \emph{Sola fide from Paul (Galatians 2:16)} & defeats \\
16 \emph{Chrysostom on John 3:36 (believing is not enough without a right life)} & 2 \emph{Luke 7:50: “your faith has saved you”} & defeats \\
16 \emph{Chrysostom on John 3:36 (believing is not enough without a right life)} & 3 \emph{Sola fide from Peter (Acts 15:9–11)} & defeats \\
16 \emph{Chrysostom on John 3:36 (believing is not enough without a right life)} & 13 \emph{Sola fide from John (6:28–29)} & defeats \\
16 \emph{Chrysostom on John 3:36 (believing is not enough without a right life)} & 17 \emph{Obedience as sanctification (John 3:36, Matthew 7:21)} & defeats \\
17 \emph{Obedience as sanctification (John 3:36, Matthew 7:21)} & 4 \emph{Trent, against ‘not by works’} & defeats \\
17 \emph{Obedience as sanctification (John 3:36, Matthew 7:21)} & 11 \emph{Trent, against ‘not by works’, read as what Paul’s word means} & defeats \\
17 \emph{Obedience as sanctification (John 3:36, Matthew 7:21)} & 16 \emph{Chrysostom on John 3:36 (believing is not enough without a right life)} & defeats \\
4 \emph{Trent, against ‘not by works’} & 11 \emph{Trent, against ‘not by works’, read as what Paul’s word means} & is part of \\
7 \emph{Variegated nomism (Gathercole, Carson et al.)} & 1 \emph{Sola fide from Paul (Galatians 2:16)} & is part of \\
1 \emph{Sola fide from Paul (Galatians 2:16)} & 11 \emph{Trent, against ‘not by works’, read as what Paul’s word means} & attack on a part — already a defeat \\
2 \emph{Luke 7:50: “your faith has saved you”} & 11 \emph{Trent, against ‘not by works’, read as what Paul’s word means} & attack on a part — already a defeat \\
3 \emph{Sola fide from Peter (Acts 15:9–11)} & 11 \emph{Trent, against ‘not by works’, read as what Paul’s word means} & attack on a part — already a defeat \\
5 \emph{Apocalyptic reading, against faith as the condition} & 11 \emph{Trent, against ‘not by works’, read as what Paul’s word means} & attack on a part — already a defeat \\
6 \emph{Covenantal nomism (Sanders, Dunn)} & 1 \emph{Sola fide from Paul (Galatians 2:16)} & attack on a part — already a defeat \\
9 \emph{Paul’s gospel (Galatians 1), against justification as renewal} & 11 \emph{Trent, against ‘not by works’, read as what Paul’s word means} & attack on a part — already a defeat \\
13 \emph{Sola fide from John (6:28–29)} & 11 \emph{Trent, against ‘not by works’, read as what Paul’s word means} & attack on a part — already a defeat \\
17 \emph{Obedience as sanctification (John 3:36, Matthew 7:21)} & 11 \emph{Trent, against ‘not by works’, read as what Paul’s word means} & attack on a part — already a defeat \\
\bottomrule
\end{longtable}
\noindent The defeats are the cells of \texttt{sola\allowbreak{}Fide\allowbreak{}Dispute\_\allowbreak{}defeats}, the supports the cells of \texttt{sola\allowbreak{}Fide\allowbreak{}Dispute\_\allowbreak{}supports} and the parts the cells of \texttt{sola\allowbreak{}Fide\allowbreak{}Dispute\_\allowbreak{}part\allowbreak{}Of}, each computed from the parties' premises and checked by the kernel. The attacks derived through support and through parts are reported, not counted: every verdict is computed from the defeats alone. Every derived attack is already a defeat.
\par\medskip\noindent\textbf{\texttt{reformed\_\allowbreak{}strands\_\allowbreak{}stand\_\allowbreak{}with\_\allowbreak{}the\_\allowbreak{}critics}}\par\nobreak\noindent
Paul’s, Luke’s and Peter’s cases can be held together with the critics’, on the reading where Sanders’ account of Judaism is denied and everything else holds.

{\small\begin{verbatim}
theorem reformed_strands_stand_with_the_critics :
    solaFideDispute.StandTogether [Party.pauline, Party.dominical,
    Party.apostolic, Party.critics]
-- axioms: propext, Classical.choice, Quot.sound
\end{verbatim}}
\par\medskip\noindent\textbf{\texttt{paul\_\allowbreak{}and\_\allowbreak{}luke\_\allowbreak{}stand\_\allowbreak{}with\_\allowbreak{}jervell}}\par\nobreak\noindent
Paul’s case and Luke’s can be held together with the critics’ and with Jervell’s.

{\small\begin{verbatim}
theorem paul_and_luke_stand_with_jervell : solaFideDispute.StandTogether
    [Party.pauline, Party.dominical, Party.jervell, Party.critics]
-- axioms: propext, Classical.choice, Quot.sound
\end{verbatim}}
\par\medskip\noindent\textbf{\texttt{luke\_\allowbreak{}and\_\allowbreak{}acts\_\allowbreak{}stand\_\allowbreak{}with\_\allowbreak{}sanders}}\par\nobreak\noindent
Luke’s case and Peter’s can be held together with Sanders’: neither rests on the reading of “works of the law”.

{\small\begin{verbatim}
theorem luke_and_acts_stand_with_sanders : solaFideDispute.StandTogether
    [Party.dominical, Party.apostolic, Party.sanders]
-- axioms: propext, Classical.choice, Quot.sound
\end{verbatim}}
\par\medskip\noindent\textbf{\texttt{trent\_\allowbreak{}stands\_\allowbreak{}with\_\allowbreak{}sanders\_\allowbreak{}and\_\allowbreak{}jervell}}\par\nobreak\noindent
Trent can be held together with Sanders and Jervell.

{\small\begin{verbatim}
theorem trent_stands_with_sanders_and_jervell : solaFideDispute.StandTogether
    [Party.trent, Party.sanders, Party.jervell]
-- axioms: propext, Classical.choice, Quot.sound
\end{verbatim}}
\par\medskip\noindent\textbf{\texttt{trent\_\allowbreak{}stands\_\allowbreak{}with\_\allowbreak{}the\_\allowbreak{}critics\_\allowbreak{}and\_\allowbreak{}jervell}}\par\nobreak\noindent
Trent can be held together with the critics and Jervell.

{\small\begin{verbatim}
theorem trent_stands_with_the_critics_and_jervell :
    solaFideDispute.StandTogether [Party.trent, Party.critics, Party.jervell]
-- axioms: propext, Classical.choice, Quot.sound
\end{verbatim}}
\par\medskip\noindent\textbf{\texttt{apocalyptic\_\allowbreak{}stands\_\allowbreak{}with\_\allowbreak{}sanders\_\allowbreak{}and\_\allowbreak{}jervell}}\par\nobreak\noindent
The apocalyptic reading can be held together with Sanders and Jervell.

{\small\begin{verbatim}
theorem apocalyptic_stands_with_sanders_and_jervell :
    solaFideDispute.StandTogether [Party.apocalyptic, Party.sanders,
    Party.jervell]
-- axioms: propext, Classical.choice, Quot.sound
\end{verbatim}}
\par\medskip\noindent\textbf{\texttt{apocalyptic\_\allowbreak{}stands\_\allowbreak{}with\_\allowbreak{}the\_\allowbreak{}critics\_\allowbreak{}and\_\allowbreak{}jervell}}\par\nobreak\noindent
The apocalyptic reading can be held together with the critics and Jervell.

{\small\begin{verbatim}
theorem apocalyptic_stands_with_the_critics_and_jervell :
    solaFideDispute.StandTogether [Party.apocalyptic, Party.critics,
    Party.jervell]
-- axioms: propext, Classical.choice, Quot.sound
\end{verbatim}}
\paragraph*{What the dispute decides}
\leavevmode\par\nobreak
\par\medskip\noindent\textbf{\texttt{only\allowbreak{}Pauls\allowbreak{}Word\allowbreak{}Unanswered}}\par\nobreak\noindent
How far the debate can be settled, and why no further. Nothing defeats either argument about Paul’s word, so both are accepted first. Nothing joins them, because every other position has an attacker that those two do not rule out. Each Reformed strand has Trent. Trent has Paul’s case. The apocalyptic reading has Paul’s case. Sanders has the critics, and the critics have Sanders. Jervell has Peter’s case. Paul’s gospel has Trent. Aquinas has Luther, and Luther has Aquinas. And Trent taken as a claim about Paul’s word has the argument from the word itself. Some of these are ties, not losses: Paul’s gospel answers Trent as Trent answers it.

\paragraph{What the dispute forces is exactly \emph{Paul’s word (δικαιόω, by the rule of least meaning)}, \emph{Romans 4:3–8, on Paul’s word (by the rule of least meaning)}.}
\begin{itemize}
\item Stage 1: \emph{Paul’s word (δικαιόω, by the rule of least meaning)}, \emph{Romans 4:3–8, on Paul’s word (by the rule of least meaning)}.
\begin{itemize}
\item \emph{Paul’s word (δικαιόω, by the rule of least meaning)} is defeated by nothing.
\item \emph{Romans 4:3–8, on Paul’s word (by the rule of least meaning)} is defeated by nothing.
\end{itemize}
\item Nothing else is forced:
\begin{itemize}
\item \emph{Trent, against ‘not by works’} defeats \emph{Sola fide from Paul (Galatians 2:16)}, and nothing forced defeats \emph{Trent, against ‘not by works’}.
\item \emph{Trent, against ‘not by works’} defeats \emph{Luke 7:50: “your faith has saved you”}, and nothing forced defeats \emph{Trent, against ‘not by works’}.
\item \emph{Trent, against ‘not by works’} defeats \emph{Sola fide from Peter (Acts 15:9–11)}, and nothing forced defeats \emph{Trent, against ‘not by works’}.
\item \emph{Sola fide from Paul (Galatians 2:16)} defeats \emph{Trent, against ‘not by works’}, and nothing forced defeats \emph{Sola fide from Paul (Galatians 2:16)}.
\item \emph{Sola fide from Paul (Galatians 2:16)} defeats \emph{Apocalyptic reading, against faith as the condition}, and nothing forced defeats \emph{Sola fide from Paul (Galatians 2:16)}.
\item \emph{Variegated nomism (Gathercole, Carson et al.)} defeats \emph{Covenantal nomism (Sanders, Dunn)}, and nothing forced defeats \emph{Variegated nomism (Gathercole, Carson et al.)}.
\item \emph{Covenantal nomism (Sanders, Dunn)} defeats \emph{Variegated nomism (Gathercole, Carson et al.)}, and nothing forced defeats \emph{Covenantal nomism (Sanders, Dunn)}.
\item \emph{Sola fide from Peter (Acts 15:9–11)} defeats \emph{Law-observant Luke (Jervell)}, and nothing forced defeats \emph{Sola fide from Peter (Acts 15:9–11)}.
\item \emph{Trent, against ‘not by works’} defeats \emph{Paul’s gospel (Galatians 1), against justification as renewal}, and nothing forced defeats \emph{Trent, against ‘not by works’}.
\item \emph{Paul’s word (δικαιόω, by the rule of least meaning)} defeats \emph{Trent, against ‘not by works’, read as what Paul’s word means} (\texttt{where\allowbreak{}Trents\allowbreak{}Word\allowbreak{}Reading\allowbreak{}Falls}), and nothing forced defeats \emph{Paul’s word (δικαιόω, by the rule of least meaning)}.
\item \emph{Trent, against ‘not by works’} defeats \emph{Sola fide from John (6:28–29)}, and nothing forced defeats \emph{Trent, against ‘not by works’}.
\item \emph{Luther on Galatians 3:11–12 (against faith formed by charity)} defeats \emph{Faith formed by charity (Aquinas on John 6:29)}, and nothing forced defeats \emph{Luther on Galatians 3:11–12 (against faith formed by charity)}.
\item \emph{Faith formed by charity (Aquinas on John 6:29)} defeats \emph{Luther on Galatians 3:11–12 (against faith formed by charity)}, and nothing forced defeats \emph{Faith formed by charity (Aquinas on John 6:29)}.
\item \emph{Sola fide from Paul (Galatians 2:16)} defeats \emph{Chrysostom on John 3:36 (believing is not enough without a right life)}, and nothing forced defeats \emph{Sola fide from Paul (Galatians 2:16)}.
\item \emph{Trent, against ‘not by works’} defeats \emph{Obedience as sanctification (John 3:36, Matthew 7:21)}, and nothing forced defeats \emph{Trent, against ‘not by works’}.
\end{itemize}
\end{itemize}
\par\noindent\textbf{What this rests on.} The reasons state 15 defeats and 48 absences of defeat, each a cell of the defeat table, \texttt{sola\allowbreak{}Fide\allowbreak{}Dispute\_\allowbreak{}defeats}, computed from the two parties' premises and checked by the kernel. Whether an attack survives turns on the attacker's weakest link and on the rating of what it attacks, so a changed rating can change the verdict. The weakest links:
\begin{itemize}
\item \emph{Sola fide from Paul (Galatians 2:16)}, weakest at \emph{disputed}:
\begin{itemize}
\item In Galatians 2:16 ‘faith of Christ’ (πίστις Χριστοῦ) means our faith in Christ --- \cite{dunn-once-more-pistis-christou-1991}; \cite{matlock-detheologizing-2000}; Gal 2:16; Rom 3:22; Phil 3:9
\item Second Temple Judaism: in by grace, staying in by works (encoded as a denial, so it ranks disputed) --- \cite{sanders-paul-palestinian-judaism-1977}
\item Good works are the fruit and evidence of saving faith, not its ground --- \cite[XI.2; XVI.2; XVI.5]{westminster-confession-1647}; \cite[III.xvi.1]{calvin-institutes-1960}; Jas 2:18-22; Eph 2:10; Rom 6:22; Luke 17:10
\item Justification and sanctification are inseparable but distinct --- \cite[Q. 77]{westminster-larger-catechism-1647}; 1 Cor 6:11; 1 Cor 1:30; Rom 4:6; Rom 4:8; Ezek 36:27; Heb 9:13-14; Rom 3:24-25; Rom 6:6; Rom 6:14; \cite[XI.1]{westminster-confession-1647}; \cite[III.xi.6]{calvin-institutes-1960}
\item an inference step --- \cite[Trent, Session VI (1547), canon 9]{tanner-decrees-1990}
\item an inference step --- \cite{gathercole-where-boasting-2002}; \cite{carson-variegated-nomism-1-2001}
\end{itemize}
\item \emph{Paul’s word (δικαιόω, by the rule of least meaning)}, weakest at \emph{well supported}:
\begin{itemize}
\item In Paul, ‘justify’ (δικαιόω) is a verdict: to declare righteous --- \cite[s.v. δικαιόω]{bauer-danker-lexicon-2000}; \cite{morris-apostolic-preaching-1965}; \cite[p. 80]{moo-romans-1996}; \cite{irons-righteousness-of-god-2015}; \cite{wright-justification-2009}; \cite[ad loc. Rom 3:24]{fitzmyer-romans-1993}
\item A word contributes the least meaning its context requires; one occurrence does not carry the whole concept --- \cite{silva-biblical-words-1994}; \cite{joos-semantic-axiom-1972}; \cite{barr-semantics-2004}
\item an inference step --- \cite{silva-biblical-words-1994}; \cite[p. 80]{moo-romans-1996}; \cite{barr-semantics-2004}
\end{itemize}
\item \emph{Luke 7:50: “your faith has saved you”}, weakest at \emph{disputed}:
\begin{itemize}
\item Luke 7:47 — her love is the evidence of her forgiveness, not its ground --- \cite[ad loc. Luke 7:47]{marshall-luke-1978}; \cite[54--54]{padilla-narrative-criticism-2021}; \cite[Apology IV, on love and the fulfilling of the law: Luke 7:47–50]{kolb-wengert-book-of-concord-2000}; Luke 7:41-43; Luke 7:47; Luke 7:47-50
\item an inference step --- \cite[Trent, Session VI (1547), canon 9]{tanner-decrees-1990}
\end{itemize}
\item \emph{Sola fide from Peter (Acts 15:9–11)}, weakest at \emph{disputed}:
\begin{itemize}
\item The yoke refused at Acts 15:10 is the whole law as a condition of salvation --- \cite[ad loc. Acts 15:10]{bruce-acts-1988}; Acts 15:10; Gal 5:1; Jas 2:10; \cite{das-paul-law-covenant-2001}
\item Good works are the fruit and evidence of saving faith, not its ground --- \cite[XI.2; XVI.2; XVI.5]{westminster-confession-1647}; \cite[III.xvi.1]{calvin-institutes-1960}; Jas 2:18-22; Eph 2:10; Rom 6:22; Luke 17:10
\item Justification and sanctification are inseparable but distinct --- \cite[Q. 77]{westminster-larger-catechism-1647}; 1 Cor 6:11; 1 Cor 1:30; Rom 4:6; Rom 4:8; Ezek 36:27; Heb 9:13-14; Rom 3:24-25; Rom 6:6; Rom 6:14; \cite[XI.1]{westminster-confession-1647}; \cite[III.xi.6]{calvin-institutes-1960}
\item an inference step --- \cite[Trent, Session VI (1547), canon 9]{tanner-decrees-1990}
\end{itemize}
\item \emph{Trent, against ‘not by works’}, weakest at \emph{disputed}:
\begin{itemize}
\item Justification is not remission of sins only, but renewal of the inward man --- \cite[Trent, Session VI (1547), Decree on Justification, ch. 7]{tanner-decrees-1990}; \cite[Trent, Session VI (1547), canon 11]{tanner-decrees-1990}
\end{itemize}
\item \emph{Apocalyptic reading, against faith as the condition}, weakest at \emph{disputed}:
\begin{itemize}
\item ‘Faith of Christ’ is Christ’s own faithfulness, the means of God’s rescue (encoded as a denial, so it ranks disputed) --- \cite{campbell-deliverance-god-2009}; \cite[ad loc. Gal 2:16]{martyn-galatians-1997}; \cite{hays-faith-jesus-christ-2002}
\item ‘The righteousness of God’ (δικαιοσύνη θεοῦ) names God’s act of rescue in Christ --- \cite{campbell-deliverance-god-2009}; \cite{martyn-galatians-1997}; Rom 3:21-26
\end{itemize}
\item \emph{Covenantal nomism (Sanders, Dunn)}, weakest at \emph{disputed}:
\begin{itemize}
\item Second Temple Judaism: in by grace, staying in by works --- \cite{sanders-paul-palestinian-judaism-1977}
\item an inference step --- \cite{dunn-new-perspective-2005}; \cite{wright-what-paul-said-1997}
\end{itemize}
\item \emph{Variegated nomism (Gathercole, Carson et al.)}, weakest at \emph{disputed}:
\begin{itemize}
\item Second Temple Judaism also held final vindication according to works (encoded as a denial, so it ranks disputed) --- \cite{gathercole-where-boasting-2002}; \cite{carson-variegated-nomism-1-2001}
\item an inference step --- \cite{gathercole-where-boasting-2002}; \cite{carson-variegated-nomism-1-2001}
\end{itemize}
\item \emph{Law-observant Luke (Jervell)}, weakest at \emph{disputed}:
\begin{itemize}
\item an inference step --- \cite{jervell-luke-people-god-1972}
\end{itemize}
\item \emph{Paul’s gospel (Galatians 1), against justification as renewal}, weakest at \emph{disputed}:
\begin{itemize}
\item an inference step --- \cite[ad loc. Gal 1:6]{moo-galatians-2013}; \cite[XI.1]{westminster-confession-1647}
\item an inference step --- \cite[XI.1]{westminster-confession-1647}; \cite[III.xi.2]{calvin-institutes-1960}; \cite{morris-apostolic-preaching-1965}
\end{itemize}
\item \emph{Trent, against ‘not by works’, read as what Paul’s word means}, weakest at \emph{disputed}:
\begin{itemize}
\item Justification is not remission of sins only, but renewal of the inward man --- \cite[Trent, Session VI (1547), Decree on Justification, ch. 7]{tanner-decrees-1990}; \cite[Trent, Session VI (1547), canon 11]{tanner-decrees-1990}
\item an inference step --- \cite[Trent, Session VI (1547), Decree on Justification, ch. 8]{tanner-decrees-1990}; \cite[Trent, Session VI (1547), Decree on Justification, ch. 7]{tanner-decrees-1990}; \cite[On the Spirit and the Letter, ch. 45 (26.45)]{npnf1-05-1887}; \cite{mcgrath-iustitia-dei-2005}
\end{itemize}
\item \emph{Romans 4:3–8, on Paul’s word (by the rule of least meaning)}, weakest at \emph{well supported}:
\begin{itemize}
\item Romans 4:4–5 contrasts wages owed with a gift reckoned to the one who believes --- Rom 4:4-5; \cite[III.xi.18]{calvin-institutes-1960}
\item ‘Count’ (λογίζομαι) in Romans 4:3–11 is the language of an account: credited to one --- \cite[s.v. λογίζομαι]{bauer-danker-lexicon-2000}; \cite[ad loc. Rom 4:3]{moo-romans-1996}
\item A word contributes the least meaning its context requires; one occurrence does not carry the whole concept --- \cite{silva-biblical-words-1994}; \cite{joos-semantic-axiom-1972}; \cite{barr-semantics-2004}
\item an inference step --- \cite[ad loc. Rom 4:6]{moo-romans-1996}; \cite[s.v. λογίζομαι]{bauer-danker-lexicon-2000}
\item an inference step --- \cite{silva-biblical-words-1994}; \cite[p. 80]{moo-romans-1996}; \cite{barr-semantics-2004}
\end{itemize}
\item \emph{Sola fide from John (6:28–29)}, weakest at \emph{disputed}:
\begin{itemize}
\item The believing of John 6:29 is trust, which brings nothing and receives Christ --- \cite[ad loc. John 6:29]{calvin-commentary-john-1847}; John 6:28-29
\item Good works are the fruit and evidence of saving faith, not its ground --- \cite[XI.2; XVI.2; XVI.5]{westminster-confession-1647}; \cite[III.xvi.1]{calvin-institutes-1960}; Jas 2:18-22; Eph 2:10; Rom 6:22; Luke 17:10
\item Justification and sanctification are inseparable but distinct --- \cite[Q. 77]{westminster-larger-catechism-1647}; 1 Cor 6:11; 1 Cor 1:30; Rom 4:6; Rom 4:8; Ezek 36:27; Heb 9:13-14; Rom 3:24-25; Rom 6:6; Rom 6:14; \cite[XI.1]{westminster-confession-1647}; \cite[III.xi.6]{calvin-institutes-1960}
\item an inference step --- \cite[Trent, Session VI (1547), canon 9]{tanner-decrees-1990}
\end{itemize}
\item \emph{Faith formed by charity (Aquinas on John 6:29)}, weakest at \emph{disputed}:
\begin{itemize}
\item The believing of John 6:29 is faith living through charity, the source of works --- \cite[cap. 6, lect. 3, n. 901]{aquinas-commentary-john-2012}; \cite[Trent, Session VI (1547), Decree on Justification, ch. 7]{tanner-decrees-1990}
\end{itemize}
\item \emph{Luther on Galatians 3:11–12 (against faith formed by charity)}, weakest at \emph{disputed}:
\begin{itemize}
\item an inference step --- \cite[ad loc. Gal 3:12]{luther-commentary-galatians-1998}; \cite[ad loc. Gal 3:11]{luther-commentary-galatians-1998}
\end{itemize}
\item \emph{Chrysostom on John 3:36 (believing is not enough without a right life)}, weakest at \emph{disputed}:
\begin{itemize}
\item an inference step --- \cite[Homily 31, §1]{chrysostom-homilies-john}; \cite{stanley-did-jesus-teach-2006}
\end{itemize}
\item \emph{Obedience as sanctification (John 3:36, Matthew 7:21)}, weakest at \emph{disputed}:
\begin{itemize}
\item The obedience of John 3:36 and Matthew 7:21 is sanctification, the fruit of faith --- \cite[ad loc. John 3:36]{ridderbos-john-1997}; \cite[ad loc. John 3:36]{carson-john-1991}; \cite[ad loc. Matt 7:21]{morris-matthew-1992}; \cite[55--55]{schreiner-justification-by-works-2015}; \cite{hoekema-saved-by-grace-1994}; \cite{murray-redemption-1984}; \cite[XVI.2]{westminster-confession-1647}; \cite[XIII.1]{westminster-confession-1647}; \cite[ad loc. John 3:36]{calvin-commentary-john-1847}
\item Justification and sanctification are inseparable but distinct --- \cite[Q. 77]{westminster-larger-catechism-1647}; 1 Cor 6:11; 1 Cor 1:30; Rom 4:6; Rom 4:8; Ezek 36:27; Heb 9:13-14; Rom 3:24-25; Rom 6:6; Rom 6:14; \cite[XI.1]{westminster-confession-1647}; \cite[III.xi.6]{calvin-institutes-1960}
\item an inference step --- \cite[55--55]{schreiner-justification-by-works-2015}; \cite[ad loc. John 3:36]{calvin-commentary-john-1847}; \cite[XVI.2]{westminster-confession-1647}; \cite[XIII.1]{westminster-confession-1647}
\end{itemize}
\end{itemize}
\noindent Why \emph{Paul’s word (δικαιόω, by the rule of least meaning)} defeats \emph{Trent, against ‘not by works’, read as what Paul’s word means}, at the level of the claims: \texttt{where\allowbreak{}Trents\allowbreak{}Word\allowbreak{}Reading\allowbreak{}Falls}.
\par\medskip\noindent\textbf{\texttt{only\_\allowbreak{}pauls\_\allowbreak{}word\_\allowbreak{}prevails}}\par\nobreak\noindent
\textbf{Only the meaning of Paul’s word prevails outright.} What the debate forces, before any choice between rivals, is exactly this: Paul’s “justify” (δικαιόω) does not mean the renewal of the inward person. It is argued two ways, from how the word is used and from Romans 4.

Nothing defeats either argument. The one position that contradicts them, Trent taken as a claim about Paul’s word, rests on a definition rated \texttt{disputed}, and its replies fail against arguments whose every link is rated \texttt{well\allowbreak{}Supported} or better.

It does not show that sola fide prevails, or that Trent’s teaching is refuted. Nothing else is forced. Every other position is attacked by one the debate cannot rule out. One loses, Trent taken as a claim about Paul’s word, and every other is in a tie with a rival the debate cannot choose between.

{\small\begin{verbatim}
theorem only_pauls_word_prevails : Framework.grounded solaFideDispute.defeats
    = {Party.lexical, Party.romansOnTheWord}
-- axioms: propext, Classical.choice, Quot.sound
\end{verbatim}}
\par\medskip\noindent\textbf{\texttt{trents\allowbreak{}Word\allowbreak{}Reading\allowbreak{}Answered}}\par\nobreak\noindent
Why Trent, taken as a claim about Paul’s word, cannot be defended: the argument from the word defeats it, and nothing defeats that argument. Romans 4, read for the word, defeats it too; one unanswered attacker is enough.

\paragraph{\emph{Trent, against ‘not by works’, read as what Paul’s word means} cannot be defended: no admissible position holds it.}
\begin{itemize}
\item \emph{Paul’s word (δικαιόω, by the rule of least meaning)} defeats \emph{Trent, against ‘not by works’, read as what Paul’s word means} (\texttt{where\allowbreak{}Trents\allowbreak{}Word\allowbreak{}Reading\allowbreak{}Falls}), and nothing defeats \emph{Paul’s word (δικαιόω, by the rule of least meaning)}.
\end{itemize}
\par\noindent\textbf{What this rests on.} The reasons state 1 defeat and 17 absences of defeat, each a cell of the defeat table, \texttt{sola\allowbreak{}Fide\allowbreak{}Dispute\_\allowbreak{}defeats}, computed from the two parties' premises and checked by the kernel. Whether an attack survives turns on the attacker's weakest link and on the rating of what it attacks, so a changed rating can change the verdict. The weakest links:
\begin{itemize}
\item \emph{Paul’s word (δικαιόω, by the rule of least meaning)}, weakest at \emph{well supported}:
\begin{itemize}
\item In Paul, ‘justify’ (δικαιόω) is a verdict: to declare righteous --- \cite[s.v. δικαιόω]{bauer-danker-lexicon-2000}; \cite{morris-apostolic-preaching-1965}; \cite[p. 80]{moo-romans-1996}; \cite{irons-righteousness-of-god-2015}; \cite{wright-justification-2009}; \cite[ad loc. Rom 3:24]{fitzmyer-romans-1993}
\item A word contributes the least meaning its context requires; one occurrence does not carry the whole concept --- \cite{silva-biblical-words-1994}; \cite{joos-semantic-axiom-1972}; \cite{barr-semantics-2004}
\item an inference step --- \cite{silva-biblical-words-1994}; \cite[p. 80]{moo-romans-1996}; \cite{barr-semantics-2004}
\end{itemize}
\item \emph{Trent, against ‘not by works’, read as what Paul’s word means}, weakest at \emph{disputed}:
\begin{itemize}
\item Justification is not remission of sins only, but renewal of the inward man --- \cite[Trent, Session VI (1547), Decree on Justification, ch. 7]{tanner-decrees-1990}; \cite[Trent, Session VI (1547), canon 11]{tanner-decrees-1990}
\item an inference step --- \cite[Trent, Session VI (1547), Decree on Justification, ch. 8]{tanner-decrees-1990}; \cite[Trent, Session VI (1547), Decree on Justification, ch. 7]{tanner-decrees-1990}; \cite[On the Spirit and the Letter, ch. 45 (26.45)]{npnf1-05-1887}; \cite{mcgrath-iustitia-dei-2005}
\end{itemize}
\item \emph{Sola fide from Paul (Galatians 2:16)}, weakest at \emph{disputed}:
\begin{itemize}
\item In Galatians 2:16 ‘faith of Christ’ (πίστις Χριστοῦ) means our faith in Christ --- \cite{dunn-once-more-pistis-christou-1991}; \cite{matlock-detheologizing-2000}; Gal 2:16; Rom 3:22; Phil 3:9
\item Second Temple Judaism: in by grace, staying in by works (encoded as a denial, so it ranks disputed) --- \cite{sanders-paul-palestinian-judaism-1977}
\item Good works are the fruit and evidence of saving faith, not its ground --- \cite[XI.2; XVI.2; XVI.5]{westminster-confession-1647}; \cite[III.xvi.1]{calvin-institutes-1960}; Jas 2:18-22; Eph 2:10; Rom 6:22; Luke 17:10
\item Justification and sanctification are inseparable but distinct --- \cite[Q. 77]{westminster-larger-catechism-1647}; 1 Cor 6:11; 1 Cor 1:30; Rom 4:6; Rom 4:8; Ezek 36:27; Heb 9:13-14; Rom 3:24-25; Rom 6:6; Rom 6:14; \cite[XI.1]{westminster-confession-1647}; \cite[III.xi.6]{calvin-institutes-1960}
\item an inference step --- \cite[Trent, Session VI (1547), canon 9]{tanner-decrees-1990}
\item an inference step --- \cite{gathercole-where-boasting-2002}; \cite{carson-variegated-nomism-1-2001}
\end{itemize}
\item \emph{Luke 7:50: “your faith has saved you”}, weakest at \emph{disputed}:
\begin{itemize}
\item Luke 7:47 — her love is the evidence of her forgiveness, not its ground --- \cite[ad loc. Luke 7:47]{marshall-luke-1978}; \cite[54--54]{padilla-narrative-criticism-2021}; \cite[Apology IV, on love and the fulfilling of the law: Luke 7:47–50]{kolb-wengert-book-of-concord-2000}; Luke 7:41-43; Luke 7:47; Luke 7:47-50
\item an inference step --- \cite[Trent, Session VI (1547), canon 9]{tanner-decrees-1990}
\end{itemize}
\item \emph{Sola fide from Peter (Acts 15:9–11)}, weakest at \emph{disputed}:
\begin{itemize}
\item The yoke refused at Acts 15:10 is the whole law as a condition of salvation --- \cite[ad loc. Acts 15:10]{bruce-acts-1988}; Acts 15:10; Gal 5:1; Jas 2:10; \cite{das-paul-law-covenant-2001}
\item Good works are the fruit and evidence of saving faith, not its ground --- \cite[XI.2; XVI.2; XVI.5]{westminster-confession-1647}; \cite[III.xvi.1]{calvin-institutes-1960}; Jas 2:18-22; Eph 2:10; Rom 6:22; Luke 17:10
\item Justification and sanctification are inseparable but distinct --- \cite[Q. 77]{westminster-larger-catechism-1647}; 1 Cor 6:11; 1 Cor 1:30; Rom 4:6; Rom 4:8; Ezek 36:27; Heb 9:13-14; Rom 3:24-25; Rom 6:6; Rom 6:14; \cite[XI.1]{westminster-confession-1647}; \cite[III.xi.6]{calvin-institutes-1960}
\item an inference step --- \cite[Trent, Session VI (1547), canon 9]{tanner-decrees-1990}
\end{itemize}
\item \emph{Trent, against ‘not by works’}, weakest at \emph{disputed}:
\begin{itemize}
\item Justification is not remission of sins only, but renewal of the inward man --- \cite[Trent, Session VI (1547), Decree on Justification, ch. 7]{tanner-decrees-1990}; \cite[Trent, Session VI (1547), canon 11]{tanner-decrees-1990}
\end{itemize}
\item \emph{Apocalyptic reading, against faith as the condition}, weakest at \emph{disputed}:
\begin{itemize}
\item ‘Faith of Christ’ is Christ’s own faithfulness, the means of God’s rescue (encoded as a denial, so it ranks disputed) --- \cite{campbell-deliverance-god-2009}; \cite[ad loc. Gal 2:16]{martyn-galatians-1997}; \cite{hays-faith-jesus-christ-2002}
\item ‘The righteousness of God’ (δικαιοσύνη θεοῦ) names God’s act of rescue in Christ --- \cite{campbell-deliverance-god-2009}; \cite{martyn-galatians-1997}; Rom 3:21-26
\end{itemize}
\item \emph{Covenantal nomism (Sanders, Dunn)}, weakest at \emph{disputed}:
\begin{itemize}
\item Second Temple Judaism: in by grace, staying in by works --- \cite{sanders-paul-palestinian-judaism-1977}
\item an inference step --- \cite{dunn-new-perspective-2005}; \cite{wright-what-paul-said-1997}
\end{itemize}
\item \emph{Variegated nomism (Gathercole, Carson et al.)}, weakest at \emph{disputed}:
\begin{itemize}
\item Second Temple Judaism also held final vindication according to works (encoded as a denial, so it ranks disputed) --- \cite{gathercole-where-boasting-2002}; \cite{carson-variegated-nomism-1-2001}
\item an inference step --- \cite{gathercole-where-boasting-2002}; \cite{carson-variegated-nomism-1-2001}
\end{itemize}
\item \emph{Law-observant Luke (Jervell)}, weakest at \emph{disputed}:
\begin{itemize}
\item an inference step --- \cite{jervell-luke-people-god-1972}
\end{itemize}
\item \emph{Paul’s gospel (Galatians 1), against justification as renewal}, weakest at \emph{disputed}:
\begin{itemize}
\item an inference step --- \cite[ad loc. Gal 1:6]{moo-galatians-2013}; \cite[XI.1]{westminster-confession-1647}
\item an inference step --- \cite[XI.1]{westminster-confession-1647}; \cite[III.xi.2]{calvin-institutes-1960}; \cite{morris-apostolic-preaching-1965}
\end{itemize}
\item \emph{Romans 4:3–8, on Paul’s word (by the rule of least meaning)}, weakest at \emph{well supported}:
\begin{itemize}
\item Romans 4:4–5 contrasts wages owed with a gift reckoned to the one who believes --- Rom 4:4-5; \cite[III.xi.18]{calvin-institutes-1960}
\item ‘Count’ (λογίζομαι) in Romans 4:3–11 is the language of an account: credited to one --- \cite[s.v. λογίζομαι]{bauer-danker-lexicon-2000}; \cite[ad loc. Rom 4:3]{moo-romans-1996}
\item A word contributes the least meaning its context requires; one occurrence does not carry the whole concept --- \cite{silva-biblical-words-1994}; \cite{joos-semantic-axiom-1972}; \cite{barr-semantics-2004}
\item an inference step --- \cite[ad loc. Rom 4:6]{moo-romans-1996}; \cite[s.v. λογίζομαι]{bauer-danker-lexicon-2000}
\item an inference step --- \cite{silva-biblical-words-1994}; \cite[p. 80]{moo-romans-1996}; \cite{barr-semantics-2004}
\end{itemize}
\item \emph{Sola fide from John (6:28–29)}, weakest at \emph{disputed}:
\begin{itemize}
\item The believing of John 6:29 is trust, which brings nothing and receives Christ --- \cite[ad loc. John 6:29]{calvin-commentary-john-1847}; John 6:28-29
\item Good works are the fruit and evidence of saving faith, not its ground --- \cite[XI.2; XVI.2; XVI.5]{westminster-confession-1647}; \cite[III.xvi.1]{calvin-institutes-1960}; Jas 2:18-22; Eph 2:10; Rom 6:22; Luke 17:10
\item Justification and sanctification are inseparable but distinct --- \cite[Q. 77]{westminster-larger-catechism-1647}; 1 Cor 6:11; 1 Cor 1:30; Rom 4:6; Rom 4:8; Ezek 36:27; Heb 9:13-14; Rom 3:24-25; Rom 6:6; Rom 6:14; \cite[XI.1]{westminster-confession-1647}; \cite[III.xi.6]{calvin-institutes-1960}
\item an inference step --- \cite[Trent, Session VI (1547), canon 9]{tanner-decrees-1990}
\end{itemize}
\item \emph{Faith formed by charity (Aquinas on John 6:29)}, weakest at \emph{disputed}:
\begin{itemize}
\item The believing of John 6:29 is faith living through charity, the source of works --- \cite[cap. 6, lect. 3, n. 901]{aquinas-commentary-john-2012}; \cite[Trent, Session VI (1547), Decree on Justification, ch. 7]{tanner-decrees-1990}
\end{itemize}
\item \emph{Luther on Galatians 3:11–12 (against faith formed by charity)}, weakest at \emph{disputed}:
\begin{itemize}
\item an inference step --- \cite[ad loc. Gal 3:12]{luther-commentary-galatians-1998}; \cite[ad loc. Gal 3:11]{luther-commentary-galatians-1998}
\end{itemize}
\item \emph{Chrysostom on John 3:36 (believing is not enough without a right life)}, weakest at \emph{disputed}:
\begin{itemize}
\item an inference step --- \cite[Homily 31, §1]{chrysostom-homilies-john}; \cite{stanley-did-jesus-teach-2006}
\end{itemize}
\item \emph{Obedience as sanctification (John 3:36, Matthew 7:21)}, weakest at \emph{disputed}:
\begin{itemize}
\item The obedience of John 3:36 and Matthew 7:21 is sanctification, the fruit of faith --- \cite[ad loc. John 3:36]{ridderbos-john-1997}; \cite[ad loc. John 3:36]{carson-john-1991}; \cite[ad loc. Matt 7:21]{morris-matthew-1992}; \cite[55--55]{schreiner-justification-by-works-2015}; \cite{hoekema-saved-by-grace-1994}; \cite{murray-redemption-1984}; \cite[XVI.2]{westminster-confession-1647}; \cite[XIII.1]{westminster-confession-1647}; \cite[ad loc. John 3:36]{calvin-commentary-john-1847}
\item Justification and sanctification are inseparable but distinct --- \cite[Q. 77]{westminster-larger-catechism-1647}; 1 Cor 6:11; 1 Cor 1:30; Rom 4:6; Rom 4:8; Ezek 36:27; Heb 9:13-14; Rom 3:24-25; Rom 6:6; Rom 6:14; \cite[XI.1]{westminster-confession-1647}; \cite[III.xi.6]{calvin-institutes-1960}
\item an inference step --- \cite[55--55]{schreiner-justification-by-works-2015}; \cite[ad loc. John 3:36]{calvin-commentary-john-1847}; \cite[XVI.2]{westminster-confession-1647}; \cite[XIII.1]{westminster-confession-1647}
\end{itemize}
\end{itemize}
\noindent Why \emph{Paul’s word (δικαιόω, by the rule of least meaning)} defeats \emph{Trent, against ‘not by works’, read as what Paul’s word means}, at the level of the claims: \texttt{where\allowbreak{}Trents\allowbreak{}Word\allowbreak{}Reading\allowbreak{}Falls}.
\par\medskip\noindent\textbf{\texttt{trent\_\allowbreak{}on\_\allowbreak{}pauls\_\allowbreak{}word\_\allowbreak{}indefensible}}\par\nobreak\noindent
\textbf{Trent, taken as a claim about Paul’s word, cannot be defended.} No consistent set of positions that answers its attackers can include it. The two arguments about the word defeat it, and nothing answers either.

Trent’s definition can be taken two ways. This is the first, weighed in the debate (\texttt{what\allowbreak{}Trents\allowbreak{}Definition\allowbreak{}Claims}): if the definition says what Paul’s “justify” means, it cannot be defended here, for a stated and cited reason, at a step rated \texttt{well\allowbreak{}Supported}.

It does not show that Trent’s definition is false as a statement of what God does. That is the other way of taking it, and it can be defended (\texttt{trent\_\allowbreak{}defensible}). Nor that the position contradicts itself: taken alone it can be held without contradiction.

{\small\begin{verbatim}
theorem trent_on_pauls_word_indefensible : ∀ (S : Set Party),
    Framework.Admissible solaFideDispute.defeats S → Party.trentOnPaulsWord ∉
    S
-- axioms: propext, Classical.choice, Quot.sound
\end{verbatim}}
\par\medskip\noindent\textbf{\texttt{trent\allowbreak{}Stands\allowbreak{}With\allowbreak{}Sanders\allowbreak{}And\allowbreak{}Jervell}}\par\nobreak\noindent
Why Trent can be defended: it can be held with Sanders and Jervell, and it answers every position that attacks it. Those are each Reformed strand, Paul’s gospel, and the apocalyptic reading, whose “not by works” it denies.

\paragraph{\emph{Trent, against ‘not by works’} is accepted on some resolution.}
\begin{itemize}
\item A position holding \emph{Trent, against ‘not by works’}, \emph{Covenantal nomism (Sanders, Dunn)}, \emph{Law-observant Luke (Jervell)} can be held.
\item None of \emph{Trent, against ‘not by works’}, \emph{Covenantal nomism (Sanders, Dunn)}, \emph{Law-observant Luke (Jervell)} defeats another, and each attack on them is answered from among them:
\begin{itemize}
\item \emph{Sola fide from Paul (Galatians 2:16)} defeats \emph{Trent, against ‘not by works’}, and \emph{Trent, against ‘not by works’} defeats \emph{Sola fide from Paul (Galatians 2:16)}.
\item \emph{Luke 7:50: “your faith has saved you”} defeats \emph{Trent, against ‘not by works’} (\texttt{why\allowbreak{}The\allowbreak{}Dominical\allowbreak{}Case\allowbreak{}Stands\allowbreak{}Against\allowbreak{}Trent}), and \emph{Trent, against ‘not by works’} defeats \emph{Luke 7:50: “your faith has saved you”}.
\item \emph{Sola fide from Peter (Acts 15:9–11)} defeats \emph{Trent, against ‘not by works’}, and \emph{Trent, against ‘not by works’} defeats \emph{Sola fide from Peter (Acts 15:9–11)}.
\item \emph{Apocalyptic reading, against faith as the condition} defeats \emph{Trent, against ‘not by works’}, and \emph{Trent, against ‘not by works’} defeats \emph{Apocalyptic reading, against faith as the condition}.
\item \emph{Paul’s gospel (Galatians 1), against justification as renewal} defeats \emph{Trent, against ‘not by works’} (\texttt{where\allowbreak{}Trent\allowbreak{}Parts\allowbreak{}From\allowbreak{}Paul}), and \emph{Trent, against ‘not by works’} defeats \emph{Paul’s gospel (Galatians 1), against justification as renewal}.
\item \emph{Sola fide from John (6:28–29)} defeats \emph{Trent, against ‘not by works’}, and \emph{Trent, against ‘not by works’} defeats \emph{Sola fide from John (6:28–29)}.
\item \emph{Obedience as sanctification (John 3:36, Matthew 7:21)} defeats \emph{Trent, against ‘not by works’}, and \emph{Trent, against ‘not by works’} defeats \emph{Obedience as sanctification (John 3:36, Matthew 7:21)}.
\item \emph{Sola fide from Paul (Galatians 2:16)} defeats \emph{Covenantal nomism (Sanders, Dunn)}, and \emph{Trent, against ‘not by works’} defeats \emph{Sola fide from Paul (Galatians 2:16)}.
\item \emph{Variegated nomism (Gathercole, Carson et al.)} defeats \emph{Covenantal nomism (Sanders, Dunn)}, and \emph{Covenantal nomism (Sanders, Dunn)} defeats \emph{Variegated nomism (Gathercole, Carson et al.)}.
\item \emph{Sola fide from Peter (Acts 15:9–11)} defeats \emph{Law-observant Luke (Jervell)}, and \emph{Trent, against ‘not by works’} defeats \emph{Sola fide from Peter (Acts 15:9–11)}.
\end{itemize}
\end{itemize}
\par\noindent\textbf{What this rests on.} The reasons state 18 defeats and 41 absences of defeat, each a cell of the defeat table, \texttt{sola\allowbreak{}Fide\allowbreak{}Dispute\_\allowbreak{}defeats}, computed from the two parties' premises and checked by the kernel. Whether an attack survives turns on the attacker's weakest link and on the rating of what it attacks, so a changed rating can change the verdict. The weakest links:
\begin{itemize}
\item \emph{Trent, against ‘not by works’}, weakest at \emph{disputed}:
\begin{itemize}
\item Justification is not remission of sins only, but renewal of the inward man --- \cite[Trent, Session VI (1547), Decree on Justification, ch. 7]{tanner-decrees-1990}; \cite[Trent, Session VI (1547), canon 11]{tanner-decrees-1990}
\end{itemize}
\item \emph{Covenantal nomism (Sanders, Dunn)}, weakest at \emph{disputed}:
\begin{itemize}
\item Second Temple Judaism: in by grace, staying in by works --- \cite{sanders-paul-palestinian-judaism-1977}
\item an inference step --- \cite{dunn-new-perspective-2005}; \cite{wright-what-paul-said-1997}
\end{itemize}
\item \emph{Law-observant Luke (Jervell)}, weakest at \emph{disputed}:
\begin{itemize}
\item an inference step --- \cite{jervell-luke-people-god-1972}
\end{itemize}
\item \emph{Sola fide from Paul (Galatians 2:16)}, weakest at \emph{disputed}:
\begin{itemize}
\item In Galatians 2:16 ‘faith of Christ’ (πίστις Χριστοῦ) means our faith in Christ --- \cite{dunn-once-more-pistis-christou-1991}; \cite{matlock-detheologizing-2000}; Gal 2:16; Rom 3:22; Phil 3:9
\item Second Temple Judaism: in by grace, staying in by works (encoded as a denial, so it ranks disputed) --- \cite{sanders-paul-palestinian-judaism-1977}
\item Good works are the fruit and evidence of saving faith, not its ground --- \cite[XI.2; XVI.2; XVI.5]{westminster-confession-1647}; \cite[III.xvi.1]{calvin-institutes-1960}; Jas 2:18-22; Eph 2:10; Rom 6:22; Luke 17:10
\item Justification and sanctification are inseparable but distinct --- \cite[Q. 77]{westminster-larger-catechism-1647}; 1 Cor 6:11; 1 Cor 1:30; Rom 4:6; Rom 4:8; Ezek 36:27; Heb 9:13-14; Rom 3:24-25; Rom 6:6; Rom 6:14; \cite[XI.1]{westminster-confession-1647}; \cite[III.xi.6]{calvin-institutes-1960}
\item an inference step --- \cite[Trent, Session VI (1547), canon 9]{tanner-decrees-1990}
\item an inference step --- \cite{gathercole-where-boasting-2002}; \cite{carson-variegated-nomism-1-2001}
\end{itemize}
\item \emph{Luke 7:50: “your faith has saved you”}, weakest at \emph{disputed}:
\begin{itemize}
\item Luke 7:47 — her love is the evidence of her forgiveness, not its ground --- \cite[ad loc. Luke 7:47]{marshall-luke-1978}; \cite[54--54]{padilla-narrative-criticism-2021}; \cite[Apology IV, on love and the fulfilling of the law: Luke 7:47–50]{kolb-wengert-book-of-concord-2000}; Luke 7:41-43; Luke 7:47; Luke 7:47-50
\item an inference step --- \cite[Trent, Session VI (1547), canon 9]{tanner-decrees-1990}
\end{itemize}
\item \emph{Sola fide from Peter (Acts 15:9–11)}, weakest at \emph{disputed}:
\begin{itemize}
\item The yoke refused at Acts 15:10 is the whole law as a condition of salvation --- \cite[ad loc. Acts 15:10]{bruce-acts-1988}; Acts 15:10; Gal 5:1; Jas 2:10; \cite{das-paul-law-covenant-2001}
\item Good works are the fruit and evidence of saving faith, not its ground --- \cite[XI.2; XVI.2; XVI.5]{westminster-confession-1647}; \cite[III.xvi.1]{calvin-institutes-1960}; Jas 2:18-22; Eph 2:10; Rom 6:22; Luke 17:10
\item Justification and sanctification are inseparable but distinct --- \cite[Q. 77]{westminster-larger-catechism-1647}; 1 Cor 6:11; 1 Cor 1:30; Rom 4:6; Rom 4:8; Ezek 36:27; Heb 9:13-14; Rom 3:24-25; Rom 6:6; Rom 6:14; \cite[XI.1]{westminster-confession-1647}; \cite[III.xi.6]{calvin-institutes-1960}
\item an inference step --- \cite[Trent, Session VI (1547), canon 9]{tanner-decrees-1990}
\end{itemize}
\item \emph{Apocalyptic reading, against faith as the condition}, weakest at \emph{disputed}:
\begin{itemize}
\item ‘Faith of Christ’ is Christ’s own faithfulness, the means of God’s rescue (encoded as a denial, so it ranks disputed) --- \cite{campbell-deliverance-god-2009}; \cite[ad loc. Gal 2:16]{martyn-galatians-1997}; \cite{hays-faith-jesus-christ-2002}
\item ‘The righteousness of God’ (δικαιοσύνη θεοῦ) names God’s act of rescue in Christ --- \cite{campbell-deliverance-god-2009}; \cite{martyn-galatians-1997}; Rom 3:21-26
\end{itemize}
\item \emph{Variegated nomism (Gathercole, Carson et al.)}, weakest at \emph{disputed}:
\begin{itemize}
\item Second Temple Judaism also held final vindication according to works (encoded as a denial, so it ranks disputed) --- \cite{gathercole-where-boasting-2002}; \cite{carson-variegated-nomism-1-2001}
\item an inference step --- \cite{gathercole-where-boasting-2002}; \cite{carson-variegated-nomism-1-2001}
\end{itemize}
\item \emph{Paul’s gospel (Galatians 1), against justification as renewal}, weakest at \emph{disputed}:
\begin{itemize}
\item an inference step --- \cite[ad loc. Gal 1:6]{moo-galatians-2013}; \cite[XI.1]{westminster-confession-1647}
\item an inference step --- \cite[XI.1]{westminster-confession-1647}; \cite[III.xi.2]{calvin-institutes-1960}; \cite{morris-apostolic-preaching-1965}
\end{itemize}
\item \emph{Paul’s word (δικαιόω, by the rule of least meaning)}, weakest at \emph{well supported}:
\begin{itemize}
\item In Paul, ‘justify’ (δικαιόω) is a verdict: to declare righteous --- \cite[s.v. δικαιόω]{bauer-danker-lexicon-2000}; \cite{morris-apostolic-preaching-1965}; \cite[p. 80]{moo-romans-1996}; \cite{irons-righteousness-of-god-2015}; \cite{wright-justification-2009}; \cite[ad loc. Rom 3:24]{fitzmyer-romans-1993}
\item A word contributes the least meaning its context requires; one occurrence does not carry the whole concept --- \cite{silva-biblical-words-1994}; \cite{joos-semantic-axiom-1972}; \cite{barr-semantics-2004}
\item an inference step --- \cite{silva-biblical-words-1994}; \cite[p. 80]{moo-romans-1996}; \cite{barr-semantics-2004}
\end{itemize}
\item \emph{Trent, against ‘not by works’, read as what Paul’s word means}, weakest at \emph{disputed}:
\begin{itemize}
\item Justification is not remission of sins only, but renewal of the inward man --- \cite[Trent, Session VI (1547), Decree on Justification, ch. 7]{tanner-decrees-1990}; \cite[Trent, Session VI (1547), canon 11]{tanner-decrees-1990}
\item an inference step --- \cite[Trent, Session VI (1547), Decree on Justification, ch. 8]{tanner-decrees-1990}; \cite[Trent, Session VI (1547), Decree on Justification, ch. 7]{tanner-decrees-1990}; \cite[On the Spirit and the Letter, ch. 45 (26.45)]{npnf1-05-1887}; \cite{mcgrath-iustitia-dei-2005}
\end{itemize}
\item \emph{Romans 4:3–8, on Paul’s word (by the rule of least meaning)}, weakest at \emph{well supported}:
\begin{itemize}
\item Romans 4:4–5 contrasts wages owed with a gift reckoned to the one who believes --- Rom 4:4-5; \cite[III.xi.18]{calvin-institutes-1960}
\item ‘Count’ (λογίζομαι) in Romans 4:3–11 is the language of an account: credited to one --- \cite[s.v. λογίζομαι]{bauer-danker-lexicon-2000}; \cite[ad loc. Rom 4:3]{moo-romans-1996}
\item A word contributes the least meaning its context requires; one occurrence does not carry the whole concept --- \cite{silva-biblical-words-1994}; \cite{joos-semantic-axiom-1972}; \cite{barr-semantics-2004}
\item an inference step --- \cite[ad loc. Rom 4:6]{moo-romans-1996}; \cite[s.v. λογίζομαι]{bauer-danker-lexicon-2000}
\item an inference step --- \cite{silva-biblical-words-1994}; \cite[p. 80]{moo-romans-1996}; \cite{barr-semantics-2004}
\end{itemize}
\item \emph{Sola fide from John (6:28–29)}, weakest at \emph{disputed}:
\begin{itemize}
\item The believing of John 6:29 is trust, which brings nothing and receives Christ --- \cite[ad loc. John 6:29]{calvin-commentary-john-1847}; John 6:28-29
\item Good works are the fruit and evidence of saving faith, not its ground --- \cite[XI.2; XVI.2; XVI.5]{westminster-confession-1647}; \cite[III.xvi.1]{calvin-institutes-1960}; Jas 2:18-22; Eph 2:10; Rom 6:22; Luke 17:10
\item Justification and sanctification are inseparable but distinct --- \cite[Q. 77]{westminster-larger-catechism-1647}; 1 Cor 6:11; 1 Cor 1:30; Rom 4:6; Rom 4:8; Ezek 36:27; Heb 9:13-14; Rom 3:24-25; Rom 6:6; Rom 6:14; \cite[XI.1]{westminster-confession-1647}; \cite[III.xi.6]{calvin-institutes-1960}
\item an inference step --- \cite[Trent, Session VI (1547), canon 9]{tanner-decrees-1990}
\end{itemize}
\item \emph{Faith formed by charity (Aquinas on John 6:29)}, weakest at \emph{disputed}:
\begin{itemize}
\item The believing of John 6:29 is faith living through charity, the source of works --- \cite[cap. 6, lect. 3, n. 901]{aquinas-commentary-john-2012}; \cite[Trent, Session VI (1547), Decree on Justification, ch. 7]{tanner-decrees-1990}
\end{itemize}
\item \emph{Luther on Galatians 3:11–12 (against faith formed by charity)}, weakest at \emph{disputed}:
\begin{itemize}
\item an inference step --- \cite[ad loc. Gal 3:12]{luther-commentary-galatians-1998}; \cite[ad loc. Gal 3:11]{luther-commentary-galatians-1998}
\end{itemize}
\item \emph{Chrysostom on John 3:36 (believing is not enough without a right life)}, weakest at \emph{disputed}:
\begin{itemize}
\item an inference step --- \cite[Homily 31, §1]{chrysostom-homilies-john}; \cite{stanley-did-jesus-teach-2006}
\end{itemize}
\item \emph{Obedience as sanctification (John 3:36, Matthew 7:21)}, weakest at \emph{disputed}:
\begin{itemize}
\item The obedience of John 3:36 and Matthew 7:21 is sanctification, the fruit of faith --- \cite[ad loc. John 3:36]{ridderbos-john-1997}; \cite[ad loc. John 3:36]{carson-john-1991}; \cite[ad loc. Matt 7:21]{morris-matthew-1992}; \cite[55--55]{schreiner-justification-by-works-2015}; \cite{hoekema-saved-by-grace-1994}; \cite{murray-redemption-1984}; \cite[XVI.2]{westminster-confession-1647}; \cite[XIII.1]{westminster-confession-1647}; \cite[ad loc. John 3:36]{calvin-commentary-john-1847}
\item Justification and sanctification are inseparable but distinct --- \cite[Q. 77]{westminster-larger-catechism-1647}; 1 Cor 6:11; 1 Cor 1:30; Rom 4:6; Rom 4:8; Ezek 36:27; Heb 9:13-14; Rom 3:24-25; Rom 6:6; Rom 6:14; \cite[XI.1]{westminster-confession-1647}; \cite[III.xi.6]{calvin-institutes-1960}
\item an inference step --- \cite[55--55]{schreiner-justification-by-works-2015}; \cite[ad loc. John 3:36]{calvin-commentary-john-1847}; \cite[XVI.2]{westminster-confession-1647}; \cite[XIII.1]{westminster-confession-1647}
\end{itemize}
\end{itemize}
\noindent Why \emph{Luke 7:50: “your faith has saved you”} defeats \emph{Trent, against ‘not by works’}, at the level of the claims: \texttt{why\allowbreak{}The\allowbreak{}Dominical\allowbreak{}Case\allowbreak{}Stands\allowbreak{}Against\allowbreak{}Trent}.
\noindent Why \emph{Paul’s gospel (Galatians 1), against justification as renewal} defeats \emph{Trent, against ‘not by works’}, at the level of the claims: \texttt{where\allowbreak{}Trent\allowbreak{}Parts\allowbreak{}From\allowbreak{}Paul}.
\par\medskip\noindent\textbf{\texttt{trent\_\allowbreak{}defensible}}\par\nobreak\noindent
\textbf{Trent can be defended.} A consistent set of positions includes it, with Sanders and Jervell, and answers every attack. Every position that defeats Trent, Trent defeats in return: Paul’s case, Luke’s, Peter’s, John’s, Paul’s gospel, the apocalyptic reading, and the Reformed answer to Chrysostom. The ratings, all at the bottom, break no tie.

The arguments about Paul’s word do not attack it. Trent’s definition, taken as it stands, says nothing about what “justify” means.

It does not show that Trent prevails. It is not forced (\texttt{trent\_\allowbreak{}not\_\allowbreak{}forced}). It is in a tie with each of those seven positions.

{\small\begin{verbatim}
theorem trent_defensible : Framework.CredulouslyAccepted
    solaFideDispute.defeats Party.trent
-- axioms: propext, Classical.choice, Quot.sound
\end{verbatim}}
\par\medskip\noindent\textbf{\texttt{trent\allowbreak{}Answered\allowbreak{}By\allowbreak{}Paul}}\par\nobreak\noindent
Why Trent is not forced: a set of positions that can be defended includes Paul’s case, and Paul’s case defeats Trent.

\paragraph{\emph{Trent, against ‘not by works’} is not accepted on every resolution.}
\begin{itemize}
\item A position holding \emph{Sola fide from Paul (Galatians 2:16)}, \emph{Luke 7:50: “your faith has saved you”}, \emph{Sola fide from Peter (Acts 15:9–11)}, \emph{Variegated nomism (Gathercole, Carson et al.)} can be held, and it holds \emph{Sola fide from Paul (Galatians 2:16)}; \emph{Sola fide from Paul (Galatians 2:16)} conflicts with \emph{Trent, against ‘not by works’}: \emph{Sola fide from Paul (Galatians 2:16)} defeats \emph{Trent, against ‘not by works’}, and \emph{Trent, against ‘not by works’} defeats \emph{Sola fide from Paul (Galatians 2:16)}.
\item None of \emph{Sola fide from Paul (Galatians 2:16)}, \emph{Luke 7:50: “your faith has saved you”}, \emph{Sola fide from Peter (Acts 15:9–11)}, \emph{Variegated nomism (Gathercole, Carson et al.)} defeats another, and each attack on them is answered from among them:
\begin{itemize}
\item \emph{Trent, against ‘not by works’} defeats \emph{Sola fide from Paul (Galatians 2:16)}, and \emph{Sola fide from Paul (Galatians 2:16)} defeats \emph{Trent, against ‘not by works’}.
\item \emph{Apocalyptic reading, against faith as the condition} defeats \emph{Sola fide from Paul (Galatians 2:16)}, and \emph{Sola fide from Paul (Galatians 2:16)} defeats \emph{Apocalyptic reading, against faith as the condition}.
\item \emph{Covenantal nomism (Sanders, Dunn)} defeats \emph{Sola fide from Paul (Galatians 2:16)}, and \emph{Sola fide from Paul (Galatians 2:16)} defeats \emph{Covenantal nomism (Sanders, Dunn)}.
\item \emph{Trent, against ‘not by works’, read as what Paul’s word means} defeats \emph{Sola fide from Paul (Galatians 2:16)}, and \emph{Sola fide from Paul (Galatians 2:16)} defeats \emph{Trent, against ‘not by works’, read as what Paul’s word means}.
\item \emph{Chrysostom on John 3:36 (believing is not enough without a right life)} defeats \emph{Sola fide from Paul (Galatians 2:16)}, and \emph{Sola fide from Paul (Galatians 2:16)} defeats \emph{Chrysostom on John 3:36 (believing is not enough without a right life)}.
\item \emph{Trent, against ‘not by works’} defeats \emph{Luke 7:50: “your faith has saved you”}, and \emph{Sola fide from Paul (Galatians 2:16)} defeats \emph{Trent, against ‘not by works’}.
\item \emph{Trent, against ‘not by works’, read as what Paul’s word means} defeats \emph{Luke 7:50: “your faith has saved you”}, and \emph{Sola fide from Paul (Galatians 2:16)} defeats \emph{Trent, against ‘not by works’, read as what Paul’s word means}.
\item \emph{Chrysostom on John 3:36 (believing is not enough without a right life)} defeats \emph{Luke 7:50: “your faith has saved you”}, and \emph{Sola fide from Paul (Galatians 2:16)} defeats \emph{Chrysostom on John 3:36 (believing is not enough without a right life)}.
\item \emph{Trent, against ‘not by works’} defeats \emph{Sola fide from Peter (Acts 15:9–11)}, and \emph{Sola fide from Paul (Galatians 2:16)} defeats \emph{Trent, against ‘not by works’}.
\item \emph{Law-observant Luke (Jervell)} defeats \emph{Sola fide from Peter (Acts 15:9–11)}, and \emph{Sola fide from Peter (Acts 15:9–11)} defeats \emph{Law-observant Luke (Jervell)}.
\item \emph{Trent, against ‘not by works’, read as what Paul’s word means} defeats \emph{Sola fide from Peter (Acts 15:9–11)}, and \emph{Sola fide from Paul (Galatians 2:16)} defeats \emph{Trent, against ‘not by works’, read as what Paul’s word means}.
\item \emph{Chrysostom on John 3:36 (believing is not enough without a right life)} defeats \emph{Sola fide from Peter (Acts 15:9–11)}, and \emph{Sola fide from Paul (Galatians 2:16)} defeats \emph{Chrysostom on John 3:36 (believing is not enough without a right life)}.
\item \emph{Covenantal nomism (Sanders, Dunn)} defeats \emph{Variegated nomism (Gathercole, Carson et al.)}, and \emph{Sola fide from Paul (Galatians 2:16)} defeats \emph{Covenantal nomism (Sanders, Dunn)}.
\end{itemize}
\end{itemize}
\par\noindent\textbf{What this rests on.} The reasons state 19 defeats and 55 absences of defeat, each a cell of the defeat table, \texttt{sola\allowbreak{}Fide\allowbreak{}Dispute\_\allowbreak{}defeats}, computed from the two parties' premises and checked by the kernel. Whether an attack survives turns on the attacker's weakest link and on the rating of what it attacks, so a changed rating can change the verdict. The weakest links:
\begin{itemize}
\item \emph{Sola fide from Paul (Galatians 2:16)}, weakest at \emph{disputed}:
\begin{itemize}
\item In Galatians 2:16 ‘faith of Christ’ (πίστις Χριστοῦ) means our faith in Christ --- \cite{dunn-once-more-pistis-christou-1991}; \cite{matlock-detheologizing-2000}; Gal 2:16; Rom 3:22; Phil 3:9
\item Second Temple Judaism: in by grace, staying in by works (encoded as a denial, so it ranks disputed) --- \cite{sanders-paul-palestinian-judaism-1977}
\item Good works are the fruit and evidence of saving faith, not its ground --- \cite[XI.2; XVI.2; XVI.5]{westminster-confession-1647}; \cite[III.xvi.1]{calvin-institutes-1960}; Jas 2:18-22; Eph 2:10; Rom 6:22; Luke 17:10
\item Justification and sanctification are inseparable but distinct --- \cite[Q. 77]{westminster-larger-catechism-1647}; 1 Cor 6:11; 1 Cor 1:30; Rom 4:6; Rom 4:8; Ezek 36:27; Heb 9:13-14; Rom 3:24-25; Rom 6:6; Rom 6:14; \cite[XI.1]{westminster-confession-1647}; \cite[III.xi.6]{calvin-institutes-1960}
\item an inference step --- \cite[Trent, Session VI (1547), canon 9]{tanner-decrees-1990}
\item an inference step --- \cite{gathercole-where-boasting-2002}; \cite{carson-variegated-nomism-1-2001}
\end{itemize}
\item \emph{Trent, against ‘not by works’}, weakest at \emph{disputed}:
\begin{itemize}
\item Justification is not remission of sins only, but renewal of the inward man --- \cite[Trent, Session VI (1547), Decree on Justification, ch. 7]{tanner-decrees-1990}; \cite[Trent, Session VI (1547), canon 11]{tanner-decrees-1990}
\end{itemize}
\item \emph{Luke 7:50: “your faith has saved you”}, weakest at \emph{disputed}:
\begin{itemize}
\item Luke 7:47 — her love is the evidence of her forgiveness, not its ground --- \cite[ad loc. Luke 7:47]{marshall-luke-1978}; \cite[54--54]{padilla-narrative-criticism-2021}; \cite[Apology IV, on love and the fulfilling of the law: Luke 7:47–50]{kolb-wengert-book-of-concord-2000}; Luke 7:41-43; Luke 7:47; Luke 7:47-50
\item an inference step --- \cite[Trent, Session VI (1547), canon 9]{tanner-decrees-1990}
\end{itemize}
\item \emph{Sola fide from Peter (Acts 15:9–11)}, weakest at \emph{disputed}:
\begin{itemize}
\item The yoke refused at Acts 15:10 is the whole law as a condition of salvation --- \cite[ad loc. Acts 15:10]{bruce-acts-1988}; Acts 15:10; Gal 5:1; Jas 2:10; \cite{das-paul-law-covenant-2001}
\item Good works are the fruit and evidence of saving faith, not its ground --- \cite[XI.2; XVI.2; XVI.5]{westminster-confession-1647}; \cite[III.xvi.1]{calvin-institutes-1960}; Jas 2:18-22; Eph 2:10; Rom 6:22; Luke 17:10
\item Justification and sanctification are inseparable but distinct --- \cite[Q. 77]{westminster-larger-catechism-1647}; 1 Cor 6:11; 1 Cor 1:30; Rom 4:6; Rom 4:8; Ezek 36:27; Heb 9:13-14; Rom 3:24-25; Rom 6:6; Rom 6:14; \cite[XI.1]{westminster-confession-1647}; \cite[III.xi.6]{calvin-institutes-1960}
\item an inference step --- \cite[Trent, Session VI (1547), canon 9]{tanner-decrees-1990}
\end{itemize}
\item \emph{Variegated nomism (Gathercole, Carson et al.)}, weakest at \emph{disputed}:
\begin{itemize}
\item Second Temple Judaism also held final vindication according to works (encoded as a denial, so it ranks disputed) --- \cite{gathercole-where-boasting-2002}; \cite{carson-variegated-nomism-1-2001}
\item an inference step --- \cite{gathercole-where-boasting-2002}; \cite{carson-variegated-nomism-1-2001}
\end{itemize}
\item \emph{Apocalyptic reading, against faith as the condition}, weakest at \emph{disputed}:
\begin{itemize}
\item ‘Faith of Christ’ is Christ’s own faithfulness, the means of God’s rescue (encoded as a denial, so it ranks disputed) --- \cite{campbell-deliverance-god-2009}; \cite[ad loc. Gal 2:16]{martyn-galatians-1997}; \cite{hays-faith-jesus-christ-2002}
\item ‘The righteousness of God’ (δικαιοσύνη θεοῦ) names God’s act of rescue in Christ --- \cite{campbell-deliverance-god-2009}; \cite{martyn-galatians-1997}; Rom 3:21-26
\end{itemize}
\item \emph{Covenantal nomism (Sanders, Dunn)}, weakest at \emph{disputed}:
\begin{itemize}
\item Second Temple Judaism: in by grace, staying in by works --- \cite{sanders-paul-palestinian-judaism-1977}
\item an inference step --- \cite{dunn-new-perspective-2005}; \cite{wright-what-paul-said-1997}
\end{itemize}
\item \emph{Law-observant Luke (Jervell)}, weakest at \emph{disputed}:
\begin{itemize}
\item an inference step --- \cite{jervell-luke-people-god-1972}
\end{itemize}
\item \emph{Paul’s gospel (Galatians 1), against justification as renewal}, weakest at \emph{disputed}:
\begin{itemize}
\item an inference step --- \cite[ad loc. Gal 1:6]{moo-galatians-2013}; \cite[XI.1]{westminster-confession-1647}
\item an inference step --- \cite[XI.1]{westminster-confession-1647}; \cite[III.xi.2]{calvin-institutes-1960}; \cite{morris-apostolic-preaching-1965}
\end{itemize}
\item \emph{Paul’s word (δικαιόω, by the rule of least meaning)}, weakest at \emph{well supported}:
\begin{itemize}
\item In Paul, ‘justify’ (δικαιόω) is a verdict: to declare righteous --- \cite[s.v. δικαιόω]{bauer-danker-lexicon-2000}; \cite{morris-apostolic-preaching-1965}; \cite[p. 80]{moo-romans-1996}; \cite{irons-righteousness-of-god-2015}; \cite{wright-justification-2009}; \cite[ad loc. Rom 3:24]{fitzmyer-romans-1993}
\item A word contributes the least meaning its context requires; one occurrence does not carry the whole concept --- \cite{silva-biblical-words-1994}; \cite{joos-semantic-axiom-1972}; \cite{barr-semantics-2004}
\item an inference step --- \cite{silva-biblical-words-1994}; \cite[p. 80]{moo-romans-1996}; \cite{barr-semantics-2004}
\end{itemize}
\item \emph{Trent, against ‘not by works’, read as what Paul’s word means}, weakest at \emph{disputed}:
\begin{itemize}
\item Justification is not remission of sins only, but renewal of the inward man --- \cite[Trent, Session VI (1547), Decree on Justification, ch. 7]{tanner-decrees-1990}; \cite[Trent, Session VI (1547), canon 11]{tanner-decrees-1990}
\item an inference step --- \cite[Trent, Session VI (1547), Decree on Justification, ch. 8]{tanner-decrees-1990}; \cite[Trent, Session VI (1547), Decree on Justification, ch. 7]{tanner-decrees-1990}; \cite[On the Spirit and the Letter, ch. 45 (26.45)]{npnf1-05-1887}; \cite{mcgrath-iustitia-dei-2005}
\end{itemize}
\item \emph{Romans 4:3–8, on Paul’s word (by the rule of least meaning)}, weakest at \emph{well supported}:
\begin{itemize}
\item Romans 4:4–5 contrasts wages owed with a gift reckoned to the one who believes --- Rom 4:4-5; \cite[III.xi.18]{calvin-institutes-1960}
\item ‘Count’ (λογίζομαι) in Romans 4:3–11 is the language of an account: credited to one --- \cite[s.v. λογίζομαι]{bauer-danker-lexicon-2000}; \cite[ad loc. Rom 4:3]{moo-romans-1996}
\item A word contributes the least meaning its context requires; one occurrence does not carry the whole concept --- \cite{silva-biblical-words-1994}; \cite{joos-semantic-axiom-1972}; \cite{barr-semantics-2004}
\item an inference step --- \cite[ad loc. Rom 4:6]{moo-romans-1996}; \cite[s.v. λογίζομαι]{bauer-danker-lexicon-2000}
\item an inference step --- \cite{silva-biblical-words-1994}; \cite[p. 80]{moo-romans-1996}; \cite{barr-semantics-2004}
\end{itemize}
\item \emph{Sola fide from John (6:28–29)}, weakest at \emph{disputed}:
\begin{itemize}
\item The believing of John 6:29 is trust, which brings nothing and receives Christ --- \cite[ad loc. John 6:29]{calvin-commentary-john-1847}; John 6:28-29
\item Good works are the fruit and evidence of saving faith, not its ground --- \cite[XI.2; XVI.2; XVI.5]{westminster-confession-1647}; \cite[III.xvi.1]{calvin-institutes-1960}; Jas 2:18-22; Eph 2:10; Rom 6:22; Luke 17:10
\item Justification and sanctification are inseparable but distinct --- \cite[Q. 77]{westminster-larger-catechism-1647}; 1 Cor 6:11; 1 Cor 1:30; Rom 4:6; Rom 4:8; Ezek 36:27; Heb 9:13-14; Rom 3:24-25; Rom 6:6; Rom 6:14; \cite[XI.1]{westminster-confession-1647}; \cite[III.xi.6]{calvin-institutes-1960}
\item an inference step --- \cite[Trent, Session VI (1547), canon 9]{tanner-decrees-1990}
\end{itemize}
\item \emph{Faith formed by charity (Aquinas on John 6:29)}, weakest at \emph{disputed}:
\begin{itemize}
\item The believing of John 6:29 is faith living through charity, the source of works --- \cite[cap. 6, lect. 3, n. 901]{aquinas-commentary-john-2012}; \cite[Trent, Session VI (1547), Decree on Justification, ch. 7]{tanner-decrees-1990}
\end{itemize}
\item \emph{Luther on Galatians 3:11–12 (against faith formed by charity)}, weakest at \emph{disputed}:
\begin{itemize}
\item an inference step --- \cite[ad loc. Gal 3:12]{luther-commentary-galatians-1998}; \cite[ad loc. Gal 3:11]{luther-commentary-galatians-1998}
\end{itemize}
\item \emph{Chrysostom on John 3:36 (believing is not enough without a right life)}, weakest at \emph{disputed}:
\begin{itemize}
\item an inference step --- \cite[Homily 31, §1]{chrysostom-homilies-john}; \cite{stanley-did-jesus-teach-2006}
\end{itemize}
\item \emph{Obedience as sanctification (John 3:36, Matthew 7:21)}, weakest at \emph{disputed}:
\begin{itemize}
\item The obedience of John 3:36 and Matthew 7:21 is sanctification, the fruit of faith --- \cite[ad loc. John 3:36]{ridderbos-john-1997}; \cite[ad loc. John 3:36]{carson-john-1991}; \cite[ad loc. Matt 7:21]{morris-matthew-1992}; \cite[55--55]{schreiner-justification-by-works-2015}; \cite{hoekema-saved-by-grace-1994}; \cite{murray-redemption-1984}; \cite[XVI.2]{westminster-confession-1647}; \cite[XIII.1]{westminster-confession-1647}; \cite[ad loc. John 3:36]{calvin-commentary-john-1847}
\item Justification and sanctification are inseparable but distinct --- \cite[Q. 77]{westminster-larger-catechism-1647}; 1 Cor 6:11; 1 Cor 1:30; Rom 4:6; Rom 4:8; Ezek 36:27; Heb 9:13-14; Rom 3:24-25; Rom 6:6; Rom 6:14; \cite[XI.1]{westminster-confession-1647}; \cite[III.xi.6]{calvin-institutes-1960}
\item an inference step --- \cite[55--55]{schreiner-justification-by-works-2015}; \cite[ad loc. John 3:36]{calvin-commentary-john-1847}; \cite[XVI.2]{westminster-confession-1647}; \cite[XIII.1]{westminster-confession-1647}
\end{itemize}
\end{itemize}
\par\medskip\noindent\textbf{\texttt{trent\_\allowbreak{}not\_\allowbreak{}forced}}\par\nobreak\noindent
\textbf{Nor is Trent forced.} A consistent set of positions that answers every attack holds Paul’s case with Luke’s and Peter’s, and none of them can be held with Trent.

{\small\begin{verbatim}
theorem trent_not_forced : ¬Framework.SkepticallyAccepted
    solaFideDispute.defeats Party.trent
-- axioms: propext, Classical.choice, Quot.sound
\end{verbatim}}
\par\medskip\noindent\textbf{\texttt{apocalyptic\allowbreak{}Stands\allowbreak{}With\allowbreak{}Luke}}\par\nobreak\noindent
Why the apocalyptic reading can be defended: it can be held with Luke’s case, Peter’s and the critics’, and it answers every position that defeats it. Those are Paul’s case, John’s, Trent, and Trent taken as a claim about Paul’s word.

\paragraph{\emph{Apocalyptic reading, against faith as the condition} is accepted on some resolution.}
\begin{itemize}
\item A position holding \emph{Luke 7:50: “your faith has saved you”}, \emph{Sola fide from Peter (Acts 15:9–11)}, \emph{Apocalyptic reading, against faith as the condition}, \emph{Variegated nomism (Gathercole, Carson et al.)} can be held.
\item None of \emph{Luke 7:50: “your faith has saved you”}, \emph{Sola fide from Peter (Acts 15:9–11)}, \emph{Apocalyptic reading, against faith as the condition}, \emph{Variegated nomism (Gathercole, Carson et al.)} defeats another, and each attack on them is answered from among them:
\begin{itemize}
\item \emph{Trent, against ‘not by works’} defeats \emph{Luke 7:50: “your faith has saved you”}, and \emph{Luke 7:50: “your faith has saved you”} defeats \emph{Trent, against ‘not by works’} (\texttt{why\allowbreak{}The\allowbreak{}Dominical\allowbreak{}Case\allowbreak{}Stands\allowbreak{}Against\allowbreak{}Trent}).
\item \emph{Trent, against ‘not by works’, read as what Paul’s word means} defeats \emph{Luke 7:50: “your faith has saved you”}, and \emph{Luke 7:50: “your faith has saved you”} defeats \emph{Trent, against ‘not by works’, read as what Paul’s word means}.
\item \emph{Chrysostom on John 3:36 (believing is not enough without a right life)} defeats \emph{Luke 7:50: “your faith has saved you”}, and \emph{Luke 7:50: “your faith has saved you”} defeats \emph{Chrysostom on John 3:36 (believing is not enough without a right life)}.
\item \emph{Trent, against ‘not by works’} defeats \emph{Sola fide from Peter (Acts 15:9–11)}, and \emph{Luke 7:50: “your faith has saved you”} defeats \emph{Trent, against ‘not by works’} (\texttt{why\allowbreak{}The\allowbreak{}Dominical\allowbreak{}Case\allowbreak{}Stands\allowbreak{}Against\allowbreak{}Trent}).
\item \emph{Law-observant Luke (Jervell)} defeats \emph{Sola fide from Peter (Acts 15:9–11)}, and \emph{Sola fide from Peter (Acts 15:9–11)} defeats \emph{Law-observant Luke (Jervell)}.
\item \emph{Trent, against ‘not by works’, read as what Paul’s word means} defeats \emph{Sola fide from Peter (Acts 15:9–11)}, and \emph{Luke 7:50: “your faith has saved you”} defeats \emph{Trent, against ‘not by works’, read as what Paul’s word means}.
\item \emph{Chrysostom on John 3:36 (believing is not enough without a right life)} defeats \emph{Sola fide from Peter (Acts 15:9–11)}, and \emph{Luke 7:50: “your faith has saved you”} defeats \emph{Chrysostom on John 3:36 (believing is not enough without a right life)}.
\item \emph{Sola fide from Paul (Galatians 2:16)} defeats \emph{Apocalyptic reading, against faith as the condition}, and \emph{Apocalyptic reading, against faith as the condition} defeats \emph{Sola fide from Paul (Galatians 2:16)}.
\item \emph{Trent, against ‘not by works’} defeats \emph{Apocalyptic reading, against faith as the condition}, and \emph{Luke 7:50: “your faith has saved you”} defeats \emph{Trent, against ‘not by works’} (\texttt{why\allowbreak{}The\allowbreak{}Dominical\allowbreak{}Case\allowbreak{}Stands\allowbreak{}Against\allowbreak{}Trent}).
\item \emph{Trent, against ‘not by works’, read as what Paul’s word means} defeats \emph{Apocalyptic reading, against faith as the condition}, and \emph{Luke 7:50: “your faith has saved you”} defeats \emph{Trent, against ‘not by works’, read as what Paul’s word means}.
\item \emph{Sola fide from John (6:28–29)} defeats \emph{Apocalyptic reading, against faith as the condition}, and \emph{Apocalyptic reading, against faith as the condition} defeats \emph{Sola fide from John (6:28–29)}.
\item \emph{Covenantal nomism (Sanders, Dunn)} defeats \emph{Variegated nomism (Gathercole, Carson et al.)}, and \emph{Variegated nomism (Gathercole, Carson et al.)} defeats \emph{Covenantal nomism (Sanders, Dunn)}.
\end{itemize}
\end{itemize}
\par\noindent\textbf{What this rests on.} The reasons state 19 defeats and 56 absences of defeat, each a cell of the defeat table, \texttt{sola\allowbreak{}Fide\allowbreak{}Dispute\_\allowbreak{}defeats}, computed from the two parties' premises and checked by the kernel. Whether an attack survives turns on the attacker's weakest link and on the rating of what it attacks, so a changed rating can change the verdict. The weakest links:
\begin{itemize}
\item \emph{Luke 7:50: “your faith has saved you”}, weakest at \emph{disputed}:
\begin{itemize}
\item Luke 7:47 — her love is the evidence of her forgiveness, not its ground --- \cite[ad loc. Luke 7:47]{marshall-luke-1978}; \cite[54--54]{padilla-narrative-criticism-2021}; \cite[Apology IV, on love and the fulfilling of the law: Luke 7:47–50]{kolb-wengert-book-of-concord-2000}; Luke 7:41-43; Luke 7:47; Luke 7:47-50
\item an inference step --- \cite[Trent, Session VI (1547), canon 9]{tanner-decrees-1990}
\end{itemize}
\item \emph{Sola fide from Peter (Acts 15:9–11)}, weakest at \emph{disputed}:
\begin{itemize}
\item The yoke refused at Acts 15:10 is the whole law as a condition of salvation --- \cite[ad loc. Acts 15:10]{bruce-acts-1988}; Acts 15:10; Gal 5:1; Jas 2:10; \cite{das-paul-law-covenant-2001}
\item Good works are the fruit and evidence of saving faith, not its ground --- \cite[XI.2; XVI.2; XVI.5]{westminster-confession-1647}; \cite[III.xvi.1]{calvin-institutes-1960}; Jas 2:18-22; Eph 2:10; Rom 6:22; Luke 17:10
\item Justification and sanctification are inseparable but distinct --- \cite[Q. 77]{westminster-larger-catechism-1647}; 1 Cor 6:11; 1 Cor 1:30; Rom 4:6; Rom 4:8; Ezek 36:27; Heb 9:13-14; Rom 3:24-25; Rom 6:6; Rom 6:14; \cite[XI.1]{westminster-confession-1647}; \cite[III.xi.6]{calvin-institutes-1960}
\item an inference step --- \cite[Trent, Session VI (1547), canon 9]{tanner-decrees-1990}
\end{itemize}
\item \emph{Apocalyptic reading, against faith as the condition}, weakest at \emph{disputed}:
\begin{itemize}
\item ‘Faith of Christ’ is Christ’s own faithfulness, the means of God’s rescue (encoded as a denial, so it ranks disputed) --- \cite{campbell-deliverance-god-2009}; \cite[ad loc. Gal 2:16]{martyn-galatians-1997}; \cite{hays-faith-jesus-christ-2002}
\item ‘The righteousness of God’ (δικαιοσύνη θεοῦ) names God’s act of rescue in Christ --- \cite{campbell-deliverance-god-2009}; \cite{martyn-galatians-1997}; Rom 3:21-26
\end{itemize}
\item \emph{Variegated nomism (Gathercole, Carson et al.)}, weakest at \emph{disputed}:
\begin{itemize}
\item Second Temple Judaism also held final vindication according to works (encoded as a denial, so it ranks disputed) --- \cite{gathercole-where-boasting-2002}; \cite{carson-variegated-nomism-1-2001}
\item an inference step --- \cite{gathercole-where-boasting-2002}; \cite{carson-variegated-nomism-1-2001}
\end{itemize}
\item \emph{Sola fide from Paul (Galatians 2:16)}, weakest at \emph{disputed}:
\begin{itemize}
\item In Galatians 2:16 ‘faith of Christ’ (πίστις Χριστοῦ) means our faith in Christ --- \cite{dunn-once-more-pistis-christou-1991}; \cite{matlock-detheologizing-2000}; Gal 2:16; Rom 3:22; Phil 3:9
\item Second Temple Judaism: in by grace, staying in by works (encoded as a denial, so it ranks disputed) --- \cite{sanders-paul-palestinian-judaism-1977}
\item Good works are the fruit and evidence of saving faith, not its ground --- \cite[XI.2; XVI.2; XVI.5]{westminster-confession-1647}; \cite[III.xvi.1]{calvin-institutes-1960}; Jas 2:18-22; Eph 2:10; Rom 6:22; Luke 17:10
\item Justification and sanctification are inseparable but distinct --- \cite[Q. 77]{westminster-larger-catechism-1647}; 1 Cor 6:11; 1 Cor 1:30; Rom 4:6; Rom 4:8; Ezek 36:27; Heb 9:13-14; Rom 3:24-25; Rom 6:6; Rom 6:14; \cite[XI.1]{westminster-confession-1647}; \cite[III.xi.6]{calvin-institutes-1960}
\item an inference step --- \cite[Trent, Session VI (1547), canon 9]{tanner-decrees-1990}
\item an inference step --- \cite{gathercole-where-boasting-2002}; \cite{carson-variegated-nomism-1-2001}
\end{itemize}
\item \emph{Trent, against ‘not by works’}, weakest at \emph{disputed}:
\begin{itemize}
\item Justification is not remission of sins only, but renewal of the inward man --- \cite[Trent, Session VI (1547), Decree on Justification, ch. 7]{tanner-decrees-1990}; \cite[Trent, Session VI (1547), canon 11]{tanner-decrees-1990}
\end{itemize}
\item \emph{Covenantal nomism (Sanders, Dunn)}, weakest at \emph{disputed}:
\begin{itemize}
\item Second Temple Judaism: in by grace, staying in by works --- \cite{sanders-paul-palestinian-judaism-1977}
\item an inference step --- \cite{dunn-new-perspective-2005}; \cite{wright-what-paul-said-1997}
\end{itemize}
\item \emph{Law-observant Luke (Jervell)}, weakest at \emph{disputed}:
\begin{itemize}
\item an inference step --- \cite{jervell-luke-people-god-1972}
\end{itemize}
\item \emph{Paul’s gospel (Galatians 1), against justification as renewal}, weakest at \emph{disputed}:
\begin{itemize}
\item an inference step --- \cite[ad loc. Gal 1:6]{moo-galatians-2013}; \cite[XI.1]{westminster-confession-1647}
\item an inference step --- \cite[XI.1]{westminster-confession-1647}; \cite[III.xi.2]{calvin-institutes-1960}; \cite{morris-apostolic-preaching-1965}
\end{itemize}
\item \emph{Paul’s word (δικαιόω, by the rule of least meaning)}, weakest at \emph{well supported}:
\begin{itemize}
\item In Paul, ‘justify’ (δικαιόω) is a verdict: to declare righteous --- \cite[s.v. δικαιόω]{bauer-danker-lexicon-2000}; \cite{morris-apostolic-preaching-1965}; \cite[p. 80]{moo-romans-1996}; \cite{irons-righteousness-of-god-2015}; \cite{wright-justification-2009}; \cite[ad loc. Rom 3:24]{fitzmyer-romans-1993}
\item A word contributes the least meaning its context requires; one occurrence does not carry the whole concept --- \cite{silva-biblical-words-1994}; \cite{joos-semantic-axiom-1972}; \cite{barr-semantics-2004}
\item an inference step --- \cite{silva-biblical-words-1994}; \cite[p. 80]{moo-romans-1996}; \cite{barr-semantics-2004}
\end{itemize}
\item \emph{Trent, against ‘not by works’, read as what Paul’s word means}, weakest at \emph{disputed}:
\begin{itemize}
\item Justification is not remission of sins only, but renewal of the inward man --- \cite[Trent, Session VI (1547), Decree on Justification, ch. 7]{tanner-decrees-1990}; \cite[Trent, Session VI (1547), canon 11]{tanner-decrees-1990}
\item an inference step --- \cite[Trent, Session VI (1547), Decree on Justification, ch. 8]{tanner-decrees-1990}; \cite[Trent, Session VI (1547), Decree on Justification, ch. 7]{tanner-decrees-1990}; \cite[On the Spirit and the Letter, ch. 45 (26.45)]{npnf1-05-1887}; \cite{mcgrath-iustitia-dei-2005}
\end{itemize}
\item \emph{Romans 4:3–8, on Paul’s word (by the rule of least meaning)}, weakest at \emph{well supported}:
\begin{itemize}
\item Romans 4:4–5 contrasts wages owed with a gift reckoned to the one who believes --- Rom 4:4-5; \cite[III.xi.18]{calvin-institutes-1960}
\item ‘Count’ (λογίζομαι) in Romans 4:3–11 is the language of an account: credited to one --- \cite[s.v. λογίζομαι]{bauer-danker-lexicon-2000}; \cite[ad loc. Rom 4:3]{moo-romans-1996}
\item A word contributes the least meaning its context requires; one occurrence does not carry the whole concept --- \cite{silva-biblical-words-1994}; \cite{joos-semantic-axiom-1972}; \cite{barr-semantics-2004}
\item an inference step --- \cite[ad loc. Rom 4:6]{moo-romans-1996}; \cite[s.v. λογίζομαι]{bauer-danker-lexicon-2000}
\item an inference step --- \cite{silva-biblical-words-1994}; \cite[p. 80]{moo-romans-1996}; \cite{barr-semantics-2004}
\end{itemize}
\item \emph{Sola fide from John (6:28–29)}, weakest at \emph{disputed}:
\begin{itemize}
\item The believing of John 6:29 is trust, which brings nothing and receives Christ --- \cite[ad loc. John 6:29]{calvin-commentary-john-1847}; John 6:28-29
\item Good works are the fruit and evidence of saving faith, not its ground --- \cite[XI.2; XVI.2; XVI.5]{westminster-confession-1647}; \cite[III.xvi.1]{calvin-institutes-1960}; Jas 2:18-22; Eph 2:10; Rom 6:22; Luke 17:10
\item Justification and sanctification are inseparable but distinct --- \cite[Q. 77]{westminster-larger-catechism-1647}; 1 Cor 6:11; 1 Cor 1:30; Rom 4:6; Rom 4:8; Ezek 36:27; Heb 9:13-14; Rom 3:24-25; Rom 6:6; Rom 6:14; \cite[XI.1]{westminster-confession-1647}; \cite[III.xi.6]{calvin-institutes-1960}
\item an inference step --- \cite[Trent, Session VI (1547), canon 9]{tanner-decrees-1990}
\end{itemize}
\item \emph{Faith formed by charity (Aquinas on John 6:29)}, weakest at \emph{disputed}:
\begin{itemize}
\item The believing of John 6:29 is faith living through charity, the source of works --- \cite[cap. 6, lect. 3, n. 901]{aquinas-commentary-john-2012}; \cite[Trent, Session VI (1547), Decree on Justification, ch. 7]{tanner-decrees-1990}
\end{itemize}
\item \emph{Luther on Galatians 3:11–12 (against faith formed by charity)}, weakest at \emph{disputed}:
\begin{itemize}
\item an inference step --- \cite[ad loc. Gal 3:12]{luther-commentary-galatians-1998}; \cite[ad loc. Gal 3:11]{luther-commentary-galatians-1998}
\end{itemize}
\item \emph{Chrysostom on John 3:36 (believing is not enough without a right life)}, weakest at \emph{disputed}:
\begin{itemize}
\item an inference step --- \cite[Homily 31, §1]{chrysostom-homilies-john}; \cite{stanley-did-jesus-teach-2006}
\end{itemize}
\item \emph{Obedience as sanctification (John 3:36, Matthew 7:21)}, weakest at \emph{disputed}:
\begin{itemize}
\item The obedience of John 3:36 and Matthew 7:21 is sanctification, the fruit of faith --- \cite[ad loc. John 3:36]{ridderbos-john-1997}; \cite[ad loc. John 3:36]{carson-john-1991}; \cite[ad loc. Matt 7:21]{morris-matthew-1992}; \cite[55--55]{schreiner-justification-by-works-2015}; \cite{hoekema-saved-by-grace-1994}; \cite{murray-redemption-1984}; \cite[XVI.2]{westminster-confession-1647}; \cite[XIII.1]{westminster-confession-1647}; \cite[ad loc. John 3:36]{calvin-commentary-john-1847}
\item Justification and sanctification are inseparable but distinct --- \cite[Q. 77]{westminster-larger-catechism-1647}; 1 Cor 6:11; 1 Cor 1:30; Rom 4:6; Rom 4:8; Ezek 36:27; Heb 9:13-14; Rom 3:24-25; Rom 6:6; Rom 6:14; \cite[XI.1]{westminster-confession-1647}; \cite[III.xi.6]{calvin-institutes-1960}
\item an inference step --- \cite[55--55]{schreiner-justification-by-works-2015}; \cite[ad loc. John 3:36]{calvin-commentary-john-1847}; \cite[XVI.2]{westminster-confession-1647}; \cite[XIII.1]{westminster-confession-1647}
\end{itemize}
\end{itemize}
\noindent Why \emph{Luke 7:50: “your faith has saved you”} defeats \emph{Trent, against ‘not by works’}, at the level of the claims: \texttt{why\allowbreak{}The\allowbreak{}Dominical\allowbreak{}Case\allowbreak{}Stands\allowbreak{}Against\allowbreak{}Trent}.
\par\medskip\noindent\textbf{\texttt{apocalyptic\_\allowbreak{}defensible}}\par\nobreak\noindent
\textbf{The apocalyptic reading can be defended.} A consistent set of positions includes it and answers every attack, and that set holds Luke’s case and Peter’s too. They claim that faith saves, and the apocalyptic reading grants that, denying only that faith is the condition. Of the cases from Paul, Luke and Acts, only Paul’s defeats it, and it defeats Paul’s in return, over “faith of Christ”. Trent it defeats in return over “not by works”.

So the disagreement between the Reformed reading and the apocalyptic one lies where scholars place it: between Paul and Campbell, over “faith of Christ” (πίστις Χριστοῦ). It does not lie between Campbell and Luke.

{\small\begin{verbatim}
theorem apocalyptic_defensible : Framework.CredulouslyAccepted
    solaFideDispute.defeats Party.apocalyptic
-- axioms: propext, Classical.choice, Quot.sound
\end{verbatim}}
\par\medskip\noindent\textbf{\texttt{apocalyptic\allowbreak{}Answered\allowbreak{}By\allowbreak{}Paul}}\par\nobreak\noindent
Why the apocalyptic reading is not forced: a set of positions that can be defended includes Paul’s case, and Paul’s case defeats it.

\paragraph{\emph{Apocalyptic reading, against faith as the condition} is not accepted on every resolution.}
\begin{itemize}
\item A position holding \emph{Sola fide from Paul (Galatians 2:16)}, \emph{Luke 7:50: “your faith has saved you”}, \emph{Sola fide from Peter (Acts 15:9–11)}, \emph{Variegated nomism (Gathercole, Carson et al.)} can be held, and it holds \emph{Sola fide from Paul (Galatians 2:16)}; \emph{Sola fide from Paul (Galatians 2:16)} conflicts with \emph{Apocalyptic reading, against faith as the condition}: \emph{Sola fide from Paul (Galatians 2:16)} defeats \emph{Apocalyptic reading, against faith as the condition}, and \emph{Apocalyptic reading, against faith as the condition} defeats \emph{Sola fide from Paul (Galatians 2:16)}.
\item None of \emph{Sola fide from Paul (Galatians 2:16)}, \emph{Luke 7:50: “your faith has saved you”}, \emph{Sola fide from Peter (Acts 15:9–11)}, \emph{Variegated nomism (Gathercole, Carson et al.)} defeats another, and each attack on them is answered from among them:
\begin{itemize}
\item \emph{Trent, against ‘not by works’} defeats \emph{Sola fide from Paul (Galatians 2:16)}, and \emph{Sola fide from Paul (Galatians 2:16)} defeats \emph{Trent, against ‘not by works’}.
\item \emph{Apocalyptic reading, against faith as the condition} defeats \emph{Sola fide from Paul (Galatians 2:16)}, and \emph{Sola fide from Paul (Galatians 2:16)} defeats \emph{Apocalyptic reading, against faith as the condition}.
\item \emph{Covenantal nomism (Sanders, Dunn)} defeats \emph{Sola fide from Paul (Galatians 2:16)}, and \emph{Sola fide from Paul (Galatians 2:16)} defeats \emph{Covenantal nomism (Sanders, Dunn)}.
\item \emph{Trent, against ‘not by works’, read as what Paul’s word means} defeats \emph{Sola fide from Paul (Galatians 2:16)}, and \emph{Sola fide from Paul (Galatians 2:16)} defeats \emph{Trent, against ‘not by works’, read as what Paul’s word means}.
\item \emph{Chrysostom on John 3:36 (believing is not enough without a right life)} defeats \emph{Sola fide from Paul (Galatians 2:16)}, and \emph{Sola fide from Paul (Galatians 2:16)} defeats \emph{Chrysostom on John 3:36 (believing is not enough without a right life)}.
\item \emph{Trent, against ‘not by works’} defeats \emph{Luke 7:50: “your faith has saved you”}, and \emph{Sola fide from Paul (Galatians 2:16)} defeats \emph{Trent, against ‘not by works’}.
\item \emph{Trent, against ‘not by works’, read as what Paul’s word means} defeats \emph{Luke 7:50: “your faith has saved you”}, and \emph{Sola fide from Paul (Galatians 2:16)} defeats \emph{Trent, against ‘not by works’, read as what Paul’s word means}.
\item \emph{Chrysostom on John 3:36 (believing is not enough without a right life)} defeats \emph{Luke 7:50: “your faith has saved you”}, and \emph{Sola fide from Paul (Galatians 2:16)} defeats \emph{Chrysostom on John 3:36 (believing is not enough without a right life)}.
\item \emph{Trent, against ‘not by works’} defeats \emph{Sola fide from Peter (Acts 15:9–11)}, and \emph{Sola fide from Paul (Galatians 2:16)} defeats \emph{Trent, against ‘not by works’}.
\item \emph{Law-observant Luke (Jervell)} defeats \emph{Sola fide from Peter (Acts 15:9–11)}, and \emph{Sola fide from Peter (Acts 15:9–11)} defeats \emph{Law-observant Luke (Jervell)}.
\item \emph{Trent, against ‘not by works’, read as what Paul’s word means} defeats \emph{Sola fide from Peter (Acts 15:9–11)}, and \emph{Sola fide from Paul (Galatians 2:16)} defeats \emph{Trent, against ‘not by works’, read as what Paul’s word means}.
\item \emph{Chrysostom on John 3:36 (believing is not enough without a right life)} defeats \emph{Sola fide from Peter (Acts 15:9–11)}, and \emph{Sola fide from Paul (Galatians 2:16)} defeats \emph{Chrysostom on John 3:36 (believing is not enough without a right life)}.
\item \emph{Covenantal nomism (Sanders, Dunn)} defeats \emph{Variegated nomism (Gathercole, Carson et al.)}, and \emph{Sola fide from Paul (Galatians 2:16)} defeats \emph{Covenantal nomism (Sanders, Dunn)}.
\end{itemize}
\end{itemize}
\par\noindent\textbf{What this rests on.} The reasons state 19 defeats and 55 absences of defeat, each a cell of the defeat table, \texttt{sola\allowbreak{}Fide\allowbreak{}Dispute\_\allowbreak{}defeats}, computed from the two parties' premises and checked by the kernel. Whether an attack survives turns on the attacker's weakest link and on the rating of what it attacks, so a changed rating can change the verdict. The weakest links:
\begin{itemize}
\item \emph{Sola fide from Paul (Galatians 2:16)}, weakest at \emph{disputed}:
\begin{itemize}
\item In Galatians 2:16 ‘faith of Christ’ (πίστις Χριστοῦ) means our faith in Christ --- \cite{dunn-once-more-pistis-christou-1991}; \cite{matlock-detheologizing-2000}; Gal 2:16; Rom 3:22; Phil 3:9
\item Second Temple Judaism: in by grace, staying in by works (encoded as a denial, so it ranks disputed) --- \cite{sanders-paul-palestinian-judaism-1977}
\item Good works are the fruit and evidence of saving faith, not its ground --- \cite[XI.2; XVI.2; XVI.5]{westminster-confession-1647}; \cite[III.xvi.1]{calvin-institutes-1960}; Jas 2:18-22; Eph 2:10; Rom 6:22; Luke 17:10
\item Justification and sanctification are inseparable but distinct --- \cite[Q. 77]{westminster-larger-catechism-1647}; 1 Cor 6:11; 1 Cor 1:30; Rom 4:6; Rom 4:8; Ezek 36:27; Heb 9:13-14; Rom 3:24-25; Rom 6:6; Rom 6:14; \cite[XI.1]{westminster-confession-1647}; \cite[III.xi.6]{calvin-institutes-1960}
\item an inference step --- \cite[Trent, Session VI (1547), canon 9]{tanner-decrees-1990}
\item an inference step --- \cite{gathercole-where-boasting-2002}; \cite{carson-variegated-nomism-1-2001}
\end{itemize}
\item \emph{Apocalyptic reading, against faith as the condition}, weakest at \emph{disputed}:
\begin{itemize}
\item ‘Faith of Christ’ is Christ’s own faithfulness, the means of God’s rescue (encoded as a denial, so it ranks disputed) --- \cite{campbell-deliverance-god-2009}; \cite[ad loc. Gal 2:16]{martyn-galatians-1997}; \cite{hays-faith-jesus-christ-2002}
\item ‘The righteousness of God’ (δικαιοσύνη θεοῦ) names God’s act of rescue in Christ --- \cite{campbell-deliverance-god-2009}; \cite{martyn-galatians-1997}; Rom 3:21-26
\end{itemize}
\item \emph{Luke 7:50: “your faith has saved you”}, weakest at \emph{disputed}:
\begin{itemize}
\item Luke 7:47 — her love is the evidence of her forgiveness, not its ground --- \cite[ad loc. Luke 7:47]{marshall-luke-1978}; \cite[54--54]{padilla-narrative-criticism-2021}; \cite[Apology IV, on love and the fulfilling of the law: Luke 7:47–50]{kolb-wengert-book-of-concord-2000}; Luke 7:41-43; Luke 7:47; Luke 7:47-50
\item an inference step --- \cite[Trent, Session VI (1547), canon 9]{tanner-decrees-1990}
\end{itemize}
\item \emph{Sola fide from Peter (Acts 15:9–11)}, weakest at \emph{disputed}:
\begin{itemize}
\item The yoke refused at Acts 15:10 is the whole law as a condition of salvation --- \cite[ad loc. Acts 15:10]{bruce-acts-1988}; Acts 15:10; Gal 5:1; Jas 2:10; \cite{das-paul-law-covenant-2001}
\item Good works are the fruit and evidence of saving faith, not its ground --- \cite[XI.2; XVI.2; XVI.5]{westminster-confession-1647}; \cite[III.xvi.1]{calvin-institutes-1960}; Jas 2:18-22; Eph 2:10; Rom 6:22; Luke 17:10
\item Justification and sanctification are inseparable but distinct --- \cite[Q. 77]{westminster-larger-catechism-1647}; 1 Cor 6:11; 1 Cor 1:30; Rom 4:6; Rom 4:8; Ezek 36:27; Heb 9:13-14; Rom 3:24-25; Rom 6:6; Rom 6:14; \cite[XI.1]{westminster-confession-1647}; \cite[III.xi.6]{calvin-institutes-1960}
\item an inference step --- \cite[Trent, Session VI (1547), canon 9]{tanner-decrees-1990}
\end{itemize}
\item \emph{Variegated nomism (Gathercole, Carson et al.)}, weakest at \emph{disputed}:
\begin{itemize}
\item Second Temple Judaism also held final vindication according to works (encoded as a denial, so it ranks disputed) --- \cite{gathercole-where-boasting-2002}; \cite{carson-variegated-nomism-1-2001}
\item an inference step --- \cite{gathercole-where-boasting-2002}; \cite{carson-variegated-nomism-1-2001}
\end{itemize}
\item \emph{Trent, against ‘not by works’}, weakest at \emph{disputed}:
\begin{itemize}
\item Justification is not remission of sins only, but renewal of the inward man --- \cite[Trent, Session VI (1547), Decree on Justification, ch. 7]{tanner-decrees-1990}; \cite[Trent, Session VI (1547), canon 11]{tanner-decrees-1990}
\end{itemize}
\item \emph{Covenantal nomism (Sanders, Dunn)}, weakest at \emph{disputed}:
\begin{itemize}
\item Second Temple Judaism: in by grace, staying in by works --- \cite{sanders-paul-palestinian-judaism-1977}
\item an inference step --- \cite{dunn-new-perspective-2005}; \cite{wright-what-paul-said-1997}
\end{itemize}
\item \emph{Law-observant Luke (Jervell)}, weakest at \emph{disputed}:
\begin{itemize}
\item an inference step --- \cite{jervell-luke-people-god-1972}
\end{itemize}
\item \emph{Paul’s gospel (Galatians 1), against justification as renewal}, weakest at \emph{disputed}:
\begin{itemize}
\item an inference step --- \cite[ad loc. Gal 1:6]{moo-galatians-2013}; \cite[XI.1]{westminster-confession-1647}
\item an inference step --- \cite[XI.1]{westminster-confession-1647}; \cite[III.xi.2]{calvin-institutes-1960}; \cite{morris-apostolic-preaching-1965}
\end{itemize}
\item \emph{Paul’s word (δικαιόω, by the rule of least meaning)}, weakest at \emph{well supported}:
\begin{itemize}
\item In Paul, ‘justify’ (δικαιόω) is a verdict: to declare righteous --- \cite[s.v. δικαιόω]{bauer-danker-lexicon-2000}; \cite{morris-apostolic-preaching-1965}; \cite[p. 80]{moo-romans-1996}; \cite{irons-righteousness-of-god-2015}; \cite{wright-justification-2009}; \cite[ad loc. Rom 3:24]{fitzmyer-romans-1993}
\item A word contributes the least meaning its context requires; one occurrence does not carry the whole concept --- \cite{silva-biblical-words-1994}; \cite{joos-semantic-axiom-1972}; \cite{barr-semantics-2004}
\item an inference step --- \cite{silva-biblical-words-1994}; \cite[p. 80]{moo-romans-1996}; \cite{barr-semantics-2004}
\end{itemize}
\item \emph{Trent, against ‘not by works’, read as what Paul’s word means}, weakest at \emph{disputed}:
\begin{itemize}
\item Justification is not remission of sins only, but renewal of the inward man --- \cite[Trent, Session VI (1547), Decree on Justification, ch. 7]{tanner-decrees-1990}; \cite[Trent, Session VI (1547), canon 11]{tanner-decrees-1990}
\item an inference step --- \cite[Trent, Session VI (1547), Decree on Justification, ch. 8]{tanner-decrees-1990}; \cite[Trent, Session VI (1547), Decree on Justification, ch. 7]{tanner-decrees-1990}; \cite[On the Spirit and the Letter, ch. 45 (26.45)]{npnf1-05-1887}; \cite{mcgrath-iustitia-dei-2005}
\end{itemize}
\item \emph{Romans 4:3–8, on Paul’s word (by the rule of least meaning)}, weakest at \emph{well supported}:
\begin{itemize}
\item Romans 4:4–5 contrasts wages owed with a gift reckoned to the one who believes --- Rom 4:4-5; \cite[III.xi.18]{calvin-institutes-1960}
\item ‘Count’ (λογίζομαι) in Romans 4:3–11 is the language of an account: credited to one --- \cite[s.v. λογίζομαι]{bauer-danker-lexicon-2000}; \cite[ad loc. Rom 4:3]{moo-romans-1996}
\item A word contributes the least meaning its context requires; one occurrence does not carry the whole concept --- \cite{silva-biblical-words-1994}; \cite{joos-semantic-axiom-1972}; \cite{barr-semantics-2004}
\item an inference step --- \cite[ad loc. Rom 4:6]{moo-romans-1996}; \cite[s.v. λογίζομαι]{bauer-danker-lexicon-2000}
\item an inference step --- \cite{silva-biblical-words-1994}; \cite[p. 80]{moo-romans-1996}; \cite{barr-semantics-2004}
\end{itemize}
\item \emph{Sola fide from John (6:28–29)}, weakest at \emph{disputed}:
\begin{itemize}
\item The believing of John 6:29 is trust, which brings nothing and receives Christ --- \cite[ad loc. John 6:29]{calvin-commentary-john-1847}; John 6:28-29
\item Good works are the fruit and evidence of saving faith, not its ground --- \cite[XI.2; XVI.2; XVI.5]{westminster-confession-1647}; \cite[III.xvi.1]{calvin-institutes-1960}; Jas 2:18-22; Eph 2:10; Rom 6:22; Luke 17:10
\item Justification and sanctification are inseparable but distinct --- \cite[Q. 77]{westminster-larger-catechism-1647}; 1 Cor 6:11; 1 Cor 1:30; Rom 4:6; Rom 4:8; Ezek 36:27; Heb 9:13-14; Rom 3:24-25; Rom 6:6; Rom 6:14; \cite[XI.1]{westminster-confession-1647}; \cite[III.xi.6]{calvin-institutes-1960}
\item an inference step --- \cite[Trent, Session VI (1547), canon 9]{tanner-decrees-1990}
\end{itemize}
\item \emph{Faith formed by charity (Aquinas on John 6:29)}, weakest at \emph{disputed}:
\begin{itemize}
\item The believing of John 6:29 is faith living through charity, the source of works --- \cite[cap. 6, lect. 3, n. 901]{aquinas-commentary-john-2012}; \cite[Trent, Session VI (1547), Decree on Justification, ch. 7]{tanner-decrees-1990}
\end{itemize}
\item \emph{Luther on Galatians 3:11–12 (against faith formed by charity)}, weakest at \emph{disputed}:
\begin{itemize}
\item an inference step --- \cite[ad loc. Gal 3:12]{luther-commentary-galatians-1998}; \cite[ad loc. Gal 3:11]{luther-commentary-galatians-1998}
\end{itemize}
\item \emph{Chrysostom on John 3:36 (believing is not enough without a right life)}, weakest at \emph{disputed}:
\begin{itemize}
\item an inference step --- \cite[Homily 31, §1]{chrysostom-homilies-john}; \cite{stanley-did-jesus-teach-2006}
\end{itemize}
\item \emph{Obedience as sanctification (John 3:36, Matthew 7:21)}, weakest at \emph{disputed}:
\begin{itemize}
\item The obedience of John 3:36 and Matthew 7:21 is sanctification, the fruit of faith --- \cite[ad loc. John 3:36]{ridderbos-john-1997}; \cite[ad loc. John 3:36]{carson-john-1991}; \cite[ad loc. Matt 7:21]{morris-matthew-1992}; \cite[55--55]{schreiner-justification-by-works-2015}; \cite{hoekema-saved-by-grace-1994}; \cite{murray-redemption-1984}; \cite[XVI.2]{westminster-confession-1647}; \cite[XIII.1]{westminster-confession-1647}; \cite[ad loc. John 3:36]{calvin-commentary-john-1847}
\item Justification and sanctification are inseparable but distinct --- \cite[Q. 77]{westminster-larger-catechism-1647}; 1 Cor 6:11; 1 Cor 1:30; Rom 4:6; Rom 4:8; Ezek 36:27; Heb 9:13-14; Rom 3:24-25; Rom 6:6; Rom 6:14; \cite[XI.1]{westminster-confession-1647}; \cite[III.xi.6]{calvin-institutes-1960}
\item an inference step --- \cite[55--55]{schreiner-justification-by-works-2015}; \cite[ad loc. John 3:36]{calvin-commentary-john-1847}; \cite[XVI.2]{westminster-confession-1647}; \cite[XIII.1]{westminster-confession-1647}
\end{itemize}
\end{itemize}
\par\medskip\noindent\textbf{\texttt{apocalyptic\_\allowbreak{}not\_\allowbreak{}forced}}\par\nobreak\noindent
\textbf{Nor is it forced.} A consistent set of positions that answers every attack holds Paul’s case, faith alone from Galatians 2:16 read as faith \emph{in} Christ, and cannot hold the apocalyptic reading with it. Between Paul and Campbell the debate chooses neither: each can be defended, and neither is forced.

{\small\begin{verbatim}
theorem apocalyptic_not_forced : ¬Framework.SkepticallyAccepted
    solaFideDispute.defeats Party.apocalyptic
-- axioms: propext, Classical.choice, Quot.sound
\end{verbatim}}
\par\medskip\noindent\textbf{\texttt{luke\allowbreak{}Stands\allowbreak{}With\allowbreak{}Peter}}\par\nobreak\noindent
Why Luke’s case can be defended: it can be held with Peter’s, the apocalyptic reading and the critics’, and it answers every position that defeats it: Trent, Trent taken as a claim about Paul’s word, and Chrysostom.

\paragraph{\emph{Luke 7:50: “your faith has saved you”} is accepted on some resolution.}
\begin{itemize}
\item A position holding \emph{Luke 7:50: “your faith has saved you”}, \emph{Sola fide from Peter (Acts 15:9–11)}, \emph{Apocalyptic reading, against faith as the condition}, \emph{Variegated nomism (Gathercole, Carson et al.)} can be held.
\item None of \emph{Luke 7:50: “your faith has saved you”}, \emph{Sola fide from Peter (Acts 15:9–11)}, \emph{Apocalyptic reading, against faith as the condition}, \emph{Variegated nomism (Gathercole, Carson et al.)} defeats another, and each attack on them is answered from among them:
\begin{itemize}
\item \emph{Trent, against ‘not by works’} defeats \emph{Luke 7:50: “your faith has saved you”}, and \emph{Luke 7:50: “your faith has saved you”} defeats \emph{Trent, against ‘not by works’} (\texttt{why\allowbreak{}The\allowbreak{}Dominical\allowbreak{}Case\allowbreak{}Stands\allowbreak{}Against\allowbreak{}Trent}).
\item \emph{Trent, against ‘not by works’, read as what Paul’s word means} defeats \emph{Luke 7:50: “your faith has saved you”}, and \emph{Luke 7:50: “your faith has saved you”} defeats \emph{Trent, against ‘not by works’, read as what Paul’s word means}.
\item \emph{Chrysostom on John 3:36 (believing is not enough without a right life)} defeats \emph{Luke 7:50: “your faith has saved you”}, and \emph{Luke 7:50: “your faith has saved you”} defeats \emph{Chrysostom on John 3:36 (believing is not enough without a right life)}.
\item \emph{Trent, against ‘not by works’} defeats \emph{Sola fide from Peter (Acts 15:9–11)}, and \emph{Luke 7:50: “your faith has saved you”} defeats \emph{Trent, against ‘not by works’} (\texttt{why\allowbreak{}The\allowbreak{}Dominical\allowbreak{}Case\allowbreak{}Stands\allowbreak{}Against\allowbreak{}Trent}).
\item \emph{Law-observant Luke (Jervell)} defeats \emph{Sola fide from Peter (Acts 15:9–11)}, and \emph{Sola fide from Peter (Acts 15:9–11)} defeats \emph{Law-observant Luke (Jervell)}.
\item \emph{Trent, against ‘not by works’, read as what Paul’s word means} defeats \emph{Sola fide from Peter (Acts 15:9–11)}, and \emph{Luke 7:50: “your faith has saved you”} defeats \emph{Trent, against ‘not by works’, read as what Paul’s word means}.
\item \emph{Chrysostom on John 3:36 (believing is not enough without a right life)} defeats \emph{Sola fide from Peter (Acts 15:9–11)}, and \emph{Luke 7:50: “your faith has saved you”} defeats \emph{Chrysostom on John 3:36 (believing is not enough without a right life)}.
\item \emph{Sola fide from Paul (Galatians 2:16)} defeats \emph{Apocalyptic reading, against faith as the condition}, and \emph{Apocalyptic reading, against faith as the condition} defeats \emph{Sola fide from Paul (Galatians 2:16)}.
\item \emph{Trent, against ‘not by works’} defeats \emph{Apocalyptic reading, against faith as the condition}, and \emph{Luke 7:50: “your faith has saved you”} defeats \emph{Trent, against ‘not by works’} (\texttt{why\allowbreak{}The\allowbreak{}Dominical\allowbreak{}Case\allowbreak{}Stands\allowbreak{}Against\allowbreak{}Trent}).
\item \emph{Trent, against ‘not by works’, read as what Paul’s word means} defeats \emph{Apocalyptic reading, against faith as the condition}, and \emph{Luke 7:50: “your faith has saved you”} defeats \emph{Trent, against ‘not by works’, read as what Paul’s word means}.
\item \emph{Sola fide from John (6:28–29)} defeats \emph{Apocalyptic reading, against faith as the condition}, and \emph{Apocalyptic reading, against faith as the condition} defeats \emph{Sola fide from John (6:28–29)}.
\item \emph{Covenantal nomism (Sanders, Dunn)} defeats \emph{Variegated nomism (Gathercole, Carson et al.)}, and \emph{Variegated nomism (Gathercole, Carson et al.)} defeats \emph{Covenantal nomism (Sanders, Dunn)}.
\end{itemize}
\end{itemize}
\par\noindent\textbf{What this rests on.} The reasons state 19 defeats and 56 absences of defeat, each a cell of the defeat table, \texttt{sola\allowbreak{}Fide\allowbreak{}Dispute\_\allowbreak{}defeats}, computed from the two parties' premises and checked by the kernel. Whether an attack survives turns on the attacker's weakest link and on the rating of what it attacks, so a changed rating can change the verdict. The weakest links:
\begin{itemize}
\item \emph{Luke 7:50: “your faith has saved you”}, weakest at \emph{disputed}:
\begin{itemize}
\item Luke 7:47 — her love is the evidence of her forgiveness, not its ground --- \cite[ad loc. Luke 7:47]{marshall-luke-1978}; \cite[54--54]{padilla-narrative-criticism-2021}; \cite[Apology IV, on love and the fulfilling of the law: Luke 7:47–50]{kolb-wengert-book-of-concord-2000}; Luke 7:41-43; Luke 7:47; Luke 7:47-50
\item an inference step --- \cite[Trent, Session VI (1547), canon 9]{tanner-decrees-1990}
\end{itemize}
\item \emph{Sola fide from Peter (Acts 15:9–11)}, weakest at \emph{disputed}:
\begin{itemize}
\item The yoke refused at Acts 15:10 is the whole law as a condition of salvation --- \cite[ad loc. Acts 15:10]{bruce-acts-1988}; Acts 15:10; Gal 5:1; Jas 2:10; \cite{das-paul-law-covenant-2001}
\item Good works are the fruit and evidence of saving faith, not its ground --- \cite[XI.2; XVI.2; XVI.5]{westminster-confession-1647}; \cite[III.xvi.1]{calvin-institutes-1960}; Jas 2:18-22; Eph 2:10; Rom 6:22; Luke 17:10
\item Justification and sanctification are inseparable but distinct --- \cite[Q. 77]{westminster-larger-catechism-1647}; 1 Cor 6:11; 1 Cor 1:30; Rom 4:6; Rom 4:8; Ezek 36:27; Heb 9:13-14; Rom 3:24-25; Rom 6:6; Rom 6:14; \cite[XI.1]{westminster-confession-1647}; \cite[III.xi.6]{calvin-institutes-1960}
\item an inference step --- \cite[Trent, Session VI (1547), canon 9]{tanner-decrees-1990}
\end{itemize}
\item \emph{Apocalyptic reading, against faith as the condition}, weakest at \emph{disputed}:
\begin{itemize}
\item ‘Faith of Christ’ is Christ’s own faithfulness, the means of God’s rescue (encoded as a denial, so it ranks disputed) --- \cite{campbell-deliverance-god-2009}; \cite[ad loc. Gal 2:16]{martyn-galatians-1997}; \cite{hays-faith-jesus-christ-2002}
\item ‘The righteousness of God’ (δικαιοσύνη θεοῦ) names God’s act of rescue in Christ --- \cite{campbell-deliverance-god-2009}; \cite{martyn-galatians-1997}; Rom 3:21-26
\end{itemize}
\item \emph{Variegated nomism (Gathercole, Carson et al.)}, weakest at \emph{disputed}:
\begin{itemize}
\item Second Temple Judaism also held final vindication according to works (encoded as a denial, so it ranks disputed) --- \cite{gathercole-where-boasting-2002}; \cite{carson-variegated-nomism-1-2001}
\item an inference step --- \cite{gathercole-where-boasting-2002}; \cite{carson-variegated-nomism-1-2001}
\end{itemize}
\item \emph{Sola fide from Paul (Galatians 2:16)}, weakest at \emph{disputed}:
\begin{itemize}
\item In Galatians 2:16 ‘faith of Christ’ (πίστις Χριστοῦ) means our faith in Christ --- \cite{dunn-once-more-pistis-christou-1991}; \cite{matlock-detheologizing-2000}; Gal 2:16; Rom 3:22; Phil 3:9
\item Second Temple Judaism: in by grace, staying in by works (encoded as a denial, so it ranks disputed) --- \cite{sanders-paul-palestinian-judaism-1977}
\item Good works are the fruit and evidence of saving faith, not its ground --- \cite[XI.2; XVI.2; XVI.5]{westminster-confession-1647}; \cite[III.xvi.1]{calvin-institutes-1960}; Jas 2:18-22; Eph 2:10; Rom 6:22; Luke 17:10
\item Justification and sanctification are inseparable but distinct --- \cite[Q. 77]{westminster-larger-catechism-1647}; 1 Cor 6:11; 1 Cor 1:30; Rom 4:6; Rom 4:8; Ezek 36:27; Heb 9:13-14; Rom 3:24-25; Rom 6:6; Rom 6:14; \cite[XI.1]{westminster-confession-1647}; \cite[III.xi.6]{calvin-institutes-1960}
\item an inference step --- \cite[Trent, Session VI (1547), canon 9]{tanner-decrees-1990}
\item an inference step --- \cite{gathercole-where-boasting-2002}; \cite{carson-variegated-nomism-1-2001}
\end{itemize}
\item \emph{Trent, against ‘not by works’}, weakest at \emph{disputed}:
\begin{itemize}
\item Justification is not remission of sins only, but renewal of the inward man --- \cite[Trent, Session VI (1547), Decree on Justification, ch. 7]{tanner-decrees-1990}; \cite[Trent, Session VI (1547), canon 11]{tanner-decrees-1990}
\end{itemize}
\item \emph{Covenantal nomism (Sanders, Dunn)}, weakest at \emph{disputed}:
\begin{itemize}
\item Second Temple Judaism: in by grace, staying in by works --- \cite{sanders-paul-palestinian-judaism-1977}
\item an inference step --- \cite{dunn-new-perspective-2005}; \cite{wright-what-paul-said-1997}
\end{itemize}
\item \emph{Law-observant Luke (Jervell)}, weakest at \emph{disputed}:
\begin{itemize}
\item an inference step --- \cite{jervell-luke-people-god-1972}
\end{itemize}
\item \emph{Paul’s gospel (Galatians 1), against justification as renewal}, weakest at \emph{disputed}:
\begin{itemize}
\item an inference step --- \cite[ad loc. Gal 1:6]{moo-galatians-2013}; \cite[XI.1]{westminster-confession-1647}
\item an inference step --- \cite[XI.1]{westminster-confession-1647}; \cite[III.xi.2]{calvin-institutes-1960}; \cite{morris-apostolic-preaching-1965}
\end{itemize}
\item \emph{Paul’s word (δικαιόω, by the rule of least meaning)}, weakest at \emph{well supported}:
\begin{itemize}
\item In Paul, ‘justify’ (δικαιόω) is a verdict: to declare righteous --- \cite[s.v. δικαιόω]{bauer-danker-lexicon-2000}; \cite{morris-apostolic-preaching-1965}; \cite[p. 80]{moo-romans-1996}; \cite{irons-righteousness-of-god-2015}; \cite{wright-justification-2009}; \cite[ad loc. Rom 3:24]{fitzmyer-romans-1993}
\item A word contributes the least meaning its context requires; one occurrence does not carry the whole concept --- \cite{silva-biblical-words-1994}; \cite{joos-semantic-axiom-1972}; \cite{barr-semantics-2004}
\item an inference step --- \cite{silva-biblical-words-1994}; \cite[p. 80]{moo-romans-1996}; \cite{barr-semantics-2004}
\end{itemize}
\item \emph{Trent, against ‘not by works’, read as what Paul’s word means}, weakest at \emph{disputed}:
\begin{itemize}
\item Justification is not remission of sins only, but renewal of the inward man --- \cite[Trent, Session VI (1547), Decree on Justification, ch. 7]{tanner-decrees-1990}; \cite[Trent, Session VI (1547), canon 11]{tanner-decrees-1990}
\item an inference step --- \cite[Trent, Session VI (1547), Decree on Justification, ch. 8]{tanner-decrees-1990}; \cite[Trent, Session VI (1547), Decree on Justification, ch. 7]{tanner-decrees-1990}; \cite[On the Spirit and the Letter, ch. 45 (26.45)]{npnf1-05-1887}; \cite{mcgrath-iustitia-dei-2005}
\end{itemize}
\item \emph{Romans 4:3–8, on Paul’s word (by the rule of least meaning)}, weakest at \emph{well supported}:
\begin{itemize}
\item Romans 4:4–5 contrasts wages owed with a gift reckoned to the one who believes --- Rom 4:4-5; \cite[III.xi.18]{calvin-institutes-1960}
\item ‘Count’ (λογίζομαι) in Romans 4:3–11 is the language of an account: credited to one --- \cite[s.v. λογίζομαι]{bauer-danker-lexicon-2000}; \cite[ad loc. Rom 4:3]{moo-romans-1996}
\item A word contributes the least meaning its context requires; one occurrence does not carry the whole concept --- \cite{silva-biblical-words-1994}; \cite{joos-semantic-axiom-1972}; \cite{barr-semantics-2004}
\item an inference step --- \cite[ad loc. Rom 4:6]{moo-romans-1996}; \cite[s.v. λογίζομαι]{bauer-danker-lexicon-2000}
\item an inference step --- \cite{silva-biblical-words-1994}; \cite[p. 80]{moo-romans-1996}; \cite{barr-semantics-2004}
\end{itemize}
\item \emph{Sola fide from John (6:28–29)}, weakest at \emph{disputed}:
\begin{itemize}
\item The believing of John 6:29 is trust, which brings nothing and receives Christ --- \cite[ad loc. John 6:29]{calvin-commentary-john-1847}; John 6:28-29
\item Good works are the fruit and evidence of saving faith, not its ground --- \cite[XI.2; XVI.2; XVI.5]{westminster-confession-1647}; \cite[III.xvi.1]{calvin-institutes-1960}; Jas 2:18-22; Eph 2:10; Rom 6:22; Luke 17:10
\item Justification and sanctification are inseparable but distinct --- \cite[Q. 77]{westminster-larger-catechism-1647}; 1 Cor 6:11; 1 Cor 1:30; Rom 4:6; Rom 4:8; Ezek 36:27; Heb 9:13-14; Rom 3:24-25; Rom 6:6; Rom 6:14; \cite[XI.1]{westminster-confession-1647}; \cite[III.xi.6]{calvin-institutes-1960}
\item an inference step --- \cite[Trent, Session VI (1547), canon 9]{tanner-decrees-1990}
\end{itemize}
\item \emph{Faith formed by charity (Aquinas on John 6:29)}, weakest at \emph{disputed}:
\begin{itemize}
\item The believing of John 6:29 is faith living through charity, the source of works --- \cite[cap. 6, lect. 3, n. 901]{aquinas-commentary-john-2012}; \cite[Trent, Session VI (1547), Decree on Justification, ch. 7]{tanner-decrees-1990}
\end{itemize}
\item \emph{Luther on Galatians 3:11–12 (against faith formed by charity)}, weakest at \emph{disputed}:
\begin{itemize}
\item an inference step --- \cite[ad loc. Gal 3:12]{luther-commentary-galatians-1998}; \cite[ad loc. Gal 3:11]{luther-commentary-galatians-1998}
\end{itemize}
\item \emph{Chrysostom on John 3:36 (believing is not enough without a right life)}, weakest at \emph{disputed}:
\begin{itemize}
\item an inference step --- \cite[Homily 31, §1]{chrysostom-homilies-john}; \cite{stanley-did-jesus-teach-2006}
\end{itemize}
\item \emph{Obedience as sanctification (John 3:36, Matthew 7:21)}, weakest at \emph{disputed}:
\begin{itemize}
\item The obedience of John 3:36 and Matthew 7:21 is sanctification, the fruit of faith --- \cite[ad loc. John 3:36]{ridderbos-john-1997}; \cite[ad loc. John 3:36]{carson-john-1991}; \cite[ad loc. Matt 7:21]{morris-matthew-1992}; \cite[55--55]{schreiner-justification-by-works-2015}; \cite{hoekema-saved-by-grace-1994}; \cite{murray-redemption-1984}; \cite[XVI.2]{westminster-confession-1647}; \cite[XIII.1]{westminster-confession-1647}; \cite[ad loc. John 3:36]{calvin-commentary-john-1847}
\item Justification and sanctification are inseparable but distinct --- \cite[Q. 77]{westminster-larger-catechism-1647}; 1 Cor 6:11; 1 Cor 1:30; Rom 4:6; Rom 4:8; Ezek 36:27; Heb 9:13-14; Rom 3:24-25; Rom 6:6; Rom 6:14; \cite[XI.1]{westminster-confession-1647}; \cite[III.xi.6]{calvin-institutes-1960}
\item an inference step --- \cite[55--55]{schreiner-justification-by-works-2015}; \cite[ad loc. John 3:36]{calvin-commentary-john-1847}; \cite[XVI.2]{westminster-confession-1647}; \cite[XIII.1]{westminster-confession-1647}
\end{itemize}
\end{itemize}
\noindent Why \emph{Luke 7:50: “your faith has saved you”} defeats \emph{Trent, against ‘not by works’}, at the level of the claims: \texttt{why\allowbreak{}The\allowbreak{}Dominical\allowbreak{}Case\allowbreak{}Stands\allowbreak{}Against\allowbreak{}Trent}.
\par\medskip\noindent\textbf{\texttt{dominical\_\allowbreak{}case\_\allowbreak{}defensible}}\par\nobreak\noindent
\textbf{Luke’s case can be defended.} “Your faith has saved you”, read as Luke gives it, can be held with every position but Trent, however Trent’s definition is taken, and Chrysostom. Those are its only attackers, and Luke’s case defeats each in return. Against Trent the point at issue is what justification is: Trent’s definition denies that faith without charity suffices, and Luke’s words say her faith saved her. Against Chrysostom it is whether the right life Jesus requires elsewhere (John 3:36, Matthew 7:21) is a second condition beside faith.

{\small\begin{verbatim}
theorem dominical_case_defensible : Framework.CredulouslyAccepted
    solaFideDispute.defeats Party.dominical
-- axioms: propext, Classical.choice, Quot.sound
\end{verbatim}}
\par\medskip\noindent\textbf{\texttt{luke\allowbreak{}Answered\allowbreak{}By\allowbreak{}Trent}}\par\nobreak\noindent
Why Luke’s case is not forced: a set of positions that can be defended includes Trent, and Trent defeats it.

\paragraph{\emph{Luke 7:50: “your faith has saved you”} is not accepted on every resolution.}
\begin{itemize}
\item A position holding \emph{Trent, against ‘not by works’}, \emph{Covenantal nomism (Sanders, Dunn)}, \emph{Law-observant Luke (Jervell)} can be held, and it holds \emph{Trent, against ‘not by works’}; \emph{Trent, against ‘not by works’} conflicts with \emph{Luke 7:50: “your faith has saved you”}: \emph{Trent, against ‘not by works’} defeats \emph{Luke 7:50: “your faith has saved you”}, and \emph{Luke 7:50: “your faith has saved you”} defeats \emph{Trent, against ‘not by works’} (\texttt{why\allowbreak{}The\allowbreak{}Dominical\allowbreak{}Case\allowbreak{}Stands\allowbreak{}Against\allowbreak{}Trent}).
\item None of \emph{Trent, against ‘not by works’}, \emph{Covenantal nomism (Sanders, Dunn)}, \emph{Law-observant Luke (Jervell)} defeats another, and each attack on them is answered from among them:
\begin{itemize}
\item \emph{Sola fide from Paul (Galatians 2:16)} defeats \emph{Trent, against ‘not by works’}, and \emph{Trent, against ‘not by works’} defeats \emph{Sola fide from Paul (Galatians 2:16)}.
\item \emph{Luke 7:50: “your faith has saved you”} defeats \emph{Trent, against ‘not by works’} (\texttt{why\allowbreak{}The\allowbreak{}Dominical\allowbreak{}Case\allowbreak{}Stands\allowbreak{}Against\allowbreak{}Trent}), and \emph{Trent, against ‘not by works’} defeats \emph{Luke 7:50: “your faith has saved you”}.
\item \emph{Sola fide from Peter (Acts 15:9–11)} defeats \emph{Trent, against ‘not by works’}, and \emph{Trent, against ‘not by works’} defeats \emph{Sola fide from Peter (Acts 15:9–11)}.
\item \emph{Apocalyptic reading, against faith as the condition} defeats \emph{Trent, against ‘not by works’}, and \emph{Trent, against ‘not by works’} defeats \emph{Apocalyptic reading, against faith as the condition}.
\item \emph{Paul’s gospel (Galatians 1), against justification as renewal} defeats \emph{Trent, against ‘not by works’} (\texttt{where\allowbreak{}Trent\allowbreak{}Parts\allowbreak{}From\allowbreak{}Paul}), and \emph{Trent, against ‘not by works’} defeats \emph{Paul’s gospel (Galatians 1), against justification as renewal}.
\item \emph{Sola fide from John (6:28–29)} defeats \emph{Trent, against ‘not by works’}, and \emph{Trent, against ‘not by works’} defeats \emph{Sola fide from John (6:28–29)}.
\item \emph{Obedience as sanctification (John 3:36, Matthew 7:21)} defeats \emph{Trent, against ‘not by works’}, and \emph{Trent, against ‘not by works’} defeats \emph{Obedience as sanctification (John 3:36, Matthew 7:21)}.
\item \emph{Sola fide from Paul (Galatians 2:16)} defeats \emph{Covenantal nomism (Sanders, Dunn)}, and \emph{Trent, against ‘not by works’} defeats \emph{Sola fide from Paul (Galatians 2:16)}.
\item \emph{Variegated nomism (Gathercole, Carson et al.)} defeats \emph{Covenantal nomism (Sanders, Dunn)}, and \emph{Covenantal nomism (Sanders, Dunn)} defeats \emph{Variegated nomism (Gathercole, Carson et al.)}.
\item \emph{Sola fide from Peter (Acts 15:9–11)} defeats \emph{Law-observant Luke (Jervell)}, and \emph{Trent, against ‘not by works’} defeats \emph{Sola fide from Peter (Acts 15:9–11)}.
\end{itemize}
\end{itemize}
\par\noindent\textbf{What this rests on.} The reasons state 18 defeats and 41 absences of defeat, each a cell of the defeat table, \texttt{sola\allowbreak{}Fide\allowbreak{}Dispute\_\allowbreak{}defeats}, computed from the two parties' premises and checked by the kernel. Whether an attack survives turns on the attacker's weakest link and on the rating of what it attacks, so a changed rating can change the verdict. The weakest links:
\begin{itemize}
\item \emph{Trent, against ‘not by works’}, weakest at \emph{disputed}:
\begin{itemize}
\item Justification is not remission of sins only, but renewal of the inward man --- \cite[Trent, Session VI (1547), Decree on Justification, ch. 7]{tanner-decrees-1990}; \cite[Trent, Session VI (1547), canon 11]{tanner-decrees-1990}
\end{itemize}
\item \emph{Luke 7:50: “your faith has saved you”}, weakest at \emph{disputed}:
\begin{itemize}
\item Luke 7:47 — her love is the evidence of her forgiveness, not its ground --- \cite[ad loc. Luke 7:47]{marshall-luke-1978}; \cite[54--54]{padilla-narrative-criticism-2021}; \cite[Apology IV, on love and the fulfilling of the law: Luke 7:47–50]{kolb-wengert-book-of-concord-2000}; Luke 7:41-43; Luke 7:47; Luke 7:47-50
\item an inference step --- \cite[Trent, Session VI (1547), canon 9]{tanner-decrees-1990}
\end{itemize}
\item \emph{Covenantal nomism (Sanders, Dunn)}, weakest at \emph{disputed}:
\begin{itemize}
\item Second Temple Judaism: in by grace, staying in by works --- \cite{sanders-paul-palestinian-judaism-1977}
\item an inference step --- \cite{dunn-new-perspective-2005}; \cite{wright-what-paul-said-1997}
\end{itemize}
\item \emph{Law-observant Luke (Jervell)}, weakest at \emph{disputed}:
\begin{itemize}
\item an inference step --- \cite{jervell-luke-people-god-1972}
\end{itemize}
\item \emph{Sola fide from Paul (Galatians 2:16)}, weakest at \emph{disputed}:
\begin{itemize}
\item In Galatians 2:16 ‘faith of Christ’ (πίστις Χριστοῦ) means our faith in Christ --- \cite{dunn-once-more-pistis-christou-1991}; \cite{matlock-detheologizing-2000}; Gal 2:16; Rom 3:22; Phil 3:9
\item Second Temple Judaism: in by grace, staying in by works (encoded as a denial, so it ranks disputed) --- \cite{sanders-paul-palestinian-judaism-1977}
\item Good works are the fruit and evidence of saving faith, not its ground --- \cite[XI.2; XVI.2; XVI.5]{westminster-confession-1647}; \cite[III.xvi.1]{calvin-institutes-1960}; Jas 2:18-22; Eph 2:10; Rom 6:22; Luke 17:10
\item Justification and sanctification are inseparable but distinct --- \cite[Q. 77]{westminster-larger-catechism-1647}; 1 Cor 6:11; 1 Cor 1:30; Rom 4:6; Rom 4:8; Ezek 36:27; Heb 9:13-14; Rom 3:24-25; Rom 6:6; Rom 6:14; \cite[XI.1]{westminster-confession-1647}; \cite[III.xi.6]{calvin-institutes-1960}
\item an inference step --- \cite[Trent, Session VI (1547), canon 9]{tanner-decrees-1990}
\item an inference step --- \cite{gathercole-where-boasting-2002}; \cite{carson-variegated-nomism-1-2001}
\end{itemize}
\item \emph{Sola fide from Peter (Acts 15:9–11)}, weakest at \emph{disputed}:
\begin{itemize}
\item The yoke refused at Acts 15:10 is the whole law as a condition of salvation --- \cite[ad loc. Acts 15:10]{bruce-acts-1988}; Acts 15:10; Gal 5:1; Jas 2:10; \cite{das-paul-law-covenant-2001}
\item Good works are the fruit and evidence of saving faith, not its ground --- \cite[XI.2; XVI.2; XVI.5]{westminster-confession-1647}; \cite[III.xvi.1]{calvin-institutes-1960}; Jas 2:18-22; Eph 2:10; Rom 6:22; Luke 17:10
\item Justification and sanctification are inseparable but distinct --- \cite[Q. 77]{westminster-larger-catechism-1647}; 1 Cor 6:11; 1 Cor 1:30; Rom 4:6; Rom 4:8; Ezek 36:27; Heb 9:13-14; Rom 3:24-25; Rom 6:6; Rom 6:14; \cite[XI.1]{westminster-confession-1647}; \cite[III.xi.6]{calvin-institutes-1960}
\item an inference step --- \cite[Trent, Session VI (1547), canon 9]{tanner-decrees-1990}
\end{itemize}
\item \emph{Apocalyptic reading, against faith as the condition}, weakest at \emph{disputed}:
\begin{itemize}
\item ‘Faith of Christ’ is Christ’s own faithfulness, the means of God’s rescue (encoded as a denial, so it ranks disputed) --- \cite{campbell-deliverance-god-2009}; \cite[ad loc. Gal 2:16]{martyn-galatians-1997}; \cite{hays-faith-jesus-christ-2002}
\item ‘The righteousness of God’ (δικαιοσύνη θεοῦ) names God’s act of rescue in Christ --- \cite{campbell-deliverance-god-2009}; \cite{martyn-galatians-1997}; Rom 3:21-26
\end{itemize}
\item \emph{Variegated nomism (Gathercole, Carson et al.)}, weakest at \emph{disputed}:
\begin{itemize}
\item Second Temple Judaism also held final vindication according to works (encoded as a denial, so it ranks disputed) --- \cite{gathercole-where-boasting-2002}; \cite{carson-variegated-nomism-1-2001}
\item an inference step --- \cite{gathercole-where-boasting-2002}; \cite{carson-variegated-nomism-1-2001}
\end{itemize}
\item \emph{Paul’s gospel (Galatians 1), against justification as renewal}, weakest at \emph{disputed}:
\begin{itemize}
\item an inference step --- \cite[ad loc. Gal 1:6]{moo-galatians-2013}; \cite[XI.1]{westminster-confession-1647}
\item an inference step --- \cite[XI.1]{westminster-confession-1647}; \cite[III.xi.2]{calvin-institutes-1960}; \cite{morris-apostolic-preaching-1965}
\end{itemize}
\item \emph{Paul’s word (δικαιόω, by the rule of least meaning)}, weakest at \emph{well supported}:
\begin{itemize}
\item In Paul, ‘justify’ (δικαιόω) is a verdict: to declare righteous --- \cite[s.v. δικαιόω]{bauer-danker-lexicon-2000}; \cite{morris-apostolic-preaching-1965}; \cite[p. 80]{moo-romans-1996}; \cite{irons-righteousness-of-god-2015}; \cite{wright-justification-2009}; \cite[ad loc. Rom 3:24]{fitzmyer-romans-1993}
\item A word contributes the least meaning its context requires; one occurrence does not carry the whole concept --- \cite{silva-biblical-words-1994}; \cite{joos-semantic-axiom-1972}; \cite{barr-semantics-2004}
\item an inference step --- \cite{silva-biblical-words-1994}; \cite[p. 80]{moo-romans-1996}; \cite{barr-semantics-2004}
\end{itemize}
\item \emph{Trent, against ‘not by works’, read as what Paul’s word means}, weakest at \emph{disputed}:
\begin{itemize}
\item Justification is not remission of sins only, but renewal of the inward man --- \cite[Trent, Session VI (1547), Decree on Justification, ch. 7]{tanner-decrees-1990}; \cite[Trent, Session VI (1547), canon 11]{tanner-decrees-1990}
\item an inference step --- \cite[Trent, Session VI (1547), Decree on Justification, ch. 8]{tanner-decrees-1990}; \cite[Trent, Session VI (1547), Decree on Justification, ch. 7]{tanner-decrees-1990}; \cite[On the Spirit and the Letter, ch. 45 (26.45)]{npnf1-05-1887}; \cite{mcgrath-iustitia-dei-2005}
\end{itemize}
\item \emph{Romans 4:3–8, on Paul’s word (by the rule of least meaning)}, weakest at \emph{well supported}:
\begin{itemize}
\item Romans 4:4–5 contrasts wages owed with a gift reckoned to the one who believes --- Rom 4:4-5; \cite[III.xi.18]{calvin-institutes-1960}
\item ‘Count’ (λογίζομαι) in Romans 4:3–11 is the language of an account: credited to one --- \cite[s.v. λογίζομαι]{bauer-danker-lexicon-2000}; \cite[ad loc. Rom 4:3]{moo-romans-1996}
\item A word contributes the least meaning its context requires; one occurrence does not carry the whole concept --- \cite{silva-biblical-words-1994}; \cite{joos-semantic-axiom-1972}; \cite{barr-semantics-2004}
\item an inference step --- \cite[ad loc. Rom 4:6]{moo-romans-1996}; \cite[s.v. λογίζομαι]{bauer-danker-lexicon-2000}
\item an inference step --- \cite{silva-biblical-words-1994}; \cite[p. 80]{moo-romans-1996}; \cite{barr-semantics-2004}
\end{itemize}
\item \emph{Sola fide from John (6:28–29)}, weakest at \emph{disputed}:
\begin{itemize}
\item The believing of John 6:29 is trust, which brings nothing and receives Christ --- \cite[ad loc. John 6:29]{calvin-commentary-john-1847}; John 6:28-29
\item Good works are the fruit and evidence of saving faith, not its ground --- \cite[XI.2; XVI.2; XVI.5]{westminster-confession-1647}; \cite[III.xvi.1]{calvin-institutes-1960}; Jas 2:18-22; Eph 2:10; Rom 6:22; Luke 17:10
\item Justification and sanctification are inseparable but distinct --- \cite[Q. 77]{westminster-larger-catechism-1647}; 1 Cor 6:11; 1 Cor 1:30; Rom 4:6; Rom 4:8; Ezek 36:27; Heb 9:13-14; Rom 3:24-25; Rom 6:6; Rom 6:14; \cite[XI.1]{westminster-confession-1647}; \cite[III.xi.6]{calvin-institutes-1960}
\item an inference step --- \cite[Trent, Session VI (1547), canon 9]{tanner-decrees-1990}
\end{itemize}
\item \emph{Faith formed by charity (Aquinas on John 6:29)}, weakest at \emph{disputed}:
\begin{itemize}
\item The believing of John 6:29 is faith living through charity, the source of works --- \cite[cap. 6, lect. 3, n. 901]{aquinas-commentary-john-2012}; \cite[Trent, Session VI (1547), Decree on Justification, ch. 7]{tanner-decrees-1990}
\end{itemize}
\item \emph{Luther on Galatians 3:11–12 (against faith formed by charity)}, weakest at \emph{disputed}:
\begin{itemize}
\item an inference step --- \cite[ad loc. Gal 3:12]{luther-commentary-galatians-1998}; \cite[ad loc. Gal 3:11]{luther-commentary-galatians-1998}
\end{itemize}
\item \emph{Chrysostom on John 3:36 (believing is not enough without a right life)}, weakest at \emph{disputed}:
\begin{itemize}
\item an inference step --- \cite[Homily 31, §1]{chrysostom-homilies-john}; \cite{stanley-did-jesus-teach-2006}
\end{itemize}
\item \emph{Obedience as sanctification (John 3:36, Matthew 7:21)}, weakest at \emph{disputed}:
\begin{itemize}
\item The obedience of John 3:36 and Matthew 7:21 is sanctification, the fruit of faith --- \cite[ad loc. John 3:36]{ridderbos-john-1997}; \cite[ad loc. John 3:36]{carson-john-1991}; \cite[ad loc. Matt 7:21]{morris-matthew-1992}; \cite[55--55]{schreiner-justification-by-works-2015}; \cite{hoekema-saved-by-grace-1994}; \cite{murray-redemption-1984}; \cite[XVI.2]{westminster-confession-1647}; \cite[XIII.1]{westminster-confession-1647}; \cite[ad loc. John 3:36]{calvin-commentary-john-1847}
\item Justification and sanctification are inseparable but distinct --- \cite[Q. 77]{westminster-larger-catechism-1647}; 1 Cor 6:11; 1 Cor 1:30; Rom 4:6; Rom 4:8; Ezek 36:27; Heb 9:13-14; Rom 3:24-25; Rom 6:6; Rom 6:14; \cite[XI.1]{westminster-confession-1647}; \cite[III.xi.6]{calvin-institutes-1960}
\item an inference step --- \cite[55--55]{schreiner-justification-by-works-2015}; \cite[ad loc. John 3:36]{calvin-commentary-john-1847}; \cite[XVI.2]{westminster-confession-1647}; \cite[XIII.1]{westminster-confession-1647}
\end{itemize}
\end{itemize}
\noindent Why \emph{Luke 7:50: “your faith has saved you”} defeats \emph{Trent, against ‘not by works’}, at the level of the claims: \texttt{why\allowbreak{}The\allowbreak{}Dominical\allowbreak{}Case\allowbreak{}Stands\allowbreak{}Against\allowbreak{}Trent}.
\noindent Why \emph{Paul’s gospel (Galatians 1), against justification as renewal} defeats \emph{Trent, against ‘not by works’}, at the level of the claims: \texttt{where\allowbreak{}Trent\allowbreak{}Parts\allowbreak{}From\allowbreak{}Paul}.
\par\medskip\noindent\textbf{\texttt{dominical\_\allowbreak{}case\_\allowbreak{}not\_\allowbreak{}forced}}\par\nobreak\noindent
\textbf{Luke’s case is not forced.} A consistent set of positions that answers every attack holds Trent, and cannot hold Luke’s case with it.

Between Trent and each Reformed strand, the debate chooses neither. Their weakest links are all \texttt{disputed}, so no rating breaks a tie, and every attack between them runs in both directions.

What would decide between them is what justification \emph{is}. No rating here settles that. \texttt{why\allowbreak{}The\allowbreak{}Dominical\allowbreak{}Case\allowbreak{}Stands\allowbreak{}Against\allowbreak{}Trent} places Luke’s disagreement with Trent at Trent’s definition. \texttt{where\allowbreak{}Trent\allowbreak{}Parts\allowbreak{}From\allowbreak{}Paul}, in the chapter “Where Trent parts from Paul”, places Paul’s at one step: that a verdict on Christ’s finished work excludes the renewal worked in us. The one thing the ratings do decide is that Trent cannot claim its definition as the meaning of Paul’s word (\texttt{trent\_\allowbreak{}on\_\allowbreak{}pauls\_\allowbreak{}word\_\allowbreak{}indefensible}).

{\small\begin{verbatim}
theorem dominical_case_not_forced : ¬Framework.SkepticallyAccepted
    solaFideDispute.defeats Party.dominical
-- axioms: propext, Classical.choice, Quot.sound
\end{verbatim}}
\par\medskip\noindent\textbf{\texttt{gospel\allowbreak{}Stands\allowbreak{}With\allowbreak{}The\allowbreak{}Strands}}\par\nobreak\noindent
Why Paul’s gospel can be defended: it can be held with Paul’s case, Luke’s, Peter’s, the critics’ and the argument from Paul’s word, and it answers both positions that defeat it. Those are Trent, and Trent taken as a claim about Paul’s word.

\paragraph{\emph{Paul’s gospel (Galatians 1), against justification as renewal} is accepted on some resolution.}
\begin{itemize}
\item A position holding \emph{Paul’s gospel (Galatians 1), against justification as renewal}, \emph{Sola fide from Paul (Galatians 2:16)}, \emph{Luke 7:50: “your faith has saved you”}, \emph{Sola fide from Peter (Acts 15:9–11)}, \emph{Variegated nomism (Gathercole, Carson et al.)}, \emph{Paul’s word (δικαιόω, by the rule of least meaning)} can be held.
\item None of \emph{Paul’s gospel (Galatians 1), against justification as renewal}, \emph{Sola fide from Paul (Galatians 2:16)}, \emph{Luke 7:50: “your faith has saved you”}, \emph{Sola fide from Peter (Acts 15:9–11)}, \emph{Variegated nomism (Gathercole, Carson et al.)}, \emph{Paul’s word (δικαιόω, by the rule of least meaning)} defeats another, and each attack on them is answered from among them:
\begin{itemize}
\item \emph{Trent, against ‘not by works’} defeats \emph{Paul’s gospel (Galatians 1), against justification as renewal}, and \emph{Paul’s gospel (Galatians 1), against justification as renewal} defeats \emph{Trent, against ‘not by works’} (\texttt{where\allowbreak{}Trent\allowbreak{}Parts\allowbreak{}From\allowbreak{}Paul}).
\item \emph{Trent, against ‘not by works’, read as what Paul’s word means} defeats \emph{Paul’s gospel (Galatians 1), against justification as renewal}, and \emph{Paul’s gospel (Galatians 1), against justification as renewal} defeats \emph{Trent, against ‘not by works’, read as what Paul’s word means}.
\item \emph{Trent, against ‘not by works’} defeats \emph{Sola fide from Paul (Galatians 2:16)}, and \emph{Paul’s gospel (Galatians 1), against justification as renewal} defeats \emph{Trent, against ‘not by works’} (\texttt{where\allowbreak{}Trent\allowbreak{}Parts\allowbreak{}From\allowbreak{}Paul}).
\item \emph{Apocalyptic reading, against faith as the condition} defeats \emph{Sola fide from Paul (Galatians 2:16)}, and \emph{Sola fide from Paul (Galatians 2:16)} defeats \emph{Apocalyptic reading, against faith as the condition}.
\item \emph{Covenantal nomism (Sanders, Dunn)} defeats \emph{Sola fide from Paul (Galatians 2:16)}, and \emph{Sola fide from Paul (Galatians 2:16)} defeats \emph{Covenantal nomism (Sanders, Dunn)}.
\item \emph{Trent, against ‘not by works’, read as what Paul’s word means} defeats \emph{Sola fide from Paul (Galatians 2:16)}, and \emph{Paul’s gospel (Galatians 1), against justification as renewal} defeats \emph{Trent, against ‘not by works’, read as what Paul’s word means}.
\item \emph{Chrysostom on John 3:36 (believing is not enough without a right life)} defeats \emph{Sola fide from Paul (Galatians 2:16)}, and \emph{Sola fide from Paul (Galatians 2:16)} defeats \emph{Chrysostom on John 3:36 (believing is not enough without a right life)}.
\item \emph{Trent, against ‘not by works’} defeats \emph{Luke 7:50: “your faith has saved you”}, and \emph{Paul’s gospel (Galatians 1), against justification as renewal} defeats \emph{Trent, against ‘not by works’} (\texttt{where\allowbreak{}Trent\allowbreak{}Parts\allowbreak{}From\allowbreak{}Paul}).
\item \emph{Trent, against ‘not by works’, read as what Paul’s word means} defeats \emph{Luke 7:50: “your faith has saved you”}, and \emph{Paul’s gospel (Galatians 1), against justification as renewal} defeats \emph{Trent, against ‘not by works’, read as what Paul’s word means}.
\item \emph{Chrysostom on John 3:36 (believing is not enough without a right life)} defeats \emph{Luke 7:50: “your faith has saved you”}, and \emph{Sola fide from Paul (Galatians 2:16)} defeats \emph{Chrysostom on John 3:36 (believing is not enough without a right life)}.
\item \emph{Trent, against ‘not by works’} defeats \emph{Sola fide from Peter (Acts 15:9–11)}, and \emph{Paul’s gospel (Galatians 1), against justification as renewal} defeats \emph{Trent, against ‘not by works’} (\texttt{where\allowbreak{}Trent\allowbreak{}Parts\allowbreak{}From\allowbreak{}Paul}).
\item \emph{Law-observant Luke (Jervell)} defeats \emph{Sola fide from Peter (Acts 15:9–11)}, and \emph{Sola fide from Peter (Acts 15:9–11)} defeats \emph{Law-observant Luke (Jervell)}.
\item \emph{Trent, against ‘not by works’, read as what Paul’s word means} defeats \emph{Sola fide from Peter (Acts 15:9–11)}, and \emph{Paul’s gospel (Galatians 1), against justification as renewal} defeats \emph{Trent, against ‘not by works’, read as what Paul’s word means}.
\item \emph{Chrysostom on John 3:36 (believing is not enough without a right life)} defeats \emph{Sola fide from Peter (Acts 15:9–11)}, and \emph{Sola fide from Paul (Galatians 2:16)} defeats \emph{Chrysostom on John 3:36 (believing is not enough without a right life)}.
\item \emph{Covenantal nomism (Sanders, Dunn)} defeats \emph{Variegated nomism (Gathercole, Carson et al.)}, and \emph{Sola fide from Paul (Galatians 2:16)} defeats \emph{Covenantal nomism (Sanders, Dunn)}.
\end{itemize}
\end{itemize}
\par\noindent\textbf{What this rests on.} The reasons state 21 defeats and 87 absences of defeat, each a cell of the defeat table, \texttt{sola\allowbreak{}Fide\allowbreak{}Dispute\_\allowbreak{}defeats}, computed from the two parties' premises and checked by the kernel. Whether an attack survives turns on the attacker's weakest link and on the rating of what it attacks, so a changed rating can change the verdict. The weakest links:
\begin{itemize}
\item \emph{Paul’s gospel (Galatians 1), against justification as renewal}, weakest at \emph{disputed}:
\begin{itemize}
\item an inference step --- \cite[ad loc. Gal 1:6]{moo-galatians-2013}; \cite[XI.1]{westminster-confession-1647}
\item an inference step --- \cite[XI.1]{westminster-confession-1647}; \cite[III.xi.2]{calvin-institutes-1960}; \cite{morris-apostolic-preaching-1965}
\end{itemize}
\item \emph{Sola fide from Paul (Galatians 2:16)}, weakest at \emph{disputed}:
\begin{itemize}
\item In Galatians 2:16 ‘faith of Christ’ (πίστις Χριστοῦ) means our faith in Christ --- \cite{dunn-once-more-pistis-christou-1991}; \cite{matlock-detheologizing-2000}; Gal 2:16; Rom 3:22; Phil 3:9
\item Second Temple Judaism: in by grace, staying in by works (encoded as a denial, so it ranks disputed) --- \cite{sanders-paul-palestinian-judaism-1977}
\item Good works are the fruit and evidence of saving faith, not its ground --- \cite[XI.2; XVI.2; XVI.5]{westminster-confession-1647}; \cite[III.xvi.1]{calvin-institutes-1960}; Jas 2:18-22; Eph 2:10; Rom 6:22; Luke 17:10
\item Justification and sanctification are inseparable but distinct --- \cite[Q. 77]{westminster-larger-catechism-1647}; 1 Cor 6:11; 1 Cor 1:30; Rom 4:6; Rom 4:8; Ezek 36:27; Heb 9:13-14; Rom 3:24-25; Rom 6:6; Rom 6:14; \cite[XI.1]{westminster-confession-1647}; \cite[III.xi.6]{calvin-institutes-1960}
\item an inference step --- \cite[Trent, Session VI (1547), canon 9]{tanner-decrees-1990}
\item an inference step --- \cite{gathercole-where-boasting-2002}; \cite{carson-variegated-nomism-1-2001}
\end{itemize}
\item \emph{Luke 7:50: “your faith has saved you”}, weakest at \emph{disputed}:
\begin{itemize}
\item Luke 7:47 — her love is the evidence of her forgiveness, not its ground --- \cite[ad loc. Luke 7:47]{marshall-luke-1978}; \cite[54--54]{padilla-narrative-criticism-2021}; \cite[Apology IV, on love and the fulfilling of the law: Luke 7:47–50]{kolb-wengert-book-of-concord-2000}; Luke 7:41-43; Luke 7:47; Luke 7:47-50
\item an inference step --- \cite[Trent, Session VI (1547), canon 9]{tanner-decrees-1990}
\end{itemize}
\item \emph{Sola fide from Peter (Acts 15:9–11)}, weakest at \emph{disputed}:
\begin{itemize}
\item The yoke refused at Acts 15:10 is the whole law as a condition of salvation --- \cite[ad loc. Acts 15:10]{bruce-acts-1988}; Acts 15:10; Gal 5:1; Jas 2:10; \cite{das-paul-law-covenant-2001}
\item Good works are the fruit and evidence of saving faith, not its ground --- \cite[XI.2; XVI.2; XVI.5]{westminster-confession-1647}; \cite[III.xvi.1]{calvin-institutes-1960}; Jas 2:18-22; Eph 2:10; Rom 6:22; Luke 17:10
\item Justification and sanctification are inseparable but distinct --- \cite[Q. 77]{westminster-larger-catechism-1647}; 1 Cor 6:11; 1 Cor 1:30; Rom 4:6; Rom 4:8; Ezek 36:27; Heb 9:13-14; Rom 3:24-25; Rom 6:6; Rom 6:14; \cite[XI.1]{westminster-confession-1647}; \cite[III.xi.6]{calvin-institutes-1960}
\item an inference step --- \cite[Trent, Session VI (1547), canon 9]{tanner-decrees-1990}
\end{itemize}
\item \emph{Variegated nomism (Gathercole, Carson et al.)}, weakest at \emph{disputed}:
\begin{itemize}
\item Second Temple Judaism also held final vindication according to works (encoded as a denial, so it ranks disputed) --- \cite{gathercole-where-boasting-2002}; \cite{carson-variegated-nomism-1-2001}
\item an inference step --- \cite{gathercole-where-boasting-2002}; \cite{carson-variegated-nomism-1-2001}
\end{itemize}
\item \emph{Paul’s word (δικαιόω, by the rule of least meaning)}, weakest at \emph{well supported}:
\begin{itemize}
\item In Paul, ‘justify’ (δικαιόω) is a verdict: to declare righteous --- \cite[s.v. δικαιόω]{bauer-danker-lexicon-2000}; \cite{morris-apostolic-preaching-1965}; \cite[p. 80]{moo-romans-1996}; \cite{irons-righteousness-of-god-2015}; \cite{wright-justification-2009}; \cite[ad loc. Rom 3:24]{fitzmyer-romans-1993}
\item A word contributes the least meaning its context requires; one occurrence does not carry the whole concept --- \cite{silva-biblical-words-1994}; \cite{joos-semantic-axiom-1972}; \cite{barr-semantics-2004}
\item an inference step --- \cite{silva-biblical-words-1994}; \cite[p. 80]{moo-romans-1996}; \cite{barr-semantics-2004}
\end{itemize}
\item \emph{Trent, against ‘not by works’}, weakest at \emph{disputed}:
\begin{itemize}
\item Justification is not remission of sins only, but renewal of the inward man --- \cite[Trent, Session VI (1547), Decree on Justification, ch. 7]{tanner-decrees-1990}; \cite[Trent, Session VI (1547), canon 11]{tanner-decrees-1990}
\end{itemize}
\item \emph{Apocalyptic reading, against faith as the condition}, weakest at \emph{disputed}:
\begin{itemize}
\item ‘Faith of Christ’ is Christ’s own faithfulness, the means of God’s rescue (encoded as a denial, so it ranks disputed) --- \cite{campbell-deliverance-god-2009}; \cite[ad loc. Gal 2:16]{martyn-galatians-1997}; \cite{hays-faith-jesus-christ-2002}
\item ‘The righteousness of God’ (δικαιοσύνη θεοῦ) names God’s act of rescue in Christ --- \cite{campbell-deliverance-god-2009}; \cite{martyn-galatians-1997}; Rom 3:21-26
\end{itemize}
\item \emph{Covenantal nomism (Sanders, Dunn)}, weakest at \emph{disputed}:
\begin{itemize}
\item Second Temple Judaism: in by grace, staying in by works --- \cite{sanders-paul-palestinian-judaism-1977}
\item an inference step --- \cite{dunn-new-perspective-2005}; \cite{wright-what-paul-said-1997}
\end{itemize}
\item \emph{Law-observant Luke (Jervell)}, weakest at \emph{disputed}:
\begin{itemize}
\item an inference step --- \cite{jervell-luke-people-god-1972}
\end{itemize}
\item \emph{Trent, against ‘not by works’, read as what Paul’s word means}, weakest at \emph{disputed}:
\begin{itemize}
\item Justification is not remission of sins only, but renewal of the inward man --- \cite[Trent, Session VI (1547), Decree on Justification, ch. 7]{tanner-decrees-1990}; \cite[Trent, Session VI (1547), canon 11]{tanner-decrees-1990}
\item an inference step --- \cite[Trent, Session VI (1547), Decree on Justification, ch. 8]{tanner-decrees-1990}; \cite[Trent, Session VI (1547), Decree on Justification, ch. 7]{tanner-decrees-1990}; \cite[On the Spirit and the Letter, ch. 45 (26.45)]{npnf1-05-1887}; \cite{mcgrath-iustitia-dei-2005}
\end{itemize}
\item \emph{Romans 4:3–8, on Paul’s word (by the rule of least meaning)}, weakest at \emph{well supported}:
\begin{itemize}
\item Romans 4:4–5 contrasts wages owed with a gift reckoned to the one who believes --- Rom 4:4-5; \cite[III.xi.18]{calvin-institutes-1960}
\item ‘Count’ (λογίζομαι) in Romans 4:3–11 is the language of an account: credited to one --- \cite[s.v. λογίζομαι]{bauer-danker-lexicon-2000}; \cite[ad loc. Rom 4:3]{moo-romans-1996}
\item A word contributes the least meaning its context requires; one occurrence does not carry the whole concept --- \cite{silva-biblical-words-1994}; \cite{joos-semantic-axiom-1972}; \cite{barr-semantics-2004}
\item an inference step --- \cite[ad loc. Rom 4:6]{moo-romans-1996}; \cite[s.v. λογίζομαι]{bauer-danker-lexicon-2000}
\item an inference step --- \cite{silva-biblical-words-1994}; \cite[p. 80]{moo-romans-1996}; \cite{barr-semantics-2004}
\end{itemize}
\item \emph{Sola fide from John (6:28–29)}, weakest at \emph{disputed}:
\begin{itemize}
\item The believing of John 6:29 is trust, which brings nothing and receives Christ --- \cite[ad loc. John 6:29]{calvin-commentary-john-1847}; John 6:28-29
\item Good works are the fruit and evidence of saving faith, not its ground --- \cite[XI.2; XVI.2; XVI.5]{westminster-confession-1647}; \cite[III.xvi.1]{calvin-institutes-1960}; Jas 2:18-22; Eph 2:10; Rom 6:22; Luke 17:10
\item Justification and sanctification are inseparable but distinct --- \cite[Q. 77]{westminster-larger-catechism-1647}; 1 Cor 6:11; 1 Cor 1:30; Rom 4:6; Rom 4:8; Ezek 36:27; Heb 9:13-14; Rom 3:24-25; Rom 6:6; Rom 6:14; \cite[XI.1]{westminster-confession-1647}; \cite[III.xi.6]{calvin-institutes-1960}
\item an inference step --- \cite[Trent, Session VI (1547), canon 9]{tanner-decrees-1990}
\end{itemize}
\item \emph{Faith formed by charity (Aquinas on John 6:29)}, weakest at \emph{disputed}:
\begin{itemize}
\item The believing of John 6:29 is faith living through charity, the source of works --- \cite[cap. 6, lect. 3, n. 901]{aquinas-commentary-john-2012}; \cite[Trent, Session VI (1547), Decree on Justification, ch. 7]{tanner-decrees-1990}
\end{itemize}
\item \emph{Luther on Galatians 3:11–12 (against faith formed by charity)}, weakest at \emph{disputed}:
\begin{itemize}
\item an inference step --- \cite[ad loc. Gal 3:12]{luther-commentary-galatians-1998}; \cite[ad loc. Gal 3:11]{luther-commentary-galatians-1998}
\end{itemize}
\item \emph{Chrysostom on John 3:36 (believing is not enough without a right life)}, weakest at \emph{disputed}:
\begin{itemize}
\item an inference step --- \cite[Homily 31, §1]{chrysostom-homilies-john}; \cite{stanley-did-jesus-teach-2006}
\end{itemize}
\item \emph{Obedience as sanctification (John 3:36, Matthew 7:21)}, weakest at \emph{disputed}:
\begin{itemize}
\item The obedience of John 3:36 and Matthew 7:21 is sanctification, the fruit of faith --- \cite[ad loc. John 3:36]{ridderbos-john-1997}; \cite[ad loc. John 3:36]{carson-john-1991}; \cite[ad loc. Matt 7:21]{morris-matthew-1992}; \cite[55--55]{schreiner-justification-by-works-2015}; \cite{hoekema-saved-by-grace-1994}; \cite{murray-redemption-1984}; \cite[XVI.2]{westminster-confession-1647}; \cite[XIII.1]{westminster-confession-1647}; \cite[ad loc. John 3:36]{calvin-commentary-john-1847}
\item Justification and sanctification are inseparable but distinct --- \cite[Q. 77]{westminster-larger-catechism-1647}; 1 Cor 6:11; 1 Cor 1:30; Rom 4:6; Rom 4:8; Ezek 36:27; Heb 9:13-14; Rom 3:24-25; Rom 6:6; Rom 6:14; \cite[XI.1]{westminster-confession-1647}; \cite[III.xi.6]{calvin-institutes-1960}
\item an inference step --- \cite[55--55]{schreiner-justification-by-works-2015}; \cite[ad loc. John 3:36]{calvin-commentary-john-1847}; \cite[XVI.2]{westminster-confession-1647}; \cite[XIII.1]{westminster-confession-1647}
\end{itemize}
\end{itemize}
\noindent Why \emph{Paul’s gospel (Galatians 1), against justification as renewal} defeats \emph{Trent, against ‘not by works’}, at the level of the claims: \texttt{where\allowbreak{}Trent\allowbreak{}Parts\allowbreak{}From\allowbreak{}Paul}.
\par\medskip\noindent\textbf{\texttt{gospel\_\allowbreak{}defensible}}\par\nobreak\noindent
\textbf{Paul’s gospel can be defended}, together with Paul’s, Luke’s and Peter’s cases and with the argument from Paul’s word. A consistent set of positions that answers every attack holds them all.

{\small\begin{verbatim}
theorem gospel_defensible : Framework.CredulouslyAccepted
    solaFideDispute.defeats Party.gospel
-- axioms: propext, Classical.choice, Quot.sound
\end{verbatim}}
\par\medskip\noindent\textbf{\texttt{gospel\allowbreak{}Answered\allowbreak{}By\allowbreak{}Trent}}\par\nobreak\noindent
Why Paul’s gospel is not forced: a set of positions that can be defended includes Trent, and Trent defeats it.

\paragraph{\emph{Paul’s gospel (Galatians 1), against justification as renewal} is not accepted on every resolution.}
\begin{itemize}
\item A position holding \emph{Trent, against ‘not by works’}, \emph{Covenantal nomism (Sanders, Dunn)}, \emph{Law-observant Luke (Jervell)} can be held, and it holds \emph{Trent, against ‘not by works’}; \emph{Trent, against ‘not by works’} conflicts with \emph{Paul’s gospel (Galatians 1), against justification as renewal}: \emph{Trent, against ‘not by works’} defeats \emph{Paul’s gospel (Galatians 1), against justification as renewal}, and \emph{Paul’s gospel (Galatians 1), against justification as renewal} defeats \emph{Trent, against ‘not by works’} (\texttt{where\allowbreak{}Trent\allowbreak{}Parts\allowbreak{}From\allowbreak{}Paul}).
\item None of \emph{Trent, against ‘not by works’}, \emph{Covenantal nomism (Sanders, Dunn)}, \emph{Law-observant Luke (Jervell)} defeats another, and each attack on them is answered from among them:
\begin{itemize}
\item \emph{Sola fide from Paul (Galatians 2:16)} defeats \emph{Trent, against ‘not by works’}, and \emph{Trent, against ‘not by works’} defeats \emph{Sola fide from Paul (Galatians 2:16)}.
\item \emph{Luke 7:50: “your faith has saved you”} defeats \emph{Trent, against ‘not by works’} (\texttt{why\allowbreak{}The\allowbreak{}Dominical\allowbreak{}Case\allowbreak{}Stands\allowbreak{}Against\allowbreak{}Trent}), and \emph{Trent, against ‘not by works’} defeats \emph{Luke 7:50: “your faith has saved you”}.
\item \emph{Sola fide from Peter (Acts 15:9–11)} defeats \emph{Trent, against ‘not by works’}, and \emph{Trent, against ‘not by works’} defeats \emph{Sola fide from Peter (Acts 15:9–11)}.
\item \emph{Apocalyptic reading, against faith as the condition} defeats \emph{Trent, against ‘not by works’}, and \emph{Trent, against ‘not by works’} defeats \emph{Apocalyptic reading, against faith as the condition}.
\item \emph{Paul’s gospel (Galatians 1), against justification as renewal} defeats \emph{Trent, against ‘not by works’} (\texttt{where\allowbreak{}Trent\allowbreak{}Parts\allowbreak{}From\allowbreak{}Paul}), and \emph{Trent, against ‘not by works’} defeats \emph{Paul’s gospel (Galatians 1), against justification as renewal}.
\item \emph{Sola fide from John (6:28–29)} defeats \emph{Trent, against ‘not by works’}, and \emph{Trent, against ‘not by works’} defeats \emph{Sola fide from John (6:28–29)}.
\item \emph{Obedience as sanctification (John 3:36, Matthew 7:21)} defeats \emph{Trent, against ‘not by works’}, and \emph{Trent, against ‘not by works’} defeats \emph{Obedience as sanctification (John 3:36, Matthew 7:21)}.
\item \emph{Sola fide from Paul (Galatians 2:16)} defeats \emph{Covenantal nomism (Sanders, Dunn)}, and \emph{Trent, against ‘not by works’} defeats \emph{Sola fide from Paul (Galatians 2:16)}.
\item \emph{Variegated nomism (Gathercole, Carson et al.)} defeats \emph{Covenantal nomism (Sanders, Dunn)}, and \emph{Covenantal nomism (Sanders, Dunn)} defeats \emph{Variegated nomism (Gathercole, Carson et al.)}.
\item \emph{Sola fide from Peter (Acts 15:9–11)} defeats \emph{Law-observant Luke (Jervell)}, and \emph{Trent, against ‘not by works’} defeats \emph{Sola fide from Peter (Acts 15:9–11)}.
\end{itemize}
\end{itemize}
\par\noindent\textbf{What this rests on.} The reasons state 18 defeats and 41 absences of defeat, each a cell of the defeat table, \texttt{sola\allowbreak{}Fide\allowbreak{}Dispute\_\allowbreak{}defeats}, computed from the two parties' premises and checked by the kernel. Whether an attack survives turns on the attacker's weakest link and on the rating of what it attacks, so a changed rating can change the verdict. The weakest links:
\begin{itemize}
\item \emph{Trent, against ‘not by works’}, weakest at \emph{disputed}:
\begin{itemize}
\item Justification is not remission of sins only, but renewal of the inward man --- \cite[Trent, Session VI (1547), Decree on Justification, ch. 7]{tanner-decrees-1990}; \cite[Trent, Session VI (1547), canon 11]{tanner-decrees-1990}
\end{itemize}
\item \emph{Paul’s gospel (Galatians 1), against justification as renewal}, weakest at \emph{disputed}:
\begin{itemize}
\item an inference step --- \cite[ad loc. Gal 1:6]{moo-galatians-2013}; \cite[XI.1]{westminster-confession-1647}
\item an inference step --- \cite[XI.1]{westminster-confession-1647}; \cite[III.xi.2]{calvin-institutes-1960}; \cite{morris-apostolic-preaching-1965}
\end{itemize}
\item \emph{Covenantal nomism (Sanders, Dunn)}, weakest at \emph{disputed}:
\begin{itemize}
\item Second Temple Judaism: in by grace, staying in by works --- \cite{sanders-paul-palestinian-judaism-1977}
\item an inference step --- \cite{dunn-new-perspective-2005}; \cite{wright-what-paul-said-1997}
\end{itemize}
\item \emph{Law-observant Luke (Jervell)}, weakest at \emph{disputed}:
\begin{itemize}
\item an inference step --- \cite{jervell-luke-people-god-1972}
\end{itemize}
\item \emph{Sola fide from Paul (Galatians 2:16)}, weakest at \emph{disputed}:
\begin{itemize}
\item In Galatians 2:16 ‘faith of Christ’ (πίστις Χριστοῦ) means our faith in Christ --- \cite{dunn-once-more-pistis-christou-1991}; \cite{matlock-detheologizing-2000}; Gal 2:16; Rom 3:22; Phil 3:9
\item Second Temple Judaism: in by grace, staying in by works (encoded as a denial, so it ranks disputed) --- \cite{sanders-paul-palestinian-judaism-1977}
\item Good works are the fruit and evidence of saving faith, not its ground --- \cite[XI.2; XVI.2; XVI.5]{westminster-confession-1647}; \cite[III.xvi.1]{calvin-institutes-1960}; Jas 2:18-22; Eph 2:10; Rom 6:22; Luke 17:10
\item Justification and sanctification are inseparable but distinct --- \cite[Q. 77]{westminster-larger-catechism-1647}; 1 Cor 6:11; 1 Cor 1:30; Rom 4:6; Rom 4:8; Ezek 36:27; Heb 9:13-14; Rom 3:24-25; Rom 6:6; Rom 6:14; \cite[XI.1]{westminster-confession-1647}; \cite[III.xi.6]{calvin-institutes-1960}
\item an inference step --- \cite[Trent, Session VI (1547), canon 9]{tanner-decrees-1990}
\item an inference step --- \cite{gathercole-where-boasting-2002}; \cite{carson-variegated-nomism-1-2001}
\end{itemize}
\item \emph{Luke 7:50: “your faith has saved you”}, weakest at \emph{disputed}:
\begin{itemize}
\item Luke 7:47 — her love is the evidence of her forgiveness, not its ground --- \cite[ad loc. Luke 7:47]{marshall-luke-1978}; \cite[54--54]{padilla-narrative-criticism-2021}; \cite[Apology IV, on love and the fulfilling of the law: Luke 7:47–50]{kolb-wengert-book-of-concord-2000}; Luke 7:41-43; Luke 7:47; Luke 7:47-50
\item an inference step --- \cite[Trent, Session VI (1547), canon 9]{tanner-decrees-1990}
\end{itemize}
\item \emph{Sola fide from Peter (Acts 15:9–11)}, weakest at \emph{disputed}:
\begin{itemize}
\item The yoke refused at Acts 15:10 is the whole law as a condition of salvation --- \cite[ad loc. Acts 15:10]{bruce-acts-1988}; Acts 15:10; Gal 5:1; Jas 2:10; \cite{das-paul-law-covenant-2001}
\item Good works are the fruit and evidence of saving faith, not its ground --- \cite[XI.2; XVI.2; XVI.5]{westminster-confession-1647}; \cite[III.xvi.1]{calvin-institutes-1960}; Jas 2:18-22; Eph 2:10; Rom 6:22; Luke 17:10
\item Justification and sanctification are inseparable but distinct --- \cite[Q. 77]{westminster-larger-catechism-1647}; 1 Cor 6:11; 1 Cor 1:30; Rom 4:6; Rom 4:8; Ezek 36:27; Heb 9:13-14; Rom 3:24-25; Rom 6:6; Rom 6:14; \cite[XI.1]{westminster-confession-1647}; \cite[III.xi.6]{calvin-institutes-1960}
\item an inference step --- \cite[Trent, Session VI (1547), canon 9]{tanner-decrees-1990}
\end{itemize}
\item \emph{Apocalyptic reading, against faith as the condition}, weakest at \emph{disputed}:
\begin{itemize}
\item ‘Faith of Christ’ is Christ’s own faithfulness, the means of God’s rescue (encoded as a denial, so it ranks disputed) --- \cite{campbell-deliverance-god-2009}; \cite[ad loc. Gal 2:16]{martyn-galatians-1997}; \cite{hays-faith-jesus-christ-2002}
\item ‘The righteousness of God’ (δικαιοσύνη θεοῦ) names God’s act of rescue in Christ --- \cite{campbell-deliverance-god-2009}; \cite{martyn-galatians-1997}; Rom 3:21-26
\end{itemize}
\item \emph{Variegated nomism (Gathercole, Carson et al.)}, weakest at \emph{disputed}:
\begin{itemize}
\item Second Temple Judaism also held final vindication according to works (encoded as a denial, so it ranks disputed) --- \cite{gathercole-where-boasting-2002}; \cite{carson-variegated-nomism-1-2001}
\item an inference step --- \cite{gathercole-where-boasting-2002}; \cite{carson-variegated-nomism-1-2001}
\end{itemize}
\item \emph{Paul’s word (δικαιόω, by the rule of least meaning)}, weakest at \emph{well supported}:
\begin{itemize}
\item In Paul, ‘justify’ (δικαιόω) is a verdict: to declare righteous --- \cite[s.v. δικαιόω]{bauer-danker-lexicon-2000}; \cite{morris-apostolic-preaching-1965}; \cite[p. 80]{moo-romans-1996}; \cite{irons-righteousness-of-god-2015}; \cite{wright-justification-2009}; \cite[ad loc. Rom 3:24]{fitzmyer-romans-1993}
\item A word contributes the least meaning its context requires; one occurrence does not carry the whole concept --- \cite{silva-biblical-words-1994}; \cite{joos-semantic-axiom-1972}; \cite{barr-semantics-2004}
\item an inference step --- \cite{silva-biblical-words-1994}; \cite[p. 80]{moo-romans-1996}; \cite{barr-semantics-2004}
\end{itemize}
\item \emph{Trent, against ‘not by works’, read as what Paul’s word means}, weakest at \emph{disputed}:
\begin{itemize}
\item Justification is not remission of sins only, but renewal of the inward man --- \cite[Trent, Session VI (1547), Decree on Justification, ch. 7]{tanner-decrees-1990}; \cite[Trent, Session VI (1547), canon 11]{tanner-decrees-1990}
\item an inference step --- \cite[Trent, Session VI (1547), Decree on Justification, ch. 8]{tanner-decrees-1990}; \cite[Trent, Session VI (1547), Decree on Justification, ch. 7]{tanner-decrees-1990}; \cite[On the Spirit and the Letter, ch. 45 (26.45)]{npnf1-05-1887}; \cite{mcgrath-iustitia-dei-2005}
\end{itemize}
\item \emph{Romans 4:3–8, on Paul’s word (by the rule of least meaning)}, weakest at \emph{well supported}:
\begin{itemize}
\item Romans 4:4–5 contrasts wages owed with a gift reckoned to the one who believes --- Rom 4:4-5; \cite[III.xi.18]{calvin-institutes-1960}
\item ‘Count’ (λογίζομαι) in Romans 4:3–11 is the language of an account: credited to one --- \cite[s.v. λογίζομαι]{bauer-danker-lexicon-2000}; \cite[ad loc. Rom 4:3]{moo-romans-1996}
\item A word contributes the least meaning its context requires; one occurrence does not carry the whole concept --- \cite{silva-biblical-words-1994}; \cite{joos-semantic-axiom-1972}; \cite{barr-semantics-2004}
\item an inference step --- \cite[ad loc. Rom 4:6]{moo-romans-1996}; \cite[s.v. λογίζομαι]{bauer-danker-lexicon-2000}
\item an inference step --- \cite{silva-biblical-words-1994}; \cite[p. 80]{moo-romans-1996}; \cite{barr-semantics-2004}
\end{itemize}
\item \emph{Sola fide from John (6:28–29)}, weakest at \emph{disputed}:
\begin{itemize}
\item The believing of John 6:29 is trust, which brings nothing and receives Christ --- \cite[ad loc. John 6:29]{calvin-commentary-john-1847}; John 6:28-29
\item Good works are the fruit and evidence of saving faith, not its ground --- \cite[XI.2; XVI.2; XVI.5]{westminster-confession-1647}; \cite[III.xvi.1]{calvin-institutes-1960}; Jas 2:18-22; Eph 2:10; Rom 6:22; Luke 17:10
\item Justification and sanctification are inseparable but distinct --- \cite[Q. 77]{westminster-larger-catechism-1647}; 1 Cor 6:11; 1 Cor 1:30; Rom 4:6; Rom 4:8; Ezek 36:27; Heb 9:13-14; Rom 3:24-25; Rom 6:6; Rom 6:14; \cite[XI.1]{westminster-confession-1647}; \cite[III.xi.6]{calvin-institutes-1960}
\item an inference step --- \cite[Trent, Session VI (1547), canon 9]{tanner-decrees-1990}
\end{itemize}
\item \emph{Faith formed by charity (Aquinas on John 6:29)}, weakest at \emph{disputed}:
\begin{itemize}
\item The believing of John 6:29 is faith living through charity, the source of works --- \cite[cap. 6, lect. 3, n. 901]{aquinas-commentary-john-2012}; \cite[Trent, Session VI (1547), Decree on Justification, ch. 7]{tanner-decrees-1990}
\end{itemize}
\item \emph{Luther on Galatians 3:11–12 (against faith formed by charity)}, weakest at \emph{disputed}:
\begin{itemize}
\item an inference step --- \cite[ad loc. Gal 3:12]{luther-commentary-galatians-1998}; \cite[ad loc. Gal 3:11]{luther-commentary-galatians-1998}
\end{itemize}
\item \emph{Chrysostom on John 3:36 (believing is not enough without a right life)}, weakest at \emph{disputed}:
\begin{itemize}
\item an inference step --- \cite[Homily 31, §1]{chrysostom-homilies-john}; \cite{stanley-did-jesus-teach-2006}
\end{itemize}
\item \emph{Obedience as sanctification (John 3:36, Matthew 7:21)}, weakest at \emph{disputed}:
\begin{itemize}
\item The obedience of John 3:36 and Matthew 7:21 is sanctification, the fruit of faith --- \cite[ad loc. John 3:36]{ridderbos-john-1997}; \cite[ad loc. John 3:36]{carson-john-1991}; \cite[ad loc. Matt 7:21]{morris-matthew-1992}; \cite[55--55]{schreiner-justification-by-works-2015}; \cite{hoekema-saved-by-grace-1994}; \cite{murray-redemption-1984}; \cite[XVI.2]{westminster-confession-1647}; \cite[XIII.1]{westminster-confession-1647}; \cite[ad loc. John 3:36]{calvin-commentary-john-1847}
\item Justification and sanctification are inseparable but distinct --- \cite[Q. 77]{westminster-larger-catechism-1647}; 1 Cor 6:11; 1 Cor 1:30; Rom 4:6; Rom 4:8; Ezek 36:27; Heb 9:13-14; Rom 3:24-25; Rom 6:6; Rom 6:14; \cite[XI.1]{westminster-confession-1647}; \cite[III.xi.6]{calvin-institutes-1960}
\item an inference step --- \cite[55--55]{schreiner-justification-by-works-2015}; \cite[ad loc. John 3:36]{calvin-commentary-john-1847}; \cite[XVI.2]{westminster-confession-1647}; \cite[XIII.1]{westminster-confession-1647}
\end{itemize}
\end{itemize}
\noindent Why \emph{Paul’s gospel (Galatians 1), against justification as renewal} defeats \emph{Trent, against ‘not by works’}, at the level of the claims: \texttt{where\allowbreak{}Trent\allowbreak{}Parts\allowbreak{}From\allowbreak{}Paul}.
\noindent Why \emph{Luke 7:50: “your faith has saved you”} defeats \emph{Trent, against ‘not by works’}, at the level of the claims: \texttt{why\allowbreak{}The\allowbreak{}Dominical\allowbreak{}Case\allowbreak{}Stands\allowbreak{}Against\allowbreak{}Trent}.
\par\medskip\noindent\textbf{\texttt{gospel\_\allowbreak{}not\_\allowbreak{}forced}}\par\nobreak\noindent
\textbf{Nor is it forced.} Trent as Trent states it can still be defended, and Paul’s gospel cannot be held with it. What would decide between them is the step \texttt{where\allowbreak{}Trent\allowbreak{}Parts\allowbreak{}From\allowbreak{}Paul} names, that a verdict on Christ’s finished work excludes renewal, and that step is rated \texttt{disputed}. The argument from Paul’s word does not decide it. Taken as a claim about what God does, Trent’s definition is not reached by anything the word’s meaning shows (\texttt{lexical\_\allowbreak{}case\_\allowbreak{}does\_\allowbreak{}not\_\allowbreak{}reach\_\allowbreak{}what\_\allowbreak{}god\_\allowbreak{}does}).

{\small\begin{verbatim}
theorem gospel_not_forced : ¬Framework.SkepticallyAccepted
    solaFideDispute.defeats Party.gospel
-- axioms: propext, Classical.choice, Quot.sound
\end{verbatim}}
\paragraph*{Where Luke’s case and Trent part}
\leavevmode\par\nobreak
The verdicts above say that Luke’s case and Trent defeat each other, and that the debate chooses neither. This says \emph{why}, claim by claim. The two part over what justification \emph{is}. Trent defines it as the renewal of the inward person by charity poured in, and from that definition it follows that faith without charity does not suffice. Luke’s words say that her faith saved her.

\par\medskip\noindent\textbf{\texttt{why\allowbreak{}The\allowbreak{}Dominical\allowbreak{}Case\allowbreak{}Stands\allowbreak{}Against\allowbreak{}Trent}}\par\nobreak\noindent
\textbf{Why Luke’s case stands against Trent.} Everything turns on Luke’s own step: from “your faith has saved you”, with “saved” (σῴζω) read as saved from sin and her love as the evidence of forgiveness, to the conclusion that faith suffices.

Trent cannot hold that step together with Luke’s words. Two things in Trent’s case contradict it, and both are needed: its definition of justification as the renewal of the inward person (Session VI, ch. 7), and the step from that definition to denying that faith suffices (canon 9). Nothing else in Trent’s case does: not its account of merit, and not its reading of works.

So the disagreement is not over what Luke says. It is over what justification is. Grant Trent’s definition, and “your faith has saved you” cannot mean that faith sufficed. Grant that it did, and Trent’s definition cannot stand.

\paragraph{Why \emph{Luke 7:50: “your faith has saved you”} stands against \emph{Trent, against ‘not by works’}.}
\begin{itemize}
\item \textbf{The crux:} \((\pvar{TestimonyArguments.SolaFide}{17} \land \pvar{TestimonyArguments.SolaFide}{18} \land \pvar{TestimonyArguments.SolaFide}{19}) \rightarrow \pvar{TestimonyArguments.SolaFide}{26}\), a premise of \emph{Luke 7:50: “your faith has saved you”}.
\item \textbf{What it does:} its conclusion does not follow without it.
\item \textbf{Granted:} \(\pvar{TestimonyArguments.SolaFide}{17}\), \(\pvar{TestimonyArguments.SolaFide}{18}\), \(\pvar{TestimonyArguments.SolaFide}{19}\), which the break also needs.
\item \textbf{Where the rival breaks:} \(\pvar{TestimonyArguments.SolaFide}{33}\), \(\pvar{TestimonyArguments.SolaFide}{33} \rightarrow \lnot \pvar{TestimonyArguments.SolaFide}{26}\) cannot be held together with \((\pvar{TestimonyArguments.SolaFide}{17} \land \pvar{TestimonyArguments.SolaFide}{18} \land \pvar{TestimonyArguments.SolaFide}{19}) \rightarrow \pvar{TestimonyArguments.SolaFide}{26}\) and \(\pvar{TestimonyArguments.SolaFide}{17}\), \(\pvar{TestimonyArguments.SolaFide}{18}\), \(\pvar{TestimonyArguments.SolaFide}{19}\); each is needed for the break, and without the crux they stand.
\item \textbf{What this rests on:}
\begin{itemize}
\item \(\pvar{TestimonyArguments.SolaFide}{17}\) Luke 7:50 — Jesus says ‘your faith has saved you’ after declaring sins forgiven --- \emph{consensus}: Luke 7:50; \cite[ad loc. Luke 7:50]{marshall-luke-1978}
\item \(\pvar{TestimonyArguments.SolaFide}{18}\) ‘Saved’ (σῴζω) at Luke 7:50 means saved from sin, not healed --- \emph{well supported}: \cite[Apology IV, on love and the fulfilling of the law: Luke 7:47–50]{kolb-wengert-book-of-concord-2000}; \cite[54--55]{padilla-narrative-criticism-2021}; Luke 7:47-50; \cite[ad loc. Luke 7:50]{marshall-luke-1978}
\item \(\pvar{TestimonyArguments.SolaFide}{19}\) Luke 7:47 — her love is the evidence of her forgiveness, not its ground --- \emph{disputed}: \cite[ad loc. Luke 7:47]{marshall-luke-1978}; \cite[54--54]{padilla-narrative-criticism-2021}; \cite[Apology IV, on love and the fulfilling of the law: Luke 7:47–50]{kolb-wengert-book-of-concord-2000}; Luke 7:41-43; Luke 7:47; Luke 7:47-50
\item the step itself --- \emph{disputed}: \cite[Trent, Session VI (1547), canon 9]{tanner-decrees-1990}
\end{itemize}
\end{itemize}
\paragraph*{The Reformation debate, heard alone}
\leavevmode\par\nobreak
\par\medskip\noindent\textbf{\texttt{reformation\allowbreak{}Alone}}\par\nobreak\noindent
The positions of the sixteenth-century debate: the four Reformed strands, Paul’s gospel, the two arguments about Paul’s word, Trent as it states its definition and Trent taken as a claim about Paul’s word, the exchange over John 6:29 between Aquinas and Luther, and Chrysostom on John 3:36 with the Reformed answer to him. Not heard: the apocalyptic reading, Sanders and his critics, and Jervell.

{\small\begin{verbatim}
def reformationAlone : Dispute Claim { i // i ∉ [Party.apocalyptic,
    Party.sanders, Party.critics, Party.jervell] } :=
  solaFideDispute.restrict fun x =>
    x ∉
      [Party.apocalyptic, Party.sanders,
        Party.critics, Party.jervell]
\end{verbatim}}
\par\medskip\noindent\textbf{\texttt{reformation\allowbreak{}Settles\allowbreak{}The\allowbreak{}Word}}\par\nobreak\noindent
How far the Reformation debate can be settled. Nothing defeats either argument about Paul’s word, so both are accepted first, and nothing joins them. Each Reformed position is defeated by Trent. Trent is defeated by Paul’s case, and by Luke’s, Peter’s and Paul’s gospel. Trent taken as a claim about Paul’s word is defeated by the argument from the word itself.

\paragraph{What the dispute forces is exactly \emph{Paul’s word (δικαιόω, by the rule of least meaning)}, \emph{Romans 4:3–8, on Paul’s word (by the rule of least meaning)}.}
\begin{itemize}
\item Stage 1: \emph{Paul’s word (δικαιόω, by the rule of least meaning)}, \emph{Romans 4:3–8, on Paul’s word (by the rule of least meaning)}.
\begin{itemize}
\item \emph{Paul’s word (δικαιόω, by the rule of least meaning)} is defeated by nothing.
\item \emph{Romans 4:3–8, on Paul’s word (by the rule of least meaning)} is defeated by nothing.
\end{itemize}
\item Nothing else is forced:
\begin{itemize}
\item \emph{Trent, against ‘not by works’} defeats \emph{Sola fide from Paul (Galatians 2:16)}, and nothing forced defeats \emph{Trent, against ‘not by works’}.
\item \emph{Trent, against ‘not by works’} defeats \emph{Luke 7:50: “your faith has saved you”}, and nothing forced defeats \emph{Trent, against ‘not by works’}.
\item \emph{Trent, against ‘not by works’} defeats \emph{Sola fide from Peter (Acts 15:9–11)}, and nothing forced defeats \emph{Trent, against ‘not by works’}.
\item \emph{Sola fide from Paul (Galatians 2:16)} defeats \emph{Trent, against ‘not by works’}, and nothing forced defeats \emph{Sola fide from Paul (Galatians 2:16)}.
\item \emph{Trent, against ‘not by works’} defeats \emph{Paul’s gospel (Galatians 1), against justification as renewal}, and nothing forced defeats \emph{Trent, against ‘not by works’}.
\item \emph{Paul’s word (δικαιόω, by the rule of least meaning)} defeats \emph{Trent, against ‘not by works’, read as what Paul’s word means} (\texttt{where\allowbreak{}Trents\allowbreak{}Word\allowbreak{}Reading\allowbreak{}Falls}), and nothing forced defeats \emph{Paul’s word (δικαιόω, by the rule of least meaning)}.
\item \emph{Trent, against ‘not by works’} defeats \emph{Sola fide from John (6:28–29)}, and nothing forced defeats \emph{Trent, against ‘not by works’}.
\item \emph{Luther on Galatians 3:11–12 (against faith formed by charity)} defeats \emph{Faith formed by charity (Aquinas on John 6:29)}, and nothing forced defeats \emph{Luther on Galatians 3:11–12 (against faith formed by charity)}.
\item \emph{Faith formed by charity (Aquinas on John 6:29)} defeats \emph{Luther on Galatians 3:11–12 (against faith formed by charity)}, and nothing forced defeats \emph{Faith formed by charity (Aquinas on John 6:29)}.
\item \emph{Sola fide from Paul (Galatians 2:16)} defeats \emph{Chrysostom on John 3:36 (believing is not enough without a right life)}, and nothing forced defeats \emph{Sola fide from Paul (Galatians 2:16)}.
\item \emph{Trent, against ‘not by works’} defeats \emph{Obedience as sanctification (John 3:36, Matthew 7:21)}, and nothing forced defeats \emph{Trent, against ‘not by works’}.
\end{itemize}
\end{itemize}
\par\noindent\textbf{What this rests on.} The reasons state 11 defeats and 34 absences of defeat, each a cell of the defeat table, \texttt{sola\allowbreak{}Fide\allowbreak{}Dispute\_\allowbreak{}defeats}, computed from the two parties' premises and checked by the kernel. Whether an attack survives turns on the attacker's weakest link and on the rating of what it attacks, so a changed rating can change the verdict. The weakest links:
\begin{itemize}
\item \emph{Sola fide from Paul (Galatians 2:16)}, weakest at \emph{disputed}:
\begin{itemize}
\item In Galatians 2:16 ‘faith of Christ’ (πίστις Χριστοῦ) means our faith in Christ --- \cite{dunn-once-more-pistis-christou-1991}; \cite{matlock-detheologizing-2000}; Gal 2:16; Rom 3:22; Phil 3:9
\item Second Temple Judaism: in by grace, staying in by works (encoded as a denial, so it ranks disputed) --- \cite{sanders-paul-palestinian-judaism-1977}
\item Good works are the fruit and evidence of saving faith, not its ground --- \cite[XI.2; XVI.2; XVI.5]{westminster-confession-1647}; \cite[III.xvi.1]{calvin-institutes-1960}; Jas 2:18-22; Eph 2:10; Rom 6:22; Luke 17:10
\item Justification and sanctification are inseparable but distinct --- \cite[Q. 77]{westminster-larger-catechism-1647}; 1 Cor 6:11; 1 Cor 1:30; Rom 4:6; Rom 4:8; Ezek 36:27; Heb 9:13-14; Rom 3:24-25; Rom 6:6; Rom 6:14; \cite[XI.1]{westminster-confession-1647}; \cite[III.xi.6]{calvin-institutes-1960}
\item an inference step --- \cite[Trent, Session VI (1547), canon 9]{tanner-decrees-1990}
\item an inference step --- \cite{gathercole-where-boasting-2002}; \cite{carson-variegated-nomism-1-2001}
\end{itemize}
\item \emph{Paul’s word (δικαιόω, by the rule of least meaning)}, weakest at \emph{well supported}:
\begin{itemize}
\item In Paul, ‘justify’ (δικαιόω) is a verdict: to declare righteous --- \cite[s.v. δικαιόω]{bauer-danker-lexicon-2000}; \cite{morris-apostolic-preaching-1965}; \cite[p. 80]{moo-romans-1996}; \cite{irons-righteousness-of-god-2015}; \cite{wright-justification-2009}; \cite[ad loc. Rom 3:24]{fitzmyer-romans-1993}
\item A word contributes the least meaning its context requires; one occurrence does not carry the whole concept --- \cite{silva-biblical-words-1994}; \cite{joos-semantic-axiom-1972}; \cite{barr-semantics-2004}
\item an inference step --- \cite{silva-biblical-words-1994}; \cite[p. 80]{moo-romans-1996}; \cite{barr-semantics-2004}
\end{itemize}
\item \emph{Luke 7:50: “your faith has saved you”}, weakest at \emph{disputed}:
\begin{itemize}
\item Luke 7:47 — her love is the evidence of her forgiveness, not its ground --- \cite[ad loc. Luke 7:47]{marshall-luke-1978}; \cite[54--54]{padilla-narrative-criticism-2021}; \cite[Apology IV, on love and the fulfilling of the law: Luke 7:47–50]{kolb-wengert-book-of-concord-2000}; Luke 7:41-43; Luke 7:47; Luke 7:47-50
\item an inference step --- \cite[Trent, Session VI (1547), canon 9]{tanner-decrees-1990}
\end{itemize}
\item \emph{Sola fide from Peter (Acts 15:9–11)}, weakest at \emph{disputed}:
\begin{itemize}
\item The yoke refused at Acts 15:10 is the whole law as a condition of salvation --- \cite[ad loc. Acts 15:10]{bruce-acts-1988}; Acts 15:10; Gal 5:1; Jas 2:10; \cite{das-paul-law-covenant-2001}
\item Good works are the fruit and evidence of saving faith, not its ground --- \cite[XI.2; XVI.2; XVI.5]{westminster-confession-1647}; \cite[III.xvi.1]{calvin-institutes-1960}; Jas 2:18-22; Eph 2:10; Rom 6:22; Luke 17:10
\item Justification and sanctification are inseparable but distinct --- \cite[Q. 77]{westminster-larger-catechism-1647}; 1 Cor 6:11; 1 Cor 1:30; Rom 4:6; Rom 4:8; Ezek 36:27; Heb 9:13-14; Rom 3:24-25; Rom 6:6; Rom 6:14; \cite[XI.1]{westminster-confession-1647}; \cite[III.xi.6]{calvin-institutes-1960}
\item an inference step --- \cite[Trent, Session VI (1547), canon 9]{tanner-decrees-1990}
\end{itemize}
\item \emph{Trent, against ‘not by works’}, weakest at \emph{disputed}:
\begin{itemize}
\item Justification is not remission of sins only, but renewal of the inward man --- \cite[Trent, Session VI (1547), Decree on Justification, ch. 7]{tanner-decrees-1990}; \cite[Trent, Session VI (1547), canon 11]{tanner-decrees-1990}
\end{itemize}
\item \emph{Paul’s gospel (Galatians 1), against justification as renewal}, weakest at \emph{disputed}:
\begin{itemize}
\item an inference step --- \cite[ad loc. Gal 1:6]{moo-galatians-2013}; \cite[XI.1]{westminster-confession-1647}
\item an inference step --- \cite[XI.1]{westminster-confession-1647}; \cite[III.xi.2]{calvin-institutes-1960}; \cite{morris-apostolic-preaching-1965}
\end{itemize}
\item \emph{Trent, against ‘not by works’, read as what Paul’s word means}, weakest at \emph{disputed}:
\begin{itemize}
\item Justification is not remission of sins only, but renewal of the inward man --- \cite[Trent, Session VI (1547), Decree on Justification, ch. 7]{tanner-decrees-1990}; \cite[Trent, Session VI (1547), canon 11]{tanner-decrees-1990}
\item an inference step --- \cite[Trent, Session VI (1547), Decree on Justification, ch. 8]{tanner-decrees-1990}; \cite[Trent, Session VI (1547), Decree on Justification, ch. 7]{tanner-decrees-1990}; \cite[On the Spirit and the Letter, ch. 45 (26.45)]{npnf1-05-1887}; \cite{mcgrath-iustitia-dei-2005}
\end{itemize}
\item \emph{Romans 4:3–8, on Paul’s word (by the rule of least meaning)}, weakest at \emph{well supported}:
\begin{itemize}
\item Romans 4:4–5 contrasts wages owed with a gift reckoned to the one who believes --- Rom 4:4-5; \cite[III.xi.18]{calvin-institutes-1960}
\item ‘Count’ (λογίζομαι) in Romans 4:3–11 is the language of an account: credited to one --- \cite[s.v. λογίζομαι]{bauer-danker-lexicon-2000}; \cite[ad loc. Rom 4:3]{moo-romans-1996}
\item A word contributes the least meaning its context requires; one occurrence does not carry the whole concept --- \cite{silva-biblical-words-1994}; \cite{joos-semantic-axiom-1972}; \cite{barr-semantics-2004}
\item an inference step --- \cite[ad loc. Rom 4:6]{moo-romans-1996}; \cite[s.v. λογίζομαι]{bauer-danker-lexicon-2000}
\item an inference step --- \cite{silva-biblical-words-1994}; \cite[p. 80]{moo-romans-1996}; \cite{barr-semantics-2004}
\end{itemize}
\item \emph{Sola fide from John (6:28–29)}, weakest at \emph{disputed}:
\begin{itemize}
\item The believing of John 6:29 is trust, which brings nothing and receives Christ --- \cite[ad loc. John 6:29]{calvin-commentary-john-1847}; John 6:28-29
\item Good works are the fruit and evidence of saving faith, not its ground --- \cite[XI.2; XVI.2; XVI.5]{westminster-confession-1647}; \cite[III.xvi.1]{calvin-institutes-1960}; Jas 2:18-22; Eph 2:10; Rom 6:22; Luke 17:10
\item Justification and sanctification are inseparable but distinct --- \cite[Q. 77]{westminster-larger-catechism-1647}; 1 Cor 6:11; 1 Cor 1:30; Rom 4:6; Rom 4:8; Ezek 36:27; Heb 9:13-14; Rom 3:24-25; Rom 6:6; Rom 6:14; \cite[XI.1]{westminster-confession-1647}; \cite[III.xi.6]{calvin-institutes-1960}
\item an inference step --- \cite[Trent, Session VI (1547), canon 9]{tanner-decrees-1990}
\end{itemize}
\item \emph{Faith formed by charity (Aquinas on John 6:29)}, weakest at \emph{disputed}:
\begin{itemize}
\item The believing of John 6:29 is faith living through charity, the source of works --- \cite[cap. 6, lect. 3, n. 901]{aquinas-commentary-john-2012}; \cite[Trent, Session VI (1547), Decree on Justification, ch. 7]{tanner-decrees-1990}
\end{itemize}
\item \emph{Luther on Galatians 3:11–12 (against faith formed by charity)}, weakest at \emph{disputed}:
\begin{itemize}
\item an inference step --- \cite[ad loc. Gal 3:12]{luther-commentary-galatians-1998}; \cite[ad loc. Gal 3:11]{luther-commentary-galatians-1998}
\end{itemize}
\item \emph{Chrysostom on John 3:36 (believing is not enough without a right life)}, weakest at \emph{disputed}:
\begin{itemize}
\item an inference step --- \cite[Homily 31, §1]{chrysostom-homilies-john}; \cite{stanley-did-jesus-teach-2006}
\end{itemize}
\item \emph{Obedience as sanctification (John 3:36, Matthew 7:21)}, weakest at \emph{disputed}:
\begin{itemize}
\item The obedience of John 3:36 and Matthew 7:21 is sanctification, the fruit of faith --- \cite[ad loc. John 3:36]{ridderbos-john-1997}; \cite[ad loc. John 3:36]{carson-john-1991}; \cite[ad loc. Matt 7:21]{morris-matthew-1992}; \cite[55--55]{schreiner-justification-by-works-2015}; \cite{hoekema-saved-by-grace-1994}; \cite{murray-redemption-1984}; \cite[XVI.2]{westminster-confession-1647}; \cite[XIII.1]{westminster-confession-1647}; \cite[ad loc. John 3:36]{calvin-commentary-john-1847}
\item Justification and sanctification are inseparable but distinct --- \cite[Q. 77]{westminster-larger-catechism-1647}; 1 Cor 6:11; 1 Cor 1:30; Rom 4:6; Rom 4:8; Ezek 36:27; Heb 9:13-14; Rom 3:24-25; Rom 6:6; Rom 6:14; \cite[XI.1]{westminster-confession-1647}; \cite[III.xi.6]{calvin-institutes-1960}
\item an inference step --- \cite[55--55]{schreiner-justification-by-works-2015}; \cite[ad loc. John 3:36]{calvin-commentary-john-1847}; \cite[XVI.2]{westminster-confession-1647}; \cite[XIII.1]{westminster-confession-1647}
\end{itemize}
\end{itemize}
\noindent Why \emph{Paul’s word (δικαιόω, by the rule of least meaning)} defeats \emph{Trent, against ‘not by works’, read as what Paul’s word means}, at the level of the claims: \texttt{where\allowbreak{}Trents\allowbreak{}Word\allowbreak{}Reading\allowbreak{}Falls}.
\par\medskip\noindent\textbf{\texttt{reformation\_\allowbreak{}dispute\_\allowbreak{}forces\_\allowbreak{}only\_\allowbreak{}pauls\_\allowbreak{}word}}\par\nobreak\noindent
\textbf{What the Reformation debate forces: the meaning of Paul’s word, and nothing else.} Heard as Rome and the Reformers alone, the forced positions are exactly the two arguments about Paul’s word. Trent cannot claim that its definition is what Paul’s “justify” (δικαιόω) means.

Nothing else is forced. Sola fide from each strand, Paul’s gospel, and Trent’s definition as a claim about what God does each have an attacker they cannot remove.

{\small\begin{verbatim}
theorem reformation_dispute_forces_only_pauls_word : Framework.grounded
    reformationAlone.defeats = Solver.toSet (Witness.listSub (fun x => x ∉
    [Party.apocalyptic, Party.sanders, Party.critics, Party.jervell])
    [Party.lexical, Party.romansOnTheWord])
-- axioms: propext, Classical.choice, Quot.sound
\end{verbatim}}
\par\medskip\noindent\textbf{\texttt{trent\allowbreak{}Stands\allowbreak{}With\allowbreak{}Pauls\allowbreak{}Word}}\par\nobreak\noindent
Why Trent can be defended in the Reformation debate: it can be held with the arguments about Paul’s word, and it answers every position that attacks it.

\paragraph{\emph{Trent, against ‘not by works’} is accepted on some resolution.}
\begin{itemize}
\item A position holding \emph{Trent, against ‘not by works’}, \emph{Paul’s word (δικαιόω, by the rule of least meaning)}, \emph{Romans 4:3–8, on Paul’s word (by the rule of least meaning)} can be held.
\item None of \emph{Trent, against ‘not by works’}, \emph{Paul’s word (δικαιόω, by the rule of least meaning)}, \emph{Romans 4:3–8, on Paul’s word (by the rule of least meaning)} defeats another, and each attack on them is answered from among them:
\begin{itemize}
\item \emph{Sola fide from Paul (Galatians 2:16)} defeats \emph{Trent, against ‘not by works’}, and \emph{Trent, against ‘not by works’} defeats \emph{Sola fide from Paul (Galatians 2:16)}.
\item \emph{Luke 7:50: “your faith has saved you”} defeats \emph{Trent, against ‘not by works’} (\texttt{why\allowbreak{}The\allowbreak{}Dominical\allowbreak{}Case\allowbreak{}Stands\allowbreak{}Against\allowbreak{}Trent}), and \emph{Trent, against ‘not by works’} defeats \emph{Luke 7:50: “your faith has saved you”}.
\item \emph{Sola fide from Peter (Acts 15:9–11)} defeats \emph{Trent, against ‘not by works’}, and \emph{Trent, against ‘not by works’} defeats \emph{Sola fide from Peter (Acts 15:9–11)}.
\item \emph{Paul’s gospel (Galatians 1), against justification as renewal} defeats \emph{Trent, against ‘not by works’} (\texttt{where\allowbreak{}Trent\allowbreak{}Parts\allowbreak{}From\allowbreak{}Paul}), and \emph{Trent, against ‘not by works’} defeats \emph{Paul’s gospel (Galatians 1), against justification as renewal}.
\item \emph{Sola fide from John (6:28–29)} defeats \emph{Trent, against ‘not by works’}, and \emph{Trent, against ‘not by works’} defeats \emph{Sola fide from John (6:28–29)}.
\item \emph{Obedience as sanctification (John 3:36, Matthew 7:21)} defeats \emph{Trent, against ‘not by works’}, and \emph{Trent, against ‘not by works’} defeats \emph{Obedience as sanctification (John 3:36, Matthew 7:21)}.
\end{itemize}
\end{itemize}
\par\noindent\textbf{What this rests on.} The reasons state 12 defeats and 33 absences of defeat, each a cell of the defeat table, \texttt{sola\allowbreak{}Fide\allowbreak{}Dispute\_\allowbreak{}defeats}, computed from the two parties' premises and checked by the kernel. Whether an attack survives turns on the attacker's weakest link and on the rating of what it attacks, so a changed rating can change the verdict. The weakest links:
\begin{itemize}
\item \emph{Trent, against ‘not by works’}, weakest at \emph{disputed}:
\begin{itemize}
\item Justification is not remission of sins only, but renewal of the inward man --- \cite[Trent, Session VI (1547), Decree on Justification, ch. 7]{tanner-decrees-1990}; \cite[Trent, Session VI (1547), canon 11]{tanner-decrees-1990}
\end{itemize}
\item \emph{Paul’s word (δικαιόω, by the rule of least meaning)}, weakest at \emph{well supported}:
\begin{itemize}
\item In Paul, ‘justify’ (δικαιόω) is a verdict: to declare righteous --- \cite[s.v. δικαιόω]{bauer-danker-lexicon-2000}; \cite{morris-apostolic-preaching-1965}; \cite[p. 80]{moo-romans-1996}; \cite{irons-righteousness-of-god-2015}; \cite{wright-justification-2009}; \cite[ad loc. Rom 3:24]{fitzmyer-romans-1993}
\item A word contributes the least meaning its context requires; one occurrence does not carry the whole concept --- \cite{silva-biblical-words-1994}; \cite{joos-semantic-axiom-1972}; \cite{barr-semantics-2004}
\item an inference step --- \cite{silva-biblical-words-1994}; \cite[p. 80]{moo-romans-1996}; \cite{barr-semantics-2004}
\end{itemize}
\item \emph{Romans 4:3–8, on Paul’s word (by the rule of least meaning)}, weakest at \emph{well supported}:
\begin{itemize}
\item Romans 4:4–5 contrasts wages owed with a gift reckoned to the one who believes --- Rom 4:4-5; \cite[III.xi.18]{calvin-institutes-1960}
\item ‘Count’ (λογίζομαι) in Romans 4:3–11 is the language of an account: credited to one --- \cite[s.v. λογίζομαι]{bauer-danker-lexicon-2000}; \cite[ad loc. Rom 4:3]{moo-romans-1996}
\item A word contributes the least meaning its context requires; one occurrence does not carry the whole concept --- \cite{silva-biblical-words-1994}; \cite{joos-semantic-axiom-1972}; \cite{barr-semantics-2004}
\item an inference step --- \cite[ad loc. Rom 4:6]{moo-romans-1996}; \cite[s.v. λογίζομαι]{bauer-danker-lexicon-2000}
\item an inference step --- \cite{silva-biblical-words-1994}; \cite[p. 80]{moo-romans-1996}; \cite{barr-semantics-2004}
\end{itemize}
\item \emph{Sola fide from Paul (Galatians 2:16)}, weakest at \emph{disputed}:
\begin{itemize}
\item In Galatians 2:16 ‘faith of Christ’ (πίστις Χριστοῦ) means our faith in Christ --- \cite{dunn-once-more-pistis-christou-1991}; \cite{matlock-detheologizing-2000}; Gal 2:16; Rom 3:22; Phil 3:9
\item Second Temple Judaism: in by grace, staying in by works (encoded as a denial, so it ranks disputed) --- \cite{sanders-paul-palestinian-judaism-1977}
\item Good works are the fruit and evidence of saving faith, not its ground --- \cite[XI.2; XVI.2; XVI.5]{westminster-confession-1647}; \cite[III.xvi.1]{calvin-institutes-1960}; Jas 2:18-22; Eph 2:10; Rom 6:22; Luke 17:10
\item Justification and sanctification are inseparable but distinct --- \cite[Q. 77]{westminster-larger-catechism-1647}; 1 Cor 6:11; 1 Cor 1:30; Rom 4:6; Rom 4:8; Ezek 36:27; Heb 9:13-14; Rom 3:24-25; Rom 6:6; Rom 6:14; \cite[XI.1]{westminster-confession-1647}; \cite[III.xi.6]{calvin-institutes-1960}
\item an inference step --- \cite[Trent, Session VI (1547), canon 9]{tanner-decrees-1990}
\item an inference step --- \cite{gathercole-where-boasting-2002}; \cite{carson-variegated-nomism-1-2001}
\end{itemize}
\item \emph{Luke 7:50: “your faith has saved you”}, weakest at \emph{disputed}:
\begin{itemize}
\item Luke 7:47 — her love is the evidence of her forgiveness, not its ground --- \cite[ad loc. Luke 7:47]{marshall-luke-1978}; \cite[54--54]{padilla-narrative-criticism-2021}; \cite[Apology IV, on love and the fulfilling of the law: Luke 7:47–50]{kolb-wengert-book-of-concord-2000}; Luke 7:41-43; Luke 7:47; Luke 7:47-50
\item an inference step --- \cite[Trent, Session VI (1547), canon 9]{tanner-decrees-1990}
\end{itemize}
\item \emph{Sola fide from Peter (Acts 15:9–11)}, weakest at \emph{disputed}:
\begin{itemize}
\item The yoke refused at Acts 15:10 is the whole law as a condition of salvation --- \cite[ad loc. Acts 15:10]{bruce-acts-1988}; Acts 15:10; Gal 5:1; Jas 2:10; \cite{das-paul-law-covenant-2001}
\item Good works are the fruit and evidence of saving faith, not its ground --- \cite[XI.2; XVI.2; XVI.5]{westminster-confession-1647}; \cite[III.xvi.1]{calvin-institutes-1960}; Jas 2:18-22; Eph 2:10; Rom 6:22; Luke 17:10
\item Justification and sanctification are inseparable but distinct --- \cite[Q. 77]{westminster-larger-catechism-1647}; 1 Cor 6:11; 1 Cor 1:30; Rom 4:6; Rom 4:8; Ezek 36:27; Heb 9:13-14; Rom 3:24-25; Rom 6:6; Rom 6:14; \cite[XI.1]{westminster-confession-1647}; \cite[III.xi.6]{calvin-institutes-1960}
\item an inference step --- \cite[Trent, Session VI (1547), canon 9]{tanner-decrees-1990}
\end{itemize}
\item \emph{Paul’s gospel (Galatians 1), against justification as renewal}, weakest at \emph{disputed}:
\begin{itemize}
\item an inference step --- \cite[ad loc. Gal 1:6]{moo-galatians-2013}; \cite[XI.1]{westminster-confession-1647}
\item an inference step --- \cite[XI.1]{westminster-confession-1647}; \cite[III.xi.2]{calvin-institutes-1960}; \cite{morris-apostolic-preaching-1965}
\end{itemize}
\item \emph{Trent, against ‘not by works’, read as what Paul’s word means}, weakest at \emph{disputed}:
\begin{itemize}
\item Justification is not remission of sins only, but renewal of the inward man --- \cite[Trent, Session VI (1547), Decree on Justification, ch. 7]{tanner-decrees-1990}; \cite[Trent, Session VI (1547), canon 11]{tanner-decrees-1990}
\item an inference step --- \cite[Trent, Session VI (1547), Decree on Justification, ch. 8]{tanner-decrees-1990}; \cite[Trent, Session VI (1547), Decree on Justification, ch. 7]{tanner-decrees-1990}; \cite[On the Spirit and the Letter, ch. 45 (26.45)]{npnf1-05-1887}; \cite{mcgrath-iustitia-dei-2005}
\end{itemize}
\item \emph{Sola fide from John (6:28–29)}, weakest at \emph{disputed}:
\begin{itemize}
\item The believing of John 6:29 is trust, which brings nothing and receives Christ --- \cite[ad loc. John 6:29]{calvin-commentary-john-1847}; John 6:28-29
\item Good works are the fruit and evidence of saving faith, not its ground --- \cite[XI.2; XVI.2; XVI.5]{westminster-confession-1647}; \cite[III.xvi.1]{calvin-institutes-1960}; Jas 2:18-22; Eph 2:10; Rom 6:22; Luke 17:10
\item Justification and sanctification are inseparable but distinct --- \cite[Q. 77]{westminster-larger-catechism-1647}; 1 Cor 6:11; 1 Cor 1:30; Rom 4:6; Rom 4:8; Ezek 36:27; Heb 9:13-14; Rom 3:24-25; Rom 6:6; Rom 6:14; \cite[XI.1]{westminster-confession-1647}; \cite[III.xi.6]{calvin-institutes-1960}
\item an inference step --- \cite[Trent, Session VI (1547), canon 9]{tanner-decrees-1990}
\end{itemize}
\item \emph{Faith formed by charity (Aquinas on John 6:29)}, weakest at \emph{disputed}:
\begin{itemize}
\item The believing of John 6:29 is faith living through charity, the source of works --- \cite[cap. 6, lect. 3, n. 901]{aquinas-commentary-john-2012}; \cite[Trent, Session VI (1547), Decree on Justification, ch. 7]{tanner-decrees-1990}
\end{itemize}
\item \emph{Luther on Galatians 3:11–12 (against faith formed by charity)}, weakest at \emph{disputed}:
\begin{itemize}
\item an inference step --- \cite[ad loc. Gal 3:12]{luther-commentary-galatians-1998}; \cite[ad loc. Gal 3:11]{luther-commentary-galatians-1998}
\end{itemize}
\item \emph{Chrysostom on John 3:36 (believing is not enough without a right life)}, weakest at \emph{disputed}:
\begin{itemize}
\item an inference step --- \cite[Homily 31, §1]{chrysostom-homilies-john}; \cite{stanley-did-jesus-teach-2006}
\end{itemize}
\item \emph{Obedience as sanctification (John 3:36, Matthew 7:21)}, weakest at \emph{disputed}:
\begin{itemize}
\item The obedience of John 3:36 and Matthew 7:21 is sanctification, the fruit of faith --- \cite[ad loc. John 3:36]{ridderbos-john-1997}; \cite[ad loc. John 3:36]{carson-john-1991}; \cite[ad loc. Matt 7:21]{morris-matthew-1992}; \cite[55--55]{schreiner-justification-by-works-2015}; \cite{hoekema-saved-by-grace-1994}; \cite{murray-redemption-1984}; \cite[XVI.2]{westminster-confession-1647}; \cite[XIII.1]{westminster-confession-1647}; \cite[ad loc. John 3:36]{calvin-commentary-john-1847}
\item Justification and sanctification are inseparable but distinct --- \cite[Q. 77]{westminster-larger-catechism-1647}; 1 Cor 6:11; 1 Cor 1:30; Rom 4:6; Rom 4:8; Ezek 36:27; Heb 9:13-14; Rom 3:24-25; Rom 6:6; Rom 6:14; \cite[XI.1]{westminster-confession-1647}; \cite[III.xi.6]{calvin-institutes-1960}
\item an inference step --- \cite[55--55]{schreiner-justification-by-works-2015}; \cite[ad loc. John 3:36]{calvin-commentary-john-1847}; \cite[XVI.2]{westminster-confession-1647}; \cite[XIII.1]{westminster-confession-1647}
\end{itemize}
\end{itemize}
\noindent Why \emph{Luke 7:50: “your faith has saved you”} defeats \emph{Trent, against ‘not by works’}, at the level of the claims: \texttt{why\allowbreak{}The\allowbreak{}Dominical\allowbreak{}Case\allowbreak{}Stands\allowbreak{}Against\allowbreak{}Trent}.
\noindent Why \emph{Paul’s gospel (Galatians 1), against justification as renewal} defeats \emph{Trent, against ‘not by works’}, at the level of the claims: \texttt{where\allowbreak{}Trent\allowbreak{}Parts\allowbreak{}From\allowbreak{}Paul}.
\par\medskip\noindent\textbf{\texttt{trent\_\allowbreak{}defensible\_\allowbreak{}with\_\allowbreak{}pauls\_\allowbreak{}word}}\par\nobreak\noindent
\textbf{Trent, as Trent states it, can be defended in the Reformation debate, together with the arguments about Paul’s word.} A consistent set of positions that answers every attack holds Trent and both of those arguments.

This is what the two ways of taking Trent’s definition come to when weighed: accepting everything the word’s meaning shows costs Trent nothing, so long as its definition is a claim about what God does.

{\small\begin{verbatim}
theorem trent_defensible_with_pauls_word : Framework.CredulouslyAccepted
    reformationAlone.defeats ⟨Party.trent, trentStandsWithPaulsWord._proof_2⟩
-- axioms: propext, Classical.choice, Quot.sound
\end{verbatim}}
\par\medskip\noindent\textbf{\texttt{reformation\allowbreak{}Answered\allowbreak{}By\allowbreak{}Trent}}\par\nobreak\noindent
Why sola fide is not forced in the Reformation debate: a set of positions that can be defended includes Trent, and Trent defeats Paul’s case.

\paragraph{\emph{Sola fide from Paul (Galatians 2:16)} is not accepted on every resolution.}
\begin{itemize}
\item A position holding \emph{Trent, against ‘not by works’}, \emph{Paul’s word (δικαιόω, by the rule of least meaning)}, \emph{Romans 4:3–8, on Paul’s word (by the rule of least meaning)} can be held, and it holds \emph{Trent, against ‘not by works’}; \emph{Trent, against ‘not by works’} conflicts with \emph{Sola fide from Paul (Galatians 2:16)}: \emph{Trent, against ‘not by works’} defeats \emph{Sola fide from Paul (Galatians 2:16)}, and \emph{Sola fide from Paul (Galatians 2:16)} defeats \emph{Trent, against ‘not by works’}.
\item None of \emph{Trent, against ‘not by works’}, \emph{Paul’s word (δικαιόω, by the rule of least meaning)}, \emph{Romans 4:3–8, on Paul’s word (by the rule of least meaning)} defeats another, and each attack on them is answered from among them:
\begin{itemize}
\item \emph{Sola fide from Paul (Galatians 2:16)} defeats \emph{Trent, against ‘not by works’}, and \emph{Trent, against ‘not by works’} defeats \emph{Sola fide from Paul (Galatians 2:16)}.
\item \emph{Luke 7:50: “your faith has saved you”} defeats \emph{Trent, against ‘not by works’} (\texttt{why\allowbreak{}The\allowbreak{}Dominical\allowbreak{}Case\allowbreak{}Stands\allowbreak{}Against\allowbreak{}Trent}), and \emph{Trent, against ‘not by works’} defeats \emph{Luke 7:50: “your faith has saved you”}.
\item \emph{Sola fide from Peter (Acts 15:9–11)} defeats \emph{Trent, against ‘not by works’}, and \emph{Trent, against ‘not by works’} defeats \emph{Sola fide from Peter (Acts 15:9–11)}.
\item \emph{Paul’s gospel (Galatians 1), against justification as renewal} defeats \emph{Trent, against ‘not by works’} (\texttt{where\allowbreak{}Trent\allowbreak{}Parts\allowbreak{}From\allowbreak{}Paul}), and \emph{Trent, against ‘not by works’} defeats \emph{Paul’s gospel (Galatians 1), against justification as renewal}.
\item \emph{Sola fide from John (6:28–29)} defeats \emph{Trent, against ‘not by works’}, and \emph{Trent, against ‘not by works’} defeats \emph{Sola fide from John (6:28–29)}.
\item \emph{Obedience as sanctification (John 3:36, Matthew 7:21)} defeats \emph{Trent, against ‘not by works’}, and \emph{Trent, against ‘not by works’} defeats \emph{Obedience as sanctification (John 3:36, Matthew 7:21)}.
\end{itemize}
\end{itemize}
\par\noindent\textbf{What this rests on.} The reasons state 12 defeats and 33 absences of defeat, each a cell of the defeat table, \texttt{sola\allowbreak{}Fide\allowbreak{}Dispute\_\allowbreak{}defeats}, computed from the two parties' premises and checked by the kernel. Whether an attack survives turns on the attacker's weakest link and on the rating of what it attacks, so a changed rating can change the verdict. The weakest links:
\begin{itemize}
\item \emph{Trent, against ‘not by works’}, weakest at \emph{disputed}:
\begin{itemize}
\item Justification is not remission of sins only, but renewal of the inward man --- \cite[Trent, Session VI (1547), Decree on Justification, ch. 7]{tanner-decrees-1990}; \cite[Trent, Session VI (1547), canon 11]{tanner-decrees-1990}
\end{itemize}
\item \emph{Sola fide from Paul (Galatians 2:16)}, weakest at \emph{disputed}:
\begin{itemize}
\item In Galatians 2:16 ‘faith of Christ’ (πίστις Χριστοῦ) means our faith in Christ --- \cite{dunn-once-more-pistis-christou-1991}; \cite{matlock-detheologizing-2000}; Gal 2:16; Rom 3:22; Phil 3:9
\item Second Temple Judaism: in by grace, staying in by works (encoded as a denial, so it ranks disputed) --- \cite{sanders-paul-palestinian-judaism-1977}
\item Good works are the fruit and evidence of saving faith, not its ground --- \cite[XI.2; XVI.2; XVI.5]{westminster-confession-1647}; \cite[III.xvi.1]{calvin-institutes-1960}; Jas 2:18-22; Eph 2:10; Rom 6:22; Luke 17:10
\item Justification and sanctification are inseparable but distinct --- \cite[Q. 77]{westminster-larger-catechism-1647}; 1 Cor 6:11; 1 Cor 1:30; Rom 4:6; Rom 4:8; Ezek 36:27; Heb 9:13-14; Rom 3:24-25; Rom 6:6; Rom 6:14; \cite[XI.1]{westminster-confession-1647}; \cite[III.xi.6]{calvin-institutes-1960}
\item an inference step --- \cite[Trent, Session VI (1547), canon 9]{tanner-decrees-1990}
\item an inference step --- \cite{gathercole-where-boasting-2002}; \cite{carson-variegated-nomism-1-2001}
\end{itemize}
\item \emph{Paul’s word (δικαιόω, by the rule of least meaning)}, weakest at \emph{well supported}:
\begin{itemize}
\item In Paul, ‘justify’ (δικαιόω) is a verdict: to declare righteous --- \cite[s.v. δικαιόω]{bauer-danker-lexicon-2000}; \cite{morris-apostolic-preaching-1965}; \cite[p. 80]{moo-romans-1996}; \cite{irons-righteousness-of-god-2015}; \cite{wright-justification-2009}; \cite[ad loc. Rom 3:24]{fitzmyer-romans-1993}
\item A word contributes the least meaning its context requires; one occurrence does not carry the whole concept --- \cite{silva-biblical-words-1994}; \cite{joos-semantic-axiom-1972}; \cite{barr-semantics-2004}
\item an inference step --- \cite{silva-biblical-words-1994}; \cite[p. 80]{moo-romans-1996}; \cite{barr-semantics-2004}
\end{itemize}
\item \emph{Romans 4:3–8, on Paul’s word (by the rule of least meaning)}, weakest at \emph{well supported}:
\begin{itemize}
\item Romans 4:4–5 contrasts wages owed with a gift reckoned to the one who believes --- Rom 4:4-5; \cite[III.xi.18]{calvin-institutes-1960}
\item ‘Count’ (λογίζομαι) in Romans 4:3–11 is the language of an account: credited to one --- \cite[s.v. λογίζομαι]{bauer-danker-lexicon-2000}; \cite[ad loc. Rom 4:3]{moo-romans-1996}
\item A word contributes the least meaning its context requires; one occurrence does not carry the whole concept --- \cite{silva-biblical-words-1994}; \cite{joos-semantic-axiom-1972}; \cite{barr-semantics-2004}
\item an inference step --- \cite[ad loc. Rom 4:6]{moo-romans-1996}; \cite[s.v. λογίζομαι]{bauer-danker-lexicon-2000}
\item an inference step --- \cite{silva-biblical-words-1994}; \cite[p. 80]{moo-romans-1996}; \cite{barr-semantics-2004}
\end{itemize}
\item \emph{Luke 7:50: “your faith has saved you”}, weakest at \emph{disputed}:
\begin{itemize}
\item Luke 7:47 — her love is the evidence of her forgiveness, not its ground --- \cite[ad loc. Luke 7:47]{marshall-luke-1978}; \cite[54--54]{padilla-narrative-criticism-2021}; \cite[Apology IV, on love and the fulfilling of the law: Luke 7:47–50]{kolb-wengert-book-of-concord-2000}; Luke 7:41-43; Luke 7:47; Luke 7:47-50
\item an inference step --- \cite[Trent, Session VI (1547), canon 9]{tanner-decrees-1990}
\end{itemize}
\item \emph{Sola fide from Peter (Acts 15:9–11)}, weakest at \emph{disputed}:
\begin{itemize}
\item The yoke refused at Acts 15:10 is the whole law as a condition of salvation --- \cite[ad loc. Acts 15:10]{bruce-acts-1988}; Acts 15:10; Gal 5:1; Jas 2:10; \cite{das-paul-law-covenant-2001}
\item Good works are the fruit and evidence of saving faith, not its ground --- \cite[XI.2; XVI.2; XVI.5]{westminster-confession-1647}; \cite[III.xvi.1]{calvin-institutes-1960}; Jas 2:18-22; Eph 2:10; Rom 6:22; Luke 17:10
\item Justification and sanctification are inseparable but distinct --- \cite[Q. 77]{westminster-larger-catechism-1647}; 1 Cor 6:11; 1 Cor 1:30; Rom 4:6; Rom 4:8; Ezek 36:27; Heb 9:13-14; Rom 3:24-25; Rom 6:6; Rom 6:14; \cite[XI.1]{westminster-confession-1647}; \cite[III.xi.6]{calvin-institutes-1960}
\item an inference step --- \cite[Trent, Session VI (1547), canon 9]{tanner-decrees-1990}
\end{itemize}
\item \emph{Paul’s gospel (Galatians 1), against justification as renewal}, weakest at \emph{disputed}:
\begin{itemize}
\item an inference step --- \cite[ad loc. Gal 1:6]{moo-galatians-2013}; \cite[XI.1]{westminster-confession-1647}
\item an inference step --- \cite[XI.1]{westminster-confession-1647}; \cite[III.xi.2]{calvin-institutes-1960}; \cite{morris-apostolic-preaching-1965}
\end{itemize}
\item \emph{Trent, against ‘not by works’, read as what Paul’s word means}, weakest at \emph{disputed}:
\begin{itemize}
\item Justification is not remission of sins only, but renewal of the inward man --- \cite[Trent, Session VI (1547), Decree on Justification, ch. 7]{tanner-decrees-1990}; \cite[Trent, Session VI (1547), canon 11]{tanner-decrees-1990}
\item an inference step --- \cite[Trent, Session VI (1547), Decree on Justification, ch. 8]{tanner-decrees-1990}; \cite[Trent, Session VI (1547), Decree on Justification, ch. 7]{tanner-decrees-1990}; \cite[On the Spirit and the Letter, ch. 45 (26.45)]{npnf1-05-1887}; \cite{mcgrath-iustitia-dei-2005}
\end{itemize}
\item \emph{Sola fide from John (6:28–29)}, weakest at \emph{disputed}:
\begin{itemize}
\item The believing of John 6:29 is trust, which brings nothing and receives Christ --- \cite[ad loc. John 6:29]{calvin-commentary-john-1847}; John 6:28-29
\item Good works are the fruit and evidence of saving faith, not its ground --- \cite[XI.2; XVI.2; XVI.5]{westminster-confession-1647}; \cite[III.xvi.1]{calvin-institutes-1960}; Jas 2:18-22; Eph 2:10; Rom 6:22; Luke 17:10
\item Justification and sanctification are inseparable but distinct --- \cite[Q. 77]{westminster-larger-catechism-1647}; 1 Cor 6:11; 1 Cor 1:30; Rom 4:6; Rom 4:8; Ezek 36:27; Heb 9:13-14; Rom 3:24-25; Rom 6:6; Rom 6:14; \cite[XI.1]{westminster-confession-1647}; \cite[III.xi.6]{calvin-institutes-1960}
\item an inference step --- \cite[Trent, Session VI (1547), canon 9]{tanner-decrees-1990}
\end{itemize}
\item \emph{Faith formed by charity (Aquinas on John 6:29)}, weakest at \emph{disputed}:
\begin{itemize}
\item The believing of John 6:29 is faith living through charity, the source of works --- \cite[cap. 6, lect. 3, n. 901]{aquinas-commentary-john-2012}; \cite[Trent, Session VI (1547), Decree on Justification, ch. 7]{tanner-decrees-1990}
\end{itemize}
\item \emph{Luther on Galatians 3:11–12 (against faith formed by charity)}, weakest at \emph{disputed}:
\begin{itemize}
\item an inference step --- \cite[ad loc. Gal 3:12]{luther-commentary-galatians-1998}; \cite[ad loc. Gal 3:11]{luther-commentary-galatians-1998}
\end{itemize}
\item \emph{Chrysostom on John 3:36 (believing is not enough without a right life)}, weakest at \emph{disputed}:
\begin{itemize}
\item an inference step --- \cite[Homily 31, §1]{chrysostom-homilies-john}; \cite{stanley-did-jesus-teach-2006}
\end{itemize}
\item \emph{Obedience as sanctification (John 3:36, Matthew 7:21)}, weakest at \emph{disputed}:
\begin{itemize}
\item The obedience of John 3:36 and Matthew 7:21 is sanctification, the fruit of faith --- \cite[ad loc. John 3:36]{ridderbos-john-1997}; \cite[ad loc. John 3:36]{carson-john-1991}; \cite[ad loc. Matt 7:21]{morris-matthew-1992}; \cite[55--55]{schreiner-justification-by-works-2015}; \cite{hoekema-saved-by-grace-1994}; \cite{murray-redemption-1984}; \cite[XVI.2]{westminster-confession-1647}; \cite[XIII.1]{westminster-confession-1647}; \cite[ad loc. John 3:36]{calvin-commentary-john-1847}
\item Justification and sanctification are inseparable but distinct --- \cite[Q. 77]{westminster-larger-catechism-1647}; 1 Cor 6:11; 1 Cor 1:30; Rom 4:6; Rom 4:8; Ezek 36:27; Heb 9:13-14; Rom 3:24-25; Rom 6:6; Rom 6:14; \cite[XI.1]{westminster-confession-1647}; \cite[III.xi.6]{calvin-institutes-1960}
\item an inference step --- \cite[55--55]{schreiner-justification-by-works-2015}; \cite[ad loc. John 3:36]{calvin-commentary-john-1847}; \cite[XVI.2]{westminster-confession-1647}; \cite[XIII.1]{westminster-confession-1647}
\end{itemize}
\end{itemize}
\noindent Why \emph{Luke 7:50: “your faith has saved you”} defeats \emph{Trent, against ‘not by works’}, at the level of the claims: \texttt{why\allowbreak{}The\allowbreak{}Dominical\allowbreak{}Case\allowbreak{}Stands\allowbreak{}Against\allowbreak{}Trent}.
\noindent Why \emph{Paul’s gospel (Galatians 1), against justification as renewal} defeats \emph{Trent, against ‘not by works’}, at the level of the claims: \texttt{where\allowbreak{}Trent\allowbreak{}Parts\allowbreak{}From\allowbreak{}Paul}.
\par\medskip\noindent\textbf{\texttt{sola\_\allowbreak{}fide\_\allowbreak{}not\_\allowbreak{}forced\_\allowbreak{}in\_\allowbreak{}the\_\allowbreak{}reformation\_\allowbreak{}dispute}}\par\nobreak\noindent
\textbf{Nor is sola fide forced there.} Rome and the Reformers, heard alone, tie. Each Reformed strand and Trent defeat each other over Trent’s definition, and every rating on both sides is \texttt{disputed}. What would decide it is the step from a verdict to “renewal is no part of justification”, which is \texttt{disputed} in Galatians and again in Romans 4 (see \texttt{what\allowbreak{}Trents\allowbreak{}Definition\allowbreak{}Claims}).

{\small\begin{verbatim}
theorem sola_fide_not_forced_in_the_reformation_dispute :
    ¬Framework.SkepticallyAccepted reformationAlone.defeats ⟨Party.pauline,
    reformationAnsweredByTrent._proof_2⟩
-- axioms: propext, Classical.choice, Quot.sound
\end{verbatim}}
\paragraph*{The debate without Trent}
\leavevmode\par\nobreak
\par\medskip\noindent\textbf{\texttt{without\allowbreak{}Trent}}\par\nobreak\noindent
Every position but Trent’s two: Trent as it states its definition, and Trent taken as a claim about Paul’s word. Augustine’s explanation of “justify” goes with the second, since that is what it rests on.

{\small\begin{verbatim}
def withoutTrent : Dispute Claim { i // i ∉ [Party.trent,
    Party.trentOnPaulsWord] } :=
  solaFideDispute.restrict fun x =>
    x ∉ [Party.trent, Party.trentOnPaulsWord]
\end{verbatim}}
\par\medskip\noindent\textbf{\texttt{without\allowbreak{}Trent\allowbreak{}Settles}}\par\nobreak\noindent
What is settled without Trent. Nothing now defeats Paul’s gospel or either argument about Paul’s word, so those three are accepted first, and nothing joins them. Luke’s case is defeated by Chrysostom, and Chrysostom by Luke’s case. Paul’s case is defeated by the apocalyptic reading, Peter’s by Jervell, John’s by Aquinas. Sanders is defeated by the critics and the critics by Sanders, the apocalyptic reading by Paul’s case, Jervell by Peter’s case, Aquinas by Luther and Luther by Aquinas, and the Reformed answer to Chrysostom by Chrysostom.

\paragraph{What the dispute forces is exactly \emph{Paul’s gospel (Galatians 1), against justification as renewal}, \emph{Paul’s word (δικαιόω, by the rule of least meaning)}, \emph{Romans 4:3–8, on Paul’s word (by the rule of least meaning)}.}
\begin{itemize}
\item Stage 1: \emph{Paul’s gospel (Galatians 1), against justification as renewal}, \emph{Paul’s word (δικαιόω, by the rule of least meaning)}, \emph{Romans 4:3–8, on Paul’s word (by the rule of least meaning)}.
\begin{itemize}
\item \emph{Paul’s gospel (Galatians 1), against justification as renewal} is defeated by nothing.
\item \emph{Paul’s word (δικαιόω, by the rule of least meaning)} is defeated by nothing.
\item \emph{Romans 4:3–8, on Paul’s word (by the rule of least meaning)} is defeated by nothing.
\end{itemize}
\item Nothing else is forced:
\begin{itemize}
\item \emph{Apocalyptic reading, against faith as the condition} defeats \emph{Sola fide from Paul (Galatians 2:16)}, and nothing forced defeats \emph{Apocalyptic reading, against faith as the condition}.
\item \emph{Chrysostom on John 3:36 (believing is not enough without a right life)} defeats \emph{Luke 7:50: “your faith has saved you”}, and nothing forced defeats \emph{Chrysostom on John 3:36 (believing is not enough without a right life)}.
\item \emph{Law-observant Luke (Jervell)} defeats \emph{Sola fide from Peter (Acts 15:9–11)}, and nothing forced defeats \emph{Law-observant Luke (Jervell)}.
\item \emph{Sola fide from Paul (Galatians 2:16)} defeats \emph{Apocalyptic reading, against faith as the condition}, and nothing forced defeats \emph{Sola fide from Paul (Galatians 2:16)}.
\item \emph{Variegated nomism (Gathercole, Carson et al.)} defeats \emph{Covenantal nomism (Sanders, Dunn)}, and nothing forced defeats \emph{Variegated nomism (Gathercole, Carson et al.)}.
\item \emph{Covenantal nomism (Sanders, Dunn)} defeats \emph{Variegated nomism (Gathercole, Carson et al.)}, and nothing forced defeats \emph{Covenantal nomism (Sanders, Dunn)}.
\item \emph{Sola fide from Peter (Acts 15:9–11)} defeats \emph{Law-observant Luke (Jervell)}, and nothing forced defeats \emph{Sola fide from Peter (Acts 15:9–11)}.
\item \emph{Faith formed by charity (Aquinas on John 6:29)} defeats \emph{Sola fide from John (6:28–29)}, and nothing forced defeats \emph{Faith formed by charity (Aquinas on John 6:29)}.
\item \emph{Luther on Galatians 3:11–12 (against faith formed by charity)} defeats \emph{Faith formed by charity (Aquinas on John 6:29)}, and nothing forced defeats \emph{Luther on Galatians 3:11–12 (against faith formed by charity)}.
\item \emph{Faith formed by charity (Aquinas on John 6:29)} defeats \emph{Luther on Galatians 3:11–12 (against faith formed by charity)}, and nothing forced defeats \emph{Faith formed by charity (Aquinas on John 6:29)}.
\item \emph{Luke 7:50: “your faith has saved you”} defeats \emph{Chrysostom on John 3:36 (believing is not enough without a right life)}, and nothing forced defeats \emph{Luke 7:50: “your faith has saved you”}.
\item \emph{Chrysostom on John 3:36 (believing is not enough without a right life)} defeats \emph{Obedience as sanctification (John 3:36, Matthew 7:21)}, and nothing forced defeats \emph{Chrysostom on John 3:36 (believing is not enough without a right life)}.
\end{itemize}
\end{itemize}
\par\noindent\textbf{What this rests on.} The reasons state 12 defeats and 75 absences of defeat, each a cell of the defeat table, \texttt{sola\allowbreak{}Fide\allowbreak{}Dispute\_\allowbreak{}defeats}, computed from the two parties' premises and checked by the kernel. Whether an attack survives turns on the attacker's weakest link and on the rating of what it attacks, so a changed rating can change the verdict. The weakest links:
\begin{itemize}
\item \emph{Sola fide from Paul (Galatians 2:16)}, weakest at \emph{disputed}:
\begin{itemize}
\item In Galatians 2:16 ‘faith of Christ’ (πίστις Χριστοῦ) means our faith in Christ --- \cite{dunn-once-more-pistis-christou-1991}; \cite{matlock-detheologizing-2000}; Gal 2:16; Rom 3:22; Phil 3:9
\item Second Temple Judaism: in by grace, staying in by works (encoded as a denial, so it ranks disputed) --- \cite{sanders-paul-palestinian-judaism-1977}
\item Good works are the fruit and evidence of saving faith, not its ground --- \cite[XI.2; XVI.2; XVI.5]{westminster-confession-1647}; \cite[III.xvi.1]{calvin-institutes-1960}; Jas 2:18-22; Eph 2:10; Rom 6:22; Luke 17:10
\item Justification and sanctification are inseparable but distinct --- \cite[Q. 77]{westminster-larger-catechism-1647}; 1 Cor 6:11; 1 Cor 1:30; Rom 4:6; Rom 4:8; Ezek 36:27; Heb 9:13-14; Rom 3:24-25; Rom 6:6; Rom 6:14; \cite[XI.1]{westminster-confession-1647}; \cite[III.xi.6]{calvin-institutes-1960}
\item an inference step --- \cite[Trent, Session VI (1547), canon 9]{tanner-decrees-1990}
\item an inference step --- \cite{gathercole-where-boasting-2002}; \cite{carson-variegated-nomism-1-2001}
\end{itemize}
\item \emph{Paul’s gospel (Galatians 1), against justification as renewal}, weakest at \emph{disputed}:
\begin{itemize}
\item an inference step --- \cite[ad loc. Gal 1:6]{moo-galatians-2013}; \cite[XI.1]{westminster-confession-1647}
\item an inference step --- \cite[XI.1]{westminster-confession-1647}; \cite[III.xi.2]{calvin-institutes-1960}; \cite{morris-apostolic-preaching-1965}
\end{itemize}
\item \emph{Luke 7:50: “your faith has saved you”}, weakest at \emph{disputed}:
\begin{itemize}
\item Luke 7:47 — her love is the evidence of her forgiveness, not its ground --- \cite[ad loc. Luke 7:47]{marshall-luke-1978}; \cite[54--54]{padilla-narrative-criticism-2021}; \cite[Apology IV, on love and the fulfilling of the law: Luke 7:47–50]{kolb-wengert-book-of-concord-2000}; Luke 7:41-43; Luke 7:47; Luke 7:47-50
\item an inference step --- \cite[Trent, Session VI (1547), canon 9]{tanner-decrees-1990}
\end{itemize}
\item \emph{Sola fide from Peter (Acts 15:9–11)}, weakest at \emph{disputed}:
\begin{itemize}
\item The yoke refused at Acts 15:10 is the whole law as a condition of salvation --- \cite[ad loc. Acts 15:10]{bruce-acts-1988}; Acts 15:10; Gal 5:1; Jas 2:10; \cite{das-paul-law-covenant-2001}
\item Good works are the fruit and evidence of saving faith, not its ground --- \cite[XI.2; XVI.2; XVI.5]{westminster-confession-1647}; \cite[III.xvi.1]{calvin-institutes-1960}; Jas 2:18-22; Eph 2:10; Rom 6:22; Luke 17:10
\item Justification and sanctification are inseparable but distinct --- \cite[Q. 77]{westminster-larger-catechism-1647}; 1 Cor 6:11; 1 Cor 1:30; Rom 4:6; Rom 4:8; Ezek 36:27; Heb 9:13-14; Rom 3:24-25; Rom 6:6; Rom 6:14; \cite[XI.1]{westminster-confession-1647}; \cite[III.xi.6]{calvin-institutes-1960}
\item an inference step --- \cite[Trent, Session VI (1547), canon 9]{tanner-decrees-1990}
\end{itemize}
\item \emph{Apocalyptic reading, against faith as the condition}, weakest at \emph{disputed}:
\begin{itemize}
\item ‘Faith of Christ’ is Christ’s own faithfulness, the means of God’s rescue (encoded as a denial, so it ranks disputed) --- \cite{campbell-deliverance-god-2009}; \cite[ad loc. Gal 2:16]{martyn-galatians-1997}; \cite{hays-faith-jesus-christ-2002}
\item ‘The righteousness of God’ (δικαιοσύνη θεοῦ) names God’s act of rescue in Christ --- \cite{campbell-deliverance-god-2009}; \cite{martyn-galatians-1997}; Rom 3:21-26
\end{itemize}
\item \emph{Covenantal nomism (Sanders, Dunn)}, weakest at \emph{disputed}:
\begin{itemize}
\item Second Temple Judaism: in by grace, staying in by works --- \cite{sanders-paul-palestinian-judaism-1977}
\item an inference step --- \cite{dunn-new-perspective-2005}; \cite{wright-what-paul-said-1997}
\end{itemize}
\item \emph{Variegated nomism (Gathercole, Carson et al.)}, weakest at \emph{disputed}:
\begin{itemize}
\item Second Temple Judaism also held final vindication according to works (encoded as a denial, so it ranks disputed) --- \cite{gathercole-where-boasting-2002}; \cite{carson-variegated-nomism-1-2001}
\item an inference step --- \cite{gathercole-where-boasting-2002}; \cite{carson-variegated-nomism-1-2001}
\end{itemize}
\item \emph{Law-observant Luke (Jervell)}, weakest at \emph{disputed}:
\begin{itemize}
\item an inference step --- \cite{jervell-luke-people-god-1972}
\end{itemize}
\item \emph{Paul’s word (δικαιόω, by the rule of least meaning)}, weakest at \emph{well supported}:
\begin{itemize}
\item In Paul, ‘justify’ (δικαιόω) is a verdict: to declare righteous --- \cite[s.v. δικαιόω]{bauer-danker-lexicon-2000}; \cite{morris-apostolic-preaching-1965}; \cite[p. 80]{moo-romans-1996}; \cite{irons-righteousness-of-god-2015}; \cite{wright-justification-2009}; \cite[ad loc. Rom 3:24]{fitzmyer-romans-1993}
\item A word contributes the least meaning its context requires; one occurrence does not carry the whole concept --- \cite{silva-biblical-words-1994}; \cite{joos-semantic-axiom-1972}; \cite{barr-semantics-2004}
\item an inference step --- \cite{silva-biblical-words-1994}; \cite[p. 80]{moo-romans-1996}; \cite{barr-semantics-2004}
\end{itemize}
\item \emph{Romans 4:3–8, on Paul’s word (by the rule of least meaning)}, weakest at \emph{well supported}:
\begin{itemize}
\item Romans 4:4–5 contrasts wages owed with a gift reckoned to the one who believes --- Rom 4:4-5; \cite[III.xi.18]{calvin-institutes-1960}
\item ‘Count’ (λογίζομαι) in Romans 4:3–11 is the language of an account: credited to one --- \cite[s.v. λογίζομαι]{bauer-danker-lexicon-2000}; \cite[ad loc. Rom 4:3]{moo-romans-1996}
\item A word contributes the least meaning its context requires; one occurrence does not carry the whole concept --- \cite{silva-biblical-words-1994}; \cite{joos-semantic-axiom-1972}; \cite{barr-semantics-2004}
\item an inference step --- \cite[ad loc. Rom 4:6]{moo-romans-1996}; \cite[s.v. λογίζομαι]{bauer-danker-lexicon-2000}
\item an inference step --- \cite{silva-biblical-words-1994}; \cite[p. 80]{moo-romans-1996}; \cite{barr-semantics-2004}
\end{itemize}
\item \emph{Sola fide from John (6:28–29)}, weakest at \emph{disputed}:
\begin{itemize}
\item The believing of John 6:29 is trust, which brings nothing and receives Christ --- \cite[ad loc. John 6:29]{calvin-commentary-john-1847}; John 6:28-29
\item Good works are the fruit and evidence of saving faith, not its ground --- \cite[XI.2; XVI.2; XVI.5]{westminster-confession-1647}; \cite[III.xvi.1]{calvin-institutes-1960}; Jas 2:18-22; Eph 2:10; Rom 6:22; Luke 17:10
\item Justification and sanctification are inseparable but distinct --- \cite[Q. 77]{westminster-larger-catechism-1647}; 1 Cor 6:11; 1 Cor 1:30; Rom 4:6; Rom 4:8; Ezek 36:27; Heb 9:13-14; Rom 3:24-25; Rom 6:6; Rom 6:14; \cite[XI.1]{westminster-confession-1647}; \cite[III.xi.6]{calvin-institutes-1960}
\item an inference step --- \cite[Trent, Session VI (1547), canon 9]{tanner-decrees-1990}
\end{itemize}
\item \emph{Faith formed by charity (Aquinas on John 6:29)}, weakest at \emph{disputed}:
\begin{itemize}
\item The believing of John 6:29 is faith living through charity, the source of works --- \cite[cap. 6, lect. 3, n. 901]{aquinas-commentary-john-2012}; \cite[Trent, Session VI (1547), Decree on Justification, ch. 7]{tanner-decrees-1990}
\end{itemize}
\item \emph{Luther on Galatians 3:11–12 (against faith formed by charity)}, weakest at \emph{disputed}:
\begin{itemize}
\item an inference step --- \cite[ad loc. Gal 3:12]{luther-commentary-galatians-1998}; \cite[ad loc. Gal 3:11]{luther-commentary-galatians-1998}
\end{itemize}
\item \emph{Chrysostom on John 3:36 (believing is not enough without a right life)}, weakest at \emph{disputed}:
\begin{itemize}
\item an inference step --- \cite[Homily 31, §1]{chrysostom-homilies-john}; \cite{stanley-did-jesus-teach-2006}
\end{itemize}
\item \emph{Obedience as sanctification (John 3:36, Matthew 7:21)}, weakest at \emph{disputed}:
\begin{itemize}
\item The obedience of John 3:36 and Matthew 7:21 is sanctification, the fruit of faith --- \cite[ad loc. John 3:36]{ridderbos-john-1997}; \cite[ad loc. John 3:36]{carson-john-1991}; \cite[ad loc. Matt 7:21]{morris-matthew-1992}; \cite[55--55]{schreiner-justification-by-works-2015}; \cite{hoekema-saved-by-grace-1994}; \cite{murray-redemption-1984}; \cite[XVI.2]{westminster-confession-1647}; \cite[XIII.1]{westminster-confession-1647}; \cite[ad loc. John 3:36]{calvin-commentary-john-1847}
\item Justification and sanctification are inseparable but distinct --- \cite[Q. 77]{westminster-larger-catechism-1647}; 1 Cor 6:11; 1 Cor 1:30; Rom 4:6; Rom 4:8; Ezek 36:27; Heb 9:13-14; Rom 3:24-25; Rom 6:6; Rom 6:14; \cite[XI.1]{westminster-confession-1647}; \cite[III.xi.6]{calvin-institutes-1960}
\item an inference step --- \cite[55--55]{schreiner-justification-by-works-2015}; \cite[ad loc. John 3:36]{calvin-commentary-john-1847}; \cite[XVI.2]{westminster-confession-1647}; \cite[XIII.1]{westminster-confession-1647}
\end{itemize}
\end{itemize}
\par\medskip\noindent\textbf{\texttt{without\_\allowbreak{}trent\_\allowbreak{}pauls\_\allowbreak{}gospel\_\allowbreak{}and\_\allowbreak{}word\_\allowbreak{}prevail}}\par\nobreak\noindent
\textbf{Without Trent, two things are forced.} That justification is not the renewal of the inward person, from Paul’s gospel. And that Paul’s “justify” (δικαιόω) does not mean that renewal, argued two ways. With Trent not heard, nothing attacks either.

Luke’s “faith suffices” is no longer among them. Chrysostom, reading John 3:36 and Matthew 7:21, concludes that believing is not enough without a right life, and so stands against it (\texttt{luke\_\allowbreak{}not\_\allowbreak{}forced\_\allowbreak{}without\_\allowbreak{}trent}).

It does not show that \emph{faith alone}, or even “faith suffices”, is refuted. Each is in a tie; see the results that follow.

{\small\begin{verbatim}
theorem without_trent_pauls_gospel_and_word_prevail : Framework.grounded
    withoutTrent.defeats = Solver.toSet (Witness.listSub (fun x => x ∉
    [Party.trent, Party.trentOnPaulsWord]) [Party.gospel, Party.lexical,
    Party.romansOnTheWord])
-- axioms: propext, Classical.choice, Quot.sound
\end{verbatim}}
\par\medskip\noindent\textbf{\texttt{luke\allowbreak{}Answered\allowbreak{}By\allowbreak{}Chrysostom}}\par\nobreak\noindent
Why Luke’s case is not forced without Trent: a set of positions that can be defended includes Chrysostom, who defeats it.

\paragraph{\emph{Luke 7:50: “your faith has saved you”} is not accepted on every resolution.}
\begin{itemize}
\item A position holding \emph{Apocalyptic reading, against faith as the condition}, \emph{Faith formed by charity (Aquinas on John 6:29)}, \emph{Chrysostom on John 3:36 (believing is not enough without a right life)}, \emph{Variegated nomism (Gathercole, Carson et al.)}, \emph{Paul’s gospel (Galatians 1), against justification as renewal}, \emph{Law-observant Luke (Jervell)}, \emph{Paul’s word (δικαιόω, by the rule of least meaning)}, \emph{Romans 4:3–8, on Paul’s word (by the rule of least meaning)} can be held, and it holds \emph{Chrysostom on John 3:36 (believing is not enough without a right life)}; \emph{Chrysostom on John 3:36 (believing is not enough without a right life)} conflicts with \emph{Luke 7:50: “your faith has saved you”}: \emph{Chrysostom on John 3:36 (believing is not enough without a right life)} defeats \emph{Luke 7:50: “your faith has saved you”}, and \emph{Luke 7:50: “your faith has saved you”} defeats \emph{Chrysostom on John 3:36 (believing is not enough without a right life)}.
\item None of \emph{Apocalyptic reading, against faith as the condition}, \emph{Faith formed by charity (Aquinas on John 6:29)}, \emph{Chrysostom on John 3:36 (believing is not enough without a right life)}, \emph{Variegated nomism (Gathercole, Carson et al.)}, \emph{Paul’s gospel (Galatians 1), against justification as renewal}, \emph{Law-observant Luke (Jervell)}, \emph{Paul’s word (δικαιόω, by the rule of least meaning)}, \emph{Romans 4:3–8, on Paul’s word (by the rule of least meaning)} defeats another, and each attack on them is answered from among them:
\begin{itemize}
\item \emph{Sola fide from Paul (Galatians 2:16)} defeats \emph{Apocalyptic reading, against faith as the condition}, and \emph{Apocalyptic reading, against faith as the condition} defeats \emph{Sola fide from Paul (Galatians 2:16)}.
\item \emph{Sola fide from John (6:28–29)} defeats \emph{Apocalyptic reading, against faith as the condition}, and \emph{Apocalyptic reading, against faith as the condition} defeats \emph{Sola fide from John (6:28–29)}.
\item \emph{Sola fide from John (6:28–29)} defeats \emph{Faith formed by charity (Aquinas on John 6:29)}, and \emph{Apocalyptic reading, against faith as the condition} defeats \emph{Sola fide from John (6:28–29)}.
\item \emph{Luther on Galatians 3:11–12 (against faith formed by charity)} defeats \emph{Faith formed by charity (Aquinas on John 6:29)}, and \emph{Faith formed by charity (Aquinas on John 6:29)} defeats \emph{Luther on Galatians 3:11–12 (against faith formed by charity)}.
\item \emph{Sola fide from Paul (Galatians 2:16)} defeats \emph{Chrysostom on John 3:36 (believing is not enough without a right life)}, and \emph{Apocalyptic reading, against faith as the condition} defeats \emph{Sola fide from Paul (Galatians 2:16)}.
\item \emph{Luke 7:50: “your faith has saved you”} defeats \emph{Chrysostom on John 3:36 (believing is not enough without a right life)}, and \emph{Chrysostom on John 3:36 (believing is not enough without a right life)} defeats \emph{Luke 7:50: “your faith has saved you”}.
\item \emph{Sola fide from Peter (Acts 15:9–11)} defeats \emph{Chrysostom on John 3:36 (believing is not enough without a right life)}, and \emph{Chrysostom on John 3:36 (believing is not enough without a right life)} defeats \emph{Sola fide from Peter (Acts 15:9–11)}.
\item \emph{Sola fide from John (6:28–29)} defeats \emph{Chrysostom on John 3:36 (believing is not enough without a right life)}, and \emph{Apocalyptic reading, against faith as the condition} defeats \emph{Sola fide from John (6:28–29)}.
\item \emph{Obedience as sanctification (John 3:36, Matthew 7:21)} defeats \emph{Chrysostom on John 3:36 (believing is not enough without a right life)}, and \emph{Chrysostom on John 3:36 (believing is not enough without a right life)} defeats \emph{Obedience as sanctification (John 3:36, Matthew 7:21)}.
\item \emph{Covenantal nomism (Sanders, Dunn)} defeats \emph{Variegated nomism (Gathercole, Carson et al.)}, and \emph{Variegated nomism (Gathercole, Carson et al.)} defeats \emph{Covenantal nomism (Sanders, Dunn)}.
\item \emph{Sola fide from Peter (Acts 15:9–11)} defeats \emph{Law-observant Luke (Jervell)}, and \emph{Chrysostom on John 3:36 (believing is not enough without a right life)} defeats \emph{Sola fide from Peter (Acts 15:9–11)}.
\end{itemize}
\end{itemize}
\par\noindent\textbf{What this rests on.} The reasons state 18 defeats and 109 absences of defeat, each a cell of the defeat table, \texttt{sola\allowbreak{}Fide\allowbreak{}Dispute\_\allowbreak{}defeats}, computed from the two parties' premises and checked by the kernel. Whether an attack survives turns on the attacker's weakest link and on the rating of what it attacks, so a changed rating can change the verdict. The weakest links:
\begin{itemize}
\item \emph{Chrysostom on John 3:36 (believing is not enough without a right life)}, weakest at \emph{disputed}:
\begin{itemize}
\item an inference step --- \cite[Homily 31, §1]{chrysostom-homilies-john}; \cite{stanley-did-jesus-teach-2006}
\end{itemize}
\item \emph{Luke 7:50: “your faith has saved you”}, weakest at \emph{disputed}:
\begin{itemize}
\item Luke 7:47 — her love is the evidence of her forgiveness, not its ground --- \cite[ad loc. Luke 7:47]{marshall-luke-1978}; \cite[54--54]{padilla-narrative-criticism-2021}; \cite[Apology IV, on love and the fulfilling of the law: Luke 7:47–50]{kolb-wengert-book-of-concord-2000}; Luke 7:41-43; Luke 7:47; Luke 7:47-50
\item an inference step --- \cite[Trent, Session VI (1547), canon 9]{tanner-decrees-1990}
\end{itemize}
\item \emph{Apocalyptic reading, against faith as the condition}, weakest at \emph{disputed}:
\begin{itemize}
\item ‘Faith of Christ’ is Christ’s own faithfulness, the means of God’s rescue (encoded as a denial, so it ranks disputed) --- \cite{campbell-deliverance-god-2009}; \cite[ad loc. Gal 2:16]{martyn-galatians-1997}; \cite{hays-faith-jesus-christ-2002}
\item ‘The righteousness of God’ (δικαιοσύνη θεοῦ) names God’s act of rescue in Christ --- \cite{campbell-deliverance-god-2009}; \cite{martyn-galatians-1997}; Rom 3:21-26
\end{itemize}
\item \emph{Faith formed by charity (Aquinas on John 6:29)}, weakest at \emph{disputed}:
\begin{itemize}
\item The believing of John 6:29 is faith living through charity, the source of works --- \cite[cap. 6, lect. 3, n. 901]{aquinas-commentary-john-2012}; \cite[Trent, Session VI (1547), Decree on Justification, ch. 7]{tanner-decrees-1990}
\end{itemize}
\item \emph{Variegated nomism (Gathercole, Carson et al.)}, weakest at \emph{disputed}:
\begin{itemize}
\item Second Temple Judaism also held final vindication according to works (encoded as a denial, so it ranks disputed) --- \cite{gathercole-where-boasting-2002}; \cite{carson-variegated-nomism-1-2001}
\item an inference step --- \cite{gathercole-where-boasting-2002}; \cite{carson-variegated-nomism-1-2001}
\end{itemize}
\item \emph{Paul’s gospel (Galatians 1), against justification as renewal}, weakest at \emph{disputed}:
\begin{itemize}
\item an inference step --- \cite[ad loc. Gal 1:6]{moo-galatians-2013}; \cite[XI.1]{westminster-confession-1647}
\item an inference step --- \cite[XI.1]{westminster-confession-1647}; \cite[III.xi.2]{calvin-institutes-1960}; \cite{morris-apostolic-preaching-1965}
\end{itemize}
\item \emph{Law-observant Luke (Jervell)}, weakest at \emph{disputed}:
\begin{itemize}
\item an inference step --- \cite{jervell-luke-people-god-1972}
\end{itemize}
\item \emph{Paul’s word (δικαιόω, by the rule of least meaning)}, weakest at \emph{well supported}:
\begin{itemize}
\item In Paul, ‘justify’ (δικαιόω) is a verdict: to declare righteous --- \cite[s.v. δικαιόω]{bauer-danker-lexicon-2000}; \cite{morris-apostolic-preaching-1965}; \cite[p. 80]{moo-romans-1996}; \cite{irons-righteousness-of-god-2015}; \cite{wright-justification-2009}; \cite[ad loc. Rom 3:24]{fitzmyer-romans-1993}
\item A word contributes the least meaning its context requires; one occurrence does not carry the whole concept --- \cite{silva-biblical-words-1994}; \cite{joos-semantic-axiom-1972}; \cite{barr-semantics-2004}
\item an inference step --- \cite{silva-biblical-words-1994}; \cite[p. 80]{moo-romans-1996}; \cite{barr-semantics-2004}
\end{itemize}
\item \emph{Romans 4:3–8, on Paul’s word (by the rule of least meaning)}, weakest at \emph{well supported}:
\begin{itemize}
\item Romans 4:4–5 contrasts wages owed with a gift reckoned to the one who believes --- Rom 4:4-5; \cite[III.xi.18]{calvin-institutes-1960}
\item ‘Count’ (λογίζομαι) in Romans 4:3–11 is the language of an account: credited to one --- \cite[s.v. λογίζομαι]{bauer-danker-lexicon-2000}; \cite[ad loc. Rom 4:3]{moo-romans-1996}
\item A word contributes the least meaning its context requires; one occurrence does not carry the whole concept --- \cite{silva-biblical-words-1994}; \cite{joos-semantic-axiom-1972}; \cite{barr-semantics-2004}
\item an inference step --- \cite[ad loc. Rom 4:6]{moo-romans-1996}; \cite[s.v. λογίζομαι]{bauer-danker-lexicon-2000}
\item an inference step --- \cite{silva-biblical-words-1994}; \cite[p. 80]{moo-romans-1996}; \cite{barr-semantics-2004}
\end{itemize}
\item \emph{Sola fide from Paul (Galatians 2:16)}, weakest at \emph{disputed}:
\begin{itemize}
\item In Galatians 2:16 ‘faith of Christ’ (πίστις Χριστοῦ) means our faith in Christ --- \cite{dunn-once-more-pistis-christou-1991}; \cite{matlock-detheologizing-2000}; Gal 2:16; Rom 3:22; Phil 3:9
\item Second Temple Judaism: in by grace, staying in by works (encoded as a denial, so it ranks disputed) --- \cite{sanders-paul-palestinian-judaism-1977}
\item Good works are the fruit and evidence of saving faith, not its ground --- \cite[XI.2; XVI.2; XVI.5]{westminster-confession-1647}; \cite[III.xvi.1]{calvin-institutes-1960}; Jas 2:18-22; Eph 2:10; Rom 6:22; Luke 17:10
\item Justification and sanctification are inseparable but distinct --- \cite[Q. 77]{westminster-larger-catechism-1647}; 1 Cor 6:11; 1 Cor 1:30; Rom 4:6; Rom 4:8; Ezek 36:27; Heb 9:13-14; Rom 3:24-25; Rom 6:6; Rom 6:14; \cite[XI.1]{westminster-confession-1647}; \cite[III.xi.6]{calvin-institutes-1960}
\item an inference step --- \cite[Trent, Session VI (1547), canon 9]{tanner-decrees-1990}
\item an inference step --- \cite{gathercole-where-boasting-2002}; \cite{carson-variegated-nomism-1-2001}
\end{itemize}
\item \emph{Sola fide from Peter (Acts 15:9–11)}, weakest at \emph{disputed}:
\begin{itemize}
\item The yoke refused at Acts 15:10 is the whole law as a condition of salvation --- \cite[ad loc. Acts 15:10]{bruce-acts-1988}; Acts 15:10; Gal 5:1; Jas 2:10; \cite{das-paul-law-covenant-2001}
\item Good works are the fruit and evidence of saving faith, not its ground --- \cite[XI.2; XVI.2; XVI.5]{westminster-confession-1647}; \cite[III.xvi.1]{calvin-institutes-1960}; Jas 2:18-22; Eph 2:10; Rom 6:22; Luke 17:10
\item Justification and sanctification are inseparable but distinct --- \cite[Q. 77]{westminster-larger-catechism-1647}; 1 Cor 6:11; 1 Cor 1:30; Rom 4:6; Rom 4:8; Ezek 36:27; Heb 9:13-14; Rom 3:24-25; Rom 6:6; Rom 6:14; \cite[XI.1]{westminster-confession-1647}; \cite[III.xi.6]{calvin-institutes-1960}
\item an inference step --- \cite[Trent, Session VI (1547), canon 9]{tanner-decrees-1990}
\end{itemize}
\item \emph{Covenantal nomism (Sanders, Dunn)}, weakest at \emph{disputed}:
\begin{itemize}
\item Second Temple Judaism: in by grace, staying in by works --- \cite{sanders-paul-palestinian-judaism-1977}
\item an inference step --- \cite{dunn-new-perspective-2005}; \cite{wright-what-paul-said-1997}
\end{itemize}
\item \emph{Sola fide from John (6:28–29)}, weakest at \emph{disputed}:
\begin{itemize}
\item The believing of John 6:29 is trust, which brings nothing and receives Christ --- \cite[ad loc. John 6:29]{calvin-commentary-john-1847}; John 6:28-29
\item Good works are the fruit and evidence of saving faith, not its ground --- \cite[XI.2; XVI.2; XVI.5]{westminster-confession-1647}; \cite[III.xvi.1]{calvin-institutes-1960}; Jas 2:18-22; Eph 2:10; Rom 6:22; Luke 17:10
\item Justification and sanctification are inseparable but distinct --- \cite[Q. 77]{westminster-larger-catechism-1647}; 1 Cor 6:11; 1 Cor 1:30; Rom 4:6; Rom 4:8; Ezek 36:27; Heb 9:13-14; Rom 3:24-25; Rom 6:6; Rom 6:14; \cite[XI.1]{westminster-confession-1647}; \cite[III.xi.6]{calvin-institutes-1960}
\item an inference step --- \cite[Trent, Session VI (1547), canon 9]{tanner-decrees-1990}
\end{itemize}
\item \emph{Luther on Galatians 3:11–12 (against faith formed by charity)}, weakest at \emph{disputed}:
\begin{itemize}
\item an inference step --- \cite[ad loc. Gal 3:12]{luther-commentary-galatians-1998}; \cite[ad loc. Gal 3:11]{luther-commentary-galatians-1998}
\end{itemize}
\item \emph{Obedience as sanctification (John 3:36, Matthew 7:21)}, weakest at \emph{disputed}:
\begin{itemize}
\item The obedience of John 3:36 and Matthew 7:21 is sanctification, the fruit of faith --- \cite[ad loc. John 3:36]{ridderbos-john-1997}; \cite[ad loc. John 3:36]{carson-john-1991}; \cite[ad loc. Matt 7:21]{morris-matthew-1992}; \cite[55--55]{schreiner-justification-by-works-2015}; \cite{hoekema-saved-by-grace-1994}; \cite{murray-redemption-1984}; \cite[XVI.2]{westminster-confession-1647}; \cite[XIII.1]{westminster-confession-1647}; \cite[ad loc. John 3:36]{calvin-commentary-john-1847}
\item Justification and sanctification are inseparable but distinct --- \cite[Q. 77]{westminster-larger-catechism-1647}; 1 Cor 6:11; 1 Cor 1:30; Rom 4:6; Rom 4:8; Ezek 36:27; Heb 9:13-14; Rom 3:24-25; Rom 6:6; Rom 6:14; \cite[XI.1]{westminster-confession-1647}; \cite[III.xi.6]{calvin-institutes-1960}
\item an inference step --- \cite[55--55]{schreiner-justification-by-works-2015}; \cite[ad loc. John 3:36]{calvin-commentary-john-1847}; \cite[XVI.2]{westminster-confession-1647}; \cite[XIII.1]{westminster-confession-1647}
\end{itemize}
\end{itemize}
\par\medskip\noindent\textbf{\texttt{luke\_\allowbreak{}not\_\allowbreak{}forced\_\allowbreak{}without\_\allowbreak{}trent}}\par\nobreak\noindent
\textbf{Without Trent, Luke’s case is not forced either.} A consistent set of positions that answers every attack holds Chrysostom, and cannot hold Luke’s case with it. Chrysostom reads “whoever does not obey the Son shall not see life” (John 3:36) and “the one who does the will of my Father” (Matthew 7:21) as requiring a right life beside believing. Luke’s case reads “your faith has saved you” (Luke 7:50) as saying that faith sufficed.

A reading of two other sayings of Jesus, rated \texttt{disputed}, ties with Luke’s. It does not show that Chrysostom outranks Luke, or that any text of Luke’s is contradicted. Luke’s case defeats Chrysostom in return, and Chrysostom accepts every word of Luke 7:50; what he denies is the step from it to “faith suffices”.

{\small\begin{verbatim}
theorem luke_not_forced_without_trent : ¬Framework.SkepticallyAccepted
    withoutTrent.defeats ⟨Party.dominical, lukeAnsweredByChrysostom._proof_2⟩
-- axioms: propext, Classical.choice, Quot.sound
\end{verbatim}}
\par\medskip\noindent\textbf{\texttt{paul\allowbreak{}Answered\allowbreak{}By\allowbreak{}Campbell}}\par\nobreak\noindent
Why Paul’s case is not forced without Trent: a set of positions that can be defended includes the apocalyptic reading, which defeats Paul’s case over “faith of Christ”.

\paragraph{\emph{Sola fide from Paul (Galatians 2:16)} is not accepted on every resolution.}
\begin{itemize}
\item A position holding \emph{Apocalyptic reading, against faith as the condition}, \emph{Sola fide from Peter (Acts 15:9–11)}, \emph{Variegated nomism (Gathercole, Carson et al.)}, \emph{Luke 7:50: “your faith has saved you”}, \emph{Paul’s gospel (Galatians 1), against justification as renewal}, \emph{Paul’s word (δικαιόω, by the rule of least meaning)}, \emph{Romans 4:3–8, on Paul’s word (by the rule of least meaning)} can be held, and it holds \emph{Apocalyptic reading, against faith as the condition}; \emph{Apocalyptic reading, against faith as the condition} conflicts with \emph{Sola fide from Paul (Galatians 2:16)}: \emph{Apocalyptic reading, against faith as the condition} defeats \emph{Sola fide from Paul (Galatians 2:16)}, and \emph{Sola fide from Paul (Galatians 2:16)} defeats \emph{Apocalyptic reading, against faith as the condition}.
\item None of \emph{Apocalyptic reading, against faith as the condition}, \emph{Sola fide from Peter (Acts 15:9–11)}, \emph{Variegated nomism (Gathercole, Carson et al.)}, \emph{Luke 7:50: “your faith has saved you”}, \emph{Paul’s gospel (Galatians 1), against justification as renewal}, \emph{Paul’s word (δικαιόω, by the rule of least meaning)}, \emph{Romans 4:3–8, on Paul’s word (by the rule of least meaning)} defeats another, and each attack on them is answered from among them:
\begin{itemize}
\item \emph{Sola fide from Paul (Galatians 2:16)} defeats \emph{Apocalyptic reading, against faith as the condition}, and \emph{Apocalyptic reading, against faith as the condition} defeats \emph{Sola fide from Paul (Galatians 2:16)}.
\item \emph{Sola fide from John (6:28–29)} defeats \emph{Apocalyptic reading, against faith as the condition}, and \emph{Apocalyptic reading, against faith as the condition} defeats \emph{Sola fide from John (6:28–29)}.
\item \emph{Law-observant Luke (Jervell)} defeats \emph{Sola fide from Peter (Acts 15:9–11)}, and \emph{Sola fide from Peter (Acts 15:9–11)} defeats \emph{Law-observant Luke (Jervell)}.
\item \emph{Chrysostom on John 3:36 (believing is not enough without a right life)} defeats \emph{Sola fide from Peter (Acts 15:9–11)}, and \emph{Sola fide from Peter (Acts 15:9–11)} defeats \emph{Chrysostom on John 3:36 (believing is not enough without a right life)}.
\item \emph{Covenantal nomism (Sanders, Dunn)} defeats \emph{Variegated nomism (Gathercole, Carson et al.)}, and \emph{Variegated nomism (Gathercole, Carson et al.)} defeats \emph{Covenantal nomism (Sanders, Dunn)}.
\item \emph{Chrysostom on John 3:36 (believing is not enough without a right life)} defeats \emph{Luke 7:50: “your faith has saved you”}, and \emph{Sola fide from Peter (Acts 15:9–11)} defeats \emph{Chrysostom on John 3:36 (believing is not enough without a right life)}.
\end{itemize}
\end{itemize}
\par\noindent\textbf{What this rests on.} The reasons state 11 defeats and 99 absences of defeat, each a cell of the defeat table, \texttt{sola\allowbreak{}Fide\allowbreak{}Dispute\_\allowbreak{}defeats}, computed from the two parties' premises and checked by the kernel. Whether an attack survives turns on the attacker's weakest link and on the rating of what it attacks, so a changed rating can change the verdict. The weakest links:
\begin{itemize}
\item \emph{Apocalyptic reading, against faith as the condition}, weakest at \emph{disputed}:
\begin{itemize}
\item ‘Faith of Christ’ is Christ’s own faithfulness, the means of God’s rescue (encoded as a denial, so it ranks disputed) --- \cite{campbell-deliverance-god-2009}; \cite[ad loc. Gal 2:16]{martyn-galatians-1997}; \cite{hays-faith-jesus-christ-2002}
\item ‘The righteousness of God’ (δικαιοσύνη θεοῦ) names God’s act of rescue in Christ --- \cite{campbell-deliverance-god-2009}; \cite{martyn-galatians-1997}; Rom 3:21-26
\end{itemize}
\item \emph{Sola fide from Paul (Galatians 2:16)}, weakest at \emph{disputed}:
\begin{itemize}
\item In Galatians 2:16 ‘faith of Christ’ (πίστις Χριστοῦ) means our faith in Christ --- \cite{dunn-once-more-pistis-christou-1991}; \cite{matlock-detheologizing-2000}; Gal 2:16; Rom 3:22; Phil 3:9
\item Second Temple Judaism: in by grace, staying in by works (encoded as a denial, so it ranks disputed) --- \cite{sanders-paul-palestinian-judaism-1977}
\item Good works are the fruit and evidence of saving faith, not its ground --- \cite[XI.2; XVI.2; XVI.5]{westminster-confession-1647}; \cite[III.xvi.1]{calvin-institutes-1960}; Jas 2:18-22; Eph 2:10; Rom 6:22; Luke 17:10
\item Justification and sanctification are inseparable but distinct --- \cite[Q. 77]{westminster-larger-catechism-1647}; 1 Cor 6:11; 1 Cor 1:30; Rom 4:6; Rom 4:8; Ezek 36:27; Heb 9:13-14; Rom 3:24-25; Rom 6:6; Rom 6:14; \cite[XI.1]{westminster-confession-1647}; \cite[III.xi.6]{calvin-institutes-1960}
\item an inference step --- \cite[Trent, Session VI (1547), canon 9]{tanner-decrees-1990}
\item an inference step --- \cite{gathercole-where-boasting-2002}; \cite{carson-variegated-nomism-1-2001}
\end{itemize}
\item \emph{Sola fide from Peter (Acts 15:9–11)}, weakest at \emph{disputed}:
\begin{itemize}
\item The yoke refused at Acts 15:10 is the whole law as a condition of salvation --- \cite[ad loc. Acts 15:10]{bruce-acts-1988}; Acts 15:10; Gal 5:1; Jas 2:10; \cite{das-paul-law-covenant-2001}
\item Good works are the fruit and evidence of saving faith, not its ground --- \cite[XI.2; XVI.2; XVI.5]{westminster-confession-1647}; \cite[III.xvi.1]{calvin-institutes-1960}; Jas 2:18-22; Eph 2:10; Rom 6:22; Luke 17:10
\item Justification and sanctification are inseparable but distinct --- \cite[Q. 77]{westminster-larger-catechism-1647}; 1 Cor 6:11; 1 Cor 1:30; Rom 4:6; Rom 4:8; Ezek 36:27; Heb 9:13-14; Rom 3:24-25; Rom 6:6; Rom 6:14; \cite[XI.1]{westminster-confession-1647}; \cite[III.xi.6]{calvin-institutes-1960}
\item an inference step --- \cite[Trent, Session VI (1547), canon 9]{tanner-decrees-1990}
\end{itemize}
\item \emph{Variegated nomism (Gathercole, Carson et al.)}, weakest at \emph{disputed}:
\begin{itemize}
\item Second Temple Judaism also held final vindication according to works (encoded as a denial, so it ranks disputed) --- \cite{gathercole-where-boasting-2002}; \cite{carson-variegated-nomism-1-2001}
\item an inference step --- \cite{gathercole-where-boasting-2002}; \cite{carson-variegated-nomism-1-2001}
\end{itemize}
\item \emph{Luke 7:50: “your faith has saved you”}, weakest at \emph{disputed}:
\begin{itemize}
\item Luke 7:47 — her love is the evidence of her forgiveness, not its ground --- \cite[ad loc. Luke 7:47]{marshall-luke-1978}; \cite[54--54]{padilla-narrative-criticism-2021}; \cite[Apology IV, on love and the fulfilling of the law: Luke 7:47–50]{kolb-wengert-book-of-concord-2000}; Luke 7:41-43; Luke 7:47; Luke 7:47-50
\item an inference step --- \cite[Trent, Session VI (1547), canon 9]{tanner-decrees-1990}
\end{itemize}
\item \emph{Paul’s gospel (Galatians 1), against justification as renewal}, weakest at \emph{disputed}:
\begin{itemize}
\item an inference step --- \cite[ad loc. Gal 1:6]{moo-galatians-2013}; \cite[XI.1]{westminster-confession-1647}
\item an inference step --- \cite[XI.1]{westminster-confession-1647}; \cite[III.xi.2]{calvin-institutes-1960}; \cite{morris-apostolic-preaching-1965}
\end{itemize}
\item \emph{Paul’s word (δικαιόω, by the rule of least meaning)}, weakest at \emph{well supported}:
\begin{itemize}
\item In Paul, ‘justify’ (δικαιόω) is a verdict: to declare righteous --- \cite[s.v. δικαιόω]{bauer-danker-lexicon-2000}; \cite{morris-apostolic-preaching-1965}; \cite[p. 80]{moo-romans-1996}; \cite{irons-righteousness-of-god-2015}; \cite{wright-justification-2009}; \cite[ad loc. Rom 3:24]{fitzmyer-romans-1993}
\item A word contributes the least meaning its context requires; one occurrence does not carry the whole concept --- \cite{silva-biblical-words-1994}; \cite{joos-semantic-axiom-1972}; \cite{barr-semantics-2004}
\item an inference step --- \cite{silva-biblical-words-1994}; \cite[p. 80]{moo-romans-1996}; \cite{barr-semantics-2004}
\end{itemize}
\item \emph{Romans 4:3–8, on Paul’s word (by the rule of least meaning)}, weakest at \emph{well supported}:
\begin{itemize}
\item Romans 4:4–5 contrasts wages owed with a gift reckoned to the one who believes --- Rom 4:4-5; \cite[III.xi.18]{calvin-institutes-1960}
\item ‘Count’ (λογίζομαι) in Romans 4:3–11 is the language of an account: credited to one --- \cite[s.v. λογίζομαι]{bauer-danker-lexicon-2000}; \cite[ad loc. Rom 4:3]{moo-romans-1996}
\item A word contributes the least meaning its context requires; one occurrence does not carry the whole concept --- \cite{silva-biblical-words-1994}; \cite{joos-semantic-axiom-1972}; \cite{barr-semantics-2004}
\item an inference step --- \cite[ad loc. Rom 4:6]{moo-romans-1996}; \cite[s.v. λογίζομαι]{bauer-danker-lexicon-2000}
\item an inference step --- \cite{silva-biblical-words-1994}; \cite[p. 80]{moo-romans-1996}; \cite{barr-semantics-2004}
\end{itemize}
\item \emph{Covenantal nomism (Sanders, Dunn)}, weakest at \emph{disputed}:
\begin{itemize}
\item Second Temple Judaism: in by grace, staying in by works --- \cite{sanders-paul-palestinian-judaism-1977}
\item an inference step --- \cite{dunn-new-perspective-2005}; \cite{wright-what-paul-said-1997}
\end{itemize}
\item \emph{Law-observant Luke (Jervell)}, weakest at \emph{disputed}:
\begin{itemize}
\item an inference step --- \cite{jervell-luke-people-god-1972}
\end{itemize}
\item \emph{Sola fide from John (6:28–29)}, weakest at \emph{disputed}:
\begin{itemize}
\item The believing of John 6:29 is trust, which brings nothing and receives Christ --- \cite[ad loc. John 6:29]{calvin-commentary-john-1847}; John 6:28-29
\item Good works are the fruit and evidence of saving faith, not its ground --- \cite[XI.2; XVI.2; XVI.5]{westminster-confession-1647}; \cite[III.xvi.1]{calvin-institutes-1960}; Jas 2:18-22; Eph 2:10; Rom 6:22; Luke 17:10
\item Justification and sanctification are inseparable but distinct --- \cite[Q. 77]{westminster-larger-catechism-1647}; 1 Cor 6:11; 1 Cor 1:30; Rom 4:6; Rom 4:8; Ezek 36:27; Heb 9:13-14; Rom 3:24-25; Rom 6:6; Rom 6:14; \cite[XI.1]{westminster-confession-1647}; \cite[III.xi.6]{calvin-institutes-1960}
\item an inference step --- \cite[Trent, Session VI (1547), canon 9]{tanner-decrees-1990}
\end{itemize}
\item \emph{Faith formed by charity (Aquinas on John 6:29)}, weakest at \emph{disputed}:
\begin{itemize}
\item The believing of John 6:29 is faith living through charity, the source of works --- \cite[cap. 6, lect. 3, n. 901]{aquinas-commentary-john-2012}; \cite[Trent, Session VI (1547), Decree on Justification, ch. 7]{tanner-decrees-1990}
\end{itemize}
\item \emph{Luther on Galatians 3:11–12 (against faith formed by charity)}, weakest at \emph{disputed}:
\begin{itemize}
\item an inference step --- \cite[ad loc. Gal 3:12]{luther-commentary-galatians-1998}; \cite[ad loc. Gal 3:11]{luther-commentary-galatians-1998}
\end{itemize}
\item \emph{Chrysostom on John 3:36 (believing is not enough without a right life)}, weakest at \emph{disputed}:
\begin{itemize}
\item an inference step --- \cite[Homily 31, §1]{chrysostom-homilies-john}; \cite{stanley-did-jesus-teach-2006}
\end{itemize}
\item \emph{Obedience as sanctification (John 3:36, Matthew 7:21)}, weakest at \emph{disputed}:
\begin{itemize}
\item The obedience of John 3:36 and Matthew 7:21 is sanctification, the fruit of faith --- \cite[ad loc. John 3:36]{ridderbos-john-1997}; \cite[ad loc. John 3:36]{carson-john-1991}; \cite[ad loc. Matt 7:21]{morris-matthew-1992}; \cite[55--55]{schreiner-justification-by-works-2015}; \cite{hoekema-saved-by-grace-1994}; \cite{murray-redemption-1984}; \cite[XVI.2]{westminster-confession-1647}; \cite[XIII.1]{westminster-confession-1647}; \cite[ad loc. John 3:36]{calvin-commentary-john-1847}
\item Justification and sanctification are inseparable but distinct --- \cite[Q. 77]{westminster-larger-catechism-1647}; 1 Cor 6:11; 1 Cor 1:30; Rom 4:6; Rom 4:8; Ezek 36:27; Heb 9:13-14; Rom 3:24-25; Rom 6:6; Rom 6:14; \cite[XI.1]{westminster-confession-1647}; \cite[III.xi.6]{calvin-institutes-1960}
\item an inference step --- \cite[55--55]{schreiner-justification-by-works-2015}; \cite[ad loc. John 3:36]{calvin-commentary-john-1847}; \cite[XVI.2]{westminster-confession-1647}; \cite[XIII.1]{westminster-confession-1647}
\end{itemize}
\end{itemize}
\par\medskip\noindent\textbf{\texttt{sola\_\allowbreak{}fide\_\allowbreak{}from\_\allowbreak{}paul\_\allowbreak{}not\_\allowbreak{}forced\_\allowbreak{}without\_\allowbreak{}trent}}\par\nobreak\noindent
\textbf{Without Trent, sola fide from Paul is still not forced.} A consistent set of positions that answers every attack holds Campbell’s apocalyptic reading, on which “faith of Christ” (πίστις Χριστοῦ) is Christ’s own faithfulness, and cannot hold Paul’s case with it. What stands in the way is a question about Paul’s Greek, not about Rome.

{\small\begin{verbatim}
theorem sola_fide_from_paul_not_forced_without_trent :
    ¬Framework.SkepticallyAccepted withoutTrent.defeats ⟨Party.pauline,
    paulAnsweredByCampbell._proof_2⟩
-- axioms: propext, Classical.choice, Quot.sound
\end{verbatim}}
\par\medskip\noindent\textbf{\texttt{peter\allowbreak{}Answered\allowbreak{}By\allowbreak{}Jervell}}\par\nobreak\noindent
Why Peter’s case is not forced without Trent: a set of positions that can be defended includes Jervell, who defeats it over the yoke.

\paragraph{\emph{Sola fide from Peter (Acts 15:9–11)} is not accepted on every resolution.}
\begin{itemize}
\item A position holding \emph{Variegated nomism (Gathercole, Carson et al.)}, \emph{Luke 7:50: “your faith has saved you”}, \emph{Paul’s gospel (Galatians 1), against justification as renewal}, \emph{Law-observant Luke (Jervell)}, \emph{Paul’s word (δικαιόω, by the rule of least meaning)}, \emph{Sola fide from Paul (Galatians 2:16)}, \emph{Romans 4:3–8, on Paul’s word (by the rule of least meaning)} can be held, and it holds \emph{Law-observant Luke (Jervell)}; \emph{Law-observant Luke (Jervell)} conflicts with \emph{Sola fide from Peter (Acts 15:9–11)}: \emph{Law-observant Luke (Jervell)} defeats \emph{Sola fide from Peter (Acts 15:9–11)}, and \emph{Sola fide from Peter (Acts 15:9–11)} defeats \emph{Law-observant Luke (Jervell)}.
\item None of \emph{Variegated nomism (Gathercole, Carson et al.)}, \emph{Luke 7:50: “your faith has saved you”}, \emph{Paul’s gospel (Galatians 1), against justification as renewal}, \emph{Law-observant Luke (Jervell)}, \emph{Paul’s word (δικαιόω, by the rule of least meaning)}, \emph{Sola fide from Paul (Galatians 2:16)}, \emph{Romans 4:3–8, on Paul’s word (by the rule of least meaning)} defeats another, and each attack on them is answered from among them:
\begin{itemize}
\item \emph{Covenantal nomism (Sanders, Dunn)} defeats \emph{Variegated nomism (Gathercole, Carson et al.)}, and \emph{Variegated nomism (Gathercole, Carson et al.)} defeats \emph{Covenantal nomism (Sanders, Dunn)}.
\item \emph{Chrysostom on John 3:36 (believing is not enough without a right life)} defeats \emph{Luke 7:50: “your faith has saved you”}, and \emph{Luke 7:50: “your faith has saved you”} defeats \emph{Chrysostom on John 3:36 (believing is not enough without a right life)}.
\item \emph{Sola fide from Peter (Acts 15:9–11)} defeats \emph{Law-observant Luke (Jervell)}, and \emph{Law-observant Luke (Jervell)} defeats \emph{Sola fide from Peter (Acts 15:9–11)}.
\item \emph{Apocalyptic reading, against faith as the condition} defeats \emph{Sola fide from Paul (Galatians 2:16)}, and \emph{Sola fide from Paul (Galatians 2:16)} defeats \emph{Apocalyptic reading, against faith as the condition}.
\item \emph{Covenantal nomism (Sanders, Dunn)} defeats \emph{Sola fide from Paul (Galatians 2:16)}, and \emph{Variegated nomism (Gathercole, Carson et al.)} defeats \emph{Covenantal nomism (Sanders, Dunn)}.
\item \emph{Chrysostom on John 3:36 (believing is not enough without a right life)} defeats \emph{Sola fide from Paul (Galatians 2:16)}, and \emph{Luke 7:50: “your faith has saved you”} defeats \emph{Chrysostom on John 3:36 (believing is not enough without a right life)}.
\end{itemize}
\end{itemize}
\par\noindent\textbf{What this rests on.} The reasons state 10 defeats and 99 absences of defeat, each a cell of the defeat table, \texttt{sola\allowbreak{}Fide\allowbreak{}Dispute\_\allowbreak{}defeats}, computed from the two parties' premises and checked by the kernel. Whether an attack survives turns on the attacker's weakest link and on the rating of what it attacks, so a changed rating can change the verdict. The weakest links:
\begin{itemize}
\item \emph{Law-observant Luke (Jervell)}, weakest at \emph{disputed}:
\begin{itemize}
\item an inference step --- \cite{jervell-luke-people-god-1972}
\end{itemize}
\item \emph{Sola fide from Peter (Acts 15:9–11)}, weakest at \emph{disputed}:
\begin{itemize}
\item The yoke refused at Acts 15:10 is the whole law as a condition of salvation --- \cite[ad loc. Acts 15:10]{bruce-acts-1988}; Acts 15:10; Gal 5:1; Jas 2:10; \cite{das-paul-law-covenant-2001}
\item Good works are the fruit and evidence of saving faith, not its ground --- \cite[XI.2; XVI.2; XVI.5]{westminster-confession-1647}; \cite[III.xvi.1]{calvin-institutes-1960}; Jas 2:18-22; Eph 2:10; Rom 6:22; Luke 17:10
\item Justification and sanctification are inseparable but distinct --- \cite[Q. 77]{westminster-larger-catechism-1647}; 1 Cor 6:11; 1 Cor 1:30; Rom 4:6; Rom 4:8; Ezek 36:27; Heb 9:13-14; Rom 3:24-25; Rom 6:6; Rom 6:14; \cite[XI.1]{westminster-confession-1647}; \cite[III.xi.6]{calvin-institutes-1960}
\item an inference step --- \cite[Trent, Session VI (1547), canon 9]{tanner-decrees-1990}
\end{itemize}
\item \emph{Variegated nomism (Gathercole, Carson et al.)}, weakest at \emph{disputed}:
\begin{itemize}
\item Second Temple Judaism also held final vindication according to works (encoded as a denial, so it ranks disputed) --- \cite{gathercole-where-boasting-2002}; \cite{carson-variegated-nomism-1-2001}
\item an inference step --- \cite{gathercole-where-boasting-2002}; \cite{carson-variegated-nomism-1-2001}
\end{itemize}
\item \emph{Luke 7:50: “your faith has saved you”}, weakest at \emph{disputed}:
\begin{itemize}
\item Luke 7:47 — her love is the evidence of her forgiveness, not its ground --- \cite[ad loc. Luke 7:47]{marshall-luke-1978}; \cite[54--54]{padilla-narrative-criticism-2021}; \cite[Apology IV, on love and the fulfilling of the law: Luke 7:47–50]{kolb-wengert-book-of-concord-2000}; Luke 7:41-43; Luke 7:47; Luke 7:47-50
\item an inference step --- \cite[Trent, Session VI (1547), canon 9]{tanner-decrees-1990}
\end{itemize}
\item \emph{Paul’s gospel (Galatians 1), against justification as renewal}, weakest at \emph{disputed}:
\begin{itemize}
\item an inference step --- \cite[ad loc. Gal 1:6]{moo-galatians-2013}; \cite[XI.1]{westminster-confession-1647}
\item an inference step --- \cite[XI.1]{westminster-confession-1647}; \cite[III.xi.2]{calvin-institutes-1960}; \cite{morris-apostolic-preaching-1965}
\end{itemize}
\item \emph{Paul’s word (δικαιόω, by the rule of least meaning)}, weakest at \emph{well supported}:
\begin{itemize}
\item In Paul, ‘justify’ (δικαιόω) is a verdict: to declare righteous --- \cite[s.v. δικαιόω]{bauer-danker-lexicon-2000}; \cite{morris-apostolic-preaching-1965}; \cite[p. 80]{moo-romans-1996}; \cite{irons-righteousness-of-god-2015}; \cite{wright-justification-2009}; \cite[ad loc. Rom 3:24]{fitzmyer-romans-1993}
\item A word contributes the least meaning its context requires; one occurrence does not carry the whole concept --- \cite{silva-biblical-words-1994}; \cite{joos-semantic-axiom-1972}; \cite{barr-semantics-2004}
\item an inference step --- \cite{silva-biblical-words-1994}; \cite[p. 80]{moo-romans-1996}; \cite{barr-semantics-2004}
\end{itemize}
\item \emph{Sola fide from Paul (Galatians 2:16)}, weakest at \emph{disputed}:
\begin{itemize}
\item In Galatians 2:16 ‘faith of Christ’ (πίστις Χριστοῦ) means our faith in Christ --- \cite{dunn-once-more-pistis-christou-1991}; \cite{matlock-detheologizing-2000}; Gal 2:16; Rom 3:22; Phil 3:9
\item Second Temple Judaism: in by grace, staying in by works (encoded as a denial, so it ranks disputed) --- \cite{sanders-paul-palestinian-judaism-1977}
\item Good works are the fruit and evidence of saving faith, not its ground --- \cite[XI.2; XVI.2; XVI.5]{westminster-confession-1647}; \cite[III.xvi.1]{calvin-institutes-1960}; Jas 2:18-22; Eph 2:10; Rom 6:22; Luke 17:10
\item Justification and sanctification are inseparable but distinct --- \cite[Q. 77]{westminster-larger-catechism-1647}; 1 Cor 6:11; 1 Cor 1:30; Rom 4:6; Rom 4:8; Ezek 36:27; Heb 9:13-14; Rom 3:24-25; Rom 6:6; Rom 6:14; \cite[XI.1]{westminster-confession-1647}; \cite[III.xi.6]{calvin-institutes-1960}
\item an inference step --- \cite[Trent, Session VI (1547), canon 9]{tanner-decrees-1990}
\item an inference step --- \cite{gathercole-where-boasting-2002}; \cite{carson-variegated-nomism-1-2001}
\end{itemize}
\item \emph{Romans 4:3–8, on Paul’s word (by the rule of least meaning)}, weakest at \emph{well supported}:
\begin{itemize}
\item Romans 4:4–5 contrasts wages owed with a gift reckoned to the one who believes --- Rom 4:4-5; \cite[III.xi.18]{calvin-institutes-1960}
\item ‘Count’ (λογίζομαι) in Romans 4:3–11 is the language of an account: credited to one --- \cite[s.v. λογίζομαι]{bauer-danker-lexicon-2000}; \cite[ad loc. Rom 4:3]{moo-romans-1996}
\item A word contributes the least meaning its context requires; one occurrence does not carry the whole concept --- \cite{silva-biblical-words-1994}; \cite{joos-semantic-axiom-1972}; \cite{barr-semantics-2004}
\item an inference step --- \cite[ad loc. Rom 4:6]{moo-romans-1996}; \cite[s.v. λογίζομαι]{bauer-danker-lexicon-2000}
\item an inference step --- \cite{silva-biblical-words-1994}; \cite[p. 80]{moo-romans-1996}; \cite{barr-semantics-2004}
\end{itemize}
\item \emph{Apocalyptic reading, against faith as the condition}, weakest at \emph{disputed}:
\begin{itemize}
\item ‘Faith of Christ’ is Christ’s own faithfulness, the means of God’s rescue (encoded as a denial, so it ranks disputed) --- \cite{campbell-deliverance-god-2009}; \cite[ad loc. Gal 2:16]{martyn-galatians-1997}; \cite{hays-faith-jesus-christ-2002}
\item ‘The righteousness of God’ (δικαιοσύνη θεοῦ) names God’s act of rescue in Christ --- \cite{campbell-deliverance-god-2009}; \cite{martyn-galatians-1997}; Rom 3:21-26
\end{itemize}
\item \emph{Covenantal nomism (Sanders, Dunn)}, weakest at \emph{disputed}:
\begin{itemize}
\item Second Temple Judaism: in by grace, staying in by works --- \cite{sanders-paul-palestinian-judaism-1977}
\item an inference step --- \cite{dunn-new-perspective-2005}; \cite{wright-what-paul-said-1997}
\end{itemize}
\item \emph{Sola fide from John (6:28–29)}, weakest at \emph{disputed}:
\begin{itemize}
\item The believing of John 6:29 is trust, which brings nothing and receives Christ --- \cite[ad loc. John 6:29]{calvin-commentary-john-1847}; John 6:28-29
\item Good works are the fruit and evidence of saving faith, not its ground --- \cite[XI.2; XVI.2; XVI.5]{westminster-confession-1647}; \cite[III.xvi.1]{calvin-institutes-1960}; Jas 2:18-22; Eph 2:10; Rom 6:22; Luke 17:10
\item Justification and sanctification are inseparable but distinct --- \cite[Q. 77]{westminster-larger-catechism-1647}; 1 Cor 6:11; 1 Cor 1:30; Rom 4:6; Rom 4:8; Ezek 36:27; Heb 9:13-14; Rom 3:24-25; Rom 6:6; Rom 6:14; \cite[XI.1]{westminster-confession-1647}; \cite[III.xi.6]{calvin-institutes-1960}
\item an inference step --- \cite[Trent, Session VI (1547), canon 9]{tanner-decrees-1990}
\end{itemize}
\item \emph{Faith formed by charity (Aquinas on John 6:29)}, weakest at \emph{disputed}:
\begin{itemize}
\item The believing of John 6:29 is faith living through charity, the source of works --- \cite[cap. 6, lect. 3, n. 901]{aquinas-commentary-john-2012}; \cite[Trent, Session VI (1547), Decree on Justification, ch. 7]{tanner-decrees-1990}
\end{itemize}
\item \emph{Luther on Galatians 3:11–12 (against faith formed by charity)}, weakest at \emph{disputed}:
\begin{itemize}
\item an inference step --- \cite[ad loc. Gal 3:12]{luther-commentary-galatians-1998}; \cite[ad loc. Gal 3:11]{luther-commentary-galatians-1998}
\end{itemize}
\item \emph{Chrysostom on John 3:36 (believing is not enough without a right life)}, weakest at \emph{disputed}:
\begin{itemize}
\item an inference step --- \cite[Homily 31, §1]{chrysostom-homilies-john}; \cite{stanley-did-jesus-teach-2006}
\end{itemize}
\item \emph{Obedience as sanctification (John 3:36, Matthew 7:21)}, weakest at \emph{disputed}:
\begin{itemize}
\item The obedience of John 3:36 and Matthew 7:21 is sanctification, the fruit of faith --- \cite[ad loc. John 3:36]{ridderbos-john-1997}; \cite[ad loc. John 3:36]{carson-john-1991}; \cite[ad loc. Matt 7:21]{morris-matthew-1992}; \cite[55--55]{schreiner-justification-by-works-2015}; \cite{hoekema-saved-by-grace-1994}; \cite{murray-redemption-1984}; \cite[XVI.2]{westminster-confession-1647}; \cite[XIII.1]{westminster-confession-1647}; \cite[ad loc. John 3:36]{calvin-commentary-john-1847}
\item Justification and sanctification are inseparable but distinct --- \cite[Q. 77]{westminster-larger-catechism-1647}; 1 Cor 6:11; 1 Cor 1:30; Rom 4:6; Rom 4:8; Ezek 36:27; Heb 9:13-14; Rom 3:24-25; Rom 6:6; Rom 6:14; \cite[XI.1]{westminster-confession-1647}; \cite[III.xi.6]{calvin-institutes-1960}
\item an inference step --- \cite[55--55]{schreiner-justification-by-works-2015}; \cite[ad loc. John 3:36]{calvin-commentary-john-1847}; \cite[XVI.2]{westminster-confession-1647}; \cite[XIII.1]{westminster-confession-1647}
\end{itemize}
\end{itemize}
\par\medskip\noindent\textbf{\texttt{sola\_\allowbreak{}fide\_\allowbreak{}from\_\allowbreak{}peter\_\allowbreak{}not\_\allowbreak{}forced\_\allowbreak{}without\_\allowbreak{}trent}}\par\nobreak\noindent
\textbf{Nor from Peter.} A consistent set of positions that answers every attack holds Jervell’s reading of Acts 15, on which the yoke is not the law as a condition of salvation, and cannot hold Peter’s case with it.

So without Trent it is forced that faith suffices (Luke) and that justification is not renewal (Paul’s gospel). \emph{Faith alone} still depends on questions about what the texts say: “faith of Christ”, “works of the law”, the yoke, and, for John, what believing is (\texttt{sola\_\allowbreak{}fide\_\allowbreak{}from\_\allowbreak{}john\_\allowbreak{}not\_\allowbreak{}forced\_\allowbreak{}without\_\allowbreak{}trent}).

{\small\begin{verbatim}
theorem sola_fide_from_peter_not_forced_without_trent :
    ¬Framework.SkepticallyAccepted withoutTrent.defeats ⟨Party.apostolic,
    peterAnsweredByJervell._proof_2⟩
-- axioms: propext, Classical.choice, Quot.sound
\end{verbatim}}
\par\medskip\noindent\textbf{\texttt{john\allowbreak{}Answered\allowbreak{}By\allowbreak{}Aquinas}}\par\nobreak\noindent
Why John’s case is not forced without Trent: a set of positions that can be defended includes Aquinas, who defeats it over what believing is.

\paragraph{\emph{Sola fide from John (6:28–29)} is not accepted on every resolution.}
\begin{itemize}
\item A position holding \emph{Faith formed by charity (Aquinas on John 6:29)}, \emph{Luke 7:50: “your faith has saved you”}, \emph{Paul’s gospel (Galatians 1), against justification as renewal}, \emph{Paul’s word (δικαιόω, by the rule of least meaning)}, \emph{Romans 4:3–8, on Paul’s word (by the rule of least meaning)} can be held, and it holds \emph{Faith formed by charity (Aquinas on John 6:29)}; \emph{Faith formed by charity (Aquinas on John 6:29)} conflicts with \emph{Sola fide from John (6:28–29)}: \emph{Faith formed by charity (Aquinas on John 6:29)} defeats \emph{Sola fide from John (6:28–29)}, and \emph{Sola fide from John (6:28–29)} defeats \emph{Faith formed by charity (Aquinas on John 6:29)}.
\item None of \emph{Faith formed by charity (Aquinas on John 6:29)}, \emph{Luke 7:50: “your faith has saved you”}, \emph{Paul’s gospel (Galatians 1), against justification as renewal}, \emph{Paul’s word (δικαιόω, by the rule of least meaning)}, \emph{Romans 4:3–8, on Paul’s word (by the rule of least meaning)} defeats another, and each attack on them is answered from among them:
\begin{itemize}
\item \emph{Sola fide from John (6:28–29)} defeats \emph{Faith formed by charity (Aquinas on John 6:29)}, and \emph{Faith formed by charity (Aquinas on John 6:29)} defeats \emph{Sola fide from John (6:28–29)}.
\item \emph{Luther on Galatians 3:11–12 (against faith formed by charity)} defeats \emph{Faith formed by charity (Aquinas on John 6:29)}, and \emph{Faith formed by charity (Aquinas on John 6:29)} defeats \emph{Luther on Galatians 3:11–12 (against faith formed by charity)}.
\item \emph{Chrysostom on John 3:36 (believing is not enough without a right life)} defeats \emph{Luke 7:50: “your faith has saved you”}, and \emph{Luke 7:50: “your faith has saved you”} defeats \emph{Chrysostom on John 3:36 (believing is not enough without a right life)}.
\end{itemize}
\end{itemize}
\par\noindent\textbf{What this rests on.} The reasons state 6 defeats and 72 absences of defeat, each a cell of the defeat table, \texttt{sola\allowbreak{}Fide\allowbreak{}Dispute\_\allowbreak{}defeats}, computed from the two parties' premises and checked by the kernel. Whether an attack survives turns on the attacker's weakest link and on the rating of what it attacks, so a changed rating can change the verdict. The weakest links:
\begin{itemize}
\item \emph{Faith formed by charity (Aquinas on John 6:29)}, weakest at \emph{disputed}:
\begin{itemize}
\item The believing of John 6:29 is faith living through charity, the source of works --- \cite[cap. 6, lect. 3, n. 901]{aquinas-commentary-john-2012}; \cite[Trent, Session VI (1547), Decree on Justification, ch. 7]{tanner-decrees-1990}
\end{itemize}
\item \emph{Sola fide from John (6:28–29)}, weakest at \emph{disputed}:
\begin{itemize}
\item The believing of John 6:29 is trust, which brings nothing and receives Christ --- \cite[ad loc. John 6:29]{calvin-commentary-john-1847}; John 6:28-29
\item Good works are the fruit and evidence of saving faith, not its ground --- \cite[XI.2; XVI.2; XVI.5]{westminster-confession-1647}; \cite[III.xvi.1]{calvin-institutes-1960}; Jas 2:18-22; Eph 2:10; Rom 6:22; Luke 17:10
\item Justification and sanctification are inseparable but distinct --- \cite[Q. 77]{westminster-larger-catechism-1647}; 1 Cor 6:11; 1 Cor 1:30; Rom 4:6; Rom 4:8; Ezek 36:27; Heb 9:13-14; Rom 3:24-25; Rom 6:6; Rom 6:14; \cite[XI.1]{westminster-confession-1647}; \cite[III.xi.6]{calvin-institutes-1960}
\item an inference step --- \cite[Trent, Session VI (1547), canon 9]{tanner-decrees-1990}
\end{itemize}
\item \emph{Luke 7:50: “your faith has saved you”}, weakest at \emph{disputed}:
\begin{itemize}
\item Luke 7:47 — her love is the evidence of her forgiveness, not its ground --- \cite[ad loc. Luke 7:47]{marshall-luke-1978}; \cite[54--54]{padilla-narrative-criticism-2021}; \cite[Apology IV, on love and the fulfilling of the law: Luke 7:47–50]{kolb-wengert-book-of-concord-2000}; Luke 7:41-43; Luke 7:47; Luke 7:47-50
\item an inference step --- \cite[Trent, Session VI (1547), canon 9]{tanner-decrees-1990}
\end{itemize}
\item \emph{Paul’s gospel (Galatians 1), against justification as renewal}, weakest at \emph{disputed}:
\begin{itemize}
\item an inference step --- \cite[ad loc. Gal 1:6]{moo-galatians-2013}; \cite[XI.1]{westminster-confession-1647}
\item an inference step --- \cite[XI.1]{westminster-confession-1647}; \cite[III.xi.2]{calvin-institutes-1960}; \cite{morris-apostolic-preaching-1965}
\end{itemize}
\item \emph{Paul’s word (δικαιόω, by the rule of least meaning)}, weakest at \emph{well supported}:
\begin{itemize}
\item In Paul, ‘justify’ (δικαιόω) is a verdict: to declare righteous --- \cite[s.v. δικαιόω]{bauer-danker-lexicon-2000}; \cite{morris-apostolic-preaching-1965}; \cite[p. 80]{moo-romans-1996}; \cite{irons-righteousness-of-god-2015}; \cite{wright-justification-2009}; \cite[ad loc. Rom 3:24]{fitzmyer-romans-1993}
\item A word contributes the least meaning its context requires; one occurrence does not carry the whole concept --- \cite{silva-biblical-words-1994}; \cite{joos-semantic-axiom-1972}; \cite{barr-semantics-2004}
\item an inference step --- \cite{silva-biblical-words-1994}; \cite[p. 80]{moo-romans-1996}; \cite{barr-semantics-2004}
\end{itemize}
\item \emph{Romans 4:3–8, on Paul’s word (by the rule of least meaning)}, weakest at \emph{well supported}:
\begin{itemize}
\item Romans 4:4–5 contrasts wages owed with a gift reckoned to the one who believes --- Rom 4:4-5; \cite[III.xi.18]{calvin-institutes-1960}
\item ‘Count’ (λογίζομαι) in Romans 4:3–11 is the language of an account: credited to one --- \cite[s.v. λογίζομαι]{bauer-danker-lexicon-2000}; \cite[ad loc. Rom 4:3]{moo-romans-1996}
\item A word contributes the least meaning its context requires; one occurrence does not carry the whole concept --- \cite{silva-biblical-words-1994}; \cite{joos-semantic-axiom-1972}; \cite{barr-semantics-2004}
\item an inference step --- \cite[ad loc. Rom 4:6]{moo-romans-1996}; \cite[s.v. λογίζομαι]{bauer-danker-lexicon-2000}
\item an inference step --- \cite{silva-biblical-words-1994}; \cite[p. 80]{moo-romans-1996}; \cite{barr-semantics-2004}
\end{itemize}
\item \emph{Sola fide from Paul (Galatians 2:16)}, weakest at \emph{disputed}:
\begin{itemize}
\item In Galatians 2:16 ‘faith of Christ’ (πίστις Χριστοῦ) means our faith in Christ --- \cite{dunn-once-more-pistis-christou-1991}; \cite{matlock-detheologizing-2000}; Gal 2:16; Rom 3:22; Phil 3:9
\item Second Temple Judaism: in by grace, staying in by works (encoded as a denial, so it ranks disputed) --- \cite{sanders-paul-palestinian-judaism-1977}
\item Good works are the fruit and evidence of saving faith, not its ground --- \cite[XI.2; XVI.2; XVI.5]{westminster-confession-1647}; \cite[III.xvi.1]{calvin-institutes-1960}; Jas 2:18-22; Eph 2:10; Rom 6:22; Luke 17:10
\item Justification and sanctification are inseparable but distinct --- \cite[Q. 77]{westminster-larger-catechism-1647}; 1 Cor 6:11; 1 Cor 1:30; Rom 4:6; Rom 4:8; Ezek 36:27; Heb 9:13-14; Rom 3:24-25; Rom 6:6; Rom 6:14; \cite[XI.1]{westminster-confession-1647}; \cite[III.xi.6]{calvin-institutes-1960}
\item an inference step --- \cite[Trent, Session VI (1547), canon 9]{tanner-decrees-1990}
\item an inference step --- \cite{gathercole-where-boasting-2002}; \cite{carson-variegated-nomism-1-2001}
\end{itemize}
\item \emph{Sola fide from Peter (Acts 15:9–11)}, weakest at \emph{disputed}:
\begin{itemize}
\item The yoke refused at Acts 15:10 is the whole law as a condition of salvation --- \cite[ad loc. Acts 15:10]{bruce-acts-1988}; Acts 15:10; Gal 5:1; Jas 2:10; \cite{das-paul-law-covenant-2001}
\item Good works are the fruit and evidence of saving faith, not its ground --- \cite[XI.2; XVI.2; XVI.5]{westminster-confession-1647}; \cite[III.xvi.1]{calvin-institutes-1960}; Jas 2:18-22; Eph 2:10; Rom 6:22; Luke 17:10
\item Justification and sanctification are inseparable but distinct --- \cite[Q. 77]{westminster-larger-catechism-1647}; 1 Cor 6:11; 1 Cor 1:30; Rom 4:6; Rom 4:8; Ezek 36:27; Heb 9:13-14; Rom 3:24-25; Rom 6:6; Rom 6:14; \cite[XI.1]{westminster-confession-1647}; \cite[III.xi.6]{calvin-institutes-1960}
\item an inference step --- \cite[Trent, Session VI (1547), canon 9]{tanner-decrees-1990}
\end{itemize}
\item \emph{Apocalyptic reading, against faith as the condition}, weakest at \emph{disputed}:
\begin{itemize}
\item ‘Faith of Christ’ is Christ’s own faithfulness, the means of God’s rescue (encoded as a denial, so it ranks disputed) --- \cite{campbell-deliverance-god-2009}; \cite[ad loc. Gal 2:16]{martyn-galatians-1997}; \cite{hays-faith-jesus-christ-2002}
\item ‘The righteousness of God’ (δικαιοσύνη θεοῦ) names God’s act of rescue in Christ --- \cite{campbell-deliverance-god-2009}; \cite{martyn-galatians-1997}; Rom 3:21-26
\end{itemize}
\item \emph{Covenantal nomism (Sanders, Dunn)}, weakest at \emph{disputed}:
\begin{itemize}
\item Second Temple Judaism: in by grace, staying in by works --- \cite{sanders-paul-palestinian-judaism-1977}
\item an inference step --- \cite{dunn-new-perspective-2005}; \cite{wright-what-paul-said-1997}
\end{itemize}
\item \emph{Variegated nomism (Gathercole, Carson et al.)}, weakest at \emph{disputed}:
\begin{itemize}
\item Second Temple Judaism also held final vindication according to works (encoded as a denial, so it ranks disputed) --- \cite{gathercole-where-boasting-2002}; \cite{carson-variegated-nomism-1-2001}
\item an inference step --- \cite{gathercole-where-boasting-2002}; \cite{carson-variegated-nomism-1-2001}
\end{itemize}
\item \emph{Law-observant Luke (Jervell)}, weakest at \emph{disputed}:
\begin{itemize}
\item an inference step --- \cite{jervell-luke-people-god-1972}
\end{itemize}
\item \emph{Luther on Galatians 3:11–12 (against faith formed by charity)}, weakest at \emph{disputed}:
\begin{itemize}
\item an inference step --- \cite[ad loc. Gal 3:12]{luther-commentary-galatians-1998}; \cite[ad loc. Gal 3:11]{luther-commentary-galatians-1998}
\end{itemize}
\item \emph{Chrysostom on John 3:36 (believing is not enough without a right life)}, weakest at \emph{disputed}:
\begin{itemize}
\item an inference step --- \cite[Homily 31, §1]{chrysostom-homilies-john}; \cite{stanley-did-jesus-teach-2006}
\end{itemize}
\item \emph{Obedience as sanctification (John 3:36, Matthew 7:21)}, weakest at \emph{disputed}:
\begin{itemize}
\item The obedience of John 3:36 and Matthew 7:21 is sanctification, the fruit of faith --- \cite[ad loc. John 3:36]{ridderbos-john-1997}; \cite[ad loc. John 3:36]{carson-john-1991}; \cite[ad loc. Matt 7:21]{morris-matthew-1992}; \cite[55--55]{schreiner-justification-by-works-2015}; \cite{hoekema-saved-by-grace-1994}; \cite{murray-redemption-1984}; \cite[XVI.2]{westminster-confession-1647}; \cite[XIII.1]{westminster-confession-1647}; \cite[ad loc. John 3:36]{calvin-commentary-john-1847}
\item Justification and sanctification are inseparable but distinct --- \cite[Q. 77]{westminster-larger-catechism-1647}; 1 Cor 6:11; 1 Cor 1:30; Rom 4:6; Rom 4:8; Ezek 36:27; Heb 9:13-14; Rom 3:24-25; Rom 6:6; Rom 6:14; \cite[XI.1]{westminster-confession-1647}; \cite[III.xi.6]{calvin-institutes-1960}
\item an inference step --- \cite[55--55]{schreiner-justification-by-works-2015}; \cite[ad loc. John 3:36]{calvin-commentary-john-1847}; \cite[XVI.2]{westminster-confession-1647}; \cite[XIII.1]{westminster-confession-1647}
\end{itemize}
\end{itemize}
\par\medskip\noindent\textbf{\texttt{sola\_\allowbreak{}fide\_\allowbreak{}from\_\allowbreak{}john\_\allowbreak{}not\_\allowbreak{}forced\_\allowbreak{}without\_\allowbreak{}trent}}\par\nobreak\noindent
\textbf{Nor from John.} Remove both of Trent’s parties from the hearing: its case, and its definition read as a claim about Paul’s word. Aquinas remains, and Trent describes justifying faith as he does: the believing Jesus names in John 6:29 is faith formed by charity, not bare trust. A reader can hold Aquinas’s reading consistently and answer every objection raised here. John’s case cannot be held with it.

Without the Council, John 6:29 still does not compel faith alone, because what believing is remains contested. It does not show that Aquinas’s reading wins. John’s case can be held just as consistently. Luther’s argument and Aquinas’s each defeat the other, and neither is removed.

{\small\begin{verbatim}
theorem sola_fide_from_john_not_forced_without_trent :
    ¬Framework.SkepticallyAccepted withoutTrent.defeats ⟨Party.johannine,
    johnAnsweredByAquinas._proof_2⟩
-- axioms: propext, Classical.choice, Quot.sound
\end{verbatim}}
\subsection*{Why each step is rated as it is}
\addcontentsline{toc}{subsection}{Why each step is rated as it is}
A step is rated \texttt{disputed} when someone objects to it. A reader’s next question is: \textbf{is the objection an argument, or only a disagreement?} This chapter asks that of every objection written down here to several steps of the Reformed reading. An objection counts as an argument (the code calls it \emph{credible}, or a \emph{critique}) when it does all of these:

\begin{itemize}
  \item it does not contradict itself;
  \item it denies the step’s conclusion;
  \item it accepts the step’s starting points;
  \item it argues for its denial by a step of its own, and does not merely state it;
  \item it rests its own starting points on what a stated standard allows.
\end{itemize}
There are three standards, and every result is worked out under each.

\begin{itemize}
  \item The \textbf{evidence} standard allows Scripture and evidence anyone can check.
  \item The \textbf{tradition} standard also allows the word of a council or a confession.
  \item The \textbf{Scripture} standard is the evidence standard, with one addition: an objection to a step that rests on Scripture alone must itself rest on Scripture alone. Here it gives the same answer as the evidence standard every time.
\end{itemize}
\subsubsection*{James’s word: how a rating changes}
One step reads James’s “justified” (δικαιόω) in the declaring sense the word has outside Paul (\texttt{james\allowbreak{}Lexical\allowbreak{}Line}). Two objections to it are written down.

\begin{itemize}
  \item \textbf{Trent’s reading of the word} takes the increase of justice as a premise. That is a disagreement, not an argument, under any standard.
  \item \textbf{The reading through Paul} argues for it. Paul’s “justify” means renewal, as VanLandingham argues from the lexicon and as Augustine explained it. James uses the word as Paul does. And James says Abraham was justified by works after Genesis 15:6. This is the strongest objection found, and it is an argument. But one of its starting points, that James uses the word as Paul does, rests on Trent’s word alone (chapter 10). No scholar’s statement of it has been verified.
\end{itemize}
So the step’s rating depends on whether a council’s word counts (\texttt{lexical\_\allowbreak{}step\_\allowbreak{}ratings}). Under the evidence standard its support stands at \texttt{well\allowbreak{}Supported}, because every objection written down fails, the strongest among them. Under the tradition standard it is \texttt{disputed}, by the reading through Paul.

Before the reading through Paul was written down, the same checks refused any rating above \texttt{disputed} under all three standards, because the strongest known objection was missing (\texttt{lexical\_\allowbreak{}step\_\allowbreak{}before\_\allowbreak{}the\_\allowbreak{}steelman}). Writing down the strongest objection is what permits the higher rating, and only where that objection fails.

\subsubsection*{Works as fruit, not ground}
Westminster’s step (\texttt{fruit\allowbreak{}Line}) derives “fruit, not ground” from its proof texts. Its support is \texttt{well\allowbreak{}Supported}: the Lutheran World Federation and the Catholic Church confess together that good works “follow justification and are its fruits” (§37), and that “whatever in the justified precedes or follows the free gift of faith is neither the basis of justification nor merits it” (§25, \emph{Joint Declaration}). That joint statement is about the \emph{ground} of justification. It does not settle whether works do anything further. On that the Declaration records the two sides apart. Catholics hold that good works “contribute to growth in grace, so that the righteousness that comes from God is preserved and communion with Christ is deepened” (§38). Lutherans hold “a preservation of grace and a growth in grace and faith”, and view good works “as the fruits and signs of justification and not as one’s own ‘merits’” (§39). Two objections are written down, both from Trent, and both argued.

\begin{itemize}
  \item \textbf{Canon 24, applied to James 2:24}: works cause the increase of justification. That claim rests on the council’s word alone.
  \item \textbf{Chapter 16, from the texts that promise a reward} (1 Corinthians 15:58, Hebrews 6:10, Hebrews 10:35, 2 Timothy 4:8). The texts are Scripture, and Westminster accepts them (XVI.6). But they say God \emph{rewards} the works. Chapter 16 adds that the reward is “rendered to their good works and merits”, and that “we must believe” the justified “have truly merited eternal life”. The passage from reward to merit is the council’s own claim. Dulles, defending Trent’s teaching, grants that “the fact that a reward is promised does not make it merited”. In the \emph{Joint Declaration}, Catholics say that in affirming the “meritorious” character of good works “they wish to say that, according to the biblical witness, a reward in heaven is promised to these works”, and to emphasise “the responsibility of persons for their actions” (§38).
\end{itemize}
So each of Trent’s objections rests, at one claim, on the council’s word, and the step’s rating depends on whether that word counts (\texttt{fruit\_\allowbreak{}step\_\allowbreak{}ratings}). Under the evidence standard neither objection stands, and the support remains \texttt{well\allowbreak{}Supported}. Under the tradition standard both stand, and the step is \texttt{disputed}.

This shows exactly that, and no more. Trent’s denial does not stand on Scripture and checkable evidence alone. A reader who accepts the council’s word has an argued objection. A reader who does not has none. It does not show that the reward texts teach the Reformed reading of them. It shows that Trent’s objection, as written down here, needs one claim that is cited to the council alone.

\subsubsection*{The same checks, applied the other way}
A check applied in one direction only would be unfair. So the Reformed objections to Rome’s claims are assessed by the same checks, under the same three standards.

\begin{itemize}
  \item \textbf{Westminster against chapter 16’s premise about merit} (XVI.5, XVI.6). God does not forget the work of believers: that is Hebrews 6:10, which Westminster gives as a proof and which is one of Trent’s own four texts. Having done all, believers are unworthy servants who have done only their duty (Luke 17:10). Wages are counted as due to one who works, and as a gift to one who believes (Romans 4:4–5). Their sacrifices are acceptable to God through Jesus Christ (1 Peter 2:5). So the reward is of grace, and is not given to merits. Every claim Westminster states outright is a text it cites. What is its own is the step, a reading rated \texttt{disputed} (\texttt{westminster\_\allowbreak{}answers\_\allowbreak{}from\_\allowbreak{}the\_\allowbreak{}texts}).
  \item \textbf{The Reformed reconciliation of James with Paul, against canon 24’s step and against chapter 16’s.} It accepts each step’s starting points, and argues from Westminster’s proof texts that works are the fruit of faith and not its ground.
\end{itemize}
Each is an argument under \textbf{every} standard, so each of Trent’s claims is \texttt{disputed} under every standard (\texttt{rome\_\allowbreak{}steps\_\allowbreak{}ratings}). Set beside \texttt{fruit\_\allowbreak{}step\_\allowbreak{}ratings}, that is the difference the checks find. The Reformed objections to Trent are arguments under each of the three standards. Trent’s objections to the Reformed step are arguments only under the standard that allows the council’s word. Neither side’s step is shown to be right by this. What differs is what each side’s objection rests on.

One caution. The checks test an objection’s starting points, not its step. Trent’s move from reward to merit is written as a starting point, because chapter 16 states it as something to be believed. Westminster’s move from the same texts to “not given to merits” is written as its step, and is rated \texttt{disputed}. A reader who thinks the two moves are alike should weigh this result accordingly.

\subsubsection*{Trent’s definition, and whether faith suffices}
Trent’s definition, that justification is “not remission of sins merely, but also the sanctification and renewal of the inward man” (ch. 7), can be taken two ways (\texttt{what\allowbreak{}Trents\allowbreak{}Definition\allowbreak{}Claims}), and Trent draws from it that faith does not suffice (ch. 7, canon 9). The checks are applied to each, in both directions.

From the Reformed side, each of these is an argument under every standard (\texttt{trent\_\allowbreak{}definition\_\allowbreak{}ratings}):

\begin{itemize}
  \item Paul’s gospel and Romans 4, against the definition taken as a claim about what God does;
  \item Romans 4, explaining Paul’s verb from Paul’s own text, against the definition taken as a claim about Paul’s word;
  \item Luke 18, against the step to “faith does not suffice”.
\end{itemize}
So each of Trent’s three claims is \texttt{disputed} under all three standards.

From Trent’s side (\texttt{trent\_\allowbreak{}dissents\_\allowbreak{}from\_\allowbreak{}the\_\allowbreak{}definition\_\allowbreak{}steps}):

\begin{itemize}
  \item \textbf{Its reading of Paul’s word} is an argument against the step about the word, under all three standards. Trent keeps the courtroom sense, the step’s first starting point, and argues from Augustine that the word means making righteous as well, which is what the step denies. So the step is \texttt{disputed} under every standard.
  \item \textbf{Its definition, taken as a claim about what God does}, is stated, not argued, against Paul’s gospel: Trent holds it as a premise.
  \item \textbf{Its denial that faith suffices} is an argument against Luke 18’s step under the tradition standard only. Its one starting point is the definition, which rests on the council’s word.
\end{itemize}
The question of the word is argued from both sides. The argument from Paul’s word holds the courtroom sense as a premise, cited to the lexicon, and argues from it, from Paul’s other words for renewal and from the rule of least meaning that the word means no more. Trent holds that it means renewal, cited to itself and to VanLandingham, and keeps the courtroom sense beside it. Neither premise, by what it means, denies the other, so each is an argument against the other (\texttt{lexical\allowbreak{}On\allowbreak{}Word\allowbreak{}Assessment}, \texttt{trent\allowbreak{}Word\allowbreak{}On\allowbreak{}Lexical\allowbreak{}Assessment}). Romans 4 argues against Trent’s reading from the text as well.

\subsubsection*{Claims that exclude each other by their meaning}
Some claims that are listed separately exclude each other by what they mean. “Justification includes the remission of sins, and not sanctification” denies “justification includes both” in so many words. Every check here takes those exclusions as given (\texttt{Meanings.Sola\allowbreak{}Fide.joins}), in the same way for an objection on either side, and the list is proved to contain exactly the exclusions the meanings imply. Without them, a position could count as granting a starting point that, by what it means, it denies.

\subsubsection*{The debate heard under each standard}
Which objections are arguments decides who is heard, and how each step is weighed. Under a given standard, the debate over James (\texttt{james\allowbreak{}Dispute}) is heard among every position that is not an objection, together with every objection that counts as an argument under that standard, and each rated step is weighed at the rating worked out under that standard. So who is heard and how the steps are weighed come from one record, and cannot disagree about the standard.

\begin{itemize}
  \item \textbf{What James’s “justified” means depends on the standard} (\texttt{james\_\allowbreak{}word\_\allowbreak{}turns\_\allowbreak{}on\_\allowbreak{}the\_\allowbreak{}standard}). Under the evidence standard the argument from James’s word is forced, as in the full debate: the reading through Paul is not heard. Under the tradition standard it is heard, the step about the word is weighed at \texttt{disputed}, and the question is contested. The argument from James’s word and the reading through Paul can each be defended, and neither is forced.
  \item \textbf{Whether works are fruit and not ground depends on the standard too} (\texttt{fruit\_\allowbreak{}turns\_\allowbreak{}on\_\allowbreak{}the\_\allowbreak{}standard}). Under the evidence standard, the Reformed reconciliation and Westminster’s answer on the reward are forced. Neither of Trent’s objections is heard. Under the tradition standard both are heard. The reconciliation and chapter 16’s reading can each be defended, and so can Westminster’s answer and chapter 16, which defeat each other. None of them is forced.
\end{itemize}
In both cases the rating and the verdict agree. A step rated \texttt{disputed} under a standard is contested when the debate is heard under it, and a step whose support stands is forced.

The full debate still weighs every step at the rating \emph{cited} for it. For the step about the word that is \texttt{well\allowbreak{}Supported}, which is also what is worked out under the evidence standard. Under the tradition standard it is this narrower hearing, not the full debate, that gives the consistent result. Having every debate weigh worked-out ratings is the next stage of the design (\texttt{docs/src/computed-\allowbreak{}ratings.md}).

\subsubsection*{What this does not show}
An objection that counts as an argument does not show the step wrong, and an objection that fails does not show it right. These checks say only whether the disagreement is one between arguments. Which side prevails is what the hearing works out, from the ratings each side rests on.

\subsubsection*{How this is checked}
\paragraph*{The objections, assessed}
\leavevmode\par\nobreak
\par\medskip\noindent\textbf{\texttt{canon24\allowbreak{}Assessment}}\par\nobreak\noindent
\textbf{Canon 24, against “fruit, not ground”.} It is argued, and it accepts Westminster’s proof texts. But it states that good works are a cause of the increase of justification on the council’s word alone. So it is an argument under the tradition standard and not under the evidence standard.

\paragraph{Dissent.} \emph{Trent, canon 24}, against \emph{Works as fruit (Westminster XVI.2, XVI.5, and its proofs)}. \textbf{A critique.} A critique: it can be held with the step’s grounds, and it argues its denial by a step of its own.
\begin{itemize}
\item Under the evidence standard: not credible; it rests on what the standard does not accept: Good works are a cause of the increase of justification, not merely its fruits.
\item Under the tradition standard: credible.
\item Under the Scripture standard: not credible; it rests on what the standard does not accept: Good works are a cause of the increase of justification, not merely its fruits.
\end{itemize}
\par\medskip\noindent\textbf{\texttt{reward\allowbreak{}Assessment}}\par\nobreak\noindent
\textbf{Chapter 16, against “fruit, not ground”.} It is argued, and it accepts Westminster’s proof texts. Its starting points are the texts that promise a reward, which are Scripture and which Westminster accepts, and the claim that the reward is given to the works as merits. The texts do not say that; Trent states it on its own word (“we must believe”). So it is an argument under the tradition standard and not under the evidence standard.

\paragraph{Dissent.} \emph{Trent, Session VI, ch. 16}, against \emph{Works as fruit (Westminster XVI.2, XVI.5, and its proofs)}. \textbf{A critique.} A critique: it can be held with the step’s grounds, and it argues its denial by a step of its own.
\begin{itemize}
\item Under the evidence standard: not credible; it rests on what the standard does not accept: The reward God renders to the works of the justified is rendered to them as merits.
\item Under the tradition standard: credible.
\item Under the Scripture standard: not credible; it rests on what the standard does not accept: The reward God renders to the works of the justified is rendered to them as merits.
\end{itemize}
\par\medskip\noindent\textbf{\texttt{trent\allowbreak{}Word\allowbreak{}Assessment}}\par\nobreak\noindent
\textbf{Trent’s reading of James’s word, against the step about the word.} It takes as a premise that James’s “justified” means the increase of justice, and that premise alone contradicts the step’s conclusion. A disagreement, not an argument, under any standard.

\paragraph{Dissent.} \emph{Trent, Session VI, ch. 10}, against \emph{James’s word (δικαιόω, by the rule of least meaning)}. \textbf{Not a critique.} A disagreement, not a critique: what it asserts outright already contradicts the step’s conclusion, so no step of its own argues it.
\begin{itemize}
\item Under the evidence standard: not credible; it rests on what the standard does not accept: James’s ‘justified’ (δικαιόω) means the increase of the justice received.
\item Under the tradition standard: not credible.
\item Under the Scripture standard: not credible; it rests on what the standard does not accept: James’s ‘justified’ (δικαιόω) means the increase of the justice received.
\end{itemize}
\par\medskip\noindent\textbf{\texttt{renewal\allowbreak{}Assessment}}\par\nobreak\noindent
\textbf{The reading through Paul, against the step about the word.} It is argued, from Paul’s word as renewal and James’s word as Paul’s, and it accepts the step’s starting points. It is an argument under the tradition standard. Under the evidence standard it is not, because one of its starting points, that James’s word is Paul’s, rests on Trent’s word alone.

\paragraph{Dissent.} \emph{The renewal reading}, against \emph{James’s word (δικαιόω, by the rule of least meaning)}. \textbf{A critique.} A critique: it can be held with the step’s grounds, and it argues its denial by a step of its own.
\begin{itemize}
\item Under the evidence standard: not credible; it rests on what the standard does not accept: James’s ‘justify’ (δικαιόω) has the sense Paul’s has: the justification of Romans and Galatians.
\item Under the tradition standard: credible.
\item Under the Scripture standard: not credible; it rests on what the standard does not accept: James’s ‘justify’ (δικαιόω) has the sense Paul’s has: the justification of Romans and Galatians.
\end{itemize}
\par\medskip\noindent\textbf{\texttt{trent\_\allowbreak{}on\_\allowbreak{}the\_\allowbreak{}word\_\allowbreak{}is\_\allowbreak{}asserted}}\par\nobreak\noindent
\textbf{Trent’s reading of James’s word is a disagreement, not an argument.} What it states outright includes the increase of justice, which alone contradicts the step’s conclusion. It needs no step of its own to reach its denial.

{\small\begin{verbatim}
theorem trent_on_the_word_is_asserted : trentWordAssessment.findings.kind =
    Findings.Kind.asserted ∧ ∀ (s : Standard), ¬Dissent.Credible s
    trentWordAssessment.dissent
-- axioms: propext, Classical.choice, Quot.sound
\end{verbatim}}
\paragraph*{The steps, rated}
\leavevmode\par\nobreak
\par\medskip\noindent\textbf{\texttt{fruit\allowbreak{}Step}}\par\nobreak\noindent
\textbf{“Works are fruit, not ground”, with the objections to it}: canon 24’s and chapter 16’s. Chapter 16’s is the strongest known: it argues from texts both sides accept.

\paragraph{The rating, computed.} \emph{Works as fruit (Westminster XVI.2, XVI.5, and its proofs)} is cited at \emph{well supported} (\cite[XVI.2, XVI.5]{westminster-confession-1647}; \cite[§§25, 37; see §38]{lwf-catholic-joint-declaration-2000}). Dissents encoded: \emph{Trent, canon 24}, \emph{Trent, Session VI, ch. 16}. The strongest known: \emph{Trent, Session VI, ch. 16}.
\begin{itemize}
\item Under the evidence standard: \textbf{well supported} --- every encoded dissent fails a check, the strongest known among them, so the step’s support stands.
\item Under the tradition standard: \textbf{disputed} --- a credible critique: Trent, canon 24; Trent, Session VI, ch. 16.
\item Under the Scripture standard: \textbf{well supported} --- every encoded dissent fails a check, the strongest known among them, so the step’s support stands.
\end{itemize}
\par\medskip\noindent\textbf{\texttt{fruit\_\allowbreak{}step\_\allowbreak{}ratings}}\par\nobreak\noindent
\textbf{The rating of “fruit, not ground” depends on whether a council’s word counts.} Under the evidence standard its support stands at \texttt{well\allowbreak{}Supported}. Both of Trent’s objections are argued, and each rests at one claim on the council’s word alone: canon 24’s, that works cause the increase of justification; and chapter 16’s, that the reward the texts promise is given to the works as merits. Under the tradition standard, which allows the council’s word, both count as arguments, and the step is \texttt{disputed}.

It does not show that the reward texts teach the Reformed reading. It shows that Trent’s objection, as written down here, needs one claim cited to the council alone.

{\small\begin{verbatim}
theorem fruit_step_ratings : fruitStep.ratings = [(Standard.evidence,
    Computed.stands Confidence.wellSupported), (Standard.tradition,
    Computed.disputedBy ["Trent, canon 24", "Trent, Session VI, ch. 16"]),
    (Standard.scriptural, Computed.stands Confidence.wellSupported)]
-- axioms: propext
\end{verbatim}}
\par\medskip\noindent\textbf{\texttt{lexical\allowbreak{}Step}}\par\nobreak\noindent
\textbf{The step about James’s word, with the objections to it}: Trent’s reading of the word, and the reading through Paul, which is the strongest known objection because it argues for the increase from Paul’s word.

\paragraph{The rating, computed.} \emph{James’s word (δικαιόω, by the rule of least meaning)} is cited at \emph{well supported} (\cite[ad loc. Jas 2:21]{moo-james-2000}; \cite{silva-biblical-words-1994}; \cite{barr-semantics-2004}). Dissents encoded: \emph{Trent, Session VI, ch. 10}, \emph{The renewal reading}. The strongest known: \emph{The renewal reading}.
\begin{itemize}
\item Under the evidence standard: \textbf{well supported} --- every encoded dissent fails a check, the strongest known among them, so the step’s support stands.
\item Under the tradition standard: \textbf{disputed} --- a credible critique: The renewal reading.
\item Under the Scripture standard: \textbf{well supported} --- every encoded dissent fails a check, the strongest known among them, so the step’s support stands.
\end{itemize}
\par\medskip\noindent\textbf{\texttt{lexical\_\allowbreak{}step\_\allowbreak{}ratings}}\par\nobreak\noindent
\textbf{The rating of the step about James’s word depends on whether a council’s word counts.} Under the evidence standard its support stands at \texttt{well\allowbreak{}Supported}. Every objection written down fails a check, the strongest known among them, because the reading through Paul rests its claim that James’s word is Paul’s on Trent alone. Under the tradition standard, which allows Trent’s word, that reading counts as an argument and the step is \texttt{disputed}.

{\small\begin{verbatim}
theorem lexical_step_ratings : lexicalStep.ratings = [(Standard.evidence,
    Computed.stands Confidence.wellSupported), (Standard.tradition,
    Computed.disputedBy ["The renewal reading"]), (Standard.scriptural,
    Computed.stands Confidence.wellSupported)]
-- axioms: propext
\end{verbatim}}
\par\medskip\noindent\textbf{\texttt{lexical\_\allowbreak{}step\_\allowbreak{}before\_\allowbreak{}the\_\allowbreak{}steelman}}\par\nobreak\noindent
\textbf{Before the strongest objection was written down, no rating above \texttt{disputed} was permitted.} With only Trent’s reading of the word on record, no objection counts as an argument. But the strongest known one is missing, so the rating is withheld under every standard. Writing down the reading through Paul is what changes it: to \texttt{well\allowbreak{}Supported} under the evidence standard, and to \texttt{disputed} under the tradition standard.

{\small\begin{verbatim}
theorem lexical_step_before_the_steelman : (have __src := lexicalStep;
 { step
    := __src.step, support := __src.support, register := [trentWordAssessment]
    }).ratings =
 [(Standard.evidence, Computed.withheld),
    (Standard.tradition, Computed.withheld), (Standard.scriptural,
    Computed.withheld)]
-- axioms: propext
\end{verbatim}}
\paragraph*{Rome’s steps, by the same checks}
\leavevmode\par\nobreak
\par\medskip\noindent\textbf{\texttt{westminster\allowbreak{}Assessment}}\par\nobreak\noindent
\textbf{Westminster, against Trent’s premise that the reward is given to merits} (XVI.5, XVI.6). It argues from texts Westminster gives as proofs: Hebrews 6:10, one of Trent’s own four; Luke 17:10; Romans 4:4–5; 1 Peter 2:5. It argues by a step of its own, its reading of them. Every claim it states outright is a text. So it is an argument under all three standards.

\paragraph{Dissent.} \emph{Westminster XVI.5, XVI.6}, against \emph{Trent: the reward is rendered to merits (ch. 16, canon 32)}. \textbf{A critique.} A critique: it can be held with the step’s grounds, and it argues its denial by a step of its own.
\begin{itemize}
\item Under the evidence standard: credible.
\item Under the tradition standard: credible.
\item Under the Scripture standard: credible.
\end{itemize}
\par\medskip\noindent\textbf{\texttt{harmony\allowbreak{}On\allowbreak{}Canon24\allowbreak{}Assessment}}\par\nobreak\noindent
\textbf{The Reformed reconciliation, against canon 24’s step.} It accepts James 2:24, the works by which Abraham and Rahab were justified, and canon 24’s own claim that works cause the increase. It argues from Westminster’s proof texts that works are the fruit of faith and not its ground. What it states outright rests on Scripture and on readings argued from it. So it is an argument under every standard.

\paragraph{Dissent.} \emph{The Reformed harmony (Westminster XVI.2, XVI.5)}, against \emph{James 2:24, read by canon 24}. \textbf{A critique.} A critique: it can be held with the step’s grounds, and it argues its denial by a step of its own.
\begin{itemize}
\item Under the evidence standard: credible.
\item Under the tradition standard: credible.
\item Under the Scripture standard: credible.
\end{itemize}
\par\medskip\noindent\textbf{\texttt{harmony\allowbreak{}On\allowbreak{}Reward\allowbreak{}Assessment}}\par\nobreak\noindent
\textbf{The Reformed reconciliation, against chapter 16’s step.} It accepts the reward texts, and accepts for the sake of argument even that the reward is given to merits. It argues that works are the fruit of faith and not its ground. An argument under all three standards. Westminster’s own answer goes further and denies the second starting point (\texttt{westminster\allowbreak{}Assessment}).

\paragraph{Dissent.} \emph{The Reformed harmony (Westminster XVI.2, XVI.5)}, against \emph{The reward of works (Trent, Session VI, ch. 16)}. \textbf{A critique.} A critique: it can be held with the step’s grounds, and it argues its denial by a step of its own.
\begin{itemize}
\item Under the evidence standard: credible.
\item Under the tradition standard: credible.
\item Under the Scripture standard: credible.
\end{itemize}
\par\medskip\noindent\textbf{\texttt{merit\allowbreak{}Premise}}\par\nobreak\noindent
\textbf{Trent’s premise that the reward is given to merits, with the objection to it}: Westminster’s, the strongest known.

\paragraph{The rating, computed.} \emph{Trent: the reward is rendered to merits (ch. 16, canon 32)} is cited at \emph{disputed} (\cite[Trent, Session VI (1547), ch. 16; canon 32]{tanner-decrees-1990}; \cite{dulles-two-languages-1999}). Dissents encoded: \emph{Westminster XVI.5, XVI.6}. The strongest known: \emph{Westminster XVI.5, XVI.6}.
\begin{itemize}
\item Under the evidence standard: \textbf{disputed} --- a credible critique: Westminster XVI.5, XVI.6.
\item Under the tradition standard: \textbf{disputed} --- a credible critique: Westminster XVI.5, XVI.6.
\item Under the Scripture standard: \textbf{disputed} --- a credible critique: Westminster XVI.5, XVI.6.
\end{itemize}
\par\medskip\noindent\textbf{\texttt{canon24\allowbreak{}Step}}\par\nobreak\noindent
\textbf{Canon 24’s step, with the objection to it}: the Reformed reconciliation’s.

\paragraph{The rating, computed.} \emph{James 2:24, read by canon 24} is cited at \emph{disputed} (\cite[Trent, Session VI (1547), canon 24]{tanner-decrees-1990}). Dissents encoded: \emph{The Reformed harmony (Westminster XVI.2, XVI.5)}. The strongest known: \emph{The Reformed harmony (Westminster XVI.2, XVI.5)}.
\begin{itemize}
\item Under the evidence standard: \textbf{disputed} --- a credible critique: The Reformed harmony (Westminster XVI.2, XVI.5).
\item Under the tradition standard: \textbf{disputed} --- a credible critique: The Reformed harmony (Westminster XVI.2, XVI.5).
\item Under the Scripture standard: \textbf{disputed} --- a credible critique: The Reformed harmony (Westminster XVI.2, XVI.5).
\end{itemize}
\par\medskip\noindent\textbf{\texttt{reward\allowbreak{}Step}}\par\nobreak\noindent
\textbf{Chapter 16’s step, with the objection to it}: the Reformed reconciliation’s.

\paragraph{The rating, computed.} \emph{The reward of works (Trent, Session VI, ch. 16)} is cited at \emph{disputed} (\cite[Trent, Session VI (1547), ch. 16, canon 32]{tanner-decrees-1990}; \cite[XVI.5, XVI.6]{westminster-confession-1647}). Dissents encoded: \emph{The Reformed harmony (Westminster XVI.2, XVI.5)}. The strongest known: \emph{The Reformed harmony (Westminster XVI.2, XVI.5)}.
\begin{itemize}
\item Under the evidence standard: \textbf{disputed} --- a credible critique: The Reformed harmony (Westminster XVI.2, XVI.5).
\item Under the tradition standard: \textbf{disputed} --- a credible critique: The Reformed harmony (Westminster XVI.2, XVI.5).
\item Under the Scripture standard: \textbf{disputed} --- a credible critique: The Reformed harmony (Westminster XVI.2, XVI.5).
\end{itemize}
\par\medskip\noindent\textbf{\texttt{rome\_\allowbreak{}steps\_\allowbreak{}ratings}}\par\nobreak\noindent
\textbf{The same checks, applied the other way: Rome’s steps are disputed under every standard.} Westminster’s answer to Trent’s premise about merit, and the Reformed reconciliation’s answers to canon 24’s step and chapter 16’s, each argue from texts and rest on no one’s word alone. So each counts as an argument under the evidence standard as under the tradition standard, and each of Trent’s claims is \texttt{disputed} under both.

Set beside \texttt{fruit\_\allowbreak{}step\_\allowbreak{}ratings}, this is the difference the checks find. The Reformed objections to Trent’s claims are arguments under each of the three standards. Trent’s objections to the Reformed step are arguments only under the standard that allows the council’s word.

It does not show that either side’s step is right.

{\small\begin{verbatim}
theorem rome_steps_ratings : meritPremise.ratings = [(Standard.evidence,
    Computed.disputedBy ["Westminster XVI.5, XVI.6"]), (Standard.tradition,
    Computed.disputedBy ["Westminster XVI.5, XVI.6"]), (Standard.scriptural,
    Computed.disputedBy ["Westminster XVI.5, XVI.6"])] ∧ canon24Step.ratings =
    [(Standard.evidence, Computed.disputedBy ["The Reformed harmony
    (Westminster XVI.2, XVI.5)"]), (Standard.tradition, Computed.disputedBy
    ["The Reformed harmony (Westminster XVI.2, XVI.5)"]),
    (Standard.scriptural, Computed.disputedBy ["The Reformed harmony
    (Westminster XVI.2, XVI.5)"])] ∧ rewardStep.ratings = [(Standard.evidence,
    Computed.disputedBy ["The Reformed harmony (Westminster XVI.2, XVI.5)"]),
    (Standard.tradition, Computed.disputedBy ["The Reformed harmony
    (Westminster XVI.2, XVI.5)"]), (Standard.scriptural, Computed.disputedBy
    ["The Reformed harmony (Westminster XVI.2, XVI.5)"])]
-- axioms: propext
\end{verbatim}}
\par\medskip\noindent\textbf{\texttt{westminster\_\allowbreak{}answers\_\allowbreak{}from\_\allowbreak{}the\_\allowbreak{}texts}}\par\nobreak\noindent
\textbf{Westminster’s answer rests on no one’s word alone; Trent’s claim does.} Every claim that Westminster’s reading of the reward states outright is a text it gives as a proof. The one claim chapter 16 adds to its texts, that the reward is given to merits, is the council’s own.

{\small\begin{verbatim}
theorem westminster_answers_from_the_texts :
    westminsterOnReward.restingOnAuthority = [] ∧
    trentOnReward.restingOnAuthority = [Claim.rewardRenderedToMerits]
-- axioms: propext
\end{verbatim}}
\paragraph*{Trent’s definition, and whether faith suffices}
\leavevmode\par\nobreak
\par\medskip\noindent\textbf{\texttt{trent\allowbreak{}Sufficiency\allowbreak{}Source}}\par\nobreak\noindent
Trent’s step from its definition to the denial that faith suffices. If justification is the renewal of the inward person by charity poured in, then faith, “unless hope and charity be added thereto, neither unites man perfectly with Christ” (Session VI, ch. 7), and whoever says that “nothing else is required to co-operate” is anathema (canon 9). Rated \texttt{disputed}: the Reformed read Luke 18 and Luke 7:50 as a plea that sufficed (\texttt{luke18\allowbreak{}Line}).

{\small\begin{verbatim}
def trentSufficiencySource : Source :=
  {
    primary :=
      Reference.work tannerDecrees
        (Bib.Locus.sectionRef
          "Trent, Session VI (1547), Decree on Justification, ch. 7; canon 9"),
    tradition := Tradition.romanCatholic, confidence := Confidence.disputed }
\end{verbatim}}
\par\medskip\noindent\textbf{\texttt{trent\allowbreak{}Sufficiency\allowbreak{}Line}} --- Faith does not suffice (Trent, ch. 7, canon 9)\par\nobreak\noindent
\textbf{Faith does not suffice, from Trent’s definition} (ch. 7, canon 9): the definition, and the step from it. This states Trent’s denial as these pages state it throughout, as a consequence of what Trent says justification is (\texttt{trent\allowbreak{}Definition\allowbreak{}Steps}).

\begin{derivation}
  \item[(1)] \(\pvar{TestimonyArguments.SolaFide}{33}\)\gloss{Justification is not remission of sins only, but renewal of the inward man}
  \item[(2)] \(\pvar{TestimonyArguments.SolaFide}{33} \rightarrow \lnot \pvar{TestimonyArguments.SolaFide}{26}\)
  \item[\(\vdash\)] \(\lnot \pvar{TestimonyArguments.SolaFide}{26}\)\gloss{\emph{at issue, the denial of:} Faith is sufficient: whoever believes is saved}
\end{derivation}
\par\medskip\noindent\textbf{\texttt{trent\allowbreak{}On\allowbreak{}Sufficiency}} --- Faith does not suffice (Trent, ch. 7, canon 9)\par\nobreak\noindent
Trent on whether faith suffices, as a position of its own.

\begin{derivation}
  \item[(1)] \(\pvar{TestimonyArguments.SolaFide}{33}\)\gloss{Justification is not remission of sins only, but renewal of the inward man}
  \item[(2)] \(\pvar{TestimonyArguments.SolaFide}{33} \rightarrow \lnot \pvar{TestimonyArguments.SolaFide}{26}\)
  \item[\(\vdash\)] \(\lnot \pvar{TestimonyArguments.SolaFide}{26}\)\gloss{\emph{at issue, the denial of:} Faith is sufficient: whoever believes is saved}
  \keepwithnext
\end{derivation}
\emph{No premise here rests on scripture alone.}
\par\medskip\noindent\textbf{\texttt{trent\allowbreak{}Word\allowbreak{}Line}} --- Paul’s word, as the Latin West read it (Augustine; Trent, ch. 8)\par\nobreak\noindent
\textbf{Paul’s word, as the Latin West read it.} Augustine explained “being justified” as “being made righteous”, and Trent reads the Apostle so (ch. 8). So Paul’s “justify” (δικαιόω) means the renewal of the inward person. This is Trent’s definition taken as a claim about the word (\texttt{trent\allowbreak{}On\allowbreak{}Pauls\allowbreak{}Word}), set out as a line of argument.

\begin{derivation}
  \item[(1)] \(\pvar{TestimonyArguments.SolaFide}{54}\)\gloss{Augustine explains “being justified” as “being made righteous”, and the Latin West read iustificare so}
  \item[(2)] \(\pvar{TestimonyArguments.SolaFide}{54} \rightarrow \pvar{TestimonyArguments.SolaFide}{52}\)
  \item[\(\vdash\)] \(\pvar{TestimonyArguments.SolaFide}{52}\)\gloss{\emph{at issue:} Paul’s word ‘justify’ means the renewal of the inward man, not only a verdict}
\end{derivation}
\par\medskip\noindent\textbf{\texttt{trent\allowbreak{}On\allowbreak{}The\allowbreak{}Word\allowbreak{}Of\allowbreak{}Paul}} --- Paul’s word, as the Latin West read it (Augustine; Trent, ch. 8)\par\nobreak\noindent
Trent’s reading of Paul’s word, as a position of its own.

\begin{derivation}
  \item[(1)] \(\pvar{TestimonyArguments.SolaFide}{54}\)\gloss{Augustine explains “being justified” as “being made righteous”, and the Latin West read iustificare so}
  \item[(2)] \(\pvar{TestimonyArguments.SolaFide}{54} \rightarrow \pvar{TestimonyArguments.SolaFide}{52}\)
  \item[\(\vdash\)] \(\pvar{TestimonyArguments.SolaFide}{52}\)\gloss{\emph{at issue:} Paul’s word ‘justify’ means the renewal of the inward man, not only a verdict}
  \keepwithnext
\end{derivation}
\emph{No premise here rests on scripture alone.}
\par\medskip\noindent\textbf{\texttt{gospel\allowbreak{}On\allowbreak{}Definition\allowbreak{}Assessment}}\par\nobreak\noindent
\textbf{Paul’s gospel, against Trent’s definition taken as a claim about what God does.} It argues from Galatians 1:6–9, 2:21 and 5:2–4, 1 Corinthians 15:3, Romans 8:33–34 and the courtroom sense of the verb. Every claim it states outright is a text, or a reading argued from the lexicon. An argument under all three standards.

\paragraph{Dissent.} \emph{Paul’s gospel (Galatians 1; Westminster XI.1)}, against \emph{Trent: justification is the renewal of the inward man (ch. 7)}. \textbf{A critique.} A critique: it can be held with the step’s grounds, and it argues its denial by a step of its own.
\begin{itemize}
\item Under the evidence standard: credible.
\item Under the tradition standard: credible.
\item Under the Scripture standard: credible.
\end{itemize}
\par\medskip\noindent\textbf{\texttt{romans\allowbreak{}Four\allowbreak{}On\allowbreak{}Definition\allowbreak{}Assessment}}\par\nobreak\noindent
\textbf{Romans 4, against Trent’s definition.} The righteousness God counts to the ungodly is sin not counted (4:5–8), and the verb is a courtroom word. An argument under every standard.

\paragraph{Dissent.} \emph{Romans 4:3–8}, against \emph{Trent: justification is the renewal of the inward man (ch. 7)}. \textbf{A critique.} A critique: it can be held with the step’s grounds, and it argues its denial by a step of its own.
\begin{itemize}
\item Under the evidence standard: credible.
\item Under the tradition standard: credible.
\item Under the Scripture standard: credible.
\end{itemize}
\par\medskip\noindent\textbf{\texttt{romans\allowbreak{}Four\allowbreak{}On\allowbreak{}Word\allowbreak{}Assessment}}\par\nobreak\noindent
\textbf{Romans 4, on Paul’s word, against Trent’s definition taken as a claim about the word.} Paul explains the righteousness God counts to the ungodly as sin not counted (4:5–8), and by the rule of least meaning his verb means no more than that verdict. It argues from Paul’s own text, without assuming the courtroom sense. So it is an argument under all three standards, and the strongest known against taking the definition this way.

\paragraph{Dissent.} \emph{Romans 4:3–8, on Paul’s word}, against \emph{Trent: Paul’s δικαιόω denotes renewal (ch. 8, with Augustine)}. \textbf{A critique.} A critique: it can be held with the step’s grounds, and it argues its denial by a step of its own.
\begin{itemize}
\item Under the evidence standard: credible.
\item Under the tradition standard: credible.
\item Under the Scripture standard: credible.
\end{itemize}
\par\medskip\noindent\textbf{\texttt{lexical\allowbreak{}On\allowbreak{}Word\allowbreak{}Assessment}}\par\nobreak\noindent
\textbf{The argument from Paul’s word, against Trent’s definition taken as a claim about the word.} It holds as a premise that Paul’s “justify” is a verdict, cited to the lexicon, and argues from that, from Paul’s other words for renewal and from the rule of least meaning that the word does not also mean renewal. So it reaches its denial by a step, and is an argument under all three standards. Trent’s reading stands against it in the same way (\texttt{trent\allowbreak{}Word\allowbreak{}On\allowbreak{}Lexical\allowbreak{}Assessment}): each grants the verdict sense, and each argues its own answer to whether that is all the word means.

\paragraph{Dissent.} \emph{Paul’s word (by the rule of least meaning)}, against \emph{Trent: Paul’s δικαιόω denotes renewal (ch. 8, with Augustine)}. \textbf{A critique.} A critique: it can be held with the step’s grounds, and it argues its denial by a step of its own.
\begin{itemize}
\item Under the evidence standard: credible.
\item Under the tradition standard: credible.
\item Under the Scripture standard: credible.
\end{itemize}
\par\medskip\noindent\textbf{\texttt{luke18\allowbreak{}On\allowbreak{}Sufficiency\allowbreak{}Assessment}}\par\nobreak\noindent
\textbf{Luke 18, against Trent’s step to “faith does not suffice”.} The tax collector’s plea, without works, and God’s verdict. It can be held together with Trent’s definition, the step’s starting point, and argues that the plea sufficed. An argument under all three standards.

\paragraph{Dissent.} \emph{Luke 18:9–14 (the tax collector justified)}, against \emph{Faith does not suffice (Trent, ch. 7, canon 9)}. \textbf{A critique.} A critique: it can be held with the step’s grounds, and it argues its denial by a step of its own.
\begin{itemize}
\item Under the evidence standard: credible.
\item Under the tradition standard: credible.
\item Under the Scripture standard: credible.
\end{itemize}
\par\medskip\noindent\textbf{\texttt{definition\allowbreak{}Premise}}\par\nobreak\noindent
\textbf{Trent’s definition, with the objections to it}: Paul’s gospel and Romans 4.

\paragraph{The rating, computed.} \emph{Trent: justification is the renewal of the inward man (ch. 7)} is cited at \emph{disputed} (\cite[Trent, Session VI (1547), Decree on Justification, ch. 7]{tanner-decrees-1990}; \cite[Trent, Session VI (1547), canon 11]{tanner-decrees-1990}). Dissents encoded: \emph{Paul’s gospel (Galatians 1; Westminster XI.1)}, \emph{Romans 4:3–8}. The strongest known: \emph{Paul’s gospel (Galatians 1; Westminster XI.1)}.
\begin{itemize}
\item Under the evidence standard: \textbf{disputed} --- a credible critique: Paul’s gospel (Galatians 1; Westminster XI.1); Romans 4:3–8.
\item Under the tradition standard: \textbf{disputed} --- a credible critique: Paul’s gospel (Galatians 1; Westminster XI.1); Romans 4:3–8.
\item Under the Scripture standard: \textbf{disputed} --- a credible critique: Paul’s gospel (Galatians 1; Westminster XI.1); Romans 4:3–8.
\end{itemize}
\par\medskip\noindent\textbf{\texttt{word\allowbreak{}Premise}}\par\nobreak\noindent
\textbf{Trent’s reading of Paul’s word, with the objections to it}: Romans 4 on the word, the strongest known, and the argument from Paul’s word.

\paragraph{The rating, computed.} \emph{Trent: Paul’s δικαιόω denotes renewal (ch. 8, with Augustine)} is cited at \emph{disputed} (\cite[Trent, Session VI (1547), Decree on Justification, chs. 7–8]{tanner-decrees-1990}; \cite{vanlandingham-judgment-justification-2006}). Dissents encoded: \emph{Romans 4:3–8, on Paul’s word}, \emph{Paul’s word (by the rule of least meaning)}. The strongest known: \emph{Romans 4:3–8, on Paul’s word}.
\begin{itemize}
\item Under the evidence standard: \textbf{disputed} --- a credible critique: Romans 4:3–8, on Paul’s word; Paul’s word (by the rule of least meaning).
\item Under the tradition standard: \textbf{disputed} --- a credible critique: Romans 4:3–8, on Paul’s word; Paul’s word (by the rule of least meaning).
\item Under the Scripture standard: \textbf{disputed} --- a credible critique: Romans 4:3–8, on Paul’s word; Paul’s word (by the rule of least meaning).
\end{itemize}
\par\medskip\noindent\textbf{\texttt{sufficiency\allowbreak{}Step}}\par\nobreak\noindent
\textbf{Trent’s step to “faith does not suffice”, with the objection to it}: Luke 18.

\paragraph{The rating, computed.} \emph{Faith does not suffice (Trent, ch. 7, canon 9)} is cited at \emph{disputed} (\cite[Trent, Session VI (1547), Decree on Justification, ch. 7; canon 9]{tanner-decrees-1990}). Dissents encoded: \emph{Luke 18:9–14 (the tax collector justified)}. The strongest known: \emph{Luke 18:9–14 (the tax collector justified)}.
\begin{itemize}
\item Under the evidence standard: \textbf{disputed} --- a credible critique: Luke 18:9–14 (the tax collector justified).
\item Under the tradition standard: \textbf{disputed} --- a credible critique: Luke 18:9–14 (the tax collector justified).
\item Under the Scripture standard: \textbf{disputed} --- a credible critique: Luke 18:9–14 (the tax collector justified).
\end{itemize}
\par\medskip\noindent\textbf{\texttt{trent\_\allowbreak{}definition\_\allowbreak{}ratings}}\par\nobreak\noindent
\textbf{Trent’s definition, and what follows from it, are disputed under every standard.} Taken as a claim about what God does, the definition is met by Paul’s gospel and Romans 4. Taken as a claim about Paul’s word, it is met by Romans 4’s explanation of the verb. Its step to “faith does not suffice” is met by Luke 18. Each of those argues from texts and from evidence anyone can check, so each is an argument under each of the three standards.

It does not show that the definition is false. \texttt{disputed} means there is an argued objection, not that the objection wins.

{\small\begin{verbatim}
theorem trent_definition_ratings : definitionPremise.ratings =
    [(Standard.evidence, Computed.disputedBy ["Paul's gospel (Galatians 1;
    Westminster XI.1)", "Romans 4:3–8"]), (Standard.tradition,
    Computed.disputedBy ["Paul's gospel (Galatians 1; Westminster XI.1)",
    "Romans 4:3–8"]), (Standard.scriptural, Computed.disputedBy ["Paul's
    gospel (Galatians 1; Westminster XI.1)", "Romans 4:3–8"])] ∧
    wordPremise.ratings = [(Standard.evidence, Computed.disputedBy ["Romans
    4:3–8, on Paul's word", "Paul's word (by the rule of least meaning)"]),
    (Standard.tradition, Computed.disputedBy ["Romans 4:3–8, on Paul's word",
    "Paul's word (by the rule of least meaning)"]), (Standard.scriptural,
    Computed.disputedBy ["Romans 4:3–8, on Paul's word", "Paul's word (by the
    rule of least meaning)"])] ∧ sufficiencyStep.ratings =
    [(Standard.evidence, Computed.disputedBy ["Luke 18:9–14 (the tax collector
    justified)"]), (Standard.tradition, Computed.disputedBy ["Luke 18:9–14
    (the tax collector justified)"]), (Standard.scriptural,
    Computed.disputedBy ["Luke 18:9–14 (the tax collector justified)"])]
-- axioms: propext
\end{verbatim}}
\par\medskip\noindent\textbf{\texttt{trent\allowbreak{}Word\allowbreak{}On\allowbreak{}Lexical\allowbreak{}Assessment}}\par\nobreak\noindent
\textbf{Trent’s reading of Paul’s word, against the step about the word.} It argues its conclusion from Augustine’s explanation: Paul’s “justify” means making righteous. It can be held with every starting point of the step, the courtroom sense among them, because Trent keeps the verdict: the justified “are not only reputed, but are truly called, and are, just” (Session VI, ch. 7). So it grants the step’s grounds and argues against what the step draws from them, and under every standard it is an argument against the step: against the step from the verdict to “nothing more”, which is where Trent’s own words put it.

\paragraph{Dissent.} \emph{Trent, with Augustine (ch. 8)}, against \emph{Paul’s word (δικαιόω, by the rule of least meaning)}. \textbf{A critique.} A critique: it can be held with the step’s grounds, and it argues its denial by a step of its own.
\begin{itemize}
\item Under the evidence standard: credible.
\item Under the tradition standard: credible.
\item Under the Scripture standard: credible.
\end{itemize}
\par\medskip\noindent\textbf{\texttt{trent\allowbreak{}Grace\allowbreak{}On\allowbreak{}Gospel\allowbreak{}Assessment}}\par\nobreak\noindent
\textbf{Trent’s definition taken as a claim about what God does, against the step in Paul’s gospel.} It holds the definition as a premise, and that premise alone contradicts the step’s conclusion. A disagreement, not an argument, under each of the three standards.

\paragraph{Dissent.} \emph{Trent, read as what God does (ch. 7; Joint Declaration §22)}, against \emph{Paul’s gospel: a verdict on a finished work is no renewal}. \textbf{Not a critique.} A disagreement, not a critique: what it asserts outright already contradicts the step’s conclusion, so no step of its own argues it.
\begin{itemize}
\item Under the evidence standard: not credible; it rests on what the standard does not accept: Justification is not remission of sins only, but renewal of the inward man; The inward renewal of the justified grows as they do good works in grace.
\item Under the tradition standard: not credible.
\item Under the Scripture standard: not credible; it rests on what the standard does not accept: Justification is not remission of sins only, but renewal of the inward man; The inward renewal of the justified grows as they do good works in grace.
\end{itemize}
\par\medskip\noindent\textbf{\texttt{trent\allowbreak{}Sufficiency\allowbreak{}On\allowbreak{}Luke18\allowbreak{}Assessment}}\par\nobreak\noindent
\textbf{Trent’s denial that faith suffices, against Luke 18’s step.} It is argued, from its definition, and it accepts the parable. Its one starting point is the definition, which rests on the council’s word. So it is an argument under the tradition standard only.

\paragraph{Dissent.} \emph{Trent, ch. 7 and canon 9}, against \emph{Luke 18:9–14 (the tax collector justified)}. \textbf{A critique.} A critique: it can be held with the step’s grounds, and it argues its denial by a step of its own.
\begin{itemize}
\item Under the evidence standard: not credible; it rests on what the standard does not accept: Justification is not remission of sins only, but renewal of the inward man.
\item Under the tradition standard: credible.
\item Under the Scripture standard: not credible; it rests on what the standard does not accept: Justification is not remission of sins only, but renewal of the inward man.
\end{itemize}
\par\medskip\noindent\textbf{\texttt{lexical\allowbreak{}Word\allowbreak{}Step}}\par\nobreak\noindent
\textbf{The step about Paul’s word, with the objection to it}: Trent’s reading of the word, the strongest known.

\paragraph{The rating, computed.} \emph{Paul’s word (δικαιόω, by the rule of least meaning)} is cited at \emph{well supported} (\cite{silva-biblical-words-1994}; \cite[p. 80]{moo-romans-1996}; \cite{barr-semantics-2004}). Dissents encoded: \emph{Trent, with Augustine (ch. 8)}. The strongest known: \emph{Trent, with Augustine (ch. 8)}.
\begin{itemize}
\item Under the evidence standard: \textbf{disputed} --- a credible critique: Trent, with Augustine (ch. 8).
\item Under the tradition standard: \textbf{disputed} --- a credible critique: Trent, with Augustine (ch. 8).
\item Under the Scripture standard: \textbf{disputed} --- a credible critique: Trent, with Augustine (ch. 8).
\end{itemize}
\par\medskip\noindent\textbf{\texttt{trent\_\allowbreak{}dissents\_\allowbreak{}from\_\allowbreak{}the\_\allowbreak{}definition\_\allowbreak{}steps}}\par\nobreak\noindent
\textbf{The same checks, from Trent’s side.} Trent objects three times to the Reformed steps over its definition.

\begin{itemize}
  \item Its reading of Paul’s word is an argument against the step about the word, under every standard. Trent keeps the courtroom sense, since the justified “are not only reputed, but are truly called, and are, just” (Session VI, ch. 7), and argues from Augustine that the word means making righteous too. So it grants the step’s starting points and denies what the step draws from them.
  \item Its definition, taken as a claim about what God does, is stated and not argued against Paul’s gospel.
  \item Its denial that faith suffices is an argument against Luke 18’s step under the tradition standard only, because its one starting point is the definition.
\end{itemize}
So the step about Paul’s word is disputed under all three standards, by Trent’s reading of the word.

What it does not claim: that Trent’s reading is right. \texttt{disputed} means there is an argued objection, not that the objection wins. Nor does it move the dispute’s verdicts on Paul’s word. The dispute weighs this step at its cited rating, which this check does not change.

{\small\begin{verbatim}
theorem trent_dissents_from_the_definition_steps :
    trentWordOnLexicalAssessment.findings.kind = Findings.Kind.argued ∧
    trentGraceOnGospelAssessment.findings.kind = Findings.Kind.asserted ∧
    trentSufficiencyOnLuke18Assessment.profile = [(Standard.evidence, false),
    (Standard.tradition, true), (Standard.scriptural, false)] ∧
    lexicalWordStep.ratings = [(Standard.evidence, Computed.disputedBy
    ["Trent, with Augustine (ch. 8)"]), (Standard.tradition,
    Computed.disputedBy ["Trent, with Augustine (ch. 8)"]),
    (Standard.scriptural, Computed.disputedBy ["Trent, with Augustine (ch.
    8)"])]
-- axioms: propext
\end{verbatim}}
\paragraph*{The debate heard under each standard}
\leavevmode\par\nobreak
\par\medskip\noindent\textbf{\texttt{james\allowbreak{}Register}}\par\nobreak\noindent
\textbf{The record for the debate over James.} Each position that objects to a step of the Reformed reading is entered as an objection and assessed, together with the two steps they are rated against. The other positions hold those steps, or are readings of texts that the texts themselves answer.

{\small\begin{verbatim}
def jamesRegister : Register jamesDispute :=
  {
    assessmentOf := fun x =>
      match x with
      | JamesParty.trentOnTheWord => some trentWordAssessment
      | JamesParty.trentOnWorks => some canon24Assessment
      | JamesParty.trentOnReward => some rewardAssessment
      | JamesParty.renewalOnTheWord => some renewalAssessment
      | JamesParty.westminsterOnReward =>
        some westminsterAssessment
      | x => none,
    ownPosition := jamesRegister._proof_1,
    steps :=
      [fruitStep, lexicalStep, canon24Step,
        rewardStep] }
\end{verbatim}}
\par\medskip\noindent\textbf{\texttt{james\allowbreak{}Register\_\allowbreak{}defeats}}\par\nobreak\noindent
\textbf{Who defeats whom under each standard.} Under the evidence standard, as in the full debate. Under the tradition standard there is one difference: Trent’s reading of the word and the reading through Paul each defeat the argument from James’s word, now that it is weighed as \texttt{disputed}. Every answer is worked out from the positions’ premises with the ratings for that standard, and checked by machine.

{\small\begin{verbatim}
theorem jamesRegister_defeats : ∀ (s : Standard) (i j : JamesParty),
    (jamesRegister.rerated s).defeats i j ↔ (i, j) ∈
    jamesRegister.defeatsTable s
-- axioms: propext, Classical.choice, Quot.sound
\end{verbatim}}
\par\medskip\noindent\textbf{\texttt{james\allowbreak{}Under\allowbreak{}Evidence}}\par\nobreak\noindent
The debate heard under the evidence standard.

{\small\begin{verbatim}
def jamesUnderEvidence : Dispute Claim { i // jamesRegister.admits
    Standard.evidence i } :=
  jamesRegister.hearing Standard.evidence
\end{verbatim}}
\par\medskip\noindent\textbf{\texttt{james\allowbreak{}Under\allowbreak{}Tradition}}\par\nobreak\noindent
The debate heard under the tradition standard.

{\small\begin{verbatim}
def jamesUnderTradition : Dispute Claim { i // jamesRegister.admits
    Standard.tradition i } :=
  jamesRegister.hearing Standard.tradition
\end{verbatim}}
\par\medskip\noindent\textbf{\texttt{james\allowbreak{}Under\allowbreak{}Evidence\_\allowbreak{}defeats}}\par\nobreak\noindent
\textbf{Who defeats whom under the evidence standard}: as in the full debate, among the positions heard.

{\small\begin{verbatim}
theorem jamesUnderEvidence_defeats : ∀ (i j : { i // jamesRegister.admits
    Standard.evidence i }), jamesUnderEvidence.defeats i j ↔ (↑i, ↑j) ∈
    jamesRegister.defeatsTable Standard.evidence
-- axioms: propext, Classical.choice, Quot.sound
\end{verbatim}}
\par\medskip\noindent\textbf{\texttt{james\allowbreak{}Under\allowbreak{}Tradition\_\allowbreak{}defeats}}\par\nobreak\noindent
\textbf{Who defeats whom under the tradition standard}: the argument from James’s word, weighed as \texttt{disputed}, and the reading through Paul defeat each other.

{\small\begin{verbatim}
theorem jamesUnderTradition_defeats : ∀ (i j : { i // jamesRegister.admits
    Standard.tradition i }), jamesUnderTradition.defeats i j ↔ (↑i, ↑j) ∈
    jamesRegister.defeatsTable Standard.tradition
-- axioms: propext, Classical.choice, Quot.sound
\end{verbatim}}
\par\medskip\noindent\textbf{\texttt{under\allowbreak{}Evidence\allowbreak{}Settled}}\par\nobreak\noindent
How the debate is settled under the evidence standard. The argument from James’s word, the argument about James’s faith, the Reformed reconciliation and the Greek of both verses are accepted first, and nothing heard answers them. Trent’s Latin readings fall to the Greek. Every objection of Trent’s to a rated step, namely canon 24, chapter 16, the reading through Paul, and Trent’s reading of the word, either rests at some claim on the council’s word or merely states its denial, and so is not heard.

\paragraph{What the dispute forces is exactly \emph{James’s word (δικαιόω, by the rule of least meaning)}, \emph{James’s faith (the faith James calls dead)}, \emph{Harmonisation of James 2:24}, \emph{Revelation 22:11, in the Greek}, \emph{Sirach 18:22, in the Greek}, \emph{The reward of works, of grace (Westminster XVI.5, XVI.6)}.}
\begin{itemize}
\item Stage 1: \emph{James’s word (δικαιόω, by the rule of least meaning)}, \emph{James’s faith (the faith James calls dead)}, \emph{Harmonisation of James 2:24}, \emph{Revelation 22:11, in the Greek}, \emph{Sirach 18:22, in the Greek}, \emph{The reward of works, of grace (Westminster XVI.5, XVI.6)}.
\begin{itemize}
\item \emph{James’s word (δικαιόω, by the rule of least meaning)} is defeated by nothing.
\item \emph{James’s faith (the faith James calls dead)} is defeated by nothing.
\item \emph{Harmonisation of James 2:24} is defeated by nothing.
\item \emph{Revelation 22:11, in the Greek} is defeated by nothing.
\item \emph{Sirach 18:22, in the Greek} is defeated by nothing.
\item \emph{The reward of works, of grace (Westminster XVI.5, XVI.6)} is defeated by nothing.
\end{itemize}
\item Nothing else is forced:
\begin{itemize}
\item \emph{Revelation 22:11, in the Greek} defeats \emph{Revelation 22:11, as Trent quotes it}, and nothing forced defeats \emph{Revelation 22:11, in the Greek}.
\item \emph{Sirach 18:22, in the Greek} defeats \emph{Sirach 18:22, as Trent quotes it}, and nothing forced defeats \emph{Sirach 18:22, in the Greek}.
\end{itemize}
\end{itemize}
\par\noindent\textbf{What this rests on.} The reasons state 2 defeats and 48 absences of defeat, each a cell of the defeat table, \texttt{james\allowbreak{}Under\allowbreak{}Evidence\_\allowbreak{}defeats}, computed from the two parties' premises and checked by the kernel. Whether an attack survives turns on the attacker's weakest link and on the rating of what it attacks, so a changed rating can change the verdict. The weakest links:
\begin{itemize}
\item \emph{James’s word (δικαιόω, by the rule of least meaning)}, weakest at \emph{well supported}:
\begin{itemize}
\item Outside Paul, ‘justify’ (δικαιόω) declares or vindicates; it does not mean making righteous --- \cite[s.v. δικαιόω]{bauer-danker-lexicon-2000}; Deut 25:1; Matt 11:19; Luke 7:35; \cite[ad loc. Matt 11:19]{crossway-esv-2001}
\item A word contributes the least meaning its context requires; one occurrence does not carry the whole concept --- \cite{silva-biblical-words-1994}; \cite{joos-semantic-axiom-1972}; \cite{barr-semantics-2004}
\item an inference step --- \cite[ad loc. Jas 2:21]{moo-james-2000}; \cite{silva-biblical-words-1994}; \cite{barr-semantics-2004}
\end{itemize}
\item \emph{James’s faith (the faith James calls dead)}, weakest at \emph{well supported}:
\begin{itemize}
\item The faith alone that James 2:24 denies justifies is the dead faith of 2:14–26 --- Jas 2:14-26; \cite[Trent, Session VI (1547), ch. 7]{tanner-decrees-1990}; \cite[ad loc. Jas 2:24]{moo-james-2000}
\item an inference step --- \cite[ad loc. Jas 2:24]{moo-james-2000}; \cite[XI.2]{westminster-confession-1647}
\end{itemize}
\item \emph{Harmonisation of James 2:24}, weakest at \emph{disputed}:
\begin{itemize}
\item an inference step --- \cite[ad loc. Jas 2:24]{moo-james-2000}; \cite[III.xvii.11]{calvin-institutes-1960}
\end{itemize}
\item \emph{Revelation 22:11, as Trent quotes it}, weakest at \emph{disputed}:
\begin{itemize}
\item Revelation 22:11 says: he that is just, let him be justified still --- \cite[Trent, Session VI (1547), ch. 10]{tanner-decrees-1990}
\item an inference step --- \cite[Trent, Session VI (1547), ch. 10]{tanner-decrees-1990}
\end{itemize}
\item \emph{Sirach 18:22, as Trent quotes it}, weakest at \emph{disputed}:
\begin{itemize}
\item Sirach 18:22 says: be not afraid to be justified even to death --- \cite[Trent, Session VI (1547), ch. 10]{tanner-decrees-1990}
\item an inference step --- \cite[Trent, Session VI (1547), ch. 10]{tanner-decrees-1990}
\end{itemize}
\item \emph{Revelation 22:11, in the Greek}, weakest at \emph{well supported}:
\begin{itemize}
\item an inference step --- \cite[§16]{pius-divino-afflante-spiritu-1943}; \cite{holmes-sblgnt-2010}; \cite{usccb-nabre-2011}
\end{itemize}
\item \emph{Sirach 18:22, in the Greek}, weakest at \emph{well supported}:
\begin{itemize}
\item The Greek of Sirach 18:22 is about paying a vow: do not wait until death to be justified --- \cite[ad loc. Sir 18:22]{swete-old-testament-in-greek}; \cite[Sirach 18:22]{usccb-nabre-2011}
\item an inference step --- \cite[§16]{pius-divino-afflante-spiritu-1943}; \cite{holmes-sblgnt-2010}; \cite{usccb-nabre-2011}
\end{itemize}
\item \emph{The reward of works, of grace (Westminster XVI.5, XVI.6)}, weakest at \emph{disputed}:
\begin{itemize}
\item an inference step --- \cite[XVI.5, XVI.6]{westminster-confession-1647}
\end{itemize}
\end{itemize}
\par\medskip\noindent\textbf{\texttt{under\allowbreak{}Tradition\allowbreak{}Settled}}\par\nobreak\noindent
How the debate is settled under the tradition standard. The argument about James’s faith and the Greek of both verses are accepted first. The argument from James’s word and the reading through Paul defeat each other. So do the Reformed reconciliation and each of Trent’s readings of what works do. Only Trent’s reading of the word is not heard.

\paragraph{What the dispute forces is exactly \emph{James’s faith (the faith James calls dead)}, \emph{Revelation 22:11, in the Greek}, \emph{Sirach 18:22, in the Greek}.}
\begin{itemize}
\item Stage 1: \emph{James’s faith (the faith James calls dead)}, \emph{Revelation 22:11, in the Greek}, \emph{Sirach 18:22, in the Greek}.
\begin{itemize}
\item \emph{James’s faith (the faith James calls dead)} is defeated by nothing.
\item \emph{Revelation 22:11, in the Greek} is defeated by nothing.
\item \emph{Sirach 18:22, in the Greek} is defeated by nothing.
\end{itemize}
\item Nothing else is forced:
\begin{itemize}
\item \emph{James’s word, read through Paul’s (the renewal reading)} defeats \emph{James’s word (δικαιόω, by the rule of least meaning)}, and nothing forced defeats \emph{James’s word, read through Paul’s (the renewal reading)}.
\item \emph{James 2:24, read by canon 24} defeats \emph{Harmonisation of James 2:24}, and nothing forced defeats \emph{James 2:24, read by canon 24}.
\item \emph{Harmonisation of James 2:24} defeats \emph{James 2:24, read by canon 24}, and nothing forced defeats \emph{Harmonisation of James 2:24}.
\item \emph{Revelation 22:11, in the Greek} defeats \emph{Revelation 22:11, as Trent quotes it}, and nothing forced defeats \emph{Revelation 22:11, in the Greek}.
\item \emph{Sirach 18:22, in the Greek} defeats \emph{Sirach 18:22, as Trent quotes it}, and nothing forced defeats \emph{Sirach 18:22, in the Greek}.
\item \emph{Harmonisation of James 2:24} defeats \emph{The reward of works (Trent, Session VI, ch. 16)}, and nothing forced defeats \emph{Harmonisation of James 2:24}.
\item \emph{James’s word (δικαιόω, by the rule of least meaning)} defeats \emph{James’s word, read through Paul’s (the renewal reading)}, and nothing forced defeats \emph{James’s word (δικαιόω, by the rule of least meaning)}.
\item \emph{The reward of works (Trent, Session VI, ch. 16)} defeats \emph{The reward of works, of grace (Westminster XVI.5, XVI.6)}, and nothing forced defeats \emph{The reward of works (Trent, Session VI, ch. 16)}.
\end{itemize}
\end{itemize}
\par\noindent\textbf{What this rests on.} The reasons state 8 defeats and 48 absences of defeat, each a cell of the defeat table, \texttt{james\allowbreak{}Under\allowbreak{}Tradition\_\allowbreak{}defeats}, computed from the two parties' premises and checked by the kernel. Whether an attack survives turns on the attacker's weakest link and on the rating of what it attacks, so a changed rating can change the verdict. The weakest links:
\begin{itemize}
\item \emph{James’s word (δικαιόω, by the rule of least meaning)}, weakest at \emph{disputed}:
\begin{itemize}
\item an inference step --- \cite[ad loc. Jas 2:21]{moo-james-2000}; \cite{silva-biblical-words-1994}; \cite{barr-semantics-2004}
\end{itemize}
\item \emph{James’s faith (the faith James calls dead)}, weakest at \emph{well supported}:
\begin{itemize}
\item The faith alone that James 2:24 denies justifies is the dead faith of 2:14–26 --- Jas 2:14-26; \cite[Trent, Session VI (1547), ch. 7]{tanner-decrees-1990}; \cite[ad loc. Jas 2:24]{moo-james-2000}
\item an inference step --- \cite[ad loc. Jas 2:24]{moo-james-2000}; \cite[XI.2]{westminster-confession-1647}
\end{itemize}
\item \emph{Harmonisation of James 2:24}, weakest at \emph{disputed}:
\begin{itemize}
\item an inference step --- \cite[XVI.2, XVI.5]{westminster-confession-1647}; \cite[§§25, 37; see §38]{lwf-catholic-joint-declaration-2000}
\item an inference step --- \cite[ad loc. Jas 2:24]{moo-james-2000}; \cite[III.xvii.11]{calvin-institutes-1960}
\end{itemize}
\item \emph{James 2:24, read by canon 24}, weakest at \emph{disputed}:
\begin{itemize}
\item Good works are a cause of the increase of justification, not merely its fruits --- \cite[Trent, Session VI (1547), canon 24]{tanner-decrees-1990}
\item an inference step --- \cite[Trent, Session VI (1547), canon 24]{tanner-decrees-1990}
\end{itemize}
\item \emph{Revelation 22:11, as Trent quotes it}, weakest at \emph{disputed}:
\begin{itemize}
\item Revelation 22:11 says: he that is just, let him be justified still --- \cite[Trent, Session VI (1547), ch. 10]{tanner-decrees-1990}
\item an inference step --- \cite[Trent, Session VI (1547), ch. 10]{tanner-decrees-1990}
\end{itemize}
\item \emph{Sirach 18:22, as Trent quotes it}, weakest at \emph{disputed}:
\begin{itemize}
\item Sirach 18:22 says: be not afraid to be justified even to death --- \cite[Trent, Session VI (1547), ch. 10]{tanner-decrees-1990}
\item an inference step --- \cite[Trent, Session VI (1547), ch. 10]{tanner-decrees-1990}
\end{itemize}
\item \emph{Revelation 22:11, in the Greek}, weakest at \emph{well supported}:
\begin{itemize}
\item an inference step --- \cite[§16]{pius-divino-afflante-spiritu-1943}; \cite{holmes-sblgnt-2010}; \cite{usccb-nabre-2011}
\end{itemize}
\item \emph{Sirach 18:22, in the Greek}, weakest at \emph{well supported}:
\begin{itemize}
\item The Greek of Sirach 18:22 is about paying a vow: do not wait until death to be justified --- \cite[ad loc. Sir 18:22]{swete-old-testament-in-greek}; \cite[Sirach 18:22]{usccb-nabre-2011}
\item an inference step --- \cite[§16]{pius-divino-afflante-spiritu-1943}; \cite{holmes-sblgnt-2010}; \cite{usccb-nabre-2011}
\end{itemize}
\item \emph{The reward of works (Trent, Session VI, ch. 16)}, weakest at \emph{disputed}:
\begin{itemize}
\item The reward God renders to the works of the justified is rendered to them as merits --- \cite[Trent, Session VI (1547), ch. 16; canon 32]{tanner-decrees-1990}; \cite{dulles-two-languages-1999}
\item an inference step --- \cite[Trent, Session VI (1547), ch. 16, canon 32]{tanner-decrees-1990}; \cite[XVI.5, XVI.6]{westminster-confession-1647}
\end{itemize}
\item \emph{James’s word, read through Paul’s (the renewal reading)}, weakest at \emph{disputed}:
\begin{itemize}
\item Paul’s word ‘justify’ means the renewal of the inward man, not only a verdict --- \cite[Trent, Session VI (1547), Decree on Justification, chs. 7–8]{tanner-decrees-1990}; \cite{vanlandingham-judgment-justification-2006}
\item James’s ‘justify’ (δικαιόω) has the sense Paul’s has: the justification of Romans and Galatians --- \cite[Trent, Session VI (1547), ch. 10]{tanner-decrees-1990}
\item an inference step --- \cite[Trent, Session VI (1547), ch. 10]{tanner-decrees-1990}; \cite{vanlandingham-judgment-justification-2006}
\end{itemize}
\item \emph{The reward of works, of grace (Westminster XVI.5, XVI.6)}, weakest at \emph{disputed}:
\begin{itemize}
\item an inference step --- \cite[XVI.5, XVI.6]{westminster-confession-1647}
\end{itemize}
\end{itemize}
\par\medskip\noindent\textbf{\texttt{renewal\_\allowbreak{}holds\_\allowbreak{}increase}}\par\nobreak\noindent
The reading through Paul holds that James’s “justified” means the increase of justice.

{\small\begin{verbatim}
theorem renewal_holds_increase : Entails (jamesDispute.node
    JamesParty.renewalOnTheWord).premises (p
    Claim.jamesJustifyDenotesIncrease)
-- axioms: propext, Classical.choice, Quot.sound
\end{verbatim}}
\par\medskip\noindent\textbf{\texttt{lexical\_\allowbreak{}denies\_\allowbreak{}increase}}\par\nobreak\noindent
The argument from James’s word denies it.

{\small\begin{verbatim}
theorem lexical_denies_increase : Entails (jamesDispute.node
    JamesParty.lexical).premises (∼p Claim.jamesJustifyDenotesIncrease)
-- axioms: propext, Classical.choice, Quot.sound
\end{verbatim}}
\par\medskip\noindent\textbf{\texttt{renewal\allowbreak{}Heard}}\par\nobreak\noindent
Why the reading through Paul can be defended under the tradition standard: it answers the argument from James’s word, now that the two are rated alike.

\paragraph{\emph{James’s word, read through Paul’s (the renewal reading)} is accepted on some resolution.}
\begin{itemize}
\item A position holding \emph{James’s faith (the faith James calls dead)}, \emph{James’s word, read through Paul’s (the renewal reading)}, \emph{Revelation 22:11, in the Greek}, \emph{Sirach 18:22, in the Greek} can be held.
\item None of \emph{James’s faith (the faith James calls dead)}, \emph{James’s word, read through Paul’s (the renewal reading)}, \emph{Revelation 22:11, in the Greek}, \emph{Sirach 18:22, in the Greek} defeats another, and each attack on them is answered from among them:
\begin{itemize}
\item \emph{James’s word (δικαιόω, by the rule of least meaning)} defeats \emph{James’s word, read through Paul’s (the renewal reading)}, and \emph{James’s word, read through Paul’s (the renewal reading)} defeats \emph{James’s word (δικαιόω, by the rule of least meaning)}.
\item \emph{Harmonisation of James 2:24} defeats \emph{James’s word, read through Paul’s (the renewal reading)}, and \emph{James’s word, read through Paul’s (the renewal reading)} defeats \emph{Harmonisation of James 2:24}.
\end{itemize}
\end{itemize}
\par\noindent\textbf{What this rests on.} The reasons state 4 defeats and 42 absences of defeat, each a cell of the defeat table, \texttt{james\allowbreak{}Under\allowbreak{}Tradition\_\allowbreak{}defeats}, computed from the two parties' premises and checked by the kernel. Whether an attack survives turns on the attacker's weakest link and on the rating of what it attacks, so a changed rating can change the verdict. The weakest links:
\begin{itemize}
\item \emph{James’s faith (the faith James calls dead)}, weakest at \emph{well supported}:
\begin{itemize}
\item The faith alone that James 2:24 denies justifies is the dead faith of 2:14–26 --- Jas 2:14-26; \cite[Trent, Session VI (1547), ch. 7]{tanner-decrees-1990}; \cite[ad loc. Jas 2:24]{moo-james-2000}
\item an inference step --- \cite[ad loc. Jas 2:24]{moo-james-2000}; \cite[XI.2]{westminster-confession-1647}
\end{itemize}
\item \emph{James’s word, read through Paul’s (the renewal reading)}, weakest at \emph{disputed}:
\begin{itemize}
\item Paul’s word ‘justify’ means the renewal of the inward man, not only a verdict --- \cite[Trent, Session VI (1547), Decree on Justification, chs. 7–8]{tanner-decrees-1990}; \cite{vanlandingham-judgment-justification-2006}
\item James’s ‘justify’ (δικαιόω) has the sense Paul’s has: the justification of Romans and Galatians --- \cite[Trent, Session VI (1547), ch. 10]{tanner-decrees-1990}
\item an inference step --- \cite[Trent, Session VI (1547), ch. 10]{tanner-decrees-1990}; \cite{vanlandingham-judgment-justification-2006}
\end{itemize}
\item \emph{Revelation 22:11, in the Greek}, weakest at \emph{well supported}:
\begin{itemize}
\item an inference step --- \cite[§16]{pius-divino-afflante-spiritu-1943}; \cite{holmes-sblgnt-2010}; \cite{usccb-nabre-2011}
\end{itemize}
\item \emph{Sirach 18:22, in the Greek}, weakest at \emph{well supported}:
\begin{itemize}
\item The Greek of Sirach 18:22 is about paying a vow: do not wait until death to be justified --- \cite[ad loc. Sir 18:22]{swete-old-testament-in-greek}; \cite[Sirach 18:22]{usccb-nabre-2011}
\item an inference step --- \cite[§16]{pius-divino-afflante-spiritu-1943}; \cite{holmes-sblgnt-2010}; \cite{usccb-nabre-2011}
\end{itemize}
\item \emph{James’s word (δικαιόω, by the rule of least meaning)}, weakest at \emph{disputed}:
\begin{itemize}
\item an inference step --- \cite[ad loc. Jas 2:21]{moo-james-2000}; \cite{silva-biblical-words-1994}; \cite{barr-semantics-2004}
\end{itemize}
\item \emph{Harmonisation of James 2:24}, weakest at \emph{disputed}:
\begin{itemize}
\item an inference step --- \cite[XVI.2, XVI.5]{westminster-confession-1647}; \cite[§§25, 37; see §38]{lwf-catholic-joint-declaration-2000}
\item an inference step --- \cite[ad loc. Jas 2:24]{moo-james-2000}; \cite[III.xvii.11]{calvin-institutes-1960}
\end{itemize}
\item \emph{James 2:24, read by canon 24}, weakest at \emph{disputed}:
\begin{itemize}
\item Good works are a cause of the increase of justification, not merely its fruits --- \cite[Trent, Session VI (1547), canon 24]{tanner-decrees-1990}
\item an inference step --- \cite[Trent, Session VI (1547), canon 24]{tanner-decrees-1990}
\end{itemize}
\item \emph{Revelation 22:11, as Trent quotes it}, weakest at \emph{disputed}:
\begin{itemize}
\item Revelation 22:11 says: he that is just, let him be justified still --- \cite[Trent, Session VI (1547), ch. 10]{tanner-decrees-1990}
\item an inference step --- \cite[Trent, Session VI (1547), ch. 10]{tanner-decrees-1990}
\end{itemize}
\item \emph{Sirach 18:22, as Trent quotes it}, weakest at \emph{disputed}:
\begin{itemize}
\item Sirach 18:22 says: be not afraid to be justified even to death --- \cite[Trent, Session VI (1547), ch. 10]{tanner-decrees-1990}
\item an inference step --- \cite[Trent, Session VI (1547), ch. 10]{tanner-decrees-1990}
\end{itemize}
\item \emph{The reward of works (Trent, Session VI, ch. 16)}, weakest at \emph{disputed}:
\begin{itemize}
\item The reward God renders to the works of the justified is rendered to them as merits --- \cite[Trent, Session VI (1547), ch. 16; canon 32]{tanner-decrees-1990}; \cite{dulles-two-languages-1999}
\item an inference step --- \cite[Trent, Session VI (1547), ch. 16, canon 32]{tanner-decrees-1990}; \cite[XVI.5, XVI.6]{westminster-confession-1647}
\end{itemize}
\item \emph{The reward of works, of grace (Westminster XVI.5, XVI.6)}, weakest at \emph{disputed}:
\begin{itemize}
\item an inference step --- \cite[XVI.5, XVI.6]{westminster-confession-1647}
\end{itemize}
\end{itemize}
\par\medskip\noindent\textbf{\texttt{lexical\allowbreak{}Heard}}\par\nobreak\noindent
Why the argument from James’s word can be defended there: it answers the reading through Paul.

\paragraph{\emph{James’s word (δικαιόω, by the rule of least meaning)} is accepted on some resolution.}
\begin{itemize}
\item A position holding \emph{James’s word (δικαιόω, by the rule of least meaning)}, \emph{James’s faith (the faith James calls dead)}, \emph{Revelation 22:11, in the Greek}, \emph{Sirach 18:22, in the Greek} can be held.
\item None of \emph{James’s word (δικαιόω, by the rule of least meaning)}, \emph{James’s faith (the faith James calls dead)}, \emph{Revelation 22:11, in the Greek}, \emph{Sirach 18:22, in the Greek} defeats another, and each attack on them is answered from among them:
\begin{itemize}
\item \emph{James’s word, read through Paul’s (the renewal reading)} defeats \emph{James’s word (δικαιόω, by the rule of least meaning)}, and \emph{James’s word (δικαιόω, by the rule of least meaning)} defeats \emph{James’s word, read through Paul’s (the renewal reading)}.
\end{itemize}
\end{itemize}
\par\noindent\textbf{What this rests on.} The reasons state 2 defeats and 43 absences of defeat, each a cell of the defeat table, \texttt{james\allowbreak{}Under\allowbreak{}Tradition\_\allowbreak{}defeats}, computed from the two parties' premises and checked by the kernel. Whether an attack survives turns on the attacker's weakest link and on the rating of what it attacks, so a changed rating can change the verdict. The weakest links:
\begin{itemize}
\item \emph{James’s word (δικαιόω, by the rule of least meaning)}, weakest at \emph{disputed}:
\begin{itemize}
\item an inference step --- \cite[ad loc. Jas 2:21]{moo-james-2000}; \cite{silva-biblical-words-1994}; \cite{barr-semantics-2004}
\end{itemize}
\item \emph{James’s faith (the faith James calls dead)}, weakest at \emph{well supported}:
\begin{itemize}
\item The faith alone that James 2:24 denies justifies is the dead faith of 2:14–26 --- Jas 2:14-26; \cite[Trent, Session VI (1547), ch. 7]{tanner-decrees-1990}; \cite[ad loc. Jas 2:24]{moo-james-2000}
\item an inference step --- \cite[ad loc. Jas 2:24]{moo-james-2000}; \cite[XI.2]{westminster-confession-1647}
\end{itemize}
\item \emph{Revelation 22:11, in the Greek}, weakest at \emph{well supported}:
\begin{itemize}
\item an inference step --- \cite[§16]{pius-divino-afflante-spiritu-1943}; \cite{holmes-sblgnt-2010}; \cite{usccb-nabre-2011}
\end{itemize}
\item \emph{Sirach 18:22, in the Greek}, weakest at \emph{well supported}:
\begin{itemize}
\item The Greek of Sirach 18:22 is about paying a vow: do not wait until death to be justified --- \cite[ad loc. Sir 18:22]{swete-old-testament-in-greek}; \cite[Sirach 18:22]{usccb-nabre-2011}
\item an inference step --- \cite[§16]{pius-divino-afflante-spiritu-1943}; \cite{holmes-sblgnt-2010}; \cite{usccb-nabre-2011}
\end{itemize}
\item \emph{Harmonisation of James 2:24}, weakest at \emph{disputed}:
\begin{itemize}
\item an inference step --- \cite[XVI.2, XVI.5]{westminster-confession-1647}; \cite[§§25, 37; see §38]{lwf-catholic-joint-declaration-2000}
\item an inference step --- \cite[ad loc. Jas 2:24]{moo-james-2000}; \cite[III.xvii.11]{calvin-institutes-1960}
\end{itemize}
\item \emph{James 2:24, read by canon 24}, weakest at \emph{disputed}:
\begin{itemize}
\item Good works are a cause of the increase of justification, not merely its fruits --- \cite[Trent, Session VI (1547), canon 24]{tanner-decrees-1990}
\item an inference step --- \cite[Trent, Session VI (1547), canon 24]{tanner-decrees-1990}
\end{itemize}
\item \emph{Revelation 22:11, as Trent quotes it}, weakest at \emph{disputed}:
\begin{itemize}
\item Revelation 22:11 says: he that is just, let him be justified still --- \cite[Trent, Session VI (1547), ch. 10]{tanner-decrees-1990}
\item an inference step --- \cite[Trent, Session VI (1547), ch. 10]{tanner-decrees-1990}
\end{itemize}
\item \emph{Sirach 18:22, as Trent quotes it}, weakest at \emph{disputed}:
\begin{itemize}
\item Sirach 18:22 says: be not afraid to be justified even to death --- \cite[Trent, Session VI (1547), ch. 10]{tanner-decrees-1990}
\item an inference step --- \cite[Trent, Session VI (1547), ch. 10]{tanner-decrees-1990}
\end{itemize}
\item \emph{The reward of works (Trent, Session VI, ch. 16)}, weakest at \emph{disputed}:
\begin{itemize}
\item The reward God renders to the works of the justified is rendered to them as merits --- \cite[Trent, Session VI (1547), ch. 16; canon 32]{tanner-decrees-1990}; \cite{dulles-two-languages-1999}
\item an inference step --- \cite[Trent, Session VI (1547), ch. 16, canon 32]{tanner-decrees-1990}; \cite[XVI.5, XVI.6]{westminster-confession-1647}
\end{itemize}
\item \emph{James’s word, read through Paul’s (the renewal reading)}, weakest at \emph{disputed}:
\begin{itemize}
\item Paul’s word ‘justify’ means the renewal of the inward man, not only a verdict --- \cite[Trent, Session VI (1547), Decree on Justification, chs. 7–8]{tanner-decrees-1990}; \cite{vanlandingham-judgment-justification-2006}
\item James’s ‘justify’ (δικαιόω) has the sense Paul’s has: the justification of Romans and Galatians --- \cite[Trent, Session VI (1547), ch. 10]{tanner-decrees-1990}
\item an inference step --- \cite[Trent, Session VI (1547), ch. 10]{tanner-decrees-1990}; \cite{vanlandingham-judgment-justification-2006}
\end{itemize}
\item \emph{The reward of works, of grace (Westminster XVI.5, XVI.6)}, weakest at \emph{disputed}:
\begin{itemize}
\item an inference step --- \cite[XVI.5, XVI.6]{westminster-confession-1647}
\end{itemize}
\end{itemize}
\par\medskip\noindent\textbf{\texttt{james\_\allowbreak{}word\_\allowbreak{}turns\_\allowbreak{}on\_\allowbreak{}the\_\allowbreak{}standard}}\par\nobreak\noindent
\textbf{What James’s “justified” means depends on the standard, and the verdict and the rating agree.}

Under the evidence standard the argument from James’s word is forced, with the argument about James’s faith and the Greek of both verses. The reading through Paul is not heard, because it rests on Trent’s word that James’s “justified” (δικαιόω) is Paul’s.

Under the tradition standard the reading through Paul is heard, the step about the word is weighed at the rating worked out for it there, \texttt{disputed}, and whether James’s word means the increase of justice is contested. Each side can be defended, and neither is forced.

{\small\begin{verbatim}
theorem james_word_turns_on_the_standard : Framework.grounded
    jamesUnderEvidence.defeats = Solver.toSet (Witness.listSub
    (jamesRegister.admits Standard.evidence) [JamesParty.lexical,
    JamesParty.faith, JamesParty.harmony, JamesParty.revelationText,
    JamesParty.sirachText, JamesParty.westminsterOnReward]) ∧ Nonempty
    (Contested jamesUnderTradition (p Claim.jamesJustifyDenotesIncrease))
-- axioms: propext, Classical.choice, Quot.sound
\end{verbatim}}
\par\medskip\noindent\textbf{\texttt{under\_\allowbreak{}tradition\_\allowbreak{}forces}}\par\nobreak\noindent
\textbf{Under the tradition standard, what is forced is the argument about James’s faith and the Greek of the two verses}, and no more.

{\small\begin{verbatim}
theorem under_tradition_forces : Framework.grounded
    jamesUnderTradition.defeats = Solver.toSet (Witness.listSub
    (jamesRegister.admits Standard.tradition) [JamesParty.faith,
    JamesParty.revelationText, JamesParty.sirachText])
-- axioms: propext, Classical.choice, Quot.sound
\end{verbatim}}
\par\medskip\noindent\textbf{\texttt{trent\allowbreak{}On\allowbreak{}Reward\_\allowbreak{}denies\_\allowbreak{}fruit}}\par\nobreak\noindent
Trent’s reading of the reward texts denies that works are fruit and not ground.

{\small\begin{verbatim}
theorem trentOnReward_denies_fruit : Entails (jamesDispute.node
    JamesParty.trentOnReward).premises (∼p Claim.worksAreFruitNotGround)
-- axioms: propext, Classical.choice, Quot.sound
\end{verbatim}}
\par\medskip\noindent\textbf{\texttt{harmony\allowbreak{}Heard}}\par\nobreak\noindent
Why the Reformed reconciliation can be defended under the tradition standard: it answers canon 24 and chapter 16.

\paragraph{\emph{Harmonisation of James 2:24} is accepted on some resolution.}
\begin{itemize}
\item A position holding \emph{James’s word (δικαιόω, by the rule of least meaning)}, \emph{James’s faith (the faith James calls dead)}, \emph{Harmonisation of James 2:24}, \emph{Revelation 22:11, in the Greek}, \emph{Sirach 18:22, in the Greek} can be held.
\item None of \emph{James’s word (δικαιόω, by the rule of least meaning)}, \emph{James’s faith (the faith James calls dead)}, \emph{Harmonisation of James 2:24}, \emph{Revelation 22:11, in the Greek}, \emph{Sirach 18:22, in the Greek} defeats another, and each attack on them is answered from among them:
\begin{itemize}
\item \emph{James’s word, read through Paul’s (the renewal reading)} defeats \emph{James’s word (δικαιόω, by the rule of least meaning)}, and \emph{James’s word (δικαιόω, by the rule of least meaning)} defeats \emph{James’s word, read through Paul’s (the renewal reading)}.
\item \emph{James 2:24, read by canon 24} defeats \emph{Harmonisation of James 2:24}, and \emph{Harmonisation of James 2:24} defeats \emph{James 2:24, read by canon 24}.
\item \emph{The reward of works (Trent, Session VI, ch. 16)} defeats \emph{Harmonisation of James 2:24}, and \emph{Harmonisation of James 2:24} defeats \emph{The reward of works (Trent, Session VI, ch. 16)}.
\item \emph{James’s word, read through Paul’s (the renewal reading)} defeats \emph{Harmonisation of James 2:24}, and \emph{James’s word (δικαιόω, by the rule of least meaning)} defeats \emph{James’s word, read through Paul’s (the renewal reading)}.
\end{itemize}
\end{itemize}
\par\noindent\textbf{What this rests on.} The reasons state 7 defeats and 51 absences of defeat, each a cell of the defeat table, \texttt{james\allowbreak{}Under\allowbreak{}Tradition\_\allowbreak{}defeats}, computed from the two parties' premises and checked by the kernel. Whether an attack survives turns on the attacker's weakest link and on the rating of what it attacks, so a changed rating can change the verdict. The weakest links:
\begin{itemize}
\item \emph{James’s word (δικαιόω, by the rule of least meaning)}, weakest at \emph{disputed}:
\begin{itemize}
\item an inference step --- \cite[ad loc. Jas 2:21]{moo-james-2000}; \cite{silva-biblical-words-1994}; \cite{barr-semantics-2004}
\end{itemize}
\item \emph{James’s faith (the faith James calls dead)}, weakest at \emph{well supported}:
\begin{itemize}
\item The faith alone that James 2:24 denies justifies is the dead faith of 2:14–26 --- Jas 2:14-26; \cite[Trent, Session VI (1547), ch. 7]{tanner-decrees-1990}; \cite[ad loc. Jas 2:24]{moo-james-2000}
\item an inference step --- \cite[ad loc. Jas 2:24]{moo-james-2000}; \cite[XI.2]{westminster-confession-1647}
\end{itemize}
\item \emph{Harmonisation of James 2:24}, weakest at \emph{disputed}:
\begin{itemize}
\item an inference step --- \cite[XVI.2, XVI.5]{westminster-confession-1647}; \cite[§§25, 37; see §38]{lwf-catholic-joint-declaration-2000}
\item an inference step --- \cite[ad loc. Jas 2:24]{moo-james-2000}; \cite[III.xvii.11]{calvin-institutes-1960}
\end{itemize}
\item \emph{Revelation 22:11, in the Greek}, weakest at \emph{well supported}:
\begin{itemize}
\item an inference step --- \cite[§16]{pius-divino-afflante-spiritu-1943}; \cite{holmes-sblgnt-2010}; \cite{usccb-nabre-2011}
\end{itemize}
\item \emph{Sirach 18:22, in the Greek}, weakest at \emph{well supported}:
\begin{itemize}
\item The Greek of Sirach 18:22 is about paying a vow: do not wait until death to be justified --- \cite[ad loc. Sir 18:22]{swete-old-testament-in-greek}; \cite[Sirach 18:22]{usccb-nabre-2011}
\item an inference step --- \cite[§16]{pius-divino-afflante-spiritu-1943}; \cite{holmes-sblgnt-2010}; \cite{usccb-nabre-2011}
\end{itemize}
\item \emph{James 2:24, read by canon 24}, weakest at \emph{disputed}:
\begin{itemize}
\item Good works are a cause of the increase of justification, not merely its fruits --- \cite[Trent, Session VI (1547), canon 24]{tanner-decrees-1990}
\item an inference step --- \cite[Trent, Session VI (1547), canon 24]{tanner-decrees-1990}
\end{itemize}
\item \emph{Revelation 22:11, as Trent quotes it}, weakest at \emph{disputed}:
\begin{itemize}
\item Revelation 22:11 says: he that is just, let him be justified still --- \cite[Trent, Session VI (1547), ch. 10]{tanner-decrees-1990}
\item an inference step --- \cite[Trent, Session VI (1547), ch. 10]{tanner-decrees-1990}
\end{itemize}
\item \emph{Sirach 18:22, as Trent quotes it}, weakest at \emph{disputed}:
\begin{itemize}
\item Sirach 18:22 says: be not afraid to be justified even to death --- \cite[Trent, Session VI (1547), ch. 10]{tanner-decrees-1990}
\item an inference step --- \cite[Trent, Session VI (1547), ch. 10]{tanner-decrees-1990}
\end{itemize}
\item \emph{The reward of works (Trent, Session VI, ch. 16)}, weakest at \emph{disputed}:
\begin{itemize}
\item The reward God renders to the works of the justified is rendered to them as merits --- \cite[Trent, Session VI (1547), ch. 16; canon 32]{tanner-decrees-1990}; \cite{dulles-two-languages-1999}
\item an inference step --- \cite[Trent, Session VI (1547), ch. 16, canon 32]{tanner-decrees-1990}; \cite[XVI.5, XVI.6]{westminster-confession-1647}
\end{itemize}
\item \emph{James’s word, read through Paul’s (the renewal reading)}, weakest at \emph{disputed}:
\begin{itemize}
\item Paul’s word ‘justify’ means the renewal of the inward man, not only a verdict --- \cite[Trent, Session VI (1547), Decree on Justification, chs. 7–8]{tanner-decrees-1990}; \cite{vanlandingham-judgment-justification-2006}
\item James’s ‘justify’ (δικαιόω) has the sense Paul’s has: the justification of Romans and Galatians --- \cite[Trent, Session VI (1547), ch. 10]{tanner-decrees-1990}
\item an inference step --- \cite[Trent, Session VI (1547), ch. 10]{tanner-decrees-1990}; \cite{vanlandingham-judgment-justification-2006}
\end{itemize}
\item \emph{The reward of works, of grace (Westminster XVI.5, XVI.6)}, weakest at \emph{disputed}:
\begin{itemize}
\item an inference step --- \cite[XVI.5, XVI.6]{westminster-confession-1647}
\end{itemize}
\end{itemize}
\par\medskip\noindent\textbf{\texttt{reward\allowbreak{}Heard}}\par\nobreak\noindent
Why chapter 16’s reading of the reward texts can be defended there: it answers the Reformed reconciliation.

\paragraph{\emph{The reward of works (Trent, Session VI, ch. 16)} is accepted on some resolution.}
\begin{itemize}
\item A position holding \emph{James’s word (δικαιόω, by the rule of least meaning)}, \emph{James’s faith (the faith James calls dead)}, \emph{James 2:24, read by canon 24}, \emph{The reward of works (Trent, Session VI, ch. 16)}, \emph{Revelation 22:11, in the Greek}, \emph{Sirach 18:22, in the Greek} can be held.
\item None of \emph{James’s word (δικαιόω, by the rule of least meaning)}, \emph{James’s faith (the faith James calls dead)}, \emph{James 2:24, read by canon 24}, \emph{The reward of works (Trent, Session VI, ch. 16)}, \emph{Revelation 22:11, in the Greek}, \emph{Sirach 18:22, in the Greek} defeats another, and each attack on them is answered from among them:
\begin{itemize}
\item \emph{James’s word, read through Paul’s (the renewal reading)} defeats \emph{James’s word (δικαιόω, by the rule of least meaning)}, and \emph{James’s word (δικαιόω, by the rule of least meaning)} defeats \emph{James’s word, read through Paul’s (the renewal reading)}.
\item \emph{Harmonisation of James 2:24} defeats \emph{James 2:24, read by canon 24}, and \emph{James 2:24, read by canon 24} defeats \emph{Harmonisation of James 2:24}.
\item \emph{Harmonisation of James 2:24} defeats \emph{The reward of works (Trent, Session VI, ch. 16)}, and \emph{James 2:24, read by canon 24} defeats \emph{Harmonisation of James 2:24}.
\item \emph{The reward of works, of grace (Westminster XVI.5, XVI.6)} defeats \emph{The reward of works (Trent, Session VI, ch. 16)}, and \emph{The reward of works (Trent, Session VI, ch. 16)} defeats \emph{The reward of works, of grace (Westminster XVI.5, XVI.6)}.
\end{itemize}
\end{itemize}
\par\noindent\textbf{What this rests on.} The reasons state 7 defeats and 62 absences of defeat, each a cell of the defeat table, \texttt{james\allowbreak{}Under\allowbreak{}Tradition\_\allowbreak{}defeats}, computed from the two parties' premises and checked by the kernel. Whether an attack survives turns on the attacker's weakest link and on the rating of what it attacks, so a changed rating can change the verdict. The weakest links:
\begin{itemize}
\item \emph{James’s word (δικαιόω, by the rule of least meaning)}, weakest at \emph{disputed}:
\begin{itemize}
\item an inference step --- \cite[ad loc. Jas 2:21]{moo-james-2000}; \cite{silva-biblical-words-1994}; \cite{barr-semantics-2004}
\end{itemize}
\item \emph{James’s faith (the faith James calls dead)}, weakest at \emph{well supported}:
\begin{itemize}
\item The faith alone that James 2:24 denies justifies is the dead faith of 2:14–26 --- Jas 2:14-26; \cite[Trent, Session VI (1547), ch. 7]{tanner-decrees-1990}; \cite[ad loc. Jas 2:24]{moo-james-2000}
\item an inference step --- \cite[ad loc. Jas 2:24]{moo-james-2000}; \cite[XI.2]{westminster-confession-1647}
\end{itemize}
\item \emph{James 2:24, read by canon 24}, weakest at \emph{disputed}:
\begin{itemize}
\item Good works are a cause of the increase of justification, not merely its fruits --- \cite[Trent, Session VI (1547), canon 24]{tanner-decrees-1990}
\item an inference step --- \cite[Trent, Session VI (1547), canon 24]{tanner-decrees-1990}
\end{itemize}
\item \emph{The reward of works (Trent, Session VI, ch. 16)}, weakest at \emph{disputed}:
\begin{itemize}
\item The reward God renders to the works of the justified is rendered to them as merits --- \cite[Trent, Session VI (1547), ch. 16; canon 32]{tanner-decrees-1990}; \cite{dulles-two-languages-1999}
\item an inference step --- \cite[Trent, Session VI (1547), ch. 16, canon 32]{tanner-decrees-1990}; \cite[XVI.5, XVI.6]{westminster-confession-1647}
\end{itemize}
\item \emph{Revelation 22:11, in the Greek}, weakest at \emph{well supported}:
\begin{itemize}
\item an inference step --- \cite[§16]{pius-divino-afflante-spiritu-1943}; \cite{holmes-sblgnt-2010}; \cite{usccb-nabre-2011}
\end{itemize}
\item \emph{Sirach 18:22, in the Greek}, weakest at \emph{well supported}:
\begin{itemize}
\item The Greek of Sirach 18:22 is about paying a vow: do not wait until death to be justified --- \cite[ad loc. Sir 18:22]{swete-old-testament-in-greek}; \cite[Sirach 18:22]{usccb-nabre-2011}
\item an inference step --- \cite[§16]{pius-divino-afflante-spiritu-1943}; \cite{holmes-sblgnt-2010}; \cite{usccb-nabre-2011}
\end{itemize}
\item \emph{Harmonisation of James 2:24}, weakest at \emph{disputed}:
\begin{itemize}
\item an inference step --- \cite[XVI.2, XVI.5]{westminster-confession-1647}; \cite[§§25, 37; see §38]{lwf-catholic-joint-declaration-2000}
\item an inference step --- \cite[ad loc. Jas 2:24]{moo-james-2000}; \cite[III.xvii.11]{calvin-institutes-1960}
\end{itemize}
\item \emph{Revelation 22:11, as Trent quotes it}, weakest at \emph{disputed}:
\begin{itemize}
\item Revelation 22:11 says: he that is just, let him be justified still --- \cite[Trent, Session VI (1547), ch. 10]{tanner-decrees-1990}
\item an inference step --- \cite[Trent, Session VI (1547), ch. 10]{tanner-decrees-1990}
\end{itemize}
\item \emph{Sirach 18:22, as Trent quotes it}, weakest at \emph{disputed}:
\begin{itemize}
\item Sirach 18:22 says: be not afraid to be justified even to death --- \cite[Trent, Session VI (1547), ch. 10]{tanner-decrees-1990}
\item an inference step --- \cite[Trent, Session VI (1547), ch. 10]{tanner-decrees-1990}
\end{itemize}
\item \emph{James’s word, read through Paul’s (the renewal reading)}, weakest at \emph{disputed}:
\begin{itemize}
\item Paul’s word ‘justify’ means the renewal of the inward man, not only a verdict --- \cite[Trent, Session VI (1547), Decree on Justification, chs. 7–8]{tanner-decrees-1990}; \cite{vanlandingham-judgment-justification-2006}
\item James’s ‘justify’ (δικαιόω) has the sense Paul’s has: the justification of Romans and Galatians --- \cite[Trent, Session VI (1547), ch. 10]{tanner-decrees-1990}
\item an inference step --- \cite[Trent, Session VI (1547), ch. 10]{tanner-decrees-1990}; \cite{vanlandingham-judgment-justification-2006}
\end{itemize}
\item \emph{The reward of works, of grace (Westminster XVI.5, XVI.6)}, weakest at \emph{disputed}:
\begin{itemize}
\item an inference step --- \cite[XVI.5, XVI.6]{westminster-confession-1647}
\end{itemize}
\end{itemize}
\par\medskip\noindent\textbf{\texttt{fruit\_\allowbreak{}turns\_\allowbreak{}on\_\allowbreak{}the\_\allowbreak{}standard}}\par\nobreak\noindent
\textbf{Whether works are fruit and not ground depends on whether a council’s word counts, and the verdict and the rating agree.}

Under the evidence standard the Reformed reconciliation is forced. It derives “fruit, not ground” from Westminster’s proof texts. Neither of Trent’s objections is heard: each rests at one claim on the council’s word.

Under the tradition standard both are heard, the step is weighed at \texttt{disputed}, and the question is contested. The reconciliation and chapter 16’s reading of the reward texts can each be defended, and neither is forced.

{\small\begin{verbatim}
theorem fruit_turns_on_the_standard : Framework.grounded
    jamesUnderEvidence.defeats = Solver.toSet (Witness.listSub
    (jamesRegister.admits Standard.evidence) [JamesParty.lexical,
    JamesParty.faith, JamesParty.harmony, JamesParty.revelationText,
    JamesParty.sirachText, JamesParty.westminsterOnReward]) ∧ Nonempty
    (Contested jamesUnderTradition (p Claim.worksAreFruitNotGround))
-- axioms: propext, Classical.choice, Quot.sound
\end{verbatim}}
\subsection*{Appendix: the claims and their sources}
\addcontentsline{toc}{subsection}{Appendix: the claims and their sources}
Every claim the chapters use is listed here once, with the number every formula calls it by: the verses, the disputed readings of them, who wrote the disputed letters, what the rival positions hold, the premises that answer James, and the three parts of the conclusion. The table gives each claim’s kind, who holds it and how firmly, and the sources that carry it. Behind the table, each claim’s own description, and the reasoning for each source and rating, are kept for a reader who wants to check one.

\subsubsection*{A source for every claim}
Every claim in the argument is given a source here: a passage of Scripture, or a work and the place in it. Each also has a rating that says how well supported it is among the sources cited.

The claims are shared by all the positions. The \emph{source} for a disputed reading is not, because each side of a disagreement is cited to its own scholars. The Reformed and the New Perspective part over what Paul’s “works of the law” (ἔργα νόμου) means. Dunn and Hays part over “faith of Christ” (πίστις Χριστοῦ). The apocalyptic reading parts from both.

So there are several lists of sources and not one. \texttt{base\allowbreak{}Cite} gives the default source for every claim. \texttt{reformed\allowbreak{}Cite}, \texttt{new\allowbreak{}Perspective\allowbreak{}Cite}, \texttt{subjective\allowbreak{}Genitive\allowbreak{}Cite}, \texttt{apocalyptic\allowbreak{}Cite} and \texttt{law\allowbreak{}Observant\allowbreak{}Luke\allowbreak{}Cite} each replace the one claim where that position parts from the default, and \texttt{critical\allowbreak{}Authorship\allowbreak{}Cite} replaces the four claims about who wrote the disputed letters. The disagreement between positions therefore shows up as a difference in the list of sources each one prints.

\subsubsection*{How the steps are rated}
When positions are weighed against each other, the steps a position takes are weighed as well as its premises. A step is rated \texttt{disputed} when a cited source accepts its starting points and denies its conclusion. These are the ratings the positions in the debate carry.

\subsubsection*{The claims, numbered}
\begin{longtable}{@{}r >{\raggedright\arraybackslash}p{0.30\textwidth} l >{\raggedright\arraybackslash}p{0.19\textwidth} >{\raggedright\arraybackslash}p{0.21\textwidth}@{}}
\toprule
 & claim & kind & held & source \\
\midrule
\endhead
\pvartarget{TestimonyArguments.SolaFide}{1}$P_{1}$ & Ephesians 2:8–9 teaches salvation by grace through faith, not of works & textual & Reformed Protestant, well supported & Eph 2:8-9; \cite[III.xi.7]{calvin-institutes-1960}; \cite[ad loc. Eph 2:9]{lincoln-ephesians-1990} \\
\pvartarget{TestimonyArguments.SolaFide}{2}$P_{2}$ & Romans 3:28 teaches justification by faith apart from works of the law & textual & Reformed Protestant, well supported & Rom 3:28; \cite[III.xi.19]{calvin-institutes-1960}; \cite[ad loc. Rom 3:28]{fitzmyer-romans-1993} \\
\pvartarget{TestimonyArguments.SolaFide}{3}$P_{3}$ & Galatians 2:16 teaches that no one is justified by works of the law & textual & Reformed Protestant, well supported & Gal 2:16; \cite[III.xi.19]{calvin-institutes-1960}; \cite[ad loc. Gal 2:16]{martyn-galatians-1997} \\
\pvartarget{TestimonyArguments.SolaFide}{4}$P_{4}$ & Romans 4:4–5 contrasts wages owed with a gift reckoned to the one who believes & textual & Reformed Protestant, well supported & Rom 4:4-5; \cite[III.xi.18]{calvin-institutes-1960} \\
\pvartarget{TestimonyArguments.SolaFide}{5}$P_{5}$ & Titus 3:5 teaches that God saved us not by works done in righteousness & textual & Reformed Protestant, well supported & Titus 3:5; \cite[III.xiv.5]{calvin-institutes-1960} \\
\pvartarget{TestimonyArguments.SolaFide}{6}$P_{6}$ & Paul’s ‘works of the law’ means human works in general, not boundary markers & linguistic & Reformed Protestant, disputed & \cite[III.xi.19]{calvin-institutes-1960}; \cite{westerholm-perspectives-2003}; \cite{das-paul-law-covenant-2001}; \cite{gathercole-where-boasting-2002}; \cite{carson-variegated-nomism-1-2001}; \cite{piper-future-justification-2007}; Rom 9:11-12; Rom 11:6; Gal 3:10; Gal 5:3; Jas 2:10 \\
\pvartarget{TestimonyArguments.SolaFide}{7}$P_{7}$ & In Galatians 2:16 ‘faith of Christ’ (πίστις Χριστοῦ) means our faith in Christ & linguistic & critical scholarship, disputed & \cite{dunn-once-more-pistis-christou-1991}; \cite{matlock-detheologizing-2000}; Gal 2:16; Rom 3:22; Phil 3:9 \\
\pvartarget{TestimonyArguments.SolaFide}{8}$P_{8}$ & Galatians opposes requiring circumcision, besides faith, for justification & interpretive & critical scholarship, well supported & Gal 2:3-5; Gal 5:2-4; \cite[ad loc. Gal 5:2]{martyn-galatians-1997}; \cite[ad loc. Gal 5:2]{moo-galatians-2013}; Acts 15:1; Acts 15:5; Acts 15:19; \cite[ad loc. Acts 15:1]{bruce-acts-1988} \\
\pvartarget{TestimonyArguments.SolaFide}{9}$P_{9}$ & Acts 15:9–11 — Peter: hearts cleansed by faith; saved through grace & textual & Christian, hist.-gramm., well supported & Acts 15:9-11; \cite[ad loc. Acts 15:9]{bruce-acts-1988} \\
\pvartarget{TestimonyArguments.SolaFide}{10}$P_{10}$ & The yoke refused at Acts 15:10 is the whole law as a condition of salvation & interpretive & Christian, hist.-gramm., disputed & \cite[ad loc. Acts 15:10]{bruce-acts-1988}; Acts 15:10; Gal 5:1; Jas 2:10; \cite{das-paul-law-covenant-2001} \\
\pvartarget{TestimonyArguments.SolaFide}{11}$P_{11}$ & Acts 15:20–21, 21:20–24: Jewish believers keep the law; gentiles keep part & interpretive & critical scholarship, well supported & Acts 15:20-21; Acts 21:20-24; \cite{jervell-luke-people-god-1972} \\
\pvartarget{TestimonyArguments.SolaFide}{12}$P_{12}$ & Ephesians was written by Paul & historical & Christian, hist.-gramm., disputed & \cite{hoehner-ephesians-2002}; Eph 1:1 \\
\pvartarget{TestimonyArguments.SolaFide}{13}$P_{13}$ & Titus was written by Paul & historical & Christian, hist.-gramm., disputed & \cite{mounce-pastoral-epistles-2000}; Titus 1:1 \\
\pvartarget{TestimonyArguments.SolaFide}{14}$P_{14}$ & 1 Peter was written by the apostle Peter & historical & Christian, hist.-gramm., disputed & \cite{jobes-1-peter-2005}; \cite{schreiner-peter-jude-2003}; 1 Pet 1:1 \\
\pvartarget{TestimonyArguments.SolaFide}{15}$P_{15}$ & 2 Peter was written by the apostle Peter & historical & Christian, hist.-gramm., disputed & \cite{schreiner-peter-jude-2003}; 2 Pet 1:1 \\
\pvartarget{TestimonyArguments.SolaFide}{16}$P_{16}$ & ‘The righteousness of God’ (δικαιοσύνη θεοῦ) names God’s act of rescue in Christ & interpretive & critical scholarship, disputed & \cite{campbell-deliverance-god-2009}; \cite{martyn-galatians-1997}; Rom 3:21-26 \\
\pvartarget{TestimonyArguments.SolaFide}{17}$P_{17}$ & Luke 7:50 — Jesus says ‘your faith has saved you’ after declaring sins forgiven & textual & Christian, hist.-gramm., consensus & Luke 7:50; \cite[ad loc. Luke 7:50]{marshall-luke-1978} \\
\pvartarget{TestimonyArguments.SolaFide}{18}$P_{18}$ & ‘Saved’ (σῴζω) at Luke 7:50 means saved from sin, not healed & linguistic & Christian, hist.-gramm., well supported & \cite[Apology IV, on love and the fulfilling of the law: Luke 7:47–50]{kolb-wengert-book-of-concord-2000}; \cite[54--55]{padilla-narrative-criticism-2021}; Luke 7:47-50; \cite[ad loc. Luke 7:50]{marshall-luke-1978} \\
\pvartarget{TestimonyArguments.SolaFide}{19}$P_{19}$ & Luke 7:47 — her love is the evidence of her forgiveness, not its ground & interpretive & Christian, hist.-gramm., disputed & \cite[ad loc. Luke 7:47]{marshall-luke-1978}; \cite[54--54]{padilla-narrative-criticism-2021}; \cite[Apology IV, on love and the fulfilling of the law: Luke 7:47–50]{kolb-wengert-book-of-concord-2000}; Luke 7:41-43; Luke 7:47; Luke 7:47-50 \\
\pvartarget{TestimonyArguments.SolaFide}{20}$P_{20}$ & James 2:14–26 targets a barren faith — mere assent — not Paul’s doctrine & interpretive & Reformed Protestant, well supported & \cite[ad loc. Jas 2:14]{moo-james-2000}; \cite[ad loc. Jas 2:24]{johnson-james-1995}; Jas 2:19 \\
\pvartarget{TestimonyArguments.SolaFide}{21}$P_{21}$ & Good works are the fruit and evidence of saving faith, not its ground & theological & Reformed Protestant, disputed & \cite[XI.2; XVI.2; XVI.5]{westminster-confession-1647}; \cite[III.xvi.1]{calvin-institutes-1960}; Jas 2:18-22; Eph 2:10; Rom 6:22; Luke 17:10 \\
\pvartarget{TestimonyArguments.SolaFide}{22}$P_{22}$ & James 2:24 is compatible with Paul, using ‘justify’ in a different sense & interpretive & Reformed Protestant, disputed & Jas 2:24; \cite[ad loc. Jas 2:24]{moo-james-2000}; \cite[III.xvii.11]{calvin-institutes-1960} \\
\pvartarget{TestimonyArguments.SolaFide}{23}$P_{23}$ & James 2:24 says a person is justified by works and not by faith alone & textual & Christian, hist.-gramm., consensus & Jas 2:24; \cite[ad loc. Jas 2:24]{moo-james-2000}; \cite[ad loc. Jas 2:24]{johnson-james-1995} \\
\pvartarget{TestimonyArguments.SolaFide}{24}$P_{24}$ & Scripture does not contradict itself & theological & Reformed Protestant, well supported & \cite[I.vii]{calvin-institutes-1960}; 2 Pet 3:15-16; \cite[ad loc. 2 Pet 3:16]{schreiner-peter-jude-2003} \\
\pvartarget{TestimonyArguments.SolaFide}{25}$P_{25}$ & Justification is by faith alone & theological & Reformed Protestant, well supported & \cite[III.xi.1]{calvin-institutes-1960}; \cite{schreiner-faith-alone-2015}; \cite{mannermaa-christ-present-faith-2005}; \cite[1 Clement 32.4]{roberts-ante-nicene-fathers-1-1885}; \cite[8.6]{hilary-commentary-matthew-2013}; \cite[Homily 20, Of Humility]{basil-ascetical-works-1999} \\
\pvartarget{TestimonyArguments.SolaFide}{26}$P_{26}$ & Faith is sufficient: whoever believes is saved & theological & Reformed Protestant, disputed & Luke 7:50; Acts 15:9-11; \cite{schreiner-faith-alone-2015} \\
\pvartarget{TestimonyArguments.SolaFide}{27}$P_{27}$ & Salvation is by grace, a gift and not wages owed & theological & Reformed Protestant, well supported & \cite[III.xi–xviii]{calvin-institutes-1960}; 1 Pet 1:18-19; \cite[ad loc. 1 Pet 1:18]{jobes-1-peter-2005} \\
\pvartarget{TestimonyArguments.SolaFide}{28}$P_{28}$ & Salvation is not by works: no work is its ground & theological & Reformed Protestant, well supported & \cite[III.xi–xviii]{calvin-institutes-1960}; \cite{schreiner-faith-alone-2015} \\
\pvartarget{TestimonyArguments.SolaFide}{29}$P_{29}$ & Salvation is received through faith & theological & Reformed Protestant, well supported & \cite[III.xi–xviii]{calvin-institutes-1960}; 1 Pet 1:3-5; \cite[ad loc. 1 Pet 1:5]{jobes-1-peter-2005} \\
\pvartarget{TestimonyArguments.SolaFide}{30}$P_{30}$ & Justification is a verdict only: pardon and righteousness counted, not poured in & theological & Reformed Protestant, well supported & \cite[XI.1]{westminster-confession-1647}; \cite[III.xi.2]{calvin-institutes-1960}; \cite[Epistle to Diognetus 9]{roberts-ante-nicene-fathers-1-1885} \\
\pvartarget{TestimonyArguments.SolaFide}{31}$P_{31}$ & Second Temple Judaism: in by grace, staying in by works & historical & critical scholarship, disputed & \cite{sanders-paul-palestinian-judaism-1977} \\
\pvartarget{TestimonyArguments.SolaFide}{32}$P_{32}$ & Works performed in grace merit an increase of justification & theological & Roman Catholic, well supported & \cite[Trent, Session VI (1547), Decree on Justification, ch. 16]{tanner-decrees-1990} \\
\pvartarget{TestimonyArguments.SolaFide}{33}$P_{33}$ & Justification is not remission of sins only, but renewal of the inward man & theological & Roman Catholic, disputed & \cite[Trent, Session VI (1547), Decree on Justification, ch. 7]{tanner-decrees-1990}; \cite[Trent, Session VI (1547), canon 11]{tanner-decrees-1990} \\
\pvartarget{TestimonyArguments.SolaFide}{34}$P_{34}$ & Justification and sanctification are inseparable but distinct & theological & Reformed Protestant, disputed & \cite[Q. 77]{westminster-larger-catechism-1647}; 1 Cor 6:11; 1 Cor 1:30; Rom 4:6; Rom 4:8; Ezek 36:27; Heb 9:13-14; Rom 3:24-25; Rom 6:6; Rom 6:14; \cite[XI.1]{westminster-confession-1647}; \cite[III.xi.6]{calvin-institutes-1960} \\
\pvartarget{TestimonyArguments.SolaFide}{35}$P_{35}$ & The inward renewal of the justified grows as they do good works in grace & theological & Roman Catholic, well supported & \cite[Trent, Session VI (1547), Decree on Justification, ch. 10]{tanner-decrees-1990}; \cite[XIII.1]{westminster-confession-1647} \\
\pvartarget{TestimonyArguments.SolaFide}{36}$P_{36}$ & John 6:28–29 — the work God requires is that you believe in him whom he sent & textual & Christian, hist.-gramm., consensus & John 6:28-29; \cite[ad loc. John 6:29]{calvin-commentary-john-1847}; \cite[cap. 6, lect. 3, n. 901]{aquinas-commentary-john-2012} \\
\pvartarget{TestimonyArguments.SolaFide}{37}$P_{37}$ & John 3:16–18, 3:36, 5:24, 20:31 — eternal life through believing in the Son\textsuperscript{\dag} & textual & Christian, hist.-gramm., consensus & John 3:16-18; John 3:36; John 5:24; John 20:31 \\
\pvartarget{TestimonyArguments.SolaFide}{38}$P_{38}$ & John 3:36 — whoever believes has life; whoever does not obey shall not see life & textual & Christian, hist.-gramm., consensus & John 3:36; \cite[62]{westcott-john-1892}; \cite[ad loc. John 3:36]{calvin-commentary-john-1847} \\
\pvartarget{TestimonyArguments.SolaFide}{39}$P_{39}$ & The obedience of John 3:36 and Matthew 7:21 is sanctification, the fruit of faith & interpretive & Reformed Protestant, disputed & \cite[ad loc. John 3:36]{ridderbos-john-1997}; \cite[ad loc. John 3:36]{carson-john-1991}; \cite[ad loc. Matt 7:21]{morris-matthew-1992}; \cite[55--55]{schreiner-justification-by-works-2015}; \cite{hoekema-saved-by-grace-1994}; \cite{murray-redemption-1984}; \cite[XVI.2]{westminster-confession-1647}; \cite[XIII.1]{westminster-confession-1647}; \cite[ad loc. John 3:36]{calvin-commentary-john-1847} \\
\pvartarget{TestimonyArguments.SolaFide}{40}$P_{40}$ & The believing of John 6:29 is trust, which brings nothing and receives Christ & interpretive & Reformed Protestant, disputed & \cite[ad loc. John 6:29]{calvin-commentary-john-1847}; John 6:28-29 \\
\pvartarget{TestimonyArguments.SolaFide}{41}$P_{41}$ & The believing of John 6:29 is faith living through charity, the source of works & interpretive & Roman Catholic, disputed & \cite[cap. 6, lect. 3, n. 901]{aquinas-commentary-john-2012}; \cite[Trent, Session VI (1547), Decree on Justification, ch. 7]{tanner-decrees-1990} \\
\pvartarget{TestimonyArguments.SolaFide}{42}$P_{42}$ & Galatians 3:11–12 — none is justified by the law; the law is not of faith & textual & Christian, hist.-gramm., consensus & Gal 3:11-12; \cite[ad loc. Gal 3:11]{luther-commentary-galatians-1998} \\
\pvartarget{TestimonyArguments.SolaFide}{43}$P_{43}$ & Love of God and neighbour is what the law commands & textual & Christian, hist.-gramm., consensus & Deut 6:5; Matt 22:37-40; \cite[ad loc. Gal 3:12]{luther-commentary-galatians-1998} \\
\pvartarget{TestimonyArguments.SolaFide}{44}$P_{44}$ & Galatians 1:6–9: whoever preaches a gospel contrary to the one received is accursed & textual & Christian, hist.-gramm., consensus & Gal 1:6-9; \cite[ad loc. Gal 1:6]{moo-galatians-2013} \\
\pvartarget{TestimonyArguments.SolaFide}{45}$P_{45}$ & 1 Corinthians 15:3: the gospel of first importance — Christ died for our sins\textsuperscript{\dag} & textual & Christian, hist.-gramm., consensus & 1 Cor 15:3 \\
\pvartarget{TestimonyArguments.SolaFide}{46}$P_{46}$ & Galatians 2:21: if righteousness were through the law, Christ died for no purpose & textual & Christian, hist.-gramm., consensus & Gal 2:21; \cite[ad loc. Gal 2:21]{moo-galatians-2013} \\
\pvartarget{TestimonyArguments.SolaFide}{47}$P_{47}$ & Galatians 5:2–4: you who would be justified by the law are severed from Christ & textual & Christian, hist.-gramm., consensus & Gal 5:2-4; \cite[ad loc. Gal 5:2]{moo-galatians-2013} \\
\pvartarget{TestimonyArguments.SolaFide}{48}$P_{48}$ & Romans 8:33–34: it is God who justifies — who is to condemn?\textsuperscript{\dag} & textual & Christian, hist.-gramm., consensus & Rom 8:33-34 \\
\pvartarget{TestimonyArguments.SolaFide}{49}$P_{49}$ & In Paul, ‘justify’ (δικαιόω) is a verdict: to declare righteous & linguistic & Christian, hist.-gramm., well supported & \cite[s.v. δικαιόω]{bauer-danker-lexicon-2000}; \cite{morris-apostolic-preaching-1965}; \cite[p. 80]{moo-romans-1996}; \cite{irons-righteousness-of-god-2015}; \cite{wright-justification-2009}; \cite[ad loc. Rom 3:24]{fitzmyer-romans-1993} \\
\pvartarget{TestimonyArguments.SolaFide}{50}$P_{50}$ & Paul names renewal with other words (ἀνακαίνωσις, ἁγιασμός) and sets them beside ‘justify’\textsuperscript{\dag} & textual & Christian, hist.-gramm., consensus & Titus 3:5-7; Rom 12:2; Rom 6:19-22; 1 Cor 6:11 \\
\pvartarget{TestimonyArguments.SolaFide}{51}$P_{51}$ & A word contributes the least meaning its context requires; one occurrence does not carry the whole concept & linguistic & Christian, hist.-gramm., well supported & \cite{silva-biblical-words-1994}; \cite{joos-semantic-axiom-1972}; \cite{barr-semantics-2004} \\
\pvartarget{TestimonyArguments.SolaFide}{52}$P_{52}$ & Paul’s word ‘justify’ means the renewal of the inward man, not only a verdict & linguistic & Roman Catholic, disputed & \cite[Trent, Session VI (1547), Decree on Justification, chs. 7–8]{tanner-decrees-1990}; \cite{vanlandingham-judgment-justification-2006} \\
\pvartarget{TestimonyArguments.SolaFide}{53}$P_{53}$ & When God justifies, he also renews: pardon and renewal are given together, not to be separated & theological & Christian, hist.-gramm., well supported & \cite[§22]{lwf-catholic-joint-declaration-2000}; \cite[Trent, Session VI (1547), Decree on Justification, ch. 7]{tanner-decrees-1990}; \cite[III.xi.6]{calvin-institutes-1960}; \cite{gorman-cruciform-god-2009} \\
\pvartarget{TestimonyArguments.SolaFide}{54}$P_{54}$ & Augustine explains “being justified” as “being made righteous”, and the Latin West read iustificare so & historical & Christian, hist.-gramm., consensus & \cite[On the Spirit and the Letter, ch. 45 (26.45)]{npnf1-05-1887}; \cite{mcgrath-iustitia-dei-2005}; \cite[Trent, Session IV (1546), Decree on the Vulgate]{tanner-decrees-1990} \\
\pvartarget{TestimonyArguments.SolaFide}{55}$P_{55}$ & Romans 4:6–8: righteousness counted apart from works — sins forgiven, sin not counted & textual & Christian, hist.-gramm., consensus & Rom 4:6-8; \cite[ad loc. Rom 4:6]{moo-romans-1996} \\
\pvartarget{TestimonyArguments.SolaFide}{56}$P_{56}$ & ‘Count’ (λογίζομαι) in Romans 4:3–11 is the language of an account: credited to one & linguistic & Christian, hist.-gramm., well supported & \cite[s.v. λογίζομαι]{bauer-danker-lexicon-2000}; \cite[ad loc. Rom 4:3]{moo-romans-1996} \\
\pvartarget{TestimonyArguments.SolaFide}{57}$P_{57}$ & In Romans 4:6–8 Paul explains righteousness counted apart from works as sin not counted & interpretive & Christian, hist.-gramm., well supported & \cite[ad loc. Rom 4:6]{moo-romans-1996}; Rom 4:6-8 \\
\pvartarget{TestimonyArguments.SolaFide}{58}$P_{58}$ & Christ’s death for our sins is the whole ground of justification; to add to it is another gospel & interpretive & Reformed Protestant, disputed & \cite[XI.1]{westminster-confession-1647}; Gal 2:21; \cite[ad loc. Gal 2:21]{moo-galatians-2013}; \cite[Epistle to Diognetus 9]{roberts-ante-nicene-fathers-1-1885} \\
\pvartarget{TestimonyArguments.SolaFide}{59}$P_{59}$ & Luke 18:9–14 — the tax collector pleads for mercy and goes home justified, rather than the Pharisee who listed his works & textual & Christian, hist.-gramm., consensus & Luke 18:9-14; \cite[201--207]{calvin-harmony-evangelists-2-1845} \\
\pvartarget{TestimonyArguments.SolaFide}{60}$P_{60}$ & ‘Justified’ at Luke 18:14 is God’s verdict: accepted as righteous before God & linguistic & Christian, hist.-gramm., well supported & \cite[206--207]{calvin-harmony-evangelists-2-1845}; Luke 18:14 \\
\pvartarget{TestimonyArguments.SolaFide}{61}$P_{61}$ & The tax collector brings no works that merit, only a humble and contrite plea for mercy & interpretive & Christian, hist.-gramm., consensus & Luke 18:9-14; \cite[206--207]{calvin-harmony-evangelists-2-1845}; \cite[36]{waterworth-trent-1848} \\
\pvartarget{TestimonyArguments.SolaFide}{62}$P_{62}$ & Matthew 7:21–23 — the criterion at the judgment is doing the Father’s will, not profession or mighty works & textual & Christian, hist.-gramm., well supported & Matt 7:21-23; \cite[367--368]{calvin-harmony-evangelists-1845}; \cite[2826]{catholic-church-catechism-2000} \\
\pvartarget{TestimonyArguments.SolaFide}{63}$P_{63}$ & In Matthew 7:21–23, doing the Father’s will is a condition of entering the kingdom in its own right, not only the fruit of faith & interpretive & Roman Catholic, disputed & \cite[2826]{catholic-church-catechism-2000}; \cite[1821]{catholic-church-catechism-2000} \\
\pvartarget{TestimonyArguments.SolaFide}{64}$P_{64}$ & In Matthew 7:21–23, doing the Father’s will includes believing in Christ, and “I never knew you” means no union with him by faith & interpretive & Reformed Protestant, disputed & \cite[367--368]{calvin-harmony-evangelists-1845} \\
\pvartarget{TestimonyArguments.SolaFide}{65}$P_{65}$ & Matthew 19:17 — “If you would enter life, keep the commandments” & textual & Christian, hist.-gramm., consensus & Matt 19:16-22; \cite[35]{waterworth-trent-1848} \\
\pvartarget{TestimonyArguments.SolaFide}{66}$P_{66}$ & In Matthew 19:17, keeping the commandments is the way to life, made possible by grace & interpretive & Roman Catholic, disputed & \cite[12]{john-paul-ii-veritatis-splendor-1993}; \cite[2052–2053]{catholic-church-catechism-2000}; \cite[35]{waterworth-trent-1848} \\
\pvartarget{TestimonyArguments.SolaFide}{67}$P_{67}$ & In Matthew 19:17, Christ answers on the law’s own terms, so that the young man may see he has not kept it and turn to faith & interpretive & Reformed Protestant, disputed & \cite[393--394]{calvin-harmony-evangelists-2-1845} \\
\pvartarget{TestimonyArguments.SolaFide}{68}$P_{68}$ & Water baptism is the instrumental cause of justification & theological & Roman Catholic, disputed & \cite[34--35]{waterworth-trent-1848}; \cite[XI.2]{westminster-confession-1647} \\
\pvartarget{TestimonyArguments.SolaFide}{69}$P_{69}$ & The washing itself is necessary for salvation, though desire may justify & theological & Roman Catholic, disputed & \cite{feeney-bread-of-life-1952} \\
\pvartarget{TestimonyArguments.SolaFide}{70}$P_{70}$ & Justification cannot be effected without the laver of regeneration, or the desire thereof & textual & Roman Catholic, consensus & \cite[32]{waterworth-trent-1848}; \cite[III q.68 a.2]{aquinas-summa-1920} \\
\pvartarget{TestimonyArguments.SolaFide}{71}$P_{71}$ & The desire for baptism brings about its fruits without being a sacrament & theological & Roman Catholic, consensus & \cite[1258]{catholic-church-catechism-2000}; \cite[III q.69 a.4 ad 2]{aquinas-summa-1920} \\
\pvartarget{TestimonyArguments.SolaFide}{72}$P_{72}$ & Luke 23:43 — the thief, unbaptised, is promised paradise & textual & Christian, hist.-gramm., consensus & Luke 23:43; \cite[On Baptism IV.23.31]{npnf1-04-1887} \\
\pvartarget{TestimonyArguments.SolaFide}{73}$P_{73}$ & The thief’s case bears on salvation under the Gospel & interpretive & Christian, hist.-gramm., well supported & \cite[On Baptism IV.22.30–23.31]{npnf1-04-1887}; \cite[III q.66 a.11]{aquinas-summa-1920} \\
\pvartarget{TestimonyArguments.SolaFide}{74}$P_{74}$ & Acts 8:14–17 — the Samaritans, baptized, receive the Spirit afterwards & textual & Christian, hist.-gramm., consensus & Acts 8:14-17; \cite[338--339]{calvin-commentary-acts-1844} \\
\pvartarget{TestimonyArguments.SolaFide}{75}$P_{75}$ & Acts 10:44–48 — the Spirit falls on Cornelius’s household before baptism & textual & Christian, hist.-gramm., consensus & Acts 10:44-48; \cite[453--454]{calvin-commentary-acts-1844} \\
\pvartarget{TestimonyArguments.SolaFide}{76}$P_{76}$ & God is not bound to water baptism: he gives the Spirit before it and after it & theological & Christian, hist.-gramm., consensus & \cite[1257]{catholic-church-catechism-2000}; \cite[454]{calvin-commentary-acts-1844}; \cite[III q.69 a.4 ad 2]{aquinas-summa-1920} \\
\pvartarget{TestimonyArguments.SolaFide}{77}$P_{77}$ & John 3:5 — unless one is born of water and the Spirit, he cannot enter the kingdom & textual & Christian, hist.-gramm., consensus & John 3:5; \cite{oliver-water-john-2022} \\
\pvartarget{TestimonyArguments.SolaFide}{78}$P_{78}$ & John 3:5 requires water baptism for entry into the kingdom & interpretive & Roman Catholic, disputed & \cite[56]{waterworth-trent-1848}; \cite[1257]{catholic-church-catechism-2000} \\
\pvartarget{TestimonyArguments.SolaFide}{79}$P_{79}$ & The water of John 3:5 is baptismal water & interpretive & Roman Catholic, disputed & \cite[56]{waterworth-trent-1848}; \cite[110]{calvin-commentary-john-1847} \\
\pvartarget{TestimonyArguments.SolaFide}{80}$P_{80}$ & The water of John 3:5 is natural birth & interpretive & critical scholarship, disputed & \cite{oliver-water-john-2022} \\
\pvartarget{TestimonyArguments.SolaFide}{81}$P_{81}$ & The water of John 3:5 is the cleansing promised in Ezekiel 36:25–27, not baptism & interpretive & Christian, hist.-gramm., disputed & \cite{carson-born-of-water-2019}; \cite{oliver-water-john-2022} \\
\pvartarget{TestimonyArguments.SolaFide}{82}$P_{82}$ & The water of John 3:5 is the Spirit’s own cleansing: water and Spirit are one & interpretive & Reformed Protestant, disputed & \cite[110--111]{calvin-commentary-john-1847}; \cite[IV.xvi.25]{calvin-institutes-1960} \\
\pvartarget{TestimonyArguments.SolaFide}{83}$P_{83}$ & James 2:14–26: faith without works cannot save, and is dead (2:17, 2:26) & textual & Christian, hist.-gramm., consensus & Jas 2:14-26; \cite[ad loc. Jas 2:14]{moo-james-2000}; \cite[ad loc. Jas 2:14]{johnson-james-1995} \\
\pvartarget{TestimonyArguments.SolaFide}{84}$P_{84}$ & James 2:21–25: Abraham justified by works at Genesis 22, fulfilling Genesis 15:6; Rahab justified by works\textsuperscript{\dag} & textual & Christian, hist.-gramm., consensus & Jas 2:21-25; Gen 15:6; Gen 22:9-12 \\
\pvartarget{TestimonyArguments.SolaFide}{85}$P_{85}$ & Outside Paul, ‘justify’ (δικαιόω) declares or vindicates; it does not mean making righteous & linguistic & Christian, hist.-gramm., well supported & \cite[s.v. δικαιόω]{bauer-danker-lexicon-2000}; Deut 25:1; Matt 11:19; Luke 7:35; \cite[ad loc. Matt 11:19]{crossway-esv-2001} \\
\pvartarget{TestimonyArguments.SolaFide}{86}$P_{86}$ & James’s ‘justified’ (δικαιόω) means the increase of the justice received & linguistic & Roman Catholic, disputed & \cite[Trent, Session VI (1547), ch. 10]{tanner-decrees-1990} \\
\pvartarget{TestimonyArguments.SolaFide}{87}$P_{87}$ & The faith alone that James 2:24 denies justifies is the dead faith of 2:14–26 & interpretive & Christian, hist.-gramm., well supported & Jas 2:14-26; \cite[Trent, Session VI (1547), ch. 7]{tanner-decrees-1990}; \cite[ad loc. Jas 2:24]{moo-james-2000} \\
\pvartarget{TestimonyArguments.SolaFide}{88}$P_{88}$ & The faith the Reformers say alone justifies is no dead faith, but works by love & historical & Reformed Protestant, consensus & \cite[XI.2]{westminster-confession-1647}; \cite[III.xvi.1]{calvin-institutes-1960} \\
\pvartarget{TestimonyArguments.SolaFide}{89}$P_{89}$ & The faith James 2:24 denies justifies is not the faith of the Reformers’ formula & interpretive & Christian, hist.-gramm., well supported & \cite[ad loc. Jas 2:24]{moo-james-2000} \\
\pvartarget{TestimonyArguments.SolaFide}{90}$P_{90}$ & Scripture teaches an increase of justification by good works & theological & Roman Catholic, disputed & \cite[Trent, Session VI (1547), ch. 10]{tanner-decrees-1990} \\
\pvartarget{TestimonyArguments.SolaFide}{91}$P_{91}$ & Good works are a cause of the increase of justification, not merely its fruits & theological & Roman Catholic, disputed & \cite[Trent, Session VI (1547), canon 24]{tanner-decrees-1990} \\
\pvartarget{TestimonyArguments.SolaFide}{92}$P_{92}$ & Revelation 22:11 says: he that is just, let him be justified still & textual & Roman Catholic, disputed & \cite[Trent, Session VI (1547), ch. 10]{tanner-decrees-1990} \\
\pvartarget{TestimonyArguments.SolaFide}{93}$P_{93}$ & The Greek of Revelation 22:11 says: let the righteous still do righteousness & textual & Christian, hist.-gramm., consensus & \cite[app. Rev 22:11]{holmes-sblgnt-2010}; \cite[Revelation 22:11]{usccb-nabre-2011} \\
\pvartarget{TestimonyArguments.SolaFide}{94}$P_{94}$ & Sirach 18:22 says: be not afraid to be justified even to death & textual & Roman Catholic, disputed & \cite[Trent, Session VI (1547), ch. 10]{tanner-decrees-1990} \\
\pvartarget{TestimonyArguments.SolaFide}{95}$P_{95}$ & The Greek of Sirach 18:22 is about paying a vow: do not wait until death to be justified & textual & Christian, hist.-gramm., well supported & \cite[ad loc. Sir 18:22]{swete-old-testament-in-greek}; \cite[Sirach 18:22]{usccb-nabre-2011} \\
\pvartarget{TestimonyArguments.SolaFide}{96}$P_{96}$ & The original text, written by the inspired author, outweighs any translation: a reading it lacks is not his word & linguistic & Christian, hist.-gramm., consensus & \cite[§16]{pius-divino-afflante-spiritu-1943}; \cite[I.8]{westminster-confession-1647} \\
\pvartarget{TestimonyArguments.SolaFide}{97}$P_{97}$ & The Vulgate, as the Church understands it, is free from error in faith and morals: its authenticity is juridical & theological & Roman Catholic, disputed & \cite[§21]{pius-divino-afflante-spiritu-1943}; \cite[Trent, Session IV (1546)]{tanner-decrees-1990} \\
\pvartarget{TestimonyArguments.SolaFide}{98}$P_{98}$ & Trent quotes the Vulgate’s Revelation 22:11 and Sirach 18:22 as juridically authentic, not as the inspired authors’ words & interpretive & Roman Catholic, disputed & \cite[§§20–21]{pius-divino-afflante-spiritu-1943} \\
\pvartarget{TestimonyArguments.SolaFide}{99}$P_{99}$ & A doctrine taught from the Vulgate is to be confirmed from the original texts & theological & Roman Catholic, well supported & \cite[§22]{pius-divino-afflante-spiritu-1943}; \cite[I.8]{westminster-confession-1647} \\
\pvartarget{TestimonyArguments.SolaFide}{100}$P_{100}$ & Revelation 22:11 and Sirach 18:22 confirm the increase of justification from the original texts & textual & Roman Catholic, disputed & \cite[Trent, Session VI (1547), ch. 10]{tanner-decrees-1990} \\
\pvartarget{TestimonyArguments.SolaFide}{101}$P_{101}$ & Ephesians 2:10: created in Christ Jesus for good works, after ‘not a result of works’ & textual & Christian, hist.-gramm., consensus & Eph 2:8-10; \cite[XVI.2]{westminster-confession-1647} \\
\pvartarget{TestimonyArguments.SolaFide}{102}$P_{102}$ & James 2:18, 22: I will show you my faith by my works; faith completed by works & textual & Christian, hist.-gramm., consensus & Jas 2:18-22; \cite[XVI.2]{westminster-confession-1647} \\
\pvartarget{TestimonyArguments.SolaFide}{103}$P_{103}$ & Romans 6:22: the fruit you get leads to sanctification, and its end is eternal life & textual & Christian, hist.-gramm., consensus & Rom 6:22; \cite[XVI.2]{westminster-confession-1647} \\
\pvartarget{TestimonyArguments.SolaFide}{104}$P_{104}$ & Luke 17:10: having done all that was commanded, we are unworthy servants & textual & Christian, hist.-gramm., consensus & Luke 17:10; \cite[XVI.5]{westminster-confession-1647} \\
\pvartarget{TestimonyArguments.SolaFide}{105}$P_{105}$ & 1 Cor 15:58; Heb 6:10; 10:35; 2 Tim 4:8: God rewards the labor of believers & textual & Christian, hist.-gramm., consensus & 1 Cor 15:58; Heb 6:10; Heb 10:35; 2 Tim 4:8; \cite[Trent, Session VI (1547), ch. 16]{tanner-decrees-1990}; \cite[XVI.6]{westminster-confession-1647} \\
\pvartarget{TestimonyArguments.SolaFide}{106}$P_{106}$ & James’s ‘justify’ (δικαιόω) has the sense Paul’s has: the justification of Romans and Galatians & linguistic & Roman Catholic, disputed & \cite[Trent, Session VI (1547), ch. 10]{tanner-decrees-1990} \\
\pvartarget{TestimonyArguments.SolaFide}{107}$P_{107}$ & The reward God renders to the works of the justified is rendered to them as merits & theological & Roman Catholic, disputed & \cite[Trent, Session VI (1547), ch. 16; canon 32]{tanner-decrees-1990}; \cite{dulles-two-languages-1999} \\
\pvartarget{TestimonyArguments.SolaFide}{108}$P_{108}$ & 1 Peter 2:5: spiritual sacrifices acceptable to God through Jesus Christ & textual & Christian, hist.-gramm., consensus & 1 Pet 2:5; \cite[XVI.6]{westminster-confession-1647} \\
\pvartarget{TestimonyArguments.SolaFide}{109}$P_{109}$ & Justification has degrees: the justified are still further justified & theological & Roman Catholic, well supported & \cite[37]{waterworth-trent-1848}; \cite[47]{waterworth-trent-1848} \\
\pvartarget{TestimonyArguments.SolaFide}{110}$P_{110}$ & Deuteronomy 25:1, Proverbs 17:15: judges acquit the innocent and condemn the guilty & textual & Christian, hist.-gramm., consensus & Deut 25:1; Prov 17:15; \cite[ad loc. Deut 25:1]{crossway-esv-2001} \\
\pvartarget{TestimonyArguments.SolaFide}{111}$P_{111}$ & A verdict acquits or condemns, and has no degrees & interpretive & Reformed Protestant, well supported & \cite[251]{buchanan-justification-1867}; \cite[116]{calvin-antidote-trent-1851}; \cite[Topic XVI, q. 10]{turretin-institutes-2-1993} \\
\pvartarget{TestimonyArguments.SolaFide}{112}$P_{112}$ & Genesis 15:6: Abram believed the LORD, and he counted it to him as righteousness & textual & Christian, hist.-gramm., consensus & Gen 15:6; \cite[ad loc. Gen 15:6]{crossway-esv-2001} \\
\pvartarget{TestimonyArguments.SolaFide}{113}$P_{113}$ & Genesis 15:12, 17: Abram asleep; a fire pot and a torch pass between the pieces & textual & Christian, hist.-gramm., consensus & Gen 15:12; Gen 15:17; \cite[ad loc. Gen 15:17]{crossway-esv-2001} \\
\pvartarget{TestimonyArguments.SolaFide}{114}$P_{114}$ & Romans 4:10: righteousness was counted to Abraham before he was circumcised & textual & Christian, hist.-gramm., consensus & Rom 4:10; \cite[ad loc. Rom 4:10]{crossway-esv-2001} \\
\pvartarget{TestimonyArguments.SolaFide}{115}$P_{115}$ & The covenant of Genesis 15 binds God alone, by an oath on himself & interpretive & Reformed Protestant, disputed & Jer 34:18; Heb 6:13-18; \cite{robertson-christ-covenants-1980}; \cite{kline-kingdom-prologue-2006}; \cite{wenham-genesis-1-15-1987}; \cite{weinfeld-covenant-of-grant-1970} \\
\pvartarget{TestimonyArguments.SolaFide}{116}$P_{116}$ & Abraham was counted righteous on believing, with no work as a condition & theological & Reformed Protestant, disputed & Rom 4:3-5; \cite[III.xi.18]{calvin-institutes-1960} \\
\pvartarget{TestimonyArguments.SolaFide}{117}$P_{117}$ & Genesis 22:16–18, 26:5: blessed because Abraham obeyed & textual & Christian, hist.-gramm., consensus & Gen 22:16-18; Gen 26:5; \cite[ad loc. Gen 22:16]{crossway-esv-2001} \\
\pvartarget{TestimonyArguments.SolaFide}{118}$P_{118}$ & 1 Maccabees 2:52, Sirach 44:20: Abraham found faithful in trial, and it was credited to him as righteousness & textual & Second Temple Jewish, consensus & \cite[ad loc. 1 Macc 2:52]{usccb-nabre-2011}; \cite[ad loc. Sir 44:20]{usccb-nabre-2011} \\
\pvartarget{TestimonyArguments.SolaFide}{119}$P_{119}$ & A royal grant rewarded loyalty already shown & historical & critical scholarship, disputed & \cite[185]{weinfeld-covenant-of-grant-1970}; \cite[672]{knoppers-royal-grants-1996} \\
\pvartarget{TestimonyArguments.SolaFide}{120}$P_{120}$ & The grant to Abraham rewarded his obedience (Genesis 22:18; 26:5) & interpretive & critical scholarship, disputed & \cite{weinfeld-covenant-of-grant-1970}; Gen 22:18; Gen 26:5 \\
\pvartarget{TestimonyArguments.SolaFide}{121}$P_{121}$ & The righteousness counted at Genesis 15:6 was credited for Abraham’s faithfulness when tested & interpretive & Second Temple Jewish, disputed & \cite[ad loc. 1 Macc 2:52]{usccb-nabre-2011} \\
\pvartarget{TestimonyArguments.SolaFide}{122}$P_{122}$ & The blessing of Genesis 22:16–18 was augmented because of Abraham’s obedience & interpretive & Christian, hist.-gramm., disputed & \cite[507]{kuruvilla-aqedah-2012} \\
\pvartarget{TestimonyArguments.SolaFide}{123}$P_{123}$ & The oath of Genesis 22:16–18 was merited by the binding of Isaac & interpretive & Roman Catholic, disputed & \cite[93]{hahn-covenant-oath-aqedah-2005} \\
\bottomrule
\end{longtable}

\noindent\footnotesize \textsuperscript{\dag} grounded in scripture alone: the reading is assumed, not argued for.
\normalsize
\subsubsection*{How this is checked}
\par\medskip\noindent\textbf{\texttt{trent\allowbreak{}Against\allowbreak{}Faith\allowbreak{}Alone}}\par\nobreak\noindent
\textbf{Trent contests the step in every Reformed strand.} Trent accepts what each strand reads: Romans 3:28 and Galatians 2:16, that Jesus’ words save, that salvation is through grace. It condemns the conclusion each draws, justification by faith alone (Session VI, canon 9).

{\small\begin{verbatim}
def trentAgainstFaithAlone : Source :=
  {
    primary :=
      Reference.work tannerDecrees
        (Bib.Locus.sectionRef "Trent, Session VI (1547), canon 9"),
    tradition := Tradition.romanCatholic, confidence := Confidence.disputed }
\end{verbatim}}
\par\medskip\noindent\textbf{\texttt{trent\allowbreak{}On\allowbreak{}Merit}}\par\nobreak\noindent
Trent’s own step, from merit to the denial of “not by works”. Rated \texttt{consensus}: both sides accept the step. It is the reason the Reformed deny merit. The disagreement is over the premise, not the step.

{\small\begin{verbatim}
def trentOnMerit : Source :=
  {
    primary :=
      Reference.work tannerDecrees
        (Bib.Locus.sectionRef
          "Trent, Session VI (1547), Decree on Justification, ch. 16"),
    tradition := Tradition.romanCatholic, confidence := Confidence.consensus }
\end{verbatim}}
\par\medskip\noindent\textbf{\texttt{trent\allowbreak{}Against\allowbreak{}The\allowbreak{}Plea\allowbreak{}Alone}}\par\nobreak\noindent
\textbf{Trent, applied to Luke 18, contests the step from it.} Trent never cites Luke 18, and accepts that nothing before justification, “whether faith or works”, merits it (ch. 8). But it requires the dispositions it lists, fear, hope, beginning to love God, and penitence (ch. 6), and it condemns whoever says “nothing else is required to co-operate” (canon 9). Read by those chapters, the tax collector’s humility and contrition are dispositions his justification required, and not a plea that was enough alone. So the step from the plea to “faith suffices” is \texttt{disputed}.

{\small\begin{verbatim}
def trentAgainstThePleaAlone : Source :=
  { primary := Reference.work waterworthTrent (Bib.Locus.page 45),
    supporting := [Reference.work waterworthTrent (Bib.Locus.pages 33 34)],
    tradition := Tradition.romanCatholic, confidence := Confidence.disputed }
\end{verbatim}}
\par\medskip\noindent\textbf{\texttt{calvin\allowbreak{}Against\allowbreak{}Obedience\allowbreak{}As\allowbreak{}Ground}}\par\nobreak\noindent
\textbf{Calvin contests the Catholic step from Matthew 7:21.} He accepts the text, and that entry requires doing the Father’s will. He denies that the doing is more than the fruit of faith: “These words, therefore, do not exclude faith, but presuppose it as the principle from which other good works flow.”

{\small\begin{verbatim}
def calvinAgainstObedienceAsGround : Source :=
  { primary := Reference.work calvinHarmony1 (Bib.Locus.pages 367 368),
    tradition := Tradition.reformedProtestant,
    confidence := Confidence.disputed }
\end{verbatim}}
\par\medskip\noindent\textbf{\texttt{catechism\allowbreak{}Against\allowbreak{}Doing\allowbreak{}As\allowbreak{}Believing}}\par\nobreak\noindent
\textbf{The Catechism contests the Reformed step from Matthew 7:21.} It accepts the text and reads it the other way: “one enters the kingdom of heaven not by speaking words, but by doing ‘the will of my Father in heaven’” (§2826).

{\small\begin{verbatim}
def catechismAgainstDoingAsBelieving : Source :=
  {
    primary :=
      Reference.work catechismCatholicChurch (Bib.Locus.sectionRef "2826"),
    tradition := Tradition.romanCatholic, confidence := Confidence.disputed }
\end{verbatim}}
\par\medskip\noindent\textbf{\texttt{calvin\allowbreak{}Against\allowbreak{}The\allowbreak{}Commandments\allowbreak{}As\allowbreak{}The\allowbreak{}Way}}\par\nobreak\noindent
\textbf{Calvin contests the Catholic step from Matthew 19:17.} He accepts the text: “we have no right, therefore, to deny that the keeping of the law is righteousness”. He denies that it is the way to life for any sinner: “this passage was erroneously interpreted by some of the ancients, whom the Papists have followed, as if Christ taught that, by keeping the law, we may merit eternal life.” Calvin is answering a claim to merit. The reading written down here claims less, the way made possible by grace, and he denies that too, since for him no sinner keeps the law.

{\small\begin{verbatim}
def calvinAgainstTheCommandmentsAsTheWay : Source :=
  { primary := Reference.work calvinHarmony2 (Bib.Locus.page 394),
    tradition := Tradition.reformedProtestant,
    confidence := Confidence.disputed }
\end{verbatim}}
\par\medskip\noindent\textbf{\texttt{veritatis\allowbreak{}Splendor\allowbreak{}Against\allowbreak{}The\allowbreak{}Legal\allowbreak{}Reply}}\par\nobreak\noindent
\textbf{John Paul II contests the Reformed step from Matthew 19:17.} He accepts the text and reads it as the real path: Jesus “proposes them to us as the way and condition of salvation” (\emph{Veritatis Splendor} 12).

{\small\begin{verbatim}
def veritatisSplendorAgainstTheLegalReply : Source :=
  { primary := Reference.work veritatisSplendor (Bib.Locus.sectionRef "12"),
    tradition := Tradition.romanCatholic, confidence := Confidence.disputed }
\end{verbatim}}
\par\medskip\noindent\textbf{\texttt{calvin\allowbreak{}Against\allowbreak{}Baptism\allowbreak{}In\allowbreak{}John\allowbreak{}Three\allowbreak{}Five}}\par\nobreak\noindent
\textbf{Calvin contests Trent’s step from John 3:5.} He accepts the verse and denies that it requires water baptism: “I cannot bring myself to believe that Christ speaks of baptism” (on John 3:5, 110).

{\small\begin{verbatim}
def calvinAgainstBaptismInJohnThreeFive : Source :=
  { primary := Reference.work calvinJohn (Bib.Locus.page 110),
    tradition := Tradition.reformedProtestant,
    confidence := Confidence.disputed }
\end{verbatim}}
\par\medskip\noindent\textbf{\texttt{readings\allowbreak{}Without\allowbreak{}The\allowbreak{}Water}}\par\nobreak\noindent
\textbf{The steps from each non-baptismal reading of John 3:5 are well supported.} If the water is natural birth, or Ezekiel’s cleansing and not baptism, or the Spirit’s own cleansing, the verse does not require water baptism. No reader was found who holds one of these readings and still requires baptism from the verse. Those who require it deny the reading (Trent, Session VII, On Baptism, canon 2).

{\small\begin{verbatim}
def readingsWithoutTheWater : Source :=
  { primary := Reference.work oliverWaterJohn,
    supporting :=
      [Reference.work Bib.calvinInstitutes (Bib.Locus.sectionRef "IV.xvi.25")],
    tradition := Tradition.christianHistoricalGrammatical,
    confidence := Confidence.wellSupported }
\end{verbatim}}
\par\medskip\noindent\textbf{\texttt{augustine\allowbreak{}On\allowbreak{}The\allowbreak{}Thief}}\par\nobreak\noindent
\textbf{The step from the thief is well supported.} Grant that he was promised paradise unbaptised, and that his case bears on salvation under the Gospel, and the washing itself is not necessary for salvation: “as in the thief, to whom the material administration of the sacrament was necessarily wanting, the salvation was complete” (Augustine, \emph{On Baptism} IV.23.31). Aquinas quotes the passage with approval (III q.66 a.11). Feeney denies the conclusion. How he answers the thief could not be verified.

{\small\begin{verbatim}
def augustineOnTheThief : Source :=
  {
    primary :=
      Reference.work npnfAugustineManichaeans
        (Bib.Locus.sectionRef "On Baptism IV.23.31"),
    supporting :=
      [Reference.work aquinasSumma (Bib.Locus.sectionRef "III q.66 a.11")],
    tradition := Tradition.christianHistoricalGrammatical,
    confidence := Confidence.wellSupported }
\end{verbatim}}
\par\medskip\noindent\textbf{\texttt{catechism\allowbreak{}On\allowbreak{}Desire}}\par\nobreak\noindent
\textbf{The step from the desire to the salvation it assures is well supported.} If the desire for baptism brings about its fruits, the washing itself is not necessary for salvation: the Catechism says it “assures them the salvation that they were not able to receive through the sacrament” (1259). Feeney accepts that the desire justifies and denies that it saves. So what he denies is the starting point, this saving fruit as the Catechism states it, and not the step.

{\small\begin{verbatim}
def catechismOnDesire : Source :=
  {
    primary :=
      Reference.work catechismCatholicChurch (Bib.Locus.sectionRef "1259"),
    tradition := Tradition.romanCatholic,
    confidence := Confidence.wellSupported }
\end{verbatim}}
\par\medskip\noindent\textbf{\texttt{westminster\allowbreak{}On\allowbreak{}The\allowbreak{}Desire}}\par\nobreak\noindent
\textbf{The Reformed step from the desire, as Westminster holds it.} Grace and salvation “are not so inseparably annexed unto it, as that no person can be regenerated, or saved, without it” (XXVIII.5), and faith is “the alone instrument of justification” (XI.2).

{\small\begin{verbatim}
def westminsterOnTheDesire : Source :=
  {
    primary :=
      Reference.work westminsterConfession (Bib.Locus.sectionRef "XXVIII.5"),
    supporting :=
      [Reference.work westminsterConfession (Bib.Locus.sectionRef "XI.2")],
    tradition := Tradition.reformedProtestant,
    confidence := Confidence.disputed }
\end{verbatim}}
\par\medskip\noindent\textbf{\texttt{aquinas\allowbreak{}On\allowbreak{}The\allowbreak{}Desire}}\par\nobreak\noindent
\textbf{Aquinas contests the Reformed step from the desire.} He accepts that Cornelius and others “receive grace and virtues through their faith in Christ and their desire for Baptism”. He denies that the sacrament is therefore not the instrument: the desire is for the sacrament, and “when baptized, they receive a yet greater fulness of grace and virtues” (III q.69 a.4 ad 2).

{\small\begin{verbatim}
def aquinasOnTheDesire : Source :=
  {
    primary :=
      Reference.work aquinasSumma (Bib.Locus.sectionRef "III q.69 a.4 ad 2"),
    tradition := Tradition.romanCatholic, confidence := Confidence.disputed }
\end{verbatim}}
\par\medskip\noindent\textbf{\texttt{calvin\allowbreak{}On\allowbreak{}Acts\allowbreak{}Ten}}\par\nobreak\noindent
\textbf{Where the Spirit comes before the water, and after it.} Both sides accept that God is not bound to the water; the Catechism says God “himself is not bound by his sacraments” (1257). Calvin’s own sentence on Acts 10:47 goes further, “the Spirit is not included in baptism”, and that further claim is not what is accepted here.

{\small\begin{verbatim}
def calvinOnActsTen : Source :=
  { primary := Reference.work calvinActs1 (Bib.Locus.page 454),
    supporting :=
      [Reference.work catechismCatholicChurch (Bib.Locus.sectionRef "1257")],
    tradition := Tradition.christianHistoricalGrammatical,
    confidence := Confidence.consensus }
\end{verbatim}}
\par\medskip\noindent\textbf{\texttt{melanchthon\allowbreak{}On\allowbreak{}Luke7}}\par\nobreak\noindent
Where Melanchthon reads Luke 7:47 in the light of 7:50, in the \emph{Apology of the Augsburg Confession} (1531): the part of Article IV on love and the fulfilling of the law. The \emph{Concordia Triglotta} numbers it Article III, 31–34. The reference is given by article and topic because paragraph numbers differ between editions.

{\small\begin{verbatim}
def melanchthonOnLuke7 : Bib.Locus :=
  Bib.Locus.sectionRef
    "Apology IV, on love and the fulfilling of the law: Luke 7:47–50"
\end{verbatim}}
\par\medskip\noindent\textbf{\texttt{campbell\allowbreak{}On\allowbreak{}Deliverance}}\par\nobreak\noindent
The step in the apocalyptic reading: from Christ’s own faithfulness and God’s act of rescue to the denial that faith is the condition. Rated \texttt{plausible}: no source cited here accepts both starting points and still holds that faith is the condition of justification.

{\small\begin{verbatim}
def campbellOnDeliverance : Source :=
  { primary := Reference.work campbellDeliveranceOfGod,
    supporting := [Reference.work martynGalatians],
    tradition := Tradition.criticalScholarship,
    confidence := Confidence.plausible }
\end{verbatim}}
\par\medskip\noindent\textbf{\texttt{reformed\allowbreak{}Works\allowbreak{}Of\allowbreak{}Law\allowbreak{}Source}}\par\nobreak\noindent
The Reformed reading of “works of the law” (ἔργα νόμου), cited to the tradition that holds it and to those who defend it today against the New Perspective. It is \texttt{disputed} because the New Perspective denies it; see \texttt{new\allowbreak{}Perspective\allowbreak{}Cite}.

{\small\begin{verbatim}
def reformedWorksOfLawSource : Source
\end{verbatim}}
\par\medskip\noindent\textbf{\texttt{base\allowbreak{}Cite}}\par\nobreak\noindent
The sources shared by every position: for the verses, and for the premises about how to read them, whose wording none of the positions disputes.

{\small\begin{verbatim}
def baseCite : Claim → AtomMeta
\end{verbatim}}
\par\medskip\noindent\textbf{\texttt{reformed\allowbreak{}Cite}}\par\nobreak\noindent
The Reformed sources for the disputed readings.

{\small\begin{verbatim}
def reformedCite : Claim → AtomMeta :=
  fun x =>
    match x with
    | Claim.worksOfLawMeansWorksGenerally =>
      {
        label :=
          "Paul's ‘works of the law’ means human works in general, not boundary markers",
        kind := PremiseKind.linguistic,
        source := reformedWorksOfLawSource }
    | c => baseCite c
\end{verbatim}}
\par\medskip\noindent\textbf{\texttt{new\allowbreak{}Perspective\allowbreak{}Cite}}\par\nobreak\noindent
The New Perspective’s source for its reading: “works of the law” (ἔργα νόμου) means the practices that marked Israel off, so Paul is not speaking of human effort in general.

{\small\begin{verbatim}
def newPerspectiveCite : Claim → AtomMeta :=
  fun x =>
    match x with
    | Claim.worksOfLawMeansWorksGenerally =>
      {
        label :=
          "Paul's ‘works of the law’ means the Jewish boundary markers, not works in general",
        kind := PremiseKind.linguistic,
        source :=
          { primary := Reference.work dunnNewPerspective,
            supporting :=
              [Reference.work sandersPaulPalestinianJudaism,
                Reference.work wrightWhatPaulSaid,
                Reference.work wrightJustification],
            tradition := Tradition.criticalScholarship,
            confidence := Confidence.disputed } }
    | c => baseCite c
\end{verbatim}}
\par\medskip\noindent\textbf{\texttt{subjective\allowbreak{}Genitive\allowbreak{}Cite}}\par\nobreak\noindent
Hays’s reading: “faith of Christ” (πίστις Χριστοῦ) is Christ’s own faithfulness. Everything else as in \texttt{base\allowbreak{}Cite}.

{\small\begin{verbatim}
def subjectiveGenitiveCite : Claim → AtomMeta
\end{verbatim}}
\par\medskip\noindent\textbf{\texttt{law\allowbreak{}Observant\allowbreak{}Luke\allowbreak{}Cite}}\par\nobreak\noindent
Jervell’s reading of Luke and Acts: what Acts 15 does not lay on gentiles is Israel’s law as a mark of belonging, and the law is not refused as a means of salvation. Everything else as in \texttt{base\allowbreak{}Cite}.

{\small\begin{verbatim}
def lawObservantLukeCite : Claim → AtomMeta
\end{verbatim}}
\par\medskip\noindent\textbf{\texttt{critical\allowbreak{}Authorship\allowbreak{}Cite}}\par\nobreak\noindent
The critical view of authorship: Ephesians and Titus are not by Paul, and 1 and 2 Peter are not by Peter. Each claim is cited to a scholar who holds that view of that letter, with Ehrman for the case across all four. Marshall’s middle view of the Pastoral letters, that they were written by Paul’s circle after his death without intent to deceive, is cited with the critical side, because it too denies that Paul wrote Titus.

{\small\begin{verbatim}
def criticalAuthorshipCite : Claim → AtomMeta
\end{verbatim}}
\par\medskip\noindent\textbf{\texttt{critics\allowbreak{}Cite}}\par\nobreak\noindent
The critics of the New Perspective: Judaism in Paul’s day was not uniformly “in by grace, staying in by obedience”, because final vindication according to works was widely held. Gathercole is the main source, with the volume edited by Carson, O’Brien and Seifrid in support.

{\small\begin{verbatim}
def criticsCite : Claim → AtomMeta :=
  fun x =>
    match x with
    | Claim.secondTempleCovenantalNomism =>
      {
        label :=
          "Second Temple Judaism also held final vindication according to works",
        kind := PremiseKind.historical,
        source :=
          { primary := Reference.work gathercoleWhereIsBoasting,
            supporting := [Reference.work carsonVariegatedNomism1],
            tradition := Tradition.reformedProtestant,
            confidence := Confidence.disputed } }
    | c => reformedCite c
\end{verbatim}}
\par\medskip\noindent\textbf{\texttt{tridentine\allowbreak{}Cite}}\par\nobreak\noindent
Trent’s account of what justification is: “not remission of sins merely, but also the sanctification and renewal of the inward man”. Everything else as in \texttt{base\allowbreak{}Cite}.

{\small\begin{verbatim}
def tridentineCite : Claim → AtomMeta :=
  fun x =>
    match x with
    | Claim.justificationIsForensicOnly =>
      {
        label :=
          "Justification is not a verdict only: it renews the inward man",
        kind := PremiseKind.theological,
        source :=
          {
            primary :=
              Reference.work tannerDecrees
                (Bib.Locus.sectionRef
                  "Trent, Session VI (1547), Decree on Justification, ch. 7"),
            tradition := Tradition.romanCatholic,
            confidence := Confidence.wellSupported } }
    | c => baseCite c
\end{verbatim}}
\par\medskip\noindent\textbf{\texttt{finnish\allowbreak{}Cite}}\par\nobreak\noindent
The Finnish reading of Luther: justification is not a verdict only, because Christ himself is present in faith and is the believer’s righteousness. Everything else as in \texttt{reformed\allowbreak{}Cite}.

{\small\begin{verbatim}
def finnishCite : Claim → AtomMeta :=
  fun x =>
    match x with
    | Claim.justificationIsForensicOnly =>
      {
        label :=
          "Justification is not a verdict only: Christ is present in faith",
        kind := PremiseKind.theological,
        source :=
          { primary := Reference.work mannermaaChristPresentInFaith,
            tradition := Tradition.criticalScholarship,
            confidence := Confidence.disputed } }
    | c => reformedCite c
\end{verbatim}}
\par\medskip\noindent\textbf{\texttt{apocalyptic\allowbreak{}Cite}}\par\nobreak\noindent
The apocalyptic reading: “faith of Christ” is Christ’s own faithfulness. The source recorded is Campbell’s reason for reading it so.

{\small\begin{verbatim}
def apocalypticCite : Claim → AtomMeta
\end{verbatim}}

\bibliographystyle{plain}
\bibliography{references}
\end{document}
